% !TEX program = xelatex
\documentclass[11pt,openany,fontset=fandol]{ctexbook}
\usepackage[paperwidth=176mm,paperheight=250mm,inner=21mm,outer=18mm,top=20mm,bottom=20mm,headheight=14pt,headsep=6mm,footskip=10mm]{geometry}
\usepackage{amsmath,amssymb,amsthm,bm,mathtools,booktabs,array,enumitem,float}
\usepackage{tikz,fancyhdr,hyperref}
\usetikzlibrary{arrows.meta,calc,patterns,angles,quotes,decorations.markings}
\hypersetup{hidelinks,pdftitle={运动和静力进阶：五讲合订本},pdfauthor={}}
\setlength{\parindent}{2em}\setlength{\parskip}{3pt}
\setlength{\emergencystretch}{2em}\linespread{1.18}\raggedbottom
\setcounter{tocdepth}{1}
\ctexset{chapter={name={第,讲},number=\chinese{chapter},format=\LARGE\bfseries\centering,beforeskip=16pt,afterskip=22pt},section={format=\Large\bfseries},subsection={format=\large\bfseries}}
\pagestyle{fancy}\fancyhf{}\fancyhead[LE,RO]{\small\thepage}
\fancyhead[LO]{\small 运动和静力进阶}\fancyhead[RE]{\small\nouppercase{\leftmark}}
\renewcommand{\headrulewidth}{0pt}
\fancypagestyle{plain}{\fancyhf{}\fancyfoot[C]{\small\thepage}\renewcommand{\headrulewidth}{0pt}}
\theoremstyle{definition}\newtheorem{example}{例}[chapter]\newtheorem{exercise}{习题}[chapter]
\newcommand{\LectureNeedspace}[1]{\par\begingroup\dimen0=#1\relax\ifdim\pagegoal<\maxdimen\ifdim\dimexpr\pagegoal-\pagetotal\relax<\dimen0\newpage\fi\fi\endgroup}
\newcommand{\SolutionStart}{\par\medskip\ifdim\pagetotal>0pt\ifdim\dimexpr\pagegoal-\pagetotal\relax<3\baselineskip\newpage\fi\fi\noindent\textbf{解}\quad\ignorespaces}
\newenvironment{solution}{\SolutionStart}{\par\medskip}
\newcommand{\dd}{\mathop{}\!\mathrm d}\newcommand{\ee}{\mathrm e}\newcommand{\vct}[1]{\bm{#1}}
\newcommand{\Prob}{\mathbb P}\newcommand{\avg}[1]{\left\langle#1\right\rangle}
\tikzset{axis/.style={->,>=Stealth,line width=.4pt},curve/.style={line width=.7pt},guide/.style={dashed,line width=.35pt},every node/.style={font=\small}}



\tikzset{flow/.style={curve,postaction={decorate},decoration={markings,mark=at position .22 with {\arrow{Stealth}},mark=at position .68 with {\arrow{Stealth}}}}}



% 分册与十编合订本共用的定位页码；局部章、图、题号保持原有对应。
\newcommand{\CoursePartNumber}{2}
\geometry{inner=19.5mm,outer=19.5mm,headheight=16pt}
\AtBeginDocument{%
  \fancyhead{}\fancyhead[C]{\small 运动和静力进阶}%
  \fancyfoot[C]{\small\CoursePartNumber-\thepage}%
  \fancypagestyle{plain}{\fancyhf{}\fancyfoot[C]{\small\CoursePartNumber-\thepage}\renewcommand{\headrulewidth}{0pt}}%
}

\begin{document}\frontmatter
\begin{titlepage}
\centering\vspace*{30mm}
{\Large 物理竞赛二轮讲义\quad 第二编\par}\vspace{8mm}

{\Huge\bfseries 运动和静力进阶\par}
\vspace{12mm}{\Large 物理竞赛二轮讲义\par}
\vspace{10mm}{\large 第一讲至第五讲\par}
\vfill{\large 2026\par}
\end{titlepage}
\begingroup\small\setlength{\parskip}{0pt}\linespread{1}\selectfont\tableofcontents\endgroup\mainmatter
\setcounter{chapter}{0}


\chapter{微元法进阶}

细杆的质量由各小段相加，落点概率由相应的方向微元积累，束流推力则来自每次碰撞的动量转移。列积分前，需要确定微元覆盖的对象和密度的定义；换元时还须保留全部分支。本讲沿这些关系讨论连续分布、空间方向和碰撞散射。


\section{微元选取与变量变换}
坐标改变后，同一批物质或事件对应的区间宽度也会改变。密度变换由微元的物理含义和几何关系确定，支持区域、分支数与总量须彼此一致。

\subsection{微元表示的量与积分区域}
对细杆、薄片和体分布，质量微元分别写为
\[
\dd m=\lambda\dd s,\qquad \dd m=\sigma\dd S,\qquad \dd m=\rho\dd V.
\]
这里 $\lambda,\sigma,\rho$ 的单位依次为 $\mathrm{kg/m}$、$\mathrm{kg/m^2}$、$\mathrm{kg/m^3}$。密度的定义已包含微元的种类。若把弧长 $\dd s$ 换成水平长度 $\dd x$，新的线密度一般也随之改变。

设一条平面曲线写为 $y=y(x)$。图\ref{mp1:fig:micro}左图中的直角微元给出
\[
\dd s=\sqrt{1+y'^2}\,|\dd x|.
\]
因而对沿曲线均匀的细丝，$\dd m=\lambda\sqrt{1+y'^2}|\dd x|$。它沿水平投影的分布通常不均匀。若曲线用参数 $u$ 表示，则
\[
\dd s=\left|\frac{\dd\vct r}{\dd u}\right||\dd u|.
\]
积分的上下限应覆盖实际存在的曲线段。参数在往返过程中重复经过同一点时，还要区分所求的是几何长度还是沿运动历程累计的路程。

\begin{figure}[H]\centering

\begin{tikzpicture}[scale=.9]
\begin{scope}
\draw[axis](0,0)--(2.8,0)node[right]{$x$};\draw[axis](0,0)--(0,2.1)node[above]{$y$};
\draw[curve,domain=.2:2.3]plot(\x,{.25*(\x)^2});\fill(1.2,.36)circle(1.3pt)node[above left]{$P$};
\draw[curve,domain=1.2:1.8]plot(\x,{.25*(\x)^2});\draw[guide](1.2,.36)--(1.8,.36)--(1.8,.81);
\node[below,fill=white,inner sep=1pt]at(1.5,.36){$\dd x$};\node[right]at(1.8,.58){$\dd y$};\node[above left]at(1.58,.65){$\dd s$};
\end{scope}
\begin{scope}[shift={(4.5,.3)}]
\fill[black!8](0,0)--(2,.3)--(2.7,1.5)--(.7,1.2)--cycle;
\draw[axis,thick](0,0)--(2,.3)node[below right]{$\vct r_u\dd u$};\draw[axis,thick](0,0)--(.7,1.2)node[left]{$\vct r_v\dd v$};
\draw[guide](2,.3)--(2.7,1.5)--(.7,1.2);\draw[axis](1.3,.75)--(.9,2.1)node[above]{$\vct n$};\node at(1.6,.8){$\dd S$};
\end{scope}
\begin{scope}[shift={(9,.3)}]
\draw(0,0)rectangle(1.5,1);\draw(0,1)--(.5,1.55)--(2,1.55)--(2,.55)--(1.5,0);\draw(1.5,1)--(2,1.55);
\draw[guide](0,0)--(.5,.55)--(2,.55) (.5,.55)--(.5,1.55);\node at(.8,.55){$\dd V$};
\end{scope}
\end{tikzpicture}
\caption{参数变化对应的线元、面积元与体积元。曲面微元由两条切向边张成，体积微元还包含第三个方向。}\label{mp1:fig:micro}
\end{figure}

对参数曲面 $\vct r(u,v)$，图\ref{mp1:fig:micro}中两条切向边为 $\vct r_u\dd u$ 与 $\vct r_v\dd v$。平行四边形面积为
\[
\dd S=|\vct r_u\times\vct r_v|\,|\dd u\dd v|.
\]
若只需要曲面在某个坐标平面的投影面积，就应取相应的法向分量。真实面积与投影面积是不同的积分元。

空间参数变换 $\vct r=\vct r(u,v,w)$ 的体积微元由混合积给出：
\[
\dd V=\left|\vct r_u\cdot(\vct r_v\times\vct r_w)\right||\dd u\dd v\dd w|.
\]
它等于雅可比行列式的绝对值。绝对值使质量、面积和概率等标量微元为正；在做功或通量积分中，方向则由位移矢量或有向面积元另外记录。

平面极坐标的微小两边长为 $\dd r,r\dd\phi$，故 $\dd S=r\dd r\dd\phi$。把整圈合并成环带，就得到 $\dd S=2\pi r\dd r$。选用环带的前提是被积量在同一圆周上相同；若密度依赖方位角，仍须保留角积分。

\subsection{分布密度随坐标改变}
设随机变量 $X$ 的密度为 $f_X(x)$，其含义是
\[
\Prob\{x<X<x+\dd x\}=f_X(x)\dd x,
\qquad \int_{\mathcal D_X}f_X(x)\dd x=1.
\]
密度有单位，数值可以大于1；一个小区间的概率由密度乘区间宽度得到。若 $Y=g(X)$ 在所讨论区间内单调，$y=g(x)$，同一批对象的概率满足
\[
f_Y(y)|\dd y|=f_X(x)|\dd x|,
\]
于是
\begin{equation}\label{mp1:eq:transform-one}
f_Y(y)=f_X(x(y))\left|\frac{\dd x}{\dd y}\right|.
\end{equation}
原变量均匀，经过非线性变换后通常不再均匀。变换后的允许区域也要由原区域映射得到，不能只改被积式。

当同一个 $y$ 对应多个原像 $x_i$ 时，各原像附近的对象都进入同一段 $\dd y$，应将概率相加：
\begin{equation}\label{mp1:eq:transform-many}
f_Y(y)=\sum_{i}\frac{f_X(x_i(y))}{|g'(x_i(y))|}.
\end{equation}
此式在各分支导数非零的内部使用。导数为零的点附近，密度可能发散；是否仍有有限概率，要看相应的反常积分。


\subsection{驻点映射与可积的密度峰}
多分支公式在$g'(x)=0$处不能逐点直接使用，但可先在驻点两侧换元，再取极限。设原密度在$x_c$附近连续且非零，映射$g$在该邻域三阶连续可导，且$g'(x_c)=0$、$g''(x_c)=A\ne0$，记$y_c=g(x_c)$。Taylor展开给
\[y-y_c=\frac A2(x-x_c)^2+O((x-x_c)^3).
\]
在有原像的一侧，令$\delta=|y-y_c|$，两个局部分支分别为
\[x_\pm-x_c=\pm\sqrt{\frac{2\delta}{|A|}}+O(\delta),
\qquad |g'(x_\pm)|=\sqrt{2|A|\delta}+O(\delta).
\]
若两支都属于原来的允许区域，则应先把两支概率相加：
\[f_Y(y)=\frac{f_X(x_+)}{|g'(x_+)|}
+\frac{f_X(x_-)}{|g'(x_-)|}
\sim\frac{\sqrt2 f_X(x_c)}{\sqrt{|A|\delta}}.
\]
图\ref{mp1:fig:series-density-fold}左图的同一水平线截出两个原像，它们共同产生右图的密度峰。若物理区域只含驻点的一侧，则只保留一个分支，首项系数减半；指数仍为$-1/2$。

\begin{figure}[H]\centering
\begin{tikzpicture}[>=Stealth,line cap=round,line join=round]
\begin{scope}[xshift=-2.5cm]
\draw[axis](-1.8,0)--(1.9,0)node[right]{$x-x_c$};\draw[axis](0,-.15)--(0,1.8)node[above]{$y-y_c$};
\draw[curve,domain=-1.5:1.5]plot(\x,{.6*(\x)^2});\draw[guide](-1.5,.6)--(1.5,.6);
\fill(-1,.6)circle(1.5pt)node[above left]{$x_-$};\fill(1,.6)circle(1.5pt)node[above right]{$x_+$};\end{scope}
\begin{scope}[xshift=2.5cm]
\draw[axis](0,0)--(2.8,0)node[right]{$\delta$};\draw[axis](0,0)--(0,2.8)node[above]{$f_Y$};
\draw[curve,domain=.09:2.5,samples=101]plot(\x,{.8/sqrt(\x)});\node[right]at(1,.8){$f_Y\propto\delta^{-1/2}$};\end{scope}
\end{tikzpicture}
\caption{左图取$A>0$，同一像区间对应驻点两侧的原像。右图的密度峰可积。}\label{mp1:fig:series-density-fold}
\end{figure}

密度无界不等于概率无界。在宽度$\varepsilon$的像区间内，
\[\int_0^\varepsilon\frac{\sqrt2f_X(x_c)}{\sqrt{|A|\delta}}\dd\delta
=\frac{\sqrt2f_X(x_c)}{\sqrt{|A|}}[2\sqrt\delta]_0^\varepsilon
=2f_X(x_c)\sqrt{\frac{2\varepsilon}{|A|}},
\]
仍随$\varepsilon\to0$趋零。这个积分等于原变量落在$x_c\pm\sqrt{2\varepsilon/|A|}$之间的首阶概率，说明$y_c$处的密度峰不含有限的点质量。只有当一整段非零概率的原区域被映到同一个值时，才须另加离散概率项。

方向分布、落点密度和散射角分布中的密度峰常由这种映射驻点产生。列式时应先确定原分布、分支数及允许区域，再讨论峰的形状。把发散处删掉会改变总概率；把两条原像合成一条，也会把密度系数算错。

\par\medskip\noindent\begin{minipage}{\textwidth}
\begin{example}[平方映射的两个原像]\label{mp1:ex:e1}
随机变量 $X$ 在 $[-a,a]$ 上均匀分布，$a>0$，令 $Y=X^2/a$。求 $Y$ 的密度、$\Prob(Y<a/4)$，并说明 $Y=0$ 附近的密度是否有意义。图\ref{mp1:fig:e1-q}给出两支原像。
\end{example}

\begin{figure}[H]\centering

\begin{tikzpicture}[x=2.3cm,y=2.1cm]
\draw[axis](-1.2,0)--(1.3,0)node[right]{$X/a$};\draw[axis](0,0)--(0,1.25)node[above]{$Y/a$};
\draw[curve,domain=-1:1,samples=80]plot(\x,{(\x)^2});\draw[guide](-1,1)--(1,1);\draw[guide](-.6,.36)--(.6,.36);
\fill(-.6,.36)circle(.7pt)node[left]{$X_-$};\fill(.6,.36)circle(.7pt)node[right]{$X_+$};
\node[below]at(-1,0){$-1$};\node[below]at(1,0){$1$};\node[left]at(0,1){$1$};
\end{tikzpicture}
\caption{例~\ref{mp1:ex:e1}题设：平方映射把 $X$ 的正、负两支送到同一个 $Y$。图中坐标均按 $a$ 归一化。}\label{mp1:fig:e1-q}
\end{figure}

\end{minipage}\par\medskip

\begin{solution}
原密度为 $f_X=1/(2a)$，变换后的区域为 $0\le y\le a$。对内部点 $0<y<a$，原像有两支：
\[
x_+=\sqrt{ay},\qquad x_-=-\sqrt{ay},\qquad
\left|\frac{\dd x_\pm}{\dd y}\right|=\frac{\sqrt a}{2\sqrt y}.
\]
图\ref{mp1:fig:e1-s}左图中，两段原像区间各贡献一半概率，故
\[
f_Y(y)=\frac1{2a}\frac{\sqrt a}{2\sqrt y}
+\frac1{2a}\frac{\sqrt a}{2\sqrt y}
=\frac1{2\sqrt{ay}},\qquad 0<y<a.
\]
区域外密度为零。只保留 $x_+$ 会使总概率少一半。

\begin{figure}[H]\centering

\begin{tikzpicture}[scale=.95]
\draw[axis](-2.1,0)--(2.2,0)node[right]{$X$};\draw[thick](-1.5,0)--(-1.1,0) (1.1,0)--(1.5,0);
\node[below]at(-1.3,0){$|\dd x_-|$};\node[below]at(1.3,0){$\dd x_+$};
\draw[axis](-1.3,.15)--(-.3,1.3);\draw[axis](1.3,.15)--(.3,1.3);\draw[thick](-.35,1.55)--(.35,1.55);\node[above]at(0,1.55){$\dd y$};
\begin{scope}[shift={(5,0)},x=3cm,y=1cm]
\draw[axis](0,0)--(1.15,0)node[right]{$y/a$};\draw[axis](0,0)--(0,2.6)node[above]{$a f_Y$};
\draw[curve,<-,>=Stealth,domain=.045:1,samples=100]plot(\x,{1/(2*sqrt(\x))});\draw[guide](1,0)--(1,.5);\node[below]at(1,0){$1$};\node[left]at(0,.5){$1/2$};
\end{scope}
\end{tikzpicture}
\caption{例~\ref{mp1:ex:e1}解图：同一个 $\dd y$ 接收两个原像区间。右图为 $a f_Y$；原点附近的密度发散，但面积积分有限。}\label{mp1:fig:e1-s}
\end{figure}

由密度积分，
\[
\int_0^B\frac{\dd y}{2\sqrt{ay}}
=\frac1{\sqrt a}\left[\sqrt y\right]_0^B
=\sqrt{\frac Ba},\qquad 0\le B\le a.
\]
取 $B=a$，总概率为1；取 $B=a/4$，得到 $\Prob(Y<a/4)=1/2$，也等于 $\Prob(|X|<a/2)$。图\ref{mp1:fig:e1-s}右图中的发散来自原点附近平方映射的压缩：一个很短的 $Y$ 区间接收较长的两个 $X$ 区间。对$0<\varepsilon\le a$，$\Prob(0<Y<\varepsilon)=\sqrt{\varepsilon/a}$，随$\varepsilon\to0$趋零；若$\varepsilon>a$，该概率为1。点$Y=0$没有非零的集中概率。
\end{solution}

\subsection{质量矩与分布合成}
质量微元与概率微元的变换方式相同。区别在于总质量不必为1。对质量分布，质心由一阶质量矩定义：
\[
M=\int\dd m,\qquad \vct r_C=\frac1M\int\vct r\dd m.
\]
均匀重力场中，总重力作用线通过质心。求支座反力时，可以把总重力合并，但计算质心仍须从真实分布出发。

如果只关心某一坐标，就将其余坐标积分掉。例如面密度为 $\sigma(x,y)$ 的薄片，其水平投影线密度为
\[
\lambda_x(x)=\int_{\mathcal D_y(x)}\sigma(x,y)\dd y.
\]
上下边界依赖 $x$ 时，积分限也随 $x$ 改变。这一过程称为边缘化；它保留了每个竖直条带中的总量。

\par\medskip\noindent\begin{minipage}{\textwidth}
\begin{example}[变密度细杆的质心与支承]\label{mp1:ex:e2}
图\ref{mp1:fig:e2-q}中的水平细杆长 $L$，线密度为 $\lambda(x)=\lambda_0(1+x/L)$，$0\le x\le L$。重力场均匀，杆在两端简支。求总质量、质心位置和两支座反力。
\end{example}

\begin{figure}[H]\centering

\begin{tikzpicture}[scale=1.1]
\draw[very thick](0,0)--(6,0);\draw(0,-.04)--(-.25,-.45)--(.25,-.45)--cycle;\draw(6,-.04)--(5.75,-.45)--(6.25,-.45)--cycle;
\draw[axis](-.3,-.85)--(6.5,-.85)node[right]{$x$};\node[below]at(0,-.85){$0$};\node[below]at(6,-.85){$L$};
\foreach\x in{0,.6,...,6}{\draw(\x,.08)--(\x,{.08+.4*(1+\x/6)});}\node[above]at(3,.95){$\lambda(x)=\lambda_0(1+x/L)$};\node[left]at(0,0){$A$};\node[right]at(6,0){$B$};
\end{tikzpicture}
\caption{例~\ref{mp1:ex:e2}题设：细杆两端简支，线密度沿 $x$ 增加。支座只提供竖直反力。}\label{mp1:fig:e2-q}
\end{figure}

\end{minipage}\par\medskip

\begin{solution}
坐标 $x$ 沿直杆，故弧长微元就是 $\dd x$。图\ref{mp1:fig:e2-s}左图所取微元的质量为 $\lambda_0(1+x/L)\dd x$。于是
\[
M=\lambda_0\int_0^L\left(1+\frac xL\right)\dd x
=\lambda_0\left[x+\frac{x^2}{2L}\right]_0^L
=\frac32\lambda_0L.
\]
对左端 $A$ 的质量矩为
\[
Mx_C=\lambda_0\int_0^Lx\left(1+\frac xL\right)\dd x
=\lambda_0\left[\frac{x^2}{2}+\frac{x^3}{3L}\right]_0^L
=\lambda_0\left(\frac{L^2}{2}+\frac{L^2}{3}\right)
=\frac56\lambda_0L^2.
\]
所以 $x_C=5L/9$，比几何中点略偏向线密度较大的 $B$ 端。

\begin{figure}[H]\centering

\begin{tikzpicture}[scale=.9]
\draw[thick](0,0)--(4,0);\draw[line width=4pt](1.5,0)--(1.8,0);\node[above]at(1.65,.2){$\dd m=\lambda(x)\dd x$};\draw[axis](1.65,-.1)--(1.65,-1)node[below]{$g\dd m$};\node[below left]at(1.4,0){$x$};
\begin{scope}[shift={(6,0)}]
\draw[very thick](0,0)--(4,0);\draw[axis](0,0)--(0,.8)node[above]{$R_A$};\draw[axis](4,0)--(4,1)node[above]{$R_B$};\draw[axis,thick](2.222,0)--(2.222,-1.8)node[below]{$Mg$};\node[above]at(2.222,.2){$C$};\draw[<->](0,-2.4)--(2.222,-2.4)node[midway,fill=white]{$5L/9$};
\end{scope}
\end{tikzpicture}
\caption{例~\ref{mp1:ex:e2}解图：左图从坐标 $x$ 处取质量微元；右图把总重力作用线放在 $x_C=5L/9$，据此求两支座反力。}\label{mp1:fig:e2-s}
\end{figure}

图\ref{mp1:fig:e2-s}右图把分布重力合并为 $Mg$。竖直合力与对 $A$ 的合力矩分别给出
\[
R_A+R_B=Mg,\qquad R_BL=Mgx_C.
\]
代入 $x_C$，得到
\[
R_A=\frac49Mg=\frac23\lambda_0gL,\qquad
R_B=\frac59Mg=\frac56\lambda_0gL.
\]
两个反力均向上，满足题设的简支接触；右端反力较大。
还可逐个微元求支座反力。对位置$x$处的重量$g\dd m$，两端分担量满足$\dd R_A+\dd R_B=g\dd m$及$L\dd R_B=xg\dd m$，故
\[
 R_A=\int_0^L\left(1-\frac xL\right)g\dd m,\qquad
 R_B=\int_0^L\frac xL g\dd m.
\]
图\ref{mp1:fig:e2-s}中先合并重力再列矩方程，得到同样的结果。两个线性权函数表示同一微元的重量在两支座间的分担比例。在$x=L/2$处，两端对该微元的分担相等；全杆的右端反力较大，是因为线密度向右增加，使右侧微元的总重量较大。
\end{solution}

\subsection{曲面分布与投影分布}
轴对称曲面可以按母线上的窄条旋转分割。设曲面为 $z=z(r)$，母线弧长为 $\dd s=\sqrt{1+z'^2}\dd r$。窄条旋转一周形成环带，第一阶面积为
\begin{equation}\label{mp1:eq:surface-revolution}
\dd S=2\pi r\sqrt{1+z'^2}\dd r.
\end{equation}
也可将曲面写为 $\vct R(r,\phi)=(r\cos\phi,r\sin\phi,z(r))$。两条参数切向与叉积分别为
\[
\begin{aligned}
\vct R_r&=(\cos\phi,\sin\phi,z'),\qquad
\vct R_\phi=(-r\sin\phi,r\cos\phi,0),\\
\vct R_r\times\vct R_\phi&=(-rz'\cos\phi,-rz'\sin\phi,r),\\
|\vct R_r\times\vct R_\phi|&=r\sqrt{1+z'^2}.
\end{aligned}
\]
对 $0\le\phi<2\pi$ 积分，便得到式\eqref{mp1:eq:surface-revolution}的环带面积。均匀面密度 $\sigma$ 指每单位真实面积的质量；若按水平投影 $2\pi r\dd r$ 积分，就必须把因子 $\sqrt{1+z'^2}$ 一并放入。

\par\medskip\noindent\begin{minipage}{\textwidth}
\begin{example}[抛物面薄壳的质量与重心]\label{mp1:ex:e3}
均匀薄壳由曲面 $z=r^2/(2a)$，$0\le r\le R$ 构成，$a>0$，面密度为 $\sigma$，只含侧面。求其质量和重心高度。图\ref{mp1:fig:e3-q}中 $z=0$ 位于顶点。
\end{example}

\begin{figure}[H]\centering

\begin{tikzpicture}[scale=1.3]
\draw[axis](0,-.15)--(0,2.25)node[above]{$z$};\draw[axis](-1.9,0)--(2.2,0)node[right]{$r$};
\draw[curve,domain=-1.6:1.6,samples=100]plot(\x,{.5*(\x)^2});\draw(-1.6,1.28)arc(180:360:1.6 and .28);\draw[guide](-1.6,1.28)arc(180:0:1.6 and .28);
\draw[guide](0,1.28)--(1.6,1.28)node[midway,above]{$R$};\node[below left]at(0,0){$O$};\node[right]at(1.6,.65){$z=r^2/(2a)$};\node[left]at(-1.6,.7){$\sigma$};
\end{tikzpicture}
\caption{例~\ref{mp1:ex:e3}题设：均匀抛物面薄壳只有侧面，顶点为 $O$，边缘半径为 $R$。虚线表示轴对称的背面轮廓。}\label{mp1:fig:e3-q}
\end{figure}

\end{minipage}\par\medskip

\begin{solution}
图\ref{mp1:fig:e3-s}左图中，$\dd z=(r/a)\dd r$，故 $\dd s=\sqrt{1+r^2/a^2}\dd r$。右图环带的质量为
\[
\dd m=2\pi\sigma r\sqrt{1+\frac{r^2}{a^2}}\dd r.
\]
令 $u=1+r^2/a^2$，记 $U=1+R^2/a^2$，有 $r\dd r=a^2\dd u/2$。因此
\[
M=\pi\sigma a^2\int_1^Uu^{1/2}\dd u
=\pi\sigma a^2\left[\frac23u^{3/2}\right]_1^U
=\frac{2\pi\sigma a^2}{3}(U^{3/2}-1).
\]
轴对称性使重心位于 $z$ 轴。

\begin{figure}[H]\centering

\begin{tikzpicture}[scale=1.05]
\draw[axis](0,0)--(2.5,0)node[right]{$r$};
\draw[axis](0,0)--(0,2.5)node[above]{$z$};
\draw[curve,domain=0:1.9]plot(\x,{.5*(\x)^2});
\draw[line width=1.2pt,domain=1:1.4]plot(\x,{.5*(\x)^2});
\draw[guide](1,.5)--(1.4,.5)--(1.4,.98);
\node[below]at(1.2,.5){$\dd r$};
\node[right]at(1.4,.74){$\dd z$};
\node[above left]at(1.15,.68){$\dd s$};
\begin{scope}[shift={(5,1)}]
\draw(-1.4,0)arc(180:360:1.4 and .35);
\draw[guide](-1.4,0)arc(180:0:1.4 and .35);
\draw(-1.6,.32)arc(180:360:1.6 and .4);
\draw[guide](-1.6,.32)arc(180:0:1.6 and .4);
\draw(-1.4,0)--(-1.6,.32);
\draw[line width=1.2pt](1.4,0)--(1.6,.32)node[midway,right=3pt]{$\dd s$};
\draw[axis](0,-.65)--(0,1.1)node[above]{$z$};
\draw[guide](0,0)--(1.4,0)node[midway,below]{$r$};
\node[below]at(0,-1.1){$\dd S=2\pi r\dd s$};
\end{scope}
\end{tikzpicture}
\caption{例~\ref{mp1:ex:e3}解图：左图的母线弧长 $\dd s$ 大于水平投影 $\dd r$；绕轴一周后得到右图的薄环带。}\label{mp1:fig:e3-s}
\end{figure}

高度在环带内相同，故一阶质量矩为
\begin{align*}
Mz_C&=\int_0^R\frac{r^2}{2a}\,2\pi\sigma r\sqrt{1+\frac{r^2}{a^2}}\dd r\\
&=\frac{\pi\sigma a^3}{2}\int_1^U(u-1)u^{1/2}\dd u\\
&=\frac{\pi\sigma a^3}{2}\left[\frac25u^{5/2}-\frac23u^{3/2}\right]_1^U\\
&=\pi\sigma a^3\left[\frac{U^{5/2}-1}{5}-\frac{U^{3/2}-1}{3}\right].
\end{align*}
除以质量，得到
\[
z_C=a\,\frac{3U^{5/2}-5U^{3/2}+2}{10(U^{3/2}-1)}.
\]
考察浅曲面极限，令 $\eta=R^2/a^2\ll1$，则 $U=1+\eta$。二项式展开为
\[
\begin{aligned}
(1+\eta)^{3/2}&=1+\frac32\eta+\frac38\eta^2+O(\eta^3),\\
(1+\eta)^{5/2}&=1+\frac52\eta+\frac{15}{8}\eta^2+O(\eta^3).
\end{aligned}
\]
重心公式的分子中，常数项与一次项分别相消，留下
\[
3U^{5/2}-5U^{3/2}+2=\frac{15}{4}\eta^2+O(\eta^3),
\qquad 10(U^{3/2}-1)=15\eta+O(\eta^2).
\]
因此
\[
z_C=\frac{a\eta}{4}+O(a\eta^2)\simeq\frac{R^2}{4a},
\qquad M=\frac{2\pi\sigma a^2}{3}\left[\frac32\eta+O(\eta^2)\right]
=\pi\sigma R^2[1+O(\eta)].
\]
此时真实面积接近投影面积，重心高度接近边缘高度的一半。
\end{solution}












\par\medskip\noindent\begin{minipage}{\textwidth}
\begin{example}[单调调分映射与支持区间]\label{mp1:rev:score}
把分数视为连续变量，原分数$X$在$[1,100]$上均匀分布，调整为$Y=10\sqrt X$。图\ref{mp1:fig:rev:score-q}画出映射。求调整后的分数密度与累计分布，并说明若原区间改为$[0,100]$，结果怎样改变。
\end{example}

\begin{figure}[H]\centering
\begin{tikzpicture}[>=Stealth,line cap=round,line join=round]
\draw[axis](0,0)--(4.4,0)node[right]{$X$};\draw[axis](0,0)--(0,3)node[above]{$Y$};\draw[curve,domain=.04:4,samples=80]plot(\x,{1.3*sqrt(\x)});\fill(.04,.26)circle(1.3pt);\fill(4,2.6)circle(1.3pt);\node[right]at(1.5,1.9){$Y=10\sqrt X$};\node[below]at(4,0){$100$};\node[left]at(0,2.6){$100$};\node[below right]at(.04,.26){$(1,10)$};
\end{tikzpicture}
\caption{原分数区间经过单调映射成为新分数区间}\label{mp1:fig:rev:score-q}
\end{figure}
\end{minipage}\par\medskip

\begin{solution}
原密度为$f_X=1/99$，而反函数与导数为
\[
 X=Y^2/100,\qquad\frac{\dd X}{\dd Y}=\frac Y{50}.
\]
概率微元满足$f_Y\dd Y=f_X\dd X$，故
\[
 f_Y(y)=\begin{cases}y/4950,&10<y<100,\\0,&\text{其他},\end{cases}
 \quad
 F_Y(y)=\begin{cases}0,&y<10,\\(y^2-100)/9900,&10\le y\le100,\\1,&y>100.\end{cases}
\]
在支持区间内部，累计概率由
\[
F_Y(y)=\int_{10}^y\frac{t}{4950}\dd t
=\left[\frac{t^2}{9900}\right]_{10}^y
=\frac{y^2-100}{9900}
\]
给出，取 $y=100$ 即得总概率1。图\ref{mp1:fig:rev:score-s}中的密度随新分数线性增大。原来等宽、等概率的$\dd X$小段经过映射后成为不同宽度的$\dd Y$小段；高分处映射斜率较小，新密度因而较高。

若原区间为$[0,100]$，则$f_X=1/100$、新支持区间为$[0,100]$，结果改为$f_Y=y/5000$、$F_Y=y^2/10000$。改变原分数区间时，新分数的下限和密度的归一化常数都须重新确定。

\begin{figure}[H]\centering
\begin{tikzpicture}[>=Stealth,line cap=round,line join=round]
\draw[axis](0,0)--(4.3,0)node[right]{$y$};\draw[axis](0,0)--(0,2.5)node[above]{$f_Y$};\draw[curve](.4,.2)--(4,2);\draw[guide](.4,0)--(.4,.2);\draw[guide](4,0)--(4,2);\node[below]at(.4,0){$10$};\node[below]at(4,0){$100$};\node[right]at(2.3,.8){$f_Y=y/4950$};
\end{tikzpicture}
\caption{密度在支持区间内随新分数线性增大}\label{mp1:fig:rev:score-s}
\end{figure}

\end{solution}

\section{角分布、立体角与靶面密度}
方向均匀必须说明是对平面角还是对立体角均匀。发射方向映到靶面后，距离与入射角共同影响密度，有限靶面的边界限定接收范围。

\subsection{二维方向与三维方向}
平面中的方向用一个角 $\theta$ 表示。若全方向均匀，任一角区间的概率为
\[
\dd P=\frac{\dd\theta}{2\pi}.
\]
空间方向则需两个角。取极角 $0\le\theta\le\pi$、方位角 $0\le\phi<2\pi$，极角从正 $z$ 轴量起。单位球面上两条小边的长度为 $\dd\theta$ 与 $\sin\theta\dd\phi$，故立体角微元为
\begin{equation}\label{mp1:eq:solid-angle}
\dd\Omega=\sin\theta\dd\theta\dd\phi.
\end{equation}
图\ref{mp1:fig:solid}表明，相同的 $\dd\theta$ 在赤道附近覆盖的方向比在极点附近更多。对半顶角 $\alpha$ 的圆锥，先对方位角、再对极角积分：
\[
\Omega(\alpha)=[\phi]_0^{2\pi}\int_0^\alpha\sin\theta\dd\theta
=2\pi[-\cos\theta]_0^\alpha=2\pi(1-\cos\alpha).
\]
取 $\alpha=\pi$，全空间立体角为 $4\pi$。空间各向同性指每单位立体角的密度相同，因此
\[
\dd P=\frac{\dd\Omega}{4\pi}.
\]
对方位角积分后，极角密度为 $p_\theta=\tfrac12\sin\theta$。极角并非均匀分布；方位角在各极角下均匀，也不足以保证空间各向同性。

\begin{figure}[H]\centering

\begin{tikzpicture}[scale=1.1]
\pgfmathsetmacro{\sr}{1.6}
\pgfmathsetmacro{\viewangle}{12}
\pgfmathsetmacro{\tupper}{45}
\pgfmathsetmacro{\tlower}{55}
\pgfmathsetmacro{\ru}{\sr*sin(\tupper)}
\pgfmathsetmacro{\rl}{\sr*sin(\tlower)}
\pgfmathsetmacro{\cu}{\sr*cos(\tupper)*cos(\viewangle)}
\pgfmathsetmacro{\cl}{\sr*cos(\tlower)*cos(\viewangle)}
\pgfmathsetmacro{\eu}{\ru*sin(\viewangle)}
\pgfmathsetmacro{\el}{\rl*sin(\viewangle)}
\fill[black!12](-\ru,\cu)arc[start angle=180,end angle=360,x radius=\ru,y radius=\eu]--(\rl,\cl)arc[start angle=0,end angle=-180,x radius=\rl,y radius=\el]--cycle;
\fill[black!8](-\ru,\cu)arc[start angle=180,end angle=0,x radius=\ru,y radius=\eu]--(\rl,\cl)arc[start angle=0,end angle=180,x radius=\rl,y radius=\el]--cycle;
\draw(0,0)circle(\sr);
\draw[guide](-\sr,0)arc[start angle=180,end angle=0,x radius=\sr,y radius=.333];
\draw(-\sr,0)arc[start angle=180,end angle=360,x radius=\sr,y radius=.333];
\draw[axis](0,-1.9)--(0,2.05)node[above]{$z$};
\draw[axis](0,0)--(2.05,0)node[right]{$x$};
\draw(-\ru,\cu)arc[start angle=180,end angle=360,x radius=\ru,y radius=\eu];
\draw[guide](-\ru,\cu)arc[start angle=180,end angle=0,x radius=\ru,y radius=\eu];
\draw(-\rl,\cl)arc[start angle=180,end angle=360,x radius=\rl,y radius=\el];
\draw[guide](-\rl,\cl)arc[start angle=180,end angle=0,x radius=\rl,y radius=\el];
\draw(0,0)--(\ru,\cu);
\draw[guide](0,0)--(\rl,\cl);
\pgfmathsetmacro{\angleu}{atan2(\cu,\ru)}
\pgfmathsetmacro{\anglel}{atan2(\cl,\rl)}
\draw[domain=0:\tupper,samples=30]plot({.5*sin(\x)},{.5*cos(\x)*cos(\viewangle)});
\node at(.22,.61){$\theta$};
\draw[domain=\tupper:\tlower,samples=20]plot({.91*sin(\x)},{.91*cos(\x)*cos(\viewangle)});
\node[anchor=west,fill=white,inner sep=1pt]at(.9,.63){$\dd\theta$};
\node[anchor=east]at(-1.42,1.1){方向带};
\node[below left]at(0,0){$O$};
\node[anchor=west]at(2.7,.65){$\dd\Omega=\sin\theta\dd\theta\dd\phi$};
\node[anchor=west]at(2.7,-.1){$\dd\Omega_{\rm ring}=2\pi\sin\theta\dd\theta$};
\end{tikzpicture}
\caption{单位球面上的方向带。极角为 $\theta$ 的整圈周长为 $2\pi\sin\theta$，带宽为 $\dd\theta$。}\label{mp1:fig:solid}
\end{figure}

当 $\alpha\ll1$ 时，$\Omega\simeq\pi\alpha^2$。圆锥很窄时，覆盖方向数随角宽的平方增加。

\subsection{从一个点看面积微元}
设源点为 $O$，面积微元位于 $\vct R$，其单位法线为 $\vct n$。微元在垂直射线方向上的投影面积为 $|\vct n\cdot\widehat{\vct R}|\dd S$。将投影除以距离平方，就得到张开的立体角：
\begin{equation}\label{mp1:eq:projected-solid}
\dd\Omega=\frac{|\vct n\cdot\widehat{\vct R}|}{R^2}\dd S
=\frac{|\vct n\cdot\vct R|}{R^3}\dd S.
\end{equation}
图\ref{mp1:fig:project}中的距离因子来自球面面积随半径平方增加，余弦因子来自面积的倾斜投影。二者均不能省去。

\begin{figure}[H]\centering

\begin{tikzpicture}[scale=1.05]
\coordinate(O)at(0,0);\coordinate(P)at(3.6,0);
\coordinate(A)at(2.95,-.9);\coordinate(C)at(4.25,.9);
\draw[axis](O)--(5.25,0)node[right]{$\widehat{\vct R}$};
\draw[line width=1pt](A)--(C);
\draw[line width=1pt](3.6,-.9)--(3.6,.9);
\draw[guide](A)--(3.6,-.9) (C)--(3.6,.9);
\draw[axis](P)--(4.6,-.7222)node[right]{$\vct n$};
\draw(4.15,0)arc(0:-35.84:.55);
\node at(4.42,-.22){$\gamma$};
\fill(O)circle(1.3pt)node[below left]{$O$};
\fill(P)circle(1.3pt)node[below left]{$P$};
\node[above right=2pt]at(C){$\dd S$};
\node[above left=2pt]at(3.6,.9){$\dd S_\perp$};
\draw[guide](O)--(0,-1.4) (P)--(3.6,-1.4);
\draw[<->](0,-1.3)--(3.6,-1.3)node[midway,fill=white]{$R$};
\node[anchor=west]at(5.9,.6){$\dd S_\perp=|\cos\gamma|\dd S$};
\node[anchor=west]at(5.9,-.2){$\dd\Omega=\dd S_\perp/R^2$};
\end{tikzpicture}
\caption{包含射线与法线的截面。两条粗线表示面元和投影面的截线，垂直纸面方向的宽度不变。投影沿 $\widehat{\vct R}$ 平行进行，$R=OP$，$\gamma$ 为法线与射线的夹角。}\label{mp1:fig:project}
\end{figure}

若源点的角分布写为 $\dd P=f_\Omega\dd\Omega$，则对应的面密度为
\[
g_S=f_\Omega\frac{|\vct n\cdot\vct R|}{R^3}.
\]
这个公式按射线到达微元的同一条分支使用。若表面遮挡了自身，或射线先被别处截获，须先确定可见区域。

对正对源点、位于 $z=h>0$ 的平面，有 $R^2=h^2+x^2+y^2$，$|\vct n\cdot\vct R|=h$。因此
\[
\dd\Omega=\frac{h\dd x\dd y}{(h^2+x^2+y^2)^{3/2}}.
\]
若只讨论朝向该平面的半空间，仍需说明密度是对全部发射事件定义，还是以射线朝向平面为条件定义。两者相差的常数由条件事件的概率确定。

\subsection{无条件密度与命中后的条件密度}
当部分事件不能产生所讨论的坐标时，在靶上积分得到的总概率小于1。例如向后发射的射线不会到达前方平面。靶上的密度可以先对全部发射定义，积分给出到达概率；若已知到达，再除以这个到达概率，才成为归一化到1的条件密度。

\par\medskip\noindent\begin{minipage}{\textwidth}
\begin{example}[平面随机射线与无限长靶线]\label{mp1:ex:e4}
在平面内从 $O$ 向全方向均匀发射射线，忽略重力。靶线位于 $x=h>0$，其坐标为 $y$，见图\ref{mp1:fig:e4-q}。求对全部发射事件定义的靶线密度、已知到达靶线后的条件密度，以及击中有限靶段 $|y|\le b$ 的概率。
\end{example}

\begin{figure}[H]\centering

\begin{tikzpicture}[scale=1.05]
\draw[axis](-1.3,0)--(4.2,0)node[right]{$x$};\draw[thick](3,-1.5)--(3,2.3)node[above]{$x=h$};\fill(0,0)circle(1.5pt)node[below]{$O$};\draw[axis](0,0)--(3,1.5)node[right]{$P(h,y)$};
\draw[guide](0,0)circle(.8);\draw(0.55,0)arc(0:26.565:.55);\node[right]at(.68,.18){$\theta$};\draw[<->](0,-1)--(3,-1)node[midway,fill=white]{$h$};\node[right]at(3,-1.5){靶线};
\end{tikzpicture}
\caption{例~\ref{mp1:ex:e4}题设：平面内全方向发射。只有朝向靶线的半圆方向能到达 $x=h$。}\label{mp1:fig:e4-q}
\end{figure}

\end{minipage}\par\medskip

\begin{solution}
朝向靶线的方向为 $-\pi/2<\theta<\pi/2$，占全圆的一半。对这一区域，$y=h\tan\theta$ 单调增加，有
\[
\frac{\dd y}{\dd\theta}=h\sec^2\theta,
\qquad \frac{\dd\theta}{\dd y}=\frac h{h^2+y^2}.
\]
图\ref{mp1:fig:e4-s}中同一个细角区间与靶线段相对应，故
\[
g(y)\dd y=\frac{\dd\theta}{2\pi},\qquad
g(y)=\frac h{2\pi(h^2+y^2)}.
\]
整条靶线的到达概率可由下述有限区间积分取 $b\to\infty$ 求出。

\begin{figure}[H]\centering

\begin{tikzpicture}[scale=1.1]
\fill(0,0)circle(1.5pt)node[below]{$O$};\draw[guide](0,0)--(3.2,0);\draw[thick](3,-1.9)--(3,1.9);
\draw(0,0)--(3,.65);\draw(0,0)--(3,.95);\draw[line width=2pt](3,.65)--(3,.95)node[midway,right]{$\dd y$};\node[left]at(3,.65){$y$};\draw(1.27,.275)arc(12.23:17.57:1.3);\node[above,fill=white,inner sep=1pt]at(1.4,.44){$\dd\theta$};
\draw[guide](0,0)--(3,-1.6) (0,0)--(3,1.6);\fill(3,-1.6)circle(1.5pt)node[right]{$-b$};\fill(3,1.6)circle(1.5pt)node[right]{$b$};\node[below]at(1.5,0){$h$};
\end{tikzpicture}
\caption{例~\ref{mp1:ex:e4}解图：靶上线段 $\dd y$ 与发射点处的角区间 $\dd\theta$ 对应。有限靶段由两条边界射线确定。}\label{mp1:fig:e4-s}
\end{figure}

有限靶段的边界射线满足 $\theta=\pm\arctan(b/h)$，所以对全部发射，
\[
\Prob(|y|\le b,\text{到达})=\int_{-b}^b g(y)\dd y
=\frac1{2\pi}\left[\arctan\frac yh\right]_{-b}^b
=\frac1\pi\arctan\frac bh.
\]
当 $b\to\infty$，$\arctan(b/h)\to\pi/2$，故整条靶线的到达概率为 $1/2$。于是条件密度为
\[
g(y\mid\text{到达})=\frac{g(y)}{1/2}
=\frac h{\pi(h^2+y^2)}.
\]
若已知到达前方，有限靶段的条件命中率也是无条件概率的两倍，其 $b\to\infty$ 极限为1。
\end{solution}

\par\medskip\noindent\begin{minipage}{\textwidth}
\begin{example}[空间各向同性点源与圆靶]\label{mp1:ex:e5}
点源向全空间各向同性发射直线粒子，忽略重力。圆靶位于 $z=h>0$，圆心在 $z$ 轴上，半径为 $a$，见图\ref{mp1:fig:e5-q}。求靶面单位面积密度、靶上径向密度及命中圆靶的概率。比较两种密度的最大位置。
\end{example}

\begin{figure}[H]\centering

\begin{tikzpicture}[scale=1.1]
\draw[axis](-.2,0)--(5,0)node[right]{$z$};\fill(0,0)circle(1.6pt)node[below]{$O$};\draw(3.5,0)ellipse(.45 and 1.6);\draw(0,0)--(3.5,1.6) (0,0)--(3.5,-1.6);
\draw[thick](3.5,0)--(3.5,1)node[midway,right]{$r$};\fill(3.5,1)circle(1.3pt)node[right]{$P$};\draw[guide](0,0)--(3.5,1);\draw(.8,0)arc(0:15.945:.8);\node at(1.05,.19){$\theta$};\node[below]at(1.8,0){$h$};\node[right]at(3.7,-1.6){半径 $a$};
\end{tikzpicture}
\caption{例~\ref{mp1:ex:e5}题设：空间各向同性点源与正对点源的圆靶。靶面上的径向坐标为 $r$。}\label{mp1:fig:e5-q}
\end{figure}

\end{minipage}\par\medskip

\begin{solution}
先考虑圆靶所在的整平面。对全空间各向同性源，$f_\Omega=1/(4\pi)$。在这个平面上，以圆心为原点，$r^2=x^2+y^2$，射线距离为 $\sqrt{h^2+r^2}$。由投影立体角公式，
\begin{equation}\label{mp1:eq:isotropic-plane}
g_S(r)=\frac h{4\pi(h^2+r^2)^{3/2}}.
\end{equation}
这表示每单位面积的概率。实际圆靶只保留 $0\le r\le a$，靶外不计命中。图\ref{mp1:fig:e5-s}左图中的整圈环带面积为 $2\pi r\dd r$，对整平面的径向密度为
\[
p_r(r)=2\pi r g_S(r)=\frac{hr}{2(h^2+r^2)^{3/2}},\qquad r>0.
\]
积分到半径 $a$ 时，令 $u=h^2+r^2$，$\dd u=2r\dd r$，上下限变为 $h^2,h^2+a^2$，得到
\begin{align*}
P_a&=\int_0^a p_r(r)\dd r
=\frac h4\int_{h^2}^{h^2+a^2}u^{-3/2}\dd u\\
&=\frac h4\left[-2u^{-1/2}\right]_{h^2}^{h^2+a^2}
=\frac12\left(1-\frac h{\sqrt{h^2+a^2}}\right).
\end{align*}
圆靶的半张角满足 $\cos\alpha=h/\sqrt{h^2+a^2}$，这也就是 $\Omega(\alpha)/(4\pi)$。

\begin{figure}[H]\centering

\begin{tikzpicture}[scale=.95]
\fill[even odd rule,black!12](0,0)circle(1.1) (0,0)circle(.9);
\draw(0,0)circle(1.6);\draw(0,0)circle(.9);\draw(0,0)circle(1.1);
\fill(0,0)circle(1pt);
\draw(0,0)--(-35:.9)node[midway,below left=1pt]{$r$};
\draw[|-|](45:.9)--(45:1.1)node[midway,above right=3pt]{$\dd r$};
\node[below]at(0,-2){$\dd S=2\pi r\dd r$};
\begin{scope}[shift={(4,-1)},x=1.4cm,y=4cm]
\draw[axis](0,0)--(3.1,0)node[right]{$r/h$};
\draw[axis](0,0)--(0,.5)node[above]{$2hp_r$};
\draw[curve,domain=0:3,samples=100]plot(\x,{\x/(1+(\x)^2)^1.5});
\draw[guide](.707107,0)--(.707107,.3849);
\node[below]at(.707107,0){$1/\sqrt2$};
\end{scope}
\end{tikzpicture}
\caption{例~\ref{mp1:ex:e5}解图：左图按环带积分，面积为 $2\pi r\dd r$；右图是整平面的径向密度，在 $r/h=1/\sqrt2$ 达到最大。}\label{mp1:fig:e5-s}
\end{figure}

面密度 $g_S$ 随 $r$ 增大而下降，在圆心最大；径向密度则含有环带周长因子。对它求导：
\[
\begin{aligned}
p_r'(r)&=\frac h2\left[(h^2+r^2)^{-3/2}
-3r^2(h^2+r^2)^{-5/2}\right]\\
&=\frac h2\frac{h^2-2r^2}{(h^2+r^2)^{5/2}},
\end{aligned}
\]
所以在无限靶面上，径向密度于 $r=h/\sqrt2$ 达到最大，如图\ref{mp1:fig:e5-s}右图。若圆靶半径小于该值，靶内径向密度一直增加，到边缘最大。

整平面的到达概率为 $1/2$。已知朝向平面时，上述密度均乘2；若已知击中有限圆靶，则在 $0\le r\le a$ 内除以 $P_a$。半径 $r$ 的量纲是长度，$p_r$ 的量纲为长度的倒数，$g_S$ 则为面积的倒数。
\end{solution}

\subsection{轴对称角分布的三种密度}
对绕 $z$ 轴对称的分布，按单位立体角的密度 $f_\Omega(\theta)$、按极角的密度 $p_\theta(\theta)$ 与按平面面积的密度 $g_S(r)$ 满足
\begin{align*}
p_\theta(\theta)&=2\pi\sin\theta\,f_\Omega(\theta),\\
g_S(r)&=f_\Omega(\theta)\frac h{(h^2+r^2)^{3/2}},
\qquad r=h\tan\theta.
\end{align*}
第一式合并同一极角处的一整圈方向，第二式把方向微元投影到靶面。三种密度采用不同的微元，最大值的位置也可能不同。

\par\medskip\noindent\begin{minipage}{\textwidth}
\begin{example}[上半空间的余弦分布]\label{mp1:ex:e6}
一个定向源只向上半空间发射，其每单位立体角的密度为 $f_\Omega(\theta)=C\cos\theta$，$0\le\theta\le\pi/2$。图\ref{mp1:fig:e6-q}中的圆靶半径为 $a$、距离为 $h$。求 $C$、极角密度、靶面密度和命中概率。
\end{example}

\begin{figure}[H]\centering

\begin{tikzpicture}[scale=1.05]
\draw[thick](-1.5,0)--(1.5,0);\draw[axis](0,-.15)--(0,3.5)node[above]{$z$};\fill(0,0)circle(1.5pt)node[below]{$O$};\draw(0,2.6)ellipse(1.65 and .3);\draw(0,0)--(1.65,2.6) (0,0)--(-1.65,2.6);\draw[guide](0,2.6)--(1.65,2.6)node[midway,above]{$a$};
\draw(0,.7)arc(90:57.6:.7);\node at(.3,1){$\alpha$};\node[left]at(0,1.9){$h$};\node[right]at(2,.7){$f_\Omega=C\cos\theta$};
\end{tikzpicture}
\caption{例~\ref{mp1:ex:e6}题设：只向上半空间发射的余弦分布模型。极角从法线 $z$ 轴量起。}\label{mp1:fig:e6-q}
\end{figure}

\end{minipage}\par\medskip

\begin{solution}
将整个上半空间的概率积分到1，
\[
1=2\pi C\int_0^{\pi/2}\cos\theta\sin\theta\dd\theta
=2\pi C\left[\frac12\sin^2\theta\right]_0^{\pi/2}=\pi C,
\]
故 $C=1/\pi$。合并方位角后的极角密度为
\[
p_\theta(\theta)=2\sin\theta\cos\theta=\sin2\theta.
\]
图\ref{mp1:fig:e6-s}左图中，该密度在 $\theta=\pi/4$ 最大，虽然每单位立体角的密度在轴线上最大。

\begin{figure}[H]\centering

\begin{tikzpicture}[scale=.95]
\begin{scope}[x=1.6cm,y=1.6cm]
\draw[axis](0,0)--(1.86,0)node[right]{$\theta$};
\draw[axis](0,0)--(0,1.3)node[above]{$p_\theta$};
\draw[curve,domain=0:90,samples=80]plot({pi*\x/180},{sin(2*\x)});
\draw[guide](0,1)--(.785398,1)--(.785398,0);
\node[below]at(.785398,0){$\pi/4$};
\node[below]at(1.570796,0){$\pi/2$};
\node[left]at(0,1){$1$};
\end{scope}
\begin{scope}[shift={(4.5,0)},x=1.2cm,y=2cm]
\draw[axis](0,0)--(3.3,0)node[right]{$a/h$};
\draw[axis](0,0)--(0,1.2)node[above]{$P$};
\draw[curve,domain=0:3,samples=80]plot(\x,{(\x)^2/(1+(\x)^2)});
\draw[guide](0,1)--(3,1);
\draw[guide](0,.5)--(1,.5)--(1,0);
\fill(1,.5)circle(.8pt);
\node[left]at(0,1){$1$};\node[left]at(0,.5){$1/2$};\node[below]at(1,0){$1$};
\end{scope}
\end{tikzpicture}
\caption{例~\ref{mp1:ex:e6}解图：方向环带的概率为 $2\sin\theta\cos\theta\dd\theta$。右图给出圆靶截获比例随 $a/h$ 的变化。}\label{mp1:fig:e6-s}
\end{figure}

靶面面积还要乘一次投影余弦。由于 $\cos\theta=h/\sqrt{h^2+r^2}$，
\[
g_S(r)=\frac{h^2}{\pi(h^2+r^2)^2}.
\]
圆靶的命中率为
\begin{align*}
P_a&=\int_0^a g_S(r)2\pi r\dd r
=\int_0^a\frac{2h^2r}{(h^2+r^2)^2}\dd r\\
&=h^2\int_{h^2}^{h^2+a^2}u^{-2}\dd u
=h^2\left[-\frac1u\right]_{h^2}^{h^2+a^2}\\
&=1-\frac{h^2}{h^2+a^2}
=\frac{a^2}{h^2+a^2}=\sin^2\alpha.
\end{align*}
这里仍取 $u=h^2+r^2$、$\dd u=2r\dd r$。图\ref{mp1:fig:e6-s}右图中 $a=h$ 对应截获一半粒子。$a\to\infty$ 时概率趋于1，因为本题的源本就只向靶所在的上半空间发射。
\end{solution}












\par\medskip\noindent\begin{minipage}{\textwidth}
\begin{example}[等厚球带与经纬度区域]\label{mp1:rev:earth}
半径为$R$的球面均匀光滑。证明任意两个相距相同高度的平行平面所截球带面积相同。再把地球视为半径$R$的球，求北纬$30^\circ$至$31^\circ$、东经$65^\circ$至$66^\circ$的区域面积占全球面积的比例。图\ref{mp1:fig:rev:earth-q}标出纬度与球带高度。
\end{example}

\begin{figure}[H]\centering
\begin{tikzpicture}[>=Stealth,line cap=round,line join=round]
\fill[gray!25]({1.5*cos(30)},{1.5*sin(30)})arc(30:31:1.5)--({-1.5*cos(31)},{1.5*sin(31)})arc(149:150:1.5)--cycle;
\draw(0,0)circle(1.5);\draw[axis](0,-1.75)--(0,1.85)node[above]{$z$};\draw[guide](-1.5,0)--(1.7,0);
\draw[curve]({-1.5*cos(30)},{1.5*sin(30)})--({1.5*cos(30)},{1.5*sin(30)});
\draw[curve]({-1.5*cos(31)},{1.5*sin(31)})--({1.5*cos(31)},{1.5*sin(31)});
\draw[axis](0,0)--({1.5*cos(30)},{1.5*sin(30)})node[midway,above left]{$R$};\draw(.6,0)arc(0:30:.6);\node at(.8,.14){$\lambda$};
\draw[guide]({1.5*cos(30)},{1.5*sin(30)})--(2.4,.2);\draw[guide]({1.5*cos(31)},{1.5*sin(31)})--(2.4,1.2);
\fill[gray!25](2.4,.2)rectangle(3.4,1.2);\draw[curve](2.4,.2)--(3.4,.2);\draw[curve](2.4,1.2)--(3.4,1.2);
\node[right]at(3.45,.2){$z_1=R\sin30^\circ$};\node[right]at(3.45,1.2){$z_2=R\sin31^\circ$};\node at(2.9,1.65){高度差放大};
\end{tikzpicture}
\caption{纬度从赤道量起，球带高度由$z=R\sin\lambda$确定。右侧放大高度差，以区分$30^\circ$与$31^\circ$两条边界。}\label{mp1:fig:rev:earth-q}
\end{figure}
\end{minipage}\par\medskip

\begin{solution}
设球面极角为$\theta$，球带微元面积为$\dd S=2\pi R^2\sin\theta\dd\theta$，而$|\dd z|=R\sin\theta\dd\theta$，故
\[
 \dd S=2\pi R|\dd z|,\qquad
 S=2\pi R\int_{z_1}^{z_2}\dd z
 =2\pi R[z]_{z_1}^{z_2}=2\pi R(z_2-z_1).
\]
面积只取决于两平面的高度差，与球带所在高度无关。

经度只截取整圈球带的一部分。图\ref{mp1:fig:rev:earth-s}中经度宽度为$\Delta\phi=\pi/180$，两纬线对应的高度差为$R(\sin31^\circ-\sin30^\circ)$，所以
\[
 S_{\text{区}}=R^2\int_{\lambda_1}^{\lambda_2}\cos\lambda\dd\lambda
 \int_{\phi_1}^{\phi_2}\dd\phi
 =R^2[\sin\lambda]_{\lambda_1}^{\lambda_2}[\phi]_{\phi_1}^{\phi_2},
 \]
其中 $\lambda_1=\pi/6$、$\lambda_2=31\pi/180$、$\phi_2-\phi_1=\pi/180$。代入后，
\[
 S_{\text{区}}=R^2\Delta\phi(\sin31^\circ-\sin30^\circ),
 \qquad
 \boxed{\frac{S_{\text{区}}}{4\pi R^2}
 =\frac{\sin31^\circ-\sin30^\circ}{720}
 \simeq2.0886\times10^{-5}}.
\]
比例约为$0.0020886\%$。纬度与经度的角宽均为$1^\circ$，但纬线长度随纬度变化，计算面积时须保留这一因素。

\begin{figure}[H]\centering
\begin{tikzpicture}[>=Stealth,line cap=round,line join=round]
\draw[axis](0,0)--(4.7,0)node[right]{$\phi$};\draw[axis](0,0)--(0,2.7)node[above]{$z$};\fill[gray!20](1,.8)rectangle(3.5,1.75);\draw[curve](1,.8)rectangle(3.5,1.75);\node[below]at(1,0){$65^\circ$};\node[below]at(3.5,0){$66^\circ$};\draw[guide](1,0)--(1,.8);\draw[guide](3.5,0)--(3.5,.8);\node[left]at(0,.8){$R\sin30^\circ$};\node[left]at(0,1.75){$R\sin31^\circ$};\draw[guide](0,.8)--(1,.8);\draw[guide](0,1.75)--(1,1.75);\node at(2.25,2.3){$\dd S=R\dd\phi\dd z$};
\end{tikzpicture}
\caption{以经度和高度作坐标后，球面面积微元具有常系数}\label{mp1:fig:rev:earth-s}
\end{figure}

\end{solution}

\section{连续受力与碰撞散射}
连续受力从微元平衡累加，散射计数则把入射窄带换成散射角区间。两类模型分别保留积分边界和参考系约定，再由局部关系求总量。

\subsection{连续分布的静力平衡}
若作用力连续分布，先明确载荷对什么长度或面积定义。沿弧长 $s$ 的线载荷写为 $\dd\vct F=\vct q_s\dd s$；对水平投影长度定义的载荷则写为 $\dd\vct F=\vct q_x\dd x$。即使两者的系数都是常数，它们也描述不同的物理分布。

对细长系统，总力与对点 $O$ 的力矩由各微元相加：
\[
\vct F=\int\dd\vct F,\qquad
\vct M_O=\int\vct r\times\dd\vct F.
\]
平行均匀重力可通过质心合并；非平行分布力应分别积分力与力矩。局部受力方程也可沿材料逐段积分，以得到张力、内力或压力的分布。

对于不抗弯的理想柔链，张力沿切线。隔离 $s$ 到 $s+\dd s$ 的一小段，左右端作用力方向相反，局部平衡式为
\[
\frac{\dd(T\vct t)}{\dd s}+\vct q_s=0.
\]
展开乘积导数，局部方程为
\[
\frac{\dd T}{\dd s}\vct t+T\frac{\dd\vct t}{\dd s}+\vct q_s=0.
\]
张力大小和切向量都可以随位置改变，两项分别表示大小变化与方向变化的影响。下面的均匀链条以真实弧长定义质量分布，因此重力微元为 $\lambda g\dd s$。

\par\medskip\noindent\begin{minipage}{\textwidth}
\begin{example}[均匀柔链的平衡形状]\label{mp1:ex:e7}
一条线密度 $\lambda$ 恒定、不可伸长的柔链，两端等高，水平距离为 $2a$，链长为 $\ell>2a$。重力加速度为 $g$。按图\ref{mp1:fig:e7-q}以最低点为原点，求平衡形状、两端张力及链长与垂度的关系。
\end{example}

\begin{figure}[H]\centering

\begin{tikzpicture}[scale=1.3]
\draw[axis](-2,0)--(2.2,0)node[right]{$x$};\draw[axis](0,-.2)--(0,2.3)node[above]{$y$};\draw[curve,domain=-1.6:1.6,samples=100]plot(\x,{.5*(exp(\x)+exp(-(\x)))-1});
\fill(-1.6,1.5775)circle(2pt)node[left]{$A$};\fill(1.6,1.5775)circle(2pt)node[right]{$B$};\draw[guide](-1.6,1.5775)--(1.6,1.5775);\node[below left]at(0,0){$O$};\draw[<->](1.95,0)--(1.95,1.5775)node[midway,right]{$f$};\draw[<->](0,-.4)--(1.6,-.4)node[midway,fill=white]{$a$};\node[left]at(-1.15,.9){$\lambda$};\node[above]at(-.75,1.95){链长 $\ell$};
\end{tikzpicture}
\caption{例~\ref{mp1:ex:e7}题设：两端等高的均匀柔链。以最低点为原点，$a$ 为半跨距，$f$ 为垂度。}\label{mp1:fig:e7-q}
\end{figure}

\end{minipage}\par\medskip

\begin{solution}
对图\ref{mp1:fig:e7-s}中的微元列水平、竖直平衡。由于重力竖直，张力的水平分量为常量，记为 $H>0$。若切线倾角为 $\psi$，则
\[
T_y=H\tan\psi=Hy',\qquad \dd T_y=\lambda g\dd s.
\]
在链条可写为单值函数 $y(x)$ 的区域，$\dd s=\sqrt{1+y'^2}\dd x$，从而
\begin{equation}\label{mp1:eq:catenary-ode}
Hy''=\lambda g\sqrt{1+y'^2}.
\end{equation}
令 $c=H/(\lambda g)$，$p=y'$。分离变量，以最低点的 $p(0)=0$ 为初值，积分得
\[
\int_0^p\frac{\dd u}{\sqrt{1+u^2}}
=\left[\ln\left(u+\sqrt{1+u^2}\right)\right]_0^p
=\frac xc.
\]
于是 $p+\sqrt{1+p^2}=\ee^{x/c}$，移项平方，得到 $2p\ee^{x/c}=\ee^{2x/c}-1$，即
\[
p=\frac{\ee^{x/c}-\ee^{-x/c}}2=\sinh\frac xc.
\]
再次积分，并取 $y(0)=0$，得到悬链线
\[
y(x)=\int_0^x\sinh\frac\xi c\dd\xi
=\left[c\cosh\frac\xi c\right]_0^x
=c\left(\cosh\frac xc-1\right),\qquad -a\le x\le a.
\]
这里 $\cosh u=(\ee^u+\ee^{-u})/2$，$\sinh u=(\ee^u-\ee^{-u})/2$。

\begin{figure}[H]\centering

\begin{tikzpicture}[scale=1.2]
\draw[curve,domain=.65:1.15,samples=40]plot(\x,{.5*(exp(\x)+exp(-(\x)))-1});\fill(.65,.2188)circle(1.2pt);\fill(1.15,.7374)circle(1.2pt);
\draw[axis,thick](.65,.2188)--(-.45,-.5476)node[left]{$\vct T_-$};\draw[axis,thick](1.15,.7374)--(2.25,2.3003)node[above]{$\vct T_+$};
\draw[axis](.9,.45)--(.9,-.3464)node[below]{$\lambda g\dd s$};\draw[guide](.65,.2188)--(.65,1.1) (1.15,.7374)--(1.15,1.1);\draw[<->](.65,1.1)--(1.15,1.1)node[midway,above]{$\dd x$};
\begin{scope}[shift={(4.2,.2)}]
\draw[axis](0,0)--(2,0)node[right]{$H$};\draw[axis](2,0)--(2,1.6)node[right]{$T_y$};\draw[axis,thick](0,0)--(2,1.6)node[above left]{$T$};\draw(.7,0)arc(0:38.66:.7);\node at(.85,.25){$\psi$};\node[below]at(1,-.55){$\tan\psi=y'$};
\end{scope}
\end{tikzpicture}
\caption{例~\ref{mp1:ex:e7}解图：隔离一小段链条，左右端张力沿各自切线，重力为 $\lambda g\dd s$。水平分量相同，竖直分量的差承担微元重力。}\label{mp1:fig:e7-s}
\end{figure}

由 $\dd s/\dd x=\sqrt{1+\sinh^2(x/c)}=\cosh(x/c)$，右半链的长度为
\[
s(x)=\int_0^x\cosh\frac\xi c\dd\xi
=\left[c\sinh\frac\xi c\right]_0^x=c\sinh\frac xc.
\]
因此
\[
\ell=2c\sinh\frac ac,\qquad
f=c\left(\cosh\frac ac-1\right).
\]
给定 $a,\ell$ 后，第一式确定 $c$。令 $u=a/c>0$，则 $\ell/(2a)=\sinh u/u$。求导，
\begin{align*}
\frac{\dd}{\dd u}\frac{\sinh u}{u}
&=\frac{u\cosh u-\sinh u}{u^2},\\
\frac{\dd}{\dd u}(u\cosh u-\sinh u)
&=\cosh u+u\sinh u-\cosh u\\
&=u\sinh u>0.
\end{align*}
分子从零开始增加，故 $\sinh u/u$ 从1单调增加到无穷，对任意 $\ell>2a$ 都有唯一正解。

图\ref{mp1:fig:e7-s}的张力三角形给出
\[
T(x)=\sqrt{H^2+T_y^2}=H\cosh\frac xc.
\]
所以最低点张力为 $H=\lambda gc$，两端张力均为
\[
T_A=T_B=\lambda gc\cosh\frac ac
=\sqrt{H^2+\left(\frac{\lambda g\ell}{2}\right)^2}.
\]
两端张力的竖直分量合起来承担整条链重。考察浅垂极限，令 $u=a/c\ll1$，则
\[
\cosh u-1=\frac{u^2}{2}+O(u^4),\qquad
\sinh u-u=\frac{u^3}{6}+O(u^5).
\]
分别代入垂度与总长的表达式，得到
\[
f=c(\cosh u-1)=\frac{a^2}{2c}[1+O(u^2)],\qquad
\ell-2a=2c(\sinh u-u)=\frac{a^3}{3c^2}[1+O(u^2)].
\]
由第一式取最低阶关系 $c\simeq a^2/(2f)$，代入第二式及 $H=\lambda gc$，便有
\[
f\simeq\frac{a^2}{2c},\qquad
\ell-2a\simeq\frac{a^3}{3c^2}\simeq\frac{4f^2}{3a},\qquad
H\simeq\frac{\lambda ga^2}{2f}.
\]
垂度很小时，保持给定跨距需要很大的水平张力。
\end{solution}

\subsection{入射截面与散射方向}
一束平行粒子入射时，用瞄准距离（碰撞参数）$b$ 表示粒子初始轨迹到靶心的垂直距离。若通量密度为 $j$，单位为 $\mathrm{m^{-2}s^{-1}}$，则面积 $\dd\sigma$ 内在时间 $\tau$ 中入射的粒子数为
\[
\dd N=j\tau\dd\sigma.
\]
对轴对称碰撞几何，$b$ 到 $b+\dd b$ 的整圈入射环带面积为 $2\pi b\dd b$。相应事件的截面微元就是这块入射面积。

若 $b$ 与出射极角 $\theta$ 一一对应，微元在变换中满足
\[
\dd\sigma=2\pi b\left|\frac{\dd b}{\dd\theta}\right|\dd\theta,
\qquad \dd\Omega=2\pi\sin\theta\dd\theta.
\]
从而微分截面为
\begin{equation}\label{mp1:eq:differential-cross-section}
\frac{\dd\sigma}{\dd\Omega}
=\frac b{\sin\theta}\left|\frac{\dd b}{\dd\theta}\right|.
\end{equation}
若同一方向来自多个瞄准距离分支，要把各分支的截面相加。轴线附近 $\sin\theta=0$ 时，密度的行为由分子与分母共同决定，不能单看分母。

当入射通量在碰撞截面内均匀时，以已经发生碰撞为条件，方向概率为
\[
\dd P=\frac{\dd\sigma}{\sigma_{\rm tot}},\qquad
f_\Omega=\frac1{\sigma_{\rm tot}}\frac{\dd\sigma}{\dd\Omega}.
\]
截面以面积为单位，归一化后的概率密度则按指定的方向微元定义。

\par\medskip\noindent\begin{minipage}{\textwidth}
\begin{example}[固定光滑硬球的散射与束流推力]\label{mp1:ex:e8}
半径为 $r_A$、质量为 $m$ 的入射球以速率 $v$ 撞向固定靶球，靶球半径为 $r_B$。碰撞光滑且完全弹性，忽略重力，令 $R=r_A+r_B$，均匀入射通量为 $j$。图\ref{mp1:fig:e8-q}定义瞄准距离 $b$ 和接触角 $\varphi$。求总碰撞截面、微分截面、碰撞后的极角密度、时间 $\tau$ 内整圈角带的计数，以及靶受到的平均轴向推力。
\end{example}

\begin{figure}[H]\centering

\begin{tikzpicture}[scale=1.05]
\draw[axis](-.3,0)--(5.2,0)node[right]{$z$};\draw[thick](3,0)circle(.95);\fill(3,0)circle(1.3pt)node[below]{$B$};\draw(1.7876,.7)circle(.45);\fill(1.7876,.7)circle(1.2pt)node[above left]{$A$};
\draw[axis,thick](-.7,.7)--(1.3,.7)node[midway,above]{$v$};\draw[guide](1.7876,.7)--(3,0);\draw[<->](-.8,0)--(-.8,.7)node[midway,left]{$b$};\draw[guide](1.7876,.7)--(3.4,.7);
\draw(2.1876,.7)arc(0:-30:.4);\node[right]at(2.3,.55){$\varphi$};

\draw[axis,thick](1.7876,.7)--(.7876,2.43205)node[above]{$\vct v'$};\draw(2.4876,.7)arc(0:120:.7);\node at(2.15,1.5){$\theta$};\node at(3,-1.35){靶球固定};
\end{tikzpicture}
\caption{例~\ref{mp1:ex:e8}题设：$R=r_A+r_B$ 是两球相切时的球心距，$b$ 是入射轨迹到固定靶球心的距离。}\label{mp1:fig:e8-q}
\end{figure}

\end{minipage}\par\medskip

\begin{solution}
两球相切时球心距为 $R$，能碰撞的轨迹满足 $0\le b\le R$，故总截面为 $\sigma_{\rm tot}=\pi R^2$。固定、光滑、完全弹性的碰撞使入射速度的法向分量反向，切向分量不变。图\ref{mp1:fig:e8-q}中的接触角因此满足
\[
b=R\sin\varphi,
\qquad \theta=\pi-2\varphi,
\qquad b=R\cos\frac\theta2.
\]
这在正碰 $b=0,\theta=\pi$ 与擦边 $b=R,\theta=0$ 两端都符合几何关系。图\ref{mp1:fig:e8-s}左图的入射环带映射到右图整圈方向带，得到
\begin{align*}
\frac{\dd\sigma}{\dd\Omega}
&=\frac{b}{\sin\theta}\left|\frac{\dd b}{\dd\theta}\right|
=\frac{R\cos(\theta/2)}{\sin\theta}
\frac R2\sin\frac\theta2
=\frac{R^2}{4},\qquad 0<\theta<\pi.
\end{align*}

\begin{figure}[H]\centering

\begin{tikzpicture}[scale=1.05]
\fill[even odd rule,black!12](0,0)circle(.99) (0,0)circle(.82);
\draw(0,0)circle(1.45);\draw(0,0)circle(.82);\draw(0,0)circle(.99);
\fill(0,0)circle(1pt);
\draw(0,0)--(-35:.82)node[midway,below left=1pt]{$b$};
\draw[|-|](45:.82)--(45:.99)node[midway,above right=3pt]{$|\dd b|$};
\node[below]at(0,-1.9){$\dd\sigma=2\pi b|\dd b|$};
\begin{scope}[shift={(5,0)}]
\pgfmathsetmacro{\sr}{1.45}
\pgfmathsetmacro{\tilt}{10}
\pgfmathsetmacro{\ta}{2*acos(.99/\sr)}
\pgfmathsetmacro{\tb}{2*acos(.82/\sr)}
\pgfmathsetmacro{\ca}{\sr*cos(\ta)*cos(\tilt)}
\pgfmathsetmacro{\cb}{\sr*cos(\tb)*cos(\tilt)}
\pgfmathsetmacro{\ra}{\sr*sin(\ta)}
\pgfmathsetmacro{\rb}{\sr*sin(\tb)}
\pgfmathsetmacro{\ea}{\ra*sin(\tilt)}
\pgfmathsetmacro{\eb}{\rb*sin(\tilt)}
\fill[black!12]plot[domain=-90:90,samples=60]({\ca+\ea*cos(\x)},{\ra*sin(\x)})
--plot[domain=90:-90,samples=60]({\cb+\eb*cos(\x)},{\rb*sin(\x)})--cycle;
\draw(0,0)circle(\sr);
\draw[guide](\ca,0)ellipse[x radius=\ea,y radius=\ra];
\draw[guide](\cb,0)ellipse[x radius=\eb,y radius=\rb];
\draw(\ca,-\ra)arc[start angle=-90,end angle=90,x radius=\ea,y radius=\ra];
\draw(\cb,-\rb)arc[start angle=-90,end angle=90,x radius=\eb,y radius=\rb];
\draw[axis](-1.75,0)--(1.9,0)node[right]{$z$};
\draw(0,0)--(\ca,\ra);
\draw[guide](0,0)--(\cb,\rb);
\pgfmathsetmacro{\aa}{atan2(\ra,\ca)}
\pgfmathsetmacro{\ab}{atan2(\rb,\cb)}
\draw[domain=0:\ta,samples=35]plot({.48*cos(\x)*cos(\tilt)},{.48*sin(\x)});
\node at(.36,.42){$\theta$};
\draw[domain=\ta:\tb,samples=20]plot({.8*cos(\x)*cos(\tilt)},{.8*sin(\x)});
\node[fill=white,inner sep=1pt]at(-.19,.98){$\dd\theta$};
\node[below]at(0,-1.9){$\dd\Omega=2\pi\sin\theta\dd\theta$};
\end{scope}
\end{tikzpicture}
\caption{例~\ref{mp1:ex:e8}解图：左图中入射环带的截面积是 $2\pi b\dd b$；右图把它映射为散射方向的整圈角带。}\label{mp1:fig:e8-s}
\end{figure}

以发生碰撞为条件，每单位立体角密度为 $f_\Omega=1/(4\pi)$，而按极角的密度为
\[
p_\theta(\theta)=2\pi\sin\theta f_\Omega=\frac12\sin\theta.
\]
时间 $\tau$ 内，整圈角带 $[\theta,\theta+\dd\theta]$ 的计数为
\[
\dd N=j\tau\frac{R^2}{4}\,2\pi\sin\theta\dd\theta
=\frac{\pi j\tau R^2}{2}\sin\theta\dd\theta.
\]
每个粒子的出射速率仍为 $v$，交给靶的轴向动量为 $mv(1-\cos\theta)$。横向动量因轴对称而相互抵消，所以平均推力为
\begin{align*}
F_z&=jmv\int(1-\cos\theta)\dd\sigma\\
&=jmv\frac{R^2}{4}\,2\pi\int_0^\pi(1-\cos\theta)\sin\theta\dd\theta\\
&=\frac{\pi jmvR^2}{2}\left[-\cos\theta+\frac12\cos^2\theta\right]_0^\pi\\
&=\frac{\pi jmvR^2}{2}\left(\frac32+\frac12\right)\\
&=jmv\pi R^2.
\end{align*}
固定靶的支承装置须提供大小相同、方向相反的力。推力中每秒碰撞数 $j\pi R^2$ 与平均单次轴向动量转移 $mv$ 相乘，单位为力。
\end{solution}

\subsection{探测器的形状决定所收集的方向微元}
上例计数收集了同一极角处的全部方位角。如果探测器只是一条窄带，收集区域就小得多。设探测器位于以靶心为中心、半径为 $D$ 的远处球面上，沿包含入射轴的大圆布置，宽度为 $w\ll D$。

图\ref{mp1:fig:belt}中，远离两个极点的一支弧段在极角区间 $\dd\theta$ 内的面积约为 $wD\dd\theta$，所以
\[
\dd\Omega_{\rm one}\simeq\frac wD\dd\theta.
\]
这一区域的立体角不含整圈环带的因子 $2\pi\sin\theta$。若对同一极角把大圆两侧的两支弧段都计入，结果乘2。

\begin{figure}[H]\centering

\begin{tikzpicture}[scale=1.05]
\draw(0,0)circle(1.6);
\draw[axis](-1.9,0)--(2,0)node[right]{$z$};
\fill(0,0)circle(1.3pt)node[below]{$B$};
\draw(0,0)--(60:1.6)node[midway,above left]{$D$};
\draw[guide](0,0)--(74:1.6);
\draw[line width=2pt](60:1.6)arc(60:74:1.6);
\draw(.55,0)arc(0:60:.55);\node at(.68,.32){$\theta$};
\draw(60:1.02)arc(60:74:1.02);
\node[above left]at(.43,1.02){$\dd\theta$};
\node[below]at(0,-1.9){大圆所在平面内的方向};
\begin{scope}[shift={(4.4,.4)}]
\fill[black!10](0,0)rectangle(3,.6);\draw(0,0)rectangle(3,.6);
\draw[<->](0,-.5)--(3,-.5)node[midway,fill=white]{$D\dd\theta$};
\draw[<->](3.4,0)--(3.4,.6)node[midway,right]{$w$};
\node at(1.5,1.2){单支局部微元};
\node at(1.5,-1.2){$\dd\Omega\simeq(w/D)\dd\theta$};
\end{scope}
\end{tikzpicture}
\caption{大圆窄带探测器的一支弧段。左图为大圆所在平面内的方向，粗弧长为 $D\dd\theta$；横向宽度 $w$ 垂直于该平面，在右侧局部展开图中标出。展开近似用于远离极点的位置。}\label{mp1:fig:belt}
\end{figure}

对固定硬球，一支窄带弧段的计数因此为
\[
\dd N_{\rm one}\simeq j\tau\frac{R^2}{4}\frac wD\dd\theta.
\]
按 $\theta$ 看它近似为常量，而全方位角环带的计数正比于 $\sin\theta$。两者描述同一个散射规律，差别来自探测器覆盖的面积。该窄带展开式适用于 $\sin\theta\gg w/D$ 的位置；极点附近应对真实覆盖区域积分，不能把两支近似重复相加。

\subsection{等质量自由靶的角分布与能量分布}
靶球可以运动时，碰撞运动学会改变 $b$ 与 $\theta$ 的对应。先求单次碰撞的速度，再变换微元，顺序与固定靶相同。以下每次碰撞都使用一颗初始静止的同质量靶球；统计结果针对同一初始条件下重复的独立碰撞。

\par\medskip\noindent\begin{minipage}{\textwidth}
\begin{example}[等质量自由靶与出射能量]\label{mp1:ex:e9}
将上例靶球改为可自由运动，且与入射球等质量。每次碰撞前靶球均静止，碰撞仍光滑且完全弹性；瞄准距离在半径 $R=r_A+r_B$ 的圆截面内按面积均匀分布。见图\ref{mp1:fig:e9-q}。求入射球的出射角密度、微分截面，以及出射动能占初动能的比例 $\varepsilon$ 的分布和平均值。
\end{example}

\begin{figure}[H]\centering

\begin{tikzpicture}[scale=1.05]
\draw[axis](-.3,0)--(5.2,0)node[right]{$z$};\draw[thick](3,0)circle(.95);\fill(3,0)circle(1.3pt)node[below]{$B$};\draw(1.7876,.7)circle(.45);\fill(1.7876,.7)circle(1.2pt)node[above left]{$A$};
\draw[axis,thick](-.7,.7)--(1.3,.7)node[midway,above]{$v$};\draw[guide](1.7876,.7)--(3,0);\draw[<->](-.8,0)--(-.8,.7)node[midway,left]{$b$};\draw[guide](1.7876,.7)--(3.4,.7);
\draw(2.1876,.7)arc(0:-30:.4);\node[right]at(2.3,.55){$\varphi$};

\draw[axis,thick](1.7876,.7)--(2.2876,1.566)node[above]{$\vct v_A$};\draw[axis,thick](3,0)--(4.5,-.866)node[right]{$\vct v_B$};\draw(2.3876,.7)arc(0:60:.6);\node at(2.45,1.1){$\theta$};\node at(3,-1.55){$m_A=m_B=m$，$B$ 可自由运动};
\end{tikzpicture}
\caption{例~\ref{mp1:ex:e9}题设：等质量自由靶的碰撞。小球 $A$ 的出射速度沿接触切线，$B$ 的速度沿两球心连线。}\label{mp1:fig:e9-q}
\end{figure}

\end{minipage}\par\medskip

\begin{solution}
取接触法线为 $\vct n$，切线为 $\vct t$。入射速度分解为 $v\cos\varphi\,\vct n+v\sin\varphi\,\vct t$。光滑碰撞不改变切向速度。记入射法向速度为 $u=v\cos\varphi$，两球碰后法向速度为 $u_A,u_B$，由动量和动能守恒，
\[
u_A+u_B=u,\qquad u_A^2+u_B^2=u^2.
\]
两式给出 $u_Au_B=0$。排除没有发生碰撞的原状态，取 $u_A=0,u_B=u$，即两球交换法向速度。因此
\[
\vct v_A=v\sin\varphi\,\vct t,\qquad
\vct v_B=v\cos\varphi\,\vct n.
\]
两出射速度垂直，矢量和等于入射速度。图\ref{mp1:fig:e9-s}左图的直角三角形给出
\[
\theta=\frac\pi2-\varphi,\qquad
v_A=v\cos\theta,\qquad b=R\cos\theta,
\quad 0<\theta<\frac\pi2.
\]
因此
\[
\frac{\dd\sigma}{\dd\Omega}
=\frac{b}{\sin\theta}\left|\frac{\dd b}{\dd\theta}\right|
=\frac{R\cos\theta}{\sin\theta}R\sin\theta
=R^2\cos\theta.
\]
总截面仍为 $\pi R^2$。每单位立体角密度和极角密度分别为
\[
f_\Omega=\frac{\cos\theta}{\pi},\qquad
p_\theta=2\sin\theta\cos\theta=\sin2\theta.
\]

\begin{figure}[H]\centering

\begin{tikzpicture}[scale=1.1]
\draw[axis,thick](0,0)--(3,0)node[midway,below]{$\vct v$};\draw[axis](0,0)--(.75,1.299)node[midway,left]{$\vct v_A$};\draw[axis](.75,1.299)--(3,0)node[midway,above right]{$\vct v_B$};\draw(.6,0)arc(0:60:.6);\node at(.65,.4){$\theta$};
\begin{scope}[shift={(5,0)},x=2.6cm,y=1.7cm]
\draw[axis](0,0)--(1.25,0)node[right]{$\varepsilon$};\draw[axis](0,0)--(0,1.4)node[above]{$p_\varepsilon$};\draw[curve](0,1)--(1,1);\draw[guide](1,0)--(1,1);\node[left]at(0,1){$1$};\node[below]at(1,0){$1$};
\end{scope}
\end{tikzpicture}
\caption{例~\ref{mp1:ex:e9}解图：速度直角三角形给出 $v_A=v\cos\theta$。右图表示动能比 $\varepsilon=E_A/E_0$ 的均匀分布。}\label{mp1:fig:e9-s}
\end{figure}

记初动能为 $E_0=mv^2/2$。由 $v_A=v\cos\theta$，
\[
\varepsilon=\frac{E_A}{E_0}=\cos^2\theta=\frac{b^2}{R^2}.
\]
在入射截面中，$\dd P=2b\dd b/R^2$。而 $\dd\varepsilon=2b\dd b/R^2$，所以图\ref{mp1:fig:e9-s}右图中的密度为
\[
p_\varepsilon(\varepsilon)=1,\qquad 0<\varepsilon<1.
\]
其平均值为
\[
\avg\varepsilon=\int_0^1\varepsilon p_\varepsilon(\varepsilon)\dd\varepsilon
=\left[\frac{\varepsilon^2}{2}\right]_0^1=\frac12.
\]
平均有一半初动能留在入射球中，另一半转给靶球。正碰使入射球停下，此时它没有出射方向；这一单独事件在连续截面分布中概率为零。自由靶的角分布与前述余弦源模型相同，来源则是碰撞几何与速度交换。
\end{solution}












\par\medskip\noindent\begin{minipage}{\textwidth}
\begin{example}[窄环带探测器的固定靶与自由靶计数]\label{mp1:rev:detector}
半径$r_A$的入射球与半径$r_B$的靶球质量均为$m$，入射通量密度为$j$，碰撞光滑且完全弹性。令$a_c=r_A+r_B$。以靶心为中心，在包含入射轴的大圆上布置半径$D$、窄宽度$w$的环带探测器，$D\gg a_c,w$，见图\ref{mp1:fig:rev:detector-q}。令$\chi\in(-\pi,\pi]$为沿该大圆从入射正向量起的有符号角。求时间$T$内每段$\dd\chi$的入射球计数：(1)靶固定；(2)靶自由，每次碰撞前恢复到原位并静止。
\end{example}

\begin{figure}[H]\centering
\begin{tikzpicture}[>=Stealth,line cap=round,line join=round]
\draw(0,0)circle(1.7);\draw[axis](-2.3,0)--(2.3,0)node[right]{入射轴};\fill(0,0)circle(1.5pt)node[below]{$B$};\draw[axis](0,0)--(1.3,1.1);\draw(.55,0)arc(0:40.24:.55);\node at(.8,.25){$\chi$};\node[left]at(-1.7,.7){$D$};\node[right,align=left]at(2.7,1){环带宽度$w$\\垂直于图面};
\end{tikzpicture}
\caption{环带沿包含入射轴的整条大圆布置；$\chi$逐段标记位置}\label{mp1:fig:rev:detector-q}
\end{figure}
\end{minipage}\par\medskip

\begin{solution}
用球面上相对探测器平面的纬度$\psi$和大圆角$\chi$描述方向，立体角微元为$\dd\Omega=\cos\psi\dd\psi\dd\chi$。窄带满足$|\psi|\lesssim w/(2D)$，积分后
\[
 \dd\Omega\simeq\int_{-w/(2D)}^{w/(2D)}\cos\psi\dd\psi\dd\chi
 =[\sin\psi]_{-w/(2D)}^{w/(2D)}\dd\chi
 =2\sin\frac w{2D}\dd\chi
 \simeq\frac wD\dd\chi.
\]
这里的$\chi$按整圈的有符号位置逐段计数，每个方向只出现一次，无须再乘两支弧段的因子2。

例\ref{mp1:ex:e8}的固定靶散射结果给出$\dd\sigma/\dd\Omega=a_c^2/4$，因此
\[
 \boxed{\frac{\dd N_1}{\dd\chi}\simeq\frac{jTa_c^2w}{4D},
 \quad -\pi<\chi\le\pi.}
\]
由例\ref{mp1:ex:e9}的自由等质量靶结果，$\dd\sigma/\dd\Omega=a_c^2\cos\theta$，且入射球只向前半空间出射。窄带中$\cos\theta\simeq\cos\chi$，故
\[
 \boxed{\frac{\dd N_2}{\dd\chi}\simeq
 \begin{cases}jTa_c^2(w/D)\cos\chi,&|\chi|<\pi/2,\\0,&|\chi|\ge\pi/2.\end{cases}}
\]
图\ref{mp1:fig:rev:detector-s}将两种计数除以$jTa_c^2w/D$后画在同一纵轴上。固定靶整圈近似均匀，自由靶只在前半圈有计数。总计数须分别对各自的支持区间积分：
\[
\begin{aligned}
N_1&\simeq\frac{jTa_c^2w}{4D}\int_{-\pi}^{\pi}\dd\chi
=\frac{jTa_c^2w}{4D}[\chi]_{-\pi}^{\pi}
=\frac{\pi jTa_c^2w}{2D},\\
N_2&\simeq\frac{jTa_c^2w}{D}\int_{-\pi/2}^{\pi/2}\cos\chi\dd\chi
=\frac{jTa_c^2w}{D}[\sin\chi]_{-\pi/2}^{\pi/2}
=\frac{2jTa_c^2w}{D}.
\end{aligned}
\]
它们与整圈极角带的$\sin\theta$、$\sin2\theta$分布不同，原因是探测器覆盖的方向微元不同。上述式子取窄带与远场的最低阶，有限宽度时可在$\psi$方向保留积分。

\begin{figure}[H]\centering
\begin{tikzpicture}[>=Stealth,line cap=round,line join=round]
\draw[axis](-3.5,0)--(3.5,0)node[right]{$\chi$};\draw[axis](0,0)--(0,2.5)node[above]{归一化计数};\draw[curve](-3.14,.5)--(3.14,.5);\draw[curve,domain=-1.5708:1.5708,samples=90]plot(\x,{2*cos(\x r)});\draw[guide](-1.5708,0)--(-1.5708,2.15);\draw[guide](1.5708,0)--(1.5708,2.15);\node[below]at(-1.5708,0){$-\pi/2$};\node[below]at(1.5708,0){$\pi/2$};\node[left]at(0,.5){$1/4$};\node[left]at(0,2){$1$};\node at(2.4,.85){固定靶};\node at(1.3,1.65){自由靶};
\end{tikzpicture}
\caption{固定靶计数为常量，自由靶计数在前半圈按余弦变化}\label{mp1:fig:rev:detector-s}
\end{figure}

\end{solution}

\section{习题与解答}

\par\medskip\noindent\begin{minipage}{\textwidth}
\begin{exercise}[速率分布与动能分布]\label{mp1:pr:p1}
某组同质量粒子的速率密度为 $f_v(v)=Cv^2$，$0<v<v_0$，区域外为零。求 $C$、动能 $E=mv^2/2$ 的密度及平均动能。图\ref{mp1:fig:p1-q}表示所给速率分布。
\end{exercise}

\begin{figure}[H]\centering

\begin{tikzpicture}[scale=1.05]
\fill[black!12](1.5,0)--plot[domain=1.5:1.8](\x,{.2*(\x)^2})--(1.8,0)--cycle;
\draw[axis](0,0)--(4,0)node[right]{$v$};
\draw[axis](0,0)--(0,2.25)node[above]{$f_v$};
\draw[curve,domain=0:3]plot(\x,{.2*(\x)^2});
\draw[guide](3,0)--(3,1.8)--(0,1.8);
\node[below]at(3,0){$v_0$};
\node[left]at(0,1.8){$Cv_0^2$};
\draw[|-|](1.5,-.28)--(1.8,-.28)node[midway,below]{$\dd v$};
\node[anchor=west]at(4.4,.8){$E=mv^2/2$};
\end{tikzpicture}
\caption{习题~\ref{mp1:pr:p1}题设：阴影面积为速率区间内的概率。由 $\dd E=mv\dd v$，等宽速率区间所对应的动能区间随 $v$ 而变。}\label{mp1:fig:p1-q}
\end{figure}

\end{minipage}\par\medskip

\begin{solution}
由总概率为1，
\[
1=C\int_0^{v_0}v^2\dd v=C\left[\frac{v^3}{3}\right]_0^{v_0}
=\frac{Cv_0^3}{3},
\]
得 $C=3/v_0^3$。记 $E_0=mv_0^2/2$，动能区域为 $0<E<E_0$。由 $\dd E=mv\dd v$，
\[
f_E(E)=f_v(v)\frac1{mv}=\frac{3v}{mv_0^3}
=\frac3{2E_0}\sqrt{\frac E{E_0}}.
\]
图\ref{mp1:fig:p1-s}给出这个平方根分布。曲线下总面积为1，密度的单位为能量的倒数。

\begin{figure}[H]\centering

\begin{tikzpicture}[x=4cm,y=1.6cm]
\fill[black!15](.5,0)--plot[domain=.5:.58,samples=12](\x,{1.5*sqrt(\x)})--(.58,0)--cycle;
\draw[axis](0,0)--(1.2,0)node[right]{$E/E_0$};
\draw[axis](0,0)--(0,1.9)node[above]{$E_0f_E$};
\draw[curve,domain=0:1,samples=80]plot(\x,{1.5*sqrt(\x)});
\draw[guide](1,0)--(1,1.5);
\draw[|-|](.5,-.13)--(.58,-.13)node[midway,below]{$\dd E/E_0$};
\node[below]at(1,0){$1$};\node[left]at(0,1.5){$3/2$};
\end{tikzpicture}
\caption{习题~\ref{mp1:pr:p1}解图：动能密度随 $E^{1/2}$ 增长。阴影为一个能量区间的概率微元，求平均动能时再乘该区间的能量。}\label{mp1:fig:p1-s}
\end{figure}

求平均值时必须再乘 $E$。令 $u=E/E_0$，则 $\dd E=E_0\dd u$、上下限变为0和1：
\[
\avg E=\int_0^{E_0}E f_E(E)\dd E
=\frac32 E_0\int_0^1u^{3/2}\dd u
=\frac32E_0\left[\frac25u^{5/2}\right]_0^1
=\frac35E_0.
\]
由累计概率也能求出密度。在$0<E<E_0$内，
\[
 \Prob(mv^2/2\le E)=\Prob\left(v\le v_0\sqrt{E/E_0}\right)
 =\left(\frac E{E_0}\right)^{3/2}.
\]
对 $E$ 求导，$\dd(E/E_0)^{3/2}/\dd E=(3/2)(E/E_0)^{1/2}/E_0$，正是 $f_E$。图\ref{mp1:fig:p1-s}中的阴影宽度为$\dd E$，面积等于速率图中相应$\dd v$条带的面积。由$v$改用$E$作横坐标时，纵坐标密度也须随之变换，才能保持这两个概率相等。
\end{solution}

\par\medskip\noindent\begin{minipage}{\textwidth}
\begin{exercise}[均匀相位对应的位移分布]\label{mp1:pr:p2}
简谐振动 $x=A\cos\phi$，相位 $\phi$ 在 $[0,2\pi)$ 上均匀分布。求位移密度、$\Prob(|x|\le A/2)$ 和 $\avg{x^2}$。图\ref{mp1:fig:p2-q}用圆上的投影表示位移。
\end{exercise}

\begin{figure}[H]\centering

\begin{tikzpicture}[scale=1.1]
\draw(0,0)circle(1.45);\draw[axis](-1.8,0)--(2,0)node[right]{$x/A$};\draw[axis](0,-1.7)--(0,1.8);\draw(0,0)--(.725,1.256);\draw[guide](.725,-1.256)--(.725,1.256);\fill(.725,1.256)circle(1.3pt);\fill(.725,-1.256)circle(1.3pt);\draw(.5,0)arc(0:60:.5);\node at(.6,.36){$\phi$};\node[below]at(.725,0){$x/A$};
\node[below left]at(-1.45,0){$-1$};\node[below right]at(1.45,0){$1$};
\end{tikzpicture}
\caption{习题~\ref{mp1:pr:p2}题设：相位均匀分布对应单位圆上的均匀取样。图中两点有相同的横坐标 $x/A$，给出同一位移的两支相位。}\label{mp1:fig:p2-q}
\end{figure}

\end{minipage}\par\medskip

\begin{solution}
对 $-A<x<A$，上、下半圆各有一个相位对应同一横坐标。每支都有
\[
\left|\frac{\dd\phi}{\dd x}\right|
=\frac1{A|\sin\phi|}=\frac1{\sqrt{A^2-x^2}}.
\]
两支概率相加，得 $f_x(x)=1/(\pi\sqrt{A^2-x^2})$。图\ref{mp1:fig:p2-s}中两端较高的密度对应振子在转折点附近运动较慢。

\begin{figure}[H]\centering

\begin{tikzpicture}[x=2.5cm,y=1.5cm]
\fill[black!12](-.5,0)--plot[domain=-.5:.5,samples=80](\x,{1/(pi*sqrt(1-(\x)^2))})--(.5,0)--cycle;
\draw[axis](-1.25,0)--(1.3,0)node[right]{$x/A$};
\draw[axis](0,0)--(0,2.2)node[above]{$Af_x$};
\draw[curve,<-,>=Stealth,domain=-.985:0,samples=90]plot(\x,{1/(pi*sqrt(1-(\x)^2))});
\draw[curve,->,>=Stealth,domain=0:.985,samples=90]plot(\x,{1/(pi*sqrt(1-(\x)^2))});
\draw[guide](-1,0)--(-1,2) (1,0)--(1,2);
\draw[guide](-.5,0)--(-.5,.368) (.5,0)--(.5,.368);
\node[below]at(-1,0){$-1$};\node[below]at(1,0){$1$};
\node[below]at(-.5,0){$-1/2$};\node[below]at(.5,0){$1/2$};
\end{tikzpicture}
\caption{习题~\ref{mp1:pr:p2}解图：阴影面积为 $\Prob(|x|\le A/2)=1/3$。曲线在 $x=\pm A$ 附近向上发散，端点处没有集中概率。}\label{mp1:fig:p2-s}
\end{figure}

由原函数 $\arcsin(x/A)$，
\[
\Prob(|x|\le A/2)=\frac1\pi\int_{-A/2}^{A/2}\frac{\dd x}{\sqrt{A^2-x^2}}
=\frac1\pi\left[\arcsin\frac xA\right]_{-A/2}^{A/2}
=\frac13.
\]
均方位移在相位变量中更容易积分：
\[
\avg{x^2}=\frac{A^2}{2\pi}\int_0^{2\pi}\frac{1+\cos2\phi}{2}\dd\phi
=\frac{A^2}{2\pi}\left[\frac\phi2+\frac{\sin2\phi}{4}\right]_0^{2\pi}
=\frac{A^2}{2}.
\]
相位随时间匀速增加，因此这个密度也表示长时间内各位移区间所占的时间比例。
对于一次确定的简谐振动，速度大小为$|v|=\omega\sqrt{A^2-x^2}$。一个周期$T=2\pi/\omega$内，内部小区间$[x,x+\dd x]$被经过两次，所占时间为$2\dd x/|v|$，所以
\[
 f_x(x)\dd x=\frac{2}{T}\frac{\dd x}{\omega\sqrt{A^2-x^2}}
 =\frac{\dd x}{\pi\sqrt{A^2-x^2}}.
\]
图\ref{mp1:fig:p2-q}中的两个原像分别对应振子的往程和返程。端点附近，密度正比于到端点距离的$-1/2$次方；该奇性可积，单个端点的概率仍为零。
\end{solution}

\par\medskip\noindent\begin{minipage}{\textwidth}
\begin{exercise}[实心圆锥的质量矩]\label{mp1:pr:p3}
均匀实心圆锥的底面半径为 $R$，高为 $H$，体密度为 $\rho$。按图\ref{mp1:fig:p3-q}从底面向上取坐标 $z$。求质量、重心高度和绕对称轴的转动惯量。
\end{exercise}

\begin{figure}[H]\centering

\begin{tikzpicture}[scale=1.2]
\draw[axis](0,-.3)--(0,3.1)node[above]{$z$};\draw(-1.6,0)--(0,2.6)--(1.6,0);\draw(-1.6,0)arc(180:360:1.6 and .3);\draw[guide](-1.6,0)arc(180:0:1.6 and .3);\draw[guide](0,0)--(1.6,0)node[midway,above]{$R$};\node[left]at(0,2.6){$H$};\node[right]at(1.1,1.3){$\rho$};
\end{tikzpicture}
\caption{习题~\ref{mp1:pr:p3}题设：圆锥的坐标 $z$ 从底面向顶点量起，半径随高度线性减小。}\label{mp1:fig:p3-q}
\end{figure}

\end{minipage}\par\medskip

\begin{solution}
图\ref{mp1:fig:p3-s}中的薄圆盘半径为 $r(z)=R(1-z/H)$，质量为 $\dd m=\rho\pi r^2\dd z$。设 $u=z/H$，$\dd z=H\dd u$，两端 $z=0,H$ 对应 $u=0,1$，则
\[
M=\rho\pi R^2H\int_0^1(1-u)^2\dd u
=\rho\pi R^2H\left[-\frac{(1-u)^3}{3}\right]_0^1
=\frac13\rho\pi R^2H.
\]
高度矩为
\[
Mz_C=\rho\pi R^2H^2\int_0^1u(1-u)^2\dd u
=\rho\pi R^2H^2\left[\frac{u^2}{2}-\frac{2u^3}{3}+\frac{u^4}{4}\right]_0^1
=\frac1{12}\rho\pi R^2H^2,
\]
故 $z_C=H/4$，图中 $C$ 接近质量较集中的底面。

\begin{figure}[H]\centering

\begin{tikzpicture}[scale=1.2]
\pgfmathsetmacro{\zr}{1.1}
\pgfmathsetmacro{\zt}{1.32}
\pgfmathsetmacro{\rr}{1.6*(1-\zr/2.6)}
\pgfmathsetmacro{\rt}{1.6*(1-\zt/2.6)}
\draw(-1.6,0)--(0,2.6)--(1.6,0);
\draw(-1.6,0)arc(180:360:1.6 and .3);
\draw[guide](-1.6,0)arc(180:0:1.6 and .3);
\draw[axis](0,-.2)--(0,2.9)node[above]{$z$};
\draw(-\rr,\zr)arc[start angle=180,end angle=360,x radius=\rr,y radius=.18];
\draw[guide](-\rr,\zr)arc[start angle=180,end angle=0,x radius=\rr,y radius=.18];
\draw(-\rt,\zt)arc[start angle=180,end angle=360,x radius=\rt,y radius=.16];
\draw[guide](-\rt,\zt)arc[start angle=180,end angle=0,x radius=\rt,y radius=.16];
\draw(-\rr,\zr)--(-\rt,\zt) (\rr,\zr)--(\rt,\zt);
\draw[guide](0,\zr)--(\rr,\zr)node[midway,below=2pt]{$r(z)$};
\draw[guide](\rr,\zr)--(1.55,\zr) (\rt,\zt)--(1.55,\zt);
\draw[|-|](1.45,\zr)--(1.45,\zt)node[midway,right=3pt]{$\dd z$};
\fill(0,.65)circle(1.5pt)node[left=3pt]{$C:\ H/4$};
\end{tikzpicture}
\caption{习题~\ref{mp1:pr:p3}解图：水平薄圆盘的半径为 $R(1-z/H)$。质心和轴向惯量分别用同一质量微元加权。}\label{mp1:fig:p3-s}
\end{figure}

同一薄盘绕轴的惯量可从同心环带求得：
\[
\dd I_z=\int_0^{r(z)}s^2\rho\dd z\,2\pi s\dd s
=2\pi\rho\dd z\left[\frac{s^4}{4}\right]_0^{r(z)}
=\frac12\rho\pi r^4(z)\dd z.
\]
沿高度相加，
\[
I_z=\frac{\rho\pi R^4H}{2}\int_0^1(1-u)^4\dd u
=\frac{\rho\pi R^4H}{2}\left[-\frac{(1-u)^5}{5}\right]_0^1
=\frac1{10}\rho\pi R^4H=\frac3{10}MR^2.
\]
质量、重心与惯量使用同一几何分割，只是相应的积分权重依次为 $1,z,s^2$。
\end{solution}

\par\medskip\noindent\begin{minipage}{\textwidth}
\begin{exercise}[均匀圆盘的横坐标分布]\label{mp1:pr:p4}
在半径 $R$ 的圆盘内按面积均匀取一点，求其横坐标 $X$ 的密度及 $\Prob(X>R/2)$。图\ref{mp1:fig:p4-q}中的竖直条带含有相近的横坐标。
\end{exercise}

\begin{figure}[H]\centering

\begin{tikzpicture}[scale=1.1]
\pgfmathsetmacro{\bandleft}{.85}
\pgfmathsetmacro{\bandright}{1.03}
\pgfmathsetmacro{\yleft}{sqrt(1.7^2-\bandleft^2)}
\pgfmathsetmacro{\yright}{sqrt(1.7^2-\bandright^2)}
\begin{scope}\clip(0,0)circle(1.7);
\fill[black!12](\bandleft,-1.7)rectangle(\bandright,1.7);
\end{scope}
\draw(0,0)circle(1.7);
\draw[axis](-2,0)--(2.2,0)node[right]{$x$};
\draw[axis](0,-1.9)--(0,2)node[above]{$y$};
\draw[thick](\bandleft,-\yleft)--(\bandleft,\yleft);
\draw[guide](\bandright,-\yright)--(\bandright,1.85);
\draw[guide](\bandleft,\yleft)--(\bandleft,1.85);
\draw[|-|](\bandleft,1.8)--(\bandright,1.8)node[midway,above=2pt]{$\dd x$};
\node[below,fill=white,inner sep=1pt]at(\bandleft,-.08){$x$};
\node[below right]at(1.7,0){$R$};
\end{tikzpicture}
\caption{习题~\ref{mp1:pr:p4}题设：均匀圆盘上的点按面积取样。竖直弦带内所有点有相近的横坐标。}\label{mp1:fig:p4-q}
\end{figure}

\end{minipage}\par\medskip

\begin{solution}
横坐标为 $x$ 的弦长为 $2\sqrt{R^2-x^2}$，宽为 $\dd x$ 的条带面积就是弦长乘宽。圆盘上面积密度为 $1/(\pi R^2)$，故
\[
f_X(x)=\frac{2\sqrt{R^2-x^2}}{\pi R^2},\qquad -R<x<R.
\]
图\ref{mp1:fig:p4-s}的曲线在中心最高，在两边降为零，正好对应竖直弦长的变化。

\begin{figure}[H]\centering

\begin{tikzpicture}[x=2.4cm,y=2.6cm]
\fill[black!12](.5,0)--plot[domain=.5:1,samples=80](\x,{2*sqrt(1-(\x)^2)/pi})--(1,0)--cycle;
\draw[axis](-1.25,0)--(1.3,0)node[right]{$x/R$};\draw[axis](0,0)--(0,.9)node[above]{$R f_X$};\draw[curve,domain=-1:1,samples=160]plot(\x,{2*sqrt(1-(\x)^2)/pi});\node[below]at(-1,0){$-1$};\node[below]at(1,0){$1$};\node[below]at(.5,0){$1/2$};
\end{tikzpicture}
\caption{习题~\ref{mp1:pr:p4}解图：横坐标密度正比于弦长。右侧阴影的面积给出 $X>R/2$ 的概率。}\label{mp1:fig:p4-s}
\end{figure}

令 $x=R\sin t$，则 $\dd x=R\cos t\dd t$、$\sqrt{R^2-x^2}=R\cos t$；$x=R/2,R$ 分别对应 $t=\pi/6,\pi/2$。所求概率为
\begin{align*}
\Prob(X>R/2)&=\frac2\pi\int_{\pi/6}^{\pi/2}\cos^2t\dd t\\
&=\frac1\pi\left[t+\frac12\sin2t\right]_{\pi/6}^{\pi/2}
=\frac13-\frac{\sqrt3}{4\pi}.
\end{align*}
它等于图\ref{mp1:fig:p4-s}右侧阴影面积，也等于圆盘右侧弓形面积占全圆的比例。均匀面积分布在横坐标上并不均匀。
\end{solution}

\par\medskip\noindent\begin{minipage}{\textwidth}
\begin{exercise}[极角均匀是否意味着空间均匀]\label{mp1:pr:p5}
一种空间方向取样方式令 $\theta$ 在 $[0,\pi]$ 上均匀、$\phi$ 在 $[0,2\pi)$ 上均匀，二者独立。求每单位立体角的密度、半顶角 $\alpha$ 的正向圆锥概率及 $\avg{\cos^2\theta}$，与各向同性分布比较。见图\ref{mp1:fig:p5-q}。
\end{exercise}

\begin{figure}[H]\centering

\begin{tikzpicture}[scale=1.1]
\pgfmathsetmacro{\sr}{1.6}
\pgfmathsetmacro{\tilt}{12}
\foreach\base in{30,75}{
  \pgfmathsetmacro{\ru}{\sr*sin(\base)}
  \pgfmathsetmacro{\rl}{\sr*sin(\base+8)}
  \pgfmathsetmacro{\cu}{\sr*cos(\base)*cos(\tilt)}
  \pgfmathsetmacro{\cl}{\sr*cos(\base+8)*cos(\tilt)}
  \pgfmathsetmacro{\eu}{\ru*sin(\tilt)}
  \pgfmathsetmacro{\el}{\rl*sin(\tilt)}
  \fill[black!12](-\ru,\cu)arc[start angle=180,end angle=360,x radius=\ru,y radius=\eu]--(\rl,\cl)arc[start angle=0,end angle=-180,x radius=\rl,y radius=\el]--cycle;
  \fill[black!8](-\ru,\cu)arc[start angle=180,end angle=0,x radius=\ru,y radius=\eu]--(\rl,\cl)arc[start angle=0,end angle=180,x radius=\rl,y radius=\el]--cycle;
  \draw(-\ru,\cu)arc[start angle=180,end angle=360,x radius=\ru,y radius=\eu];
  \draw[guide](-\ru,\cu)arc[start angle=180,end angle=0,x radius=\ru,y radius=\eu];
  \draw(-\rl,\cl)arc[start angle=180,end angle=360,x radius=\rl,y radius=\el];
  \draw[guide](-\rl,\cl)arc[start angle=180,end angle=0,x radius=\rl,y radius=\el];
}
\draw(0,0)circle(\sr);
\draw[axis](0,-1.9)--(0,2.05)node[above]{$z$};
\node[anchor=west]at(1.2,1.32){极区方向带};
\node[anchor=west]at(1.85,.2){赤道附近方向带};
\node[below]at(0,-2.05){两带的极角宽度均为 $\dd\theta$};
\end{tikzpicture}
\caption{习题~\ref{mp1:pr:p5}题设：本题按极角等间隔取样。两个方向带的极角宽度相同，覆盖的立体角不同。}\label{mp1:fig:p5-q}
\end{figure}

\end{minipage}\par\medskip

\begin{solution}
对两个角坐标的联合密度为 $1/(2\pi^2)$，即
\[
\dd P=\frac{\dd\theta\dd\phi}{2\pi^2}
=\frac1{2\pi^2\sin\theta}\dd\Omega.
\]
故 $f_\Omega=1/(2\pi^2\sin\theta)$，在极区较大。极点处的发散由立体角微元的 $\sin\theta$ 抵消，并无集中概率。

\begin{figure}[H]\centering

\begin{tikzpicture}[scale=1.15]
\draw[axis](0,-.2)--(0,2.7)node[above]{$z$};\draw(-1.2,2)arc(180:360:1.2 and .28);\draw[guide](-1.2,2)arc(180:0:1.2 and .28);\draw(0,0)--(-1.2,2) (0,0)--(1.2,2);\draw(0,.7)arc(90:59:.7);\node at(.28,.85){$\alpha$};\node[right]at(2,1.5){$P_1=\alpha/\pi$};\node[right]at(2,.7){$P_2=(1-\cos\alpha)/2$};
\end{tikzpicture}
\caption{习题~\ref{mp1:pr:p5}解图：取半顶角 $\alpha$ 的圆锥。按极角均匀与按立体角均匀，对这个圆锥赋予不同概率。}\label{mp1:fig:p5-s}
\end{figure}

图\ref{mp1:fig:p5-s}的圆锥对应 $0\le\theta\le\alpha$，两种取样的概率分别为
\[
\begin{aligned}
P_1&=\frac1\pi\int_0^\alpha\dd\theta
=\frac1\pi[\theta]_0^\alpha=\frac\alpha\pi,\\
P_2&=\frac12\int_0^\alpha\sin\theta\dd\theta
=\frac12[-\cos\theta]_0^\alpha=\frac{1-\cos\alpha}{2}.
\end{aligned}
\]
窄圆锥的两概率分别按 $\alpha$ 和 $\alpha^2$ 变化，因此差别在小角度处尤其明显。

本题的二阶平均值为
\[
\avg{\cos^2\theta}_1=\frac1\pi\int_0^\pi\cos^2\theta\dd\theta
=\frac1\pi\left[\frac\theta2+\frac{\sin2\theta}{4}\right]_0^\pi=\frac12.
\]
各向同性时，令 $u=\cos\theta$，$\dd u=-\sin\theta\dd\theta$；$\theta=0,\pi$ 对应 $u=1,-1$。倒转上下限后，有
\[
\avg{\cos^2\theta}_2=\frac12\int_0^\pi\cos^2\theta\sin\theta\dd\theta
=\frac12\int_{-1}^1u^2\dd u
=\frac12\left[\frac{u^3}{3}\right]_{-1}^1=\frac13.
\]
两种取样都具有绕轴对称性，但只有后一种在任意旋转下保持方向分布不变。
\end{solution}

\par\medskip\noindent\begin{minipage}{\textwidth}
\begin{exercise}[等高球带与质量矩]\label{mp1:pr:p6}
均匀薄球壳半径为 $R$，面密度为 $\sigma$。保留 $z_1\le z\le z_2$ 的球带，$-R\le z_1<z_2\le R$，见图\ref{mp1:fig:p6-q}。求球带面积、质量、重心高度和绕 $z$ 轴的转动惯量。
\end{exercise}

\begin{figure}[H]\centering

\begin{tikzpicture}[scale=1.1]
\pgfmathsetmacro{\sr}{1.7}
\pgfmathsetmacro{\tilt}{12}
\pgfmathsetmacro{\zlo}{.55}
\pgfmathsetmacro{\zhi}{1.15}
\pgfmathsetmacro{\rlo}{sqrt(\sr^2-\zlo^2)}
\pgfmathsetmacro{\rhi}{sqrt(\sr^2-\zhi^2)}
\pgfmathsetmacro{\clo}{\zlo*cos(\tilt)}
\pgfmathsetmacro{\chi}{\zhi*cos(\tilt)}
\pgfmathsetmacro{\elo}{\rlo*sin(\tilt)}
\pgfmathsetmacro{\ehi}{\rhi*sin(\tilt)}
\fill[black!12](-\rhi,\chi)arc[start angle=180,end angle=360,x radius=\rhi,y radius=\ehi]--(\rlo,\clo)arc[start angle=0,end angle=-180,x radius=\rlo,y radius=\elo]--cycle;
\fill[black!8](-\rhi,\chi)arc[start angle=180,end angle=0,x radius=\rhi,y radius=\ehi]--(\rlo,\clo)arc[start angle=0,end angle=180,x radius=\rlo,y radius=\elo]--cycle;
\draw(0,0)circle(\sr);
\draw(-\rlo,\clo)arc[start angle=180,end angle=360,x radius=\rlo,y radius=\elo];
\draw[guide](-\rlo,\clo)arc[start angle=180,end angle=0,x radius=\rlo,y radius=\elo];
\draw(-\rhi,\chi)arc[start angle=180,end angle=360,x radius=\rhi,y radius=\ehi];
\draw[guide](-\rhi,\chi)arc[start angle=180,end angle=0,x radius=\rhi,y radius=\ehi];
\draw[axis](0,-2)--(0,2.15)node[above]{$z$};
\draw[guide](0,\clo)--(2.05,\clo)node[right]{$z_1$};
\draw[guide](0,\chi)--(2.05,\chi)node[right]{$z_2$};
\draw[guide](0,0)--(1.7,0)node[midway,below]{$R$};
\end{tikzpicture}
\caption{习题~\ref{mp1:pr:p6}题设：半径 $R$ 的均匀薄球壳，保留 $z_1\le z\le z_2$ 的球带。}\label{mp1:fig:p6-q}
\end{figure}

\end{minipage}\par\medskip

\begin{solution}
球面上 $z=R\cos\theta$，$|\dd z|=R\sin\theta\dd\theta$。图\ref{mp1:fig:p6-s}中的整圈方向带面积为
\[
\dd S=2\pi R^2\sin\theta\dd\theta=2\pi R|\dd z|.
\]
因此等高的完整球带面积相同，
\[
S=2\pi R(z_2-z_1),\qquad M=2\pi\sigma R(z_2-z_1).
\]
对高度积分时，$\dd m=2\pi\sigma R\dd z$ 是均匀分布，故
\[
z_C=\frac{\int_{z_1}^{z_2}z\dd z}{z_2-z_1}
=\frac{[z^2/2]_{z_1}^{z_2}}{z_2-z_1}
=\frac{z_2^2-z_1^2}{2(z_2-z_1)}=\frac{z_1+z_2}{2}.
\]

\begin{figure}[H]\centering

\begin{tikzpicture}[scale=1.1]
\pgfmathsetmacro{\sr}{1.7}
\pgfmathsetmacro{\ta}{53}
\pgfmathsetmacro{\tb}{61}
\pgfmathsetmacro{\xp}{\sr*sin(\ta)}
\pgfmathsetmacro{\zp}{\sr*cos(\ta)}
\pgfmathsetmacro{\xq}{\sr*sin(\tb)}
\pgfmathsetmacro{\zq}{\sr*cos(\tb)}
\coordinate(P)at(\xp,\zp);\coordinate(Q)at(\xq,\zq);
\draw(0,0)circle(\sr);
\draw[axis](0,-1.9)--(0,2.1)node[above]{$z$};
\draw(0,0)--(P)node[midway,below right=1pt]{$R$};
\draw[guide](0,\zp)--(P)node[midway,above]{$r$};
\draw[line width=1.3pt](P)arc[start angle=37,end angle=29,radius=\sr];
\draw[guide](P)--(2.05,\zp) (Q)--(2.05,\zq);
\draw[|-|](1.98,\zp)--(1.98,\zq)node[midway,right=2pt]{$|\dd z|$};
\node[above right=1pt]at(1.47,1.03){$\dd s$};
\draw(0,.5)arc(90:37:.5);\node at(.27,.61){$\theta$};
\node[anchor=west]at(2.9,.6){$\dd S=2\pi r\dd s$};
\node[anchor=west]at(2.9,-.1){$\dd S=2\pi R|\dd z|$};
\end{tikzpicture}
\caption{习题~\ref{mp1:pr:p6}解图：球壳的子午截面。粗弧长为 $\dd s$，竖直高度差为 $|\dd z|$；粗弧绕 $z$ 轴旋转一周形成球带微元。}\label{mp1:fig:p6-s}
\end{figure}

高度 $z$ 处的环带到轴距离平方为 $R^2-z^2$。于是
\begin{align*}
I_z&=2\pi\sigma R\int_{z_1}^{z_2}(R^2-z^2)\dd z\\
&=2\pi\sigma R\left[R^2z-\frac{z^3}{3}\right]_{z_1}^{z_2}\\
&=2\pi\sigma R(z_2-z_1)\left[R^2-\frac{z_2^2+z_1z_2+z_1^2}{3}\right]\\
&=M\left[R^2-\frac{z_1^2+z_1z_2+z_2^2}{3}\right].
\end{align*}
取整球 $z_1=-R,z_2=R$，即得 $I_z=2MR^2/3$。薄球带极限下，上式趋近环带惯量 $M(R^2-z_C^2)$。
\end{solution}

\par\medskip\noindent\begin{minipage}{\textwidth}
\begin{exercise}[无限长条靶]\label{mp1:pr:p7}
全空间各向同性点源位于原点，靶面为 $z=h>0$ 上的无限长条 $|x|\le a$，$y$ 不限，见图\ref{mp1:fig:p7-q}。求命中长条的无条件概率，并与二维随机射线命中有限线段的结果比较。
\end{exercise}

\begin{figure}[H]\centering

\begin{tikzpicture}[scale=1.05]
\fill[black!8](-.8,-1.8)rectangle(.8,1.8);
\draw[thick](-.8,-1.8)--(-.8,1.8) (.8,-1.8)--(.8,1.8);
\draw[axis](-2.3,0)--(2.4,0)node[right]{$x$};
\draw[axis](0,-2)--(0,2.2)node[above]{$y$};
\node[below]at(-.8,0){$-a$};\node[below]at(.8,0){$a$};
\draw[axis](-.45,1.35)--(-.45,1.85);
\draw[axis](-.45,-1.35)--(-.45,-1.85);
\node[anchor=west]at(3,.8){靶面 $z=h$};
\end{tikzpicture}
\caption{习题~\ref{mp1:pr:p7}题设：无限长条靶仅限制 $x$；$y$ 方向延伸到无穷远。}\label{mp1:fig:p7-q}
\end{figure}

\end{minipage}\par\medskip

\begin{solution}
在靶面上，面积密度为 $h/[4\pi(h^2+x^2+y^2)^{3/2}]$。固定 $x$，令 $A=h^2+x^2>0$，先对 $y$ 积分。作换元 $y=\sqrt A\tan t$，其中 $-\pi/2<t<\pi/2$，有
\begin{align*}
\frac{\dd y}{(A+y^2)^{3/2}}&=\frac{\cos t}{A}\dd t,\\
\int_{-\infty}^{\infty}\frac{\dd y}{(A+y^2)^{3/2}}
&=\frac1A\int_{-\pi/2}^{\pi/2}\cos t\dd t\\
&=\frac1A[\sin t]_{-\pi/2}^{\pi/2}=\frac2A.
\end{align*}
因此，按横坐标的边缘密度为
\[
g_x(x)=\frac h{2\pi(h^2+x^2)}.
\]
图\ref{mp1:fig:p7-s}所取横坐标区间对应整个长条。

\begin{figure}[H]\centering

\begin{tikzpicture}[x=1.5cm,y=8cm]
\fill[black!12](-1,0)--plot[domain=-1:1,samples=60](\x,{1/(2*pi*(1+(\x)^2))})--(1,0)--cycle;
\draw[axis](-2.8,0)--(3,0)node[right]{$x/h$};
\draw[axis](0,0)--(0,.24)node[above]{$hg_x$};
\draw[curve,domain=-2.6:2.6,samples=100]plot(\x,{1/(2*pi*(1+(\x)^2))});
\draw[guide](-1,0)--(-1,.0795775) (1,0)--(1,.0795775);
\node[left]at(0,.159155){$1/(2\pi)$};
\node[below]at(-1,0){$-a/h$};\node[below]at(1,0){$a/h$};
\end{tikzpicture}
\caption{习题~\ref{mp1:pr:p7}解图：先沿长条方向积分，得到横坐标的边缘密度，再对 $[-a,a]$ 积分。}\label{mp1:fig:p7-s}
\end{figure}

因此
\[
P=\int_{-a}^a g_x(x)\dd x
=\frac1{2\pi}\left[\arctan\frac xh\right]_{-a}^a
=\frac1\pi\arctan\frac ah.
\]
结果与二维射线命中 $|y|\le a$ 的无条件概率相同。三维情形是先把长条方向上的所有落点合并，其方向区域也等效为一个有限角宽的二面角。若改为半径 $a$ 的圆靶，需在两个坐标方向同时限制，便不能使用长条结果。
再对所有$x$积分，$\int_{-\infty}^{\infty}g_x\dd x=1/2$，因为只有朝向靶面的一半射线能命中平面。图\ref{mp1:fig:p7-s}中阴影面积是无条件概率；若以已经命中平面为条件，须将结果除以$1/2$，密度和命中长条的比例都加倍。
\end{solution}

\par\medskip\noindent\begin{minipage}{\textwidth}
\begin{exercise}[矩形靶的立体角（选做）]\label{mp1:pr:p8}
全空间各向同性点源位于原点，矩形靶在 $z=h>0$ 平面内，占据 $|x|\le a,|y|\le b$，见图\ref{mp1:fig:p8-q}。求其立体角与命中概率，并求小靶、无限长条和源点趋近靶面的极限。
\end{exercise}

\begin{figure}[H]\centering

\begin{tikzpicture}[scale=1.1]
\draw(-1.8,1.9)--(1.4,2.5)--(2.4,1.4)--(-.8,.8)--cycle;\fill(0,-1)circle(1.6pt)node[below]{$O$};\foreach\p in{(-1.8,1.9),(1.4,2.5),(2.4,1.4),(-.8,.8)}{\draw[guide](0,-1)--\p;}
\draw[guide](0,-1)--(.3,1.65)node[midway,right]{$h$};
\fill(.3,1.65)circle(1.2pt)node[above right]{$C$};\node[above]at(-.2,2.25){$2a$};\node[right]at(2,2){$2b$};\node[left]at(-1.9,1.2){$z=h$};
\end{tikzpicture}
\caption{习题~\ref{mp1:pr:p8}题设：$C$ 为位于 $z=h$ 平面内的矩形中心，$OC=h$。矩形的两条全边长分别为 $2a$ 和 $2b$。}\label{mp1:fig:p8-q}
\end{figure}

\end{minipage}\par\medskip

\begin{solution}
图\ref{mp1:fig:p8-s}中的第一象限与其余三象限对称，故
\[
\Omega=4\int_0^a\int_0^b\frac{h\dd y\dd x}{(h^2+x^2+y^2)^{3/2}}.
\]
先对 $y$ 积分，令 $A=h^2+x^2$，有
\[
\int_0^b\frac{h\dd y}{(A+y^2)^{3/2}}
=\frac hA\left[\frac y{\sqrt{A+y^2}}\right]_0^b
=\frac{hb}{(h^2+x^2)\sqrt{h^2+x^2+b^2}}.
\]
剩下的积分取原函数 $F(x)=\arctan u(x)$，其中 $B=h^2+b^2$、$u=bx/[h\sqrt{B+x^2}]$。为核实这一原函数，逐项求导：
\[
\begin{aligned}
u'&=\frac bh\left[(B+x^2)^{-1/2}-x^2(B+x^2)^{-3/2}\right]
=\frac{bB}{h(B+x^2)^{3/2}},\\
1+u^2&=\frac{h^2(B+x^2)+b^2x^2}{h^2(B+x^2)}
=\frac{B(h^2+x^2)}{h^2(B+x^2)},\\
F'&=\frac{u'}{1+u^2}
=\frac{hb}{(h^2+x^2)\sqrt{B+x^2}}.
\end{aligned}
\]
故可用
\[
F(x)=\arctan\frac{bx}{h\sqrt{h^2+x^2+b^2}}.
\]

\begin{figure}[H]\centering

\begin{tikzpicture}[scale=1.2]
\fill[black!12](0,0)rectangle(1.6,1.1);
\draw[axis](-2,0)--(2.2,0)node[right]{$x$};\draw[axis](0,-1.5)--(0,1.8)node[above]{$y$};\draw(-1.6,-1.1)rectangle(1.6,1.1);\node[below]at(1.6,0){$a$};\node[left]at(0,1.1){$b$};\draw[guide](1,0)--(1,1.1);\draw[thick](1.02,.5)rectangle(1.22,.7);\node[right]at(1.25,.6){$\dd x\dd y$};
\end{tikzpicture}
\caption{习题~\ref{mp1:pr:p8}解图：阴影为第一象限的积分区域，面积微元为 $\dd x\dd y$。四个象限对称，总立体角等于阴影区域积分的4倍。}\label{mp1:fig:p8-s}
\end{figure}

代入 $F(a)-F(0)$，得到
\[
\Omega=4\arctan\frac{ab}{h\sqrt{h^2+a^2+b^2}},\qquad
P=\frac1\pi\arctan\frac{ab}{h\sqrt{h^2+a^2+b^2}}.
\]
当 $a,b\ll h$，$P\simeq ab/(\pi h^2)$，等于靶面积 $4ab$ 除以球面积 $4\pi h^2$。当 $b\to\infty$，$P\to\arctan(a/h)/\pi$，回到长条靶结果。固定 $a,b>0$ 而令 $h\to0^+$，$P\to1/2$；正前方矩形张开的立体角趋于整个前半空间。
\end{solution}

\par\medskip\noindent\begin{minipage}{\textwidth}
\begin{exercise}[变密度悬绳的张力与伸长]\label{mp1:pr:p9}
一条细绳原长 $L$，上端固定、下端自由，线密度按原长坐标定义为 $\lambda(s)=\lambda_0(1+s/L)$，$s$ 从下端向上量。绳的轴向拉伸刚度为常量 $K$，其微小应变满足 $\epsilon=T/K$。求静止时张力分布和总伸长的一阶近似。见图\ref{mp1:fig:p9-q}。
\end{exercise}

\begin{figure}[H]\centering

\begin{tikzpicture}[scale=1.1]
\draw[thick](-.6,2.5)--(.6,2.5);\draw[very thick](0,0)--(0,2.5);\draw[axis](1,0)--(1,2.9)node[above]{$s$};\node[right]at(1,2.5){$L$};\node[right]at(1,0){$0$};\node[left]at(-.2,1.2){$\lambda_0(1+s/L)$};\node[below]at(0,0){自由端};
\end{tikzpicture}
\caption{习题~\ref{mp1:pr:p9}题设：悬绳以未伸长的原长坐标 $s$ 标记，下端自由，上端固定。线密度是对原长定义的。}\label{mp1:fig:p9-q}
\end{figure}

\end{minipage}\par\medskip

\begin{solution}
对图\ref{mp1:fig:p9-s}中截面以下的整段列平衡，截面张力承担原长区间 $[0,s]$ 中所有质量的重力：
\[
T(s)=g\int_0^s\lambda_0\left(1+\frac uL\right)\dd u
=\lambda_0g\left[u+\frac{u^2}{2L}\right]_0^s
=\lambda_0g\left(s+\frac{s^2}{2L}\right).
\]
它满足自由端 $T(0)=0$，上端张力为 $T(L)=3\lambda_0gL/2$，等于整绳重量。

\begin{figure}[H]\centering

\begin{tikzpicture}[scale=1.1]
\draw[line width=2.4pt,black!35](0,0)--(0,1.8);\draw[axis,thick](0,1.8)--(0,2.9)node[above]{$T(s)$};\draw[axis](0,.9)--(0,-.2)node[below]{$g\int_0^s\lambda(u)\dd u$};\fill(0,.9)circle(1.1pt);\node[right]at(.2,1.8){截面 $s$};\draw[guide](-.6,1.8)--(.6,1.8);\node[right]at(2,1.3){$\dd(\Delta L)=T(s)\dd s/K$};
\end{tikzpicture}
\caption{习题~\ref{mp1:pr:p9}解图：隔离截面以下的整段，截面张力等于这段绳的重力；再对伸长微元积分。}\label{mp1:fig:p9-s}
\end{figure}

微小伸长为 $\dd(\Delta L)=\epsilon\dd s=T(s)\dd s/K$，故
\[
\Delta L=\frac{\lambda_0g}{K}\int_0^L\left(s+\frac{s^2}{2L}\right)\dd s
=\frac{\lambda_0g}{K}\left[\frac{s^2}{2}+\frac{s^3}{6L}\right]_0^L
=\frac{2\lambda_0gL^2}{3K}.
\]
其中 $K$ 的单位为力，例如均匀截面材料可取 $K=YA$。小应变条件可写为 $T(L)/K\ll1$。质量微元始终按原长定义，伸长后实际空间中的线密度会改变，但绳段质量不变。
由同一材料坐标还可求每个截面的位移。固定上端，取向下的位移为$u(s)$，则长度增加要求$u_s=-T/K$、$u(L)=0$，从而
\[
 u(s)=\int_s^L\frac{T(\xi)}K\dd\xi
 =\frac{\lambda_0g}{K}\left[\frac{\xi^2}{2}+\frac{\xi^3}{6L}\right]_s^L
 =\frac{\lambda_0g}{K}
 \left[\frac{L^2-s^2}{2}+\frac{L^3-s^3}{6L}\right].
\]
自由端位移$u(0)$等于总伸长。图\ref{mp1:fig:p9-s}中，截面以下的重量随$s$增大而增加。截面的位移则是它与固定端之间各段伸长的累积，所以$u(s)$随$s$减小而增加。
\end{solution}

\par\medskip\noindent\begin{minipage}{\textwidth}
\begin{exercise}[按水平长度均匀的悬索载荷]\label{mp1:pr:p10}
一条无自重、不可抗弯的柔索两端等高，半跨距为 $a$。它承担按水平长度均匀的竖直向下载荷，强度为 $w$，水平张力为 $H>0$。按图\ref{mp1:fig:p10-q}以最低点为原点，求索形、垂度和两端张力，并解释它为何与均匀自重链不同。
\end{exercise}

\begin{figure}[H]\centering

\begin{tikzpicture}[scale=1.1]
\draw[curve,domain=-2:2]plot(\x,{.35*(\x)^2});\draw[guide](-2,1.4)--(2,1.4);\foreach\x in{-1.8,-1.2,-.6,0,.6,1.2,1.8}{\draw[axis](\x,{.35*(\x)^2+.75})--(\x,{.35*(\x)^2+.06});}\node[above]at(0,2.3){每单位水平长度载荷 $w$};\fill(-2,1.4)circle(2pt);\fill(2,1.4)circle(2pt);\draw[axis](-2.3,0)--(2.4,0)node[right]{$x$};\node[below]at(2,0){$a$};\node[below]at(-2,0){$-a$};
\end{tikzpicture}
\caption{习题~\ref{mp1:pr:p10}题设：载荷按水平长度均匀分布：水平区间 $\dd x$ 上的竖直载荷为 $w\dd x$。}\label{mp1:fig:p10-q}
\end{figure}

\end{minipage}\par\medskip

\begin{solution}
图\ref{mp1:fig:p10-s}左图中，两端张力的水平分量大小均为 $H$。对水平宽 $\dd x$ 的微元，竖直载荷为 $w\dd x$，因此
\[
\dd T_y=w\dd x,\qquad T_y=Hy',\qquad Hy''=w.
\]
取 $y'(0)=0,y(0)=0$，连续积分两次，
\[
y'(x)=\frac wH\int_0^x\dd\xi=\frac wH[\xi]_0^x=\frac wHx,
\qquad
y(x)=\frac wH\int_0^x\xi\dd\xi=\frac wH\left[\frac{\xi^2}{2}\right]_0^x.
\]
因此索形为抛物线，
\[
y(x)=\frac w{2H}x^2,\qquad f=y(a)=\frac{wa^2}{2H}.
\]

\begin{figure}[H]\centering

\begin{tikzpicture}[scale=1.05]
\draw[curve,domain=.4:1.2,samples=30]plot(\x,{.4*(\x)^2});
\draw[axis,thick](.4,.064)--(-.5,-.224)node[left]{$\vct T_-$};
\draw[axis,thick](1.2,.576)--(2.1,1.44)node[above]{$\vct T_+$};
\draw[axis](.8,.256)--(.8,-.32)node[below]{$w\dd x$};
\draw[guide](.4,.064)--(.4,1.95) (1.2,.576)--(1.2,1.95);
\draw[|-|](.4,1.85)--(1.2,1.85)node[midway,above]{$\dd x$};
\node at(.8,2.5){微元受力};
\begin{scope}[shift={(5.6,0)}]
\draw[curve,domain=-1.5:1.5,samples=70]plot(\x,{.32*(\x)^2});
\draw[axis,thick](-1.5,.72)--(-2.3,1.488)node[above]{$\vct T_A$};
\draw[axis,thick](1.5,.72)--(2.3,1.488)node[above]{$\vct T_B$};
\draw[axis](0,.3)--(0,-1.236)node[below]{$2aw$};
\node[left]at(-1.5,.72){$A$};\node[right]at(1.5,.72){$B$};
\node at(0,2.5){整条绳受力};
\end{scope}
\end{tikzpicture}
\caption{习题~\ref{mp1:pr:p10}解图：左图的水平张力分量相互抵消，竖直分量之差平衡 $w\dd x$。右图把整条绳的分布载荷合并为 $2aw$，由两端张力共同承担。}\label{mp1:fig:p10-s}
\end{figure}

图\ref{mp1:fig:p10-s}右图中，每端张力的竖直分量为 $wa$，故
\[
T_A=T_B=\sqrt{H^2+w^2a^2}.
\]
两种索形使用不同的载荷微元：
\[
\dd G_{\text{本题}}=w\dd x,\qquad
\dd G_{\text{自重}}=\lambda g\sqrt{1+y'^2}\dd x.
\]
前者的 $y''$ 为常量，后者的 $y''$ 随斜率改变。小斜率时 $\dd s\simeq\dd x$，悬链线才可近似为抛物线。
索长也可由已经求得的形状核实。令 $u=wx/H$，$\dd x=(H/w)\dd u$，上限记为 $q=wa/H$。为算出弧长积分，先记 $J(q)=\int_0^q\sqrt{1+u^2}\dd u$。分部积分并用 $u^2=(1+u^2)-1$，得到
\[
\begin{aligned}
J(q)&=[u\sqrt{1+u^2}]_0^q-\int_0^q\frac{u^2}{\sqrt{1+u^2}}\dd u\\
&=q\sqrt{1+q^2}-J(q)+[\operatorname{arsinh}u]_0^q.
\end{aligned}
\]
移项可得 $2J=q\sqrt{1+q^2}+\operatorname{arsinh}q$，故
\[
 \ell=2\int_0^a\sqrt{1+(wx/H)^2}\dd x
 =\frac{2H}{w}J(q)
 =a\sqrt{1+q^2}+\frac Hw\operatorname{arsinh}q.
\]
若另给定索长，以上关系便可用于确定$H$，再求垂度。小斜率$q\ll1$时，$\ell=2a+w^2a^3/(3H^2)+O(aq^4)$。图\ref{mp1:fig:p10-s}中的水平分力恒定，总张力为$T(x)=\sqrt{H^2+w^2x^2}$，从最低点向两端逐渐增大。
\end{solution}

\par\medskip\noindent\begin{minipage}{\textwidth}
\begin{exercise}[非均匀束流对固定球的推力]\label{mp1:pr:p11}
固定硬球的碰撞条件同例\ref{mp1:ex:e8}，但轴对称入射通量随瞄准距离改变：$j(b)=j_0(1+\eta b^2/R^2)$，$0\le b\le R$，$\eta>-1$。图\ref{mp1:fig:p11-q}表示入射截面。求单位时间碰撞数、每单位立体角的出射速率密度及平均轴向推力。
\end{exercise}

\begin{figure}[H]\centering

\begin{tikzpicture}[scale=1.1]
\draw(0,0)circle(1.5);
\foreach\r in{.4,.8,1.2}{\draw[guide](0,0)circle(\r);}
\fill(0,0)circle(1.2pt)node[below left]{$B$};
\draw(0,0)--(.95,0)node[midway,above]{$b$};
\fill(.95,0)circle(1.2pt)node[right=2pt]{$P$};
\draw(0,0)--(-65:1.5)node[midway,left=3pt,fill=white,inner sep=1pt]{$R$};
\node[anchor=west]at(2.5,.8){$j(b)=j_0(1+\eta b^2/R^2)$};
\node[anchor=west]at(2.5,.1){$\eta>-1$};
\end{tikzpicture}
\caption{习题~\ref{mp1:pr:p11}题设：入射截面上的轴对称束流。$B$ 为靶心投影，$BP=b$；同一虚线圆周上的束流密度相同。}\label{mp1:fig:p11-q}
\end{figure}

\end{minipage}\par\medskip

\begin{solution}
先在入射截面相加，单位时间碰撞数为
\[
\Gamma=\int_0^Rj(b)2\pi b\dd b
=2\pi j_0\left[\frac{b^2}{2}+\frac{\eta b^4}{4R^2}\right]_0^R
=\pi j_0R^2\left(1+\frac\eta2\right).
\]
几何映射仍为 $b=R\cos(\theta/2)$。因此单位立体角、单位时间的出射粒子数为
\[
\frac{\dd\Gamma}{\dd\Omega}
=j(b)\frac{\dd\sigma}{\dd\Omega}
=\frac{j_0R^2}{4}\left(1+\eta\cos^2\frac\theta2\right).
\]
除以 $\Gamma$ 才得到以发生碰撞为条件的方向概率密度。

\begin{figure}[H]\centering

\begin{tikzpicture}[scale=1.05]
\pgfmathsetmacro{\plen}{1.8}
\pgfmathsetmacro{\scatterangle}{55}
\coordinate(B)at(0,0);
\coordinate(Q)at({\plen*cos(\scatterangle)},{\plen*sin(\scatterangle)});
\draw[axis](-2.2,0)--(-.4,0)node[midway,above]{$mv$};
\draw[axis](B)--(2.35,0)node[right]{$z$};
\draw[axis,thick](B)--(Q)node[above]{$mv$};
\draw[guide](Q)--({\plen*cos(\scatterangle)},0);
\fill(B)circle(1.3pt)node[below left]{$B$};
\draw(.55,0)arc(0:\scatterangle:.55);
\node at(.65,.35){$\theta$};
\node[below]at(.7,-.2){$mv\cos\theta$};
\node[anchor=west]at(3,.8){$\Delta p_z=mv(1-\cos\theta)$};
\end{tikzpicture}
\caption{习题~\ref{mp1:pr:p11}解图：固定靶弹性碰撞的动量示意。粒子碰撞前后的动量大小均为 $mv$，出射轴向分量为 $mv\cos\theta$；靶获得的轴向动量为二者之差。}\label{mp1:fig:p11-s}
\end{figure}

图\ref{mp1:fig:p11-s}中单次轴向动量转移为 $mv(1-\cos\theta)$。用 $u=b^2/R^2$，有 $\dd u=2b\dd b/R^2$、$\cos\theta=2u-1$，而 $b=0,R$ 对应 $u=0,1$，故
\begin{align*}
F_z&=\int_0^R mv(1-\cos\theta)j(b)2\pi b\dd b\\
&=\pi j_0R^2mv\int_0^1 2(1-u)(1+\eta u)\dd u\\
&=\pi j_0R^2mv\left[2u+(\eta-1)u^2-\frac{2\eta}{3}u^3\right]_0^1\\
&=\pi j_0R^2mv\left(1+\frac\eta3\right).
\end{align*}
若 $\eta>0$，额外通量偏向碰撞截面的外缘，那里多为小角散射，单次动量转移较小。因此碰撞数的增量系数为 $\eta/2$，推力的增量系数为 $\eta/3$。
\end{solution}

\par\medskip\noindent\begin{minipage}{\textwidth}
\begin{exercise}[角度筛选与出射动能]\label{mp1:pr:p12}
重复例\ref{mp1:ex:e9}的等质量自由靶实验，每次靶球初始静止。以碰撞点为近似原点，在前方距离 $D\gg R$ 处放置半径 $a$ 的圆形探测器，圆心位于入射轴上，见图\ref{mp1:fig:p12-q}。忽略有限球径和飞行过程中的外力。求入射球被收集的比例，以及被收集粒子的平均出射动能。
\end{exercise}

\begin{figure}[H]\centering

\begin{tikzpicture}[scale=1.05]
\draw[axis](0,0)--(5,0)node[right]{$z$};\fill(0,0)circle(2pt)node[below]{$B$};\draw(3.8,0)ellipse(.35 and 1.4);\draw(0,0)--(3.8,1.4) (0,0)--(3.8,-1.4);\draw[guide](3.8,0)--(3.8,1.4)node[midway,right]{$a$};\draw(.8,0)arc(0:20.225:.8);\node at(1.2,.24){$\alpha$};\node[below]at(2,0){$D$};
\end{tikzpicture}
\caption{习题~\ref{mp1:pr:p12}题设：探测器只收集前向锥 $\theta\le\alpha$ 内的入射球。探测器足够远，可忽略碰撞位置的偏移。}\label{mp1:fig:p12-q}
\end{figure}

\end{minipage}\par\medskip

\begin{solution}
探测器对应半顶角 $\alpha=\arctan(a/D)$ 的前向圆锥。入射球的极角密度为 $p_\theta=2\sin\theta\cos\theta$，故在已经碰撞的事件中，收集比例为
\[
P_{\rm det}=\int_0^\alpha2\sin\theta\cos\theta\dd\theta
=\left[\sin^2\theta\right]_0^\alpha
=\sin^2\alpha=\frac{a^2}{D^2+a^2}.
\]
图\ref{mp1:fig:p12-s}给出垂直于入射轴的截面。$B$ 为靶心投影，$P$ 为入射球心轨迹与截面的交点，$BP=b$。由 $b=R\cos\theta$，接收范围 $0\le\theta\le\alpha$ 对应阴影外环 $R\cos\alpha\le b\le R$，其面积占比为 $[\pi R^2-\pi(R\cos\alpha)^2]/(\pi R^2)=\sin^2\alpha$。

\begin{figure}[H]\centering

\begin{tikzpicture}[scale=1.1]
% The same illustrative acceptance angle is used in the preceding detector diagram.
\pgfmathsetmacro{\outerradius}{1.5}
\pgfmathsetmacro{\accepthalfangle}{atan(1.4/3.8)}
\pgfmathsetmacro{\innerradius}{\outerradius*cos(\accepthalfangle)}
\pgfmathsetmacro{\sampleradius}{(\outerradius+\innerradius)/2}
\coordinate(B)at(0,0);
\coordinate(P)at(\sampleradius,0);
\fill[even odd rule,black!15](B)circle(\outerradius) (B)circle(\innerradius);
\draw(B)circle(\outerradius);
\draw(B)circle(\innerradius);
\draw[guide](B)--(135:\innerradius);
\node[above left=5pt]at(135:\innerradius){$R\cos\alpha$};
\draw(B)--(-50:\outerradius)node[midway,below left=2pt]{$R$};
\draw(B)--(P)node[midway,above=2pt]{$b$};
\fill(B)circle(1.3pt)node[below left=2pt]{$B$};
\fill(P)circle(.8pt)node[right=5pt]{$P$};
\node[anchor=west]at(2.5,.55){$\varepsilon=E_A/E_0=b^2/R^2$};
\node[anchor=west]at(2.5,-.25){$\cos^2\alpha\le\varepsilon\le1$};
\end{tikzpicture}
\caption{习题~\ref{mp1:pr:p12}解图：入射截面内的接收环带，$BP=b$。内、外半径分别为 $R\cos\alpha$ 和 $R$，阴影区内的轨迹对应出射角 $\theta\le\alpha$。}\label{mp1:fig:p12-s}
\end{figure}

由 $\varepsilon=E_A/E_0=b^2/R^2$，原本均匀的动能比分布被截取为
\[
\cos^2\alpha\le\varepsilon\le1.
\]
在这个区间内条件密度为 $1/\sin^2\alpha$，平均动能为
\begin{align*}
\avg{E_A\mid\text{收集}}
&=\frac{E_0}{\sin^2\alpha}\int_{\cos^2\alpha}^1\varepsilon\dd\varepsilon\\
&=\frac{E_0}{\sin^2\alpha}\left[\frac{\varepsilon^2}{2}\right]_{\cos^2\alpha}^1
=\frac{E_0(1-\cos^4\alpha)}{2\sin^2\alpha}\\
&=\frac{E_0}{2}(1+\cos^2\alpha)
=\frac{E_0}{2}\left(1+\frac{D^2}{D^2+a^2}\right).
\end{align*}
窄接收锥主要保留擦边、小偏转的粒子，条件平均动能较高。接收半角从0增至 $\pi/2$ 时，条件平均从 $E_0$ 降至全体出射粒子的平均值 $E_0/2$。
\end{solution}






\chapter{运动关联进阶}

绳长、杆长和接触条件限制了各部分的位置，对这些位置关系连续求导，便得到速度与加速度的关联。构件转动时，杆和绳段的方向也在改变；质点若还相对构件移动，两种运动须一并计入。瞬时图形中的速度比例，不能直接套用于加速度。


\section{绳杆约束与刚体关联}
不可伸长或长度不变的条件限制系统位形，求导后才得到速度与加速度关联。固定端、移动端及方向随构形变化的部分须在方程中分别表示。

\subsection{从位置关系连续求导}
设系统用坐标 $q_1,\ldots,q_n$ 描述，几何约束写为
\[
F(q_1,\ldots,q_n,t)=0.
\]
其中的显含时间项表示导轨、支点或绳长受到外界驱动。沿实际运动求导，得到
\[
\sum_i F_{q_i}\dot q_i+F_t=0.
\]
再求一次导数时，对每个乘积 $F_{q_i}\dot q_i$ 分别求导，并保留 $F_t$ 的全导数：
\[
\begin{aligned}
\frac{\dd}{\dd t}(F_{q_i}\dot q_i)
&=F_{q_i}\ddot q_i+
\left(\sum_jF_{q_iq_j}\dot q_j+F_{q_it}\right)\dot q_i,\\
\frac{\dd F_t}{\dd t}&=\sum_iF_{tq_i}\dot q_i+F_{tt}.
\end{aligned}
\]
假定约束函数二阶偏导连续，合并相等的混合偏导项，得到
\[
\sum_i F_{q_i}\ddot q_i
+\sum_{i,j}F_{q_iq_j}\dot q_i\dot q_j
+2\sum_iF_{q_it}\dot q_i+F_{tt}=0.
\]
第二次求导出现了速度的二次项。约束曲线的弯曲、杆的转向和活动绳段方向的改变，都会产生这类项；驱动端匀速时也须保留。

若约束只有 $F(x,y)=0$，且 $F_y\ne0$，可以在局部把 $y$ 看成 $x$ 的函数。此时
\[
\dot y=-\frac{F_x}{F_y}\dot x,
\quad
\ddot y=-\frac{F_x}{F_y}\ddot x
-\frac{F_{xx}\dot x^2+2F_{xy}\dot x\dot y+F_{yy}\dot y^2}{F_y}.
\]
分母趋零时，先回到原约束判断构形。此时选定的驱动坐标可能失效，规定的非零驱动速度也可能使其余部分的速度无界，因而无法维持到这一位置。

绳长约束须按实际路径分段。定滑轮只改变绳的方向，绕在固定半径滑轮上的恒定包角部分可以并入常数；动滑轮两侧的直绳段通常同时改变。若切点随运动移动，包角也随之改变，就不能把相应弧长随意删去。本讲采用的理想小滑轮均明确忽略半径，或保持包角不变。

\subsection{刚杆的速度与加速度关系}
两端点位置差记为 $\vct r=\vct r_B-\vct r_A$。刚杆满足 $\vct r\cdot\vct r=l^2$。两次求导分别给出
\[
2\vct r\cdot\dot{\vct r}=0,\qquad
2\dot{\vct r}\cdot\dot{\vct r}+2\vct r\cdot\ddot{\vct r}=0.
\]
代入 $\dot{\vct r}=\vct v_B-\vct v_A$、$\ddot{\vct r}=\vct a_B-\vct a_A$，得到
\begin{align*}
\vct r\cdot(\vct v_B-\vct v_A)&=0,\\
\vct r\cdot(\vct a_B-\vct a_A)&=-|\vct v_B-\vct v_A|^2.
\end{align*}
第一式说明两端沿杆方向的速度分量相同；第二式说明沿杆方向的加速度分量一般不同。

在平面运动中，取逆时针为角速度的正方向，$\vct k$ 垂直纸面向外，记 $\omega=\dot\theta$、$\alpha=\ddot\theta$。同一刚体上两点满足
\begin{equation}\label{mp2:eq:l2-rigid-v}
\vct v_B=\vct v_A+\omega\vct k\times\vct r,
\end{equation}
\begin{equation}\label{mp2:eq:l2-rigid-a}
\vct a_B=\vct a_A+\alpha\vct k\times\vct r-\omega^2\vct r.
\end{equation}
图\ref{mp2:fig:l2rigid}将相对加速度分成切向项与法向项。法向项总指向参考点 $A$，不随 $\omega$ 的正负改变；切向项随 $\alpha$ 改变方向。
式\eqref{mp2:eq:l2-rigid-a}也可直接对相对速度求导。因 $\dot{\vct r}=\omega\vct k\times\vct r$，
\[
\begin{aligned}
\vct a_B-\vct a_A
&=\alpha\vct k\times\vct r+\omega\vct k\times\dot{\vct r}\\
&=\alpha\vct k\times\vct r+\omega^2\vct k\times(\vct k\times\vct r)\\
&=\alpha\vct k\times\vct r-\omega^2\vct r.
\end{aligned}
\]

\begin{figure}[H]\centering
\begin{tikzpicture}[scale=.95]
\draw[thick](0,0)--(2.7,0);
\fill(0,0)circle(1.3pt)node[below]{$A$};\fill(2.7,0)circle(1.3pt)node[below]{$B$};
\draw[axis](2.7,0)--(2.7,1.6)node[above]{$\vct v_B-\vct v_A$};
\draw[|-|](0,-.55)--(2.7,-.55)node[midway,fill=white]{$l$};
\begin{scope}[shift={(6,0)}]
\draw[thick](0,0)--(2.7,0);\fill(0,0)circle(1.3pt)node[below]{$A$};\fill(2.7,0)circle(1.3pt)node[below right]{$B$};
\draw[axis,thick](2.7,0)--(1.2,0)node[midway,below]{$\omega^2l$};
\draw[axis](2.7,0)--(2.7,1.3)node[midway,right]{$\alpha l$};
\draw[axis](2.7,0)--(1.2,1.3)node[above left]{$\vct a_B-\vct a_A$};
\draw[guide](1.2,0)--(1.2,1.3)--(2.7,1.3);
\end{scope}
\end{tikzpicture}
\caption{同一刚体上两点的关联。左图的相对速度垂直于杆；右图的相对加速度还含沿杆指向 $A$ 的法向项。}\label{mp2:fig:l2rigid}
\end{figure}

相对速度垂直于杆，可以据此作速度图。作加速度图时，还须加入 $-\omega^2\vct r$。平面刚体的速度瞬心在当时速度为零，相应材料点的加速度却未必为零；若以该点为参考计算其他点的加速度，它自身的加速度仍须保留。

\par\medskip\noindent\begin{minipage}{\textwidth}
\begin{example}[斜绳牵引中的二阶几何项]\label{mp2:ex:l2e1}

固定小滑轮 $O$ 比水平导轨高 $h$，滑块 $A$ 在导轨上，另一端滑块 $B$ 竖直运动。绳轻、不可伸长并始终绷紧，忽略滑轮半径，见图\ref{mp2:fig:l2e1-q}。取 $A$ 到 $O$ 铅垂投影的水平距离为 $x>0$，$B$ 到滑轮的向下距离为 $y$。已知 $\dot x=u$、$\ddot x=a$，求 $\dot y,\ddot y$。当 $u>0,a=0$ 时，判断 $B$ 的运动方向及加速度方向。

\end{example}

\begin{figure}[H]\centering
\begin{tikzpicture}[scale=1.05]
\coordinate(O)at(0,2.7);\coordinate(A)at(3.6,0);\coordinate(B)at(0,.6);
\draw[thick](-.5,0)--(4.3,0);\draw[thick](O)--(A);\draw[thick](O)--(B);
\draw[fill=white](O)circle(.12);\draw[fill=white](3.35,-.08)rectangle(3.85,.25);
\draw[fill=white](-.22,.4)rectangle(.22,.8);
\node[above left]at(O){$O$};\node[right]at(3.95,.1){$A$};\node[left]at(-.32,.6){$B$};
\draw[guide](0,0)--(O);\draw[|-|](-.7,0)--(-.7,2.7)node[midway,left]{$h$};
\draw[|-|](0,-.5)--(3.6,-.5)node[midway,fill=white]{$x$};
\draw[axis](.65,2.7)--(.65,.6)node[midway,right]{$y$};
\draw(3.05,0)arc(180:143.13:.55);\node at(2.83,.3){$\theta$};
\draw[axis](3.6,.55)--(4.4,.55)node[right]{$u$};
\end{tikzpicture}
\caption{例~\ref{mp2:ex:l2e1}题设：绳经固定小滑轮 $O$，一端为水平滑块 $A$，另一端为竖直滑块 $B$。滑轮半径忽略。}\label{mp2:fig:l2e1-q}
\end{figure}
\end{minipage}\par\medskip

\begin{solution}
图\ref{mp2:fig:l2e1-s}中的斜绳长为 $s=\sqrt{x^2+h^2}$，绳长给出 $s+y=L$。坐标 $y$ 向下为正，求导得
\[
\dot s=\frac{x}{s}u,\qquad \dot y=-\frac{x}{s}u.
\]
再求导时，$x/s$ 也在改变：
\[
\frac{\dd}{\dd t}\left(\frac xs\right)
=\frac{us-x\dot s}{s^2}
=\frac{h^2}{s^3}u.
\]
所以
\begin{equation}\label{mp2:eq:l2-rope-acc}
\ddot y=-\frac xs a-\frac{h^2}{s^3}u^2.
\end{equation}

\begin{figure}[H]\centering
\begin{tikzpicture}[scale=1.05]
\draw[thick](0,2.4)--(3.2,0);\draw[guide](0,2.4)--(0,0)--(3.2,0);
\node[above left]at(0,2.4){$O$};\node[below right]at(3.2,0){$A$};
\node[left]at(0,1.2){$h$};\node[below]at(1.6,0){$x$};\node[above right]at(1.6,1.2){$s$};
\draw(2.6,0)arc(180:143.13:.6);\node at(2.36,.3){$\theta$};
\begin{scope}[shift={(5,.7)}]
\node[anchor=west]at(0,1.25){$\dot s=(x/s)u$};
\node[anchor=west]at(0,.45){$\ddot s=(x/s)a+h^2u^2/s^3$};
\node[anchor=west]at(0,-.35){$\ddot y=-\ddot s$};
\end{scope}
\end{tikzpicture}
\caption{例~\ref{mp2:ex:l2e1}解图：$s=OA$，$\cos\theta=x/s$。右图保留了绳段转向产生的二阶几何项。}\label{mp2:fig:l2e1-s}
\end{figure}

图中 $\cos\theta=x/s$，速度关系也可写成 $\dot y=-u\cos\theta$；但把它再次求导时，必须保留 $\dot\theta$。匀速拉动 $A$ 时，$\dot y<0$ 且 $\ddot y<0$，故 $B$ 向上运动并加速。$A$ 离滑轮投影越远，水平速度在斜绳方向的投影越接近 $u$。

在 $x=0$ 的瞬间，$\dot y=0$，而 $\ddot y=-u^2/h$。斜绳长在这一位置取得极小值，其一阶变化为零，二阶变化仍非零。以上关系以绳保持绷紧为前提，实际运动中还须由张力判断这一条件能否维持。

\end{solution}

\par\medskip\noindent\begin{minipage}{\textwidth}
\begin{example}[沿直角导轨滑动的刚杆]\label{mp2:ex:l2e2}

长为 $l$ 的刚杆两端 $A,B$ 分别在水平、竖直导轨上运动，见图\ref{mp2:fig:l2e2-q}。取 $A=(x,0)$、$B=(0,y)$，且 $x,y>0$。若 $A$ 以恒定速率 $u>0$ 向右运动，求 $B$ 的速度、加速度，杆的角速度与角加速度，以及杆中点的速度。

\end{example}

\begin{figure}[H]\centering
\begin{tikzpicture}[scale=1.05]
\pgfmathsetmacro{\ang}{45}
\pgfmathsetmacro{\xx}{3.5*cos(\ang)}\pgfmathsetmacro{\yy}{3.5*sin(\ang)}
\draw[thick](-.2,0)--(3.9,0);\draw[thick](0,-.15)--(0,3.8);
\draw[line width=1pt](0,\yy)--(\xx,0);
\fill(0,\yy)circle(1.5pt)node[left]{$B$};\fill(\xx,0)circle(1.5pt)node[below right]{$A$};
\node[below left]at(0,0){$O$};\node[above right]at({\xx/2},{\yy/2}){$l$};
\draw[|-|](0,-.5)--(\xx,-.5)node[midway,fill=white]{$x$};
\node[left]at(0,{\yy/2}){$y$};
\end{tikzpicture}
\caption{例~\ref{mp2:ex:l2e2}题设：刚杆两端沿互相垂直的光滑导轨运动，$OA=x$，$OB=y$。}\label{mp2:fig:l2e2-q}
\end{figure}
\end{minipage}\par\medskip

\begin{solution}
由杆长不变，$x^2+y^2=l^2$。连续两次求导：
\[
xu+y\dot y=0,\qquad u^2+\dot y^2+y\ddot y=0.
\]
于是
\[
\dot y=-\frac xy u,\qquad
\ddot y=-\frac{u^2+\dot y^2}{y}=-\frac{l^2u^2}{y^3}.
\]
图\ref{mp2:fig:l2e2-s}中，上端的速度、加速度均向下。下端虽匀速，上端的速率却不断增加。

\begin{figure}[H]\centering
\begin{tikzpicture}[scale=.95]
\pgfmathsetmacro{\ang}{45}
\pgfmathsetmacro{\xx}{3.5*cos(\ang)}\pgfmathsetmacro{\yy}{3.5*sin(\ang)}
\draw[thick](-.2,0)--(3.9,0);\draw[thick](0,-.15)--(0,3.8);
\draw[line width=1pt](0,\yy)--(\xx,0);
\fill(0,\yy)circle(1.5pt)node[left]{$B$};\fill(\xx,0)circle(1.5pt)node[below right]{$A$};
\node[below left]at(0,0){$O$};\node[above right]at({\xx/2},{\yy/2}){$l$};
\draw[|-|](0,-.5)--(\xx,-.5)node[midway,fill=white]{$x$};
\node[left]at(0,{\yy/2}){$y$};

\draw[axis](\xx,.25)--({\xx+.8},.25)node[right]{$u$};
\draw[axis,thick](0,\yy)--(0,{\yy-.85})node[midway,left]{$|\dot y|$};
\node at(1.7,4.15){速度};
\begin{scope}[shift={(6.1,\yy)}]
\fill(0,0)circle(1.5pt)node[above]{$B$};
\draw[axis,thick](0,0)--(0,-1.6)node[midway,right]{$|\ddot y|=l^2u^2/y^3$};
\end{scope}
\node at(7.1,4.15){加速度};
\end{tikzpicture}
\caption{例~\ref{mp2:ex:l2e2}解图：左图标出两端速度，右图为 $B$ 的加速度。$B$ 的速度与加速度均向下；箭头长度仅作示意。}\label{mp2:fig:l2e2-s}
\end{figure}

沿杆从 $A$ 指向 $B$，有 $\vct r_{B/A}=(-x,y)$。代入速度式，
\[
(0,\dot y)=(u,0)+\omega(-y,-x),
\]
故 $\omega=u/y>0$，杆作逆时针转动。角加速度为
\[
\alpha=-\frac{u\dot y}{y^2}=\frac{xu^2}{y^3}>0.
\]
中点 $M=(x/2,y/2)$ 的速度是
\[
\vct v_M=\frac u2\vct e_x-\frac{xu}{2y}\vct e_y.
\]
在 $y\to0^+$ 时，$\omega$、$|\dot y|$ 和 $|\ddot y|$ 均无界。这说明给定的恒定水平驱动不能以有限运动量持续到完全水平构形。若杆只是靠在墙上，还可能在到达这一几何极限以前失去接触；题设的导轨连接则排除了脱离。

\end{solution}

\par\medskip\noindent\begin{minipage}{\textwidth}
\begin{example}[两连杆端点的竖直驱动]\label{mp2:ex:l2e3}

两根等长刚杆 $OA,AB$ 均长 $l$，$O$ 固定，$B$ 沿竖直导轨运动。某时刻 $OB=l$ 且水平，$A$ 位于上方，$B$ 的速度、加速度分别为 $v\vct e_y$、$a\vct e_y$，其中 $v>0$，$a$ 为有符号量，见图\ref{mp2:fig:l2e3-q}。取两杆绝对角速度、角加速度均以逆时针为正，求 $\omega_1,\omega_2,\alpha_1,\alpha_2$。

\end{example}

\begin{figure}[H]\centering
\begin{tikzpicture}[scale=1.05]
\coordinate(O)at(0,0);\coordinate(A)at(1.4,{1.4*sqrt(3)});\coordinate(B)at(2.8,0);
\draw[thick](-.15,-.4)--(-.15,2.7);\draw[guide](2.8,-.3)--(2.8,3.1);
\draw[line width=1pt](O)--(A)--(B);
\fill(O)circle(1.4pt)node[below left]{$O$};\fill(A)circle(1.4pt)node[above]{$A$};\fill(B)circle(1.4pt)node[below right]{$B$};
\draw[guide](O)--(B);\node[below]at(1.4,0){$l$};
\node[left]at(.6,1.2){$l$};\node[right]at(2.2,1.2){$l$};
\draw(.55,0)arc(0:60:.55);\node at(.7,.37){$\pi/3$};
\draw[axis](3.25,.15)--(3.25,1.4)node[right]{$v$};
\end{tikzpicture}
\caption{例~\ref{mp2:ex:l2e3}题设：$OA=AB=OB=l$ 的瞬时构形，端点 $B$ 沿竖直导轨向上运动。}\label{mp2:fig:l2e3-q}
\end{figure}
\end{minipage}\par\medskip

\begin{solution}
以 $O$ 为原点，在所给瞬间
\[
\vct r_A=\left(\frac l2,\frac{\sqrt3l}2\right),\qquad
\vct r_{B/A}=\left(\frac l2,-\frac{\sqrt3l}2\right).
\]
图\ref{mp2:fig:l2e3-s}右图是按这两个方向作出的速度三角形。先由杆 $OA$，再由杆 $AB$，有
\begin{align*}
\vct v_A&=l\omega_1\left(-\frac{\sqrt3}2,\frac12\right),\\
\vct v_B-\vct v_A&=l\omega_2\left(\frac{\sqrt3}2,\frac12\right).
\end{align*}
水平分量为零给出 $\omega_1=\omega_2$，竖直分量给出
\[
\omega_1=\omega_2=\frac vl.
\]

\begin{figure}[H]\centering
\begin{tikzpicture}[scale=.95]
\coordinate(O)at(0,0);\coordinate(A)at(1.4,{1.4*sqrt(3)});\coordinate(B)at(2.8,0);
\draw[thick](-.15,-.4)--(-.15,2.7);\draw[guide](2.8,-.3)--(2.8,3.1);
\draw[line width=1pt](O)--(A)--(B);
\fill(O)circle(1.4pt)node[below left]{$O$};\fill(A)circle(1.4pt)node[above]{$A$};\fill(B)circle(1.4pt)node[below right]{$B$};
\draw[guide](O)--(B);\node[below]at(1.4,0){$l$};
\node[left]at(.6,1.2){$l$};\node[right]at(2.2,1.2){$l$};
\draw(.55,0)arc(0:60:.55);\node at(.7,.37){$\pi/3$};

\begin{scope}[shift={(5.2,.1)}]
\coordinate(S)at(0,0);\coordinate(T)at({-1.2*sqrt(3)/2},.6);\coordinate(U)at(0,1.2);
\draw[axis,thick](S)--(T)node[midway,left]{$\vct v_A$};
\draw[axis,thick](T)--(U)node[midway,above left]{$\vct v_{B/A}$};
\draw[axis](S)--(U)node[midway,right]{$\vct v_B$};
\node at(-.3,-.55){$\omega_1=\omega_2=v/l$};
\end{scope}
\end{tikzpicture}
\caption{例~\ref{mp2:ex:l2e3}解图：右侧速度三角形中两条斜边水平分量相消。加速度关联另含两杆各自的法向项，不能照抄速度三角形。}\label{mp2:fig:l2e3-s}
\end{figure}

两角速度只在所给瞬间相等。求角加速度时，应重新写出两杆的加速度关系。记 $w=v/l$，加速度关联为
\[
\vct a_A=\alpha_1\vct k\times\vct r_A-w^2\vct r_A,
\]
\[
\vct a_B=\vct a_A+\alpha_2\vct k\times\vct r_{B/A}-w^2\vct r_{B/A}.
\]
两式相加并投影到坐标轴：
\[
0=\frac{\sqrt3l}{2}(\alpha_2-\alpha_1)-lw^2,
\qquad
a=\frac l2(\alpha_1+\alpha_2).
\]
解得
\[
\alpha_1=\frac al-\frac{v^2}{\sqrt3l^2},\qquad
\alpha_2=\frac al+\frac{v^2}{\sqrt3l^2}.
\]
取 $a=0$，两杆的角加速度一负一正，尽管这一瞬间的角速度相等。两杆各自的法向加速度使速度三角形不能直接用于加速度。

\end{solution}

\par\medskip\noindent\begin{minipage}{\textwidth}
\begin{example}[滑动方板的质心与顶点轨迹]\label{mp2:ex:l2e4}

边长为 $l$ 的均匀正方形薄板 $ABED$，相邻顶点 $A,B$ 分别沿竖直、水平导轨运动，见图\ref{mp2:fig:l2e4-q}。$C$ 为质心，$D$ 为与 $B$ 相对的顶点。设边 $AB$ 与向下竖直方向的夹角为 $0\le\theta\le\pi/2$，$\dot\theta=\omega$、$\ddot\theta=\alpha$。求 $C$ 的实际轨迹、速度和加速度，以及 $D$ 的轨迹；讨论 $\theta=\pi/4$ 时的运动。

\end{example}

\begin{figure}[H]\centering
\begin{tikzpicture}[scale=1.05]
\pgfmathsetmacro{\ang}{30}\pgfmathsetmacro{\ll}{2.5}
\coordinate(O)at(0,0);\coordinate(A)at(0,{\ll*cos(\ang)});\coordinate(B)at({\ll*sin(\ang)},0);
\coordinate(D)at({\ll*cos(\ang)},{\ll*(cos(\ang)+sin(\ang))});
\coordinate(E)at({\ll*(sin(\ang)+cos(\ang))},{\ll*sin(\ang)});
\coordinate(C)at({\ll*(sin(\ang)+cos(\ang))/2},{\ll*(sin(\ang)+cos(\ang))/2});
\fill[black!5](A)--(B)--(E)--(D)--cycle;\draw[thick](A)--(B)--(E)--(D)--cycle;
\draw[axis](-.2,0)--(4.1,0)node[right]{$x$};\draw[axis](0,-.2)--(0,4.2)node[above]{$y$};
\foreach\p/\pos in{A/left,B/below,E/right,D/above,C/above left}{\fill(\p)circle(1.2pt)node[\pos]{$\p$};}
\node[below left]at(O){$O$};\draw[guide](A)--(E) (B)--(D);
\draw(0,{\ll*cos(\ang)-.6})arc(270:300:.6);\node at(.25,{\ll*cos(\ang)-.85}){$\theta$};
\end{tikzpicture}
\caption{例~\ref{mp2:ex:l2e4}题设：正方形 $ABED$ 的 $A,B$ 两顶点沿直角导轨滑动，$C$ 是板的质心。$\theta$ 为边 $AB$ 与向下竖直方向的夹角。}\label{mp2:fig:l2e4-q}
\end{figure}
\end{minipage}\par\medskip

\begin{solution}
由图\ref{mp2:fig:l2e4-s}中的边长和方位，
\[
A=(0,l\cos\theta),\qquad B=(l\sin\theta,0),
\]
\[
D=(l\cos\theta,l\cos\theta+l\sin\theta).
\]
质心是对角线 $BD$ 的中点，所以
\[
x_C=y_C=\frac l2(\sin\theta+\cos\theta).
\]
在给定角度范围内，$\sin\theta+\cos\theta$ 从1增至 $\sqrt2$，再降至1。实际轨迹是直线 $y=x$ 上的一段：
\[
\frac l2\le x_C=y_C\le\frac l{\sqrt2}.
\]
同一内点通常对应两个板的姿态。图\ref{mp2:fig:l2e4-s}中，$C_0$ 为 $\theta=0,\pi/2$ 时的质心位置，$C_*$ 为 $\theta=\pi/4$ 时的位置。随 $\theta$ 从0增至 $\pi/2$，质心先从 $C_0$ 到 $C_*$，再沿原路返回。

\begin{figure}[H]\centering
\begin{tikzpicture}[scale=.85]
\pgfmathsetmacro{\ang}{30}\pgfmathsetmacro{\ll}{2.5}
\coordinate(O)at(0,0);\coordinate(A)at(0,{\ll*cos(\ang)});\coordinate(B)at({\ll*sin(\ang)},0);
\coordinate(D)at({\ll*cos(\ang)},{\ll*(cos(\ang)+sin(\ang))});
\coordinate(E)at({\ll*(sin(\ang)+cos(\ang))},{\ll*sin(\ang)});
\coordinate(C)at({\ll*(sin(\ang)+cos(\ang))/2},{\ll*(sin(\ang)+cos(\ang))/2});
\fill[black!5](A)--(B)--(E)--(D)--cycle;\draw[thick](A)--(B)--(E)--(D)--cycle;
\draw[axis](-.2,0)--(4.1,0)node[right]{$x$};\draw[axis](0,-.2)--(0,4.2)node[above]{$y$};
\foreach\p/\pos in{A/left,B/below,E/right,D/above,C/above left}{\fill(\p)circle(1.2pt)node[\pos]{$\p$};}
\node[below left]at(O){$O$};\draw[guide](A)--(E) (B)--(D);
\draw(0,{\ll*cos(\ang)-.6})arc(270:300:.6);\node at(.25,{\ll*cos(\ang)-.85}){$\theta$};

\begin{scope}[shift={(6.1,.25)},scale=2.25]
\draw[axis](0,0)--(1.55,0)node[right]{$x/l$};
\draw[axis](0,0)--(0,1.85)node[above]{$y/l$};
\node[below left]at(0,0){$0$};
\draw[guide](0,1)--(1,1)--(1,0);\node[below]at(1,0){$1$};
\draw[line width=1pt,domain=0:90,samples=100]plot({cos(\x)},{cos(\x)+sin(\x)});
\draw[axis,line width=.8pt,domain=35:55,samples=24]plot({cos(\x)},{cos(\x)+sin(\x)});
\fill(1,1)circle(1.2pt)node[right]{$D_0$};
\fill(0,1)circle(1.2pt)node[left]{$D_1$};
\node at(.65,1.66){$D$ 的轨迹};
\draw[line width=1.3pt](.5,.5)--({1/sqrt(2)},{1/sqrt(2)});
\fill(.5,.5)circle(1.2pt)node[below left]{$C_0$};
\fill({1/sqrt(2)},{1/sqrt(2)})circle(1.2pt)node[above left]{$C_*$};
\end{scope}
\end{tikzpicture}
\caption{例~\ref{mp2:ex:l2e4}解图：右图坐标均以 $l$ 为单位，只保留题设允许的轨迹。箭头表示 $\theta$ 增大的方向：$D$ 由 $D_0$ 到 $D_1$，$C$ 则沿短线段按 $C_0\to C_*\to C_0$ 往返。}\label{mp2:fig:l2e4-s}
\end{figure}

对参数式求导，
\[
\vct v_C=\frac{l\omega}{2}(\cos\theta-\sin\theta)(\vct e_x+\vct e_y),
\]
\[
\vct a_C=\frac l2\bigl[\alpha(\cos\theta-\sin\theta)
-\omega^2(\sin\theta+\cos\theta)\bigr](\vct e_x+\vct e_y).
\]
在 $\theta=\pi/4$ 时，$\vct v_C=0$，但若 $\omega\ne0$，则
\[
\vct a_C=-\frac{l\omega^2}{\sqrt2}(\vct e_x+\vct e_y)\ne0.
\]
此时板的速度瞬心恰在 $C$，而质心仍沿直线轨迹加速。

对顶点 $D$，消去参数得
\[
x_D^2+(y_D-x_D)^2=l^2.
\]
这是椭圆方程，实际轨迹还须满足 $0\le\theta\le\pi/2$。图\ref{mp2:fig:l2e4-s}保留的粗弧两端为 $D_0=(l,l)$、$D_1=(0,l)$，最高点在 $\theta=\pi/4$，坐标为 $(l/\sqrt2,\sqrt2l)$。完整椭圆的其余部分超出了题设的角度范围。

\end{solution}











\section{接触与滚动关联}
无滑动联系瞬时接触的两个材料点。速度条件写在同一参考系中；求加速度时还要考虑接触位置沿表面的迁移。

\subsection{材料点、接触位置与接触条件}
接触问题中，材料点与几何接触位置应分别定义。$A$ 是物体甲上当时位于接触处的材料点，$B$ 是物体乙上当时位于接触处的材料点，$Q(t)$ 则表示随时间寻找出来的几何接触位置。它们在某时刻的位置重合，下一时刻通常已不是同一对材料点接触。

设公法线单位向量为 $\vct n$。保持接触时，材料点的相对法向速度为零：
\[
(\vct v_A-\vct v_B)\cdot\vct n=0.
\]
光滑接触允许切向相对滑动。若还无滑动，才有完整的速度条件 $\vct v_A=\vct v_B$。由此不能直接推出 $\vct a_A=\vct a_B$，因为沿接触过程求导时，两侧被选中的材料点在不断更换。

对单侧接触，间隙须非负，法向支持力也不能为拉力。若列出的运动学轨迹需要负的支持力，接触阶段便已结束；若假设无滑动，还须有足够的静摩擦力。以下各题在明确给定的接触阶段内研究关联关系。

\subsection{滚动速度与接触点加速度}
图\ref{mp2:fig:l2e5-q}中的圆盘半径为 $R$，圆心速度为 $V\vct e_x$，水平带速为 $U\vct e_x$。角速度以逆时针为正，轮底相对圆心的位矢为 $-R\vct e_y$，故无滑动给出
\[
V+\omega R=U.
\]
当带固定、圆盘向右运动时，$\omega=-V/R$，转向为顺时针。若把顺时针转过的角另记为 $\theta>0$，则常写 $x_C=R\theta+\mathrm{const}$；两种写法使用的角度正方向不同。

\par\medskip\noindent\begin{minipage}{\textwidth}
\begin{example}[移动传送带上的纯滚动]\label{mp2:ex:l2e5}

半径 $R$ 的圆盘与水平传送带保持接触且无滑动。圆心速度为 $V(t)\vct e_x$，带速为 $U(t)\vct e_x$，圆心高度不变。按图\ref{mp2:fig:l2e5-q}的定义，求圆盘的 $\omega,\alpha$，任意固定轮上材料点 $P$ 的速度和加速度，以及当时接触的材料点 $A,B$ 和几何接触位置 $Q$ 的速度、加速度。最后讨论固定地面上圆心匀速的情形。

\end{example}

\begin{figure}[H]\centering
\begin{tikzpicture}[scale=1.05]
\draw[thick](-2.1,0)--(3.2,0);\draw[guide](-2.1,-.18)--(3.2,-.18);
\draw(0,1.3)circle(1.3);\fill(0,1.3)circle(1.4pt)node[below left]{$C$};
\coordinate(P)at({1.3*cos(55)},{1.3+1.3*sin(55)});
\draw(0,1.3)--(P)node[midway,left]{$R$};\fill(P)circle(1.3pt)node[above right]{$P$};
\draw[axis](0,1.3)--(1.5,1.3)node[right]{$V$};
\draw[axis](-1.3,-.9)--(.1,-.9)node[right]{$U$};
\fill(0,0)circle(1.3pt);\node[below right]at(.25,-.28){$A,B,Q$ 重合};
\draw[axis]({1.6*cos(-20)},{1.3+1.6*sin(-20)})arc(-20:30:1.6);
\node[right]at(1.65,1.9){$\omega>0$};
\end{tikzpicture}
\caption{例~\ref{mp2:ex:l2e5}题设：圆盘在水平传送带上纯滚动。$A,B$ 分别是当时接触的轮上、带上材料点，$Q$ 是几何接触位置；角速度以逆时针为正。}\label{mp2:fig:l2e5-q}
\end{figure}
\end{minipage}\par\medskip

\begin{solution}
无滑动条件及其沿滚动过程的导数给出
\[
\omega=\frac{U-V}{R},\qquad \alpha=\frac{\dot U-\dot V}{R}.
\]
对任意固定材料点 $P$，以 $\vct r=\vct r_{P/C}$ 表示它相对圆心的位置，有
\[
\vct v_P=V\vct e_x+\omega\vct k\times\vct r,
\]
\[
\vct a_P=\dot V\vct e_x+\alpha\vct k\times\vct r-\omega^2\vct r.
\]
图\ref{mp2:fig:l2e5-s}左图是固定地面时的一个瞬时速度分布。速度沿轮缘的位置改变，接触处的零速度不代表圆盘整体静止。

\begin{figure}[H]\centering
\begin{tikzpicture}[scale=.95]
\foreach\dx in{0,5.3}{
\draw[thick](\dx-1.6,0)--(\dx+2.4,0);
\draw(\dx,1)circle(1);\fill(\dx,1)circle(1.2pt)node[left]{$C$};
\fill(\dx,0)circle(1.2pt)node[below left]{$A$};\fill(\dx,2)circle(1.2pt)node[above left]{$T$};
}
\draw[axis,thick](0,1)--(1,1)node[right]{$V$};
\draw[axis,thick](0,2)--(2,2)node[right]{$2V$};\node[right]at(.25,.2){$v_A=0$};
\node at(.4,2.75){速度};
\draw[axis,thick](5.3,0)--(5.3,.75);\node[left]at(4.15,.35){$V^2/R$};
\draw[axis,thick](5.3,2)--(5.3,1.25);\node[left]at(4.15,1.65){$V^2/R$};
\node[right]at(6.45,1){$\vct a_C=0$};\node at(5.7,2.75){加速度};
\end{tikzpicture}
\caption{例~\ref{mp2:ex:l2e5}解图：固定地面、圆心匀速时的速度与加速度分别作图。轮底材料点 $A$ 速度为零，加速度却指向圆心。}\label{mp2:fig:l2e5-s}
\end{figure}

在当前接触处代入 $\vct r_A=-R\vct e_y$，得到
\[
\vct v_A=U\vct e_x,\qquad
\vct a_A=\dot U\vct e_x+\omega^2R\vct e_y.
\]
传送带作水平平移，其材料点 $B$ 满足
\[
\vct v_B=U\vct e_x,\qquad \vct a_B=\dot U\vct e_x.
\]
所以两接触材料点速度相等，加速度相差 $\omega^2R\vct e_y$。几何接触位置始终在圆心正下方，$Q=(x_C,0)$，故
\[
\vct v_Q=V\vct e_x,\qquad \vct a_Q=\dot V\vct e_x.
\]
这里分别求的是轮上材料点、带上材料点和几何接触位置的运动。

固定地面、圆心匀速时，$U=0,\dot V=0$，轮底材料点有
\[
\vct v_A=0,\qquad \vct a_A=\frac{V^2}{R}\vct e_y.
\]
轮顶 $T$ 则有 $\vct v_T=2V\vct e_x$、$\vct a_T=-(V^2/R)\vct e_y$，如图\ref{mp2:fig:l2e5-s}右图。轮缘材料点的加速度处处指向圆心，但材料点在地面参考系中的轨迹并不是圆。

\end{solution}

\par\medskip\noindent\begin{minipage}{\textwidth}
\begin{example}[相切直杆上的三种接触点]\label{mp2:ex:l2e6}

半径为 $R$ 的固定半圆柱放在水平地面上，圆心为 $O$。直杆 $AB$ 始终与它相切，$A$ 在地面上以恒速 $V>0$ 向右运动，杆与向左水平线的锐角为 $0<\alpha<\pi/2$。见图\ref{mp2:fig:l2e6-q}，记当时相切的杆上材料点为 $C$、圆柱材料点为 $D$、几何切点为 $Q$。求三者的速度和加速度，并解释为什么不能对“切点速度”直接求导来得到 $C$ 的加速度。

\end{example}

\begin{figure}[H]\centering
\begin{tikzpicture}[scale=.95]
\pgfmathsetmacro{\rr}{1.3}\pgfmathsetmacro{\aa}{30}
\pgfmathsetmacro{\ss}{\rr*cos(\aa)/sin(\aa)}
\coordinate(O)at(0,0);\coordinate(A)at({\rr/sin(\aa)},0);
\coordinate(Q)at({\rr*sin(\aa)},{\rr*cos(\aa)});
\coordinate(B)at({\rr/sin(\aa)-4.3*cos(\aa)},{4.3*sin(\aa)});
\draw[thick](-1.8,0)--(3.65,0);
\draw[thick](\rr,0)arc(0:180:\rr);
\draw[line width=1pt](A)--(B);
\fill(O)circle(1.3pt)node[below]{$O$};\fill(A)circle(1.3pt)node[below]{$A$};\node[above]at(B){$B$};
\draw[guide](O)--(Q)node[midway,left]{$R$};\fill(Q)circle(1.5pt);
\draw({\rr/sin(\aa)-.55},0)arc(180:150:.55);\node at({\rr/sin(\aa)-.85},.2){$\alpha$};
\draw[axis]({\rr/sin(\aa)},.35)--({\rr/sin(\aa)+.8},.35)node[right]{$V$};

\draw[axis](Q)--++({.6*sin(\aa)},{.6*cos(\aa)})node[above]{$\vct n$};
\draw[guide](Q)--++(0,1.25)node[above]{$C,D,Q$};
\node[anchor=west]at(4.1,1.65){$C$：杆上的材料点};
\node[anchor=west]at(4.1,.95){$D$：圆柱上的材料点};
\node[anchor=west]at(4.1,.25){$Q$：几何切点};
\end{tikzpicture}
\caption{例~\ref{mp2:ex:l2e6}题设：直杆与固定半圆柱相切。杆上材料点 $C$、圆柱上材料点 $D$ 和几何切点 $Q$ 在所示瞬间重合。}\label{mp2:fig:l2e6-q}
\end{figure}
\end{minipage}\par\medskip

\begin{solution}
取沿杆由 $A$ 指向 $B$ 的单位切向量、公法线为
\[
\vct t=(-\cos\alpha,\sin\alpha),\qquad
\vct n=(\sin\alpha,\cos\alpha).
\]
于是 $\dot{\vct t}=\dot\alpha\vct n$、$\dot{\vct n}=-\dot\alpha\vct t$。由 $OQ\perp AB$，
\[
OA=\frac R{\sin\alpha},\qquad AQ=R\cot\alpha.
\]
因为 $\dot{OA}=V$，对 $OA=R/\sin\alpha$ 求导，得到
\[
V=-R\frac{\cos\alpha}{\sin^2\alpha}\dot\alpha,
\]
即
\[
\dot\alpha=-\frac VR\frac{\sin^2\alpha}{\cos\alpha},
\]
其中 $V$ 为常量。对正弦平方与余弦倒数同时求导，
\[
\begin{aligned}
\ddot\alpha
&=-\frac VR\left(2\sin\alpha+
\frac{\sin^3\alpha}{\cos^2\alpha}\right)\dot\alpha\\
&=\frac{V^2}{R^2}\frac{\sin^3\alpha(2\cos^2\alpha+\sin^2\alpha)}{\cos^3\alpha}.
\end{aligned}
\]
利用 $2\cos^2\alpha+\sin^2\alpha=1+\cos^2\alpha$，得
\[
\ddot\alpha=\frac{V^2}{R^2}
\frac{\sin^3\alpha(1+\cos^2\alpha)}{\cos^3\alpha}.
\]
杆的绝对方位角为 $\pi-\alpha$，故它的角速度、角加速度分别是 $-\dot\alpha$、$-\ddot\alpha$。

先固定某个杆上材料点，它到 $A$ 的距离记为常量 $s_C$。其位置、速度和加速度为
\begin{align*}
\vct r_C&=\vct r_A+s_C\vct t,\\
\vct v_C&=V\vct e_x+s_C\dot\alpha\vct n,\\
\vct a_C&=s_C\ddot\alpha\vct n-s_C\dot\alpha^2\vct t.
\end{align*}
求完导数后，再在当前瞬间代入 $s_C=AQ=R\cot\alpha$。利用
$\vct e_x=\sin\alpha\vct n-\cos\alpha\vct t$，得
\[
\vct v_C=-V\cos\alpha\,\vct t,
\]
\[
\vct a_C=\frac{V^2}{R}
\left[\frac{\sin^2\alpha(1+\cos^2\alpha)}{\cos^2\alpha}\vct n
-\frac{\sin^3\alpha}{\cos\alpha}\vct t\right].
\]

\begin{figure}[H]\centering
\begin{tikzpicture}[scale=1.05]
\draw[guide](-1.9,1.1)--(1.9,-1.1);
\fill(0,0)circle(1.4pt)node[below left]{$C,D,Q$};
\draw[axis,thick](0,0)--({1.5*cos(30)},{-1.5*sin(30)})node[below right]{$\vct v_C$};
\draw[axis,thick](0,0)--({-.5*cos(30)},{.5*sin(30)})node[above left]{$\vct v_Q$};
\draw[axis](0,0)--(.6,1.03923)node[above]{$\vct n$};
\node[above left]at(-1.3,.75){$\vct t$};
\node[anchor=west]at(3.1,.8){$\vct v_D=0$};
\node[anchor=west]at(3.1,0){$v_{Cn}=v_{Dn}=v_{Qn}=0$};
\end{tikzpicture}
\caption{例~\ref{mp2:ex:l2e6}解图：取 $\alpha=\pi/6$ 示意。$\vct v_C$ 沿杆向右下，$\vct v_Q$ 沿杆向左上，大小之比为3；固定圆柱材料点 $D$ 的速度为零。}\label{mp2:fig:l2e6-s}
\end{figure}

固定圆柱上的材料点 $D$ 有 $\vct v_D=\vct a_D=0$。几何点 $Q$ 则由 $\vct r_Q=R\vct n$ 定义，因此
\[
\vct v_Q=-R\dot\alpha\vct t,
\qquad
\vct a_Q=-R\ddot\alpha\vct t-R\dot\alpha^2\vct n.
\]
图\ref{mp2:fig:l2e6-s}中 $C$ 和 $Q$ 沿杆向相反方向移动，$D$ 不动；三者的法向速度都为零，符合保持接触的条件。

\begin{figure}[H]\centering
\begin{tikzpicture}[scale=1.05]
\draw[axis,guide](0,0)--(.9,1.55885)node[above]{$\vct n$};
\draw[axis,guide](0,0)--(-1.55885,.9)node[left]{$\vct t$};
\fill(0,0)circle(1.3pt)node[below left]{$C,D,Q$};
\draw[axis,thick](0,0)--(1.25,1.29904)node[above right]{$\vct a_C$};
\draw[axis,thick](0,0)--(.75,-.72169)node[below right]{$\vct a_Q$};
\node[anchor=west]at(3,.8){$\vct a_D=0$};
\node[anchor=west]at(3,0){$a_{Cn}=7V^2/(12R)$};
\node[anchor=west]at(3,-.8){$a_{Qn}=-V^2/(12R)$};
\end{tikzpicture}
\caption{例~\ref{mp2:ex:l2e6}加速度图：仍取 $\alpha=\pi/6$，两条加速度按同一比例绘制。材料点 $C$ 与几何点 $Q$ 的法向分量不相等。}\label{mp2:fig:l2e6-a}
\end{figure}

图\ref{mp2:fig:l2e6-a}显示三种加速度的区别。取 $\alpha=\pi/6$，有
\[
\vct a_C=\frac{V^2}{R}
\left(\frac7{12}\vct n-\frac1{4\sqrt3}\vct t\right),
\]
\[
\vct a_Q=\frac{V^2}{R}
\left(-\frac1{12}\vct n-\frac7{12\sqrt3}\vct t\right),
\qquad \vct a_D=0.
\]
简式 $\vct v_C=-V\cos\alpha\vct t$ 只给出每时刻处于接触位置的杆点速度，其中的 $C$ 随时间更换。直接对它求导，所得结果便不再是此前选定材料点的加速度。因此须保持 $s_C$ 不变完成求导，再代入当前接触位置。几何切点 $Q(t)$ 的定义贯穿整个接触过程，它的速度、加速度可直接由参数式求导得到。

\end{solution}

\par\medskip\noindent\begin{minipage}{\textwidth}
\begin{example}[圆盘内滚与直线轨迹]\label{mp2:ex:l2e7}

半径为 $r$ 的小圆盘沿半径为 $R>r$ 的固定大圆内壁无滑动滚动，圆心 $C$ 以恒定角速度 $\Omega>0$ 绕大圆心 $O$ 逆时针运动。小盘边缘材料点 $P$ 起初与最右侧接触点重合，见图\ref{mp2:fig:l2e7-q}。求小盘的绝对自转角速度与 $P$ 的轨迹参数式。若在速度、加速度均非零时，$P$ 的两者始终共线，求 $R/r$ 并说明运动形式。

\end{example}

\begin{figure}[H]\centering
\begin{tikzpicture}[scale=1.05]
\pgfmathsetmacro{\RR}{2.6}\pgfmathsetmacro{\rr}{0.9}\pgfmathsetmacro{\ph}{40}
\pgfmathsetmacro{\aa}{\RR-\rr}\pgfmathsetmacro{\qq}{\aa/\rr}
\coordinate(O)at(0,0);\coordinate(C)at({\aa*cos(\ph)},{\aa*sin(\ph)});
\coordinate(K)at({\RR*cos(\ph)},{\RR*sin(\ph)});
\coordinate(P)at({\aa*cos(\ph)+\rr*cos(-\qq*\ph)},{\aa*sin(\ph)+\rr*sin(-\qq*\ph)});
\draw(O)circle(\RR);\draw[thick](C)circle(\rr);
\draw[guide](O)--(K);\draw(C)--(P);
\fill(O)circle(1.2pt)node[below left]{$O$};\fill(C)circle(1.3pt)node[above left]{$C$};
\fill(K)circle(1.3pt)node[above right]{$K$};\fill(P)circle(1.3pt)node[below right]{$P$};
\draw[guide](0,0)--(\RR,0);\draw(.65,0)arc(0:\ph:.65);\node at(.85,.28){$\phi$};
\node[above left]at({.5*\aa*cos(\ph)},{.5*\aa*sin(\ph)}){$R-r$};
\node[right]at({(\aa+.5*\rr)*cos(\ph)},{(\aa+.5*\rr)*sin(\ph)}){$r$};
\draw(O)--(-120:\RR)node[midway,left]{$R$};
\end{tikzpicture}
\caption{例~\ref{mp2:ex:l2e7}题设：小圆盘沿固定大圆内壁纯滚动，$K$ 是当时接触点。小盘上标记的材料点 $P$ 起初位于最右侧接触位置。}\label{mp2:fig:l2e7-q}
\end{figure}
\end{minipage}\par\medskip

\begin{solution}
记 $b=R-r$、$q=b/r$，公转角 $\phi=\Omega t$。圆心速度为 $b\Omega\vct e_\phi$。接触处相对圆心的位矢为 $r\vct e_r$，无滑动要求该材料点速度为零，因此
\[
b\Omega+\omega r=0,\qquad \omega=-q\Omega.
\]
这里的 $\omega$ 是在地面参考系中量得的绝对自转角速度，不是小盘相对连心线的转速。

标记点起初位于 $\vct r_{P/C}=r\vct e_x$，故
\begin{align*}
x_P&=b\cos\phi+r\cos(q\phi),\\
y_P&=b\sin\phi-r\sin(q\phi).
\end{align*}
为求速度与加速度共线的条件，改用随连心线转动的基矢计算两者。

以连心线为径向，记 $CP$ 与 $OC$ 的有向夹角为 $\psi$，则
\[
\vct r_{P/C}=r(\cos\psi\vct e_r+\sin\psi\vct e_\phi).
\]
由刚体运动式，
\[
\vct v_P=b\Omega\left[\sin\psi\vct e_r+(1-\cos\psi)\vct e_\phi\right],
\]
\[
\vct a_P=-b\Omega^2
\left[(1+q\cos\psi)\vct e_r+q\sin\psi\vct e_\phi\right].
\]
两向量的平面叉积为
\begin{align*}
(\vct v_P\times\vct a_P)\cdot\vct k
&=b^2\Omega^3\bigl[(1-\cos\psi)(1+q\cos\psi)-q\sin^2\psi\bigr]\\
&=b^2\Omega^3(1-q)(1-\cos\psi).
\end{align*}
若 $1-\cos\psi=0$，$P$ 恰在接触位置且速度为零，不属于所问的非零速度情形。因此共线条件给出 $q=1$，即 $R=2r$。

\begin{figure}[H]\centering
\begin{tikzpicture}[scale=1.05]
\pgfmathsetmacro{\RR}{2.4}\pgfmathsetmacro{\rr}{1.2}\pgfmathsetmacro{\ph}{60}
\pgfmathsetmacro{\aa}{\RR-\rr}\pgfmathsetmacro{\qq}{\aa/\rr}
\coordinate(O)at(0,0);\coordinate(C)at({\aa*cos(\ph)},{\aa*sin(\ph)});
\coordinate(K)at({\RR*cos(\ph)},{\RR*sin(\ph)});
\coordinate(P)at({\aa*cos(\ph)+\rr*cos(-\qq*\ph)},{\aa*sin(\ph)+\rr*sin(-\qq*\ph)});
\draw(O)circle(\RR);\draw[thick](C)circle(\rr);
\draw[guide](O)--(K);\draw(C)--(P);
\fill(O)circle(1.2pt)node[below left]{$O$};\fill(C)circle(1.3pt)node[above left]{$C$};
\fill(K)circle(1.3pt)node[above right]{$K$};\fill(P)circle(1.3pt)node[below right]{$P$};
\draw[guide](0,0)--(\RR,0);\draw(.65,0)arc(0:\ph:.65);\node at(.85,.28){$\phi$};
\node[above left]at({.5*\aa*cos(\ph)},{.5*\aa*sin(\ph)}){$R-r$};
\node[right]at({(\aa+.5*\rr)*cos(\ph)},{(\aa+.5*\rr)*sin(\ph)}){$r$};


\draw[line width=1pt](-\RR,0)--(\RR,0);
\draw[guide](P)--({\RR*cos(\ph)},-.8);
\draw[axis]({\RR*cos(\ph)},-.35)--({\RR*cos(\ph)-1.4},-.35)node[below left]{$\vct v_P$};
\draw[axis]({\RR*cos(\ph)},-.8)--({\RR*cos(\ph)-.8},-.8)node[left]{$\vct a_P$};
\end{tikzpicture}
\caption{例~\ref{mp2:ex:l2e7}解图：$R=2r$ 时，$P$ 始终位于水平直径上。图示 $\phi=\pi/3$，速度和加速度的方向都向左；在别的位置，两者也可能反向。}\label{mp2:fig:l2e7-s}
\end{figure}

代入参数式，图\ref{mp2:fig:l2e7-s}中的点满足
\[
x_P=2r\cos\phi,\qquad y_P=0,
\]
\[
\vct v_P=-2r\Omega\sin\phi\,\vct e_x,
\qquad
\vct a_P=-2r\Omega^2\cos\phi\,\vct e_x.
\]
所以 $P$ 沿一条直径作简谐运动。速度与加速度共线时可以同向，也可以反向；在直径端点速度为零，在圆心处加速度为零，不能给零向量指定方向。若标记点换为另一初始轮缘点，所得直线直径的方位也随之改变。

\end{solution}











\section{复杂机构与转动参考系}
交点和瞬心通常是随构形变化的几何位置，未必对应固定材料点。位形约束确定其轨迹，相对运动关系再区分机构的运动与坐标系的转动。

\subsection{转动基矢的求导}
把一个矢量写成随坐标系转动的基矢分量，求导时必须同时求基矢的导数。以平面转动为例，
\[
\vct e_1=(\cos\theta,\sin\theta),\quad
\vct e_2=(-\sin\theta,\cos\theta),
\]
便有 $\dot{\vct e}_1=\omega\vct e_2$、$\dot{\vct e}_2=-\omega\vct e_1$。图\ref{mp2:fig:l2frame}左图的两基矢随坐标系转动；即使某矢量在这组基矢中的分量恒定，它在地面系中的方向仍可改变。

对一般以角速度 $\vct\Omega$ 转动的坐标系，任一矢量 $\vct A$ 满足输运公式
\[
\left(\frac{\dd\vct A}{\dd t}\right)_{\!\text{地}}
=\left(\frac{\dd\vct A}{\dd t}\right)_{\!\text{转}}
+\vct\Omega\times\vct A.
\]
下标“转”表示只对转动系中的分量求导。若其原点 $O$ 还在平移，将质点位置写为 $\vct r_P=\vct r_O+\vct\rho$，并定义
\[
\vct v_{\rm rel}=\left(\frac{\dd\vct\rho}{\dd t}\right)_{\!\text{转}},
\qquad
\vct a_{\rm rel}=\left(\frac{\dd\vct v_{\rm rel}}{\dd t}\right)_{\!\text{转}}.
\]
则绝对速度为
\[
\vct v_P=\vct v_O+\vct\Omega\times\vct\rho+\vct v_{\rm rel}.
\]
再次在地面系求导。两个含转动的项分别产生
\[
\begin{aligned}
\left(\frac{\dd}{\dd t}(\vct\Omega\times\vct\rho)\right)_{\!\text{地}}
&=\dot{\vct\Omega}\times\vct\rho
+\vct\Omega\times(\vct v_{\rm rel}+\vct\Omega\times\vct\rho),\\
\left(\frac{\dd\vct v_{\rm rel}}{\dd t}\right)_{\!\text{地}}
&=\vct a_{\rm rel}+\vct\Omega\times\vct v_{\rm rel}.
\end{aligned}
\]
与原点加速度相加，得到
\begin{equation}\label{mp2:eq:l2-transport}
\begin{split}
\vct a_P={}&\vct a_O+\dot{\vct\Omega}\times\vct\rho
+\vct\Omega\times(\vct\Omega\times\vct\rho)\\
&+2\vct\Omega\times\vct v_{\rm rel}+\vct a_{\rm rel}.
\end{split}
\end{equation}
系数2来自两个相同的交叉项，分别在对 $\vct\Omega\times\vct\rho$ 和 $\vct v_{\rm rel}$ 求导时出现。这里各项都是加速度；若要在转动系中写动力学方程，还须结合牛顿定律处理惯性力。

\begin{figure}[H]\centering
\begin{tikzpicture}[scale=.9]
\draw[axis](-.2,0)--(3.4,0)node[right]{$x$};\draw[axis](0,-.2)--(0,2.8)node[above]{$y$};
\draw[axis,thick](0,0)--({2.3*cos(30)},{2.3*sin(30)})node[right]{$\vct e_1$};
\draw[axis,thick](0,0)--({2.3*cos(120)},{2.3*sin(120)})node[above]{$\vct e_2$};
\draw(.65,0)arc(0:30:.65);\node at(.9,.22){$\theta$};\node[below left]at(0,0){$O$};
\begin{scope}[shift={(6.6,0)}]
\draw[thick](-.3,-.1732)--(2.8,1.6166);
\coordinate(P)at(1.73205,1);\fill(P)circle(1.4pt)node[below right]{$P$};
\draw[axis](P)--++(.86603,.5)node[right]{$\dot s\vct e_r$};
\draw[axis](P)--++(-.6,1.03923)node[above]{$2\omega\dot s\vct e_\theta$};
\fill(0,0)circle(1.3pt)node[below left]{$O$};\node[below]at(.85,.4){$s$};
\end{scope}
\end{tikzpicture}
\caption{转动基矢与相对运动。左图的两个单位基矢随坐标系一起转动。右图取 $\omega>0,\dot s>0$，沿槽箭头表示相对速度，横向箭头表示加速度项 $2\omega\dot s\vct e_\theta$。}\label{mp2:fig:l2frame}
\end{figure}

若原点固定，质点沿转动直槽运动，取 $\vct\rho=s\vct e_r$，则
\[
\vct v_P=\dot s\vct e_r+s\omega\vct e_\theta,
\]
\begin{equation}\label{mp2:eq:l2-polar-a}
\vct a_P=(\ddot s-s\omega^2)\vct e_r
+(s\alpha+2\dot s\omega)\vct e_\theta.
\end{equation}
式中 $s$ 可按槽方向取有符号坐标。$\alpha$ 是槽的绝对角加速度；若换到另一转动构件上，必须重新计算相对角速度，并把不随该构件固定的支点坐标一同变换。

\subsection{闭合机构中的速度与加速度}
多杆机构可以从固定铰出发逐段写刚体关系，也可以先列矢量闭合式再求导。同一个活动铰沿两条链算出的速度、加速度必须一致，这就给出了求从动杆转动量的方程。

\par\medskip\noindent\begin{minipage}{\textwidth}
\begin{example}[四连杆的两条运动链]\label{mp2:ex:l2e8}

三根杆通过 $A,B,C,D$ 四处铰接，$A,D$ 固定。图\ref{mp2:fig:l2e8-q}给出某瞬间的构形：$A=(0,0)$、$B=(0,l)$、$C=(l,2l)$、$D=(3l,0)$。杆 $AB$ 以恒定角速度大小 $\Omega>0$ 顺时针转动。求这一瞬间的 $\vct v_C,\vct a_C$，以及杆 $BC,DC$ 的绝对角速度和角加速度，后两杆均取逆时针为正。

\end{example}

\begin{figure}[H]\centering
\begin{tikzpicture}[scale=1.05]
\pgfmathsetmacro{\ll}{1.2}\pgfmathsetmacro{\kk}{2}
\coordinate(A)at(0,0);\coordinate(B)at(0,\ll);\coordinate(C)at(\ll,{2*\ll});\coordinate(D)at({(1+\kk)*\ll},0);
\draw[thick](-.4,0)--({(1+\kk)*\ll+.4},0);
\draw[line width=1pt](A)--(B)--(C)--(D);
\foreach\p/\pos in{A/below,B/left,C/above,D/below}{\fill(\p)circle(1.4pt)node[\pos]{$\p$};}
\draw[guide](C)--(\ll,0) (0,{2*\ll})--(C);
\draw[|-|](-.55,0)--(-.55,\ll)node[midway,left]{$l$};
\draw[|-|](-.55,\ll)--(-.55,{2*\ll})node[midway,left]{$l$};
\draw[|-|](0,-.6)--(\ll,-.6)node[midway,fill=white]{$l$};
\draw[|-|](\ll,-.6)--({(1+\kk)*\ll},-.6)node[midway,fill=white]{$2l$};
\draw[axis](70:.7)arc(70:20:.7);\node at(.95,.55){$\Omega$};
\end{tikzpicture}
\caption{例~\ref{mp2:ex:l2e8}题设：两固定铰为 $A,D$。所示构形中 $AB=l$，$B=(0,l)$、$C=(l,2l)$、$D=(3l,0)$。}\label{mp2:fig:l2e8-q}
\end{figure}
\end{minipage}\par\medskip

\begin{solution}
杆 $AB$ 绕固定点 $A$ 匀速转动，故
\[
\vct v_B=\Omega l\vct e_x,\qquad
\vct a_B=-\Omega^2l\vct e_y.
\]
以 $\omega_2,\alpha_2$ 表示 $BC$ 的转动，以 $\omega_3,\alpha_3$ 表示 $DC$ 的转动。两相对位矢是
\[
\vct r_{C/B}=l(1,1),\qquad \vct r_{C/D}=l(-2,2).
\]
图\ref{mp2:fig:l2e8-s}左图中，$C$ 的速度从两条链写为
\[
\frac{\vct v_C}{l}=(\Omega-\omega_2,\omega_2)
=(-2\omega_3,-2\omega_3).
\]
故
\[
\omega_2=\frac\Omega2,\qquad \omega_3=-\frac\Omega4,
\qquad \vct v_C=\frac{\Omega l}{2}(\vct e_x+\vct e_y).
\]

\begin{figure}[H]\centering
\begin{tikzpicture}[scale=.9]
\pgfmathsetmacro{\ll}{1.2}\pgfmathsetmacro{\kk}{2}
\coordinate(A)at(0,0);\coordinate(B)at(0,\ll);\coordinate(C)at(\ll,{2*\ll});\coordinate(D)at({(1+\kk)*\ll},0);
\draw[thick](-.4,0)--({(1+\kk)*\ll+.4},0);
\draw[line width=1pt](A)--(B)--(C)--(D);
\foreach\p/\pos in{A/below,B/left,C/above,D/below}{\fill(\p)circle(1.4pt)node[\pos]{$\p$};}
\draw[guide](C)--(\ll,0) (0,{2*\ll})--(C);
\draw[|-|](-.55,0)--(-.55,\ll)node[midway,left]{$l$};
\draw[|-|](-.55,\ll)--(-.55,{2*\ll})node[midway,left]{$l$};
\draw[|-|](0,-.6)--(\ll,-.6)node[midway,fill=white]{$l$};
\draw[|-|](\ll,-.6)--({(1+\kk)*\ll},-.6)node[midway,fill=white]{$2l$};
\draw[axis](70:.7)arc(70:20:.7);\node at(.95,.55){$\Omega$};

\draw[axis,thick](B)--++(1,0)node[midway,below]{$\vct v_B$};
\draw[axis,thick](C)--++(.5,.5)node[above]{$\vct v_C$};
\node at(1.6,4){速度};
\begin{scope}[shift={(6.6,2.4)}]
\fill(0,0)circle(1.3pt)node[right]{$C$};
\draw[axis](0,0)--(-1.25,0)node[midway,above]{$5\Omega^2l/8$};
\draw[axis](-1.25,0)--(-1.25,-1.75)node[midway,left]{$7\Omega^2l/8$};
\draw[axis,thick](0,0)--(-1.25,-1.75)node[midway,right]{$\vct a_C$};
\end{scope}
\node at(6,4){加速度};
\end{tikzpicture}
\caption{例~\ref{mp2:ex:l2e8}解图：左图表示速度，右图将 $C$ 的加速度分解为向左、向下两项，旁注为各分量的大小。}\label{mp2:fig:l2e8-s}
\end{figure}

各杆的法向项使用自身的绝对角速度。加速度的两种表达是
\[
\frac{\vct a_C}{l}
=(-\alpha_2-\omega_2^2,-\Omega^2+\alpha_2-\omega_2^2),
\]
\[
\frac{\vct a_C}{l}
=(-2\alpha_3+2\omega_3^2,-2\alpha_3-2\omega_3^2).
\]
代入刚求得的角速度，逐分量相等：
\begin{align*}
-\alpha_2-\frac{\Omega^2}{4}&=-2\alpha_3+\frac{\Omega^2}{8},\\
-\Omega^2+\alpha_2-\frac{\Omega^2}{4}&=-2\alpha_3-\frac{\Omega^2}{8}.
\end{align*}
解得 $\alpha_2=\alpha_3=3\Omega^2/8$，于是
\[
\vct a_C=-\frac58\Omega^2l\vct e_x-\frac78\Omega^2l\vct e_y.
\]
右图按分量画出合加速度。当前构形中，两从动杆的角速度大小不同、转向相反，角加速度却相等；这一结果只适用于题给尺寸和当前瞬间。

\end{solution}

\par\medskip\noindent\begin{minipage}{\textwidth}
\begin{example}[曲柄经斜绳牵引重物]\label{mp2:ex:l2e9}

固定点 $O,B$ 的竖直距离为 $H$，曲柄 $OA=R<H$，以恒定角速度大小 $\Omega>0$ 顺时针转动。轻绳从 $A$ 经固定小滑轮 $B,C$ 后竖直悬挂 $D$，见图\ref{mp2:fig:l2e9-q}。忽略滑轮半径，绳始终绷紧。令 $OA$ 与向上竖直方向的夹角为 $\theta$，$0<\theta<\pi$，$D$ 的向下坐标为 $y$。求 $v_D,a_D$，并求这一阶段 $D$ 的最大速率及相应构形。

\end{example}

\begin{figure}[H]\centering
\begin{tikzpicture}[scale=1.05]
\pgfmathsetmacro{\rr}{1.4}\pgfmathsetmacro{\ang}{50}
\coordinate(O)at(0,0);\coordinate(B)at(0,3.6);\coordinate(C)at(-1.8,3.6);
\coordinate(A)at({\rr*sin(\ang)},{\rr*cos(\ang)});
\draw[line width=1pt](O)--(A);\draw[thick](A)--(B)--(C)--(-1.8,1.3);
\draw[fill=white](B)circle(.08);\draw[fill=white](C)circle(.08);
\draw[fill=white](-2,1.05)rectangle(-1.6,1.55);\node[left]at(-2,1.3){$D$};
\fill(O)circle(1.4pt)node[below]{$O$};\fill(A)circle(1.3pt)node[below right]{$A$};
\node[above right]at(B){$B$};\node[above left]at(C){$C$};
\node[right]at({.5*\rr*sin(\ang)},{.5*\rr*cos(\ang)}){$R$};
\draw[guide](O)--(B);\draw(0,.6)arc(90:40:.6);\node at(.32,.76){$\theta$};
\draw[|-|](-.65,0)--(-.65,3.6)node[midway,left]{$H$};
\draw[axis](-2.7,3.6)--(-2.7,1.3)node[midway,left]{$y$};
\end{tikzpicture}
\caption{例~\ref{mp2:ex:l2e9}题设：曲柄端 $A$ 通过绳段 $AB$、固定小滑轮 $B,C$ 牵引 $D$。$OB=H>R$，$OA=R$，$y$ 竖直向下。}\label{mp2:fig:l2e9-q}
\end{figure}
\end{minipage}\par\medskip

\begin{solution}
取 $O=(0,0)$、$B=(0,H)$，则 $A=(R\sin\theta,R\cos\theta)$，$\dot\theta=\Omega$。变长绳段为
\[
s=AB=\sqrt{H^2+R^2-2HR\cos\theta}.
\]
其余固定绳段并入常数，绳长给出 $s+y=L$。先对 $s^2=H^2+R^2-2HR\cos\theta$ 求导，有
\[
2s\dot s=2HR\Omega\sin\theta,\qquad
\dot s=\frac{HR\Omega\sin\theta}{s}.
\]
再次求导时，分子与分母都随时间改变：
\[
\ddot s=HR\Omega\left(\frac{\Omega\cos\theta}{s}
-\frac{\sin\theta\dot s}{s^2}\right).
\]
所以
\[
v_D=-\dot s=-\frac{HR\Omega\sin\theta}{s},
\]
\[
a_D=-HR\Omega^2
\left(\frac{\cos\theta}{s}-\frac{HR\sin^2\theta}{s^3}\right).
\]
在给定阶段内 $\sin\theta>0$，$D$ 向上运动；其加速度方向随构形改变。

\begin{figure}[H]\centering
\begin{tikzpicture}[scale=1.05]
\pgfmathsetmacro{\rr}{1.5}\pgfmathsetmacro{\ang}{60}
\coordinate(O)at(0,0);\coordinate(B)at(0,3);
\coordinate(A)at({\rr*sin(\ang)},{\rr*cos(\ang)});
\draw[guide](O)circle(\rr);\draw[guide](O)--(B);
\draw[line width=1pt](O)--(A);\draw[thick](A)--(B);
\fill(O)circle(1.3pt)node[below]{$O$};\fill(B)circle(1.3pt)node[above]{$B$};\fill(A)circle(1.3pt)node[below left]{$A$};
\draw[axis,thick](A)--++(.65,-1.12583)node[right]{$\vct v_A$};
\draw({\rr*sin(\ang)-.173205},{\rr*cos(\ang)-.1})--({\rr*sin(\ang)-.273205},{\rr*cos(\ang)+.073205})--({\rr*sin(\ang)-.1},{\rr*cos(\ang)+.173205});
\draw(0,.65)arc(90:30:.65);\node at(.35,.75){$\theta_*$};
\node[anchor=west]at(3,2){$\cos\theta_*=R/H$};
\node[anchor=west]at(3,1.1){$|v_D|_{\max}=R\Omega$};
\node[anchor=west]at(3,.2){$a_D=0$};
\end{tikzpicture}
\caption{例~\ref{mp2:ex:l2e9}解图：图示 $H=2R$，最大牵引速率出现在 $OA\perp AB$ 时；此时曲柄端速度沿绳段 $BA$。}\label{mp2:fig:l2e9-s}
\end{figure}

最大速率处 $\ddot s=0$，故
\[
\cos\theta\,(H^2+R^2-2HR\cos\theta)-HR\sin^2\theta=0.
\]
记 $c=\cos\theta$，整理为
\[
HRc^2-(H^2+R^2)c+HR=0.
\]
两根为 $R/H$ 与 $H/R$。由于 $H>R$，只有
\[
\cos\theta_* =\frac RH
\]
位于物理范围内。在该位置，$s=\sqrt{H^2-R^2}$，所以
\[
|v_D|_{\max}=R\Omega.
\]
图\ref{mp2:fig:l2e9-s}中 $OA\perp AB$，曲柄端的切向速度恰沿绳段，全部传递为绳长变化率。也可用 $|\dot s|\le|\vct v_A|=R\Omega$ 先得到这个上界，再由垂直条件判断何时取到等号。这里的 $a_D=0$ 只是速率达到极值的瞬间，不表示重物在整个运动中匀速。

若改用斜绳$AB$与竖直方向的锐角$\alpha$描述构形，几何关系为
\[
 \sin\alpha=\frac{R\sin\theta}{s},\qquad
 \cos\alpha=\frac{H-R\cos\theta}{s},\qquad
 0\le\alpha\le\arcsin(R/H).
\]
因此 $v_D=-H\Omega\sin\alpha$。同一个 $\alpha$ 一般对应两个曲柄位置，加速度还与所在分支有关。记 $\Delta=\sqrt{R^2-H^2\sin^2\alpha}$。由三角形的边长关系 $R^2=H^2+s^2-2Hs\cos\alpha$，解得 $s=H\cos\alpha+\eta\Delta$，$\eta=\pm1$。再对
$\tan\alpha=R\sin\theta/(H-R\cos\theta)$ 求导，
\[
\sec^2\alpha\,\dot\alpha
=\frac{R\Omega[\cos\theta(H-R\cos\theta)-R\sin^2\theta]}
{(H-R\cos\theta)^2}.
\]
用 $\cos\alpha=(H-R\cos\theta)/s$ 约去因子，再用上述二次关系，
\[
\dot\alpha=\frac{R\Omega(H\cos\theta-R)}{s^2}
=\frac{\Omega(H\cos\alpha-s)}s
=-\eta\frac{\Omega\Delta}s.
\]
代入 $a_D=-H\Omega\cos\alpha\dot\alpha$，两支加速度为
\[
 a_D=\eta\frac{H\Omega^2\cos\alpha\,\Delta}{s},\qquad\eta=\pm1.
\]
第一阶段从$\theta=0$起对应$\eta=-1$；通过斜绳的切线构形后转为$\eta=+1$。图\ref{mp2:fig:l2e9-s}的最大速率构形恰有$\Delta=0$，两支在此相接，$a_D=0$。给定$\alpha$可以确定速度；加速度还取决于曲柄所在的构形分支。
\end{solution}

\par\medskip\noindent\begin{minipage}{\textwidth}
\begin{example}[两转动直槽的几何交点]\label{mp2:ex:l2e10}

两无限长直槽分别绕固定点 $O=(0,0)$、$A=(0,-h)$ 转动，其绝对方位角为 $\theta_1,\theta_2$，交点为 $P$。各角速度 $\omega_1,\omega_2$ 和角加速度 $\alpha_1,\alpha_2$ 均以逆时针为正。某瞬间 $\theta_1=0$、$\theta_2=\pi/4$，如图\ref{mp2:fig:l2e10-q}。求 $P$ 的速度和加速度。$P$ 是两个槽的交点，不是任一槽上的固定材料点。

\end{example}

\begin{figure}[H]\centering
\begin{tikzpicture}[scale=1.05]
\coordinate(O)at(0,0);\coordinate(A)at(0,-2);\coordinate(P)at(2,0);
\draw[line width=1pt](-.5,0)--(4,0)node[right]{槽1};
\draw[line width=1pt](-.3,-2.3)--(3.7,1.7)node[above]{槽2};
\draw[guide](O)--(A);\draw[|-|](-1.25,0)--(-1.25,-2)node[midway,left]{$h$};
\fill(O)circle(1.4pt)node[above left]{$O$};\fill(A)circle(1.4pt)node[below left]{$A$};\fill(P)circle(1.4pt)node[below right]{$P$};
\draw[guide](A)--(1.1,-2);\draw(.55,-2)arc(0:45:.55);\node at(.85,-1.7){$\pi/4$};
\draw[axis](30:.75)arc(30:85:.75)node[above]{$\omega_1$};
\draw[axis]({1.1*cos(60)},{-2+1.1*sin(60)})arc(60:115:1.1)node[left]{$\omega_2$};
\end{tikzpicture}
\caption{例~\ref{mp2:ex:l2e10}题设：两无限长直槽分别绕固定点 $O,A$ 转动，交点为 $P$。图示 $\theta_1=0$、$\theta_2=\pi/4$ 的瞬间。}\label{mp2:fig:l2e10-q}
\end{figure}
\end{minipage}\par\medskip

\begin{solution}
先保留邻近时刻的位置关系，不能把图中的特定角度提前代入再求导。写
\[
\vct r_P=s(\cos\theta_1,\sin\theta_1),
\qquad \delta=\theta_2-\theta_1.
\]
用 $\vct r_P-\vct r_A$ 与槽2平行，或在三角形 $OAP$ 中作投影，得
\[
s=\frac{h\cos\theta_2}{\sin(\theta_2-\theta_1)}.
\]
图\ref{mp2:fig:l2e10-s}标出这一位置关系。当两槽趋于平行时，交点可远离铰点；分母为零的构形必须另外判断是否还有唯一交点。

\begin{figure}[H]\centering
\begin{tikzpicture}[scale=.95]
\pgfmathsetmacro{\th}{20}\pgfmathsetmacro{\ph}{50}\pgfmathsetmacro{\hh}{2}
\pgfmathsetmacro{\ss}{\hh*cos(\ph)/sin(\ph-\th)}
\coordinate(O)at(0,0);\coordinate(A)at(0,-\hh);\coordinate(P)at({\ss*cos(\th)},{\ss*sin(\th)});
\draw[axis](O)--({3.4*cos(\th)},{3.4*sin(\th)})node[right]{$x'$};
\draw[axis](O)--({1.6*cos(\th+90)},{1.6*sin(\th+90)})node[above]{$y'$};
\draw[thick](A)--({4.3*cos(\ph)},{-\hh+4.3*sin(\ph)});
\draw[guide](O)--(A) (A)--(1,-\hh) (O)--(1.1,0);
\fill(O)circle(1.3pt)node[left]{$O$};\fill(A)circle(1.3pt)node[below left]{$A$};\fill(P)circle(1.3pt)node[above]{$P$};
\draw(1.3,0)arc(0:\th:1.3);\node[font=\footnotesize]at(.95,.16){$\theta_1$};
\draw(.5,-\hh)arc(0:\ph:.5);\node at(.85,{-\hh+.35}){$\theta_2$};
\node[above left]at({.5*\ss*cos(\th)},{.5*\ss*sin(\th)}){$s$};
\node[anchor=west]at(4.2,1){$s=h\cos\theta_2/\sin\delta$};
\node[anchor=west]at(4.2,.15){$A'=(-h\sin\theta_1,-h\cos\theta_1)$};
\end{tikzpicture}
\caption{例~\ref{mp2:ex:l2e10}解图：在邻近构形中，$\delta=\theta_2-\theta_1$。地面上固定的点 $A$ 在随槽1转动的坐标系中运动。}\label{mp2:fig:l2e10-s}
\end{figure}

为保留求导中的各项，令 $f=\cos\theta_2$、$g=\sin\delta$，将位置式写成 $sg=hf$。连续求导，得到
\[
\dot s\,g+s\dot g=h\dot f,
\qquad
\ddot s\,g+2\dot s\dot g+s\ddot g=h\ddot f.
\]
在所给瞬间，$s=h$、$f=g=1/\sqrt2$，且
\[
\begin{aligned}
\dot f&=-\frac{\omega_2}{\sqrt2},&
\ddot f&=-\frac{\alpha_2+\omega_2^2}{\sqrt2},\\
\dot g&=\frac{\omega_2-\omega_1}{\sqrt2},&
\ddot g&=\frac{\alpha_2-\alpha_1-(\omega_2-\omega_1)^2}{\sqrt2}.
\end{aligned}
\]
代入一阶式，得到
\[
\dot s=h(\omega_1-2\omega_2),
\]
再代入二阶式，
\[
\ddot s=-h(\alpha_2+\omega_2^2)-2\dot s(\omega_2-\omega_1)
-h[\alpha_2-\alpha_1-(\omega_2-\omega_1)^2].
\]
将一阶结果代入并合并，得到
\[
\ddot s=h(\alpha_1-2\alpha_2+3\omega_1^2-8\omega_1\omega_2+4\omega_2^2).
\]
这里的位置式也可直接由两条直线约束联立求得：
\[
y\cos\theta_1-x\sin\theta_1=0,
\qquad (y+h)\cos\theta_2-x\sin\theta_2=0.
\]

选择随槽1转动的坐标系，$P$ 的相对运动沿槽，应用式\eqref{mp2:eq:l2-polar-a}。由于当前 $\vct e_r=\vct e_x$、$\vct e_\theta=\vct e_y$，
\[
\vct v_P=h(\omega_1-2\omega_2)\vct e_x+h\omega_1\vct e_y,
\]
\begin{align*}
a_{Px}&=h(\alpha_1-2\alpha_2+2\omega_1^2-8\omega_1\omega_2+4\omega_2^2),\\
a_{Py}&=h(\alpha_1+2\omega_1^2-4\omega_1\omega_2).
\end{align*}

\begin{figure}[H]\centering
\begin{tikzpicture}[scale=1.05]
\coordinate(O)at(0,0);\coordinate(A)at(0,-2);\coordinate(P)at(2,0);\coordinate(M)at(1,-1);
\draw[guide](M)circle(1.414214);\draw[thick](-.2,0)--(2.7,0) (-.2,-2.2)--(2.8,.8);
\fill(O)circle(1.3pt)node[above left]{$O$};\fill(A)circle(1.3pt)node[below left]{$A$};\fill(P)circle(1.3pt)node[right]{$P$};\fill(M)circle(1.3pt)node[below right]{$M$};
\draw[axis](P)--(1,0)node[midway,above]{$a_r$};
\draw[axis](1,0)--(M)node[midway,left]{$a_\theta$};
\draw[axis,thick](P)--(M)node[midway,right]{$\vct a_P$};
\draw[axis](P)--(1,1)node[above]{$\vct v_P$};
\node[anchor=west]at(4,.1){$a_r=-2h\omega^2$};
\node[anchor=west]at(4,-.75){$a_\theta=2\omega\dot s=-2h\omega^2$};
\end{tikzpicture}
\caption{例~\ref{mp2:ex:l2e10}加速度图：$\omega_1=\omega_2=\omega$ 为常量时，$\dot s=-h\omega$。向下的分量来自 $2\omega\dot s$，总加速度指向圆轨迹中心 $M$。}\label{mp2:fig:l2e10-a}
\end{figure}

图\ref{mp2:fig:l2e10-a}取 $\omega_1=\omega_2=\omega$ 且两者恒定，便有
\[
\vct v_P=(-h\omega,h\omega),\qquad
\vct a_P=(-2h\omega^2,-2h\omega^2).
\]
相对角速度虽然为零，$P$ 在槽1中的相对速度仍为 $\dot s=-h\omega$，所以横向项 $2\omega\dot s$ 不为零。原因是地面上固定的 $A$ 在该转动系中运动，其坐标为
\[
A'=(-h\sin\theta_1,-h\cos\theta_1).
\]
只把两槽角速度相减，却仍把 $A$ 当成转动系中的固定点，会漏掉这一运动。

\end{solution}











\section{接触点迁移与交点运动}


\subsection{材料标签与几何接触位置的求导}
接触处重合的材料点与几何位置需要用不同的变量追踪。设刚体$A$上的材料标签为$\lambda$，其位置写作$\vct r_A(\lambda,t)$。固定$\lambda$求导得到该材料点的速度与加速度；若接触位置在刚体表面移动，则它对应的标签是$\lambda(t)$，几何位置为
\[\vct r_P(t)=\vct r_A(\lambda(t),t).
\]
链式求导给
\[\dot{\vct r}_P=\vct v_A+\vct r_{A,\lambda}\dot\lambda,
\]
其中$\vct v_A=\partial_t\vct r_A$是当前接触材料点的速度。第二项沿表面切向，记录接触位置相对刚体的迁移；它不属于该材料点自身的速度。

若另一刚体$B$也在$P$接触，使用其材料标签$\mu(t)$，同一个几何位置还须满足
\[\dot{\vct r}_P=\vct v_B+\vct r_{B,\mu}\dot\mu.
\]
保持接触只要求法向速度相等；无滑动再要求材料点的切向速度相等。即使$\vct v_A=\vct v_B$，两者也未必等于$\dot{\vct r}_P$。图\ref{mp2:fig:series-contact-label}把有滑动时的材料点速度与几何点速度分别画出，它们有共同的法向分量。

\begin{figure}[H]\centering
\begin{tikzpicture}[>=Stealth,line cap=round,line join=round]
\draw[curve,domain=-1.7:1.7]plot(\x,{.28*(\x)^2});
\draw[curve,domain=-1.7:1.7]plot(\x,{-.22*(\x)^2});
\draw[guide](-2,0)--(2,0)node[right]{$\vct t$};\draw[axis](0,0)--(0,1.7)node[above]{$\vct n$};
\fill(0,0)circle(1.7pt)node[below right]{$P$};
\draw[->,thick](0,0)--(1.1,.7)node[right]{$\vct v_A$};
\draw[->,thick](0,0)--(-.8,.7)node[left]{$\vct v_B$};
\draw[->](0,0)--(.2,.7)node[above right]{$\dot{\vct r}_P$};
\node[right]at(1.7,.82){$A$};\node[right]at(1.7,-.64){$B$};
\end{tikzpicture}
\caption{同一接触位置的三个追踪对象；图示材料点间允许切向滑动，法向速度相同。}\label{mp2:fig:series-contact-label}
\end{figure}

对几何位置再求一次导数，须沿变化的标签 $\lambda(t)$ 对两项各求导：
\[
\begin{aligned}
\frac{\dd\vct v_A}{\dd t}
&=\vct a_A+\vct r_{A,t\lambda}\dot\lambda,\\
\frac{\dd}{\dd t}(\vct r_{A,\lambda}\dot\lambda)
&=(\vct r_{A,\lambda t}+\vct r_{A,\lambda\lambda}\dot\lambda)\dot\lambda
+\vct r_{A,\lambda}\ddot\lambda.
\end{aligned}
\]
这里 $\vct a_A=\partial_{tt}\vct r_A$ 是固定材料标签求得的加速度。合并混合偏导后，
\[\ddot{\vct r}_P=\vct a_A
+2\vct r_{A,t\lambda}\dot\lambda
+\vct r_{A,\lambda\lambda}\dot\lambda^2
+\vct r_{A,\lambda}\ddot\lambda.
\]
后三项依次来自表面速度随标签的变化、表面曲率和标签迁移加速度。跟踪固定材料点时，$\dot\lambda=\ddot\lambda=0$，这些项全部消失；跟踪几何接触位置时则须保留。因此两种追踪方式得到的加速度一般不同。

在下文的转杆夹圆柱模型中，$M,N,P$瞬时重合。$M,N$分别是两个刚体的材料点，$P$则由相切条件确定。先用位置约束求$P(t)$，再分别用两个刚体的平移、转动关系求$M,N$，才能把三种加速度分清。是否纯滚动还应按题设逐个接触界面判断。


\par\medskip\noindent\begin{minipage}{\textwidth}
\begin{example}[转动杆夹圆柱的四种运动]\label{mp2:rev:wedgecyl}
固定水平杆$OA$与绕$O$顺时针匀速转动的杆$OB$之间夹一个半径$r$的圆柱，转动角速度大小为$\omega$。圆柱在$OA$上纯滚动，与$OB$保持接触，见图\ref{mp2:fig:rev:wedgecyl-q}。当两杆夹角为$60^\circ$时，求圆心$C$、圆柱上接触材料点$M$、杆上接触材料点$N$及几何接触位置$P$的速度和加速度。$M,N,P$在该时刻重合；不要求圆柱与$OB$无滑动。
\end{example}

\begin{figure}[H]\centering
\begin{tikzpicture}[>=Stealth,line cap=round,line join=round]
\coordinate(O)at(0,0);\coordinate(C)at(1.732,1);\coordinate(P)at(.866,1.5);\draw[thick](O)--(4,0)node[right]{$A$};\draw[thick](O)--(1.7,2.944)node[above]{$B$};\draw[curve](C)circle(1);\fill(C)circle(1.5pt)node[right]{$C$};\fill(P)circle(1.5pt)node[above left]{$M,N,P$};\draw[guide](C)--(P)node[midway,right]{$r$};\draw(.6,0)arc(0:60:.6);\node at(.75,.3){$60^\circ$};\node[below left]at(O){$O$};\draw[axis](1.45,2.45)arc(60:25:1.2);\node[right]at(1.9,2.4){$\omega$};
\end{tikzpicture}
\caption{圆柱同时与两杆相切，上接触处的材料点与几何位置须分开求导}\label{mp2:fig:rev:wedgecyl-q}
\end{figure}
\end{minipage}\par\medskip

\begin{solution}
以$O$为原点，水平向右、竖直向上为正。先保留一般夹角$\theta$，$\dot\theta=-\omega$。相切几何给出
\[
 C=(s,r),\qquad OP=s=r\cot(\theta/2).
\]
由 $\dd(\cot u)/\dd u=-\csc^2u$ 及 $\dot\theta=-\omega$，有
\[
\begin{aligned}
\dot s&=-\frac r2\csc^2\frac\theta2\,\dot\theta
=\frac{r\omega}{2}\csc^2\frac\theta2,\\
\ddot s&=\frac{r\omega}{2}\left(-2\csc^2\frac\theta2\cot\frac\theta2\right)
\frac{\dot\theta}{2}
=\frac{r\omega^2}{2}\csc^2\frac\theta2\cot\frac\theta2.
\end{aligned}
\]
在 $\theta=\pi/3$ 时，
\[
 \boxed{\vct v_C=(2,0)r\omega,\qquad
 \vct a_C=(2\sqrt3,0)r\omega^2.}
\]
以逆时针为正，底部纯滚动给出圆柱$\omega_C=-\dot s/r=-2\omega$、$\alpha_C=-\ddot s/r=-2\sqrt3\omega^2$。当前$\vct r_{M/C}=(-\sqrt3/2,1/2)r$。对固定材料点$M$，
\[
 \vct v_M=\vct v_C+\omega_C\vct k\times\vct r_{M/C},\qquad
 \vct a_M=\vct a_C+\alpha_C\vct k\times\vct r_{M/C}
                -\omega_C^2\vct r_{M/C}.
\]
加速度的平动、切向和法向三项分别为$(2\sqrt3,0)$、$(\sqrt3,3)$和$(2\sqrt3,-2)$，单位均为$r\omega^2$。相加得到
\[
 \boxed{\vct v_M=(3,\sqrt3)r\omega,\qquad
 \vct a_M=(5\sqrt3,1)r\omega^2.}
\]
杆上$N$绕固定$O$匀速转动，当前$\vct r_N=(\sqrt3/2,3/2)r$，故
\[
 \boxed{\vct v_N=(3/2,-\sqrt3/2)r\omega,\qquad
 \vct a_N=(-\sqrt3/2,-3/2)r\omega^2.}
\]
对几何位置$P$，记沿杆的单位矢量为$\vct e_\theta=(\cos\theta,\sin\theta)$，其逆时针法向为$\vct e_\theta^\perp=(-\sin\theta,\cos\theta)$。位置$P=s\vct e_\theta$中的$s$随时间改变，因此
\[
 \vct v_P=\dot s\vct e_\theta-\omega s\vct e_\theta^\perp,
 \quad\vct a_P=(\ddot s-\omega^2s)\vct e_\theta-2\omega\dot s\vct e_\theta^\perp.
\]
代入当前构形，
\[
 \boxed{\vct v_P=(5/2,\sqrt3/2)r\omega,\qquad
 \vct a_P=(5\sqrt3/2,-1/2)r\omega^2.}
\]
图\ref{mp2:fig:rev:wedgecyl-s}左图将三个接触速度从同一点作出，右图列出加速度分量。三种速度不同，但在上杆公法线$\vct n=(-\sqrt3/2,1/2)$上的投影均为$-\sqrt3r\omega$，符合保持接触的条件。求材料点的加速度时，应始终跟踪同一个点；几何接触位置不断更换材料点，其速度的时间导数属于另一种运动量。

\begin{figure}[H]\centering
\begin{tikzpicture}[>=Stealth,line cap=round,line join=round]
\fill(0,0)circle(1.5pt)node[below left]{$P$};\draw[axis](0,0)--(2.1,1.212)node[right]{$\vct v_M$};\draw[axis](0,0)--(1.05,-.606)node[right]{$\vct v_N$};\draw[axis](0,0)--(1.75,.606)node[right]{$\vct v_P$};\draw[guide](-.3,-.52)--(1,1.732)node[above]{$OB$};\node[right,align=left]at(3,.9){$\vct a_C/(r\omega^2)=(2\sqrt3,0)$\\[5pt]$\vct a_M/(r\omega^2)=(5\sqrt3,1)$\\[5pt]$\vct a_N/(r\omega^2)=(-\sqrt3/2,-3/2)$\\[5pt]$\vct a_P/(r\omega^2)=(5\sqrt3/2,-1/2)$};
\end{tikzpicture}
\caption{同一接触处的三个速度，以及分别按定义求出的四种加速度}\label{mp2:fig:rev:wedgecyl-s}
\end{figure}

\end{solution}

\par\medskip\noindent\begin{minipage}{\textwidth}
\begin{example}[两转杆的垂足构形]\label{mp2:rev:intersection}
两无限长直杆分别绕固定点$O,A$转动，几何交点为$P$。当前$OA=l$，$OP$水平向右，$OA$在水平线下方与它成$60^\circ$，且$AP\perp OA$，见图\ref{mp2:fig:rev:intersection-q}。杆$AP$和$OP$的绝对角速度、角加速度分别为$\omega_1,\beta_1$和$\omega_2,\beta_2$，均以逆时针为正。求$\vct v_P,\vct a_P$；$P$不是杆上的固定材料点。
\end{example}

\begin{figure}[H]\centering
\begin{tikzpicture}[>=Stealth,line cap=round,line join=round]
\coordinate(O)at(0,0);\coordinate(A)at(.8,-1.386);\coordinate(P)at(3.2,0);\draw[thick](-.4,0)--(4.1,0);\draw[thick](.3,-1.675)--(4, .462);\draw[guide](O)--(A)node[midway,left]{$l$};\fill(O)circle(1.5pt)node[above left]{$O$};\fill(A)circle(1.5pt)node[below]{$A$};\fill(P)circle(1.5pt)node[above left]{$P$};\draw(.55,0)arc(0:-60:.55);\node at(.75,-.25){$60^\circ$};\draw(.64,-1.109)--(.917,-.949)--(1.077,-1.226);\node[above]at(3.8,.45){$\omega_1,\beta_1$};\node[below]at(3.9,-.1){$\omega_2,\beta_2$};
\end{tikzpicture}
\caption{固定垂线$OA$与两转杆构成的瞬时直角三角形}\label{mp2:fig:rev:intersection-q}
\end{figure}
\end{minipage}\par\medskip

\begin{solution}
当前坐标为$O=(0,0)$、$A=(l/2,-\sqrt3l/2)$、$P=(2l,0)$。杆$AP$的单位方向为$\vct e_1=(\sqrt3/2,1/2)$，杆$OP$的方向为$\vct e_2=(1,0)$。对两条直线约束分别求导，法向速度关系为
\[
 -\frac12v_x+\frac{\sqrt3}{2}v_y=\sqrt3l\omega_1,
 \qquad v_y=2l\omega_2.
\]
故
\[
 \boxed{v_x=2\sqrt3l(\omega_2-\omega_1),\qquad v_y=2l\omega_2.}
\]
对绕固定点 $A_0$ 转动的直线，若单位方向为 $\vct e$、逆时针法线为 $\vct n$、交点到铰点的有向距离为 $s$，则
\[
\vct r_P-\vct r_{A_0}=s\vct e,\qquad
\dot{\vct e}=\omega\vct n,\quad \dot{\vct n}=-\omega\vct e.
\]
对法向约束 $\vct n\cdot(\vct r_P-\vct r_{A_0})=0$ 求导，得到 $\vct n\cdot\vct v=\omega s$。同时，$\dot s=\vct e\cdot\vct v$。再对前一式求导，
\[
-\omega\vct e\cdot\vct v+\vct n\cdot\vct a
=\beta s+\omega\dot s
=\beta s+\omega\vct e\cdot\vct v.
\]
移项后得到
\[
 \vct n\cdot\vct a=2\omega\vct e\cdot\vct v+\beta s.
\]
这里的$2\omega$项包含交点相对直杆的滑动。分别用于两杆，
\[
 a_y=2\omega_2v_x+2l\beta_2,
 \quad-\frac12a_x+\frac{\sqrt3}{2}a_y
 =2\omega_1\left(\frac{\sqrt3}{2}v_x+\frac12v_y\right)+\sqrt3l\beta_1.
\]
代入速度，得到
\[
 \boxed{\begin{aligned}
 a_x&=l[12\omega_1^2-28\omega_1\omega_2+12\omega_2^2
       +2\sqrt3(\beta_2-\beta_1)],\\
 a_y&=l[4\sqrt3\omega_2(\omega_2-\omega_1)+2\beta_2].
 \end{aligned}}
\]
图\ref{mp2:fig:rev:intersection-s}取两杆均以相同恒定角速度$\omega$转动，此时交点速度为$(0,2l\omega)$，加速度为$(-4l\omega^2,0)$。其轨迹是圆心$(l,0)$、半径$l$的圆，绕该圆心的角速度为$2\omega$，独立验证了分量结果。若某杆实际顺时针转动，应给该杆的有符号角速度代入负值。

\begin{figure}[H]\centering
\begin{tikzpicture}[>=Stealth,line cap=round,line join=round]
\draw[curve](1.5,0)circle(1.5);\fill(0,0)circle(1.5pt)node[left]{$O$};\fill(1.5,0)circle(1.5pt)node[below]{$K=(l,0)$};\fill(3,0)circle(1.5pt)node[below right]{$P$};\draw[axis](3,0)--(3,1.5)node[above]{$2l\omega$};\draw[axis,thick](3,0)--(1.7,0)node[midway,above]{$4l\omega^2$};\node[right,align=left]at(3.8,-.4){$\omega_1=\omega_2=\omega$\\$\beta_1=\beta_2=0$};
\end{tikzpicture}
\caption{同角速度极限下，交点以两倍角速度沿圆运动}\label{mp2:fig:rev:intersection-s}
\end{figure}

\end{solution}


\setcounter{chapter}{2}

\section{习题与解答}

\par\medskip\noindent\begin{minipage}{\textwidth}
\begin{exercise}[对称双斜绳与零速度瞬间]\label{mp2:pr:l2p1}

图\ref{mp2:fig:l2p1-q}中，$F=(-a,h)$、$O=(a,h)$ 固定，动滑轮 $P=(x,0)$ 在水平导轨上运动。一根不可伸长轻绳由 $F$ 出发，绕过 $P$、$O$ 后竖直悬挂 $B$。忽略滑轮半径，绳始终绷紧。取悬挂段向下长度为 $y$。已知 $P$ 以恒速 $u$ 向右运动，求 $\dot y,\ddot y$，并给出 $x=0$ 时的结果。

\end{exercise}

\begin{figure}[H]\centering
\begin{tikzpicture}[scale=1.05]
\coordinate(F)at(-2,2.7);\coordinate(O)at(2,2.7);\coordinate(P)at(.5,0);
\draw[thick](-2.4,2.85)--(2.4,2.85);\draw[guide](-2,0)--(2,0);
\draw[thick](F)--(P)--(O)--(2,.7);
\draw[fill=white](P)circle(.13);\draw[fill=white](O)circle(.13);
\fill(F)circle(1.5pt)node[above left]{$F$};\node[above right]at(O){$O$};\node[below]at(P){$P$};
\draw[fill=white](1.78,.45)rectangle(2.22,.95);\node[right]at(2.25,.7){$B$};
\draw[guide](-2,2.7)--(-2,0) (2,2.7)--(2,0);
\draw[axis](-2.3,-.6)--(2.5,-.6)node[right]{$x$};\node[below]at(0,-.6){$0$};
\node[below]at(-2,-.6){$-a$};\node[below]at(2,-.6){$a$};
\node[left]at(-2,1.35){$h$};\draw[axis](2.75,2.7)--(2.75,.7)node[midway,right]{$y$};
\end{tikzpicture}
\caption{习题~\ref{mp2:pr:l2p1}题设：两段绳从固定点 $F$ 和固定滑轮 $O$ 通向动滑轮 $P$，绳的另一端竖直悬挂 $B$。$P$ 沿水平导轨运动。}\label{mp2:fig:l2p1-q}
\end{figure}
\end{minipage}\par\medskip

\begin{solution}
记图\ref{mp2:fig:l2p1-s}中两斜绳长为
\[
s_-=\sqrt{(x+a)^2+h^2},\qquad s_+=\sqrt{(x-a)^2+h^2}.
\]
绳长条件为 $s_-+s_++y=L$。因 $\dot x=u$、$\ddot x=0$，
对每一根斜绳分别求导，有
\[
\dot s_-=\frac{(x+a)u}{s_-},\qquad
\ddot s_-=\frac{u^2}{s_-}-\frac{(x+a)^2u^2}{s_-^3}
=\frac{h^2u^2}{s_-^3}.
\]
$s_+$ 的结果将 $a$ 换为 $-a$ 即得。两项相加后，
\[
\dot y=-u\left(\frac{x+a}{s_-}+\frac{x-a}{s_+}\right),
\]
\[
\ddot y=-h^2u^2\left(\frac1{s_-^3}+\frac1{s_+^3}\right).
\]

\begin{figure}[H]\centering
\begin{tikzpicture}[scale=1.05]
\draw[thick](-2,2.4)--(0,0)--(2,2.4);\draw[guide](-2,2.4)--(.5,0)--(2,2.4);
\fill(0,0)circle(1.4pt)node[below]{$P_0$};\fill(.5,0)circle(1.2pt)node[below]{$P$};
\node[left]at(-1.45,1.5){$s_-$};\node at(.85,1.6){$s_+$};
\draw[axis](0,-.6)--(.9,-.6)node[right]{$u$};
\node[anchor=west]at(3,1.8){$\dot y=0\quad(x=0)$};
\node[anchor=west]at(3,1){$\ddot y=-2h^2u^2/(a^2+h^2)^{3/2}$};
\end{tikzpicture}
\caption{习题~\ref{mp2:pr:l2p1}解图：实线为中线位置 $P_0$ 的绳路，虚线为右移至 $P$ 后的绳路。右侧给出 $x=0$ 时的结果。}\label{mp2:fig:l2p1-s}
\end{figure}

在对称位置 $x=0$，两个一阶项大小相等、符号相反，所以
\[
\dot y=0,\qquad \ddot y=-\frac{2h^2u^2}{(a^2+h^2)^{3/2}}.
\]
这时斜绳总长取得极小值，悬挂段最长，$B$ 的位置处在向下运动转为向上运动的转折点。图中的两条虚线表示 $P$ 略向右移动后的绳段；一段变短、一段变长，而总长的二阶增量为正。

在$x=0$附近，令$s_0=\sqrt{a^2+h^2}$，绳段总长的展开为
\[
 s_-+s_+=2s_0+\frac{h^2}{s_0^3}x^2+O(x^4),
 \qquad y-y_0=-\frac{h^2}{s_0^3}x^2+O(x^4).
\]
以通过对称位置的时刻为 $t=0$，代入 $x=ut$，即可读出转折点附近的局部时间规律。图\ref{mp2:fig:l2p1-s}中，一次项相消，二次项仍保留，所以速度为零的瞬间可以具有非零加速度。
\end{solution}

\par\medskip\noindent\begin{minipage}{\textwidth}
\begin{exercise}[两级动滑轮]\label{mp2:pr:l2p2}

图\ref{mp2:fig:l2p2-q}中的动滑轮 $P,Q$ 由竖直轻杆刚性连接，两滑轮中心均沿竖直方向运动，连接杆保持竖直。上绳一端固定，绕过 $P$ 和若干固定小滑轮后悬挂 $C$；下绳绕过 $Q$，两端连接 $A,B$。各绳不可伸长且始终绷紧，包角不变。所有竖直坐标均向下为正。某时刻 $v_A=u,v_B=2u,a_A=0,a_B=a$，求 $v_P,v_Q,v_C,a_P,a_Q,a_C$。

\end{exercise}

\begin{figure}[H]\centering
\begin{tikzpicture}[scale=.95]
\draw[thick](-.4,3.5)--(4.3,3.5);
\draw[thick](0,3.5)--(0,2.3)arc(180:360:.35)--(.7,3.02)arc(180:90:.08)--(3,3.1)arc(90:0:.3)--(3.3,.3);
\draw[fill=white](.35,2.3)circle(.35);\node[left]at(0,2.3){$P$};
\draw[fill=white](3,2.8)circle(.3);
\draw[fill=white](.78,3.02)circle(.08);\draw[thick](.78,3.5)--(.78,3.1) (3,3.5)--(3,3.1);
\draw[thick](.35,1.95)--(.35,1.2);\draw[fill=white](.35,.9)circle(.3);\node[right]at(.7,.9){$Q$};
\draw[thick](.05,.9)--(.05,.1) (.65,.9)--(.65,-.5);
\draw[fill=white](-.12,-.15)rectangle(.22,.1);\node[left]at(-.15,-.05){$A$};
\draw[fill=white](.48,-.75)rectangle(.82,-.5);\node[right]at(.9,-.6){$B$};
\draw[fill=white](3.1,.05)rectangle(3.5,.3);\node[right]at(3.55,.2){$C$};
\node[anchor=west]at(4.2,2.4){$y_P=y_Q+\mathrm{const}$};
\draw[axis](4.5,1.6)--(4.5,.3)node[right]{$y$};
\draw[line width=.9pt](0,2.3)arc(180:360:.35);
\draw[line width=.9pt](.05,.9)arc(180:0:.3);
\end{tikzpicture}
\caption{习题~\ref{mp2:pr:l2p2}题设：刚性连接的两动滑轮 $P,Q$ 分别由两条绳约束，各竖直坐标向下为正。}\label{mp2:fig:l2p2-q}
\end{figure}
\end{minipage}\par\medskip

\begin{solution}
各固定滑轮之间的绳长并入常数。上绳与下绳分别给出
\[
2y_P+y_C=L_1,\qquad y_A+y_B-2y_Q=L_2.
\]
刚杆连接给出 $y_P-y_Q=\mathrm{const}$。消去两滑轮坐标，还可把悬挂端的位置关系写为 $y_C=\text{常量}-y_A-y_B$。对这些约束求导，得到图\ref{mp2:fig:l2p2-s}中的速度关系：
\[
v_P=v_Q=\frac{v_A+v_B}{2}=\frac{3u}{2},\qquad v_C=-3u.
\]

\begin{figure}[H]\centering
\begin{tikzpicture}[scale=1.05]
\draw[thick](0,0)circle(.35);\draw[thick](0,-1.3)circle(.35);\draw[thick](0,-.35)--(0,-.95);
\node[right]at(.4,0){$P$};\node[right]at(.4,-1.3){$Q$};
\draw[axis](-.9,0)--(-.9,-.75)node[left]{$\dot y_P=\dot y_Q$};
\node[anchor=west]at(2.3,.2){$2\dot y_P+\dot y_C=0$};
\node[anchor=west]at(2.3,-.65){$\dot y_A+\dot y_B-2\dot y_Q=0$};
\node[anchor=west]at(2.3,-1.5){$\dot y_A+\dot y_B+\dot y_C=0$};
\end{tikzpicture}
\caption{习题~\ref{mp2:pr:l2p2}解图：两条绳分别列长约束；刚性连接使两滑轮中心具有相同的速度和加速度。}\label{mp2:fig:l2p2-s}
\end{figure}

绳长式对各坐标都是线性的，本题再求导没有新的几何二次项。因此
\[
a_P=a_Q=\frac{a_A+a_B}{2}=\frac a2,\qquad a_C=-a.
\]
当 $u,a>0$ 时，$P,Q$ 向下运动，$C$ 向上运动。这里的“二倍”来自同一动滑轮两侧绳段同时变化；若绳段倾斜且方向改变，就须像上一题那样保留长度函数的二阶导数。

刚杆使两个滑轮中心的位移相同。$A,B$ 两物块的位移可以不同，下滑轮的平移由它们的平均值决定。
\end{solution}

\par\medskip\noindent\begin{minipage}{\textwidth}
\begin{exercise}[接近共线构形的两连杆]\label{mp2:pr:l2p3}

两根长度均为 $l$ 的杆在 $A$ 处铰接，$O$ 固定，$B$ 沿水平导轨运动，取 $A$ 在上方分支，见图\ref{mp2:fig:l2p3-q}。记 $OB=d=2l\cos\phi$，$0<\phi<\pi/2$。若 $\dot d=-U$、$U>0$ 为常量，求两杆的角速度、角加速度及 $A$ 的速度和加速度。讨论 $\phi\to0^+$ 的极限。

\end{exercise}

\begin{figure}[H]\centering
\begin{tikzpicture}[scale=1.05]
\pgfmathsetmacro{\drawangle}{18}\pgfmathsetmacro{\ll}{2.6}
\coordinate(O)at(0,0);\coordinate(A)at({\ll*cos(\drawangle)},{\ll*sin(\drawangle)});\coordinate(B)at({2*\ll*cos(\drawangle)},0);
\draw[thick](O)--(A)--(B);\draw[guide](-.2,0)--(5.4,0);
\fill(O)circle(1.3pt)node[below]{$O$};\fill(A)circle(1.3pt)node[above]{$A$};\fill(B)circle(1.3pt)node[below]{$B$};
\draw(.7,0)arc(0:\drawangle:.7);\node at(.95,.2){$\phi$};
\draw[axis,thick](B)--++(-1,0)node[midway,below=9pt]{$U$};\node[below]at(2.5,-.2){$d=OB=2l\cos\phi$};
\end{tikzpicture}
\caption{习题~\ref{mp2:pr:l2p3}题设：两等长杆在近共线位置由端点 $B$ 的水平运动驱动。$A$ 取上方分支。}\label{mp2:fig:l2p3-q}
\end{figure}
\end{minipage}\par\medskip

\begin{solution}
对 $d=2l\cos\phi$ 求导，$-U=-2l\sin\phi\dot\phi$。再次求导时，$U$ 为常量，故
\[
\dot\phi=\frac{U}{2l\sin\phi},\qquad
\ddot\phi=-\frac{U\cos\phi}{2l\sin^2\phi}\dot\phi
=-\frac{U^2\cos\phi}{4l^2\sin^3\phi}.
\]
杆 $OA$ 的绝对转角为 $\phi$，杆 $AB$ 从 $A$ 指向 $B$ 的绝对转角为 $-\phi$，故两杆的角速度、角加速度分别是上述量及其相反数。

\begin{figure}[H]\centering
\begin{tikzpicture}[scale=1.05]
\draw[axis](0,0)--(4.8,0)node[right]{$d$};\draw[axis](0,0)--(0,2.4)node[above]{$y_A$};
\draw[curve,domain=0:90,samples=90]plot({4*sin(\x)},{2*cos(\x)});
\draw[guide](4,0)--(4,2.1);\node[below]at(4,0){$2l$};\node[left]at(0,2){$l$};
\node[anchor=west]at(5,.9){$\dot\phi=U/(2l\sin\phi)$};
\end{tikzpicture}
\caption{习题~\ref{mp2:pr:l2p3}解图：上方分支 $y_A=\sqrt{l^2-d^2/4}$ 在 $d=2l$ 处有竖直切线，非零水平驱动速度对应无界的角速度。}\label{mp2:fig:l2p3-s}
\end{figure}

由 $A=(d/2,l\sin\phi)$，横向速度为 $-U/2$，横向加速度为零。竖直分量为
\[
v_{Ay}=l\cos\phi\dot\phi=\frac U2\cot\phi,\qquad
a_{Ay}=-\frac U2\csc^2\phi\dot\phi
=-\frac{U^2}{4l\sin^3\phi}.
\]
合写为
\[
\vct v_A=-\frac U2\vct e_x+\frac U2\cot\phi\,\vct e_y,
\qquad
\vct a_A=-\frac{U^2}{4l\sin^3\phi}\vct e_y.
\]
图\ref{mp2:fig:l2p3-s}显示位置分支在 $d=2l$ 处的竖直切线。保持非零 $U$ 而令 $\phi\to0$，角速度和加速度均无界，说明这种驱动不能从完全伸直状态以有限角速度直接起步。若要平滑通过该构形，水平端点速度必须同时趋零；其具体趋零阶数还要与加速度条件相容。

用端点距的缺额$\delta=2l-d$表示接近伸直的程度，有$\delta\simeq l\phi^2$。于是
\[
 \dot\phi\simeq\frac{U}{2\sqrt{l\delta}},\qquad
 v_{Ay}\simeq\frac U2\sqrt{\frac l\delta},\qquad
 a_{Ay}\simeq-\frac{U^2\sqrt l}{4\delta^{3/2}}.
\]
图\ref{mp2:fig:l2p3-s}的竖直切线使同样的水平位移在铰点处产生越来越大的竖直位移。若改由有限角速度驱动，端点速度应满足$\dot d=-2l\sin\phi\dot\phi\to0$；这是改变驱动规律后的另一种运动，不能与题设恒定非零$U$同时使用。
\end{solution}

\par\medskip\noindent\begin{minipage}{\textwidth}
\begin{exercise}[刚杆上任意定点的轨迹]\label{mp2:pr:l2p4}

沿直角导轨运动的长杆如图\ref{mp2:fig:l2p4-q}，$A=(x,0)$、$B=(0,y)$，$x^2+y^2=l^2$。杆上定点 $P$ 满足 $AP=\lambda l$，其中 $0<\lambda<1$。求 $P$ 的轨迹。若 $A$ 以恒速 $u$ 向右运动，求 $P$ 的速度、加速度；当 $\lambda=1/2$ 时，$P$ 是否作匀速圆周运动？

\end{exercise}

\begin{figure}[H]\centering
\begin{tikzpicture}[scale=1.05]
\pgfmathsetmacro{\ang}{38}
\pgfmathsetmacro{\xx}{3.5*cos(\ang)}\pgfmathsetmacro{\yy}{3.5*sin(\ang)}
\draw[thick](-.2,0)--(3.9,0);\draw[thick](0,-.15)--(0,3.8);
\draw[line width=1pt](0,\yy)--(\xx,0);
\fill(0,\yy)circle(1.5pt)node[left]{$B$};\fill(\xx,0)circle(1.5pt)node[below right]{$A$};
\node[below left]at(0,0){$O$};\node[above right]at({\xx/2},{\yy/2}){$l$};
\draw[|-|](0,-.5)--(\xx,-.5)node[midway,fill=white]{$x$};
\node[left]at(0,{\yy/2}){$y$};

\coordinate(P)at({(1-.35)*\xx},{.35*\yy});\fill(P)circle(1.5pt)node[above right]{$P$};
\node[right]at(3.8,2.5){$0<\lambda<1$};
\end{tikzpicture}
\caption{习题~\ref{mp2:pr:l2p4}题设：刚杆两端沿直角导轨运动；$P$ 固定在杆上，$AP=\lambda l$。}\label{mp2:fig:l2p4-q}
\end{figure}
\end{minipage}\par\medskip

\begin{solution}
图\ref{mp2:fig:l2p4-q}中，$P$是杆上的固定分点，$AP=\lambda l$，故
\[
 x_P=(1-\lambda)x,\qquad y_P=\lambda y.
\]
将$x=x_P/(1-\lambda)$、$y=y_P/\lambda$代入$x^2+y^2=l^2$，便得图\ref{mp2:fig:l2p4-s}中的椭圆弧：
\[
 \boxed{\frac{x_P^2}{(1-\lambda)^2l^2}
 +\frac{y_P^2}{\lambda^2l^2}=1},\qquad x_P,y_P\ge0.
\]
\begin{figure}[H]\centering
\begin{tikzpicture}[scale=1.05]
\draw[axis](0,0)--(4.4,0)node[right]{$x_P$};\draw[axis](0,0)--(0,2.7)node[above]{$y_P$};
\draw[curve,domain=0:90,samples=90]plot({3.25*cos(\x)},{1.75*sin(\x)});
\node[below]at(3.25,0){$(1-\lambda)l$};\node[left]at(0,1.75){$\lambda l$};
\fill({3.25*cos(38)},{1.75*sin(38)})circle(1.4pt)node[above right]{$P$};
\end{tikzpicture}
\caption{习题~\ref{mp2:pr:l2p4}解图：$P=(1-\lambda)A+\lambda B$。其轨迹为半轴分别为 $(1-\lambda)l$ 与 $\lambda l$ 的椭圆弧。}\label{mp2:fig:l2p4-s}
\end{figure}
端点$A$以恒速$u$向右运动，因而$\dot x=u$、$\ddot x=0$。杆长约束连续求导两次：
\[
 x\dot x+y\dot y=0,\qquad
 \dot x^2+x\ddot x+\dot y^2+y\ddot y=0.
\]
在$y>0$的运动阶段，第一式给$\dot y=-xu/y$；代入第二式，
\[
 \ddot y=-\frac{u^2+\dot y^2}{y}
 =-\frac{u^2(x^2+y^2)}{y^3}
 =-\frac{l^2u^2}{y^3}.
\]
再对$P$的坐标求导，得到
\[
 \boxed{\vct v_P=(1-\lambda)u\vct e_x
 -\lambda\frac{xu}{y}\vct e_y},\qquad
 \boxed{\vct a_P=-\lambda\frac{l^2u^2}{y^3}\vct e_y}.
\]
当$\lambda=1/2$时，椭圆成为半径$l/2$的四分之一圆，而
\[
 v_P=\frac u2\sqrt{1+\frac{x^2}{y^2}}=\frac{ul}{2y}.
\]
速率随$y$减小而增大，因此不是匀速圆周运动。$y\to0^+$时速度和加速度发散，表明恒定水平端速的理想约束在杆接近水平时出现奇性；上述公式不包含$y=0$端点。
\end{solution}

\par\medskip\noindent\begin{minipage}{\textwidth}
\begin{exercise}[两条传送带之间的圆盘]\label{mp2:pr:l2p5}

半径 $R$ 的圆盘夹在上下两条水平传送带之间，圆心保持固定高度，两处均无滑动。见图\ref{mp2:fig:l2p5-q}。上、下带速度分别为 $3u,u$，加速度分别为 $3a,a$，方向均向右，$u>0$。求圆心速度、加速度及圆盘的角速度、角加速度；求当时轮顶材料点 $T$、轮底材料点 $B$ 的加速度。

\end{exercise}

\begin{figure}[H]\centering
\begin{tikzpicture}[scale=1.05]
\draw[thick](-1.8,0)--(3,0) (-1.8,2.4)--(3,2.4);
\draw(0,1.2)circle(1.2);\fill(0,1.2)circle(1.2pt)node[left]{$C$};
\fill(0,2.4)circle(1.2pt)node[above left]{$T$};\fill(0,0)circle(1.2pt)node[below left]{$B$};
\draw[axis,thick](0,2.85)--(2.1,2.85)node[right]{$3u$};
\draw[axis,thick](0,-.45)--(.7,-.45)node[right]{$u$};
\draw[guide](0,1.2)--(1.2,1.2)node[midway,above]{$R$};
\end{tikzpicture}
\caption{习题~\ref{mp2:pr:l2p5}题设：圆盘夹在两条平行传送带之间，两个接触处都无滑动。两带速度分别为 $3u$ 与 $u$。}\label{mp2:fig:l2p5-q}
\end{figure}
\end{minipage}\par\medskip

\begin{solution}
以逆时针为正，轮顶、轮底的切向速度条件分别是
\[
V-\omega R=3u,\qquad V+\omega R=u.
\]
相加、相减得 $V=2u$、$\omega=-u/R$。再沿运动过程求导，
\[
\dot V=2a,\qquad \alpha=-a/R.
\]
图\ref{mp2:fig:l2p5-s}中的速度分布符合这两条关联。

\begin{figure}[H]\centering
\begin{tikzpicture}[scale=.9]
\draw(0,1.2)circle(1.2);\draw[thick](-1.6,0)--(3,0) (-1.6,2.4)--(3,2.4);
\draw[axis,thick](0,2.4)--(2.1,2.4)node[above]{$3u$};
\draw[axis,thick](0,1.2)--(1.4,1.2)node[right]{$2u$};
\draw[axis,thick](0,0)--(.7,0)node[below]{$u$};
\node[above left]at(0,2.4){$T$};\node[left]at(0,1.2){$C$};\node[below left]at(0,0){$B$};
\node at(1,3){速度};\node at(1,-.65){$\omega=-u/R$};
\draw(5.5,1.2)circle(1.2);
\fill(5.5,2.4)circle(1.3pt)node[above left]{$T$};
\fill(5.5,1.2)circle(1.3pt)node[left]{$C$};
\fill(5.5,0)circle(1.3pt)node[below left]{$B$};
\draw[axis,thick](5.5,2.4)--(5.5,1.7);
\draw[axis,thick](5.5,0)--(5.5,.7);
\node at(6.4,3){加速度的竖直分量};
\node[anchor=west]at(7.1,2.1){$a_{Ty}=-u^2/R$};
\node[anchor=west]at(7.1,1.2){$a_{Cy}=0$};
\node[anchor=west]at(7.1,.3){$a_{By}=+u^2/R$};
\end{tikzpicture}
\caption{习题~\ref{mp2:pr:l2p5}解图：左图三条速度箭头的长度比为 $3:2:1$。右图只画加速度的竖直分量，取 $y$ 轴向上；轮顶向下，轮底向上。}\label{mp2:fig:l2p5-s}
\end{figure}

由 $\vct r_{T/C}=R\vct e_y$、$\vct r_{B/C}=-R\vct e_y$，
\[
\vct a_T=3a\vct e_x-\frac{u^2}{R}\vct e_y,
\qquad
\vct a_B=a\vct e_x+\frac{u^2}{R}\vct e_y.
\]
它们的水平分量分别等于相接触传送带的加速度，竖直分量却非零，并且都指向圆心。无滑动约束联系的是每时刻接触的一对材料点的速度，不能将两侧材料点的完整加速度直接相等。

若不预先代入数值，两带的速度为$U_T,U_B$时，
\[
 V=\frac{U_T+U_B}{2},\qquad
 \omega=\frac{U_B-U_T}{2R}.
\]
固定高度条件使圆心没有竖直加速度。轮缘材料点则有法向项$-\omega^2\vct r$，这正是图\ref{mp2:fig:l2p5-s}右图上、下两支箭头的来源。两带匀速时，轮缘材料点仍随圆盘转动，经过接触位置后继续运动，因而保留向心加速度。
\end{solution}

\par\medskip\noindent\begin{minipage}{\textwidth}
\begin{exercise}[外接圆盘的接触点]\label{mp2:pr:l2p6}

半径分别为 $r_1,r_2$ 的两圆盘中心固定，外缘保持无滑动接触，见图\ref{mp2:fig:l2p6-q}。两盘角速度均以逆时针为正。已知左盘 $\omega_1>0$、角加速度 $\alpha_1$，求右盘的角速度和角加速度。若两盘均匀速转动，比较当前接触处两材料点的速度、加速度。

\end{exercise}

\begin{figure}[H]\centering
\begin{tikzpicture}[scale=1.05]
\coordinate(C1)at(0,0);\coordinate(C2)at(2.25,0);\coordinate(K)at(.9,0);
\draw(C1)circle(.9);\draw(C2)circle(1.35);
\fill(C1)circle(1.3pt)node[below]{$C_1$};\fill(C2)circle(1.3pt)node[below]{$C_2$};\fill(K)circle(1.3pt)node[below]{$K$};
\draw[guide](C1)--(K) (K)--(C2);\node[above]at(.45,0){$r_1$};\node[above]at(1.6,0){$r_2$};
\draw[axis](70:.6)arc(70:150:.6);\node[above left]at(-.4,.8){$\omega_1$};
\draw[axis]({2.25+.9*cos(45)},{.9*sin(45)})arc(45:-35:.9);\node[right]at(3.2,.8){$\omega_2$};
\end{tikzpicture}
\caption{习题~\ref{mp2:pr:l2p6}题设：两圆盘的中心固定，外缘在 $K$ 处无滑动接触。图示 $r_2/r_1=3/2$。}\label{mp2:fig:l2p6-q}
\end{figure}
\end{minipage}\par\medskip

\begin{solution}
左盘在接触处的位矢为 $r_1\vct e_x$，右盘为 $-r_2\vct e_x$。切向速度相等给出
\[
r_1\omega_1=-r_2\omega_2,
\qquad \omega_2=-\frac{r_1}{r_2}\omega_1.
\]
半径不变，故 $\alpha_2=-(r_1/r_2)\alpha_1$。两盘转向相反。

\begin{figure}[H]\centering
\begin{tikzpicture}[scale=1.05]
\fill(0,0)circle(1.4pt)node[below]{$K$};
\draw[axis,thick](0,0)--(-1.5,0)node[left]{$r_1\omega_1^2$};
\draw[axis,thick](0,0)--(1,0)node[right]{$r_2\omega_2^2$};
\draw[axis](0,0)--(0,1.3)node[above]{$\vct v_1=\vct v_2$};
\node at(-1.5,-.75){指向 $C_1$};\node at(1.7,-.75){指向 $C_2$};
\end{tikzpicture}
\caption{习题~\ref{mp2:pr:l2p6}解图：经过 $K$ 的两盘材料点速度相同，记为 $\vct v_1=\vct v_2$；水平箭头为各自的向心加速度。几何接触位置 $K$ 固定。}\label{mp2:fig:l2p6-s}
\end{figure}

匀速转动时，图\ref{mp2:fig:l2p6-s}中两材料点速度都为 $r_1\omega_1\vct e_y$，而
\[
\vct a_1=-r_1\omega_1^2\vct e_x,
\qquad \vct a_2=+r_2\omega_2^2\vct e_x.
\]
两者方向相反，大小之比为
\[
\frac{|\vct a_1|}{|\vct a_2|}=\frac{r_2}{r_1}.
\]
两圆心固定，几何接触位置也固定，故这一位置的速度和加速度均为零；经过它的两轮缘材料点则具有上述速度和加速度。

角加速度不为零时，两材料点的加速度分别为
\[
 \vct a_1=-r_1\omega_1^2\vct e_x+r_1\alpha_1\vct e_y,
 \qquad
 \vct a_2=r_2\omega_2^2\vct e_x-r_2\alpha_2\vct e_y.
\]
由$r_1\alpha_1=-r_2\alpha_2$可见，切向加速度相同，法向加速度仍反向。图\ref{mp2:fig:l2p6-s}只画匀速时的法向项；一般情况下应在两支法向箭头上叠加同一个切向分量。
\end{solution}

\par\medskip\noindent\begin{minipage}{\textwidth}
\begin{exercise}[摆线尖点处的运动]\label{mp2:pr:l2p7}

半径 $R$ 的圆盘在固定地面上以圆心恒速 $V>0$ 向右纯滚，材料点 $P$ 在 $t=0$ 时接触地面，见图\ref{mp2:fig:l2p7-q}。取该接触位置为坐标原点，令 $\theta=Vt/R$。求 $P$ 的轨迹、速度和加速度，并求尖点附近 $x,y$ 之间的最低阶关系。

\end{exercise}

\begin{figure}[H]\centering
\begin{tikzpicture}[scale=1.05]
\draw[axis](-1.35,0)--(7,0)node[right]{$x$};
\draw[axis](0,0)--(0,2.65)node[above]{$y$};
\draw[guide,black!55](0,1)circle(1);
\draw[line width=.55pt](pi,1)circle(1);
\draw[line width=.9pt,domain=0:360,samples=240]plot({pi*\x/180-sin(\x)},{1-cos(\x)});
\fill(0,1)circle(1.3pt)node[left]{$C(0)$};
\fill(pi,1)circle(1.3pt)node[below=2pt]{$C(t_1)$};
\fill(0,0)circle(1.3pt)node[below left]{$P(0)$};
\fill(pi,2)circle(1.3pt)node[above]{$P(t_1)$};
\fill({2*pi},0)circle(1.3pt)node[above left]{$P(2t_1)$};
\draw(pi,1)--(pi,2)node[midway,right]{$R$};
\draw[axis,thick](pi,1)--({pi+1.3},1)node[right]{$V$};
\draw[axis]({pi-.65},1)arc(180:115:.65);
\node at(pi,2.85){$t_1=\pi R/V$};
\draw[guide](0,0)--(0,-.75) (pi,0)--(pi,-.75);
\draw[|-|](0,-.6)--(pi,-.6)node[midway,fill=white]{$\pi R$};
\node[below]at({2*pi},0){$2\pi R$};
\end{tikzpicture}
\caption{习题~\ref{mp2:pr:l2p7}题设：虚线圆盘表示初始位置，实线圆盘表示 $t_1=\pi R/V$ 时的位置。半周内圆心向右移动 $\pi R$，标记点从轮底转到轮顶。粗曲线是 $P$ 在 $0\le t\le2t_1$ 内走过的一拱摆线。}\label{mp2:fig:l2p7-q}
\end{figure}
\end{minipage}\par\medskip

\begin{solution}
图\ref{mp2:fig:l2p7-q}中，$C(0)$ 与 $C(t_1)$ 相距 $\pi R$，而 $P(0)$、$P(t_1)$ 分别位于当时的轮底和轮顶。任意时刻的圆心坐标为 $(R\theta,R)$；将初始位矢 $(0,-R)$ 顺时针转过 $\theta$，再加上圆心位置，得
\[
x=R(\theta-\sin\theta),\qquad y=R(1-\cos\theta).
\]
连续求导得
\[
\vct v_P=V\bigl[(1-\cos\theta)\vct e_x+\sin\theta\vct e_y\bigr],
\]
\[
\vct a_P=\frac{V^2}{R}
\bigl[\sin\theta\vct e_x+\cos\theta\vct e_y\bigr].
\]

\begin{figure}[H]\centering
\begin{tikzpicture}[scale=1.05]
\draw[axis](-2,0)--(2.5,0)node[right]{$x$};\draw[axis](0,0)--(0,2.2)node[above]{$y$};
\draw[curve,domain=-125:125,samples=180]plot({1.5*(pi*\x/180-sin(\x))},{1.5*(1-cos(\x))});
\fill(0,0)circle(1.4pt)node[below right]{$P(0)$};
\draw[axis,thick](0,0)--(0,1.3);
\node[anchor=west]at(3,1.65){$\vct a_P(0)=(V^2/R)\vct e_y$};
\node[anchor=west]at(3,.9){$x\sim R\theta^3/6$};
\node[anchor=west]at(3,.15){$y\sim R\theta^2/2$};
\end{tikzpicture}
\caption{习题~\ref{mp2:pr:l2p7}解图：摆线在接触位置形成尖点。$P$ 的速度为零，而加速度竖直向上。}\label{mp2:fig:l2p7-s}
\end{figure}

所以 $\vct v_P(0)=0$、$\vct a_P(0)=(V^2/R)\vct e_y$，与图\ref{mp2:fig:l2p7-s}一致。尖点附近的展开为
\[
x=\frac R6\theta^3+O(\theta^5),\qquad
y=\frac R2\theta^2+O(\theta^4).
\]
消去 $\theta$，最低阶有
\[
x^2\sim\frac{2}{9R}y^3.
\]
在原点，两支曲线相接形成尖点，不能用普通圆弧近似。材料点只在接触瞬间速度为零，随后便向上离开原来的接触位置。

\end{solution}

\par\medskip\noindent\begin{minipage}{\textwidth}
\begin{exercise}[三尖点内摆线（选做）]\label{mp2:pr:l2p8}

重复例\ref{mp2:ex:l2e7}的内滚实验，取 $R=3r$，圆心公转角速度为常量 $\Omega>0$，点 $P$ 起初在最右接触处，见图\ref{mp2:fig:l2p8-q}。求 $P$ 的轨迹参数式和速率，确定一周中的尖点位置，并求 $\phi=\pi/3$ 时的速度、加速度。

\end{exercise}

\begin{figure}[H]\centering
\begin{tikzpicture}[scale=.9]
\pgfmathsetmacro{\RR}{3}\pgfmathsetmacro{\rr}{1}\pgfmathsetmacro{\ph}{30}
\pgfmathsetmacro{\aa}{\RR-\rr}\pgfmathsetmacro{\qq}{\aa/\rr}
\coordinate(O)at(0,0);\coordinate(C)at({\aa*cos(\ph)},{\aa*sin(\ph)});
\coordinate(K)at({\RR*cos(\ph)},{\RR*sin(\ph)});
\coordinate(P)at({\aa*cos(\ph)+\rr*cos(-\qq*\ph)},{\aa*sin(\ph)+\rr*sin(-\qq*\ph)});
\draw(O)circle(\RR);\draw[thick](C)circle(\rr);
\draw[guide](O)--(K);\draw(C)--(P);
\fill(O)circle(1.2pt)node[below left]{$O$};\fill(C)circle(1.3pt)node[above left]{$C$};
\fill(K)circle(1.3pt)node[above right]{$K$};\fill(P)circle(1.3pt)node[below right]{$P$};
\draw[guide](0,0)--(\RR,0);\draw(.65,0)arc(0:\ph:.65);\node at(.85,.28){$\phi$};
\node[above left]at({.5*\aa*cos(\ph)},{.5*\aa*sin(\ph)}){$R-r$};
\node[right]at({(\aa+.5*\rr)*cos(\ph)},{(\aa+.5*\rr)*sin(\ph)}){$r$};
\draw(O)--(-120:\RR)node[midway,left]{$R$};

\draw[guide,domain=0:360,samples=240]plot({2*cos(\x)+cos(2*\x)},{2*sin(\x)-sin(2*\x)});
\end{tikzpicture}
\caption{习题~\ref{mp2:pr:l2p8}题设：大圆半径为 $3r$，小盘半径为 $r$，圆心绕 $O$ 匀速公转并在内壁纯滚。虚线是标记点的完整轨迹。}\label{mp2:fig:l2p8-q}
\end{figure}
\end{minipage}\par\medskip

\begin{solution}
此时 $R-r=2r$、$q=2$，例题的参数式化为
\[
x=r(2\cos\phi+\cos2\phi),\qquad
y=r(2\sin\phi-\sin2\phi).
\]
求导并用 $\dot\phi=\Omega$，
\[
\vct v_P=2r\Omega\bigl[(-\sin\phi-\sin2\phi)\vct e_x
+(\cos\phi-\cos2\phi)\vct e_y\bigr].
\]
对速度两分量平方相加，
\[
\begin{aligned}
v_P^2&=4r^2\Omega^2\bigl[(\sin\phi+\sin2\phi)^2
+(\cos\phi-\cos2\phi)^2\bigr]\\
&=4r^2\Omega^2(2-2\cos3\phi)
=16r^2\Omega^2\sin^2\frac{3\phi}{2}.
\end{aligned}
\]
取非负平方根，速率为 $v_P=4r\Omega|\sin(3\phi/2)|$。

\begin{figure}[H]\centering
\begin{tikzpicture}[scale=.9]
\draw[curve,domain=0:360,samples=300]plot({2*cos(\x)+cos(2*\x)},{2*sin(\x)-sin(2*\x)});
\foreach\a in{0,120,240}{\fill({3*cos(\a)},{3*sin(\a)})circle(1.4pt);}
\coordinate(P)at(.5,{sqrt(3)/2});\fill(P)circle(1.4pt)node[below right]{$P$};
\draw[axis,thick](P)--++({-sqrt(3)},1)node[above left]{$\vct v_P$};
\draw[axis,thick](P)--++(.8,{.8*sqrt(3)})node[above right]{$\vct a_P$};
\draw[axis](-2,0)--(3.7,0)node[right]{$x/r$};\draw[axis](0,-2.9)--(0,3)node[above]{$y/r$};
\end{tikzpicture}
\caption{习题~\ref{mp2:pr:l2p8}解图：三尖点内摆线。所示点对应 $\phi=\pi/3$，其速度与加速度相互垂直。}\label{mp2:fig:l2p8-s}
\end{figure}

图\ref{mp2:fig:l2p8-s}中的三个尖点对应 $\phi=0,2\pi/3,4\pi/3$，位置为
\[
(3r,0),\qquad
\left(-\frac{3r}{2},\frac{3\sqrt3r}{2}\right),\qquad
\left(-\frac{3r}{2},-\frac{3\sqrt3r}{2}\right).
\]
加速度的一般式是
\[
\vct a_P=-2r\Omega^2\bigl[(\cos\phi+2\cos2\phi)\vct e_x
+(\sin\phi-2\sin2\phi)\vct e_y\bigr].
\]
在 $\phi=\pi/3$，
\[
\vct v_P=(-2\sqrt3r\Omega,2r\Omega),\qquad
\vct a_P=(r\Omega^2,\sqrt3r\Omega^2).
\]
两者点积为零，故这一时刻相互垂直，与 $R=2r$ 时的直线运动不同。整个标记点运动的周期为 $2\pi/\Omega$；只看不带标记的圆盘形状，不能据圆的旋转对称性缩短标记点的周期。

\end{solution}

\par\medskip\noindent\begin{minipage}{\textwidth}
\begin{exercise}[匀速外移的径向滑块]\label{mp2:pr:l2p9}

水平直槽绕固定点 $O$ 以恒定角速度 $\Omega>0$ 逆时针转动，质量为 $m$ 的滑块 $P$ 相对槽以恒速 $u>0$ 向外移动，$s=s_0+ut$，见图\ref{mp2:fig:l2p9-q}。径向驱动器维持这一相对运动，槽壁无摩擦，重力由水平支承抵消。求 $P$ 的速度、加速度以及槽壁横向作用力；径向驱动器还须施加多大的径向力？

\end{exercise}

\begin{figure}[H]\centering
\begin{tikzpicture}[scale=1.05]
\pgfmathsetmacro{\th}{30}
\draw[thick](-.2,-.11547)--(3.3,1.90526);
\coordinate(P)at({2.2*cos(\th)},{2.2*sin(\th)});
\fill(0,0)circle(1.4pt)node[below left]{$O$};\fill(P)circle(2.2pt)node[below right]{$P,m$};
\draw[axis](P)--++({.9*cos(\th)},{.9*sin(\th)})node[above]{$u$};
\draw[axis](15:.7)arc(15:75:.7);\node[above]at(.45,.75){$\Omega$};
\node[below]at(.95,.45){$s$};
\end{tikzpicture}
\caption{习题~\ref{mp2:pr:l2p9}题设：水平面内转动的直槽中，滑块以给定相对速率 $u$ 外移。相对运动由径向驱动器维持。}\label{mp2:fig:l2p9-q}
\end{figure}
\end{minipage}\par\medskip

\begin{solution}
相对槽的加速度为零，绝对加速度却并不为零。代入 $\dot s=u,\ddot s=0,\alpha=0$，
\[
\vct v_P=u\vct e_r+s\Omega\vct e_\theta,
\]
\[
\vct a_P=-s\Omega^2\vct e_r+2u\Omega\vct e_\theta.
\]
图\ref{mp2:fig:l2p9-s}分别标出这两个加速度方向。

\begin{figure}[H]\centering
\begin{tikzpicture}[scale=1.05]
\pgfmathsetmacro{\th}{30}
\coordinate(P)at(0,0);
\draw[axis,guide](0,0)--({2.1*cos(\th)},{2.1*sin(\th)})node[right]{$\vct e_r$};
\draw[axis,guide](0,0)--({2.1*cos(\th+90)},{2.1*sin(\th+90)})node[above]{$\vct e_\theta$};
\draw[axis,thick](0,0)--({-1.6*cos(\th)},{-1.6*sin(\th)})node[left]{$-s\Omega^2\vct e_r$};
\draw[axis,thick](0,0)--({.8*cos(\th+90)},{.8*sin(\th+90)})node[left]{$2u\Omega\vct e_\theta$};
\fill(P)circle(1.4pt)node[below right]{$P$};
\node[anchor=west]at(3.1,.9){$N=2mu\Omega$};
\node[anchor=west]at(3.1,.1){$F_r=-ms\Omega^2$};
\end{tikzpicture}
\caption{习题~\ref{mp2:pr:l2p9}解图：相对加速度为零时，绝对加速度仍有向心项和横向项。横向项由槽壁作用力提供。}\label{mp2:fig:l2p9-s}
\end{figure}

槽壁只能施加垂直槽方向的力，故其作用力为
\[
\vct N=2mu\Omega\vct e_\theta.
\]
径向合力由驱动器承担，$F_r=-ms\Omega^2$，方向向内。驱动器维持 $\ddot s=0$，而地面系中的径向加速度含 $-s\Omega^2$。若去掉径向驱动器，不能继续假设相对速率为常量，而应由径向牛顿方程重新求 $s(t)$。

在地面系中，滑块动能为 $K=m[u^2+s^2\Omega^2]/2$。因 $u,\Omega$ 均为常量，$\dot K=(m/2)(2s\dot s\Omega^2)=msu\Omega^2$。两种力的功率分别为
\[
 \mathcal P_N=(2mu\Omega)(s\Omega)=2msu\Omega^2,
 \qquad\mathcal P_r=(-ms\Omega^2)u=-msu\Omega^2.
\]
两项功率相加等于 $\dot K$。图\ref{mp2:fig:l2p9-s}中的横向槽壁力向滑块输入能量，径向驱动器同时吸收能量。槽壁光滑只排除了沿槽摩擦，它在地面系中运动，仍可通过横向作用力做功。
\end{solution}

\par\medskip\noindent\begin{minipage}{\textwidth}
\begin{exercise}[等角速度直槽交点的完整轨迹]\label{mp2:pr:l2p10}

两无限长直槽分别绕 $O=(0,0)$、$A=(0,-h)$ 转动，恒定角速度均为 $\omega>0$，夹角 $\beta$ 恒定且 $0<\beta<\pi$。初始 $\theta_1=0,\theta_2=\beta$，见图\ref{mp2:fig:l2p10-q}。将直槽视为可延伸的直线，允许交点经过铰点两侧。求交点 $P$ 的轨迹、运动周期、速率和加速度大小。

\end{exercise}

\begin{figure}[H]\centering
\begin{tikzpicture}[scale=1.05]
\pgfmathsetmacro{\hh}{2}\pgfmathsetmacro{\bb}{60}\pgfmathsetmacro{\ss}{\hh*cos(\bb)/sin(\bb)}
\draw[thick](-.5,0)--(3.1,0);
\draw[thick](-.2,{-\hh-.2*tan(\bb)})--(2,{2*tan(\bb)-\hh});
\fill(0,0)circle(1.3pt)node[above left]{$O$};\fill(0,-2)circle(1.3pt)node[below left]{$A$};\fill(\ss,0)circle(1.4pt)node[below right]{$P$};
\draw[guide](0,-2)--(0,0);\node[left]at(0,-1){$h$};
\draw[guide](0,-2)--(1,-2);\draw(.5,-2)arc(0:60:.5);\node at(.6,-1.55){$\beta$};
\node[anchor=west]at(3.7,.5){$\dot\theta_1=\dot\theta_2=\omega$};
\node[anchor=west]at(3.7,-.35){$\theta_2-\theta_1=\beta$};
\end{tikzpicture}
\caption{习题~\ref{mp2:pr:l2p10}题设：两直槽以相同恒定角速度转动，夹角 $\beta$ 保持不变。两槽都作为无限直线，交点可以经过铰点的两侧。}\label{mp2:fig:l2p10-q}
\end{figure}
\end{minipage}\par\medskip

\begin{solution}
令 $\theta_1=\omega t$、$\theta_2=\omega t+\beta$。例\ref{mp2:ex:l2e10}的位置关系给出有符号坐标
\[
s=\frac{h\cos(\omega t+\beta)}{\sin\beta}.
\]
因 $P=s(\cos\omega t,\sin\omega t)$，用积化和差整理：
\begin{align*}
x_P&=\frac h2\cot\beta+\frac{h}{2\sin\beta}\cos(2\omega t+\beta),\\
y_P&=-\frac h2+\frac{h}{2\sin\beta}\sin(2\omega t+\beta).
\end{align*}

\begin{figure}[H]\centering
\begin{tikzpicture}[scale=1.05]
\pgfmathsetmacro{\hh}{2}\pgfmathsetmacro{\bb}{60}
\pgfmathsetmacro{\mx}{\hh*cos(\bb)/(2*sin(\bb))}\pgfmathsetmacro{\drawradius}{\hh/(2*sin(\bb))}
\coordinate(M)at(\mx,-1);\coordinate(P)at({2*\mx},0);
\draw(M)circle(\drawradius);\draw[guide](0,0)--(0,-2) (M)--(P);
\fill(M)circle(1.3pt)node[below right]{$M$};\fill(0,0)circle(1.3pt)node[above left]{$O$};\fill(0,-2)circle(1.3pt)node[below left]{$A$};\fill(P)circle(1.3pt)node[above right]{$P$};
\draw[axis](\mx+\drawradius,-1)arc(0:55:\drawradius);\node[right]at(1.85,-.2){$2\omega$};
\node[anchor=west]at(3.2,.3){$M=(\tfrac h2\cot\beta,-\tfrac h2)$};
\node[anchor=west]at(3.2,-.65){$\rho=h/(2\sin\beta)$};
\end{tikzpicture}
\caption{习题~\ref{mp2:pr:l2p10}解图：图示 $\beta=\pi/3$，交点轨迹是经过 $O,A$ 的圆，圆心为 $M$。交点绕 $M$ 的角速度是槽角速度的两倍。}\label{mp2:fig:l2p10-s}
\end{figure}

图\ref{mp2:fig:l2p10-s}中的圆心、半径分别是
\[
M=\left(\frac h2\cot\beta,-\frac h2\right),\qquad
\rho=\frac{h}{2\sin\beta}.
\]
$P$ 绕 $M$ 以角速度 $2\omega$ 作匀速圆周运动，因此
\[
T=\frac{2\pi}{2\omega}=\frac\pi\omega,\qquad
v=2\omega\rho=\frac{h\omega}{\sin\beta},\qquad
a=(2\omega)^2\rho=\frac{2h\omega^2}{\sin\beta}.
\]
加速度始终指向 $M$。每转半周，两条不分正反方向的直槽已恢复原来的两条直线，因此交点周期是槽方位角周期的一半。若实际装置使用有限杆件或只允许 $s>0$，便只能保留圆上的相应弧段，不能直接采用这里的完整周期。

\end{solution}

\par\medskip\noindent\begin{minipage}{\textwidth}
\begin{exercise}[四连杆的尺寸参数（选做）]\label{mp2:pr:l2p11}

图\ref{mp2:fig:l2p11-q}中固定点 $A=(0,0)$、$D=((1+\kappa)l,0)$，其中 $\kappa>0$。当前 $B=(0,l)$、$C=(l,2l)$，杆 $AB$ 以恒定角速度大小 $\Omega>0$ 顺时针转动。求杆 $BC,DC$ 的逆时针有符号角速度、角加速度和 $C$ 的速度、加速度，并取 $\kappa=2$ 与例\ref{mp2:ex:l2e8}比较。

\end{exercise}

\begin{figure}[H]\centering
\begin{tikzpicture}[scale=1.05]
\pgfmathsetmacro{\ll}{1.2}\pgfmathsetmacro{\kk}{3}
\coordinate(A)at(0,0);\coordinate(B)at(0,\ll);\coordinate(C)at(\ll,{2*\ll});\coordinate(D)at({(1+\kk)*\ll},0);
\draw[thick](-.4,0)--({(1+\kk)*\ll+.4},0);
\draw[line width=1pt](A)--(B)--(C)--(D);
\foreach\p/\pos in{A/below,B/left,C/above,D/below}{\fill(\p)circle(1.4pt)node[\pos]{$\p$};}
\draw[guide](C)--(\ll,0) (0,{2*\ll})--(C);
\draw[|-|](-.55,0)--(-.55,\ll)node[midway,left]{$l$};
\draw[|-|](-.55,\ll)--(-.55,{2*\ll})node[midway,left]{$l$};
\draw[|-|](0,-.6)--(\ll,-.6)node[midway,fill=white]{$l$};
\draw[|-|](\ll,-.6)--({(1+\kk)*\ll},-.6)node[midway,fill=white]{$\kappa l$};
\draw[axis](70:.7)arc(70:20:.7);\node at(.95,.55){$\Omega$};
\end{tikzpicture}
\caption{习题~\ref{mp2:pr:l2p11}题设：改变右侧固定铰的位置，记其横坐标为 $(1+\kappa)l$。}\label{mp2:fig:l2p11-q}
\end{figure}
\end{minipage}\par\medskip

\begin{solution}
记两从动杆的角速度、角加速度为 $\omega_2,\alpha_2$ 和 $\omega_3,\alpha_3$。当前相对位矢为 $l(1,1)$、$l(-\kappa,2)$。图\ref{mp2:fig:l2p11-s}中的两条速度链给出
\[
\Omega-\omega_2=-2\omega_3,\qquad \omega_2=-\kappa\omega_3.
\]
故
\[
\omega_2=\frac{\kappa}{\kappa+2}\Omega,\qquad
\omega_3=-\frac1{\kappa+2}\Omega,
\]
\[
\vct v_C=\frac{\Omega l}{\kappa+2}(2\vct e_x+\kappa\vct e_y).
\]

\begin{figure}[H]\centering
\begin{tikzpicture}[scale=.9]
\draw[axis](-.2,0)--(5.4,0)node[right]{$x$};\draw[axis](0,-.2)--(0,2.8)node[above]{$y$};
\coordinate(C)at(1.2,2.4);\coordinate(D)at(4.8,0);\coordinate(B)at(0,1.2);
\draw[thick](B)--(C)--(D);\fill(C)circle(1.4pt)node[above]{$C$};\fill(B)circle(1.3pt)node[left]{$B$};\fill(D)circle(1.3pt)node[below]{$D$};
\draw[axis,thick](C)--++(.5,.75)node[right]{$\vct v_C$};
\node[anchor=west]at(6.2,2){$v_{Cx}/l=\Omega-\omega_2=-2\omega_3$};
\node[anchor=west]at(6.2,1.1){$v_{Cy}/l=\omega_2=-\kappa\omega_3$};
\node[anchor=west]at(6.2,.2){$\kappa+2>0$};
\end{tikzpicture}
\caption{习题~\ref{mp2:pr:l2p11}解图：图示 $\kappa=3$。$C$ 的速度分别按杆 $BC$ 和杆 $DC$ 表示，两式的同一分量必须相等。}\label{mp2:fig:l2p11-s}
\end{figure}

两条加速度链逐分量相等：
\begin{align*}
-\alpha_2-\omega_2^2&=-2\alpha_3+\kappa\omega_3^2,\\
-\Omega^2+\alpha_2-\omega_2^2&=-\kappa\alpha_3-2\omega_3^2.
\end{align*}
由第一式 $\alpha_2=2\alpha_3-\omega_2^2-\kappa\omega_3^2$，代入第二式，得
\[
(\kappa+2)\alpha_3=\Omega^2+2\omega_2^2+(\kappa-2)\omega_3^2.
\]
因此
\[
\alpha_3=\frac{(\kappa+1)(3\kappa+2)}{(\kappa+2)^3}\Omega^2,
\quad
\alpha_2=-\frac{(\kappa+1)(\kappa^2-4\kappa-4)}{(\kappa+2)^3}\Omega^2.
\]
再代入由 $B$ 出发的加速度式，
\[
\begin{split}
\vct a_C=-\frac{\Omega^2l}{(\kappa+2)^3}\bigl[&(5\kappa^2+8\kappa+4)\vct e_x\\
&+(3\kappa^3+5\kappa^2+4\kappa+4)\vct e_y\bigr].
\end{split}
\]
取 $\kappa=2$，两角加速度均为 $3\Omega^2/8$，$\vct a_C=-(\Omega^2l/8)(5,7)$，回到例题。对所有 $\kappa>0$，当前速度方程的行列式与 $\kappa+2$ 成正比，不出现两从动杆共线造成的速度不唯一；这不代表机构沿整个运动历程都没有特殊构形。

\end{solution}

\par\medskip\noindent\begin{minipage}{\textwidth}
\begin{exercise}[由质心运动反求局部轨道（选做）]\label{mp2:pr:l2p12}

边长为 $l$ 的正三角形薄板在一段切线连续的固定曲线轨道上无滑动滚动，质心 $C$ 始终以非零速度沿高度 $H$ 的水平线运动。当前只有一条边与轨道相切，接触点在边内部，其他部分没有干涉，见图\ref{mp2:fig:l2p12-q}。记三角形内切圆半径为 $d=l/(2\sqrt3)$。求这一边接触阶段的局部轨道形状，取拱顶为 $x=0$；再由接触点落在边内给出必要的斜率范围，并说明为何不能任意延长这一局部解。

\end{exercise}

\begin{figure}[H]\centering
\begin{tikzpicture}[scale=1.05]
\pgfmathsetmacro{\ddr}{.8}\pgfmathsetmacro{\xx}{.4}\pgfmathsetmacro{\HH}{1.6}
\pgfmathsetmacro{\th}{atan(-sinh(\xx/\ddr))}
\pgfmathsetmacro{\nx}{-sin(\th)}\pgfmathsetmacro{\ny}{cos(\th)}
\pgfmathsetmacro{\halfedge}{sqrt(3)*\ddr}
\coordinate(C)at(\xx,\HH);\coordinate(P)at(\xx,{\HH-\ddr*cosh(\xx/\ddr)});
\coordinate(M)at({\xx-\ddr*\nx},{\HH-\ddr*\ny});
\coordinate(V1)at({\xx-\ddr*\nx-\halfedge*cos(\th)},{\HH-\ddr*\ny-\halfedge*sin(\th)});
\coordinate(V2)at({\xx-\ddr*\nx+\halfedge*cos(\th)},{\HH-\ddr*\ny+\halfedge*sin(\th)});
\coordinate(V3)at({\xx+2*\ddr*\nx},{\HH+2*\ddr*\ny});
\draw[axis](-2,0)--(2.5,0)node[right]{$x$};\draw[axis](0,-.2)--(0,3.5)node[above]{$y$};
\draw[curve,domain=-1.053566:1.053566,samples=120]plot(\x,{\HH-\ddr*cosh(\x/\ddr)});
\draw[thick](V1)--(V2)--(V3)--cycle;
\draw[guide](-1.8,\HH)--(2.2,\HH);\node[right]at(2.2,\HH){$y_C=H$};
\draw[guide](C)--(P);\draw(C)--(M);
\draw[guide]($(C)!.5!(M)$)--(-.8,1.05)node[left]{$d$};
\fill(C)circle(1.3pt)node[above right]{$C$};\fill(P)circle(1.3pt)node[below]{$P$};\fill(M)circle(1pt)node[left]{$M$};
\draw[axis](C)--++(.8,0)node[right]{$V$};
\end{tikzpicture}
\caption{习题~\ref{mp2:pr:l2p12}题设：正三角形在曲线轨道上无滑动滚动，质心 $C$ 沿水平线运动。图示 $H=2d$，只画当前一边接触的轨道段。}\label{mp2:fig:l2p12-q}
\end{figure}
\end{minipage}\par\medskip

\begin{solution}
记瞬时接触位置为 $P=(x,y)$。无滑动时，经过 $P$ 的板上材料点速度为零，故 $P$ 是板的速度瞬心。由于 $\vct v_C$ 水平，$CP$ 必须竖直，见图\ref{mp2:fig:l2p12-s}左图。这里不能把 $CP$ 当成三角形的内切圆半径；与边垂直的线段是 $CM=d$。

令边的有符号倾角为 $\theta$，$\tan\theta=y'$，$|\theta|<\pi/2$。质心到该边的垂直距离给出
\[
(H-y)\cos\theta=d,
\qquad H-y=d\sqrt{1+y'^2}.
\]
若只保留这个一阶方程，形式上还会出现不滚动的平直分支。为选出给定的非零速度滚动，继续使用瞬心速度关系。设 $\dot x=V\ne0$，则
\[
V=-\omega(H-y),\qquad \omega=\dot\theta,
\]
从而
\[
\frac{\dd\theta}{\dd x}=-\frac1{H-y}=-\frac{\cos\theta}{d}.
\]
分离变量，在拱顶取 $\theta(0)=0$，积分得
\[
\int_0^\theta\sec u\dd u
=\left[\ln(\sec u+\tan u)\right]_0^\theta
=-\frac1d\int_0^x\dd\xi=-\frac xd.
\]
令 $E=\ee^{-x/d}>0$，则 $\sec\theta+\tan\theta=E$。结合 $\sec^2\theta-\tan^2\theta=1$，有 $\sec\theta-\tan\theta=E^{-1}$，相加、相减得到
\[
\sec\theta=\frac{E+E^{-1}}2=\cosh\frac xd,\qquad
\tan\theta=\frac{E-E^{-1}}2=-\sinh\frac xd.
\]
所以局部轨道为
\[
y=H-d\cosh\frac xd.
\]
这是向下弯曲的悬链线形拱段。

\begin{figure}[H]\centering
\begin{tikzpicture}[scale=.95]
\coordinate(P)at(0,0);\coordinate(C)at(0,1.154701);\coordinate(M)at(-.5,.288675);
\draw[thick](-2,1.154701)--(1,-.57735);\draw(C)--(M)--(P)--(C);
\fill(C)circle(1.3pt)node[above]{$C$};\fill(M)circle(1.3pt)node[left]{$M$};\fill(P)circle(1.3pt)node[below]{$P$};
\node[left]at(-.25,.72){$d$};\node[right]at(0,.57){$H-y$};\node[below left]at(-.25,.14){$s$};
\draw[guide](P)--(1,0);\draw(.55,0)arc(0:-30:.55);\node at(.75,-.18){$\theta$};
\begin{scope}[shift={(6,0)},x=1.2cm,y=1cm]
\draw[axis](-1.8,0)--(2,0)node[right]{$x/d$};\draw[axis](0,-.65)--(0,1.55)node[above]{$y/d$};
\draw[guide,domain=-1.6:1.6,samples=100]plot(\x,{2-cosh(\x)});
\draw[line width=1pt,domain=-1.316958:1.316958,samples=120]plot(\x,{2-cosh(\x)});
\node[below]at(-1.316958,0){$-b/d$};\node[below]at(1.316958,0){$b/d$};
\node[above left]at(0,1){$1$};
\end{scope}
\end{tikzpicture}
\caption{习题~\ref{mp2:pr:l2p12}解图：左图中 $CM=d$、$CP=H-y$，$CM\perp MP$。右图示意 $H=2d$ 的局部轨道，粗线只保留接触点在当前边内的必要区间。}\label{mp2:fig:l2p12-s}
\end{figure}

再检查有限边长。以边中点 $M$ 为原点，沿该边的有符号距离记为 $s=MP$。图中直角三角形给出
\[
s=-d\tan\theta.
\]
接触点严格在边内部要求 $|s|<l/2$，因此
\[
|y'|<\sqrt3,\qquad
|x|<b=d\,\operatorname{arsinh}\sqrt3=d\ln(2+\sqrt3).
\]
右图的粗线标出这一必要区间。到达端点时，接触点已移到顶点，原来“同一条边内部相切”的运动阶段结束。另一边接触、顶点转换和其他部分的干涉都可能进一步限制可用轨道；这里得到的是一段局部解，不能据此任意延长曲线或直接确定整条周期路面。

还须讨论整条路面能否由这一局部解拼接而成。形式上，取$b=d\ln(2+\sqrt3)$，可写出连续的周期曲线
\[
 y(x)=H-d\cosh\frac{x-2jb}{d},\qquad
 2jb-b\le x\le2jb+b,\quad j\in\mathbb Z.
\]
相邻段在$y=H-2d$处相接，左、右斜率分别为$-\sqrt3,+\sqrt3$，拱顶高度为$H-d$。谷底两条切线的夹角为$60^\circ$，与三角形顶角相同。能否通过谷底，还须检验顶点离开接触位置后的非穿透条件。

设谷底为 $P_0=(x_0,H-2d)$，三角形质心向右运动，瞬时速率为 $V>0$，水平加速度为 $A_h=\dot V$。先假定换边后的单侧加速度有限；匀速情形对应 $A_h=0$。若仍沿下一拱段纯滚动，令 $\xi=(x-2jb)/d$，则 $\dot\xi=V/d$。由瞬心关系，
\[
\begin{aligned}
\omega&=-\frac{V}{H-y}=-\frac{V}{d\cosh\xi},\\
\alpha&=-\frac{A_h}{d\cosh\xi}
+\frac Vd\frac{\sinh\xi}{\cosh^2\xi}\dot\xi
=-\frac{A_h}{d\cosh\xi}
+\frac{V^2}{d^2}\frac{\sinh\xi}{\cosh^2\xi}.
\end{aligned}
\]
下一拱段的左端满足 $\xi=-b/d$、$\cosh\xi=2$、$\sinh\xi=-\sqrt3$，故刚离开谷底时
\[
\omega=-\frac V{2d},\qquad
\alpha=-\frac{A_h}{2d}-\frac{\sqrt3V^2}{4d^2}.
\]
把换边瞬间位于 $P_0$ 的材料顶点记为 $K$，它相对质心的位置为 $(0,-2d)$。逐项计算刚体加速度，
\[
\begin{aligned}
\vct a_K&=A_h\vct e_x+\alpha\vct k\times(0,-2d)-\omega^2(0,-2d)\\
&=(A_h+2d\alpha,2d\omega^2).
\end{aligned}
\]
水平加速度中的 $A_h$ 项抵消，得到
\[
 \vct a_K=\left(-\frac{\sqrt3V^2}{2d},\frac{V^2}{2d}\right),
 \qquad \vct v_K=0.
\]
因而换边后短时间$\tau>0$内，
\[
 x_K=x_0-\frac{\sqrt3V^2}{4d}\tau^2+o(\tau^2),\qquad
 y_K=H-2d+\frac{V^2}{4d}\tau^2+o(\tau^2).
\]
顶点先向谷底左侧运动，而左侧路面的斜率为$-\sqrt3$。该横坐标处的路面高度为
\[
 y_{\rm road}(x_K)=H-2d+\frac{3V^2}{4d}\tau^2+o(\tau^2).
\]
所以
\[
 \boxed{y_K-y_{\rm road}(x_K)
 =-\frac{V^2}{2d}\tau^2+o(\tau^2)<0.}
\]
图\ref{mp2:fig:rev:triangle-road}右图画出这一局部关系：顶点 $K$ 落到路面以下，发生几何穿入。即使左右角速度可以连续衔接，非穿透条件仍不满足。由于 $A_h$ 已消去，允许有限的水平加速度也不能消除这一二阶穿入。因此，上述周期曲线不能在保持非零质心速率和有限单侧加速度的条件下，直接用作平滑换边的完整轨道。题设限定的同一边内部接触区间内，反悬链线解仍然成立；要延伸到换边过程，必须另行处理接触和质心运动条件。

\begin{figure}[H]\centering
\begin{tikzpicture}[>=Stealth]
\begin{scope}[scale=.7]
\def\bb{1.316958}
\foreach \j in {-1,0,1}{\draw[curve,domain=-1.316958:1.316958,samples=60]plot({\x+2*\j*\bb},{-cosh(\x)});}
\draw[guide](-4.1,0)--(4.1,0)node[right]{$y_C=H$};
\coordinate(P)at(1.316958,-2);\coordinate(C)at(1.316958,0);
\draw[thick]({1.316958-1.732051},1)--(P)--({1.316958+1.732051},1)--cycle;
\fill(P)circle(1.4pt)node[below]{$P_0$};\fill(C)circle(1.4pt)node[above]{$C$};\draw[guide](P)--(C)node[midway,right]{$2d$};
\node at(0,-3.1){形式拼接，谷底瞬时相合};
\end{scope}
\begin{scope}[shift={(7,-1.4)},scale=.85]
\draw[axis](-2,0)--(1,0)node[right]{$\Delta x$};\draw[axis](0,-.3)--(0,2.7)node[above]{$\Delta y$};
\draw[curve](-1.5,2.598)--(0,0)--(.9,1.559);\draw[guide](-1.2,0)--(-1.2,2.078);
\fill(-1.2,2.078)circle(1.4pt)node[above right]{路面};
\fill(-1.2,.693)circle(1.4pt)node[below]{$K$};\draw[axis](0,0)--(-1.2,.693);
\draw[<->](-1.5,.693)--(-1.5,2.078)node[midway,left]{$\Delta h$};
\node at(-.2,3.6){$\Delta h=V^2\tau^2/(2d)$};
\node[below right]at(0,0){$P_0$};\node at(-.5,-.8){谷底附近的二阶位移};
\end{scope}
\end{tikzpicture}
\caption{局部轨道的形式拼接及其失效原因；右图按同一比例画出二阶位移，顶点$K$低于左侧路面}\label{mp2:fig:rev:triangle-road}
\end{figure}

\end{solution}






\chapter{曲线与三维运动}

质点沿曲线运动时，速度的大小与方向都可能改变。速率的变化由切向加速度描述，方向的变化则与轨迹曲率相联系。刚体上的不同材料点、同一质点在不同参考系中的轨迹，都须据各自的位置关系计算；本讲依次讨论这些情形。


\section{曲率与自然坐标}
加速度同时描述速率变化和方向变化。切向、法向把运动时序与轨迹几何分开，曲率半径由实际空间中的弧长和转角确定。

\subsection{弧长、切线与曲率圆}
设光滑轨迹的位置为 $\vct r(q)$，在所讨论的区间内 $\vct r_q\ne0$。沿规定的正向量取弧长 $s$，有
\[
\dd s=|\vct r_q|\dd q,\qquad
\vct e_t=\frac{\dd\vct r}{\dd s}.
\]
这里先取 $q$ 随 $s$ 增大。若把 $\vct e_t$ 定为实际运动方向，则 $\dot s=v=|\vct v|>0$。在停止、反向或尖点处，须先用位置的参数方程求导，再讨论两侧极限。

单位切向量满足 $\vct e_t\cdot\vct e_t=1$，求导可得 $\vct e_t\cdot\dd\vct e_t/\dd s=0$。在曲率不为零的点，定义
\begin{equation}\label{mp3:eq:l3-curv}
\frac{\dd\vct e_t}{\dd s}=\kappa\vct e_n,\qquad
\kappa=\left|\frac{\dd\vct e_t}{\dd s}\right|,\qquad \rho=\frac1\kappa.
\end{equation}
$\vct e_n$ 是主法线方向，指向轨迹弯曲的一侧。图\ref{mp3:fig:l3-natural}中的曲率圆与轨迹在 $P$ 处相切，圆心 $K=P+\rho\vct e_n$。曲率圆与轨迹在 $P$ 附近具有相同的二阶展开，离开 $P$ 后一般不再重合。

\begin{figure}[H]\centering
\begin{tikzpicture}[scale=.95]

\draw[curve,domain=-1.6:1.6,samples=70]plot(\x,{.32*\x*\x+.015*\x*\x*\x*\x});
\draw[guide](0,1.5625)circle(1.5625);
\fill(0,0)circle(1.3pt)node[below]{$P$};\fill(0,1.5625)circle(1.3pt)node[left]{$K$};
\draw[axis](0,0)--(1.3,0)node[right]{$\vct e_t$};
\draw[axis](0,0)--(0,1.3)node[right]{$\vct e_n$};
\node[left]at(-.05,.8){$\rho$};
\fill(1,.335)circle(1.3pt)node[right]{$Q$};
\node at(1.1,1.7){曲率圆};
\begin{scope}[shift={(5.8,.35)}]
\draw[axis](0,0)--(1.8,0)node[right]{$\vct e_t(s)$};
\draw[axis](0,0)--({1.8*cos(35)},{1.8*sin(35)})node[above]{$\vct e_t(s+\dd s)$};
\draw[axis](1.8,0)--({1.8*cos(35)},{1.8*sin(35)})node[midway,right]{$\dd\vct e_t$};
\draw(.65,0)arc(0:35:.65);\node at(.95,.27){$\dd\psi$};
\node[below]at(.8,-.5){$|\dd\vct e_t|\simeq|\dd\psi|$};
\end{scope}
\end{tikzpicture}
\caption{弧长与切线的转角。左图以局部坐标表示轨迹及曲率圆；右图把两个相邻切向单位向量移至同一起点。}\label{mp3:fig:l3-natural}
\end{figure}


对平面曲线，若切线与固定 $x$ 轴的夹角为 $\psi$，则 $\kappa=|\dd\psi/\dd s|$。有向曲率 $\dd\psi/\dd s$ 可以正、可以负，式\eqref{mp3:eq:l3-curv}中的 $\kappa$ 则非负。主法线始终朝凹侧。平面轨迹经过拐点时，凹向可以改变，而 $\kappa$ 在拐点降为零。

\subsection{加速度的自然分解}
由 $\vct v=v\vct e_t$ 使用乘积与链式法则，
\[
\vct a=\dot v\vct e_t+v\frac{\dd\vct e_t}{\dd s}\dot s
=\dot v\vct e_t+v^2\kappa\vct e_n.
\]
用 $\rho=1/\kappa$ 改写，
\begin{equation}\label{mp3:eq:l3-natural-a}
\vct a=\dot v\,\vct e_t+\frac{v^2}{\rho}\vct e_n
=a_t\vct e_t+a_n\vct e_n.
\end{equation}
其中 $a_t=\dot v$ 是带符号的分量，减速时为负；$a_n=v^2/\rho\ge0$。切向力改变速率，法向合力使速度转向。因此
\[
F_t=m v\frac{\dd v}{\dd s},\qquad F_n=m\frac{v^2}{\rho}.
\]
$F_n$ 表示各真实力在主法线上的分量之和；有接触支持力时，重力等其他力的法向分量也须计入。

若已知 $\vct v,\vct a$ 而未显式求出轨迹，先对 $v^2=\vct v\cdot\vct v$ 求导，得到 $2v\dot v=2\vct v\cdot\vct a$。扣除切向投影，余下的矢量即为法向加速度；同时 $|\vct v\times\vct a|=va_n=v^3\kappa$。由此得到
\begin{align}
a_t&=\frac{\vct v\cdot\vct a}{v},&
\vct a_\perp&=\vct a-\frac{\vct v\cdot\vct a}{v^2}\vct v,\label{mp3:eq:l3-proj}\\
\kappa&=\frac{|\vct v\times\vct a|}{v^3},&
\rho&=\frac{v^3}{|\vct v\times\vct a|}.\label{mp3:eq:l3-vacurv}
\end{align}
当 $v>0$ 且 $\vct a_\perp\ne0$ 时，曲率中心还可直接写为
\[
\vct r_K=\vct r+\frac{v^2}{|\vct a_\perp|^2}\vct a_\perp.
\]
若 $\vct v\times\vct a=0$，该点曲率为零；它可能是直线运动中的一点，也可能是曲线的拐点。

\subsection{不同轨迹方程下的曲率}
设平面轨迹为 $\vct r(q)=(x(q),y(q))$，记 $w=|\vct r_q|>0$。由 $\dd s/\dd q=w$ 及 $\vct e_t=\vct r_q/w$，
\[
\begin{aligned}
w_q&=\frac{\vct r_q\cdot\vct r_{qq}}w,\\
\frac{\dd\vct e_t}{\dd s}
&=\frac1w\frac{\dd}{\dd q}\left(\frac{\vct r_q}w\right)
=\frac{\vct r_{qq}}{w^2}
-\frac{\vct r_q(\vct r_q\cdot\vct r_{qq})}{w^4}.
\end{aligned}
\]
取模，并用 $|\vct r_q\times\vct r_{qq}|^2=w^2|\vct r_{qq}|^2-(\vct r_q\cdot\vct r_{qq})^2$，得到
\begin{equation}\label{mp3:eq:l3-param-curv}
\kappa=\frac{|x_qy_{qq}-y_qx_{qq}|}{(x_q^2+y_q^2)^{3/2}}.
\end{equation}
参数 $q$ 可以是时间、转角或任何正则参数。若轨迹写为 $y=f(x)$，令 $q=x$ 即得
\[
\rho=\frac{[1+(f')^2]^{3/2}}{|f''|}.
\]
轨迹近乎竖直时，$y=f(x)$ 可能失效，此时应交换坐标或保留参数形式。

极坐标中的两个单位向量随 $\varphi$ 转动，满足
$\dd\vct e_r/\dd\varphi=\vct e_\varphi$、
$\dd\vct e_\varphi/\dd\varphi=-\vct e_r$。故对轨迹 $r=r(\varphi)$，
\[
\vct r_\varphi=r'\vct e_r+r\vct e_\varphi,\quad
\vct r_{\varphi\varphi}=(r''-r)\vct e_r+2r'\vct e_\varphi.
\]
两次导数的叉积为
\[
(\vct r_\varphi\times\vct r_{\varphi\varphi})\cdot\vct k
=2(r')^2-r(r''-r)=r^2+2(r')^2-rr''.
\]
代入式\eqref{mp3:eq:l3-param-curv}的向量形式，
\begin{equation}\label{mp3:eq:l3-polar-curv}
\rho=\frac{[r^2+(r')^2]^{3/2}}
{|r^2+2(r')^2-rr''|}.
\end{equation}
式中的撇号表示对 $\varphi$ 求导。分母为零而分子非零时，曲率为零；若分子也为零，则参数不正则，不能用这个商直接判断。

\par\medskip\noindent\begin{minipage}{\textwidth}
\begin{example}[抛体轨迹的曲率与曲率中心]\label{mp3:example:l3-e1}

质点从原点以初速 $v_0$、仰角 $\theta$ 抛出，$0<\theta<\pi/2$，忽略空气阻力，重力加速度为 $g$。运动区域不设地面，轨迹可延伸至发射点下方。图\ref{mp3:fig:l3-e1-q}中 $u=v_0\cos\theta$ 为不变的水平分速度。求任意时刻的曲率半径，确定其最小值，并给出最高点与下降过程中 $v_y=-u$ 时的曲率中心。

\end{example}
\begin{figure}[H]\centering
\begin{tikzpicture}[scale=.95]

\draw[axis](-.3,0)--(5.4,0)node[right]{$x$};\draw[axis](0,-.2)--(0,2.8)node[above]{$y$};
\draw[curve,domain=0:4,samples=80]plot(\x,{\x-.25*\x*\x});
\fill(0,0)circle(1.3pt)node[below left]{$O$};
\fill(2,1)circle(1.3pt)node[above]{$H$};
\fill(3,.75)circle(1.3pt)node[above right]{$P$};
\draw[axis](0,0)--(1,1)node[above right]{$\vct v_0$};
\draw[axis](3,.75)--(4,.25)node[right]{$\vct v$};
\draw[axis](3,.75)--(3,-.3)node[below]{$\vct g$};
\draw(.6,0)arc(0:45:.6);\node at(.85,.3){$\theta$};
\node at(4.7,1.8){$v_x=u$};

\end{tikzpicture}
\caption{抛体在最高点 $H$ 与下降段 $P$ 的速度方向。}\label{mp3:fig:l3-e1-q}
\end{figure}
\end{minipage}\par\medskip
\begin{solution}
抛体的速度、加速度分别为
\[
\vct v=(u,w),\quad w=v_0\sin\theta-gt,\qquad \vct a=(0,-g).
\]
所以 $|\vct v\times\vct a|=gu$，式\eqref{mp3:eq:l3-vacurv}给出
\[
\rho(t)=\frac{(u^2+w^2)^{3/2}}{gu}.
\]
在整条抛物线上，$u$ 不变，$w=0$ 时速率最小，故
\[
\rho_{\min}=\rho_H=\frac{u^2}{g}.
\]
图\ref{mp3:fig:l3-e1-s}左图将 $H$ 平移至局部原点。此时速度水平向右，加速度竖直向下，主法线指向下方，曲率中心在 $H$ 的正下方，距离为 $u^2/g$。

\begin{figure}[H]\centering
\begin{tikzpicture}[scale=.95]

\draw[curve,domain=-1.35:1.35,samples=60]plot(\x,{-.5*\x*\x});
\draw[guide](0,-1)circle(1);
\fill(0,0)circle(1.3pt)node[above left]{$H$};
\fill(0,-1)circle(1.3pt)node[left]{$K_H$};
\draw[axis](0,0)--(1.2,0)node[right]{$\vct v_H$};
\draw[axis](0,0)--(0,-.65)node[right]{$\vct e_n$};
\node at(0,-2.55){$\rho_H=u^2/g$};
\begin{scope}[shift={(5.7,0)}]
\fill(0,0)circle(1.3pt)node[above]{$P$};
\draw[axis](0,0)--(1.1,-1.1)node[right]{$\vct v$};
\draw[axis](0,0)--(-.85,-.85)node[left]{$\vct a_\perp$};
\draw[axis](0,0)--(.85,-.85)node[above right]{$a_t\vct e_t$};
\draw[axis,thick](0,0)--(0,-1.7)node[right]{$\vct g$};
\draw[guide](-.85,-.85)--(0,-1.7)--(.85,-.85);
\draw[guide](0,0)--(-1.65,-1.65);
\fill(-1.65,-1.65)circle(1.3pt)node[left]{$K$};
\node at(.1,-2.55){$w=-u$};
\end{scope}

\end{tikzpicture}
\caption{最高点的曲率圆，以及 $w=-u$ 时重力的自然分解。右图各加速度箭头按同一比例绘制。}\label{mp3:fig:l3-e1-s}
\end{figure}


一般位置的法向加速度矢量为
\[
\vct a_\perp=\frac{gu}{u^2+w^2}(w,-u).
\]
由曲率中心的位置公式，
\[
\vct r_K-\vct r_P
=\left(\frac{(u^2+w^2)w}{gu},-\frac{u^2+w^2}{g}\right).
\]
在 $w=-u$ 时，$\rho=2\sqrt2\,u^2/g$，且
\[
\vct r_K-\vct r_P
=\left(-\frac{2u^2}{g},-\frac{2u^2}{g}\right).
\]
右图的 $PK$ 朝左下方，与下降速度垂直。重力在 $PK$ 方向的分量决定 $v^2/\rho$，沿速度方向的分量则使小球继续加速。

\end{solution}

\par\medskip\noindent\begin{minipage}{\textwidth}
\begin{example}[抛物线形槽中的支持力]\label{mp3:example:l3-e2}

质量为 $m$ 的小珠在固定、光滑的抛物线形轨道内侧运动，轨道方程为 $y=x^2/(2p)$，$p>0$。小珠从最低点 $O$ 以速率 $u$ 向右出发，见图\ref{mp3:fig:l3-e2-q}。求速度、曲率半径与支持力随位置的变化，判断小珠能否在运动中脱离轨道。

\end{example}
\begin{figure}[H]\centering
\begin{tikzpicture}[scale=.95]

\draw[axis](-2.3,0)--(2.7,0)node[right]{$x$};\draw[axis](0,-.2)--(0,2.5)node[above]{$y$};
\draw[curve,domain=-2:2,samples=70]plot(\x,{.45*\x*\x});
\fill(0,0)circle(1.3pt)node[below left]{$O$};
\fill(1.35,.8201)circle(2.3pt)node[right]{$P(x,y)$};
\draw[axis](0,0)--(1,0)node[midway,below]{$u$};
\draw[guide](1.35,0)--(1.35,.8201);
\node at(-1.2,2.25){$y=x^2/(2p)$};

\end{tikzpicture}
\caption{小珠由槽底沿右支上升；轨道从小珠下方提供支持。}\label{mp3:fig:l3-e2-q}
\end{figure}
\end{minipage}\par\medskip
\begin{solution}
从 $O$ 到 $P$，支持力垂直于位移，不做功。机械能守恒给出
\[
v^2=u^2-2gy=u^2-\frac{gx^2}{p}.
\]
令 $q=x/p$。轨道的一、二阶导数为 $y'=x/p=q$、$y''=1/p$，代入函数图像的曲率公式，
\[
\rho=\frac{(1+y'^2)^{3/2}}{|y''|}=p(1+q^2)^{3/2},\qquad
\vct e_n=\frac{(-q,1)}{\sqrt{1+q^2}}.
\]
图\ref{mp3:fig:l3-e2-s}中 $\vct e_n$ 指向槽内凹侧。重力在该方向的分量为 $-mg/\sqrt{1+q^2}$，所以
\[
N-\frac{mg}{\sqrt{1+q^2}}=\frac{mv^2}{p(1+q^2)^{3/2}}.
\]
代入能量关系，含 $x^2$ 的两项相消，得到
\[
N=\frac{m(u^2+gp)}{p(1+q^2)^{3/2}}>0.
\]

\begin{figure}[H]\centering
\begin{tikzpicture}[scale=.95]

\draw[curve,domain=-.3:2,samples=60]plot(\x,{.45*\x*\x});
\coordinate(P)at(1,.45);\fill(P)circle(2pt)node[below right]{$P$};
\draw[axis](P)--({1-1.05*.9/sqrt(1.81)},{.45+1.05/sqrt(1.81)})node[above left]{$\vct e_n,\ N$};
\draw[axis](P)--({1+1.05/sqrt(1.81)},{.45+1.05*.9/sqrt(1.81)})node[right]{$\vct e_t$};
\draw[axis](P)--(1,-.8)node[below]{$mg$};
\node at(4.9,1.5){$N>0$};
\draw[axis](3.5,0)--(6.3,0)node[right]{$x$};
\draw[axis](3.5,0)--(3.5,2.3)node[above]{$N$};
\draw[curve,domain=0:2.3,samples=60]plot({3.5+\x},{2/(1+.5*\x*\x)^1.5});\fill(5.8,{2/(1+.5*2.3*2.3)^1.5})circle(1.3pt);\draw[guide](5.8,0)--(5.8,{2/(1+.5*2.3*2.3)^1.5});\node[below]at(5.8,0){$x_{\max}$};
\node[left]at(3.5,2){$N(0)$};
\node at(5,-.9){可达区间内始终接触};

\end{tikzpicture}
\caption{支持力沿主法线指向槽内。右图为支持力随 $x\ge0$ 的变化示意，曲线只取小珠可达部分。}\label{mp3:fig:l3-e2-s}
\end{figure}


小珠到达 $x_{\max}=u\sqrt{p/g}$ 时停止上升，随后沿原路返回。转向点的曲率半径有限，但法向加速度为零，支持力恰好平衡重力的法向分量。整个可达区间内 $N>0$，接触得以维持。

上述接触按轨道只能推、不能拉处理。若用细环将小珠套在曲线上，约束可以双侧受力，便无须另加 $N\ge0$ 的条件。

\end{solution}

\par\medskip\noindent\begin{minipage}{\textwidth}
\begin{example}[沿对数螺线运动]\label{mp3:example:l3-e3}

平面轨迹为 $r=r_0\ee^{b\varphi}$，其中 $r_0>0,b>0$，质点满足 $\dot\varphi=\Omega>0$。见图\ref{mp3:fig:l3-e3-q}。求速度、加速度的极坐标分量，以及切向、法向加速度和曲率半径；确定主法线相对径向的方向。

\end{example}
\begin{figure}[H]\centering
\begin{tikzpicture}[scale=.95]

\draw[axis](-2.4,0)--(2.6,0)node[right]{$x$};\draw[axis](0,-1.9)--(0,2.4)node[above]{$y$};
\draw[curve,domain=0:300,samples=160]plot({.6*exp(.22*\x*pi/180)*cos(\x)},{.6*exp(.22*\x*pi/180)*sin(\x)});
\coordinate(P)at({.6*exp(.22*60*pi/180)*cos(60)},{.6*exp(.22*60*pi/180)*sin(60)});
\draw[guide](0,0)--(P);\fill(P)circle(1.5pt)node[above right]{$P$};
\node[below left]at(0,0){$O$};
\draw(.4,0)arc(0:60:.4);\node at(.6,.24){$\varphi$};
\node at(4,1.1){$r=r_0\ee^{b\varphi}$};
\node at(4,.3){$\dot\varphi=\Omega$};

\end{tikzpicture}
\caption{对数螺线的一段；随 $\varphi$ 增大，质点向外运动。}\label{mp3:fig:l3-e3-q}
\end{figure}
\end{minipage}\par\medskip
\begin{solution}
对 $r=r_0\ee^{b\varphi}$ 使用链式法则，
\[
\dot r=br_0\ee^{b\varphi}\dot\varphi=b\Omega r,
\qquad \ddot r=b\Omega\dot r=b^2\Omega^2r.
\]
再对位置矢量 $r\vct e_r$ 求导，代入极坐标速度、加速度式，得到
\[
\vct v=\Omega r(b\vct e_r+\vct e_\varphi),\qquad
\vct a=\Omega^2r[(b^2-1)\vct e_r+2b\vct e_\varphi].
\]
速率为 $v=\Omega r\sqrt{1+b^2}$，故
\[
a_t=\dot v=\Omega\sqrt{1+b^2}\dot r
=b\Omega^2r\sqrt{1+b^2}.
\]
图\ref{mp3:fig:l3-e3-s}以同一点 $P$ 为起点画出两组单位向量。与速度同向的单位切向量和向内弯曲的主法线分别为
\[
\vct e_t=\frac{b\vct e_r+\vct e_\varphi}{\sqrt{1+b^2}},
\qquad
\vct e_n=\frac{-\vct e_r+b\vct e_\varphi}{\sqrt{1+b^2}}.
\]

\begin{figure}[H]\centering
\begin{tikzpicture}[scale=.95]

\fill(0,0)circle(1.5pt)node[below]{$P$};
\draw[axis](0,0)--(1.6,0)node[right]{$\vct e_r$};
\draw[axis](0,0)--(0,1.6)node[above]{$\vct e_\varphi$};
\draw[axis,thick](0,0)--({1.6*.5/sqrt(1.25)},{1.6/sqrt(1.25)})node[above right]{$\vct e_t$};
\draw[axis,thick](0,0)--({-1.6/sqrt(1.25)},{1.6*.5/sqrt(1.25)})node[left]{$\vct e_n$};
\coordinate(K)at({-2.4/sqrt(1.25)},{2.4*.5/sqrt(1.25)});
\draw[guide]({-1.6/sqrt(1.25)},{1.6*.5/sqrt(1.25)})--(K);
\fill(K)circle(1.5pt)node[above]{$K$};
\node at(4,1.25){$a_t=b\Omega^2r\sqrt{1+b^2}$};
\node at(4,.3){$a_n=\Omega^2r\sqrt{1+b^2}$};
\node at(4,-.6){$PK=\rho=r\sqrt{1+b^2}$};

\end{tikzpicture}
\caption{局部方向图取 $b=1/2$。四个单位向量等长；$K$ 在主法线的延长线上。}\label{mp3:fig:l3-e3-s}
\end{figure}


将加速度投影到 $\vct e_n$，
\[
a_n=\frac{\Omega^2r[-(b^2-1)+2b^2]}{\sqrt{1+b^2}}
=\Omega^2r\sqrt{1+b^2}.
\]
因此
\[
\rho=\frac{v^2}{a_n}=r\sqrt{1+b^2}.
\]
主法线含沿 $+\vct e_\varphi$ 的分量，曲率中心不在径向线上。也可在式\eqref{mp3:eq:l3-polar-curv}中代入 $r'=br,r''=b^2r$，得到同一结果。当 $b\to0$ 时，轨迹趋于半径为 $r_0$ 的圆，$a_t\to0$、$\rho\to r_0$；当 $b>0$ 时，即使角速度不变，速率仍随 $r$ 增大。

\end{solution}

\section{刚体平面运动}
各点速度由参考点的平动和绕它的转动共同决定。瞬心简化某一时刻的速度计算，加速度仍从完整的刚体运动公式求出。

\subsection{从转角到各点的轨迹}
对平面刚体取一个材料点 $A$ 为参考点。另一材料点 $P$ 在刚体坐标中的位置为常矢量 $\vct \xi$，则
\[
\vct r_P(t)=\vct r_A(t)+Q(\theta)\vct \xi,\qquad
Q(\theta)=
\begin{pmatrix}\cos\theta&-\sin\theta\\\sin\theta&\cos\theta\end{pmatrix}.
\]
令 $J=\begin{pmatrix}0&-1\\1&0\end{pmatrix}$，直接对矩阵各元素求导，有 $\dot Q=\omega JQ$、$J^2=-I$。再次求导，
\[
\ddot Q=\alpha JQ+\omega J\dot Q
=\alpha JQ+\omega^2J^2Q=(\alpha J-\omega^2I)Q.
\]
将 $J$ 的作用写为 $\vct k\times$，并令 $\vct r=\vct r_P-\vct r_A$，位置式的两次导数便成为
\begin{align}
\vct v_P&=\vct v_A+\omega\vct k\times\vct r,\label{mp3:eq:l3-rigid-v}\\
\vct a_P&=\vct a_A+\alpha\vct k\times\vct r-\omega^2\vct r.\label{mp3:eq:l3-rigid-a}
\end{align}
其中 $\omega=\dot\theta$、$\alpha=\ddot\theta$，逆时针为正。式\eqref{mp3:eq:l3-rigid-v}给出瞬时速度分布。材料点的完整轨迹由 $\vct r_A(t)$ 和转角 $\theta(t)$ 共同确定。

\subsection{速度瞬心及其适用范围}
当 $\omega\ne0$，令平面延拓中的某点 $I$ 的瞬时速度为零，由式\eqref{mp3:eq:l3-rigid-v}解得
\[
\vct r_I-\vct r_A=\frac{\vct k\times\vct v_A}{\omega},\qquad
\vct v_P=\omega\vct k\times(\vct r_P-\vct r_I).
\]
因此 $IP$ 垂直于 $\vct v_P$，且 $v_P=|\omega|IP$。已知两点的非平行速度方向时，过两点分别作速度的垂线，其交点便是瞬心。

瞬心可以位于刚体内部、外部或刚体的几何延拓上。若 $\omega=0$ 而刚体在运动，各点速度相同，瞬心退至无穷远；若此时各点速度也为零，则不能由速度场唯一确定瞬心。

瞬心的位置随时间变化。把每一时刻不同的瞬心连起来，并不能把它们当作同一材料点。即便某一材料点此刻恰在瞬心处，它的加速度一般也不为零。计算加速度时须使用式\eqref{mp3:eq:l3-rigid-a}，或直接对固定材料点的位置求二阶导数。

\begin{figure}[H]\centering
\begin{tikzpicture}[scale=.95]

\coordinate(A)at(0,0);\coordinate(B)at(2,1.7);\coordinate(I)at(0,1.7);
\draw[thick](A)--(B);\draw[guide](A)--(I)--(B);
\fill(A)circle(1.4pt)node[below]{$A$};\fill(B)circle(1.4pt)node[above right]{$B$};
\fill(I)circle(1.4pt)node[above]{$I$};
\draw[axis](A)--(1,0)node[below]{$\vct v_A$};
\draw[axis](B)--(2,{1.7+2/1.7})node[above]{$\vct v_B$};
\begin{scope}[shift={(6,1)}]
\draw[curve](0,0)circle(1);\draw[thick](-1.6,-1)--(1.8,-1);
\fill(0,0)circle(1.4pt)node[above left]{$C$};\fill(0,-1)circle(1.4pt)node[below]{$I$};
\draw[axis](0,0)--(1.4,0)node[right]{$V$};
\draw[axis,thick](0,-1)--(0,-.2)node[midway,left]{$\vct a_I$};
\draw[axis](.4,.4)arc(45:-90:.57);
\node at(.8,-.55){$\omega$};
\end{scope}
\end{tikzpicture}
\caption{速度瞬心的作图与加速度的区别。右图为匀速滚动轮，瞬时接触材料点 $I$ 的速度为零，加速度向上。}\label{mp3:fig:l3-ic}
\end{figure}


\subsection{由刚体运动求轨迹曲率}
求材料点 $P$ 的曲率，可先用式\eqref{mp3:eq:l3-rigid-v}、\eqref{mp3:eq:l3-rigid-a}求出 $\vct v_P,\vct a_P$，再代入式\eqref{mp3:eq:l3-vacurv}。瞬心到 $P$ 的距离满足 $v_P=|\omega|IP$；轨迹曲率还取决于 $P$ 的法向加速度，所以曲率半径一般不同于 $IP$。

纯滚动还需写接触处的速度条件。对圆轮沿固定直线滚动，轮心位移与自转角满足 $x_C=R\vartheta$；对内滚、外滚，接触点到转动中心的方向不同，角速度关系的符号也随之改变。下列例题从接触速度推导，不预设轮廓上的点走圆周。

\par\medskip\noindent\begin{minipage}{\textwidth}
\begin{example}[轮内标记点的轨迹与拐点]\label{mp3:example:l3-e4}

半径为 $R$ 的圆轮以轮心速度 $V=R\Omega$ 向右匀速纯滚动，$\Omega>0$。材料点 $P$ 到轮心的距离为 $d$，$0<d\le R$，初始位于轮心正下方。取顺时针转过的角度为 $\vartheta=\Omega t$，见图\ref{mp3:fig:l3-e4-q}。求 $P$ 的轨迹、速度、加速度与曲率半径；当 $d<R$ 时，确定拐点的位置。

\end{example}
\begin{figure}[H]\centering
\begin{tikzpicture}[scale=.85]

\draw[thick](-.3,0)--(6,0);\draw[axis](-.3,-.35)--(6,-.35)node[right]{$x$};
\coordinate(C)at(2.5,1.3);\coordinate(P)at(1.65,1.3);
\draw[curve](C)circle(1.3);\draw[guide](2.5,0)--(C);
\draw[thick](C)--(P);
\fill(C)circle(1.4pt)node[above right]{$C$};\fill(P)circle(1.8pt)node[left]{$P$};
\fill(2.5,0)circle(1.4pt)node[below right]{$I$};
\draw[axis](C)--(4.3,1.3)node[right]{$V$};
\draw(2.5,.8)arc(-90:-180:.5);\node at(1.95,.75){$\vartheta$};
\node[above]at(2.02,1.55){$d$};\node[right]at(2.55,.55){$R$};
\node at(5,2.6){$\vartheta=\pi/2$};

\end{tikzpicture}
\caption{图示为 $\vartheta=\pi/2$ 的构形；$P$ 是固定在轮内的材料点，$I$ 是当前接触位置。}\label{mp3:fig:l3-e4-q}
\end{figure}
\end{minipage}\par\medskip
\begin{solution}
轮心位置为 $(R\vartheta,R)$。从初始向下的半径顺时针转过 $\vartheta$，得到
\begin{align*}
x_P&=R\vartheta-d\sin\vartheta,& y_P&=R-d\cos\vartheta,\\
\vct v_P&=\Omega(R-d\cos\vartheta,\ d\sin\vartheta),&
\vct a_P&=d\Omega^2(\sin\vartheta,\ \cos\vartheta).
\end{align*}
加速度指向当前轮心 $C$，因为轮心匀速平移而 $P$ 相对轮心匀速转动。但速度一般不垂直于 $CP$，所以 $C$ 一般不是轨迹的曲率中心。

\begin{figure}[H]\centering
\begin{tikzpicture}[scale=.85]

\draw[axis](0,0)--(6.9,0)node[right]{$x$};
\draw[curve,domain=0:360,samples=180]plot({.95*(\x*pi/180-.65*sin(\x))},{.95*(1-.65*cos(\x))});
\fill(0,{.95*(1-.65)})circle(1.3pt);
\fill({.95*2*pi},{.95*(1-.65)})circle(1.3pt);
\draw[guide](0,0)--(0,{.95*(1-.65)});
\draw[guide]({.95*2*pi},0)--({.95*2*pi},{.95*(1-.65)});
\node[below]at(0,0){$0$};\node[below]at({.95*2*pi},0){$2\pi R$};
\draw[axis]({.95*(105*pi/180-.65*sin(105))},{.95*(1-.65*cos(105))})--({.95*(115*pi/180-.65*sin(115))},{.95*(1-.65*cos(115))});
\pgfmathsetmacro{\drawangle}{acos(.65)}
\fill({.95*(\drawangle*pi/180-.65*sin(\drawangle))},{.95*(1-.65*cos(\drawangle))})circle(1.5pt)node[left]{$F_1$};
\fill({.95*((360-\drawangle)*pi/180-.65*sin(360-\drawangle))},{.95*(1-.65*cos(360-\drawangle))})circle(1.5pt)node[right]{$F_2$};
\node at(2.9,2.15){$d<R$：短幅摆线};
\begin{scope}[shift={(7.7,.45)},scale=.75]
\fill(0,0)circle(1.5pt)node[below]{$P$};
\draw[axis](0,0)--(1.3,.845)node[above right]{$\vct v_P$};
\draw[axis](0,0)--(1.3,0)node[midway,below]{$\vct a_P$};
\draw[guide](0,0)--(1.8,0);\fill(1.8,0)circle(1.5pt)node[above right]{$C$};
\node at(.4,1.9){$\vartheta=\pi/2$};
\end{scope}

\end{tikzpicture}
\caption{左图两端对应 $\vartheta=0$ 与 $2\pi$，箭头表示参数增大；右图在 $\vartheta=\pi/2$ 比较速度和加速度，后者指向轮心。}\label{mp3:fig:l3-e4-s}
\end{figure}


由图\ref{mp3:fig:l3-e4-s}中的分量，
\[
v^2=\Omega^2(R^2+d^2-2Rd\cos\vartheta),
\]
\[
(\vct v_P\times\vct a_P)\cdot\vct k
=d\Omega^3[(R-d\cos\vartheta)\cos\vartheta-d\sin^2\vartheta]
=d\Omega^3(R\cos\vartheta-d).
\]
在分母不为零且 $v>0$ 的位置，
\[
\rho=
\frac{(R^2+d^2-2Rd\cos\vartheta)^{3/2}}
{d|R\cos\vartheta-d|}.
\]
当 $d<R$，$v\ge\Omega(R-d)>0$。每一周期的两个拐点满足
$\cos\vartheta=d/R$，图中以 $F_1,F_2$ 标出；在那里加速度与速度平行，曲率为零，$\rho\to\infty$。越过每个 $F_i$ 时，$\vct v_P\times\vct a_P$ 的有向分量变号，轨迹凹向随之改变。

当 $d=R$，轨迹成为普通摆线。在 $0<\vartheta<2\pi$ 内，
\[
v=2R\Omega\sin\frac\vartheta2,\qquad
\rho=4R\sin\frac\vartheta2.
\]
顶点 $\vartheta=\pi$ 处 $\rho=4R$，而速度瞬心到 $P$ 的距离仅为 $2R$。端点 $\vartheta=0,2\pi$ 处 $v=0$，是尖点，只能讨论邻域的曲率极限。

\end{solution}

\par\medskip\noindent\begin{minipage}{\textwidth}
\begin{example}[滑梯中点的曲率与瞬心]\label{mp3:example:l3-e5}

长为 $l$ 的刚杆 $AB$ 两端沿互相垂直的固定导轨滑动。$A=(l\cos\theta,0)$，$B=(0,l\sin\theta)$，$0<\theta<\pi/2$。端点 $A$ 以恒定速率 $u$ 向右运动。求杆的角速度、角加速度，以及中点 $M$ 的轨迹、速率、切向和法向加速度，见图\ref{mp3:fig:l3-e5-q}。

\end{example}
\begin{figure}[H]\centering
\begin{tikzpicture}[scale=.95]

\draw[thick](-.2,0)--(4,0);\draw[thick](0,-.2)--(0,3.4);
\coordinate(A)at(3,0);\coordinate(B)at(0,2.4);\coordinate(M)at(1.5,1.2);
\draw[very thick](A)--(B);\draw[guide](A)--(3,2.4)--(B);
\fill(A)circle(1.5pt)node[below]{$A$};\fill(B)circle(1.5pt)node[left]{$B$};
\fill(M)circle(1.5pt)node[above right]{$M$};\fill(3,2.4)circle(1.5pt)node[above]{$I$};
\draw[axis](A)--(4,0)node[above]{$u$};
\draw[axis](B)--(0,{2.4-3/2.4})node[left]{$\vct v_B$};
\draw(.6,0)arc(0:38.66:.6);\node at(.9,.3){$\theta$};
\draw[guide](0,0)--(M);\node[below left]at(0,0){$O$};
\node at(3.3,3.1){$AB=l$};

\end{tikzpicture}
\caption{端点速度的垂线相交于 $I=(l\cos\theta,l\sin\theta)$。图中 $\theta$ 是 $OM$ 与水平线的夹角。}\label{mp3:fig:l3-e5-q}
\end{figure}
\end{minipage}\par\medskip
\begin{solution}
由 $\dot x_A=-l\sin\theta\,\dot\theta=u$，
\[
\dot\theta=-\frac{u}{l\sin\theta},\qquad
\ddot\theta=\frac{u\cos\theta}{l\sin^2\theta}\dot\theta
=-\frac{u^2\cos\theta}{l^2\sin^3\theta}.
\]
杆的有向姿态角可取 $\pi-\theta$，故逆时针角速度为
$\omega=-\dot\theta=u/(l\sin\theta)$，角加速度为
$\alpha=-\ddot\theta=u^2\cos\theta/(l^2\sin^3\theta)$。

中点位置满足
\[
\vct r_M=\frac l2(\cos\theta,\sin\theta),\qquad OM=\frac l2.
\]
所以 $M$ 沿第一象限的四分之一圆向右下方运动，轨迹曲率半径恒为 $l/2$。图\ref{mp3:fig:l3-e5-s}中的 $\vct e_t=(\sin\theta,-\cos\theta)$、$\vct e_n=(-\cos\theta,-\sin\theta)$ 分别沿实际运动与指向 $O$。

\begin{figure}[H]\centering
\begin{tikzpicture}[scale=.95]

\draw[axis](-.2,0)--(3.5,0)node[right]{$x$};\draw[axis](0,-.2)--(0,3.2)node[above]{$y$};
\draw[curve](2.1,0)arc(0:90:2.1);
\coordinate(M)at({2.1*cos(40)},{2.1*sin(40)});
\draw[guide](0,0)--({3.2*cos(40)},{3.2*sin(40)});
\fill(M)circle(1.5pt)node[above right]{$M$};
\node[below left]at(0,0){$O$};
\draw[axis](M)--++({1.1*sin(40)},{-1.1*cos(40)})node[right]{$\vct e_t$};
\draw[axis](M)--++({-1.1*cos(40)},{-1.1*sin(40)})node[left]{$\vct e_n$};
\begin{scope}[shift={(5.6,1.6)}]
\fill(0,0)circle(1.5pt)node[above]{$M$};
\draw[axis](0,0)--(.8,-.953)node[right]{$a_t\vct e_t$};
\draw[axis](0,0)--(-.8,-.671)node[left]{$a_n\vct e_n$};
\draw[axis,thick](0,0)--(0,-1.624)node[below]{$\vct a_M$};
\draw[guide](.8,-.953)--(0,-1.624)--(-.8,-.671);
\end{scope}

\end{tikzpicture}
\caption{左图为 $M$ 的实际轨迹；右图按 $\theta=40^\circ$ 的分量比例合成加速度。}\label{mp3:fig:l3-e5-s}
\end{figure}


直接求导，中点的竖直速度为 $(l/2)\cos\theta\dot\theta=-(u/2)\cot\theta$。再次求导，$a_{My}=(u/2)\csc^2\theta\dot\theta$，所以
\[
\vct v_M=\left(\frac u2,-\frac u2\cot\theta\right),\qquad
\vct a_M=\left(0,-\frac{u^2}{2l\sin^3\theta}\right).
\]
于是
\[
\begin{aligned}
v_M&=\frac{u}{2\sin\theta},\\
a_t=\dot v_M&=-\frac{u\cos\theta}{2\sin^2\theta}\dot\theta
=\frac{u^2\cos\theta}{2l\sin^3\theta},\\
a_n=\frac{v_M^2}{l/2}&=\frac{u^2}{2l\sin^2\theta}.
\end{aligned}
\]
两分量相加时，水平分量正好相消，合加速度竖直向下。

速度瞬心 $I$ 到 $M$ 的距离也为 $l/2$。它位于 $OM$ 的延长线上，与曲率中心 $O$ 分居 $M$ 两侧。若把当前位于 $I$ 的刚体延拓材料点记为 $I^\ast$，由式\eqref{mp3:eq:l3-rigid-a}，
\[
\begin{aligned}
\vct a_{I^\ast}
&=\alpha\vct k\times(0,l\sin\theta)-\omega^2(0,l\sin\theta)\\
&=(-\alpha l\sin\theta,-\omega^2l\sin\theta)\\
&=\left(-\frac{u^2\cos\theta}{l\sin^2\theta},
-\frac{u^2}{l\sin\theta}\right)\ne0.
\end{aligned}
\]
在该式中参考点为 $A$，$\vct r_{I^\ast}-\vct r_A=(0,l\sin\theta)$。该材料点虽在此刻静止，下一瞬间仍有非零的速度变化。

\end{solution}

\par\medskip\noindent\begin{minipage}{\textwidth}
\begin{example}[内滚小圆产生直线运动]\label{mp3:example:l3-e6}

半径为 $b$ 的小圆在半径为 $a$ 的固定大圆内无滑动滚动，$a>b$。大圆心为 $O$，小圆心为 $C$，$\angle xOC=\vartheta=\Omega t$，$\Omega>0$。初始时材料点 $P$ 与大圆最右端接触，见图\ref{mp3:fig:l3-e6-q}。求 $P$ 的轨迹；证明 $a=2b$ 时 $P$ 始终沿一条直径运动，并求其速度与加速度。

\end{example}
\begin{figure}[H]\centering
\begin{tikzpicture}[scale=.95]

\draw[curve](0,0)circle(2.1);\draw[axis](-2.5,0)--(3.1,0)node[right]{$x$};
\coordinate(C)at({1.05*cos(45)},{1.05*sin(45)});\coordinate(I)at({2.1*cos(45)},{2.1*sin(45)});
\draw[curve](C)circle(1.05);\draw[guide](0,0)--(I);
\fill(0,0)circle(1.5pt)node[below]{$O$};\fill(C)circle(1.5pt)node[above]{$C$};
\fill(I)circle(1.5pt)node[above right]{$I$};
\fill(2.1,0)circle(1.5pt)node[below right]{$P(0)$};
\draw(.55,0)arc(0:45:.55);\node at(.8,.25){$\vartheta$};
\node at(3.6,1.5){$OC=a-b$};
\node at(3.6,.7){$CI=b$};

\end{tikzpicture}
\caption{内滚接触点 $I$ 位于 $OC$ 的外延线上。图示半径比取 $a=2b$。}\label{mp3:fig:l3-e6-q}
\end{figure}
\end{minipage}\par\medskip
\begin{solution}
令 $c=a-b$，则 $\vct r_C=c(\cos\vartheta,\sin\vartheta)$，
$\vct v_C=c\Omega\vct e_\varphi$。接触材料点相对轮心的矢量为 $b\vct e_r$。固定大圆要求
\[
\vct v_I=\vct v_C+\omega\vct k\times(b\vct e_r)=0,
\quad \omega=-\frac cb\Omega.
\]
小圆的自转方向与公转方向相反。若 $\psi$ 为材料半径 $CP$ 与固定 $x$ 轴的夹角，由 $\psi(0)=0$ 得 $\psi=-c\vartheta/b$。所以
\[
x_P=c\cos\vartheta+b\cos\frac{c\vartheta}{b},\qquad
y_P=c\sin\vartheta-b\sin\frac{c\vartheta}{b}.
\]
这组参数式在整个滚动过程中有效，闭合与否取决于 $c/b$ 的有理性。

\begin{figure}[H]\centering
\begin{tikzpicture}[scale=.95]

\draw[curve](0,0)circle(2);\draw[thick](-2,0)--(2,0);
\foreach \drawangle in {35,90,145}{
\coordinate(C)at({cos(\drawangle)},{sin(\drawangle)});
\coordinate(P)at({2*cos(\drawangle)},0);
\draw[guide](C)circle(1);\draw[guide](0,0)--(C);\draw[thick](C)--(P);
\fill(C)circle(1.3pt);\fill(P)circle(1.5pt);
}
\node[below]at(0,0){$O$};\node[above left,fill=white,inner sep=1pt]at(0,1){$C_2$};
\node[below]at(1.638,0){$P_1$};\node[below]at(-1.638,0){$P_3$};
\node at(5.2,1.4){$OC=CP=b$};
\node at(5.2,.4){$y_P=b\sin\vartheta-b\sin\vartheta=0$};
\node at(5.2,-.6){$x_P=2b\cos\vartheta$};

\end{tikzpicture}
\caption{三个时刻的构形叠画。每条细圆都是对应时刻的小圆；三个标记位置 $P_1,O,P_3$ 共线。}\label{mp3:fig:l3-e6-s}
\end{figure}


当 $a=2b$，$c=b$，两个正弦项相消，
\[
x_P=2b\cos\vartheta,\qquad y_P=0,
\]
\[
\vct v_P=(-2b\Omega\sin\vartheta,0),\qquad
\vct a_P=(-2b\Omega^2\cos\vartheta,0).
\]
图\ref{mp3:fig:l3-e6-s}中，每一构形都有 $CP$ 的竖直分量恰好抵消 $OC$ 的竖直分量。点 $P$ 沿直径作简谐运动，运动区间为 $[-2b,2b]$。除两端转向点外，轨迹曲率为零；转向点速度为零，不能由式\eqref{mp3:eq:l3-vacurv}定义该瞬间的单位切向量。

\end{solution}

\section{空间运动}
空间转动的各轴一般不在同一平面内。转轴、角速度和位置用矢量统一表示，再由投影得到题目要求的运动分量。

\subsection{柱坐标与随位置转动的基矢}
空间点用柱坐标 $(r,\varphi,z)$ 表示，位置矢量为
$\vct r=r\vct e_r+z\vct e_z$。其中 $\vct e_z$ 固定，
$\dot{\vct e}_r=\dot\varphi\vct e_\varphi$、
$\dot{\vct e}_\varphi=-\dot\varphi\vct e_r$。对位置逐项求导得到速度，再对速度中的三项分别求导：
\[
\begin{aligned}
\frac{\dd}{\dd t}(\dot r\vct e_r)
&=\ddot r\vct e_r+\dot r\dot\varphi\vct e_\varphi,\\
\frac{\dd}{\dd t}(r\dot\varphi\vct e_\varphi)
&=(\dot r\dot\varphi+r\ddot\varphi)\vct e_\varphi
-r\dot\varphi^2\vct e_r,\\
\frac{\dd}{\dd t}(\dot z\vct e_z)&=\ddot z\vct e_z.
\end{aligned}
\]
合并相同方向的项，
\begin{align}
\vct v&=\dot r\vct e_r+r\dot\varphi\vct e_\varphi+\dot z\vct e_z,\label{mp3:eq:l3-cyl-v}\\
\vct a&=(\ddot r-r\dot\varphi^2)\vct e_r
+(r\ddot\varphi+2\dot r\dot\varphi)\vct e_\varphi+\ddot z\vct e_z.\label{mp3:eq:l3-cyl-a}
\end{align}
图\ref{mp3:fig:l3-cylinder}中 $\vct e_r,\vct e_\varphi$ 位于经过 $P$ 的水平面内。$-\!r\dot\varphi^2\vct e_r$ 指向竖直轴，不一定指向三维曲线的曲率中心。若 $\dot z$ 或 $\dot r$ 不为零，还须对完整的 $\vct v,\vct a$ 作自然分解。

\begin{figure}[H]\centering
\begin{tikzpicture}[scale=.95]

\begin{scope}[x={(.95cm,-.2cm)},y={(.25cm,.65cm)},z={(0cm,1cm)}]
\draw[axis](0,0,0)--(3.6,0,0)node[right]{$x$};
\draw[axis](0,0,0)--(0,3,0)node[right]{$y$};
\draw[axis](0,0,0)--(0,0,3.3)node[above]{$z$};
\coordinate(H)at(2,1.4,0);\coordinate(P)at(2,1.4,2.1);
\draw[guide](0,0,0)--(H)--(P)--(0,0,2.1);
\draw[thick](0,0,0)--(P);
\fill(H)circle(1.3pt)node[below]{$H$};\fill(P)circle(1.5pt)node[below right]{$P$};
\draw[axis](P)--++(.819,.573,0)node[right]{$\vct e_r$};
\draw[axis](P)--++(-.573,.819,0)node[above left]{$\vct e_\varphi$};
\draw[axis](P)--++(0,0,1)node[above]{$\vct e_z$};
\node[above]at(1.3,.91,0){$r$};\node[left]at(0,0,2.1){$z$};
\draw[domain=0:35,samples=30]plot({.7*cos(\x)},{.7*sin(\x)},0);
\node at(1,.25,0){$\varphi$};\node[below left]at(0,0,0){$O$};
\end{scope}
\end{tikzpicture}
\caption{柱坐标的投影关系。三维图采用固定斜投影；标出的三个基矢在实际空间中互相垂直且均为单位向量。}\label{mp3:fig:l3-cylinder}
\end{figure}


\subsection{球坐标的速度与加速度}
以下以 $r$ 表示到原点的距离，$\theta$ 为从 $+z$ 轴量起的极角，$\varphi$ 为方位角。它与柱坐标的径向距离不同，使用时由上下文区分。球坐标基矢为
\begin{align*}
\vct e_r&=(\sin\theta\cos\varphi,\sin\theta\sin\varphi,\cos\theta),\\
\vct e_\theta&=(\cos\theta\cos\varphi,\cos\theta\sin\varphi,-\sin\theta),\\
\vct e_\varphi&=(-\sin\varphi,\cos\varphi,0).
\end{align*}
对这些笛卡尔分量求导，得到
\begin{align*}
\dot{\vct e}_r&=\dot\theta\vct e_\theta+\dot\varphi\sin\theta\vct e_\varphi,\\
\dot{\vct e}_\theta&=-\dot\theta\vct e_r+\dot\varphi\cos\theta\vct e_\varphi,\\
\dot{\vct e}_\varphi&=-\dot\varphi(\sin\theta\vct e_r+\cos\theta\vct e_\theta).
\end{align*}
由 $\vct r=r\vct e_r$，
\[
\vct v=\dot r\vct e_r+r\dot\theta\vct e_\theta
+r\sin\theta\dot\varphi\vct e_\varphi.
\]
再求导时，三个标量系数和三根基矢都随时间变化。分别展开为
\[
\begin{aligned}
\frac{\dd}{\dd t}(\dot r\vct e_r)
&=\ddot r\vct e_r+\dot r\dot\theta\vct e_\theta
+\dot r\sin\theta\dot\varphi\vct e_\varphi,\\
\frac{\dd}{\dd t}(r\dot\theta\vct e_\theta)
&=(\dot r\dot\theta+r\ddot\theta)\vct e_\theta
-r\dot\theta^2\vct e_r
+r\dot\theta\cos\theta\dot\varphi\vct e_\varphi,\\
\frac{\dd}{\dd t}(r\sin\theta\dot\varphi\vct e_\varphi)
&=(\dot r\sin\theta\dot\varphi+r\cos\theta\dot\theta\dot\varphi
+r\sin\theta\ddot\varphi)\vct e_\varphi\\
&\quad-r\sin\theta\dot\varphi^2
(\sin\theta\vct e_r+\cos\theta\vct e_\theta).
\end{aligned}
\]
按三基矢归并，
\begin{align}
a_r&=\ddot r-r\dot\theta^2-r\sin^2\theta\dot\varphi^2,\nonumber\\
a_\theta&=r\ddot\theta+2\dot r\dot\theta
-r\sin\theta\cos\theta\dot\varphi^2,\label{mp3:eq:l3-spherical-a}\\
a_\varphi&=r\sin\theta\ddot\varphi
+2\dot r\sin\theta\dot\varphi
+2r\cos\theta\dot\theta\dot\varphi.\nonumber
\end{align}
球面上的运动虽有 $r=\text{常量}$，仍有两个角变量。沿纬线绕行时，$a_\theta$ 一般不为零，径向分量 $a_r$ 也不是法向加速度的全部。在两极 $\sin\theta=0$ 处方位角失去唯一性，应回到笛卡尔坐标或换一个局部坐标图。

\subsection{空间曲线的自然标架}
式\eqref{mp3:eq:l3-curv}与\eqref{mp3:eq:l3-natural-a}在三维中仍成立。在 $\kappa\ne0$ 处定义副法向量
$\vct e_b=\vct e_t\times\vct e_n$。副法线表示局部密切平面的法向，其转动用挠率 $\tau$ 描述：
\[
\frac{\dd\vct e_b}{\dd s}=-\tau\vct e_n.
\]
由单位长度可知 $\vct e_n\cdot\vct e_{n,s}=0$。再分别对 $\vct e_t\cdot\vct e_n=0$、$\vct e_b\cdot\vct e_n=0$ 求导，
\[
\vct e_t\cdot\vct e_{n,s}=-\vct e_{t,s}\cdot\vct e_n=-\kappa,
\qquad
\vct e_b\cdot\vct e_{n,s}=-\vct e_{b,s}\cdot\vct e_n=\tau.
\]
这确定了 $\vct e_{n,s}$ 的三个分量，故全部导数为
\[
\frac{\dd\vct e_t}{\dd s}=\kappa\vct e_n,\quad
\frac{\dd\vct e_n}{\dd s}=-\kappa\vct e_t+\tau\vct e_b,\quad
\frac{\dd\vct e_b}{\dd s}=-\tau\vct e_n.
\]
设 $\vct j=\dot{\vct a}$ 为加速度的时间导数。对 $\vct a=\dot v\vct e_t+v^2\kappa\vct e_n$ 求导，
\[
\begin{aligned}
\vct j={}&\ddot v\vct e_t+\dot v(v\kappa\vct e_n)
+(2v\dot v\kappa+v^2\dot\kappa)\vct e_n\\
&+v^2\kappa\,v(-\kappa\vct e_t+\tau\vct e_b).
\end{aligned}
\]
合并得到
\[
\vct j=(\ddot v-v^3\kappa^2)\vct e_t
+(3v\dot v\kappa+v^2\dot\kappa)\vct e_n
+v^3\kappa\tau\vct e_b.
\]
由于 $\vct v\times\vct a=v^3\kappa\vct e_b$，上式给出
\[
(\vct v\times\vct a)\cdot\vct j=v^6\kappa^2\tau,
\qquad |\vct v\times\vct a|^2=v^6\kappa^2.
\]
因此在 $v>0,\kappa>0$ 时，
\begin{equation}\label{mp3:eq:l3-torsion}
\tau=\frac{(\vct v\times\vct a)\cdot\vct j}{|\vct v\times\vct a|^2}.
\end{equation}
加速度本身没有副法向分量，三维扭转首先出现在 $\vct j$ 中。在一个连通的正则区间内，若 $\kappa>0$ 且 $\tau$ 恒为零，副法向量固定，轨迹就在一个固定平面内。恒定曲率本身不能推出圆周运动，圆柱螺线便是反例。

\par\medskip\noindent\begin{minipage}{\textwidth}
\begin{example}[匀速圆柱螺线]\label{mp3:example:l3-e7}

质点满足 $x=a\cos\Omega t,\ y=a\sin\Omega t,\ z=h\Omega t$，其中 $a,h,\Omega>0$，见图\ref{mp3:fig:l3-e7-q}。求速率、加速度、曲率半径与挠率；写出自然标架，并说明曲率圆心在什么位置。

\end{example}
\begin{figure}[H]\centering
\begin{tikzpicture}[scale=.95]

\begin{scope}[x={(.85cm,-.25cm)},y={(.55cm,.3cm)},z={(0cm,.65cm)}]
\draw[axis](0,0,0)--(2.2,0,0)node[below right]{$x$};
\draw[axis](0,0,0)--(0,2.1,0)node[right]{$y$};
\draw[axis](0,0,0)--(0,0,5.7)node[above]{$z$};
\draw[curve,domain=0:660,samples=220]plot({1.3*cos(\x)},{1.3*sin(\x)},{.007*\x});
\draw[guide,domain=0:360,samples=70]plot({1.3*cos(\x)},{1.3*sin(\x)},0);
\draw[guide](1.3,0,0)--(1.3,0,2.52);
\draw[guide](1.3,0,0)--(2.25,0,0);
\draw[guide](1.3,0,2.52)--(2.25,0,2.52);
\fill(1.3,0,0)circle(1.2pt);\fill(1.3,0,2.52)circle(1.2pt);
\draw[|-|](2.25,0,0)--(2.25,0,2.52)node[midway,right]{$2\pi h$};
\node[below]at(.65,0,0){$a$};\draw[guide](0,0,0)--(1.3,0,0);
\end{scope}
\node at(5.1,2.7){$r=a$};
\node at(5.1,1.7){$z=h\varphi$};

\end{tikzpicture}
\caption{每绕一周上升 $2\pi h$。图中的虚线圆是水平投影。}\label{mp3:fig:l3-e7-q}
\end{figure}
\end{minipage}\par\medskip
\begin{solution}
令 $L=\sqrt{a^2+h^2}$，$\varphi=\Omega t$。柱坐标公式给出
\[
\vct v=\Omega(a\vct e_\varphi+h\vct e_z),\qquad
\vct a=-a\Omega^2\vct e_r.
\]
速率 $v=\Omega L$ 恒定，故切向加速度为零，主法线为 $-\vct e_r$。于是
\[
\rho=\frac{v^2}{a\Omega^2}=\frac{a^2+h^2}{a},\qquad
\vct e_t=\frac{a\vct e_\varphi+h\vct e_z}{L},\quad
\vct e_n=-\vct e_r.
\]
由右手叉积关系，
\[
\vct e_b=\frac{a\vct e_z-h\vct e_\varphi}{L}.
\]
图\ref{mp3:fig:l3-e7-s}左图把这三个方向移至相同的水平位置比较，切向量沿螺线上升，副法线则有与方位运动反向的分量。

\begin{figure}[H]\centering
\begin{tikzpicture}[scale=.95]

\begin{scope}[x={(.9cm,-.25cm)},y={(.6cm,.35cm)},z={(0cm,1cm)}]
\coordinate(P)at(1.5,0,0);
\fill(P)circle(1.5pt)node[below right]{$P$};
\draw[axis](P)--++(0,.8,.6)node[above]{$\vct e_t$};
\draw[axis](P)--++(-1,0,0)node[left]{$\vct e_n$};
\draw[axis](P)--++(0,-.6,.8)node[left]{$\vct e_b$};
\end{scope}
\begin{scope}[shift={(5,0)}]
\draw[guide](-1.4,0)--(2.3,0);
\fill(1.7,0)circle(1.5pt)node[above]{$P$};
\fill(0,0)circle(1.5pt)node[above]{$H$};
\fill(-.95625,0)circle(1.5pt)node[above]{$K$};
\draw[|-|](0,-.55)--(1.7,-.55)node[midway,fill=white]{$a$};
\draw[|-|](-.95625,-1.1)--(1.7,-1.1)node[midway,fill=white]{$\rho=(a^2+h^2)/a$};
\draw[guide](0,-.25)--(0,1.3)node[above]{竖直轴};
\end{scope}

\end{tikzpicture}
\caption{左图取 $h/a=3/4$ 表示自然标架，三根基矢实际等长；右图在同一水平面内确定曲率中心 $K$。}\label{mp3:fig:l3-e7-s}
\end{figure}


因为 $\dd\vct e_\varphi/\dd\varphi=-\vct e_r$ 且 $\dd s/\dd\varphi=L$，
\[
\frac{\dd\vct e_b}{\dd s}
=\frac{h}{L^2}\vct e_r=-\frac{h}{L^2}\vct e_n,
\qquad
\tau=\frac{h}{a^2+h^2}.
\]
曲率中心为 $K=P-\rho\vct e_r$。在 $P$ 所在水平面内，
\[
\vct r_K=-\frac{h^2}{a}\vct e_r+z\vct e_z.
\]
图中 $PK$ 穿过竖直轴后还要延长 $h^2/a$，因此轴上的点 $H$ 不是曲率中心。仅当 $h=0$ 时，螺线退化成水平圆，才有 $K=H$、$\tau=0$。

\end{solution}

\par\medskip\noindent\begin{minipage}{\textwidth}
\begin{example}[圆锥内表面的圆周运动与小扰动]\label{mp3:example:l3-e8}

光滑固定圆锥的轴竖直，顶点为原点，半顶角为 $\beta$，$0<\beta<\pi/2$，内表面向上张开。质量 $m$ 的小珠沿内表面运动。用从顶点沿母线量起的距离 $s$ 与方位角 $\varphi$ 描述位置，见图\ref{mp3:fig:l3-e8-q}。求高度固定的水平圆周运动条件及支持力；在保持竖直角动量不变的小扰动下，求 $s$ 的小振动频率。

\end{example}
\begin{figure}[H]\centering
\begin{tikzpicture}[scale=.95]

\draw[axis](0,-.2)--(0,3.8)node[above]{$z$};
\draw[curve](0,0)--(-2.1,3.2);\draw[curve](0,0)--(2.1,3.2);
\draw[guide](0,2.2)ellipse(1.444 and .3);
\coordinate(P)at(1.444,2.2);\fill(P)circle(2pt)node[right]{$P$};
\draw[thick](0,0)--(P);\node[right]at(.72,1.1){$s$};
\node[below]at(0,0){$O$};\draw(0,.7)arc(90:56.72:.7);
\node at(.25,.9){$\beta$};
\draw[guide](P)--(1.444,2.85);
\draw[|-|](0,2.85)--(1.444,2.85)node[midway,fill=white]{$s\sin\beta$};
\node at(4,2.9){$z=s\cos\beta$};

\end{tikzpicture}
\caption{圆锥的子午截面与一条水平圆轨道。虚线椭圆表示空间水平圆的投影。}\label{mp3:fig:l3-e8-q}
\end{figure}
\end{minipage}\par\medskip
\begin{solution}
在柱坐标中 $r=s\sin\beta,z=s\cos\beta$。沿母线向上的单位向量与指向内侧的单位法向量为
\[
\vct e_s=\sin\beta\vct e_r+\cos\beta\vct e_z,\qquad
\vct n=-\cos\beta\vct e_r+\sin\beta\vct e_z.
\]
将式\eqref{mp3:eq:l3-cyl-a}投影到这两个方向和 $\vct e_\varphi$，分别得到
\begin{align*}
m(\ddot s-s\sin^2\beta\dot\varphi^2)&=-mg\cos\beta,\\
m\sin\beta(s\ddot\varphi+2\dot s\dot\varphi)&=0,\\
N-mg\sin\beta&=ms\sin\beta\cos\beta\dot\varphi^2.
\end{align*}
第二式乘以 $s\sin\beta$ 后可写成乘积的导数：
\[
\frac{\dd}{\dd t}(ms^2\sin^2\beta\dot\varphi)
=m\sin^2\beta(2s\dot s\dot\varphi+s^2\ddot\varphi)=0.
\]
因此 $\ell=ms^2\sin^2\beta\dot\varphi$ 为常量，即竖直轴角动量守恒。

\begin{figure}[H]\centering
\begin{tikzpicture}[scale=.95]

\draw[curve](0,0)--(2,3);
\coordinate(P)at(1.3,1.95);\fill(P)circle(2pt)node[right]{$P$};
\draw[axis](P)--(.2,2.683)node[above left]{$N_0\vct n$};
\draw[axis](P)--(1.3,1.217)node[below]{$mg$};
\draw[axis,thick](P)--(.2,1.95)node[left]{$m v_0^2/r_0$};
\draw[guide](.2,2.683)--(.2,1.95);
\node at(4.5,2.8){$\vct e_s=(\sin\beta,\cos\beta)$};
\node at(4.5,1.8){$\vct n=(-\cos\beta,\sin\beta)$};
\draw[axis](3.1,.2)--(6,.2)node[right]{$\eta$};
\draw[axis](4.5,-.6)--(4.5,1.05)node[above]{$\ddot\eta$};
\draw[curve](3.5,.9)--(5.5,-.5);

\end{tikzpicture}
\caption{圆周运动中的受力方向，以及扰动加速度与位移 $\eta$ 的反向关系。左图的两力合成为水平向心合力。}\label{mp3:fig:l3-e8-s}
\end{figure}


对 $s=s_0$ 的水平圆周，设角速度为 $\Omega_0$。第一式给出
\[
s_0\sin^2\beta\,\Omega_0^2=g\cos\beta.
\]
轨道半径 $r_0=s_0\sin\beta$，高度 $z_0=s_0\cos\beta$，故
\[
v_0^2=g r_0\cot\beta=g z_0,\qquad N_0=\frac{mg}{\sin\beta}.
\]
图\ref{mp3:fig:l3-e8-s}中，$N_0$ 的竖直分量平衡重力，水平分量提供向心加速度。

扰动后 $\ell$ 守恒，角速度不能继续当作固定的 $\Omega_0$。将其消去，
\[
\ddot s=\frac{\ell^2}{m^2s^3\sin^2\beta}-g\cos\beta.
\]
令 $s=s_0+\eta$，$|\eta|\ll s_0$。圆周条件给出$\ell^2/(m^2s_0^3\sin^2\beta)=g\cos\beta$，所以
\[
\ddot\eta=g\cos\beta\left[\left(1+\frac\eta{s_0}\right)^{-3}-1\right]
=-\frac{3g\cos\beta}{s_0}\eta
+O\!\left(\frac{g\eta^2}{s_0^2}\right).
\]
保留一阶项，得到
\[
\ddot\eta=-\frac{3g\cos\beta}{s_0}\eta,\qquad
\omega_s=\sqrt{\frac{3g\cos\beta}{s_0}}
=\sqrt3\,\Omega_0\sin\beta.
\]
频率为实数，表明该圆轨道对这种子午方向的小扰动稳定。一般运动中
$N=m[g\sin\beta+s\sin\beta\cos\beta\,\dot\varphi^2]>0$，
与内表面的单侧接触相容。

\end{solution}

\par\medskip\noindent\begin{minipage}{\textwidth}
\begin{example}[绕竖直直径转动的光滑圆环]\label{mp3:example:l3-e9}

半径为 $a$ 的光滑圆环绕其竖直直径以恒定角速度 $\Omega$ 转动。质量 $m$ 的小珠穿在环上，先在右半环用从最低点量起的角度 $\theta$ 描述位置，$0\le\theta\le\pi$，左半环由对称关系补出，见图\ref{mp3:fig:l3-e9-q}。求 $\theta$ 的运动方程、相对平衡位置及其稳定性，并求圆环约束力的两个分量。

\end{example}
\begin{figure}[H]\centering
\begin{tikzpicture}[scale=.95]

\draw[axis](0,-2)--(0,2.4)node[above]{$z$};
\draw[curve](0,0)circle(1.6);\draw[guide](0,0)--(0,-1.6);
\coordinate(P)at({1.6*sin(55)},{-1.6*cos(55)});
\draw[thick](0,0)--(P);\fill(P)circle(2pt)node[right]{$P$};
\fill(0,0)circle(1.4pt)node[left]{$O$};
\draw(0,-.55)arc(-90:-35:.55);\node at(.3,-.72){$\theta$};
\draw[axis,thick](.7,1.85)--(.7,2.65)node[right]{$\vct\Omega=\Omega\vct e_z$};
\node at(4,1){$r=a\sin\theta$};
\node at(4,0){$z=-a\cos\theta$};

\end{tikzpicture}
\caption{圆环在当前时刻的子午截面。$\vct\Omega$ 沿 $+z$；转向由右手定则确定。}\label{mp3:fig:l3-e9-q}
\end{figure}
\end{minipage}\par\medskip
\begin{solution}
在右半环，柱坐标为 $r=a\sin\theta,\ z=-a\cos\theta,\ \varphi=\Omega t$。对应加速度分量为
\begin{align*}
a_r&=a(\cos\theta\,\ddot\theta-\sin\theta\,\dot\theta^2-\Omega^2\sin\theta),\\
a_\varphi&=2a\Omega\cos\theta\,\dot\theta,\\
a_z&=a(\sin\theta\,\ddot\theta+\cos\theta\,\dot\theta^2).
\end{align*}
图\ref{mp3:fig:l3-e9-s}左图的环切向量为
$\vct e_\theta=\cos\theta\vct e_r+\sin\theta\vct e_z$。
圆环约束力在该方向分量为零，重力分量为 $-mg\sin\theta$，故
\[
a(\ddot\theta-\Omega^2\sin\theta\cos\theta)=-g\sin\theta,
\qquad
\ddot\theta=\sin\theta\left(\Omega^2\cos\theta-\frac ga\right).
\]

\begin{figure}[H]\centering
\begin{tikzpicture}[scale=.95]

\draw[curve](0,0)circle(1.5);
\coordinate(P)at({1.5*sin(55)},{-1.5*cos(55)});
\fill(P)circle(1.5pt)node[right]{$P$};
\draw[axis](P)--++({.9*cos(55)},{.9*sin(55)})node[above right]{$\vct e_\theta$};
\draw[axis](P)--++({-.9*sin(55)},{.9*cos(55)})node[above left]{$\vct n$};
\draw[axis](P)--++(0,-1)node[below]{$mg$};
\draw[fill=white](-1,-1.6)circle(.11);\draw(-1.06,-1.66)--(-.94,-1.54);\draw(-1.06,-1.54)--(-.94,-1.66);
\node[left]at(-1.2,-1.6){$\vct e_\varphi$};
\begin{scope}[shift={(5.4,0)}]
\draw[axis](-.2,-2)--(3,-2)node[right]{$\Omega/\sqrt{g/a}$};
\draw[axis](0,-2.2)--(0,1.9)node[above]{$\theta_\ast$};
\draw[curve](0,0)--(1,0);\draw[guide](1,0)--(2.7,0);\node[left]at(0,0){$0$};
\draw[curve,domain=1:2.7,samples=80]plot(\x,{acos(1/(\x*\x))*pi/180});
\draw[curve,domain=1:2.7,samples=80]plot(\x,{-acos(1/(\x*\x))*pi/180});
\draw[guide](1,-2)--(1,0);\node[below]at(1,-2){$1$};
\node at(2.15,1.8){稳定分支};
\node at(1.7,-2.75){完整环取有向角};
\end{scope}

\end{tikzpicture}
\caption{左图的 $\vct e_\varphi$ 指向纸面内；右图用有向角表示两个对称的非零平衡，实线稳定、虚线不稳定。}\label{mp3:fig:l3-e9-s}
\end{figure}


相对平衡满足 $\dot\theta=\ddot\theta=0$。最低点 $\theta=0$ 与最高点 $\theta=\pi$ 始终存在；当 $\Omega^2>g/a$ 时，还存在
\[
\theta_0=\arccos\frac{g}{a\Omega^2}.
\]
在完整圆环上它对应竖直轴两侧对称的两个位置。将运动方程右端记作 $f(\theta)$，用乘积法则求导，
\[
f'(\theta)=\cos\theta\left(\Omega^2\cos\theta-\frac ga\right)
-\Omega^2\sin^2\theta.
\]
平衡附近有 $\ddot\eta=f'(\theta_\ast)\eta$。最低点 $f'(0)=\Omega^2-g/a$，所以 $\Omega^2<g/a$ 时稳定，$\Omega^2>g/a$ 时不稳定；最高点 $f'(\pi)=\Omega^2+g/a>0$，始终不稳定。

非零平衡处 $f'(\theta_0)=-\Omega^2\sin^2\theta_0$，故稳定，小振动频率为 $|\Omega|\sin\theta_0$。临界转速 $\Omega^2=g/a$ 时，最低点的线性恢复项为零，但
$\ddot\theta=-\Omega^2\theta^3/2+O(\theta^5)$。要判断这一临界情形的稳定性，将运动方程乘以 $ma^2\dot\theta$，左边为 $\dd(ma^2\dot\theta^2/2)/\dd t$。右边将 $\dot\theta\dd t$ 换成 $\dd\theta$，积分得
\[
\begin{aligned}
\frac{ma^2}{2}[\dot\theta^2(t)-\dot\theta^2(0)]
&=\int_{\theta(0)}^{\theta(t)}
(ma^2\Omega^2\sin u\cos u-mga\sin u)\dd u\\
&=\left[\frac{ma^2\Omega^2}{2}\sin^2u+mga\cos u\right]_{\theta(0)}^{\theta(t)}.
\end{aligned}
\]
移项后，守恒量为 $ma^2\dot\theta^2/2+U_{\rm eff}(\theta)$，其中
\[
U_{\rm eff}=-mga\cos\theta-\frac12ma^2\Omega^2\sin^2\theta.
\]
在$g=a\Omega^2$时，
\[
U_{\rm eff}(\theta)-U_{\rm eff}(0)
=\frac12ma^2\Omega^2(1-\cos\theta)^2
=\frac18ma^2\Omega^2\theta^4+O(\theta^6)>0
\]
对充分小的非零$\theta$成立。因此最低点仍稳定，但恢复力从三次项开始，不能用线性简谐振动给出一个非零的固有频率。

环的两个法向方向可取指向环心的
$\vct n=-\sin\theta\vct e_r+\cos\theta\vct e_z$
以及垂直环面的 $\vct e_\varphi$。约束力写为 $R_n\vct n+R_\varphi\vct e_\varphi$，投影得
\[
R_n=m[a\dot\theta^2+a\Omega^2\sin^2\theta+g\cos\theta],
\quad R_\varphi=2ma\Omega\cos\theta\,\dot\theta.
\]
这里圆珠穿在细环上，可由环的不同侧面受力，$R_n$ 是带符号的分量。非平衡运动中的 $R_\varphi$ 使小珠随环面转向，须与子午截面内的 $R_n$ 一同保留。

\end{solution}

\section{动坐标与相对运动}
随机构运动的坐标系能简化相对轨迹，基矢的转动则增加速度和加速度项。这些项由基矢求导得到，再与地面系的实际运动对应。

\subsection{矢量导数的转换}
设动系原点为 $O(t)$，相对惯性系的角速度为 $\vct\Omega(t)$。动系基矢 $\vct e_i(t)$ 满足 $\dot{\vct e}_i=\vct\Omega\times\vct e_i$。对任意矢量 $\vct A=\sum_i A_i\vct e_i$，使用乘积法则，
\[
\left(\frac{\dd\vct A}{\dd t}\right)_{\!S}
=\sum_i\dot A_i\vct e_i+\sum_iA_i(\vct\Omega\times\vct e_i)
=\sum_i\dot A_i\vct e_i+\vct\Omega\times\vct A.
\]
将只对分量求导的第一项记为动系导数，便得到
\begin{equation}\label{mp3:eq:l3-transport}
\left(\frac{\dd\vct A}{\dd t}\right)_{\!S}
=\left(\frac{\dd\vct A}{\dd t}\right)_{\!S'}
+\vct\Omega\times\vct A.
\end{equation}
右边第一个导数只对动系分量 $A_i$ 求导。两个导数的差来自基矢本身的转动。

质点位置写为 $\vct r_P=\vct r_O+\vct\rho$，在式\eqref{mp3:eq:l3-transport}中先取 $\vct A=\vct\rho$，再按第二讲式\eqref{mp2:eq:l2-transport}前的乘积展开对速度求导，得到
\begin{align}
\vct v_P&=\vct v_O+\vct v_{\rm rel}+\vct\Omega\times\vct\rho,\label{mp3:eq:l3-transport-v}\\
\vct a_P&=\vct a_O+\vct a_{\rm rel}
+2\vct\Omega\times\vct v_{\rm rel}
+\dot{\vct\Omega}\times\vct\rho
+\vct\Omega\times(\vct\Omega\times\vct\rho).\label{mp3:eq:l3-transport-a}
\end{align}
两个科里奥利项分别来自相对速度基矢的转动，以及牵连速度中 $\vct\rho$ 的变化，合为系数 $2$。若 $\vct\Omega$ 自身改变方向，$\dot{\vct\Omega}$ 也须按矢量求导。它在两系中的导数相同，因为二者的差为 $\vct\Omega\times\vct\Omega=0$。

\subsection{在动系中列运动方程}
将式\eqref{mp3:eq:l3-transport-a}代入惯性系中的 $m\vct a_P=\vct F$，并移项，可写为
\begin{align}
m\vct a_{\rm rel}={}&\vct F-m\vct a_O
-2m\vct\Omega\times\vct v_{\rm rel}\nonumber\\
&-m\dot{\vct\Omega}\times\vct\rho
-m\vct\Omega\times(\vct\Omega\times\vct\rho).\label{mp3:eq:l3-moving-newton}
\end{align}
右端依次含平动惯性力、科里奥利力、欧拉惯性力和离心力。各项的方向由相应叉积确定。例如在水平面内，离心力沿径向向外，欧拉惯性力则与转动角加速度的切向作用相反。

科里奥利力与 $\vct v_{\rm rel}$ 垂直，所以对相对运动不做功，但会改变相对轨迹的方向。离心力在固定转轴匀速转动时可写为势能 $-\tfrac12m|\vct\Omega\times\vct\rho|^2$ 的负梯度；若转速随时间改变，欧拉项一般不能由时间无关的势能代替。

图\ref{mp3:fig:l3-moving}统一采用 $\vct\Omega=\Omega\vct k$ 向外、$\Omega>0$。相对速度沿 $+\vct e_r$ 时，运动学中的科里奥利加速度朝 $+\vct e_\varphi$，而移到动系方程右端的科里奥利力朝 $-\vct e_\varphi$。两者的符号用途不同。

\begin{figure}[H]\centering
\begin{tikzpicture}[scale=.85]

\foreach \drawshift in {0,5.5}{
\begin{scope}[shift={(\drawshift,0)}]
\draw[guide](0,0)--(3,0);\fill(0,0)circle(1.4pt)node[below]{$O$};
\fill(2,0)circle(1.5pt)node[below]{$P$};
\draw[axis](2,-.4)--(2.9,-.4)node[right]{$\vct v_{\rm rel}$};
\end{scope}}
\draw[axis,thick](2,0)--(.8,0)node[midway,above]{$-\Omega^2r\,\vct e_r$};
\draw[axis](2,0)--(2,1.4)node[above]{$(2\Omega\dot r+\dot\Omega r)\vct e_\varphi$};
\node at(1.5,-1){加速度中的项};
\draw[axis,thick](7.5,0)--(8.7,0)node[above]{$m\Omega^2r\,\vct e_r$};
\draw[axis](7.5,0)--(7.5,-1.4)node[below]{$-m(2\Omega\dot r+\dot\Omega r)\vct e_\varphi$};
\node at(7,1.35){动系中的惯性力};

\end{tikzpicture}
\caption{左图把同向的两项切向加速度合画，右图是移项后的惯性力。取 \mbox{$\Omega>0$}，\mbox{$\dot\Omega>0$}，\mbox{$\dot r>0$}；下方的相对速度箭头经过平移。}\label{mp3:fig:l3-moving}
\end{figure}

\par\medskip\noindent\begin{minipage}{\textwidth}
\begin{example}[匀速转动直槽中的自由滑珠]\label{mp3:example:l3-e10}

水平面内一条光滑直槽绕端点 $O$ 以恒定角速度 $\Omega>0$ 转动，质量为 $m$ 的小珠可沿槽无摩擦滑动。小珠初始距 $O$ 为 $r_0$，相对槽静止，见图\ref{mp3:fig:l3-e10-q}。求 $r(t)$、槽对珠的横向力，以及维持转速的电机对小珠所做的功。

\end{example}
\begin{figure}[H]\centering
\begin{tikzpicture}[scale=.95]

\draw[axis](-.3,0)--(4.5,0)node[right]{$x$};
\draw[very thick](0,0)--(3.6,1.8);
\draw[guide](0,0)--(3.6,0);
\fill(0,0)circle(1.5pt)node[below]{$O$};
\fill(2.3,1.15)circle(2.5pt)node[above]{$P$};
\draw[|-|](.3,-.6)--(2.6,.55)node[midway,below]{$r$};
\draw(.65,0)arc(0:26.565:.65);\node at(.95,.2){$\Omega t$};
\draw[axis](1.4,.2)arc(8:80:1.42);\node at(1.7,1.8){$\Omega$};
\node at(5,1.1){$\dot r(0)=0$};

\end{tikzpicture}
\caption{俯视水平转槽。重力与竖直支持力平衡，图中只讨论水平运动。}\label{mp3:fig:l3-e10-q}
\end{figure}
\end{minipage}\par\medskip
\begin{solution}
在惯性系中沿槽没有真实力，径向方程为
$m(\ddot r-\Omega^2r)=0$。特征方程$\lambda^2-\Omega^2=0$有两根$\pm\Omega$，故$r=C_1\ee^{\Omega t}+C_2\ee^{-\Omega t}$。初值$r(0)=r_0,\dot r(0)=0$分别给出$C_1+C_2=r_0$、$\Omega(C_1-C_2)=0$，即$C_1=C_2=r_0/2$。因此
\[
r=r_0\cosh\Omega t,\qquad \dot r=r_0\Omega\sinh\Omega t.
\]
图\ref{mp3:fig:l3-e10-s}左图中，相对运动沿槽向外，槽壁需提供沿 $+\vct e_\varphi$ 的力，维持质点方位角速度与槽相同：
\[
N_\varphi=m a_\varphi=2m\Omega\dot r
=2mr_0\Omega^2\sinh\Omega t.
\]

\begin{figure}[H]\centering
\begin{tikzpicture}[scale=.95]

\draw[thick](0,0)--(3.5,0);\fill(0,0)circle(1.4pt)node[below]{$O$};
\fill(2,0)circle(2pt)node[below]{$P$};
\draw[axis](2,0)--(2,1.3)node[above]{$N_\varphi$};
\draw[axis,guide](2,0)--(2,-1.3)node[below]{$-2m\Omega\dot r\,\vct e_\varphi$};
\draw[axis](2,-.35)--(3.3,-.35)node[right]{$\dot r$};
\node at(1.5,2){动系受力};
\begin{scope}[shift={(6.2,0)},scale=.55]
\draw[axis](-.3,0)--(3.6,0)node[right]{$x/r_0$};
\draw[axis](0,-.3)--(0,4.6)node[above]{$y/r_0$};
\draw[curve,domain=0:1.5,samples=100]plot({cosh(\x)*cos(\x*180/pi)},{cosh(\x)*sin(\x*180/pi)});
\fill(1,0)circle(1.5pt)node[below]{$P(0)$};
\node at(3.2,3.6){$r=r_0\cosh\varphi$};
\end{scope}

\end{tikzpicture}
\caption{左图为随槽转动坐标系中的横向力平衡，虚线箭头是惯性力；右图画出 $0\le\Omega t\le1.5$ 的惯性系实际轨迹。}\label{mp3:fig:l3-e10-s}
\end{figure}


在动系中，离心力提供径向相对加速度，科里奥利力朝 $-\vct e_\varphi$，与槽壁力平衡。它们与惯性系中的两条分量方程给出相同结果。

小珠对转轴的力矩为 $-rN_\varphi$，电机须补偿的力矩为 $rN_\varphi$。因此电机传给小珠的功率为
\[
\mathcal P=(rN_\varphi)\Omega
=2m\Omega^2r\dot r.
\]
从 $r_0$ 积分到 $r$，
\[
W=\int_0^t2m\Omega^2r\dot r\dd t
=2m\Omega^2\int_{r_0}^r\xi\dd\xi
=2m\Omega^2\left[\frac{\xi^2}{2}\right]_{r_0}^r
=m\Omega^2(r^2-r_0^2).
\]
惯性系中的动能是 $K=\tfrac12m(\dot r^2+\Omega^2r^2)$。由解得
$\dot r^2=\Omega^2(r^2-r_0^2)$，故
$K-K(0)=m\Omega^2(r^2-r_0^2)=W$。

槽壁的接触力沿 $\vct e_\varphi$，小珠在惯性系中有同向速度分量 $\Omega r$，所以接触力的功率为 $N_\varphi\Omega r$。电机的功同时增加径向运动的动能与随槽转动的动能；沿槽光滑只保证相对滑动方向没有摩擦。

\end{solution}

\par\medskip\noindent\begin{minipage}{\textwidth}
\begin{example}[旋转圆盘上释放的自由质点]\label{mp3:example:l3-e11}

光滑水平圆盘以角速度 $\Omega>0$ 逆时针转动。质点最初固定在距轴 $R$ 处，$t=0$ 时解除固定，随后只受竖直力，在圆盘上自由滑动。初始动、静坐标轴重合，质点位于 $(R,0)$。求惯性系与圆盘系中的轨迹，见图\ref{mp3:fig:l3-e11-q}。若圆盘半径为 $L>R$，求质点首次到达盘缘的时刻。

\end{example}
\begin{figure}[H]\centering
\begin{tikzpicture}[scale=.95]

\draw[curve](0,0)circle(1.9);\draw[axis](-2.3,0)--(2.6,0)node[right]{$x,x'$};
\draw[axis](0,-2.3)--(0,2.6)node[above]{$y,y'$};
\fill(0,0)circle(1.4pt)node[below left]{$O$};
\fill(1.15,0)circle(2pt)node[below]{$P(0)$};
\draw[axis](1.15,0)--(1.15,1.2)node[above]{$\Omega R$};
\draw[|-|](0,-.45)--(1.15,-.45)node[midway,fill=white]{$R$};
\draw[axis](0,1.3)arc(90:165:1.3);\node at(-1.6,1.05){$\Omega$};

\end{tikzpicture}
\caption{释放时质点保留切向速度 $\Omega R$，其相对圆盘的初速度为零。}\label{mp3:fig:l3-e11-q}
\end{figure}
\end{minipage}\par\medskip
\begin{solution}
释放后水平合力为零，初速度仍为释放前的切向速度 $\vct v(0)=(0,\Omega R)$。对零加速度积分，$\vct v(t)=\vct v(0)+\int_0^t\vct0\dd t=\vct v(0)$。再对常速度积分，
\[
\vct r(t)=(R,0)+\int_0^t(0,\Omega R)\dd t
=(R,0)+(0,\Omega Rt).
\]
因此惯性系中
\[
x=R,\qquad y=\Omega Rt.
\]
令 $\tau=\Omega t$。圆盘系坐标由逆旋转得到，
\[
\begin{pmatrix}x'\\y'\end{pmatrix}
=\begin{pmatrix}\cos\tau&\sin\tau\\-\sin\tau&\cos\tau\end{pmatrix}
\begin{pmatrix}R\\R\tau\end{pmatrix},
\]
\[
x'=R(\cos\tau+\tau\sin\tau),\qquad
y'=R(-\sin\tau+\tau\cos\tau).
\]
图\ref{mp3:fig:l3-e11-s}分别用各自的固定坐标画轨迹。两幅图中的坐标基准不同，不能把左图的切线当作右图的初始相对速度。

\begin{figure}[H]\centering
\begin{tikzpicture}[scale=.95]

\draw[axis](-.2,0)--(2.3,0)node[right]{$x/R$};
\draw[axis](0,-.2)--(0,2.8)node[above]{$y/R$};
\draw[curve](1,0)--(1,2.4);\fill(1,0)circle(1.5pt)node[below]{$P_0$};
\draw[axis](1,.9)--(1,1.6); \node[right]at(1,2.1){惯性系};
\begin{scope}[shift={(5.2,1.5)}]
\draw[axis](-.2,0)--(3.1,0)node[right]{$x'/R$};
\draw[axis](0,-2.5)--(0,1)node[above]{$y'/R$};
\draw[curve,domain=0:1.8,samples=100]plot({cos(\x*180/pi)+\x*sin(\x*180/pi)},{-sin(\x*180/pi)+\x*cos(\x*180/pi)});
\fill(1,0)circle(1.5pt)node[above]{$P_0$};
\fill({cos(1.8*180/pi)+1.8*sin(1.8*180/pi)},{-sin(1.8*180/pi)+1.8*cos(1.8*180/pi)})circle(1.5pt)node[below]{后续位置};
\node at(1.7,.7){圆盘系};
\end{scope}

\end{tikzpicture}
\caption{同一次自由运动在两系中的表示。右图仅取 $0\le\tau\le1.8$，从 $P_0$ 开始向外并向顺时针一侧弯曲。}\label{mp3:fig:l3-e11-s}
\end{figure}


对 $\tau\ll1$ 展开，
\[
x'=R\left(1+\frac{\tau^2}{2}+O(\tau^4)\right),\qquad
y'=-\frac R3\tau^3+O(\tau^5).
\]
因此相对速度初始为零，相对加速度初始沿径向向外；顺时针弯曲在随后出现。动系方程为
\[
\ddot x'=2\Omega\dot y'+\Omega^2x',\qquad
\ddot y'=-2\Omega\dot x'+\Omega^2y',
\]
初值代入也得到 $\ddot x'(0)=\Omega^2R,\ddot y'(0)=0$。科里奥利力在释放瞬间为零；随着向外的相对速度建立，它才逐渐使轨迹向顺时针一侧弯曲。

到轴的距离不受坐标旋转影响，
\[
r=\sqrt{x^2+y^2}=R\sqrt{1+\Omega^2t^2}.
\]
所以首次到达盘缘的时间为
\[
t_{\rm edge}=\frac1\Omega\sqrt{\frac{L^2}{R^2}-1}.
\]
以上解仅描述到达盘缘前质点仍由盘面竖直支持的阶段。

\end{solution}

\par\medskip\noindent\begin{minipage}{\textwidth}
\begin{example}[角加速圆盘上的摩擦极限]\label{mp3:example:l3-e12}

水平圆盘从静止开始以恒定角加速度 $\alpha>0$ 转动。质量为 $m$ 的小块放在距轴 $r>0$ 处，静摩擦因数为 $\mu_s$，动摩擦因数为 $\mu_k<\mu_s$。见图\ref{mp3:fig:l3-e12-q}。求小块随盘静止的条件、首次滑动时刻，并确定刚开始滑动时的相对加速度方向。

\end{example}
\begin{figure}[H]\centering
\begin{tikzpicture}[scale=.95]

\draw[curve](0,0)circle(1.9);\draw[guide](0,0)--(1.3,0);
\fill(0,0)circle(1.5pt)node[below]{$O$};
\draw[fill=white](1.1,-.18)rectangle(1.5,.18);\node[below]at(1.3,-.25){$P$};
\draw[|-|](0,-.65)--(1.3,-.65)node[midway,fill=white]{$r$};
\draw[axis](0,1.3)arc(90:175:1.3);\node at(-2.15,1.3){$\alpha$};
\node at(3.6,1){$\Omega(t)=\alpha t$};
\node at(3.6,0){$\mu_k<\mu_s$};

\end{tikzpicture}
\caption{俯视圆盘。开始转动后，静摩擦同时提供切向与径向加速度。}\label{mp3:fig:l3-e12-q}
\end{figure}
\end{minipage}\par\medskip
\begin{solution}
若小块相对圆盘静止，惯性系中的加速度为
\[
\vct a=-r\Omega^2\vct e_r+r\alpha\vct e_\varphi.
\]
竖直支持力为 $mg$，水平静摩擦力须等于 $m\vct a$。图\ref{mp3:fig:l3-e12-s}左图将两项按正交方向合成，所以不滑动的充要条件为
\[
mr\sqrt{\Omega^4+\alpha^2}\le\mu_smg.
\]
若 $r\alpha>\mu_sg$，在启动时静摩擦就不足，小块立即滑动。若 $r\alpha=\mu_sg$，启动时恰达极限，任意 $t>0$ 所需摩擦都会超过极限，首次滑动仍取 $t_\ast=0$。

\begin{figure}[H]\centering
\begin{tikzpicture}[scale=.95]

\fill(0,0)circle(1.5pt)node[below right]{$P$};
\draw[axis](0,0)--(-1.5,0)node[below]{$mr\Omega^2$};
\draw[axis](0,0)--(0,1)node[right]{$mr\alpha$};
\draw[axis,thick](0,0)--(-1.5,1)node[above left]{$\vct f_s$};
\draw[guide](-1.5,0)--(-1.5,1)--(0,1);
\begin{scope}[shift={(5.4,0)}]
\fill(0,0)circle(1.5pt)node[above left]{$P$};
\draw[axis](0,0)--(1.5,0)node[above]{$m r\Omega_\ast^2$};
\draw[axis](0,0)--(0,-1)node[left]{$m r\alpha$};
\draw[axis,thick](0,0)--(1.5,-1)node[right]{$\vct d$};
\draw[axis](0,0)--(-.85,.567)node[left]{$\vct f_k$};
\draw[guide](1.5,0)--(1.5,-1)--(0,-1);
\end{scope}

\end{tikzpicture}
\caption{左图为随盘运动所需的真实静摩擦力；右图为起滑瞬间两项惯性力的合方向及反向动摩擦。箭头长度只示意分量关系。}\label{mp3:fig:l3-e12-s}
\end{figure}


若 $r\alpha<\mu_sg$，由 $\Omega=\alpha t$，
\[
t_\ast=\frac1\alpha
\left[\left(\frac{\mu_sg}{r}\right)^2-\alpha^2\right]^{1/4}.
\]
设刚开始滑动时的转速为 $\Omega_\ast$。相对速度尚为零，科里奥利项为零；离心与欧拉惯性力的合方向为
\[
\vct d=\frac{\Omega_\ast^2\vct e_r-\alpha\vct e_\varphi}
{\sqrt{\Omega_\ast^4+\alpha^2}}.
\]
如右图，小块将相对圆盘向外并向落后于转动的一侧滑动，动摩擦与 $\vct d$ 反向。因此起滑后的右极限为
\[
\vct a_{\rm rel}(t_\ast^+)
=\left[r\sqrt{\Omega_\ast^4+\alpha^2}-\mu_kg\right]\vct d.
\]
对 $t_\ast>0$，括号等于 $(\mu_s-\mu_k)g>0$。对立即滑动情形，取 $\Omega_\ast=0$，相对加速度沿 $-\vct e_\varphi$，其大小为 $r\alpha-\mu_kg$。在启动瞬间 $\Omega=0$，静摩擦的需求已有 $mr\alpha$；之后径向分量增大，合摩擦力沿上述矢量关系变化。

\end{solution}


\subsection{转动参考系中的机构与接触运动}
\subsubsection{从一般变换到逆变换}
在式\eqref{mp3:eq:l3-transport-v}、\eqref{mp3:eq:l3-transport-a}中，将动系原点的速度、加速度记为 $\vct V_O,\vct A_O$，角加速度记为 $\vct\alpha=\dot{\vct\Omega}$，相对位置改记为 $\vct r$。这与前面的 $\vct\rho$ 表示同一位矢。
在三维运动中，由矢量三重积可知，最后一项等于$\vct\Omega(\vct\Omega\cdot\vct r)-\Omega^2\vct r$；只有$\vct r\perp\vct\Omega$时才能简写成$-\Omega^2\vct r$。

若已知地面运动，应先求相对速度，再求相对加速度：
\[
 \vct v_{\rm rel}=\vct v-\vct V_O-\vct\Omega\times\vct r,
\]
把速度逆变换代入原加速度关系，先保留两项二次叉积：
\[
\begin{aligned}
\vct a_{\rm rel}={}&\vct a-\vct A_O-\vct\alpha\times\vct r
-2\vct\Omega\times(\vct v-\vct V_O-\vct\Omega\times\vct r)\\
&-\vct\Omega\times(\vct\Omega\times\vct r)\\
={}&\vct a-\vct A_O-\vct\alpha\times\vct r
-2\vct\Omega\times(\vct v-\vct V_O)\\
&+\vct\Omega\times(\vct\Omega\times\vct r).
\end{aligned}
\]
逆变换中的科里奥利项所用的速度也已经改变，不能只把原式中的角速度换号。图\ref{mp3:fig:hw:transport}标出相对位置、速度以及转动轴。

\begin{figure}[H]\centering
\begin{tikzpicture}[>=Stealth,line cap=round,line join=round]

\draw[axis](0,0)--(3.4,0)node[right]{$x$};\draw[axis](0,0)--(0,2.7)node[above]{$z$};\draw[axis](0,0)--(-1.2,-.65)node[left]{$y$};\fill(0,0)circle(1.3pt)node[below right]{$O$};
\draw[axis,thick](0,0)--(2,1.1)node[midway,below]{$\vct r$};\fill(2,1.1)circle(1.3pt)node[below right]{$P$};\draw[axis](2,1.1)--(2.8,2.2)node[above]{$\vct v_{\rm rel}$};\draw[axis](0,0)--(.5,2.2)node[left]{$\vct\Omega$};\node[align=left,right]at(4,.9){$\vct v_{\rm rel}$在动系中求导\\$\vct\Omega$的方向可以随时间改变};

\end{tikzpicture}
\caption{三维转动参考系中的位置、相对速度与转动轴}\label{mp3:fig:hw:transport}
\end{figure}

\par\medskip\noindent\begin{minipage}{\textwidth}
\begin{example}[两根转杆的绝对转动与相对转动]\label{mp3:hw:l3-17}
图\ref{mp3:fig:hw:l3-17-q}中，$AB=2l$、$BC=l$，$A$固定。当前$AB$水平向右，$BC$竖直向上。两杆的绝对角速度、角加速度分别为$\omega_1,\beta_1$和$\omega_2,\beta_2$，均以逆时针为正。求$C$的地面速度、加速度，并在随$AB$转动的参考系中重算。
\end{example}

\begin{figure}[H]\centering
\begin{tikzpicture}[>=Stealth,line cap=round,line join=round]

\draw[axis](-.3,0)--(4,0)node[right]{$x$};\draw[axis](0,-.3)--(0,2.5)node[above]{$y$};\draw[thick](0,0)--(3,0)--(3,1.5);\fill(0,0)circle(1.5pt)node[below left]{$A$};\fill(3,0)circle(1.5pt)node[below]{$B$};\fill(3,1.5)circle(1.5pt)node[above]{$C$};\node[below]at(1.5,0){$2l$};\node[right]at(3,.8){$l$};\node at(1.2,1){$\omega_1,\beta_1$};\node[right]at(3.5,1.8){$\omega_2,\beta_2$};

\end{tikzpicture}
\caption{当前构形中的两根连杆；角速度均指相对地面的转动}\label{mp3:fig:hw:l3-17-q}
\end{figure}

\end{minipage}\par\medskip
\begin{solution}
先把$BC$相对于$B$的运动与$B$的平动相加。由$\vct r_B=(2l,0)$，
\[
 \vct v_B=(0,2l\omega_1),\qquad
 \vct a_B=(-2l\omega_1^2,2l\beta_1).
\]
再用$\vct r_{C/B}=(0,l)$，得到
\[
 \boxed{\vct v_C=(-l\omega_2,2l\omega_1)},
 \qquad
 \boxed{\vct a_C=(-2l\omega_1^2-l\beta_2,
                         2l\beta_1-l\omega_2^2)}.
\]
本题给出的$BC$角速度相对地面定义；$AB$的转动已经通过铰点的平动计入。

改在随$AB$转动、原点仍为$A$的动系中计算。$B$固定，$BC$相对动系的角速度与角加速度为$\omega_2-\omega_1$、$\beta_2-\beta_1$，因而
\[
 \vct v'_C=-l(\omega_2-\omega_1)\vct e_x,
\]
\[
 \vct a'_C=-l(\beta_2-\beta_1)\vct e_x
           -l(\omega_2-\omega_1)^2\vct e_y.
\]
图\ref{mp3:fig:hw:l3-17-s}右图将相对加速度与三项转动贡献分开。代入运输公式时，位置必须从动系原点$A$量起，即$\vct r_C=(2l,l)$。各项为
\[
 \begin{aligned}
 \vct\beta_1\times\vct r_C&=(-l\beta_1,2l\beta_1),\\
 -\omega_1^2\vct r_C&=(-2l\omega_1^2,-l\omega_1^2),\\
 2\vct\omega_1\times\vct v'_C
  &=(0,-2l\omega_1(\omega_2-\omega_1)).
 \end{aligned}
\]
纵向三个平方项合并成$-l\omega_2^2$，横向两个角加速度项合并成$-l\beta_2$，与地面系结果相同。若某杆实际顺时针运动，只须给它的有符号角速度取负值。


\begin{figure}[H]\centering
\begin{tikzpicture}[>=Stealth,line cap=round,line join=round]

\fill(0,0)circle(1.5pt)node[below left]{$B$};\draw[axis](0,0)--(0,1.7)node[above]{$\vct v_B$};\draw[axis](0,1.7)--(-1.6,1.7)node[left]{$\vct v_{C/B}$};\draw[axis,thick](0,0)--(-1.6,1.7)node[midway,left]{$\vct v_C$};
\node[right,align=left]at(1.4,1.2){随$AB$的动系：\\[5pt]相对项$\vct a'_C$\\欧拉项$\vct\beta_1\times\vct r_C$\\向心项$-\omega_1^2\vct r_C$\\科里奥利项$2\vct\omega_1\times\vct v'_C$};

\end{tikzpicture}
\caption{左图按正角速度示意速度合成；右侧列出动系计算所需的四项}\label{mp3:fig:hw:l3-17-s}
\end{figure}

\end{solution}

\par\medskip\noindent\begin{minipage}{\textwidth}
\begin{example}[绕圆周上一点转动的水平圆环]\label{mp3:hw:l3-18}
半径$R$的光滑圆环保持水平，绕圆周上的固定点$O$以恒定角速度$\omega\vct e_z$转动。质量$m$的小环套在大环上，当前相对大环的有符号速率为$v$。在随大环转动的平面坐标中，圆心为$C=(R,0)$，小环位置为$P=C+R(\cos\theta,\sin\theta)$，$v>0$指向$\theta$增大。求相对切向加速度及大环对小环的支持力，重力竖直向下。见图\ref{mp3:fig:hw:l3-18-q}。
\end{example}

\begin{figure}[H]\centering
\begin{tikzpicture}[>=Stealth,line cap=round,line join=round]

\draw[curve](1.6,0)circle(1.6);\draw[axis](-.4,0)--(3.9,0)node[right]{$x'$};\fill(0,0)circle(1.4pt)node[below left]{$O$};\fill(1.6,0)circle(1.4pt)node[below]{$C$};\coordinate(P)at(2.731,1.131);\draw[guide](1.6,0)--(P)node[midway,above left]{$R$};\fill(P)circle(1.6pt)node[above right]{$P$};\draw[axis](P)--(2.024,1.838)node[above]{$v$};\draw(2.15,0)arc(0:45:.55);\node at(2.32,.27){$\theta$};\node[left]at(-.1,1){$\omega\vct e_z$};

\end{tikzpicture}
\caption{转动圆环的平面坐标与小环的相对运动方向}\label{mp3:fig:hw:l3-18-q}
\end{figure}

\end{minipage}\par\medskip
\begin{solution}
取向外径向单位矢量$\vct n=(\cos\theta,\sin\theta)$、切向$\vct t=(-\sin\theta,\cos\theta)$。在大环系中，
\[
 \vct v'=v\vct t,\qquad
 \vct a'=\dot v\vct t-\frac{v^2}{R}\vct n.
\]
运输项的位置须从转轴量起，因此$\vct r=\overrightarrow{OP}=R(1+\cos\theta,\sin\theta)$。$\overrightarrow{CP}$只用于小环相对大环的圆周运动。其投影为
\[
 \vct r\cdot\vct t=-R\sin\theta,\qquad
 \vct r\cdot\vct n=R(1+\cos\theta).
\]
科里奥利项$2\omega\vct e_z\times v\vct t=-2\omega v\vct n$沿径向。光滑接触没有平面切向力，因此
\[
 m(\dot v+\omega^2R\sin\theta)=0,
 \qquad \boxed{\dot v=-\omega^2R\sin\theta}.
\]
图\ref{mp3:fig:hw:l3-18-s}中，平面法向的三项加速度相加后，大环作用力的水平分量为
\[
 \boxed{\vct N_h=-m\left[\frac{v^2}{R}+2\omega v
                         +\omega^2R(1+\cos\theta)\right]\vct n}.
\]
小环还受到向上的支持分量$mg\vct e_z$，从而总支持力为$\vct N=\vct N_h+mg\vct e_z$，大小为
\[
 N=m\sqrt{g^2+\left[v^2/R+2\omega v+\omega^2R(1+\cos\theta)\right]^2}.
\]
小环套在环上，水平法向力可向内也可向外；若只靠单侧接触，须另外检验允许的力向。令$\omega=0$，结果回到固定水平圆环的向心力及竖直支持。


\begin{figure}[H]\centering
\begin{tikzpicture}[>=Stealth,line cap=round,line join=round]

\fill(0,0)circle(1.5pt)node[below right]{$P$};\draw[axis](0,0)--(-2,0)node[left]{$\vct N_h$};\draw[axis](0,0)--(0,1.6)node[above]{$mg\vct e_z$};\draw[axis,thick](0,0)--(-2,1.6)node[above left]{$\vct N$};\node[right,align=left]at(1.1,.7){水平径向加速度：\\$-v^2/R$\\$-2\omega v$\\$-\omega^2R(1+\cos\theta)$};

\end{tikzpicture}
\caption{水平法向加速度的三项贡献，以及支持力的竖直分量}\label{mp3:fig:hw:l3-18-s}
\end{figure}

\end{solution}

\par\medskip\noindent\begin{minipage}{\textwidth}
\begin{example}[两根不等角速度转杆的交点]\label{mp3:hw:l3-19}
两根无限长直杆分别绕$A=(0,0)$、$B=(l,0)$转动，当前几何交点$C=(0,l)$，见图\ref{mp3:fig:hw:l3-19-q}。两杆均顺时针转动，角速度大小分别为$\omega$、$\sqrt2\omega$，顺时针角加速度分别为$\beta$、$\sqrt2\beta$。求$C$的速度与加速度；$\beta$允许为负。
\end{example}

\begin{figure}[H]\centering
\begin{tikzpicture}[>=Stealth,line cap=round,line join=round]

\draw[axis](-.4,0)--(3.3,0)node[right]{$x$};\draw[axis](0,-.3)--(0,3)node[above]{$y$};\draw[thick](0,-.15)--(0,2.8);\draw[thick](-.5,2.5)--(2.5,-.5);\fill(0,0)circle(1.4pt)node[below left]{$A$};\fill(2,0)circle(1.4pt)node[below]{$B$};\fill(0,2)circle(1.4pt)node[above right]{$C$};\node[below]at(1,0){$l$};\node[left]at(0,1){$l$};\node[right]at(.8,2.5){$\omega,\beta$};\node[right]at(2.2,1){$\sqrt2\omega,\sqrt2\beta$};\draw(1.5,0)arc(180:135:.5);\node at(1.27,.27){$45^\circ$};

\end{tikzpicture}
\caption{交点同时位于两根转杆上，两杆转向均为顺时针}\label{mp3:fig:hw:l3-19-q}
\end{figure}

\end{minipage}\par\medskip
\begin{solution}
对绕固定点转动的直线，设其单位方向为 $\vct e$、左法线为 $\vct n$，交点到铰点的有向距离为 $s$。采用第二讲例\ref{mp2:rev:intersection}从法向位置约束逐步求导所得的关系，
\[
 \vct n\cdot\vct v=\Omega s,\qquad
 \vct n\cdot\vct a=2\Omega\vct e\cdot\vct v+\dot\Omega s.
\]
这里$\Omega$以逆时针为正，所以本题两杆分别取$-\omega$、$-\sqrt2\omega$。由杆$AC$的法向速度关系，$v_x=l\omega$；由杆$BC$的关系，$v_x+v_y=2\sqrt2l\omega$。于是
\[
 \boxed{\vct v_C=l\omega(1,2\sqrt2-1)}.
\]
再对$AC$使用加速度关系，得
\[
 a_x=2\omega v_y+l\beta.
\]
对$BC$，取$\vct e=(-1,1)/\sqrt2$、$s=\sqrt2l$，得到
\[
 a_x+a_y=2\sqrt2\omega(-v_x+v_y)+2\sqrt2l\beta.
\]
消去速度分量，
\[
 \boxed{\begin{aligned}
 a_x&=l\big[(4\sqrt2-2)\omega^2+\beta\big],\\
 a_y&=l\big[(10-8\sqrt2)\omega^2+(2\sqrt2-1)\beta\big].
 \end{aligned}}
\]
图\ref{mp3:fig:hw:l3-19-s}中的两条法线分别给出一个标量方程。若把交点当成杆上的材料点，便会漏掉加速度关系中的滑动项$2\Omega\vct e\cdot\vct v$。也可让两杆方向角分别取$\pi/2-\omega t-\beta t^2/2$、$3\pi/4-\sqrt2\omega t-\sqrt2\beta t^2/2$，先求直线交点，再在$t=0$求导，独立得到上述结果。


\begin{figure}[H]\centering
\begin{tikzpicture}[>=Stealth,line cap=round,line join=round]

\fill(0,0)circle(1.5pt)node[below left]{$C$};\draw[axis](0,0)--(1.1,2.01)node[above]{$\vct v_C$};\draw[guide](0,0)--(1.1,0)--(1.1,2.01);\node[below]at(.55,0){$l\omega$};\node[right]at(1.1,1){$(2\sqrt2-1)l\omega$};\begin{scope}[shift={(4.8,0)}]\draw[thick](0,-.6)--(0,2);\draw[thick](-1.2,1.2)--(1.2,-1.2);\fill(0,0)circle(1.5pt);\draw[axis](0,0)--(-1.5,0)node[left]{$\vct n_A$};\draw[axis](0,0)--(-1,-1)node[below]{$\vct n_B$};\end{scope}

\end{tikzpicture}
\caption{速度分量与用于求解加速度的两条杆法线}\label{mp3:fig:hw:l3-19-s}
\end{figure}

\end{solution}

\par\medskip\noindent\begin{minipage}{\textwidth}
\begin{example}[绕焦点转动的椭圆与竖直顶杆]\label{mp3:hw:l3-20}
椭圆的长半轴$a=2l$、半焦距$c=l$，绕固定焦点$F$转动。$A$为半短轴端点，当前位于$F$正上方；椭圆中心在$FA$右侧，见图\ref{mp3:fig:hw:l3-20-q}。竖直导向的杆端始终位于过$F$的竖直线上，并与椭圆光滑接触。椭圆的有符号角速度、角加速度为$\omega,\beta$，以逆时针为正，求杆端速度与加速度。
\end{example}

\begin{figure}[H]\centering
\begin{tikzpicture}[>=Stealth,line cap=round,line join=round]

\coordinate(F)at(0,0);\coordinate(C)at(.866,.5);\begin{scope}[shift={(C)},rotate=30]\draw[curve](0,0)ellipse[x radius=2,y radius=1.732];\draw[guide](-2,0)--(2,0);\end{scope}\draw[thick](0,2)--(0,3.2);\draw[thick](-.35,2.7)--(-.12,2.7);\draw[thick](.12,2.7)--(.35,2.7);\fill(F)circle(1.5pt)node[below left]{$F$};\fill(C)circle(1.4pt)node[right]{$C$};\fill(0,2)circle(1.6pt)node[above left]{$A$};\draw[guide](F)--(0,2)node[midway,left]{$2l$};\node[right]at(3.4,1.7){$a=2l,\ b=\sqrt3l$};\node[right]at(3.4,.8){$\omega,\beta$：逆时针为正};

\end{tikzpicture}
\caption{半短轴端点、固定焦点与竖直导杆的瞬时位置}\label{mp3:fig:hw:l3-20-q}
\end{figure}

\end{minipage}\par\medskip
\begin{solution}
由$b^2=a^2-c^2$，得$b=\sqrt3l$。图\ref{mp3:fig:hw:l3-20-s}中，椭圆长轴方向与水平线夹角为$q_0=\pi/6$。令长轴方向为$\vct u=(\cos q,\sin q)$，短轴方向为$\vct w=(-\sin q,\cos q)$，则椭圆中心为$l\vct u$。杆端位置$(0,y)$满足
\[
 \frac{(y\sin q-l)^2}{4l^2}
 +\frac{y^2\cos^2q}{3l^2}=1.
\]
该约束给出竖直导向线与椭圆的交点。乘以$12l^2$并整理，可因式分解为
\[
(4-\sin^2q)y^2-6ly\sin q-9l^2
=[(2-\sin q)y-3l][(2+\sin q)y+3l]=0.
\]
杆端接触上方分支，故$y=3l/(2-\sin q)>0$。求导始终跟踪这一几何位置，接触的椭圆材料点则随时间更换。连续求导得
\[
y_q=\frac{3l\cos q}{(2-\sin q)^2},\qquad
y_{qq}=3l\left[-\frac{\sin q}{(2-\sin q)^2}
+\frac{2\cos^2q}{(2-\sin q)^3}\right].
\]
在当前$q=\pi/6$处，$y=2l$，且
\[
 \frac{\dd y}{\dd q}=\frac{2l}{\sqrt3},\qquad
 \frac{\dd^2y}{\dd q^2}=\frac{2l}{3}.
\]
由$\dot q=\omega$、$\ddot q=\beta$，有$\dot y=y_q\omega$、$\ddot y=y_q\beta+y_{qq}\omega^2$，因而
\[
 \boxed{\vct v=\frac{2l\omega}{\sqrt3}\vct e_y},\qquad
 \boxed{\vct a=\left(\frac{2l\beta}{\sqrt3}
                         +\frac{2l\omega^2}{3}\right)\vct e_y}.
\]
也可用图\ref{mp3:fig:hw:l3-20-s}中的接触相对运动重算。当前椭圆材料点速度为$-2l\omega\vct e_x$，所以杆端相对椭圆的速度为$\vct v'=(2l\omega,2l\omega/\sqrt3)$，沿长轴方向的切线$\vct u$。$A$处曲率半径为$\rho=a^2/b=4l/\sqrt3$，相对法向指向$-\vct w$。记相对切向加速度为$a'_t$，则
\[
\vct a'=a'_t\vct u-\frac{|\vct v'|^2}{\rho}\vct w
=a'_t\left(\frac{\sqrt3}{2},\frac12\right)
+\left(\frac{2l\omega^2}{\sqrt3},-2l\omega^2\right).
\]
原点$F$固定；欧拉项、向心项和科里奥利项依次为
\[
\beta\vct k\times\vct r=(-2l\beta,0),\qquad
-\omega^2\vct r=(0,-2l\omega^2),\qquad
2\omega\vct k\times\vct v'
=\left(-\frac{4l\omega^2}{\sqrt3},4l\omega^2\right).
\]
杆端的水平加速度为零，因此
$\sqrt3a'_t/2-2l\beta-2l\omega^2/\sqrt3=0$，解得
$a'_t=4l\beta/\sqrt3+4l\omega^2/3$。竖直方向三个法向、向心与科里奥利分量相消，剩下$a_y=a'_t/2$，与坐标求导一致。图中的速度三角形同时说明：接触处法向速度相同，材料点速度与杆端速度并不相同。


\begin{figure}[H]\centering
\begin{tikzpicture}[>=Stealth,line cap=round,line join=round]

\coordinate(A)at(0,0);\draw[guide](-1.2,-.693)--(2.8,1.617)node[right]{切线};\fill(A)circle(1.5pt)node[below right]{$A$};\draw[axis,thick](A)--(0,1.5)node[above right]{杆端速度};\draw[axis](A)--(-2.6,0)node[left]{材料点速度};\draw[axis](-2.6,0)--(0,1.5)node[midway,above left]{相对速度};\draw[guide](0,0)--(-.8,1.386)node[left]{$\vct w$};\node at(1.8,-.8){$q_0=\pi/6$};

\end{tikzpicture}
\caption{正角速度下的速度合成；相对速度沿椭圆在半短轴端点的切线}\label{mp3:fig:hw:l3-20-s}
\end{figure}

\end{solution}

\section{空间滚动与定点转动}
空间滚动同时受接触条件和刚体姿态限制。绝对角速度合成后，由接触点瞬时静止求平动；材料点加速度还包含角速度方向变化的贡献。

\par\medskip\noindent\begin{minipage}{\textwidth}
\begin{example}[沿圆周转弯的直立圆盘]\label{mp3:hw:l3-21}
半径$r$的直立圆盘在水平面上纯滚动。圆心$B$由水平杆$O'B=R$带动，绕竖直轴以恒定角速度$\omega>0$转动。图\ref{mp3:fig:hw:l3-21-q}所示瞬间，$O=(0,0,0)$、$O'=(0,0,r)$、$B=(0,R,r)$，圆盘平面为$xz$平面。求圆盘的绝对角速度，以及盘上材料点$C=(-r,R,r)$的速度、加速度。
\end{example}

\begin{figure}[H]\centering
\begin{tikzpicture}[>=Stealth,line cap=round,line join=round]

\draw[axis](0,0)--(4,0)node[right]{$y$};\draw[axis](0,0)--(-1.7,-1)node[left]{$x$};\draw[axis](0,0)--(0,3.4)node[above]{$z$};\coordinate(B)at(2.4,1.3);\draw[thick](0,1.3)--(B)node[midway,above]{$R$};\draw[curve,domain=0:360,samples=100]plot({2.4-.7*cos(\x)},{1.3-.4*cos(\x)+1.3*sin(\x)});\draw[guide](B)--(2.4,0)node[midway,right]{$r$};\fill(0,1.3)circle(1.4pt)node[left]{$O'$};\fill(B)circle(1.4pt)node[above right]{$B$};\fill(3.1,1.7)circle(1.4pt)node[right]{$C$};\fill(2.4,0)circle(1.4pt)node[below]{$D$};\node[below left]at(0,0){$O$};\node[right]at(4,2.8){$C=(-r,R,r)$};

\end{tikzpicture}
\caption{圆盘的自转轴随水平连杆转动；$D$为当时接触地面的材料点}\label{mp3:fig:hw:l3-21-q}
\end{figure}

\end{minipage}\par\medskip
\begin{solution}
圆心作水平圆周运动，当前
\[
 \vct v_B=-R\omega\vct e_x,\qquad
 \vct a_B=-R\omega^2\vct e_y.
\]
设绕连杆方向的自转角速度为$\nu\vct e_y$。接触点$D$相对$B$的位置为$-r\vct e_z$，纯滚动要求
\[
 0=\vct v_B+(\omega\vct e_z+\nu\vct e_y)\times(-r\vct e_z),
\]
所以$\nu=-R\omega/r$，绝对角速度为
\[
 \boxed{\vct W=\omega\vct e_z-\frac{R}{r}\omega\vct e_y}.
\]
图\ref{mp3:fig:hw:l3-21-s}左图把自转与转弯的角速度作矢量合成。两个分量的大小恒定，但$\vct e_y$绕竖直轴转动，$\dot{\vct e}_y=-\omega\vct e_x$，故
\[
 \boxed{\vct\alpha=\frac{R}{r}\omega^2\vct e_x}.
\]
对材料点$C$，$\vct r_{C/B}=-r\vct e_x$，先计算相对圆心的两次叉积：
\[
\vct W\times\vct r_{C/B}=-r\omega\vct e_y-R\omega\vct e_z,
\qquad
\vct W\times(\vct W\times\vct r_{C/B})
=\left(r+\frac{R^2}{r}\right)\omega^2\vct e_x.
\]
欧拉项$\vct\alpha\times\vct r_{C/B}=0$。再分别加上圆心的速度和加速度，得到
\[
 \boxed{\vct v_C=-R\omega\vct e_x-r\omega\vct e_y
                         -R\omega\vct e_z},
\]
\[
 \boxed{\vct a_C=\left(r+\frac{R^2}{r}\right)\omega^2\vct e_x
                         -R\omega^2\vct e_y}.
\]
对盘上其他点，欧拉项一般仍有贡献。本点的竖直加速度也恰为零。


\begin{figure}[H]\centering
\begin{tikzpicture}[>=Stealth,line cap=round,line join=round]

\draw[axis](0,0)--(0,1.6)node[above]{$\omega\vct e_z$};\draw[axis](0,0)--(-2,0)node[left]{$-(R/r)\omega\vct e_y$};\draw[axis,thick](0,0)--(-2,1.6)node[above left]{$\vct W$};\fill(0,0)circle(1.4pt);\begin{scope}[shift={(4.2,.2)}]\fill(0,0)circle(1.5pt)node[below right]{$C$};\draw[axis](0,0)--(1.3,.65)node[above]{$-R\omega\vct e_x$};\draw[axis](0,0)--(-1.3,0)node[left]{$-r\omega\vct e_y$};\draw[axis](0,0)--(0,-1.5)node[below]{$-R\omega\vct e_z$};\end{scope}

\end{tikzpicture}
\caption{角速度合成与材料点$C$的三个速度分量；投影方向仅作空间示意}\label{mp3:fig:hw:l3-21-s}
\end{figure}

\end{solution}

\par\medskip\noindent\begin{minipage}{\textwidth}
\begin{example}[圆锥在水平面上的定点纯滚动]\label{mp3:hw:l3-22}
圆锥轴向高度为$h$、半顶角为$\theta\in(0,\pi/2)$，锥尖铰接在水平面上的固定点$O$，圆锥在平面上纯滚动。对称轴绕竖直方向匀速转动，周期$2\pi/\Omega$。当前接触母线沿$x$正向，轴线位于$xz$平面，见图\ref{mp3:fig:hw:l3-22-q}。求底面圆心$O'$、当前最高材料点$B$的速度与加速度。
\end{example}

\begin{figure}[H]\centering
\begin{tikzpicture}[>=Stealth,line cap=round,line join=round]

\draw[axis](0,0)--(4.6,0)node[right]{$x$};\draw[axis](0,0)--(0,3.2)node[above]{$z$};\draw[axis](0,0)--(-1,-.6)node[left]{$y$};\coordinate(C)at(2.6,1.5);\coordinate(B)at(1.733,3);\coordinate(D)at(3.466,0);\draw[thick](B)--(0,0)--(D);\draw[curve,domain=0:360,samples=100]plot({2.6-.867*cos(\x)-.953*sin(\x)},{1.5+1.5*cos(\x)-.5196*sin(\x)});\draw[guide](0,0)--(C)node[midway,above]{$h$};\draw[guide](B)--(D);\fill(C)circle(1.4pt)node[right]{$O'$};\fill(B)circle(1.4pt)node[above]{$B$};\node[below left]at(0,0){$O$};\draw(.7,0)arc(0:30:.7);\node at(.95,.22){$\theta$};\node at(4.7,2.5){$r=h\tan\theta$};

\end{tikzpicture}
\caption{接触母线、圆锥轴与最高点；当前轴线与水平面的夹角为$\theta$}\label{mp3:fig:hw:l3-22-q}
\end{figure}

\end{minipage}\par\medskip
\begin{solution}
令$\vct e=(\cos\theta,0,\sin\theta)$为由锥尖指向底面圆心的轴向，$\vct n=(-\sin\theta,0,\cos\theta)$为截面内向上方向。当前位置为
\[
 \vct r_{O'}=h\vct e,\qquad
 \vct r_B=h\vct e+h\tan\theta\vct n.
\]
纯滚动时接触母线上的材料点速度均为零，故绝对角速度沿$x$轴。若相对进动系的自转角速度为$\nu\vct e$，则
\[
 \vct W=\Omega\vct e_z+\nu\vct e,
 \quad\Omega+\nu\sin\theta=0,
 \quad\boxed{\vct W=-\Omega\cot\theta\vct e_x}.
\]
图\ref{mp3:fig:hw:l3-22-s}中，瞬时轴虽然在地面上，其方向仍随进动旋转；所以
\[
 \vct\alpha=-\Omega\cot\theta\frac{\dd\vct e_x}{\dd t}
 =-\Omega\cot\theta(\Omega\vct e_y)
 =-\Omega^2\cot\theta\vct e_y.
\]
对任一固定材料点使用$\vct v=\vct W\times\vct r$、$\vct a=\vct\alpha\times\vct r+\vct W\times(\vct W\times\vct r)$。对当前位于$xz$平面内的点$\vct r=(x,0,z)$，分别展开为
\[
\begin{aligned}
\vct v&=\Omega z\cot\theta\,\vct e_y,\\
\vct\alpha\times\vct r
&=-\Omega^2z\cot\theta\,\vct e_x
+\Omega^2x\cot\theta\,\vct e_z,\\
\vct W\times(\vct W\times\vct r)
&=-\Omega^2z\cot^2\theta\,\vct e_z.
\end{aligned}
\]
圆心的$x=h\cos\theta,z=h\sin\theta$，两项竖直加速度恰好相消，于是
\[
 \boxed{\vct v_{O'}=\Omega h\cos\theta\vct e_y,
 \qquad\vct a_{O'}=-\Omega^2h\cos\theta\vct e_x}.
\]
最高点的坐标为$[h\cos(2\theta)/\cos\theta,0,2h\sin\theta]$。其欧拉项与二次叉积项分别为
\[
\begin{aligned}
\vct\alpha\times\vct r_B
&=-2\Omega^2h\cos\theta\,\vct e_x
+\frac{\Omega^2h\cos2\theta}{\sin\theta}\vct e_z,\\
\vct W\times(\vct W\times\vct r_B)
&=-\frac{2\Omega^2h\cos^2\theta}{\sin\theta}\vct e_z.
\end{aligned}
\]
竖直分量相加时，$\cos2\theta-2\cos^2\theta=-1$，因此
\[
 \boxed{\vct v_B=2\Omega h\cos\theta\vct e_y},\qquad
 \boxed{\vct a_B=-2\Omega^2h\cos\theta\vct e_x
                        -\frac{\Omega^2h}{\sin\theta}\vct e_z}.
\]
这里跟踪的是当前位于最高位置的材料点。它随后会离开该位置；各时刻的几何最高位置连成的轨迹，属于另一个求导对象。


\begin{figure}[H]\centering
\begin{tikzpicture}[>=Stealth,line cap=round,line join=round]

\draw[axis](0,0)--(3,0)node[right]{接触母线};\draw[axis,thick](0,0)--(-1.8,0)node[left]{$\vct W$};\draw[axis](0,0)--(0,2)node[above]{$\Omega\vct e_z$};\draw[axis](0,0)--(-1.8,-2)node[below left]{$\nu\vct e$};\draw[guide](0,2)--(-1.8,0);\node[right,align=left]at(4.5,1.5){$\vct\alpha=-\Omega^2\cot\theta\vct e_y$\\[6pt]$v_B=2v_{O'}$\\[6pt]$a_{Bz}<0$};

\end{tikzpicture}
\caption{进动、自转合成沿接触母线的瞬时角速度；其方向仍随时间改变}\label{mp3:fig:hw:l3-22-s}
\end{figure}

\end{solution}

\par\medskip\noindent\begin{minipage}{\textwidth}
\begin{example}[匀速进动陀螺上的两个材料点]\label{mp3:hw:l3-23}
半顶角$\alpha$、底面半径$r$的圆锥陀螺，锥尖固定于$O'$。其对称轴与竖直方向夹角为$\theta$，绕竖直轴以恒定角速度$\Omega$进动，同时相对进动系绕自身轴以恒定有符号角速度$\omega$自转。取当前轴线与竖直方向所在平面为$xz$平面，向上的轴向为$\vct e=(\sin\theta,0,\cos\theta)$，陀螺从锥尖沿$-\vct e$伸展。$A$为底面圆周上当前最低材料点，$B$为沿$-y$方向最远的材料点，见图\ref{mp3:fig:hw:l3-23-q}。求两点的速度和加速度。
\end{example}

\begin{figure}[H]\centering
\begin{tikzpicture}[>=Stealth,line cap=round,line join=round]

\coordinate(O)at(0,0);\coordinate(C)at(-1.25,-2.165);\draw[axis](O)--(2,0)node[right]{$x$};\draw[axis](O)--(0,1.4)node[above]{$z$};\draw[axis](O)--(1,.43)node[right]{$y$};\draw[curve,domain=0:360,samples=100]plot({-1.25+.9*(.866*cos(\x)+.7*sin(\x))},{-2.165+.9*(-.5*cos(\x)+.3*sin(\x))});\draw[thick](-2.029,-1.715)--(O)--(-.471,-2.615);\draw[guide](O)--(C)node[midway,left]{$h$};\fill(C)circle(1.4pt)node[right]{$C$};\fill(-.471,-2.615)circle(1.5pt)node[below right]{$A$};\fill(-1.88,-2.435)circle(1.5pt)node[below left]{$B$};\node[above left]at(O){$O'$};\draw[axis](O)--(.6,1.039)node[right]{$\vct e$};\node[right]at(2,-1.6){$h=r\cot\alpha$};

\end{tikzpicture}
\caption{陀螺的轴向、底面及两种材料点；$B$沿$-y$方向偏离纸面}\label{mp3:fig:hw:l3-23-q}
\end{figure}

\end{minipage}\par\medskip
\begin{solution}
记$h=r\cot\alpha$、$s=\sin\theta$、$c=\cos\theta$，并取截面内向下偏移的单位矢量$\vct t=(c,0,-s)$。当前
\[
 \vct r_A=-h\vct e+r\vct t=(-hs+rc,0,-hc-rs),
 \qquad\vct r_B=-h\vct e-r\vct e_y=(-hs,-r,-hc).
\]
两种角速度必须作矢量相加。图\ref{mp3:fig:hw:l3-23-s}中，绝对角速度和角加速度为
\[
 \vct W=\Omega\vct e_z+\omega\vct e,
 \qquad\vct\alpha=\Omega\omega\sin\theta\vct e_y.
\]
后一式来自$\dot{\vct e}=\Omega\vct e_z\times\vct e$，即使$\Omega,\omega$均恒定，角加速度也一般不为零。为展开叉积，暂记$p=\omega s$、$q=\Omega+\omega c$、$k=\Omega\omega s$，则$\vct W=(p,0,q)$、$\vct\alpha=(0,k,0)$。对当前位矢为$(x,y,z)$的固定材料点，
\[
\begin{aligned}
\vct v&=(-qy,\ qx-pz,\ py),\\
\vct\alpha\times\vct r&=(kz,0,-kx),\\
\vct W\times(\vct W\times\vct r)
&=(-q^2x+pqz,\ -(p^2+q^2)y,\ pqx-p^2z).
\end{aligned}
\]
所以$a_x=(k+pq)z-q^2x$、$a_y=-(p^2+q^2)y$、$a_z=(pq-k)x-p^2z$。代入$A$的坐标，并用$s^2+c^2=1$合并自转项和混合项，得到
\[
 \boxed{\vct v_A=[-\Omega hs+\Omega rc+\omega r]\vct e_y},
\]
\[
 \boxed{\vct a_A=[\Omega^2hs-\Omega^2rc-2\Omega\omega r
                     -\omega^2rc]\vct e_x
                     +\omega^2rs\vct e_z}.
\]
对$B$，代入$x=-hs,y=-r,z=-hc$。由于$pq-k=\omega^2sc$，竖直加速度为$\omega^2sc(-hs)-\omega^2s^2(-hc)=0$；其余分量为
\[
 \boxed{\vct v_B=r(\Omega+\omega c)\vct e_x
                 -\Omega hs\vct e_y-\omega rs\vct e_z},
\]
\[
 \boxed{\vct a_B=\Omega^2hs\vct e_x
       +r(\Omega^2+2\Omega\omega c+\omega^2)\vct e_y}.
\]
若改在进动系中计算，$A,B$相对该系只随自转运动，分别加入$2\Omega\vct e_z\times\vct v'$及$\Omega\vct e_z\times(\Omega\vct e_z\times\vct r)$，会得到相同结果。取$\Omega=0$可还原定轴自转；取$\omega=0$则整个锥体只有绕竖直轴的定轴转动。混合项随进动消失，自转项则在停止自转时消失。


\begin{figure}[H]\centering
\begin{tikzpicture}[>=Stealth,line cap=round,line join=round]

\draw[axis](0,0)--(0,2.3)node[above]{$\Omega\vct e_z$};\draw[axis](0,0)--(1.25,2.165)node[right]{$\omega\vct e$};\draw[axis,thick](0,0)--(1.25,4.465)node[above]{$\vct W$};\draw[guide](0,2.3)--(1.25,4.465);\draw[guide](1.25,2.165)--(1.25,4.465);\node[right,align=left]at(2.4,2.2){进动使轴向改变：\\[4pt]$\dot{\vct e}=\Omega\vct e_z\times\vct e$\\[6pt]$\vct\alpha=\omega\dot{\vct e}$};

\end{tikzpicture}
\caption{进动与自转的角速度合成；轴向变化产生角加速度}\label{mp3:fig:hw:l3-23-s}
\end{figure}

\end{solution}


\setcounter{chapter}{3}

\section{习题与解答}

\par\medskip\noindent\begin{minipage}{\textwidth}
\begin{exercise}[椭圆上的曲率极值]\label{mp3:exercise:l3-p1}

质点的位置为 $x=a\cos\omega t,\ y=b\sin\omega t$，其中 $a>b>0,\omega>0$，见图\ref{mp3:fig:l3-p1-q}。求曲率半径及其极值；判断加速度是否总沿主法线，并求第一象限内一般点的曲率中心。

\end{exercise}
\begin{figure}[H]\centering
\begin{tikzpicture}[scale=.95]

\draw[axis](-2.8,0)--(3,0)node[right]{$x$};\draw[axis](0,-1.5)--(0,1.8)node[above]{$y$};
\draw[curve](0,0)ellipse(2.3 and 1.15);
\fill(2.3,0)circle(1.3pt)node[below right]{$A(a,0)$};
\fill(0,1.15)circle(1.3pt)node[above left]{$B(0,b)$};
\fill({2.3*cos(40)},{1.15*sin(40)})circle(1.3pt)node[above right]{$P$};
\draw[guide](0,0)--({2.3*cos(40)},{1.15*sin(40)});
\node[below left]at(0,0){$O$};

\end{tikzpicture}
\caption{椭圆的参数角为 $\omega t$，一般不同于位置矢量的极角。}\label{mp3:fig:l3-p1-q}
\end{figure}
\end{minipage}\par\medskip
\begin{solution}
令 $q=\omega t$，并记 $D=\sqrt{a^2\sin^2q+b^2\cos^2q}$。先列出参数的一、二阶导数：
\[
x_q=-a\sin q,\quad y_q=b\cos q,\qquad
x_{qq}=-a\cos q,\quad y_{qq}=-b\sin q.
\]
于是
\[
x_qy_{qq}-y_qx_{qq}=ab\sin^2q+ab\cos^2q=ab,
\qquad \rho=\frac{D^3}{ab}.
\]
因为 $b\le D\le a$，曲率半径的最小值与最大值分别为
\[
\rho_{\min}=\frac{b^2}{a}\quad(q=0,\pi),\qquad
\rho_{\max}=\frac{a^2}{b}\quad(q=\pi/2,3\pi/2).
\]
图\ref{mp3:fig:l3-p1-s}画出两处极值点的曲率中心。长轴端的切线随弧长转动得更快，因此曲率较大。

\begin{figure}[H]\centering
\begin{tikzpicture}[scale=.95]

\draw[curve](0,0)ellipse(2 and 1);
\fill(2,0)circle(1.5pt)node[right]{$A$};
\fill(1.5,0)circle(1.5pt)node[above]{$K_A$};
\draw[axis](2,0)--(1.5,0);\node[below]at(1.75,-.1){$b^2/a$};
\fill(0,1)circle(1.5pt)node[above]{$B$};
\fill(0,-3)circle(1.5pt)node[left]{$K_B$};
\draw[guide](0,1)--(0,-3);\node[right]at(.05,-1.7){$a^2/b$};
\draw[axis](0,1)--(0,.1);
\node[anchor=west]at(3,0){$K_A=(a-b^2/a,0)$};
\node[anchor=west]at(3,-1){$K_B=(0,b-a^2/b)$};

\end{tikzpicture}
\caption{以 $a=2b$ 示意两处曲率中心。$K_B$ 可位于椭圆外，曲率圆只描述 $B$ 附近的弯曲。}\label{mp3:fig:l3-p1-s}
\end{figure}


速度的方向为 $(-a\sin q,b\cos q)$，主法线为
\[
\vct e_n=\frac{(-b\cos q,-a\sin q)}D.
\]
由 $\vct r_K=\vct r_P+\rho\vct e_n$，逐分量代入，
\[
\begin{aligned}
x_K&=a\cos q-\frac{D^2}{a}\cos q
=\frac{a^2-b^2}{a}\cos^3q,\\
y_K&=b\sin q-\frac{D^2}{b}\sin q
=\frac{b^2-a^2}{b}\sin^3q.
\end{aligned}
\]
加速度 $\vct a=-\omega^2(x,y)$ 始终指向 $O$，但 $O$ 一般不等于 $K$。直接计算
\[
\vct v\cdot\vct a
=\omega^3(a^2-b^2)\sin q\cos q,
\]
除四个轴端点外，椭圆上的速率随位置改变，因而加速度同时含切向和法向分量。

由曲率中心的坐标还可判断它与椭圆的位置关系。图\ref{mp3:fig:l3-p1-s}中，$K_A$位于椭圆内部；$K_B$的纵坐标为$-(a^2-b^2)/b<0$。当$(a^2-b^2)/b>b$，即$a>\sqrt2b$时，$K_B$位于椭圆外部。曲率中心由局部弯曲决定，并不要求位于封闭轨迹内部。
\end{solution}

\par\medskip\noindent\begin{minipage}{\textwidth}
\begin{exercise}[心形线的正则段与尖点（选做）]\label{mp3:exercise:l3-p2}

质点沿极坐标曲线 $r=a(1+\cos\varphi)$ 运动，$a>0$，$\dot\varphi=\Omega>0$，且 $-\pi<\varphi<\pi$。求速率、切向加速度、法向加速度和曲率半径，并说明 $\varphi\to\pm\pi$ 时直接使用自然坐标公式须注意什么。图\ref{mp3:fig:l3-p2-q}给出完整曲线的形状。

\end{exercise}
\begin{figure}[H]\centering
\begin{tikzpicture}[scale=.95]

\draw[axis](-1,0)--(4.2,0)node[right]{$x$};\draw[axis](0,-2.4)--(0,2.5)node[above]{$y$};
\draw[curve,domain=-180:180,samples=200]plot({1.6*(1+cos(\x))*cos(\x)},{1.6*(1+cos(\x))*sin(\x)});
\fill(0,0)circle(1.5pt)node[below left]{$O$};
\fill(3.2,0)circle(1.5pt)node[below right]{$A$};
\coordinate(P)at({1.6*(1+cos(55))*cos(55)},{1.6*(1+cos(55))*sin(55)});
\draw[guide](0,0)--(P);\fill(P)circle(1.5pt)node[right]{$P$};
\draw(.5,0)arc(0:55:.5);\node at(.8,.3){$\varphi$};

\end{tikzpicture}
\caption{心形线在 $O$ 有尖点。题目规定的开区间不含该点。}\label{mp3:fig:l3-p2-q}
\end{figure}
\end{minipage}\par\medskip
\begin{solution}
在给定区间内 $\cos(\varphi/2)>0$。由 $r'=-a\sin\varphi$、$r''=-a\cos\varphi$，
\[
\begin{aligned}
r^2+(r')^2
&=a^2[(1+\cos\varphi)^2+\sin^2\varphi]
=4a^2\cos^2\frac\varphi2,\\
r^2+2(r')^2-rr''
&=a^2[(1+\cos\varphi)^2+2\sin^2\varphi
+(1+\cos\varphi)\cos\varphi]\\
&=3a^2(1+\cos\varphi)=6a^2\cos^2\frac\varphi2.
\end{aligned}
\]
所以
\[
v=2a\Omega\cos\frac\varphi2,\qquad
\rho=\frac{4a}{3}\cos\frac\varphi2,
\]
\[
a_t=\dot v=2a\Omega\left(-\frac12\sin\frac\varphi2\right)\Omega
=-a\Omega^2\sin\frac\varphi2,\qquad
a_n=\frac{v^2}{\rho}=3a\Omega^2\cos\frac\varphi2.
\]
图\ref{mp3:fig:l3-p2-s}在 $A$ 处标出曲率中心：$\varphi=0$ 时 $v=2a\Omega$，主法线向左，$\rho=4a/3$，因此 $x_K=2a-4a/3=2a/3$。

\begin{figure}[H]\centering
\begin{tikzpicture}[scale=.95]

\draw[curve,domain=-180:180,samples=160]plot({1.5*(1+cos(\x))*cos(\x)},{1.5*(1+cos(\x))*sin(\x)});
\draw[guide](1,0)circle(2);
\fill(3,0)circle(1.5pt)node[right]{$A$};
\fill(1,0)circle(1.5pt)node[below]{$K$};
\draw[axis](3,0)--(1.6,0)node[midway,below]{$\vct e_n$};
\draw[axis](3,0)--(3,1.2)node[above]{$\vct v$};
\node[above]at(1,.1){$AK=4a/3$};
\node[below left]at(0,0){尖点 $O$};

\end{tikzpicture}
\caption{$A$ 处的曲率圆。圆与心形线只在局部具有相同的二阶几何性质。}\label{mp3:fig:l3-p2-s}
\end{figure}


当 $\varphi\to\pm\pi$，$v,\rho,a_n$ 均趋于零，而 $a_t$ 的两侧极限分别为 $\mp a\Omega^2$。在尖点本身，$\vct r_\varphi=0$，参数不正则，单位速度方向没有连续的双侧定义。参数位置仍可二次求导，笛卡尔加速度保持有限；不能将尖点处的 $v=0,\rho=0$ 直接代入 $v^2/\rho$。

\end{solution}

\par\medskip\noindent\begin{minipage}{\textwidth}
\begin{exercise}[凸形轨道上的接触条件]\label{mp3:exercise:l3-p3}

小球从凸形轨道顶点 $O$ 以水平速率 $u>0$ 向右运动。轨道光滑，方程 $y=-x^2/(2p)$，$p>0$，小球位于轨道上方，支持力只能把它推离轨道，见图\ref{mp3:fig:l3-p3-q}。求维持接触时的支持力，并区分 $u^2<gp,\ u^2=gp,\ u^2>gp$ 三种情形。

\end{exercise}
\begin{figure}[H]\centering
\begin{tikzpicture}[scale=.95]

\draw[axis](-.5,0)--(4.2,0)node[right]{$x$};\draw[axis](0,-2.3)--(0,1)node[above]{$y$};
\draw[curve,domain=-.5:3.1,samples=70]plot(\x,{-.2*\x*\x});
\fill(0,0)circle(1.5pt)node[above left]{$O$};
\draw[axis](0,0)--(1.1,0)node[above]{$u$};
\fill(2,-.8)circle(2pt)node[right]{$P$};
\node at(2.5,.65){$y=-x^2/(2p)$};

\end{tikzpicture}
\caption{小球在凸轨道上表面运动；向轨道凹侧的主法线朝左下方。}\label{mp3:fig:l3-p3-q}
\end{figure}
\end{minipage}\par\medskip
\begin{solution}
若接触能够维持，由机械能守恒，
\[
v^2=u^2+\frac{gx^2}{p}.
\]
记 $q=x/p$。由 $y'=-x/p=-q$、$y''=-1/p$，得 $\rho=(1+y'^2)^{3/2}/|y''|=p(1+q^2)^{3/2}$，
$\vct e_n=(-q,-1)/\sqrt{1+q^2}$。图\ref{mp3:fig:l3-p3-s}中支持力与 $\vct e_n$ 反向，故
\[
\frac{mg}{\sqrt{1+q^2}}-N=\frac{mv^2}{p(1+q^2)^{3/2}},
\quad
N=\frac{m(gp-u^2)}{p(1+q^2)^{3/2}}.
\]

\begin{figure}[H]\centering
\begin{tikzpicture}[scale=.95]

\draw[curve,domain=-.5:2.2,samples=60]plot(\x,{-.3*\x*\x});
\coordinate(P)at(1,-.3);\fill(P)circle(2pt)node[right]{$P$};
\draw[axis](P)--(1,-1.65)node[below]{$mg$};
\draw[axis](P)--(1.55,.62)node[above]{$N$};
\draw[axis,guide](P)--(.45,-1.22)node[left]{$\vct e_n$};
\begin{scope}[shift={(5,0)}]
\draw[curve,domain=0:2.8,samples=70]plot(\x,{-.22*\x*\x});
\draw[guide,domain=0:2.8,samples=70]plot(\x,{-.12*\x*\x});
\fill(0,0)circle(1.3pt)node[above left]{$O$};
\node[right]at(2.8,-.94){抛物线};
\node[right]at(2.8,-1.73){轨道};
\node at(1.4,.65){$u^2>gp$};
\end{scope}

\end{tikzpicture}
\caption{接触时的法向受力与超临界初速下的脱离轨迹。右图两条曲线具有相同的初始切线。}\label{mp3:fig:l3-p3-s}
\end{figure}


当 $u^2<gp$，所需 $N$ 始终为正，完整的无限长轨道上可一直保持接触。当 $u^2=gp$，$N=0$，自由抛体轨迹
$y=-gx^2/(2u^2)$ 恰与轨道重合；这是无压紧力的临界运动，对扰动不具有稳健接触。

当 $u^2>gp$，约束解从顶点开始就要求 $N<0$，轨道无法提供这种拉力，小球立即脱离。实际轨迹为 $y=-gx^2/(2u^2)$，在 $x>0$ 处高于轨道，见右图。负的约束力结果只用于判定脱离，不能当作实际接触力。

将脱离后的轨迹与轨道直接相减，可进一步核实相对位置：
\[
 y_{\text{自由}}-y_{\text{轨道}}
 =\frac{x^2}{2p}\left(1-\frac{gp}{u^2}\right).
\]
当$u^2>gp$，差值为正，与图\ref{mp3:fig:l3-p3-s}右图一致；轨道在小球下方，无法继续施力。当$u^2<gp$，自由轨迹反而会进入轨道内部，实际运动需由正支持力阻止这种穿入。临界值使两条抛物线完全重合，接触关系在整个轨迹上成立。
\end{solution}

\par\medskip\noindent\begin{minipage}{\textwidth}
\begin{exercise}[匀速移动支点上的转杆]\label{mp3:exercise:l3-p4}

刚杆 $AB$ 长为 $l$。支点 $A$ 沿水平线以速度 $u>0$ 匀速运动，杆相对地面以恒定逆时针角速度 $\omega>0$ 转动。取某时刻 $A=(0,0)$、$B=(l,0)$，见图\ref{mp3:fig:l3-p4-q}。求该时刻 $B$ 的速度、加速度、曲率半径及速度瞬心。

\end{exercise}
\begin{figure}[H]\centering
\begin{tikzpicture}[scale=.95]

\draw[thick](-.7,0)--(4.4,0);\draw[very thick](0,0)--(3,0);
\fill(0,0)circle(1.5pt)node[below]{$A$};\fill(3,0)circle(1.5pt)node[below]{$B$};
\draw[axis](0,.45)--(1.2,.45)node[right]{$u$};
\draw[axis](1.9,.3)arc(15:100:.7);\node at(1.7,1.3){$\omega$};
\draw[|-|](0,-.6)--(3,-.6)node[midway,fill=white]{$l$};

\end{tikzpicture}
\caption{支点在惯性系中匀速移动，杆的自转角速度恒定。}\label{mp3:fig:l3-p4-q}
\end{figure}
\end{minipage}\par\medskip
\begin{solution}
取$\vct k$垂直纸面向外，题设$u,\omega>0$。在图\ref{mp3:fig:l3-p4-q}所示瞬间，支点加速度为零，杆的角加速度为零，故式\eqref{mp3:eq:l3-rigid-v}、\eqref{mp3:eq:l3-rigid-a}给出
\[
 \vct v_B=u\vct i+\omega l\vct j,\qquad
 \vct a_B=-\omega^2l\vct i.
\]
记$S=u^2+\omega^2l^2=v_B^2$。由于
$|\vct v_B\times\vct a_B|=\omega^3l^2$，曲率半径为
\[
 \boxed{\rho_B=\frac{S^{3/2}}{\omega^3l^2}}.
\]
\begin{figure}[H]\centering
\begin{tikzpicture}[scale=.95]

\coordinate(A)at(0,0);\coordinate(B)at(3,0);\coordinate(I)at(0,1.4);
\coordinate(K)at(-.653333,1.704889);
\draw[thick](A)--(B);\draw[guide](A)--(I)--(B);
\draw[guide](I)--(K);
\fill(K)circle(1.5pt)node[above left]{$K_B$};
\fill(A)circle(1.5pt)node[below]{$A$};\fill(B)circle(1.5pt)node[below]{$B$};
\fill(I)circle(1.5pt)node[above]{$I$};
\draw[axis](B)--(3.55,1.1786)node[above right]{$\vct v_B$};
\draw[axis,thick](B)--(1.65,0)node[midway,below]{$\vct a_B$};
\draw[axis](B)--(3.55,0)node[below]{$u\vct i$};
\draw[axis](B)--(3,1.1786)node[above left]{$\omega l\vct j$};
\node[left]at(0,.65){$u/\omega$};

\end{tikzpicture}
\caption{合速度垂直于$IB$，加速度沿杆向左。曲率中心$K_B$在$BI$的延长线上，$BK_B=\rho_B>BI$。}\label{mp3:fig:l3-p4-s}
\end{figure}
图\ref{mp3:fig:l3-p4-s}中的瞬心满足
$\vct v_A+\omega\vct k\times(\vct r_I-\vct r_A)=0$，
所以$I=(0,u/\omega)$。于是
$IB=\sqrt{l^2+u^2/\omega^2}=\sqrt S/\omega$，且$v_B=\omega IB$。
但$\rho_B/IB=S/(\omega^2l^2)>1$，瞬时转动半径与轨迹曲率半径不同。

求曲率中心$K_B$时，先将加速度分成切向、法向两部分。切向投影为
\[
 \vct a_\parallel
 =\frac{\vct a_B\cdot\vct v_B}{S}\vct v_B
 =-\frac{u\omega^2l}{S}(u\vct i+\omega l\vct j).
\]
因此
\[
 \vct a_N=\vct a_B-\vct a_\parallel
 =\frac{\omega^3l^2}{S}(-\omega l\vct i+u\vct j),
 \qquad |\vct a_N|^2=\frac{\omega^6l^4}{S}.
\]
法向加速度沿主法线指向曲率中心，且$|\vct a_N|=v_B^2/\rho_B$，故
\[
 \vct r_{K_B}-\vct r_B
 =\frac{v_B^2}{|\vct a_N|^2}\vct a_N
 =\boxed{-\frac{S}{\omega^2l}\vct i
 +\frac{uS}{\omega^3l^2}\vct j}.
\]
该位移与合速度垂直，横向为负、竖向为正，因而$K_B$位于$B$的左上方，在图中$BI$的延长线上；$BK_B=\rho_B>BI$。
\end{solution}

\par\medskip\noindent\begin{minipage}{\textwidth}
\begin{exercise}[外滚等圆的轮缘轨迹（选做）]\label{mp3:exercise:l3-p5}

半径为 $a$ 的小圆在半径同为 $a$ 的固定圆外纯滚动，轮心 $C$ 绕固定圆心 $O$ 以角速度 $\Omega>0$ 逆时针公转。初始接触点为 $(a,0)$，材料点 $P$ 初始就在该处，见图\ref{mp3:fig:l3-p5-q}。求 $P$ 的参数轨迹、速率和曲率半径，并给出第一周中的尖点位置。

\end{exercise}
\begin{figure}[H]\centering
\begin{tikzpicture}[scale=.95]

\draw[curve](0,0)circle(1);\draw[axis](-1.4,0)--(4.2,0)node[right]{$x$};
\coordinate(C)at({2*cos(35)},{2*sin(35)});
\draw[curve](C)circle(1);\draw[guide](0,0)--(C);
\fill(0,0)circle(1.4pt)node[below]{$O$};\fill(C)circle(1.4pt)node[above]{$C$};
\fill({cos(35)},{sin(35)})circle(1.4pt)node[left]{$I$};
\fill(1,0)circle(1.4pt)node[below]{$P(0)$};
\draw(.45,0)arc(0:35:.45);\node at(.7,.17){$\vartheta$};
\node at(4,1.8){$OC=2a$};

\end{tikzpicture}
\caption{外滚时接触半径从 $C$ 指向 $I$，朝向固定圆心。}\label{mp3:fig:l3-p5-q}
\end{figure}
\end{minipage}\par\medskip
\begin{solution}
轮心速度为 $2a\Omega\vct e_\varphi$，而 $\overrightarrow{CI}=-a\vct e_r$。接触速度为零给出
$2a\Omega-a\omega=0$，故小圆自转角速度 $\omega=2\Omega$。初始 $\overrightarrow{CP}=-a\vct i$，因而在 $\vartheta=\Omega t$ 时
\[
x=2a\cos\vartheta-a\cos2\vartheta,\qquad
y=2a\sin\vartheta-a\sin2\vartheta.
\]
图\ref{mp3:fig:l3-p5-s}画出一整周轨迹及 $\vartheta=\pi/2$ 时的位置。此时 $C=(0,2a)$、$P=(a,2a)$，可以直接核对自转项的符号。

\begin{figure}[H]\centering
\begin{tikzpicture}[scale=.95]

\draw[axis](-3.4,0)--(2.2,0)node[right]{$x/a$};
\draw[axis](0,-2.9)--(0,3.1)node[above]{$y/a$};
\draw[curve,domain=0:360,samples=200]plot({2*cos(\x)-cos(2*\x)},{2*sin(\x)-sin(2*\x)});
\draw[guide](0,0)circle(1);
\draw[guide](0,2)circle(1);
\fill(0,0)circle(1.4pt)node[below left]{$O$};
\fill(1,0)circle(1.4pt)node[below right]{尖点};
\fill(0,2)circle(1.4pt)node[left]{$C$};
\fill(1,2)circle(1.5pt)node[right]{$P$};
\draw[thick](0,2)--(1,2);
\node at(3.3,1.5){$\vartheta=\pi/2$};
\node at(3.3,.5){$CP=a$};

\end{tikzpicture}
\caption{外摆线的完整一周。虚线圆分别为固定圆与 $\vartheta=\pi/2$ 时的小圆。}\label{mp3:fig:l3-p5-s}
\end{figure}


对参数式连续求导，得到
\[
\begin{aligned}
\vct v&=2a\Omega(-\sin\vartheta+\sin2\vartheta,\ \cos\vartheta-\cos2\vartheta),\\
\vct a&=2a\Omega^2(-\cos\vartheta+2\cos2\vartheta,\ {-\sin\vartheta}+2\sin2\vartheta).
\end{aligned}
\]
两分量平方相加时，利用 $\cos\vartheta\cos2\vartheta+\sin\vartheta\sin2\vartheta=\cos\vartheta$，得
\[
v^2=4a^2\Omega^2(2-2\cos\vartheta)=8a^2\Omega^2(1-\cos\vartheta),
\]
\[
(\vct v\times\vct a)\cdot\vct k
=4a^2\Omega^3(3-3\cos\vartheta)
=12a^2\Omega^3(1-\cos\vartheta).
\]
所以在 $0<\vartheta<2\pi$，
\[
v=4a\Omega\sin\frac\vartheta2,\qquad
\rho=\frac{8a}{3}\sin\frac\vartheta2.
\]
首尾 $\vartheta=0,2\pi$ 对应同一个尖点 $(a,0)$，其余位置 $v>0$。这一周中，连心线相对地面转过 $2\pi$，材料半径相对地面转过 $4\pi$；两者之差是小圆相对连心线转过的角度。

\end{solution}

\par\medskip\noindent\begin{minipage}{\textwidth}
\begin{exercise}[瞬时平移与不同的轨迹曲率]\label{mp3:exercise:l3-p6}

水平刚杆 $AB$ 长为 $l$。某时刻两端速度均为 $u\vct i$，其中 $u>0$；$A$ 的加速度为零，杆的角加速度为 $\beta\vct k$，$\beta>0$，见图\ref{mp3:fig:l3-p6-q}。求该时刻两端的曲率半径，并说明能否由这一瞬间的速度相等推出杆在一段时间内平移。

\end{exercise}
\begin{figure}[H]\centering
\begin{tikzpicture}[scale=.95]

\draw[very thick](0,0)--(3,0);
\fill(0,0)circle(1.4pt)node[below]{$A$};\fill(3,0)circle(1.4pt)node[below]{$B$};
\draw[axis](0,.3)--(1.1,.3)node[right]{$u$};
\draw[axis](3,.3)--(4.1,.3)node[right]{$u$};
\draw[|-|](0,-.5)--(3,-.5)node[midway,fill=white]{$l$};
\node at(1.5,1.25){$\vct\alpha=\beta\vct k$};

\end{tikzpicture}
\caption{两端瞬时速度相同，但角加速度不为零。}\label{mp3:fig:l3-p6-q}
\end{figure}
\end{minipage}\par\medskip
\begin{solution}
由两端的相对速度 $\omega\vct k\times(l\vct i)=0$ 可得 $\omega=0$。式\eqref{mp3:eq:l3-rigid-a}中的向心项此时为零，但角加速度项仍存在：
\[
\vct a_B=\beta l\vct j,\qquad \vct a_A=0.
\]
图\ref{mp3:fig:l3-p6-s}显示 $B$ 在该时刻向上弯曲，而 $A$ 的瞬时曲率为零。于是
\[
\kappa_A=0,\quad \rho_A\to\infty,\qquad
\rho_B=\frac{u^2}{\beta l}.
\]

\begin{figure}[H]\centering
\begin{tikzpicture}[scale=.95]

\draw[guide](-.5,0)--(4.5,0);
\draw[thick](0,0)--(3,0);
\fill(0,0)circle(1.4pt)node[below]{$A$};\fill(3,0)circle(1.4pt)node[below]{$B$};
\draw[axis](3,0)--(3,1.4)node[above]{$\beta l\vct j$};
\draw[curve,domain=-.55:.65,samples=50]plot({3+\x},{1.2*\x*\x});
\draw[axis](0,0)--(1,0)node[below]{$\vct v_A$};
\draw[axis](3,0)--(4,0)node[below]{$\vct v_B$};
\node at(.6,1.1){$\vct a_A=0$};
\node at(5.4,.8){$B$ 的局部轨迹};

\end{tikzpicture}
\caption{速度相同只约束轨迹的一阶变化。$B$ 的加速度垂直于速度，曲率有限。}\label{mp3:fig:l3-p6-s}
\end{figure}


以该时刻为 $t=0$，在运动足够光滑时，位置的局部展开为
\[
\vct r_A(t)=ut\,\vct i+O(t^3),\qquad
\vct r_B(t)=(l+ut)\vct i+\frac12\beta l t^2\vct j+O(t^3).
\]
相同的一阶项给出相同速度，不同的二阶项给出不同弯曲。由于 $\omega(t)=\beta t+O(t^2)$，杆随即出现转动；一段时间内的平移须要求该时间段内 $\omega$ 始终为零。

用图\ref{mp3:fig:l3-p6-s}中的局部坐标消去$t$，$B$的轨迹在二阶近似下满足
\[
 y_B=\frac{\beta l}{2u^2}(x_B-l)^2+O((x_B-l)^3).
\]
在所讨论点斜率为零，二阶导数为$\beta l/u^2$，故曲率为$\beta l/u^2$，与自然坐标法一致。$A$的加速度为零，因而此刻曲率为零。轨迹的更高阶项仍可能产生弯曲，尚不足以确定$A$在整个过程中的路径。
\end{solution}

\par\medskip\noindent\begin{minipage}{\textwidth}
\begin{exercise}[变速螺线运动]\label{mp3:exercise:l3-p7}

质点沿 $r=a,\ z=h\varphi$ 的圆柱螺线运动，$a,h>0$，但 $\varphi(t)$ 不再匀速，所讨论区间满足 $\dot\varphi>0$。见图\ref{mp3:fig:l3-p7-q}。求加速度的柱坐标分量与自然分量，说明 $\ddot\varphi$ 是否改变轨迹的曲率。

\end{exercise}
\begin{figure}[H]\centering
\begin{tikzpicture}[scale=.95]

\begin{scope}[x={(.85cm,-.2cm)},y={(.55cm,.3cm)},z={(0cm,.7cm)}]
\draw[axis](0,0,0)--(0,0,4.6)node[above]{$z$};
\draw[curve,domain=0:540,samples=160]plot({1.3*cos(\x)},{1.3*sin(\x)},{.007*\x});
\coordinate(P)at(1.3,0,2.52);\fill(P)circle(1.5pt)node[right]{$P$};
\draw[guide](0,0,2.52)--(P);
\node[above]at(.65,0,2.52){$a$};
\end{scope}
\node at(4.2,2.4){$z=h\varphi$};\node at(4.2,1.4){$\dot\varphi>0$};

\end{tikzpicture}
\caption{同一条空间轨迹可按不同的时间规律通过。}\label{mp3:fig:l3-p7-q}
\end{figure}
\end{minipage}\par\medskip
\begin{solution}
位置约束给出 $\dot r=\ddot r=0,\ \dot z=h\dot\varphi,\ \ddot z=h\ddot\varphi$，故
\[
\vct a=-a\dot\varphi^2\vct e_r
+a\ddot\varphi\vct e_\varphi+h\ddot\varphi\vct e_z.
\]
记 $L=\sqrt{a^2+h^2}$，则 $v=L\dot\varphi$，
$\vct e_t=(a\vct e_\varphi+h\vct e_z)/L$。如图\ref{mp3:fig:l3-p7-s}，后两项的合矢量恰沿切线，第一项沿主法线。因此
\[
a_t=L\ddot\varphi,\qquad a_n=a\dot\varphi^2,\qquad
\rho=\frac{L^2}{a}.
\]

\begin{figure}[H]\centering
\begin{tikzpicture}[scale=.95]

\fill(0,0)circle(1.5pt)node[below]{$P$};
\draw[axis](0,0)--(1.35,0)node[right]{$a\ddot\varphi\,\vct e_\varphi$};
\draw[axis](0,0)--(0,1.05)node[above]{$h\ddot\varphi\,\vct e_z$};
\draw[axis,thick](0,0)--(1.35,1.05)node[above right]{$a_t\vct e_t$};
\draw[guide](1.35,0)--(1.35,1.05)--(0,1.05);
\begin{scope}[shift={(5.1,.4)}]
\fill(0,0)circle(1.5pt)node[below]{$P$};
\draw[axis](0,0)--(-1.5,0)node[above]{$a_n\vct e_n=-a\dot\varphi^2\vct e_r$};
\node at(0,-1){$\vct e_n=-\vct e_r$};
\end{scope}

\end{tikzpicture}
\caption{左图在水平切向与竖直方向张成的平面内合成切向加速度，取 $\ddot\varphi>0$；右图单列水平径向的法向项。}\label{mp3:fig:l3-p7-s}
\end{figure}


角加速度改变通过轨迹的快慢，不改变轨迹本身，因而不改变 $\rho$ 或 $\tau=h/L^2$。若 $\dot\varphi$ 降至零，$a_n$ 同时为零；该瞬间须由两侧轨迹确定几何曲率，不直接计算 $v^2/a_n$。

曲率也可不经自然分量直接求出。$\vct v=\dot\varphi(a\vct e_\varphi+h\vct e_z)$，而加速度中正比于$\ddot\varphi$的部分与速度平行，不参加叉积。因此
\[
 |\vct v\times\vct a|=a\sqrt{a^2+h^2}\,\dot\varphi^3,
 \qquad
 \kappa=\frac{|\vct v\times\vct a|}{v^3}
 =\frac a{a^2+h^2}.
\]
图\ref{mp3:fig:l3-p7-s}左图两切向分量合成后与速度同向或反向；右图的法向分量只由角速度平方决定。给定的空间路径保持不变，变速只改变质点沿该路径的运动进程。
\end{solution}

\par\medskip\noindent\begin{minipage}{\textwidth}
\begin{exercise}[球面纬线的曲率中心]\label{mp3:exercise:l3-p8}

质点在半径为 $a$ 的球面上沿纬线匀速运动，球坐标满足 $r=a,\theta=\theta_0,\varphi=\Omega t$，其中 $0<\theta_0<\pi,\Omega>0$。见图\ref{mp3:fig:l3-p8-q}。求球坐标加速度分量、轨迹曲率半径及曲率中心，解释为何不能令 $\rho=a$。

\end{exercise}
\begin{figure}[H]\centering
\begin{tikzpicture}[scale=.95]

\draw[curve](0,0)circle(1.8);\draw[axis](0,-2.1)--(0,2.4)node[above]{$z$};
\draw[guide](0,.9)ellipse(1.559 and .26);
\coordinate(P)at(1.559,.9);\draw[thick](0,0)--(P);\fill(P)circle(1.5pt)node[right]{$P$};
\fill(0,.9)circle(1.4pt)node[left]{$H$};\fill(0,0)circle(1.4pt)node[below left]{$O$};
\draw(0,.6)arc(90:30:.6);\node at(.4,.65){$\theta_0$};
\node at(3.7,1){$OP=a$};
\node at(3.7,0){$PH=a\sin\theta_0$};

\end{tikzpicture}
\caption{纬线圆位于高度 $z=a\cos\theta_0$ 的水平面内。}\label{mp3:fig:l3-p8-q}
\end{figure}
\end{minipage}\par\medskip
\begin{solution}
将定值坐标代入式\eqref{mp3:eq:l3-spherical-a}，
\[
a_r=-a\Omega^2\sin^2\theta_0,\qquad
a_\theta=-a\Omega^2\sin\theta_0\cos\theta_0,\qquad
a_\varphi=0.
\]
用柱坐标的水平径向单位向量 $\vct e_R$ 表示，有
$\vct e_R=\sin\theta_0\vct e_r+\cos\theta_0\vct e_\theta$，所以
\[
\vct a=-a\Omega^2\sin\theta_0\,\vct e_R.
\]
它水平指向纬线圆心 $H$。速率为 $v=a\Omega\sin\theta_0$，故
\[
\rho=\frac{v^2}{|\vct a|}=a\sin\theta_0,\qquad K=H.
\]

\begin{figure}[H]\centering
\begin{tikzpicture}[scale=.95]

\coordinate(O)at(0,0);\coordinate(P)at(2.1,1.21);\coordinate(H)at(0,1.21);
\draw[guide](O)--(P);\draw[guide](O)--(H);
\draw[thick](H)--(P);\fill(O)circle(1.4pt)node[below]{$O$};
\fill(H)circle(1.4pt)node[above]{$H=K$};\fill(P)circle(1.4pt)node[right]{$P$};
\draw[axis](P)--++(-1.1,0)node[midway,above]{$\vct a$};
\begin{scope}[shift={(5.2,1.5)}]
\fill(0,0)circle(1.4pt)node[right]{$P$};
\draw[axis](0,0)--(-1.125,-.6495)node[below]{$a_r\vct e_r$};
\draw[axis](0,0)--(-.375,.6495)node[above]{$a_\theta\vct e_\theta$};
\draw[axis,thick](0,0)--(-1.5,0)node[left]{$\vct a$};
\draw[guide](-1.125,-.6495)--(-1.5,0)--(-.375,.6495);
\end{scope}

\end{tikzpicture}
\caption{左图确定纬线的曲率中心；右图取 $\theta_0=60^\circ$，按分量比例合成水平向心加速度。}\label{mp3:fig:l3-p8-s}
\end{figure}


图\ref{mp3:fig:l3-p8-s}中，两个球坐标分量合成后竖直分量相消。主法线是 $-\vct e_R$，不是 $-\vct e_r$。只有赤道 $\theta_0=\pi/2$ 时，纬线圆心与球心重合，才有 $\rho=a$。

几何上，纬线位于水平面$z=a\cos\theta_0$内，半径为$a\sin\theta_0$。其圆心位置为
\[
 \vct r_H=a\cos\theta_0\vct e_z,
 \qquad
 \vct r_P-\vct r_H=a\sin\theta_0\vct e_R.
\]
图\ref{mp3:fig:l3-p8-s}中的 $PH$ 与 $PO$ 分别是轨迹曲率半径和球面约束半径。纬线是一条平面圆，挠率为零；其他球面曲线则未必位于固定平面内。
\end{solution}

\par\medskip\noindent\begin{minipage}{\textwidth}
\begin{exercise}[球面轨迹的瞬时曲率与挠率（选做）]\label{mp3:exercise:l3-p9}

长为 $a$ 的杆一端固定在 $O$，杆所在竖直平面以角速度 $\Omega>0$ 绕 $z$ 轴转动，杆又在该平面内以角速度 $\omega>0$ 从水平位置向上抬起。端点 $P$ 的位置为
\[
\vct r_P=a(\cos\omega t\cos\Omega t,\ \cos\omega t\sin\Omega t,\ \sin\omega t).
\]
图\ref{mp3:fig:l3-p9-q}给出两个角的定义。求 $t=0$ 时的曲率半径和挠率，判断轨迹是否是球面上的大圆。

\end{exercise}
\begin{figure}[H]\centering
\begin{tikzpicture}[scale=.95]

\begin{scope}[x={(.9cm,-.2cm)},y={(.25cm,.65cm)},z={(0cm,.9cm)}]
\draw[axis](0,0,0)--(3.3,0,0)node[right]{$x$};
\draw[axis](0,0,0)--(0,2.8,0)node[right]{$y$};
\draw[axis](0,0,0)--(0,0,2.8)node[above]{$z$};
\coordinate(H)at(1.9,1.1,0);\coordinate(P)at(1.9,1.1,1.5);
\draw[guide](0,0,0)--(H)--(P);\draw[very thick](0,0,0)--(P);
\fill(P)circle(1.5pt)node[right]{$P$};\node[below left]at(0,0,0){$O$};
\draw[domain=0:30,samples=25]plot({.8*cos(\x)},{.8*sin(\x)},0);
\node at(1.1,.2,0){$\Omega t$};
\draw[domain=0:34.4,samples=30]plot({cos(30)*cos(\x)},{sin(30)*cos(\x)},{sin(\x)});\node[right]at(1.1,.64,.4){$\omega t$};
\node[left]at(.95,.55,.75){$a$};
\end{scope}

\end{tikzpicture}
\caption{方位角为 $\Omega t$，仰角为 $\omega t$。这两个角速度不能直接作为同向标量相加。}\label{mp3:fig:l3-p9-q}
\end{figure}
\end{minipage}\par\medskip
\begin{solution}
记 $S=\Omega^2+\omega^2$。对给定位置先作乘积求导：
\[
\begin{aligned}
\dot x&=-a[\omega\sin\omega t\cos\Omega t+\Omega\cos\omega t\sin\Omega t],\\
\dot y&=a[-\omega\sin\omega t\sin\Omega t+\Omega\cos\omega t\cos\Omega t],\\
\dot z&=a\omega\cos\omega t.
\end{aligned}
\]
再次求导，
\[
\begin{aligned}
\ddot x&=-aS\cos\omega t\cos\Omega t+2a\omega\Omega\sin\omega t\sin\Omega t,\\
\ddot y&=-aS\cos\omega t\sin\Omega t-2a\omega\Omega\sin\omega t\cos\Omega t,\\
\ddot z&=-a\omega^2\sin\omega t.
\end{aligned}
\]
在 $t=0$ 得
\[
\vct r=(a,0,0),\quad
\vct v=(0,a\Omega,a\omega),\quad
\vct a=(-aS,0,0).
\]
此时速度与加速度垂直，$v^2=a^2S$，故 $\rho=a$，曲率中心恰为 $O$。但这只是一个时刻的结果。

\begin{figure}[H]\centering
\begin{tikzpicture}[scale=.95]

\begin{scope}[x={(.9cm,-.2cm)},y={(.6cm,.3cm)},z={(0cm,1cm)}]
\coordinate(P)at(2,0,0);
\draw[guide](0,0,0)--(P);\fill(0,0,0)circle(1.4pt)node[left]{$O=K$};
\fill(P)circle(1.4pt)node[below]{$P(0)$};
\draw[axis](P)--++(-1,0,0)node[above]{$\vct e_n$};
\draw[axis](P)--++(0,.8,.6)node[above right]{$\vct e_t$};
\draw[axis](P)--++(0,-.6,.8)node[above left]{$\vct e_b$};
\end{scope}
\node at(5.2,1.6){$\vct e_t=(0,\Omega,\omega)/\sqrt S$};
\node at(5.2,.6){$\vct e_n=(-1,0,0)$};
\node at(5.2,-.4){$\vct e_b=(0,-\omega,\Omega)/\sqrt S$};

\end{tikzpicture}
\caption{$t=0$ 的自然标架，示意取 $\omega/\Omega=3/4$。副法线随时间转向，不能把当前密切平面延用为固定轨道平面。}\label{mp3:fig:l3-p9-s}
\end{figure}


图\ref{mp3:fig:l3-p9-s}给出该时刻的密切平面：它由 $\vct i$ 与 $\Omega\vct j+\omega\vct k$ 张成。对加速度再求导：$x(t)$ 为偶函数，故 $\dddot x(0)=0$；纵向两项分别给出 $\dddot y(0)=-aS\Omega-2a\omega^2\Omega$；竖向为 $\dddot z(0)=-a\omega^3$。所以三阶导数为
\[
\vct j=(0,-a\Omega(\Omega^2+3\omega^2),-a\omega^3).
\]
于是
\[
\vct v\times\vct a=a^2S(0,-\omega,\Omega),\qquad
|\vct v\times\vct a|^2=a^4S^3.
\]
数量三重积为
\[
(\vct v\times\vct a)\cdot\vct j
=a^3S\Omega\omega[(\Omega^2+3\omega^2)-\omega^2]
=a^3S\Omega\omega(\Omega^2+2\omega^2).
\]
将它除以$|\vct v\times\vct a|^2$，得到
\[
\tau(0)=
\frac{\Omega\omega(\Omega^2+2\omega^2)}
{a(\Omega^2+\omega^2)^2}>0.
\]
$\tau(0)>0$ 表明轨迹在 $P(0)$ 附近已离开当前密切平面，因此它不是大圆。此刻的 $\rho=a$ 只确定局部的二阶弯曲。

图\ref{mp3:fig:l3-p9-s}中的副法向量为$(0,-\omega,\Omega)/\sqrt S$，$\vct j$在该方向有非零投影，使密切平面开始转动。曲率半径只需要二阶运动信息，挠率还需要第三阶；两条曲线在某点具有相同的速度和加速度，仍可能有不同挠率。
\end{solution}

\par\medskip\noindent\begin{minipage}{\textwidth}
\begin{exercise}[同时平动和转动的观察坐标]\label{mp3:exercise:l3-p10}

惯性系中质点 $P$ 固定在 $(L,0)$，$L>0$。动系原点满足 $O(t)=(At^2/2,0)$，$A>0$，动系又以恒定角速度 $\Omega>0$ 逆时针转动，初始两系坐标轴平行，见图\ref{mp3:fig:l3-p10-q}。求 $P$ 的动系坐标，并用加速度转换公式核对 $t=0$ 时的相对加速度。原点经过质点时，是否出现物理奇点？

\end{exercise}
\begin{figure}[H]\centering
\begin{tikzpicture}[scale=.95]

\draw[axis](-.3,0)--(5,0)node[right]{$x$};\draw[axis](0,-.3)--(0,2.6)node[above]{$y$};
\fill(0,0)circle(1.4pt)node[below]{$O_0$};\fill(4,0)circle(1.8pt)node[above]{$P$};
\coordinate(O)at(1.3,0);\fill(O)circle(1.4pt)node[below]{$O(t)$};
\draw[axis](O)--++({2*cos(35)},{2*sin(35)})node[right]{$x'$};
\draw[axis](O)--++({-1.4*sin(35)},{1.4*cos(35)})node[above]{$y'$};
\draw[axis](1.3,-.6)--(2.5,-.6)node[right]{$\vct a_O=A\vct i$};
\draw[|-|](0,-1.1)--(4,-1.1)node[midway,fill=white]{$L$};
\draw(1.8,0)arc(0:35:.5);\node at(2.1,.24){$\Omega t$};

\end{tikzpicture}
\caption{质点在惯性系中静止，观察坐标系的原点沿水平线加速并同时转动。}\label{mp3:fig:l3-p10-q}
\end{figure}
\end{minipage}\par\medskip
\begin{solution}
记 $s=L-At^2/2$，$\tau=\Omega t$。相对位置在惯性系中的分量为 $(s,0)$，转至动系得
\[
x'=s\cos\tau,\qquad y'=-s\sin\tau.
\]
令 $\dot s=-At$、$\ddot s=-A$。对两坐标分别使用乘积法则，
\[
\begin{aligned}
\dot x'&=\dot s\cos\tau-\Omega s\sin\tau,\\
\dot y'&=-\dot s\sin\tau-\Omega s\cos\tau,\\
\ddot x'&=(\ddot s-\Omega^2s)\cos\tau-2\Omega\dot s\sin\tau,\\
\ddot y'&=-(\ddot s-\Omega^2s)\sin\tau-2\Omega\dot s\cos\tau.
\end{aligned}
\]
在 $t=0$ 代入 $s=L$、$\dot s=0$、$\tau=0$，得到
\[
\vct v_{\rm rel}(0)=-\Omega L\vct j,\qquad
\vct a_{\rm rel}(0)=-(A+\Omega^2L)\vct i.
\]
图\ref{mp3:fig:l3-p10-s}中，相对初速度向下，因而运动学科里奥利项
$2\vct\Omega\times\vct v_{\rm rel}=2\Omega^2L\vct i$ 向右。加速度转换式的四个非零项满足
\[
0=A\vct i+\vct a_{\rm rel}
+2\Omega^2L\vct i-\Omega^2L\vct i,
\]
与直接求导一致。

\begin{figure}[H]\centering
\begin{tikzpicture}[scale=.95]

\fill(0,0)circle(1.5pt)node[above]{$P$};
\draw[axis](0,0)--(0,-1.4)node[below]{$\vct v_{\rm rel}(0)$};
\draw[axis](0,.25)--(1.8,.25)node[right]{$2\Omega^2L\vct i$};
\draw[axis](0,.75)--(-1,.75)node[left]{$-\Omega^2L\vct i$};
\draw[axis](0,1.25)--(1,1.25)node[right]{$A\vct i$};
\draw[axis,thick](0,-.45)--(-2,-.45)node[left]{$\vct a_{\rm rel}(0)$};
\node at(4.7,.8){$\vct a_P=0$};
\node at(4.7,-.2){$\vct a_{\rm rel}=-(A+\Omega^2L)\vct i$};

\end{tikzpicture}
\caption{$t=0$ 时的各方向。不同高度的加速度箭头为便于辨认而平移，计算均以同一时刻的矢量相加。}\label{mp3:fig:l3-p10-s}
\end{figure}


动系原点在 $t_c=\sqrt{2L/A}$ 时经过 $P$，此时 $s=0$，所以 $x'=y'=0$。但
\[
\begin{pmatrix}\dot x'\\\dot y'\end{pmatrix}_{t_c}
=-At_c\begin{pmatrix}\cos\Omega t_c\\-\sin\Omega t_c\end{pmatrix}
\]
仍有限且非零。笛卡尔动坐标正常通过原点，不存在物理奇点；若改用相对极坐标，极角在此失去定义，问题来自坐标表示。$s<0$ 后上面的笛卡尔解照常有效，不能把带符号的 $s$ 误作非负的极径。

\end{solution}

\par\medskip\noindent\begin{minipage}{\textwidth}
\begin{exercise}[近地竖直下落的向东偏移]\label{mp3:exercise:l3-p11}

在北纬 $\lambda$ 处，从距地面高度 $H$ 的位置相对地面静止释放小球。以释放点地面的铅垂投影为原点，取东、北、上为 $x,y,z$ 正向，见图\ref{mp3:fig:l3-p11-q}。把地球视为以角速度 $\Omega$ 匀速转动，$g$ 表示当地有效重力加速度。忽略空气阻力、地球曲率及所有 $\Omega^2$ 量，求落地时的水平偏移。

\end{exercise}
\begin{figure}[H]\centering
\begin{tikzpicture}[scale=.95]

\begin{scope}[x={(1cm,0)},y={(.5cm,.35cm)},z={(0cm,1cm)}]
\draw[axis](0,0,0)--(2.5,0,0)node[right]{$x$（东）};
\draw[axis](0,0,0)--(0,2.3,0)node[right]{$y$（北）};
\draw[axis](0,0,0)--(0,0,3.2)node[above]{$z$（上）};
\fill(0,0,2.3)circle(2pt)node[above left]{$P_0$};
\draw[guide](0,0,0)--(0,0,2.3);
\draw[axis](0,0,2.3)--(0,0,1.4)node[left]{$g$};
\draw[axis,thick](0,0,0)--(0,1.3,1.5)node[right]{$\vct\Omega$};
\draw[guide](0,0,0)--(-.5,0,0);
\draw[guide](0,0,2.3)--(-.5,0,2.3);
\draw[|-|](-.5,0,0)--(-.5,0,2.3)node[midway,left]{$H$};
\draw[domain=0:49.086,samples=35]plot(0,{1.15*cos(\x)},{1.15*sin(\x)});
\node[right]at(0,.88,.47){$\lambda$};
\end{scope}
\node at(5,2.5){$\vct\Omega=(0,\Omega\cos\lambda,\Omega\sin\lambda)$};

\end{tikzpicture}
\caption{当地随地球转动的坐标系；地轴方向在当地坐标中分成北向与上向。}\label{mp3:fig:l3-p11-q}
\end{figure}
\end{minipage}\par\medskip
\begin{solution}
在东、北、上的右手坐标中，地球角速度为 $\vct\Omega=(0,\Omega\cos\lambda,\Omega\sin\lambda)$。零阶运动满足 $\ddot z=-g$、$z(0)=H$、$\dot z(0)=0$，两次积分得
\[
\dot z=\int_0^t(-g)\dd u=[-gu]_0^t=-gt,
\qquad
z=H+\int_0^t(-gu)\dd u=H-\frac12gt^2.
\]
因此 $\vct v^{(0)}=-gt\vct e_z$。角速度的竖直分量与它平行，叉积为零，只有北向分量 $\Omega\cos\lambda$ 产生一阶科里奥利加速度。计算时将零阶速度代入，得到
\[
-2\vct\Omega\times\vct v^{(0)}
=2\Omega g t\cos\lambda\,\vct e_x.
\]
因此
\[
\begin{aligned}
\ddot x&=2\Omega gt\cos\lambda,\\
\dot x&=2\Omega g\cos\lambda\int_0^t u\dd u
=2\Omega g\cos\lambda\left[\frac{u^2}{2}\right]_0^t
=\Omega gt^2\cos\lambda,\\
x&=\Omega g\cos\lambda\int_0^t u^2\dd u
=\Omega g\cos\lambda\left[\frac{u^3}{3}\right]_0^t
=\frac13\Omega gt^3\cos\lambda.
\end{aligned}
\]
图\ref{mp3:fig:l3-p11-s}中，落地偏移位于释放点铅垂线的东侧。到达地面的零阶时间为 $T=\sqrt{2H/g}$，故
\[
\Delta x=\frac{2\sqrt2}{3}
\frac{\Omega H^{3/2}\cos\lambda}{\sqrt g},\qquad
\Delta y=0\quad\text{（一阶近似）}.
\]

\begin{figure}[H]\centering
\begin{tikzpicture}[scale=.95]

\fill(0,0)circle(1.5pt)node[above]{$P$};
\draw[axis](0,0)--(0,-1.4)node[below]{$\vct v^{(0)}$};
\draw[axis](0,0)--(1.6,0)node[right]{$-2\vct\Omega\times\vct v^{(0)}$};
\node at(.4,1){向东};
\begin{scope}[shift={(5.6,0)}]
\draw[axis](-.2,-1.7)--(2.4,-1.7)node[right]{$x$};
\draw[axis](0,-1.7)--(0,1.3)node[above]{$z$};
\draw[guide](0,-1.7)--(0,.9);
\draw[curve,domain=0:1,samples=80]plot({1.5*\x*\x*\x},{.9-2.6*\x*\x});
\fill(0,.9)circle(1.5pt)node[left]{$P_0$};
\fill(1.5,-1.7)circle(1.5pt)node[above right]{落点};
\draw[|-|](0,-2.2)--(1.5,-2.2)node[midway,fill=white]{$\Delta x$};
\end{scope}

\end{tikzpicture}
\caption{左图确定科里奥利加速度的东向分量；右图用参数 $t/T$ 画出一阶落体轨迹，横向尺度夸大。}\label{mp3:fig:l3-p11-s}
\end{figure}


横向速度本身是一阶小量，再代入科里奥利项只产生二阶修正，故这一阶没有北向偏移。使用当地有效 $g$ 后，不再重复加入静态离心加速度。计算还要求 $H$ 远小于地球半径、$\Omega T\ll1$；在两极 $\cos\lambda=0$，一阶向东偏移为零。右图横向夸大绘制，不代表偏移与下落高度同量级。

图\ref{mp3:fig:l3-p11-s}中的偏移也可按下降距离 $\zeta=H-z$ 表示。用零阶下落关系消去时间，得到
\[
 x=\frac{2\sqrt2}{3}\frac{\Omega\cos\lambda}{\sqrt g}\,\zeta^{3/2}.
\]
释放点附近$\dd x/\dd\zeta\to0$，轨迹从竖直切线开始偏向东。近似的小量参数为$\Omega\sqrt{2H/g}$。在$x(t)$中代入零阶落地时间，误差属于本题舍去的高阶项。
\end{solution}

\par\medskip\noindent\begin{minipage}{\textwidth}
\begin{exercise}[转动坐标中的各向同性振子（选做）]\label{mp3:exercise:l3-p12}

水平光滑面上质点受到指向原点的线性回复力
$\vct F=-m\omega_0^2\vct r$，$\omega_0>0$。观察坐标系绕同一原点以角速度 $\Omega>0$ 匀速转动，见图\ref{mp3:fig:l3-p12-q}。推导动系运动方程及两个圆运动模态，判断 $\Omega>\omega_0$ 时是否失稳。再求惯性系初值 $\vct r(0)=(A,0)$、$\vct v(0)=0$ 对应的动系轨迹。

\end{exercise}
\begin{figure}[H]\centering
\begin{tikzpicture}[scale=.95]

\draw[axis](-.3,0)--(4.2,0)node[right]{$x$};
\draw[axis](0,-.3)--(0,2.5)node[above]{$y$};
\fill(0,0)circle(1.5pt)node[below left]{$O$};
\fill(2.4,1)circle(2pt)node[right]{$P$};
\draw[guide](0,0)--(2.4,1);\draw[axis,thick](2.4,1)--(1.2,.5)node[below right]{$-m\omega_0^2\vct r$};
\draw[axis,guide](0,0)--({2.6*cos(50)},{2.6*sin(50)})node[above]{$x'$};
\draw[axis,guide](0,0)--({-1.6*sin(50)},{1.6*cos(50)})node[above]{$y'$};
\draw(.65,0)arc(0:50:.65);\node at(.9,.45){$\Omega t$};

\end{tikzpicture}
\caption{回复力各向同性，观察坐标的旋转不改变其形式。}\label{mp3:fig:l3-p12-q}
\end{figure}
\end{minipage}\par\medskip
\begin{solution}
由式\eqref{mp3:eq:l3-moving-newton}，
\[
\ddot x'-2\Omega\dot y'+(\omega_0^2-\Omega^2)x'=0,\qquad
\ddot y'+2\Omega\dot x'+(\omega_0^2-\Omega^2)y'=0.
\]
令 $z'=x'+\mathrm i y'$，两式合为
\[
\ddot z'+2\mathrm i\Omega\dot z'+(\omega_0^2-\Omega^2)z'=0.
\]
代入 $z'\propto\ee^{\lambda t}$，特征方程配方为
\[
\lambda^2+2\mathrm i\Omega\lambda+\omega_0^2-\Omega^2=0,
\qquad (\lambda+\mathrm i\Omega)^2=-\omega_0^2.
\]
两根为 $\lambda=-\mathrm i\Omega\pm\mathrm i\omega_0$，故通解为
\[
z'=C_+\ee^{\mathrm i(\omega_0-\Omega)t}
+C_-\ee^{-\mathrm i(\omega_0+\Omega)t}.
\]
两个模态在惯性系中分别以角速度 $+\omega_0,-\omega_0$ 转动，换到动系后均减去 $\Omega$。图\ref{mp3:fig:l3-p12-s}左图表示 $\Omega<\omega_0$ 时的两个相反转向；当 $\Omega>\omega_0$，两模态在动系中都顺时针旋转，幅度仍不增长。

\begin{figure}[H]\centering
\begin{tikzpicture}[scale=.95]

\draw[curve](-1.25,0)circle(.85);\draw[curve](1.25,0)circle(.85);
\draw[axis](-.4,0)arc(0:70:.85);
\draw[axis](2.1,0)arc(0:-70:.85);
\node at(-1.25,-1.35){$\omega_0-\Omega$};
\node at(1.25,-1.35){$-(\omega_0+\Omega)$};
\node at(0,1.4){$\Omega<\omega_0$};
\begin{scope}[shift={(5.3,0)}]
\draw[axis](-1.6,0)--(1.8,0)node[right]{$x'/A$};
\draw[axis](0,-1.6)--(0,1.8)node[above]{$y'/A$};
\draw[curve,domain=0:360,samples=250]plot({1.4*cos(2*\x)*cos(\x)},{-1.4*cos(2*\x)*sin(\x)});
\fill(1.4,0)circle(1.5pt)node[above right]{$P_0$};
\node[below left]at(0,0){$O$};
\end{scope}

\end{tikzpicture}
\caption{左图将两个圆运动模态分开示意；右图由同一初值按动系参数方程作图，取 $\Omega=\omega_0/2$。}\label{mp3:fig:l3-p12-s}
\end{figure}


特征根始终为纯虚数，质点运动有界。$\omega_0^2-\Omega^2$ 与科里奥利耦合项共同决定特征根。它们组合后给出两个实的转动频率，系数本身的正负并不决定振幅是否增长。$\Omega=\omega_0$ 时一个模态在动系中静止，也不表示真实回复力消失。

初始两系坐标重合，$z'(0)=A$；但动系初速度须扣除坐标转动，$\dot z'(0)=\dot z(0)-\mathrm i\Omega z(0)=-\mathrm i\Omega A$。代入通解，得到
\[
C_++C_-=A,\qquad
\mathrm i(\omega_0-\Omega)C_+-\mathrm i(\omega_0+\Omega)C_-=-\mathrm i\Omega A.
\]
第二式除以 $\mathrm i$，再加上第一式的 $\Omega$ 倍，得到 $\omega_0(C_+-C_-)=0$。因 $\omega_0>0$，故 $C_+=C_-=A/2$。于是
\[
z'=\frac A2\ee^{-\mathrm i\Omega t}(\ee^{\mathrm i\omega_0t}+\ee^{-\mathrm i\omega_0t})
=A\cos\omega_0t\,\ee^{-\mathrm i\Omega t}.
\]
取实部、虚部，
\[
x'=A\cos\omega_0t\cos\Omega t,\qquad
y'=-A\cos\omega_0t\sin\Omega t.
\]
转回惯性系即为 $x=A\cos\omega_0t$、$y=0$。初始相对速度为 $(0,-\Omega A)$，并非零。右图取 $\Omega=\omega_0/2$，在 $0\le t\le4\pi/\omega_0$ 内画出闭合的四瓣轨迹。它在 $O$ 多次穿过，但笛卡尔坐标和速度均连续。

\end{solution}

\chapter{相图的应用}

相图以位置和速度确定一个运动状态。相轨线给出转向位置、势阱内的运动范围，以及越过势垒所需的初态条件。本讲以单自由度系统为主，由能量构造相轨线，进而讨论平衡点、分界运动和耗散运动。


\section{能量与相轨线}
相轨线同时记录位置与速度，曲线上的每一点仍须满足运动方程。能量限定可达区域，初态及时间方向选定轨线的实际分支。

\subsection{相点、方向场与转向点}
设质点在直线上运动，质量为常量 $m$，所受力只依赖位置：
\begin{equation}\label{mp4:eq:l4-autonomous}
\dot x=v,\qquad \dot v=\frac{F(x)}m.
\end{equation}
相平面横轴取 $x$，纵轴取 $v$。某一时刻的状态对应一个相点 $(x,v)$，同一次运动中相点经过的曲线称为相轨线。在力足够光滑、初值解唯一的区域，两条不同的相轨线不能相交；若经过同一点，它们属于同一个解，只相差计时起点。

相点的运动方向由 $(\dot x,\dot v)$ 决定。上半平面 $v>0$，相点向右；下半平面 $v<0$，相点向左。横轴上 $v=0$，若 $F(x)\ne0$，相点沿竖直方向穿过横轴，质点在这里转向。只有同时满足 $v=0$、$F(x)=0$ 的点才是平衡点。


\begin{figure}[H]\centering
\begin{tikzpicture}
\draw[axis](-2.5,0)--(2.5,0)node[right]{$x$};
\draw[axis](0,-1.8)--(0,1.8)node[above]{$v$};
\draw[flow] plot coordinates{(2.00000,-0.00000) (1.99846,-0.04711) (1.99383,-0.09415) (1.98614,-0.14104) (1.97538,-0.18772) (1.96157,-0.23411) (1.94474,-0.28013) (1.92491,-0.32573) (1.90211,-0.37082) (1.87638,-0.41534) (1.84776,-0.45922) (1.81629,-0.50239) (1.78201,-0.54479) (1.74499,-0.58635) (1.70528,-0.62700) (1.66294,-0.66668) (1.61803,-0.70534) (1.57063,-0.74291) (1.52081,-0.77934) (1.46865,-0.81456) (1.41421,-0.84853) (1.35760,-0.88119) (1.29890,-0.91249) (1.23819,-0.94238) (1.17557,-0.97082) (1.11114,-0.99776) (1.04500,-1.02317) (0.97724,-1.04700) (0.90798,-1.06921) (0.83732,-1.08977) (0.76537,-1.10866) (0.69223,-1.12583) (0.61803,-1.14127) (0.54288,-1.15495) (0.46689,-1.16684) (0.39018,-1.17694) (0.31287,-1.18523) (0.23507,-1.19168) (0.15692,-1.19630) (0.07852,-1.19907) (0.00000,-1.20000) (-0.07852,-1.19907) (-0.15692,-1.19630) (-0.23507,-1.19168) (-0.31287,-1.18523) (-0.39018,-1.17694) (-0.46689,-1.16684) (-0.54288,-1.15495) (-0.61803,-1.14127) (-0.69223,-1.12583) (-0.76537,-1.10866) (-0.83732,-1.08977) (-0.90798,-1.06921) (-0.97724,-1.04700) (-1.04500,-1.02317) (-1.11114,-0.99776) (-1.17557,-0.97082) (-1.23819,-0.94238) (-1.29890,-0.91249) (-1.35760,-0.88119) (-1.41421,-0.84853) (-1.46865,-0.81456) (-1.52081,-0.77934) (-1.57063,-0.74291) (-1.61803,-0.70534) (-1.66294,-0.66668) (-1.70528,-0.62700) (-1.74499,-0.58635) (-1.78201,-0.54479) (-1.81629,-0.50239) (-1.84776,-0.45922) (-1.87638,-0.41534) (-1.90211,-0.37082) (-1.92491,-0.32573) (-1.94474,-0.28013) (-1.96157,-0.23411) (-1.97538,-0.18772) (-1.98614,-0.14104) (-1.99383,-0.09415) (-1.99846,-0.04711) (-2.00000,-0.00000) (-1.99846,0.04711) (-1.99383,0.09415) (-1.98614,0.14104) (-1.97538,0.18772) (-1.96157,0.23411) (-1.94474,0.28013) (-1.92491,0.32573) (-1.90211,0.37082) (-1.87638,0.41534) (-1.84776,0.45922) (-1.81629,0.50239) (-1.78201,0.54479) (-1.74499,0.58635) (-1.70528,0.62700) (-1.66294,0.66668) (-1.61803,0.70534) (-1.57063,0.74291) (-1.52081,0.77934) (-1.46865,0.81456) (-1.41421,0.84853) (-1.35760,0.88119) (-1.29890,0.91249) (-1.23819,0.94238) (-1.17557,0.97082) (-1.11114,0.99776) (-1.04500,1.02317) (-0.97724,1.04700) (-0.90798,1.06921) (-0.83732,1.08977) (-0.76537,1.10866) (-0.69223,1.12583) (-0.61803,1.14127) (-0.54288,1.15495) (-0.46689,1.16684) (-0.39018,1.17694) (-0.31287,1.18523) (-0.23507,1.19168) (-0.15692,1.19630) (-0.07852,1.19907) (-0.00000,1.20000) (0.07852,1.19907) (0.15692,1.19630) (0.23507,1.19168) (0.31287,1.18523) (0.39018,1.17694) (0.46689,1.16684) (0.54288,1.15495) (0.61803,1.14127) (0.69223,1.12583) (0.76537,1.10866) (0.83732,1.08977) (0.90798,1.06921) (0.97724,1.04700) (1.04500,1.02317) (1.11114,0.99776) (1.17557,0.97082) (1.23819,0.94238) (1.29890,0.91249) (1.35760,0.88119) (1.41421,0.84853) (1.46865,0.81456) (1.52081,0.77934) (1.57063,0.74291) (1.61803,0.70534) (1.66294,0.66668) (1.70528,0.62700) (1.74499,0.58635) (1.78201,0.54479) (1.81629,0.50239) (1.84776,0.45922) (1.87638,0.41534) (1.90211,0.37082) (1.92491,0.32573) (1.94474,0.28013) (1.96157,0.23411) (1.97538,0.18772) (1.98614,0.14104) (1.99383,0.09415) (1.99846,0.04711) (2.00000,0.00000)};
\fill(2,0)circle(1.2pt)node[below right]{$B$};
\fill(-2,0)circle(1.2pt)node[below left]{$A$};
\fill(0,0)circle(1.2pt)node[below left]{$O$};
\draw[axis](2,0)--(2,-.65)node[right]{$\dot v<0$};\draw[axis](-2,0)--(-2,.65)node[left]{$\dot v>0$};\node at(1.3,1.55){$v>0$：向右};\node at(-1.3,-1.55){$v<0$：向左};
\end{tikzpicture}
\caption{简谐振子的相轨线方向。A、B 处速度为零，但相点仍沿竖直方向运动；O 是平衡点。}\label{mp4:fig:l4-field}
\end{figure}


在 $v\ne0$ 处消去时间，有
\[
\frac{\dd v}{\dd x}=\frac{F(x)}{mv}.
\]
转向点处的斜率可能发散，此时相轨线有竖直切线，加速度仍为有限的 $F(x)/m$。相图两轴通常具有不同量纲，线段的几何倾角依赖所选比例尺；方程给出的导数也带有相应量纲。

\subsection{由势能构造相轨线}
若 $F=-U'(x)$，将 $m\ddot x=-U'(x)$ 乘以 $\dot x=v$，左、右两项分别是 $mv\dot v=\dd(mv^2/2)/\dd t$ 和 $-U'(x)\dot x=-\dd U/\dd t$，因此
\begin{equation}\label{mp4:eq:l4-energy}
\frac{\dd}{\dd t}\left(\frac12mv^2+U(x)\right)=0,
\qquad E=\frac12mv^2+U(x).
\end{equation}
给定能量后，先在势能图上确定 $U(x)\le E$ 的可达区间，再画
\begin{equation}\label{mp4:eq:l4-energy-curve}
v=\pm\sqrt{\frac{2[E-U(x)]}{m}}.
\end{equation}
正号支向右，负号支向左。若可达区间分成互不连通的两部分，同一能量就对应两组分离的运动，质点分别被限制在各自的区间内。

设 $x_-$、$x_+$ 是两个普通转向点，二者之间 $E>U(x)$。记速度正、负支为 $v_+,v_-$，去程和回程的时间分别为
\[
T=\int_{x_-}^{x_+}\frac{\dd x}{v_+}
+\int_{x_+}^{x_-}\frac{\dd x}{v_-}
=2\sqrt{\frac m2}\int_{x_-}^{x_+}\frac{\dd x}{\sqrt{E-U(x)}}.
\]
合并常数，周期为
\begin{equation}\label{mp4:eq:l4-period}
T=\sqrt{2m}\int_{x_-}^{x_+}\frac{\dd x}{\sqrt{E-U(x)}}.
\end{equation}
这里已经计入去、回两个方向。在普通转向点附近，若 $U'(x_+)\ne0$，被积函数约为常数乘 $(x_+-x)^{-1/2}$，积分有限；若转向位置同时是势能极大点，积分性质会改变。

相轨线本身不标出时间。沿着一小段轨线，$\dd t=\dd x/v$；接近 $v=0$ 的位置是否耗时很长，须由该比值积分判断。相图中的一条闭合能量曲线，只有不经过平衡点并且绕行时间有限，才对应周期运动。

\par\medskip\noindent\begin{minipage}{\textwidth}
\begin{example}[相椭圆与到达时间]\label{mp4:example:l4-e1}
质量为 $m$ 的物体在光滑水平面上与劲度系数为 $k$ 的弹簧相连，$x$ 从平衡位置起算，$\omega=\sqrt{k/m}$。初态为 $x_0=a/2$、$v_0=\sqrt3\,\omega a/2>0$，见图\ref{mp4:fig:l4-e1-q}。求振幅、首次到达右端点的时间，以及首次回到 $x=a/2$ 时的速度和时间。
\end{example}

\begin{figure}[H]\centering
\begin{tikzpicture}
\draw[curve](-2.4,-.35)--(-2.4,.65);\draw(-2.4,.1)--(-2.15,.1)--(-1.97,.35)--(-1.7,-.15)--(-1.43,.35)--(-1.16,-.15)--(-.89,.35)--(-.62,.1)--(-.3,.1);
\draw[fill=white](-.3,-.25)rectangle(.35,.45);\node at(.03,.1){$m$};\draw(-2.6,-.3)--(2.3,-.3);
\draw[axis](-1.1,-.8)--(2.5,-.8)node[right]{$x$};\draw[guide](-1,-.3)--(-1,-1.35);\node[below left]at(-1,-.8){$0$};\node[above]at(-1.4,.55){$k$};
\draw[axis](.35,.15)--(1.6,.15)node[above]{$v_0$};
\draw[guide](.025,-.25)--(.025,-1.35);
\draw[<->,>=Stealth,line width=.35pt](-1,-1.2)--(.025,-1.2)node[midway,below]{$x_0=a/2$};
\end{tikzpicture}
\caption{弹簧振子的初始状态。箭头表示初速度，$x$ 从平衡位置起算。}\label{mp4:fig:l4-e1-q}
\end{figure}
\end{minipage}\par\medskip
\begin{solution}
弹簧的弹力为 $-kx$，运动方程是 $\ddot x+\omega^2x=0$。初态的动能和弹性势能分别为 $3ka^2/8$、$ka^2/8$，所以
\[
E=\frac12mv_0^2+\frac12kx_0^2=\frac12ka^2.
\]
在转向点速度为零，$E=kx^2/2$，故两个转向位置为 $x=\pm a$，振幅为 $a$。

消去能量，得到 $v^2=\omega^2(a^2-x^2)$。取 $q=x/a$、$w=v/(\omega a)$，相轨线为 $q^2+w^2=1$。图\ref{mp4:fig:l4-e1-s}中的初态为 $P=(1/2,\sqrt3/2)$。因为 $v_0>0$，相点先沿上支向右到达 $B=(1,0)$；在 B 处 $\dot v=-\omega^2a<0$，相点穿过横轴进入下半平面，再沿下支向左到达 $Q=(1/2,-\sqrt3/2)$。

\begin{figure}[H]\centering
\begin{tikzpicture}
\draw[axis](-2.5,0)--(2.5,0)node[right]{$q$};
\draw[axis](0,-2.2)--(0,2.2)node[above]{$w$};
\draw[flow] plot coordinates{(1.80000,-0.00000) (1.79861,-0.07067) (1.79445,-0.14123) (1.78752,-0.21157) (1.77784,-0.28158) (1.76541,-0.35116) (1.75027,-0.42020) (1.73242,-0.48859) (1.71190,-0.55623) (1.68874,-0.62301) (1.66298,-0.68883) (1.63466,-0.75359) (1.60381,-0.81718) (1.57049,-0.87952) (1.53475,-0.94050) (1.49665,-1.00003) (1.45623,-1.05801) (1.41357,-1.11437) (1.36873,-1.16901) (1.32178,-1.22184) (1.27279,-1.27279) (1.22184,-1.32178) (1.16901,-1.36873) (1.11437,-1.41357) (1.05801,-1.45623) (1.00003,-1.49665) (0.94050,-1.53475) (0.87952,-1.57049) (0.81718,-1.60381) (0.75359,-1.63466) (0.68883,-1.66298) (0.62301,-1.68874) (0.55623,-1.71190) (0.48859,-1.73242) (0.42020,-1.75027) (0.35116,-1.76541) (0.28158,-1.77784) (0.21157,-1.78752) (0.14123,-1.79445) (0.07067,-1.79861) (0.00000,-1.80000) (-0.07067,-1.79861) (-0.14123,-1.79445) (-0.21157,-1.78752) (-0.28158,-1.77784) (-0.35116,-1.76541) (-0.42020,-1.75027) (-0.48859,-1.73242) (-0.55623,-1.71190) (-0.62301,-1.68874) (-0.68883,-1.66298) (-0.75359,-1.63466) (-0.81718,-1.60381) (-0.87952,-1.57049) (-0.94050,-1.53475) (-1.00003,-1.49665) (-1.05801,-1.45623) (-1.11437,-1.41357) (-1.16901,-1.36873) (-1.22184,-1.32178) (-1.27279,-1.27279) (-1.32178,-1.22184) (-1.36873,-1.16901) (-1.41357,-1.11437) (-1.45623,-1.05801) (-1.49665,-1.00003) (-1.53475,-0.94050) (-1.57049,-0.87952) (-1.60381,-0.81718) (-1.63466,-0.75359) (-1.66298,-0.68883) (-1.68874,-0.62301) (-1.71190,-0.55623) (-1.73242,-0.48859) (-1.75027,-0.42020) (-1.76541,-0.35116) (-1.77784,-0.28158) (-1.78752,-0.21157) (-1.79445,-0.14123) (-1.79861,-0.07067) (-1.80000,-0.00000) (-1.79861,0.07067) (-1.79445,0.14123) (-1.78752,0.21157) (-1.77784,0.28158) (-1.76541,0.35116) (-1.75027,0.42020) (-1.73242,0.48859) (-1.71190,0.55623) (-1.68874,0.62301) (-1.66298,0.68883) (-1.63466,0.75359) (-1.60381,0.81718) (-1.57049,0.87952) (-1.53475,0.94050) (-1.49665,1.00003) (-1.45623,1.05801) (-1.41357,1.11437) (-1.36873,1.16901) (-1.32178,1.22184) (-1.27279,1.27279) (-1.22184,1.32178) (-1.16901,1.36873) (-1.11437,1.41357) (-1.05801,1.45623) (-1.00003,1.49665) (-0.94050,1.53475) (-0.87952,1.57049) (-0.81718,1.60381) (-0.75359,1.63466) (-0.68883,1.66298) (-0.62301,1.68874) (-0.55623,1.71190) (-0.48859,1.73242) (-0.42020,1.75027) (-0.35116,1.76541) (-0.28158,1.77784) (-0.21157,1.78752) (-0.14123,1.79445) (-0.07067,1.79861) (-0.00000,1.80000) (0.07067,1.79861) (0.14123,1.79445) (0.21157,1.78752) (0.28158,1.77784) (0.35116,1.76541) (0.42020,1.75027) (0.48859,1.73242) (0.55623,1.71190) (0.62301,1.68874) (0.68883,1.66298) (0.75359,1.63466) (0.81718,1.60381) (0.87952,1.57049) (0.94050,1.53475) (1.00003,1.49665) (1.05801,1.45623) (1.11437,1.41357) (1.16901,1.36873) (1.22184,1.32178) (1.27279,1.27279) (1.32178,1.22184) (1.36873,1.16901) (1.41357,1.11437) (1.45623,1.05801) (1.49665,1.00003) (1.53475,0.94050) (1.57049,0.87952) (1.60381,0.81718) (1.63466,0.75359) (1.66298,0.68883) (1.68874,0.62301) (1.71190,0.55623) (1.73242,0.48859) (1.75027,0.42020) (1.76541,0.35116) (1.77784,0.28158) (1.78752,0.21157) (1.79445,0.14123) (1.79861,0.07067) (1.80000,0.00000)};
\fill(0.9,1.5588457268119895)circle(1.2pt)node[above right]{$P$};
\fill(0.9,-1.5588457268119895)circle(1.2pt)node[below right]{$Q$};
\fill(1.8,0)circle(1.2pt)node[below right]{$B$};
\draw[guide](.9,-1.56)--(.9,1.56);\node[below left]at(.9,0){$1/2$};\node[below left]at(-1.8,0){$-1$};\node[above right]at(0,1.8){$1$};
\end{tikzpicture}
\caption{归一化相椭圆。P、Q 分别是向右与向左经过同一位置的状态；B 是右端转向点。}\label{mp4:fig:l4-e1-s}
\end{figure}

计算首次到达时间时，须按这段运动的实际方向选取相位。令
\[
x=a\sin\varphi(t),\qquad v=\omega a\cos\varphi(t),
\qquad \varphi(t)=\omega t+\varphi_0.
\]
由 $x_0=a/2$ 和 $v_0>0$ 同时确定 $\varphi_0=\pi/6$，而不是另一个具有相同正弦值的角 $5\pi/6$。相位的匀速变化也可由 $\dot x=a\cos\varphi\dot\varphi=v$ 得出：在转向点以外 $\dot\varphi=\omega$，连续延伸到转向点后，$\varphi(t)=\varphi_0+\int_0^t\omega\dd u=\varphi_0+\omega t$。到达 B 时相位首次增至 $\pi/2$，因此
\[
t_B=\frac{\pi/2-\pi/6}{\omega}=\frac{\pi}{3\omega}.
\]
首次返回 $x=a/2$ 对应图中的 Q，此时相位为 $5\pi/6$，所以
\[
t_Q=\frac{5\pi/6-\pi/6}{\omega}=\frac{2\pi}{3\omega},
\qquad v_Q=-\frac{\sqrt3}{2}\omega a.
\]
这里 $t_Q$ 从题设的初始时刻起算，B 到 Q 的时间则为 $\pi/(3\omega)$。

也可直接沿上支积分求第一段时间：
\[
t_B=\int_{a/2}^{a}\frac{\dd x}{\omega\sqrt{a^2-x^2}}
=\frac1\omega\left[\arcsin\frac xa\right]_{a/2}^{a}.
\]
返回段中 $v<0$、$\dd x<0$，二者之比仍给出正的时间。图中 P、Q 的横坐标相同、纵坐标相反，分别表示向右和向左经过同一位置。能量确定相轨线，初速度与箭头则确定沿轨线运动的先后次序。
\end{solution}

\par\medskip\noindent\begin{minipage}{\textwidth}
\begin{example}[四次势阱的周期]\label{mp4:example:l4-e2}
质量为$m$的质点在 $U(x)=\kappa x^4/4$ 中运动，$\kappa>0$，最大位移为 $A$。图\ref{mp4:fig:l4-e2-q}给出势阱。求相轨线、最大速率及周期随振幅的变化，并说明原点是否稳定。
\end{example}

\begin{figure}[H]\centering
\begin{tikzpicture}
\draw[axis](-2.3,0)--(2.3,0)node[right]{$x$};
\draw[axis](0,-0.25)--(0,2.4)node[above]{$U$};
\draw[curve,domain=-1.45:1.45,samples=100]plot(\x,{.45*(\x)^4});\draw[guide](-1.8,1.2)--(1.8,1.2)node[right]{$E$};\draw[guide](-1.278,0)--(-1.278,1.2);\draw[guide](1.278,0)--(1.278,1.2);\node[below]at(-1.278,0){$-A$};\node[below]at(1.278,0){$A$};
\end{tikzpicture}
\caption{四次势阱及能量直线。能量与势能的交点为两个转向位置。}\label{mp4:fig:l4-e2-q}
\end{figure}
\end{minipage}\par\medskip
\begin{solution}
图\ref{mp4:fig:l4-e2-q}中能量直线与势能曲线交于 $x=\pm A$。转向点速度为零，所以 $E=\kappa A^4/4$。由能量守恒，
\[
\frac12mv^2+\frac\kappa4x^4=\frac\kappa4A^4,
\qquad v^2=\frac\kappa{2m}(A^4-x^4).
\]
允许区间为 $[-A,A]$。速度正、负两支分别表示向右和向左的运动；经过原点时势能最小，速率达到
\[
v_{\max}=A^2\sqrt{\frac\kappa{2m}}\equiv v_*.
\]
因此令 $q=x/A$、$w=v/v_*$，得到 $w^2+q^4=1$，见图\ref{mp4:fig:l4-e2-s}。

\begin{figure}[H]\centering
\begin{tikzpicture}
\draw[axis](-2.3,0)--(2.3,0)node[right]{$q$};
\draw[axis](0,-1.8)--(0,1.8)node[above]{$w$};
\draw[flow] plot coordinates{(1.70000,-0.00000) (1.69942,-0.04442) (1.69767,-0.08876) (1.69476,-0.13294) (1.69069,-0.17691) (1.68546,-0.22056) (1.67907,-0.26385) (1.67153,-0.30669) (1.66285,-0.34900) (1.65303,-0.39074) (1.64207,-0.43181) (1.62999,-0.47217) (1.61680,-0.51175) (1.60249,-0.55048) (1.58709,-0.58832) (1.57060,-0.62521) (1.55303,-0.66109) (1.53439,-0.69593) (1.51471,-0.72967) (1.49399,-0.76228) (1.47224,-0.79373) (1.44949,-0.82397) (1.42574,-0.85299) (1.40101,-0.88076) (1.37533,-0.90727) (1.34870,-0.93249) (1.32115,-0.95642) (1.29269,-0.97906) (1.26335,-1.00040) (1.23314,-1.02046) (1.20208,-1.03923) (1.17020,-1.05674) (1.13752,-1.07300) (1.10406,-1.08804) (1.06984,-1.10188) (1.03489,-1.11456) (0.99923,-1.12611) (0.96289,-1.13657) (0.92589,-1.14599) (0.88825,-1.15442) (0.85000,-1.16190) (0.81117,-1.16848) (0.77178,-1.17424) (0.73187,-1.17921) (0.69145,-1.18346) (0.65056,-1.18706) (0.60923,-1.19006) (0.56747,-1.19253) (0.52533,-1.19452) (0.48283,-1.19609) (0.43999,-1.19730) (0.39686,-1.19822) (0.35345,-1.19888) (0.30980,-1.19934) (0.26594,-1.19964) (0.22189,-1.19983) (0.17770,-1.19993) (0.13338,-1.19998) (0.08897,-1.20000) (0.04450,-1.20000) (0.00000,-1.20000) (-0.04450,-1.20000) (-0.08897,-1.20000) (-0.13338,-1.19998) (-0.17770,-1.19993) (-0.22189,-1.19983) (-0.26594,-1.19964) (-0.30980,-1.19934) (-0.35345,-1.19888) (-0.39686,-1.19822) (-0.43999,-1.19730) (-0.48283,-1.19609) (-0.52533,-1.19452) (-0.56747,-1.19253) (-0.60923,-1.19006) (-0.65056,-1.18706) (-0.69145,-1.18346) (-0.73187,-1.17921) (-0.77178,-1.17424) (-0.81117,-1.16848) (-0.85000,-1.16190) (-0.88825,-1.15442) (-0.92589,-1.14599) (-0.96289,-1.13657) (-0.99923,-1.12611) (-1.03489,-1.11456) (-1.06984,-1.10188) (-1.10406,-1.08804) (-1.13752,-1.07300) (-1.17020,-1.05674) (-1.20208,-1.03923) (-1.23314,-1.02046) (-1.26335,-1.00040) (-1.29269,-0.97906) (-1.32115,-0.95642) (-1.34870,-0.93249) (-1.37533,-0.90727) (-1.40101,-0.88076) (-1.42574,-0.85299) (-1.44949,-0.82397) (-1.47224,-0.79373) (-1.49399,-0.76228) (-1.51471,-0.72967) (-1.53439,-0.69593) (-1.55303,-0.66109) (-1.57060,-0.62521) (-1.58709,-0.58832) (-1.60249,-0.55048) (-1.61680,-0.51175) (-1.62999,-0.47217) (-1.64207,-0.43181) (-1.65303,-0.39074) (-1.66285,-0.34900) (-1.67153,-0.30669) (-1.67907,-0.26385) (-1.68546,-0.22056) (-1.69069,-0.17691) (-1.69476,-0.13294) (-1.69767,-0.08876) (-1.69942,-0.04442) (-1.70000,-0.00000) (-1.69942,0.04442) (-1.69767,0.08876) (-1.69476,0.13294) (-1.69069,0.17691) (-1.68546,0.22056) (-1.67907,0.26385) (-1.67153,0.30669) (-1.66285,0.34900) (-1.65303,0.39074) (-1.64207,0.43181) (-1.62999,0.47217) (-1.61680,0.51175) (-1.60249,0.55048) (-1.58709,0.58832) (-1.57060,0.62521) (-1.55303,0.66109) (-1.53439,0.69593) (-1.51471,0.72967) (-1.49399,0.76228) (-1.47224,0.79373) (-1.44949,0.82397) (-1.42574,0.85299) (-1.40101,0.88076) (-1.37533,0.90727) (-1.34870,0.93249) (-1.32115,0.95642) (-1.29269,0.97906) (-1.26335,1.00040) (-1.23314,1.02046) (-1.20208,1.03923) (-1.17020,1.05674) (-1.13752,1.07300) (-1.10406,1.08804) (-1.06984,1.10188) (-1.03489,1.11456) (-0.99923,1.12611) (-0.96289,1.13657) (-0.92589,1.14599) (-0.88825,1.15442) (-0.85000,1.16190) (-0.81117,1.16848) (-0.77178,1.17424) (-0.73187,1.17921) (-0.69145,1.18346) (-0.65056,1.18706) (-0.60923,1.19006) (-0.56747,1.19253) (-0.52533,1.19452) (-0.48283,1.19609) (-0.43999,1.19730) (-0.39686,1.19822) (-0.35345,1.19888) (-0.30980,1.19934) (-0.26594,1.19964) (-0.22189,1.19983) (-0.17770,1.19993) (-0.13338,1.19998) (-0.08897,1.20000) (-0.04450,1.20000) (-0.00000,1.20000) (0.04450,1.20000) (0.08897,1.20000) (0.13338,1.19998) (0.17770,1.19993) (0.22189,1.19983) (0.26594,1.19964) (0.30980,1.19934) (0.35345,1.19888) (0.39686,1.19822) (0.43999,1.19730) (0.48283,1.19609) (0.52533,1.19452) (0.56747,1.19253) (0.60923,1.19006) (0.65056,1.18706) (0.69145,1.18346) (0.73187,1.17921) (0.77178,1.17424) (0.81117,1.16848) (0.85000,1.16190) (0.88825,1.15442) (0.92589,1.14599) (0.96289,1.13657) (0.99923,1.12611) (1.03489,1.11456) (1.06984,1.10188) (1.10406,1.08804) (1.13752,1.07300) (1.17020,1.05674) (1.20208,1.03923) (1.23314,1.02046) (1.26335,1.00040) (1.29269,0.97906) (1.32115,0.95642) (1.34870,0.93249) (1.37533,0.90727) (1.40101,0.88076) (1.42574,0.85299) (1.44949,0.82397) (1.47224,0.79373) (1.49399,0.76228) (1.51471,0.72967) (1.53439,0.69593) (1.55303,0.66109) (1.57060,0.62521) (1.58709,0.58832) (1.60249,0.55048) (1.61680,0.51175) (1.62999,0.47217) (1.64207,0.43181) (1.65303,0.39074) (1.66285,0.34900) (1.67153,0.30669) (1.67907,0.26385) (1.68546,0.22056) (1.69069,0.17691) (1.69476,0.13294) (1.69767,0.08876) (1.69942,0.04442) (1.70000,0.00000)};
\fill(0,0)circle(1.2pt)node[below left]{$O$};
\node[below right]at(1.7,0){$1$};\node[below left]at(-1.7,0){$-1$};\node[above right]at(0,1.2){$1$};
\end{tikzpicture}
\caption{四次势阱的归一化相轨线。上、下支在转向点连接，轨线围住稳定平衡点 O。}\label{mp4:fig:l4-e2-s}
\end{figure}

在右端点 $(q,w)=(1,0)$，力 $-\kappa A^3$ 向左，相点沿竖直向下的方向离开横轴。下支随后向左，到 $q=-1$ 时转入上支。轨线在 $q=0$ 附近有 $w=\pm(1-q^4/2+\cdots)$，因此上下两支在纵轴附近比椭圆更平坦。

利用势能的偶对称性，周期等于 $0$ 到 $A$ 的单程时间的四倍。取 $x=Au$，则 $\dd x=A\dd u$，$\sqrt{A^4-x^4}=A^2\sqrt{1-u^4}$；上下限 $0,A$ 对应 $0,1$。因此
\begin{align*}
\frac T4
&=\int_0^A\frac{\dd x}{\sqrt{\kappa(A^4-x^4)/(2m)}}\\
&=\frac1A\sqrt{\frac{2m}{\kappa}}
\int_0^1\frac{\dd u}{\sqrt{1-u^4}},\qquad u=\frac xA.
\end{align*}
当 $u\to1^-$，$1-u^4\sim4(1-u)$，端点奇性为 $(1-u)^{-1/2}$，积分有限。记这个无量纲积分为 $I_4$，便有
\[
T=\frac{4I_4}{A}\sqrt{\frac{2m}{\kappa}},\qquad T\propto A^{-1}.
\]
振幅减半时周期加倍。虽然每个非零振幅都对应有限周期的闭合轨线，$A\to0$ 时周期仍趋于无穷大。

原点的稳定性由势能极小值判断。若给定位移界限 $\delta_x>0$、速度界限 $\delta_v>0$，选取初态使
\[
E<\min\left(\frac\kappa4\delta_x^4,\frac m2\delta_v^2\right),
\]
则能量守恒保证始终有 $|x|<\delta_x$、$|v|<\delta_v$。因此原点在相空间中稳定。图中围住 O 的闭合轨线也表示：小扰动后会持续振动，并不渐近回到 O。这里 $U''(0)=0$，线性化方程遗漏了首个非零的回复力项 $-\kappa x^3$，不能据此判断非线性系统失稳。
\end{solution}

\par\medskip\noindent\begin{minipage}{\textwidth}
\begin{example}[指数势垒上的反射]\label{mp4:example:l4-e3}
质量为$m$的质点从 $x\to+\infty$ 一侧以速率 $v_\infty$ 入射，势能为 $U(x)=U_0\ee^{-x/a}$，$U_0,a>0$。图\ref{mp4:fig:l4-e3-q}中向左为入射方向。求转向位置及以转向时刻为 $t=0$ 的 $x(t)$、$v(t)$，判断这条开轨线是否对应逃逸。
\end{example}

\begin{figure}[H]\centering
\begin{tikzpicture}
\draw[axis](-1.6,0)--(3.2,0)node[right]{$x$};
\draw[axis](0,-0.2)--(0,2.8)node[above]{$U$};
\draw[curve,domain=-.7:3,samples=90]plot(\x,{exp(-\x)});\draw[guide](-1,1.55)--(3,1.55)node[right]{$E$};\draw[axis](2.8,.5)--(1.6,.5)node[midway,above]{$v_{\infty}$};
\end{tikzpicture}
\caption{指数势垒。质点沿 $x$ 轴负方向入射，向左运动时势能增大。}\label{mp4:fig:l4-e3-q}
\end{figure}
\end{minipage}\par\medskip
\begin{solution}
远区 $x\to+\infty$ 时势能趋于零，入射速度为 $-v_\infty$，因此总能量为 $E=mv_\infty^2/2>0$。图\ref{mp4:fig:l4-e3-q}中势能向左单调升高，且向左无界，任意有限入射能量都会遇到转向点 T。由 $U(x_T)=E$，
\[
x_T=a\ln\frac{U_0}{E},\qquad x\ge x_T.
\]
$x_T$ 可以为负；它的位置由能量与 $U_0$ 的大小关系确定。

以 T 为位置基准，能量方程化为
\[
v^2=v_\infty^2\left[1-\ee^{-(x-x_T)/a}\right].
\]
图\ref{mp4:fig:l4-e3-s}的下支取负号，表示入射；上支取正号，表示反射。T 处的真实加速度为
\[
\ddot x_T=-\frac{U'(x_T)}m=\frac E{ma}
=\frac{v_\infty^2}{2a}>0,
\]
所以相点在 T 处沿竖直向上的方向从下支转入上支。

\begin{figure}[H]\centering
\begin{tikzpicture}
\draw[axis](-0.5,0)--(4.2,0)node[right]{$(x-x_T)/a$};
\draw[axis](0,-1.9)--(0,1.9)node[above]{$v/v_\infty$};
\draw[flow] plot coordinates{(3.72586,-1.28424) (3.68389,-1.28356) (3.64194,-1.28286) (3.60001,-1.28212) (3.55811,-1.28135) (3.51623,-1.28054) (3.47438,-1.27970) (3.43256,-1.27883) (3.39076,-1.27792) (3.34900,-1.27697) (3.30727,-1.27598) (3.26557,-1.27494) (3.22391,-1.27387) (3.18228,-1.27275) (3.14069,-1.27157) (3.09914,-1.27035) (3.05763,-1.26908) (3.01616,-1.26776) (2.97474,-1.26638) (2.93336,-1.26494) (2.89203,-1.26343) (2.85075,-1.26187) (2.80953,-1.26024) (2.76835,-1.25854) (2.72724,-1.25677) (2.68618,-1.25493) (2.64519,-1.25301) (2.60425,-1.25100) (2.56339,-1.24892) (2.52260,-1.24675) (2.48187,-1.24448) (2.44123,-1.24213) (2.40066,-1.23967) (2.36017,-1.23712) (2.31977,-1.23446) (2.27946,-1.23168) (2.23924,-1.22880) (2.19912,-1.22580) (2.15909,-1.22267) (2.11917,-1.21942) (2.07936,-1.21603) (2.03966,-1.21251) (2.00008,-1.20884) (1.96063,-1.20503) (1.92130,-1.20106) (1.88210,-1.19694) (1.84304,-1.19265) (1.80412,-1.18819) (1.76535,-1.18355) (1.72673,-1.17873) (1.68828,-1.17372) (1.64999,-1.16852) (1.61188,-1.16311) (1.57395,-1.15749) (1.53620,-1.15166) (1.49865,-1.14560) (1.46130,-1.13931) (1.42416,-1.13278) (1.38723,-1.12601) (1.35054,-1.11898) (1.31407,-1.11169) (1.27785,-1.10413) (1.24188,-1.09629) (1.20617,-1.08817) (1.17074,-1.07975) (1.13558,-1.07103) (1.10071,-1.06200) (1.06614,-1.05265) (1.03189,-1.04297) (0.99795,-1.03296) (0.96435,-1.02260) (0.93109,-1.01190) (0.89819,-1.00083) (0.86566,-0.98939) (0.83350,-0.97758) (0.80174,-0.96538) (0.77039,-0.95279) (0.73945,-0.93981) (0.70894,-0.92641) (0.67888,-0.91261) (0.64928,-0.89839) (0.62015,-0.88375) (0.59150,-0.86867) (0.56335,-0.85316) (0.53572,-0.83722) (0.50862,-0.82083) (0.48206,-0.80399) (0.45605,-0.78671) (0.43062,-0.76898) (0.40578,-0.75079) (0.38154,-0.73215) (0.35791,-0.71307) (0.33492,-0.69353) (0.31257,-0.67354) (0.29088,-0.65311) (0.26987,-0.63224) (0.24955,-0.61093) (0.22993,-0.58919) (0.21103,-0.56703) (0.19286,-0.54446) (0.17544,-0.52147) (0.15877,-0.49810) (0.14287,-0.47434) (0.12776,-0.45021) (0.11344,-0.42572) (0.09993,-0.40089) (0.08723,-0.37573) (0.07536,-0.35026) (0.06433,-0.32450) (0.05415,-0.29846) (0.04482,-0.27217) (0.03636,-0.24564) (0.02876,-0.21890) (0.02205,-0.19196) (0.01621,-0.16486) (0.01127,-0.13761) (0.00722,-0.11023) (0.00406,-0.08276) (0.00181,-0.05522) (0.00045,-0.02762) (0.00000,0.00000) (0.00045,0.02762) (0.00181,0.05522) (0.00406,0.08276) (0.00722,0.11023) (0.01127,0.13761) (0.01621,0.16486) (0.02205,0.19196) (0.02876,0.21890) (0.03636,0.24564) (0.04482,0.27217) (0.05415,0.29846) (0.06433,0.32450) (0.07536,0.35026) (0.08723,0.37573) (0.09993,0.40089) (0.11344,0.42572) (0.12776,0.45021) (0.14287,0.47434) (0.15877,0.49810) (0.17544,0.52147) (0.19286,0.54446) (0.21103,0.56703) (0.22993,0.58919) (0.24955,0.61093) (0.26987,0.63224) (0.29088,0.65311) (0.31257,0.67354) (0.33492,0.69353) (0.35791,0.71307) (0.38154,0.73215) (0.40578,0.75079) (0.43062,0.76898) (0.45605,0.78671) (0.48206,0.80399) (0.50862,0.82083) (0.53572,0.83722) (0.56335,0.85316) (0.59150,0.86867) (0.62015,0.88375) (0.64928,0.89839) (0.67888,0.91261) (0.70894,0.92641) (0.73945,0.93981) (0.77039,0.95279) (0.80174,0.96538) (0.83350,0.97758) (0.86566,0.98939) (0.89819,1.00083) (0.93109,1.01190) (0.96435,1.02260) (0.99795,1.03296) (1.03189,1.04297) (1.06614,1.05265) (1.10071,1.06200) (1.13558,1.07103) (1.17074,1.07975) (1.20617,1.08817) (1.24188,1.09629) (1.27785,1.10413) (1.31407,1.11169) (1.35054,1.11898) (1.38723,1.12601) (1.42416,1.13278) (1.46130,1.13931) (1.49865,1.14560) (1.53620,1.15166) (1.57395,1.15749) (1.61188,1.16311) (1.64999,1.16852) (1.68828,1.17372) (1.72673,1.17873) (1.76535,1.18355) (1.80412,1.18819) (1.84304,1.19265) (1.88210,1.19694) (1.92130,1.20106) (1.96063,1.20503) (2.00008,1.20884) (2.03966,1.21251) (2.07936,1.21603) (2.11917,1.21942) (2.15909,1.22267) (2.19912,1.22580) (2.23924,1.22880) (2.27946,1.23168) (2.31977,1.23446) (2.36017,1.23712) (2.40066,1.23967) (2.44123,1.24213) (2.48187,1.24448) (2.52260,1.24675) (2.56339,1.24892) (2.60425,1.25100) (2.64519,1.25301) (2.68618,1.25493) (2.72724,1.25677) (2.76835,1.25854) (2.80953,1.26024) (2.85075,1.26187) (2.89203,1.26343) (2.93336,1.26494) (2.97474,1.26638) (3.01616,1.26776) (3.05763,1.26908) (3.09914,1.27035) (3.14069,1.27157) (3.18228,1.27275) (3.22391,1.27387) (3.26557,1.27494) (3.30727,1.27598) (3.34900,1.27697) (3.39076,1.27792) (3.43256,1.27883) (3.47438,1.27970) (3.51623,1.28054) (3.55811,1.28135) (3.60001,1.28212) (3.64194,1.28286) (3.68389,1.28356) (3.72586,1.28424)};
\draw[guide](0,1.3)--(3.8,1.3);\draw[guide](0,-1.3)--(3.8,-1.3);\node[left]at(0,1.3){$1$};\node[left]at(0,-1.3){$-1$};\fill(0,0)circle(1.2pt)node[below left]{$T$};

\end{tikzpicture}
\caption{反射的相轨线。横轴为到转向位置的距离除以 $a$；T 处 $v=0$，加速度向右。}\label{mp4:fig:l4-e3-s}
\end{figure}

为求时间规律，在反射段 $t\ge0$ 令
\[
y=\exp\frac{x-x_T}{2a}\ge1,
\qquad x=x_T+2a\ln y.
\]
代入速度正支后，$2a\dot y=v_\infty\sqrt{y^2-1}$。分离变量并用转向点初值 $y(0)=1$ 确定下限，
\[
\int_1^y\frac{\dd u}{\sqrt{u^2-1}}
=\left[\ln(u+\sqrt{u^2-1})\right]_1^y
=\frac{v_\infty}{2a}\int_0^t\dd s
=\frac{v_\infty t}{2a}.
\]
左边即 $\operatorname{arcosh}y$，所以
\[
\operatorname{arcosh}y=\frac{v_\infty t}{2a},
\qquad y=\cosh\frac{v_\infty t}{2a}.
\]
方程具有时间反演对称性，位置关于转向时刻是偶函数，故完整解为
\[
x=x_T+2a\ln\cosh\frac{v_\infty t}{2a},
\]
\[
v=2a\frac{\sinh(v_\infty t/2a)}{\cosh(v_\infty t/2a)}\frac{v_\infty}{2a}
=v_\infty\tanh\frac{v_\infty t}{2a}.
\]
$t<0$ 时速度为负、$t>0$ 时为正，且 $x(0)=x_T$、$v(0)=0$。由 $\tanh^2u+\operatorname{sech}^2u=1$，可直接代回能量式验证。

接近 T 时，$x-x_T=v_\infty^2t^2/(4a)+O(t^4)$，转向过程平滑；相轨线在 T 处有竖直切线并不意味着加速度发散。远区有
\[
x=x_T+v_\infty|t|-2a\ln2+o(1),
\qquad v\longrightarrow\begin{cases}-v_\infty,&t\to-\infty,\\
+v_\infty,&t\to+\infty.
\end{cases}
\]
反射后质点的位置无限增大，速度趋于正的常量，因此属于向右逃逸。图中的水平虚线表示渐近速度；无穷远对应无限时间。
\end{solution}

\subsection{坐标相关的动能系数}
若用广义坐标 $q$ 描述约束运动，动能常为 $K=\tfrac12M(q)\dot q^2$。在坐标正则且 $M(q)>0$ 的区间内，能量关系给出
\[
\dot q=\pm\sqrt{\frac{2[E-U(q)]}{M(q)}},
\]
但相图的纵轴已经是 $\dot q$，不能再把动能写成常质量的 $m\dot q^2/2$。为使运动方程在转向点 $\dot q=0$ 也适用，直接使用 $L=M(q)\dot q^2/2-U(q)$ 的拉格朗日方程。各项为
\[
\begin{aligned}
\frac{\partial L}{\partial\dot q}&=M(q)\dot q,\qquad
\frac{\partial L}{\partial q}=\frac12M'(q)\dot q^2-U'(q),\\
\frac{\dd}{\dd t}\frac{\partial L}{\partial\dot q}
&=M(q)\ddot q+M'(q)\dot q^2.
\end{aligned}
\]
两项相减，得到
\begin{equation}\label{mp4:eq:l4-variable-mass}
M(q)\ddot q+\frac12M'(q)\dot q^2+U'(q)=0.
\end{equation}
$M(q)$ 是由几何约束产生的动能系数，物体本身的质量保持不变。正则动量为 $p=M(q)\dot q$；改用 $(q,p)$ 作图后，纵坐标及相图的面积性质都会改变。因此每幅相图都须明确标出所用坐标。

\par\medskip\noindent\begin{minipage}{\textwidth}
\begin{example}[抛物线导轨中的相轨线]\label{mp4:example:l4-e4}
质量为 $m$ 的小珠在竖直平面内沿光滑固定导轨 $y=x^2/(2a)$ 运动，$y$ 竖直向上，$a>0$。图\ref{mp4:fig:l4-e4-q}中最低点为 O。取 $q=x/a$、$w=\dot x/\sqrt{ga}$，最大水平位移为 $aA$。求相轨线、周期积分及小、大振幅的周期极限。
\end{example}

\begin{figure}[H]\centering
\begin{tikzpicture}
\draw[axis](-2.2,0)--(2.2,0)node[right]{$x$};
\draw[axis](0,-0.3)--(0,2.2)node[above]{$y$};
\draw[curve,domain=-1.9:1.9,samples=80]plot(\x,{.5*(\x)^2});\fill(1.25,.78125)circle(2pt)node[above left]{$P$};\draw[guide](1.25,0)--(1.25,.78125);\node[below]at(1.25,0){$x$};\draw[axis](1.25,.78125)--(1.25,.13);\node[right]at(1.25,.53){$mg$};\node[below left]at(0,0){$O$};
\end{tikzpicture}
\caption{抛物线导轨及小珠位置 P。横坐标 $x$ 是小珠的水平位置。}\label{mp4:fig:l4-e4-q}
\end{figure}
\end{minipage}\par\medskip
\begin{solution}
图\ref{mp4:fig:l4-e4-q}中的位置约束为 $y=x^2/(2a)$，求导得 $\dot y=(x/a)\dot x$。小珠的实际速率不是一般的 $|\dot x|$，而是
\[
v_{\rm path}=\sqrt{\dot x^2+\dot y^2}
=\sqrt{1+x^2/a^2}\,|\dot x|.
\]
固定光滑导轨的支持力与速度垂直，不做功。取 O 点为重力势能零点，机械能守恒：
\[
E=\frac12m\left(1+\frac{x^2}{a^2}\right)\dot x^2
+\frac{mgx^2}{2a}.
\]
用 $q=x/a$、$w=\dot x/\sqrt{ga}$、$\varepsilon=E/(mga)$，上式成为
\[
\varepsilon=\frac12(1+q^2)w^2+\frac12q^2.
\]
在水平转向位置 $q=\pm A$，有 $\dot x=\dot y=0$，故 $2\varepsilon=A^2$。相轨线为
\begin{equation}\label{mp4:eq:l4-parabola-phase}
w^2=\frac{A^2-q^2}{1+q^2},\qquad -A\le q\le A.
\end{equation}
图\ref{mp4:fig:l4-e4-s}的纵坐标是水平速度，正支向右、负支向左。在 O 处 $|w|=A$，水平速度的大小为 $A\sqrt{ga}$。

\begin{figure}[H]\centering
\begin{tikzpicture}
\draw[axis](-2.5,0)--(2.5,0)node[right]{$q$};
\draw[axis](0,-2.5)--(0,2.5)node[above]{$w$};
\draw[flow] plot coordinates{(2.00000,-0.00000) (1.99931,-0.02342) (1.99726,-0.04686) (1.99383,-0.07035) (1.98904,-0.09390) (1.98289,-0.11755) (1.97538,-0.14131) (1.96651,-0.16521) (1.95630,-0.18926) (1.94474,-0.21351) (1.93185,-0.23796) (1.91764,-0.26265) (1.90211,-0.28760) (1.88528,-0.31283) (1.86716,-0.33839) (1.84776,-0.36429) (1.82709,-0.39056) (1.80517,-0.41723) (1.78201,-0.44434) (1.75763,-0.47192) (1.73205,-0.50000) (1.70528,-0.52861) (1.67734,-0.55780) (1.64825,-0.58759) (1.61803,-0.61803) (1.58671,-0.64916) (1.55429,-0.68101) (1.52081,-0.71363) (1.48629,-0.74705) (1.45075,-0.78133) (1.41421,-0.81650) (1.37671,-0.85260) (1.33826,-0.88967) (1.29890,-0.92775) (1.25864,-0.96688) (1.21752,-1.00708) (1.17557,-1.04838) (1.13281,-1.09080) (1.08928,-1.13434) (1.04500,-1.17900) (1.00000,-1.22474) (0.95432,-1.27154) (0.90798,-1.31931) (0.86102,-1.36796) (0.81347,-1.41736) (0.76537,-1.46731) (0.71674,-1.51761) (0.66761,-1.56796) (0.61803,-1.61803) (0.56803,-1.66741) (0.51764,-1.71563) (0.46689,-1.76214) (0.41582,-1.80635) (0.36447,-1.84762) (0.31287,-1.88526) (0.26105,-1.91859) (0.20906,-1.94695) (0.15692,-1.96973) (0.10467,-1.98641) (0.05235,-1.99658) (0.00000,-2.00000) (-0.05235,-1.99658) (-0.10467,-1.98641) (-0.15692,-1.96973) (-0.20906,-1.94695) (-0.26105,-1.91859) (-0.31287,-1.88526) (-0.36447,-1.84762) (-0.41582,-1.80635) (-0.46689,-1.76214) (-0.51764,-1.71563) (-0.56803,-1.66741) (-0.61803,-1.61803) (-0.66761,-1.56796) (-0.71674,-1.51761) (-0.76537,-1.46731) (-0.81347,-1.41736) (-0.86102,-1.36796) (-0.90798,-1.31931) (-0.95432,-1.27154) (-1.00000,-1.22474) (-1.04500,-1.17900) (-1.08928,-1.13434) (-1.13281,-1.09080) (-1.17557,-1.04838) (-1.21752,-1.00708) (-1.25864,-0.96688) (-1.29890,-0.92775) (-1.33826,-0.88967) (-1.37671,-0.85260) (-1.41421,-0.81650) (-1.45075,-0.78133) (-1.48629,-0.74705) (-1.52081,-0.71363) (-1.55429,-0.68101) (-1.58671,-0.64916) (-1.61803,-0.61803) (-1.64825,-0.58759) (-1.67734,-0.55780) (-1.70528,-0.52861) (-1.73205,-0.50000) (-1.75763,-0.47192) (-1.78201,-0.44434) (-1.80517,-0.41723) (-1.82709,-0.39056) (-1.84776,-0.36429) (-1.86716,-0.33839) (-1.88528,-0.31283) (-1.90211,-0.28760) (-1.91764,-0.26265) (-1.93185,-0.23796) (-1.94474,-0.21351) (-1.95630,-0.18926) (-1.96651,-0.16521) (-1.97538,-0.14131) (-1.98289,-0.11755) (-1.98904,-0.09390) (-1.99383,-0.07035) (-1.99726,-0.04686) (-1.99931,-0.02342) (-2.00000,-0.00000) (-1.99931,0.02342) (-1.99726,0.04686) (-1.99383,0.07035) (-1.98904,0.09390) (-1.98289,0.11755) (-1.97538,0.14131) (-1.96651,0.16521) (-1.95630,0.18926) (-1.94474,0.21351) (-1.93185,0.23796) (-1.91764,0.26265) (-1.90211,0.28760) (-1.88528,0.31283) (-1.86716,0.33839) (-1.84776,0.36429) (-1.82709,0.39056) (-1.80517,0.41723) (-1.78201,0.44434) (-1.75763,0.47192) (-1.73205,0.50000) (-1.70528,0.52861) (-1.67734,0.55780) (-1.64825,0.58759) (-1.61803,0.61803) (-1.58671,0.64916) (-1.55429,0.68101) (-1.52081,0.71363) (-1.48629,0.74705) (-1.45075,0.78133) (-1.41421,0.81650) (-1.37671,0.85260) (-1.33826,0.88967) (-1.29890,0.92775) (-1.25864,0.96688) (-1.21752,1.00708) (-1.17557,1.04838) (-1.13281,1.09080) (-1.08928,1.13434) (-1.04500,1.17900) (-1.00000,1.22474) (-0.95432,1.27154) (-0.90798,1.31931) (-0.86102,1.36796) (-0.81347,1.41736) (-0.76537,1.46731) (-0.71674,1.51761) (-0.66761,1.56796) (-0.61803,1.61803) (-0.56803,1.66741) (-0.51764,1.71563) (-0.46689,1.76214) (-0.41582,1.80635) (-0.36447,1.84762) (-0.31287,1.88526) (-0.26105,1.91859) (-0.20906,1.94695) (-0.15692,1.96973) (-0.10467,1.98641) (-0.05235,1.99658) (-0.00000,2.00000) (0.05235,1.99658) (0.10467,1.98641) (0.15692,1.96973) (0.20906,1.94695) (0.26105,1.91859) (0.31287,1.88526) (0.36447,1.84762) (0.41582,1.80635) (0.46689,1.76214) (0.51764,1.71563) (0.56803,1.66741) (0.61803,1.61803) (0.66761,1.56796) (0.71674,1.51761) (0.76537,1.46731) (0.81347,1.41736) (0.86102,1.36796) (0.90798,1.31931) (0.95432,1.27154) (1.00000,1.22474) (1.04500,1.17900) (1.08928,1.13434) (1.13281,1.09080) (1.17557,1.04838) (1.21752,1.00708) (1.25864,0.96688) (1.29890,0.92775) (1.33826,0.88967) (1.37671,0.85260) (1.41421,0.81650) (1.45075,0.78133) (1.48629,0.74705) (1.52081,0.71363) (1.55429,0.68101) (1.58671,0.64916) (1.61803,0.61803) (1.64825,0.58759) (1.67734,0.55780) (1.70528,0.52861) (1.73205,0.50000) (1.75763,0.47192) (1.78201,0.44434) (1.80517,0.41723) (1.82709,0.39056) (1.84776,0.36429) (1.86716,0.33839) (1.88528,0.31283) (1.90211,0.28760) (1.91764,0.26265) (1.93185,0.23796) (1.94474,0.21351) (1.95630,0.18926) (1.96651,0.16521) (1.97538,0.14131) (1.98289,0.11755) (1.98904,0.09390) (1.99383,0.07035) (1.99726,0.04686) (1.99931,0.02342) (2.00000,0.00000)};
\node[below right]at(2,0){$A$};\node[below left]at(-2,0){$-A$};\node[above right]at(0,2){$A$};
\end{tikzpicture}
\caption{抛物线导轨的水平坐标相图，取 $A=2$。两轴分别表示无量纲水平位置和水平速度。}\label{mp4:fig:l4-e4-s}
\end{figure}

令 $\tau=t\sqrt{g/a}$，则 $\dd q/\dd\tau=w$。由正支积分从 $q=0$ 到 $q=A$，再利用对称性，
\[
T=4\sqrt{\frac ag}\int_0^A
\sqrt{\frac{1+q^2}{A^2-q^2}}\,\dd q.
\]
这里 $\sqrt{1+q^2}$ 来自导轨的弧长因素，也就是能量中的变系数动能。作 $q=A\sin\psi$ 后，$\dd q=A\cos\psi\,\dd\psi$ 恰好消去分母中的端点根号，得到
\[
T=4\sqrt{\frac ag}\int_0^{\pi/2}
\sqrt{1+A^2\sin^2\psi}\,\dd\psi.
\]

小振幅 $A\ll1$ 时，利用
\[
\sqrt{1+A^2\sin^2\psi}
=1+\frac{A^2}{2}\sin^2\psi+O(A^4)
\]
所需的角积分为
\[
\int_0^{\pi/2}\sin^2\psi\dd\psi
=\left[\frac\psi2-\frac{\sin2\psi}{4}\right]_0^{\pi/2}=\frac\pi4.
\]
代入展开式，得到
\[
T=2\pi\sqrt{\frac ag}\left(1+\frac{A^2}{4}+O(A^4)\right).
\]
首项对应槽底附近的简谐运动，角频率为 $\sqrt{g/a}$。

大振幅时将周期除以 $4A\sqrt{a/g}$，积分变为
\[
\int_0^{\pi/2}\sqrt{A^{-2}+\sin^2\psi}\,\dd\psi
\longrightarrow\int_0^{\pi/2}\sin\psi\,\dd\psi
=[-\cos\psi]_0^{\pi/2}=1.
\]
例如 $A\ge1$ 时被积函数不超过 $\sqrt2$，可以将极限移入积分。因此 $T\sim4A\sqrt{a/g}$。物理上，此时转向高度 $y_{\max}=aA^2/2$ 很大，从该高度近似竖直落到底部的时间为 $A\sqrt{a/g}$，与四分之一周期的首项相符。

由拉格朗日方程或式\eqref{mp4:eq:l4-variable-mass}还可求出相轨线的方向：
\[
(1+q^2)q''+q(q')^2+q=0,
\qquad q''=-\frac{q(1+w^2)}{1+q^2}.
\]
图中右转向点 $q=A,w=0$ 的加速度为 $-A/(1+A^2)<0$，因此相点在此竖直向下，方向与能量支的衔接一致。
\end{solution}

\section{平衡点与分界轨线}
平衡点附近的轨线反映回复或偏离趋势。普通转折点、势垒顶与分界轨线具有不同的到达时间，临界能量不能单独保证有限时间通行。

\subsection{平衡点附近的运动}
平衡状态为 $(x_e,0)$，其中 $U'(x_e)=0$。令 $\xi=x-x_e$，在 $U''(x_e)\ne0$ 时保留一次项：
\[
\ddot\xi+\frac{U''(x_e)}m\xi=0.
\]
取 $\xi\propto\ee^{rt}$，特征方程为 $r^2+U''(x_e)/m=0$。若 $U''(x_e)>0$，两根为 $\pm\mathrm i\omega_e$，其中 $\omega_e=\sqrt{U''(x_e)/m}$，解为 $\xi=A\cos\omega_et+B\sin\omega_et$。在线性近似下，相轨线是一族闭合椭圆；原非线性系统在该势能极小点附近也有闭合轨线，但通常不是精确的椭圆。平衡点称为中心。无耗散时，附近轨线不会逐渐收缩到中心，因此这里的稳定不是渐近稳定。

若 $U''(x_e)<0$，令 $\lambda=\sqrt{-U''(x_e)/m}$，则
\[
\xi=C_+\ee^{\lambda t}+C_-\ee^{-\lambda t},\qquad
v=\lambda C_+\ee^{\lambda t}-\lambda C_-\ee^{-\lambda t}.
\]
相图在原点附近有两条特殊方向 $v=\pm\lambda\xi$。沿 $v=-\lambda\xi$，相点随时间接近平衡点；沿 $v=\lambda\xi$，相点离开平衡点。这是鞍点，任意小的初态扰动都可能激发增长项。


\begin{figure}[H]\centering
\begin{tikzpicture}
\begin{scope}[shift={(0,0)},scale=0.9]
\draw[axis](-2,0)--(2,0)node[right]{$\xi$};
\draw[axis](0,-1.6)--(0,1.6)node[above]{$v$};
\draw[flow] plot coordinates{(1.50000,-0.00000) (1.49884,-0.03926) (1.49538,-0.07846) (1.48960,-0.11754) (1.48153,-0.15643) (1.47118,-0.19509) (1.45855,-0.23345) (1.44368,-0.27144) (1.42658,-0.30902) (1.40729,-0.34612) (1.38582,-0.38268) (1.36221,-0.41866) (1.33651,-0.45399) (1.30874,-0.48862) (1.27896,-0.52250) (1.24720,-0.55557) (1.21353,-0.58779) (1.17798,-0.61909) (1.14061,-0.64945) (1.10148,-0.67880) (1.06066,-0.70711) (1.01820,-0.73432) (0.97417,-0.76041) (0.92864,-0.78532) (0.88168,-0.80902) (0.83336,-0.83147) (0.78375,-0.85264) (0.73293,-0.87250) (0.68099,-0.89101) (0.62799,-0.90814) (0.57403,-0.92388) (0.51918,-0.93819) (0.46353,-0.95106) (0.40716,-0.96246) (0.35017,-0.97237) (0.29264,-0.98079) (0.23465,-0.98769) (0.17631,-0.99307) (0.11769,-0.99692) (0.05889,-0.99923) (0.00000,-1.00000) (-0.05889,-0.99923) (-0.11769,-0.99692) (-0.17631,-0.99307) (-0.23465,-0.98769) (-0.29264,-0.98079) (-0.35017,-0.97237) (-0.40716,-0.96246) (-0.46353,-0.95106) (-0.51918,-0.93819) (-0.57403,-0.92388) (-0.62799,-0.90814) (-0.68099,-0.89101) (-0.73293,-0.87250) (-0.78375,-0.85264) (-0.83336,-0.83147) (-0.88168,-0.80902) (-0.92864,-0.78532) (-0.97417,-0.76041) (-1.01820,-0.73432) (-1.06066,-0.70711) (-1.10148,-0.67880) (-1.14061,-0.64945) (-1.17798,-0.61909) (-1.21353,-0.58779) (-1.24720,-0.55557) (-1.27896,-0.52250) (-1.30874,-0.48862) (-1.33651,-0.45399) (-1.36221,-0.41866) (-1.38582,-0.38268) (-1.40729,-0.34612) (-1.42658,-0.30902) (-1.44368,-0.27144) (-1.45855,-0.23345) (-1.47118,-0.19509) (-1.48153,-0.15643) (-1.48960,-0.11754) (-1.49538,-0.07846) (-1.49884,-0.03926) (-1.50000,-0.00000) (-1.49884,0.03926) (-1.49538,0.07846) (-1.48960,0.11754) (-1.48153,0.15643) (-1.47118,0.19509) (-1.45855,0.23345) (-1.44368,0.27144) (-1.42658,0.30902) (-1.40729,0.34612) (-1.38582,0.38268) (-1.36221,0.41866) (-1.33651,0.45399) (-1.30874,0.48862) (-1.27896,0.52250) (-1.24720,0.55557) (-1.21353,0.58779) (-1.17798,0.61909) (-1.14061,0.64945) (-1.10148,0.67880) (-1.06066,0.70711) (-1.01820,0.73432) (-0.97417,0.76041) (-0.92864,0.78532) (-0.88168,0.80902) (-0.83336,0.83147) (-0.78375,0.85264) (-0.73293,0.87250) (-0.68099,0.89101) (-0.62799,0.90814) (-0.57403,0.92388) (-0.51918,0.93819) (-0.46353,0.95106) (-0.40716,0.96246) (-0.35017,0.97237) (-0.29264,0.98079) (-0.23465,0.98769) (-0.17631,0.99307) (-0.11769,0.99692) (-0.05889,0.99923) (-0.00000,1.00000) (0.05889,0.99923) (0.11769,0.99692) (0.17631,0.99307) (0.23465,0.98769) (0.29264,0.98079) (0.35017,0.97237) (0.40716,0.96246) (0.46353,0.95106) (0.51918,0.93819) (0.57403,0.92388) (0.62799,0.90814) (0.68099,0.89101) (0.73293,0.87250) (0.78375,0.85264) (0.83336,0.83147) (0.88168,0.80902) (0.92864,0.78532) (0.97417,0.76041) (1.01820,0.73432) (1.06066,0.70711) (1.10148,0.67880) (1.14061,0.64945) (1.17798,0.61909) (1.21353,0.58779) (1.24720,0.55557) (1.27896,0.52250) (1.30874,0.48862) (1.33651,0.45399) (1.36221,0.41866) (1.38582,0.38268) (1.40729,0.34612) (1.42658,0.30902) (1.44368,0.27144) (1.45855,0.23345) (1.47118,0.19509) (1.48153,0.15643) (1.48960,0.11754) (1.49538,0.07846) (1.49884,0.03926) (1.50000,0.00000)};
\fill(0,0)circle(1.2pt)node[below left]{$C$};
\node at(0,-2){中心};\end{scope}
\begin{scope}[shift={(5,0)},scale=0.9]
\draw[axis](-2,0)--(2,0)node[right]{$\xi$};
\draw[axis](0,-1.6)--(0,1.6)node[above]{$v$};
\draw[flow,line width=.9pt] plot coordinates{(-1.50000,1.20000) (-0.08000,0.06400)};
\draw[flow,line width=.9pt] plot coordinates{(1.50000,-1.20000) (0.08000,-0.06400)};
\draw[flow,dashed] plot coordinates{(0.08000,0.06400) (1.50000,1.20000)};
\draw[flow,dashed] plot coordinates{(-0.08000,-0.06400) (-1.50000,-1.20000)};
\fill(0,0)circle(1.2pt);\node at(.35,-.85){$S$};
\node at(0,-2){鞍点};\end{scope}

\end{tikzpicture}
\caption{平衡点附近的线性相图。右图实线为趋近鞍点的稳定方向，虚线为离开鞍点的不稳定方向；箭头不经过鞍点本身。}\label{mp4:fig:l4-linear}
\end{figure}


若 $U''(x_e)=0$，线性化不能决定稳定性。须继续检查势能的首个非零项：严格极小点仍可稳定，水平拐点则未必形成势阱。此时可用附近的能量等值线判断质点是否被限制在平衡位置周围。

\subsection{势垒能量与无限到达时间}
设 $x_s$ 是非退化势能极大点，$U_s=U(x_s)$，附近有
\[
U(x)=U_s-\frac12m\lambda^2(x-x_s)^2+O((x-x_s)^3).
\]
当 $E=U_s$，$|v|\sim\lambda|x-x_s|$，所以趋近 $x_s$ 的时间包含
\[
\int\frac{\dd x}{|v|}\sim\frac1\lambda\int\frac{\dd x}{|x-x_s|}.
\]
在鞍点同一侧取距离 $\xi=|x-x_s|$，从固定距离 $d$ 趋近到 $\epsilon>0$ 的时间首项为
\[
\frac1\lambda\int_\epsilon^d\frac{\dd\xi}{\xi}
=\frac1\lambda[\ln\xi]_\epsilon^d
=\frac1\lambda\ln\frac d\epsilon\longrightarrow\infty
\qquad(\epsilon\to0^+).
\]
因而轨线只在无限时间趋近鞍点；经过鞍点的另一支须作为另一段运动处理。鞍点本身对应静止解。

经过这种临界能量的轨线常把振动、跨阱或转动分开，称为分界轨线。若某支在过去和将来都趋于同一个鞍点，称为同宿轨线；若两端趋于不同鞍点，称为异宿轨线。这两种名称都以轨线在无限时间的趋近位置为依据。

\par\medskip\noindent\begin{minipage}{\textwidth}
\begin{example}[双阱中的跨阱阈值]\label{mp4:example:l4-e5}
质量为$m$的质点处在双阱势能 $U(x)=\lambda(x^2-a^2)^2/4$ 中，$\lambda,a>0$。图\ref{mp4:fig:l4-e5-q}中的两个最低点为 $x=\pm a$。画出低于、等于和高于势垒的相轨线，求平衡点性质和小振动频率。若从右阱最低点向左给予速度，多大速度才能进入左阱？
\end{example}

\begin{figure}[H]\centering
\begin{tikzpicture}
\draw[axis](-2.6,0)--(2.6,0)node[right]{$x/a$};
\draw[axis](0,-0.2)--(0,2.1)node[above]{$U/(\lambda a^4)$};
\draw[curve,domain=-1.75:1.75,samples=140]plot({1.2*\x},{2*((\x)^2-1)^2/4});\draw[guide](-2.2,.5)--(2.2,.5)node[right]{$1/4$};\node[below]at(-1.2,0){$-1$};\node[below]at(1.2,0){$1$};\node[above left]at(0,.5){$U_s$};
\end{tikzpicture}
\caption{双阱势能。中间的极大点把两个最低点分隔开。}\label{mp4:fig:l4-e5-q}
\end{figure}
\end{minipage}\par\medskip
\begin{solution}
先由位置力求平衡点：
\[
U'(x)=\frac\lambda4\,2(x^2-a^2)\,2x
=\lambda x(x^2-a^2)=0,
\]
\[
U''(x)=\lambda[(x^2-a^2)+2x^2]=\lambda(3x^2-a^2).
\]
图\ref{mp4:fig:l4-e5-q}中的两个阱底 $x=\pm a$ 满足 $U''=2\lambda a^2>0$，是中心；势垒顶 $x=0$ 满足 $U''=-\lambda a^2<0$，是鞍点。令 $x=a+\xi$，小扰动方程为 $m\ddot\xi+2\lambda a^2\xi=0$，故两阱的小振动角频率均为
\[
\omega_c=\sqrt{\frac{2\lambda a^2}{m}}=\sqrt2\,\omega_b,
\qquad\omega_b=\sqrt{\frac{\lambda a^2}{m}}.
\]
势垒附近则有 $\ddot x-\omega_b^2x\simeq0$，扰动包含增长、衰减两个指数模态。

取 $q=x/a$、$w=v/(a\omega_b)$、$h=E/(\lambda a^4)$，能量式为
\begin{equation}\label{mp4:eq:l4-doublewell}
\frac12w^2+\frac14(q^2-1)^2=h.
\end{equation}
在图\ref{mp4:fig:l4-e5-s}中，$C_\pm=(\pm1,0)$，$S=(0,0)$；势垒高度对应 $h_b=1/4$。

\begin{figure}[H]\centering
\begin{tikzpicture}[scale=1.35]
\draw[axis](-2.9,0)--(2.9,0)node[right]{$q$};
\draw[axis](0,-1.7)--(0,1.7)node[above]{$w$};
\draw[flow,line width=.4pt] plot coordinates{(-0.88544,-0.00000) (-0.88559,-0.01134) (-0.88604,-0.02267) (-0.88680,-0.03400) (-0.88786,-0.04531) (-0.88923,-0.05662) (-0.89089,-0.06791) (-0.89285,-0.07919) (-0.89511,-0.09044) (-0.89767,-0.10167) (-0.90052,-0.11287) (-0.90367,-0.12404) (-0.90711,-0.13517) (-0.91083,-0.14627) (-0.91484,-0.15732) (-0.91914,-0.16833) (-0.92371,-0.17929) (-0.92857,-0.19020) (-0.93369,-0.20105) (-0.93909,-0.21183) (-0.94475,-0.22255) (-0.95068,-0.23319) (-0.95686,-0.24376) (-0.96330,-0.25424) (-0.96999,-0.26464) (-0.97692,-0.27494) (-0.98410,-0.28514) (-0.99151,-0.29524) (-0.99915,-0.30523) (-1.00702,-0.31510) (-1.01511,-0.32485) (-1.02341,-0.33446) (-1.03192,-0.34394) (-1.04063,-0.35328) (-1.04954,-0.36246) (-1.05865,-0.37149) (-1.06793,-0.38035) (-1.07740,-0.38904) (-1.08703,-0.39755) (-1.09684,-0.40587) (-1.10680,-0.41399) (-1.11691,-0.42191) (-1.12717,-0.42961) (-1.13756,-0.43710) (-1.14809,-0.44436) (-1.15874,-0.45138) (-1.16950,-0.45816) (-1.18037,-0.46469) (-1.19135,-0.47095) (-1.20242,-0.47695) (-1.21357,-0.48267) (-1.22481,-0.48810) (-1.23611,-0.49325) (-1.24748,-0.49809) (-1.25890,-0.50262) (-1.27037,-0.50683) (-1.28188,-0.51073) (-1.29342,-0.51429) (-1.30499,-0.51751) (-1.31657,-0.52039) (-1.32816,-0.52291) (-1.33975,-0.52508) (-1.35133,-0.52688) (-1.36289,-0.52831) (-1.37443,-0.52936) (-1.38594,-0.53004) (-1.39741,-0.53032) (-1.40884,-0.53021) (-1.42020,-0.52971) (-1.43151,-0.52880) (-1.44274,-0.52749) (-1.45390,-0.52577) (-1.46496,-0.52365) (-1.47594,-0.52110) (-1.48681,-0.51815) (-1.49758,-0.51477) (-1.50823,-0.51098) (-1.51875,-0.50677) (-1.52915,-0.50214) (-1.53940,-0.49710) (-1.54952,-0.49163) (-1.55948,-0.48575) (-1.56928,-0.47945) (-1.57892,-0.47274) (-1.58838,-0.46562) (-1.59767,-0.45809) (-1.60677,-0.45016) (-1.61568,-0.44183) (-1.62439,-0.43311) (-1.63290,-0.42400) (-1.64121,-0.41451) (-1.64929,-0.40464) (-1.65716,-0.39440) (-1.66480,-0.38380) (-1.67221,-0.37285) (-1.67939,-0.36155) (-1.68632,-0.34992) (-1.69301,-0.33795) (-1.69945,-0.32567) (-1.70564,-0.31309) (-1.71156,-0.30020) (-1.71723,-0.28703) (-1.72262,-0.27359) (-1.72775,-0.25989) (-1.73260,-0.24593) (-1.73718,-0.23174) (-1.74147,-0.21733) (-1.74548,-0.20270) (-1.74921,-0.18788) (-1.75264,-0.17288) (-1.75579,-0.15771) (-1.75864,-0.14238) (-1.76120,-0.12691) (-1.76346,-0.11132) (-1.76542,-0.09562) (-1.76709,-0.07983) (-1.76845,-0.06396) (-1.76951,-0.04802) (-1.77027,-0.03204) (-1.77072,-0.01603) (-1.77088,-0.00000) (-1.77072,0.01603) (-1.77027,0.03204) (-1.76951,0.04802) (-1.76845,0.06396) (-1.76709,0.07983) (-1.76542,0.09562) (-1.76346,0.11132) (-1.76120,0.12691) (-1.75864,0.14238) (-1.75579,0.15771) (-1.75264,0.17288) (-1.74921,0.18788) (-1.74548,0.20270) (-1.74147,0.21733) (-1.73718,0.23174) (-1.73260,0.24593) (-1.72775,0.25989) (-1.72262,0.27359) (-1.71723,0.28703) (-1.71156,0.30020) (-1.70564,0.31309) (-1.69945,0.32567) (-1.69301,0.33795) (-1.68632,0.34992) (-1.67939,0.36155) (-1.67221,0.37285) (-1.66480,0.38380) (-1.65716,0.39440) (-1.64929,0.40464) (-1.64121,0.41451) (-1.63290,0.42400) (-1.62439,0.43311) (-1.61568,0.44183) (-1.60677,0.45016) (-1.59767,0.45809) (-1.58838,0.46562) (-1.57892,0.47274) (-1.56928,0.47945) (-1.55948,0.48575) (-1.54952,0.49163) (-1.53940,0.49710) (-1.52915,0.50214) (-1.51875,0.50677) (-1.50823,0.51098) (-1.49758,0.51477) (-1.48681,0.51815) (-1.47594,0.52110) (-1.46496,0.52365) (-1.45390,0.52577) (-1.44274,0.52749) (-1.43151,0.52880) (-1.42020,0.52971) (-1.40884,0.53021) (-1.39741,0.53032) (-1.38594,0.53004) (-1.37443,0.52936) (-1.36289,0.52831) (-1.35133,0.52688) (-1.33975,0.52508) (-1.32816,0.52291) (-1.31657,0.52039) (-1.30499,0.51751) (-1.29342,0.51429) (-1.28188,0.51073) (-1.27037,0.50683) (-1.25890,0.50262) (-1.24748,0.49809) (-1.23611,0.49325) (-1.22481,0.48810) (-1.21357,0.48267) (-1.20242,0.47695) (-1.19135,0.47095) (-1.18037,0.46469) (-1.16950,0.45816) (-1.15874,0.45138) (-1.14809,0.44436) (-1.13756,0.43710) (-1.12717,0.42961) (-1.11691,0.42191) (-1.10680,0.41399) (-1.09684,0.40587) (-1.08703,0.39755) (-1.07740,0.38904) (-1.06793,0.38035) (-1.05865,0.37149) (-1.04954,0.36246) (-1.04063,0.35328) (-1.03192,0.34394) (-1.02341,0.33446) (-1.01511,0.32485) (-1.00702,0.31510) (-0.99915,0.30523) (-0.99151,0.29524) (-0.98410,0.28514) (-0.97692,0.27494) (-0.96999,0.26464) (-0.96330,0.25424) (-0.95686,0.24376) (-0.95068,0.23319) (-0.94475,0.22255) (-0.93909,0.21183) (-0.93369,0.20105) (-0.92857,0.19020) (-0.92371,0.17929) (-0.91914,0.16833) (-0.91484,0.15732) (-0.91083,0.14627) (-0.90711,0.13517) (-0.90367,0.12404) (-0.90052,0.11287) (-0.89767,0.10167) (-0.89511,0.09044) (-0.89285,0.07919) (-0.89089,0.06791) (-0.88923,0.05662) (-0.88786,0.04531) (-0.88680,0.03400) (-0.88604,0.02267) (-0.88559,0.01134) (-0.88544,0.00000)};
\draw[flow,line width=1pt] plot coordinates{(0.00000,-0.00000) (-0.00034,-0.00030) (-0.00136,-0.00121) (-0.00305,-0.00272) (-0.00542,-0.00484) (-0.00847,-0.00756) (-0.01219,-0.01088) (-0.01658,-0.01480) (-0.02163,-0.01931) (-0.02735,-0.02442) (-0.03373,-0.03011) (-0.04077,-0.03639) (-0.04845,-0.04325) (-0.05678,-0.05068) (-0.06575,-0.05867) (-0.07536,-0.06723) (-0.08559,-0.07634) (-0.09644,-0.08600) (-0.10790,-0.09619) (-0.11996,-0.10691) (-0.13263,-0.11815) (-0.14588,-0.12989) (-0.15971,-0.14213) (-0.17411,-0.15485) (-0.18906,-0.16804) (-0.20457,-0.18167) (-0.22061,-0.19575) (-0.23719,-0.21025) (-0.25427,-0.22515) (-0.27187,-0.24044) (-0.28995,-0.25609) (-0.30851,-0.27209) (-0.32754,-0.28842) (-0.34703,-0.30505) (-0.36695,-0.32196) (-0.38731,-0.33913) (-0.40807,-0.35653) (-0.42924,-0.37413) (-0.45078,-0.39192) (-0.47270,-0.40985) (-0.49497,-0.42791) (-0.51759,-0.44606) (-0.54052,-0.46428) (-0.56377,-0.48252) (-0.58730,-0.50077) (-0.61111,-0.51899) (-0.63518,-0.53715) (-0.65950,-0.55521) (-0.68404,-0.57314) (-0.70879,-0.59090) (-0.73373,-0.60847) (-0.75885,-0.62580) (-0.78413,-0.64287) (-0.80955,-0.65963) (-0.83509,-0.67605) (-0.86074,-0.69209) (-0.88647,-0.70773) (-0.91228,-0.72292) (-0.93814,-0.73762) (-0.96404,-0.75182) (-0.98995,-0.76547) (-1.01586,-0.77853) (-1.04176,-0.79098) (-1.06762,-0.80277) (-1.09343,-0.81389) (-1.11916,-0.82429) (-1.14481,-0.83396) (-1.17035,-0.84285) (-1.19577,-0.85094) (-1.22105,-0.85820) (-1.24617,-0.86461) (-1.27111,-0.87014) (-1.29586,-0.87477) (-1.32040,-0.87847) (-1.34472,-0.88123) (-1.36879,-0.88302) (-1.39260,-0.88383) (-1.41613,-0.88365) (-1.43938,-0.88244) (-1.46231,-0.88022) (-1.48492,-0.87695) (-1.50720,-0.87264) (-1.52911,-0.86727) (-1.55066,-0.86085) (-1.57183,-0.85336) (-1.59259,-0.84481) (-1.61294,-0.83519) (-1.63287,-0.82451) (-1.65235,-0.81277) (-1.67139,-0.79997) (-1.68995,-0.78614) (-1.70803,-0.77127) (-1.72563,-0.75537) (-1.74271,-0.73847) (-1.75928,-0.72058) (-1.77533,-0.70171) (-1.79084,-0.68189) (-1.80579,-0.66113) (-1.82019,-0.63946) (-1.83402,-0.61691) (-1.84727,-0.59351) (-1.85993,-0.56927) (-1.87200,-0.54424) (-1.88346,-0.51844) (-1.89431,-0.49191) (-1.90454,-0.46468) (-1.91415,-0.43679) (-1.92312,-0.40827) (-1.93145,-0.37917) (-1.93913,-0.34953) (-1.94617,-0.31939) (-1.95255,-0.28878) (-1.95827,-0.25776) (-1.96332,-0.22636) (-1.96771,-0.19464) (-1.97143,-0.16263) (-1.97448,-0.13039) (-1.97685,-0.09796) (-1.97854,-0.06539) (-1.97956,-0.03272) (-1.97990,-0.00000) (-1.97956,0.03272) (-1.97854,0.06539) (-1.97685,0.09796) (-1.97448,0.13039) (-1.97143,0.16263) (-1.96771,0.19464) (-1.96332,0.22636) (-1.95827,0.25776) (-1.95255,0.28878) (-1.94617,0.31939) (-1.93913,0.34953) (-1.93145,0.37917) (-1.92312,0.40827) (-1.91415,0.43679) (-1.90454,0.46468) (-1.89431,0.49191) (-1.88346,0.51844) (-1.87200,0.54424) (-1.85993,0.56927) (-1.84727,0.59351) (-1.83402,0.61691) (-1.82019,0.63946) (-1.80579,0.66113) (-1.79084,0.68189) (-1.77533,0.70171) (-1.75928,0.72058) (-1.74271,0.73847) (-1.72563,0.75537) (-1.70803,0.77127) (-1.68995,0.78614) (-1.67139,0.79997) (-1.65235,0.81277) (-1.63287,0.82451) (-1.61294,0.83519) (-1.59259,0.84481) (-1.57183,0.85336) (-1.55066,0.86085) (-1.52911,0.86727) (-1.50720,0.87264) (-1.48492,0.87695) (-1.46231,0.88022) (-1.43938,0.88244) (-1.41613,0.88365) (-1.39260,0.88383) (-1.36879,0.88302) (-1.34472,0.88123) (-1.32040,0.87847) (-1.29586,0.87477) (-1.27111,0.87014) (-1.24617,0.86461) (-1.22105,0.85820) (-1.19577,0.85094) (-1.17035,0.84285) (-1.14481,0.83396) (-1.11916,0.82429) (-1.09343,0.81389) (-1.06762,0.80277) (-1.04176,0.79098) (-1.01586,0.77853) (-0.98995,0.76547) (-0.96404,0.75182) (-0.93814,0.73762) (-0.91228,0.72292) (-0.88647,0.70773) (-0.86074,0.69209) (-0.83509,0.67605) (-0.80955,0.65963) (-0.78413,0.64287) (-0.75885,0.62580) (-0.73373,0.60847) (-0.70879,0.59090) (-0.68404,0.57314) (-0.65950,0.55521) (-0.63518,0.53715) (-0.61111,0.51899) (-0.58730,0.50077) (-0.56377,0.48252) (-0.54052,0.46428) (-0.51759,0.44606) (-0.49497,0.42791) (-0.47270,0.40985) (-0.45078,0.39192) (-0.42924,0.37413) (-0.40807,0.35653) (-0.38731,0.33913) (-0.36695,0.32196) (-0.34703,0.30505) (-0.32754,0.28842) (-0.30851,0.27209) (-0.28995,0.25609) (-0.27187,0.24044) (-0.25427,0.22515) (-0.23719,0.21025) (-0.22061,0.19575) (-0.20457,0.18167) (-0.18906,0.16804) (-0.17411,0.15485) (-0.15971,0.14213) (-0.14588,0.12989) (-0.13263,0.11815) (-0.11996,0.10691) (-0.10790,0.09619) (-0.09644,0.08600) (-0.08559,0.07634) (-0.07536,0.06723) (-0.06575,0.05867) (-0.05678,0.05068) (-0.04845,0.04325) (-0.04077,0.03639) (-0.03373,0.03011) (-0.02735,0.02442) (-0.02163,0.01931) (-0.01658,0.01480) (-0.01219,0.01088) (-0.00847,0.00756) (-0.00542,0.00484) (-0.00305,0.00272) (-0.00136,0.00121) (-0.00034,0.00030) (0.00000,0.00000)};
\draw[flow,line width=.4pt] plot coordinates{(1.77088,-0.00000) (1.77072,-0.01603) (1.77027,-0.03204) (1.76951,-0.04802) (1.76845,-0.06396) (1.76709,-0.07983) (1.76542,-0.09562) (1.76346,-0.11132) (1.76120,-0.12691) (1.75864,-0.14238) (1.75579,-0.15771) (1.75264,-0.17288) (1.74921,-0.18788) (1.74548,-0.20270) (1.74147,-0.21733) (1.73718,-0.23174) (1.73260,-0.24593) (1.72775,-0.25989) (1.72262,-0.27359) (1.71723,-0.28703) (1.71156,-0.30020) (1.70564,-0.31309) (1.69945,-0.32567) (1.69301,-0.33795) (1.68632,-0.34992) (1.67939,-0.36155) (1.67221,-0.37285) (1.66480,-0.38380) (1.65716,-0.39440) (1.64929,-0.40464) (1.64121,-0.41451) (1.63290,-0.42400) (1.62439,-0.43311) (1.61568,-0.44183) (1.60677,-0.45016) (1.59767,-0.45809) (1.58838,-0.46562) (1.57892,-0.47274) (1.56928,-0.47945) (1.55948,-0.48575) (1.54952,-0.49163) (1.53940,-0.49710) (1.52915,-0.50214) (1.51875,-0.50677) (1.50823,-0.51098) (1.49758,-0.51477) (1.48681,-0.51815) (1.47594,-0.52110) (1.46496,-0.52365) (1.45390,-0.52577) (1.44274,-0.52749) (1.43151,-0.52880) (1.42020,-0.52971) (1.40884,-0.53021) (1.39741,-0.53032) (1.38594,-0.53004) (1.37443,-0.52936) (1.36289,-0.52831) (1.35133,-0.52688) (1.33975,-0.52508) (1.32816,-0.52291) (1.31657,-0.52039) (1.30499,-0.51751) (1.29342,-0.51429) (1.28188,-0.51073) (1.27037,-0.50683) (1.25890,-0.50262) (1.24748,-0.49809) (1.23611,-0.49325) (1.22481,-0.48810) (1.21357,-0.48267) (1.20242,-0.47695) (1.19135,-0.47095) (1.18037,-0.46469) (1.16950,-0.45816) (1.15874,-0.45138) (1.14809,-0.44436) (1.13756,-0.43710) (1.12717,-0.42961) (1.11691,-0.42191) (1.10680,-0.41399) (1.09684,-0.40587) (1.08703,-0.39755) (1.07740,-0.38904) (1.06793,-0.38035) (1.05865,-0.37149) (1.04954,-0.36246) (1.04063,-0.35328) (1.03192,-0.34394) (1.02341,-0.33446) (1.01511,-0.32485) (1.00702,-0.31510) (0.99915,-0.30523) (0.99151,-0.29524) (0.98410,-0.28514) (0.97692,-0.27494) (0.96999,-0.26464) (0.96330,-0.25424) (0.95686,-0.24376) (0.95068,-0.23319) (0.94475,-0.22255) (0.93909,-0.21183) (0.93369,-0.20105) (0.92857,-0.19020) (0.92371,-0.17929) (0.91914,-0.16833) (0.91484,-0.15732) (0.91083,-0.14627) (0.90711,-0.13517) (0.90367,-0.12404) (0.90052,-0.11287) (0.89767,-0.10167) (0.89511,-0.09044) (0.89285,-0.07919) (0.89089,-0.06791) (0.88923,-0.05662) (0.88786,-0.04531) (0.88680,-0.03400) (0.88604,-0.02267) (0.88559,-0.01134) (0.88544,-0.00000) (0.88559,0.01134) (0.88604,0.02267) (0.88680,0.03400) (0.88786,0.04531) (0.88923,0.05662) (0.89089,0.06791) (0.89285,0.07919) (0.89511,0.09044) (0.89767,0.10167) (0.90052,0.11287) (0.90367,0.12404) (0.90711,0.13517) (0.91083,0.14627) (0.91484,0.15732) (0.91914,0.16833) (0.92371,0.17929) (0.92857,0.19020) (0.93369,0.20105) (0.93909,0.21183) (0.94475,0.22255) (0.95068,0.23319) (0.95686,0.24376) (0.96330,0.25424) (0.96999,0.26464) (0.97692,0.27494) (0.98410,0.28514) (0.99151,0.29524) (0.99915,0.30523) (1.00702,0.31510) (1.01511,0.32485) (1.02341,0.33446) (1.03192,0.34394) (1.04063,0.35328) (1.04954,0.36246) (1.05865,0.37149) (1.06793,0.38035) (1.07740,0.38904) (1.08703,0.39755) (1.09684,0.40587) (1.10680,0.41399) (1.11691,0.42191) (1.12717,0.42961) (1.13756,0.43710) (1.14809,0.44436) (1.15874,0.45138) (1.16950,0.45816) (1.18037,0.46469) (1.19135,0.47095) (1.20242,0.47695) (1.21357,0.48267) (1.22481,0.48810) (1.23611,0.49325) (1.24748,0.49809) (1.25890,0.50262) (1.27037,0.50683) (1.28188,0.51073) (1.29342,0.51429) (1.30499,0.51751) (1.31657,0.52039) (1.32816,0.52291) (1.33975,0.52508) (1.35133,0.52688) (1.36289,0.52831) (1.37443,0.52936) (1.38594,0.53004) (1.39741,0.53032) (1.40884,0.53021) (1.42020,0.52971) (1.43151,0.52880) (1.44274,0.52749) (1.45390,0.52577) (1.46496,0.52365) (1.47594,0.52110) (1.48681,0.51815) (1.49758,0.51477) (1.50823,0.51098) (1.51875,0.50677) (1.52915,0.50214) (1.53940,0.49710) (1.54952,0.49163) (1.55948,0.48575) (1.56928,0.47945) (1.57892,0.47274) (1.58838,0.46562) (1.59767,0.45809) (1.60677,0.45016) (1.61568,0.44183) (1.62439,0.43311) (1.63290,0.42400) (1.64121,0.41451) (1.64929,0.40464) (1.65716,0.39440) (1.66480,0.38380) (1.67221,0.37285) (1.67939,0.36155) (1.68632,0.34992) (1.69301,0.33795) (1.69945,0.32567) (1.70564,0.31309) (1.71156,0.30020) (1.71723,0.28703) (1.72262,0.27359) (1.72775,0.25989) (1.73260,0.24593) (1.73718,0.23174) (1.74147,0.21733) (1.74548,0.20270) (1.74921,0.18788) (1.75264,0.17288) (1.75579,0.15771) (1.75864,0.14238) (1.76120,0.12691) (1.76346,0.11132) (1.76542,0.09562) (1.76709,0.07983) (1.76845,0.06396) (1.76951,0.04802) (1.77027,0.03204) (1.77072,0.01603) (1.77088,0.00000)};
\draw[flow,line width=1pt] plot coordinates{(1.97990,-0.00000) (1.97956,-0.03272) (1.97854,-0.06539) (1.97685,-0.09796) (1.97448,-0.13039) (1.97143,-0.16263) (1.96771,-0.19464) (1.96332,-0.22636) (1.95827,-0.25776) (1.95255,-0.28878) (1.94617,-0.31939) (1.93913,-0.34953) (1.93145,-0.37917) (1.92312,-0.40827) (1.91415,-0.43679) (1.90454,-0.46468) (1.89431,-0.49191) (1.88346,-0.51844) (1.87200,-0.54424) (1.85993,-0.56927) (1.84727,-0.59351) (1.83402,-0.61691) (1.82019,-0.63946) (1.80579,-0.66113) (1.79084,-0.68189) (1.77533,-0.70171) (1.75928,-0.72058) (1.74271,-0.73847) (1.72563,-0.75537) (1.70803,-0.77127) (1.68995,-0.78614) (1.67139,-0.79997) (1.65235,-0.81277) (1.63287,-0.82451) (1.61294,-0.83519) (1.59259,-0.84481) (1.57183,-0.85336) (1.55066,-0.86085) (1.52911,-0.86727) (1.50720,-0.87264) (1.48492,-0.87695) (1.46231,-0.88022) (1.43938,-0.88244) (1.41613,-0.88365) (1.39260,-0.88383) (1.36879,-0.88302) (1.34472,-0.88123) (1.32040,-0.87847) (1.29586,-0.87477) (1.27111,-0.87014) (1.24617,-0.86461) (1.22105,-0.85820) (1.19577,-0.85094) (1.17035,-0.84285) (1.14481,-0.83396) (1.11916,-0.82429) (1.09343,-0.81389) (1.06762,-0.80277) (1.04176,-0.79098) (1.01586,-0.77853) (0.98995,-0.76547) (0.96404,-0.75182) (0.93814,-0.73762) (0.91228,-0.72292) (0.88647,-0.70773) (0.86074,-0.69209) (0.83509,-0.67605) (0.80955,-0.65963) (0.78413,-0.64287) (0.75885,-0.62580) (0.73373,-0.60847) (0.70879,-0.59090) (0.68404,-0.57314) (0.65950,-0.55521) (0.63518,-0.53715) (0.61111,-0.51899) (0.58730,-0.50077) (0.56377,-0.48252) (0.54052,-0.46428) (0.51759,-0.44606) (0.49497,-0.42791) (0.47270,-0.40985) (0.45078,-0.39192) (0.42924,-0.37413) (0.40807,-0.35653) (0.38731,-0.33913) (0.36695,-0.32196) (0.34703,-0.30505) (0.32754,-0.28842) (0.30851,-0.27209) (0.28995,-0.25609) (0.27187,-0.24044) (0.25427,-0.22515) (0.23719,-0.21025) (0.22061,-0.19575) (0.20457,-0.18167) (0.18906,-0.16804) (0.17411,-0.15485) (0.15971,-0.14213) (0.14588,-0.12989) (0.13263,-0.11815) (0.11996,-0.10691) (0.10790,-0.09619) (0.09644,-0.08600) (0.08559,-0.07634) (0.07536,-0.06723) (0.06575,-0.05867) (0.05678,-0.05068) (0.04845,-0.04325) (0.04077,-0.03639) (0.03373,-0.03011) (0.02735,-0.02442) (0.02163,-0.01931) (0.01658,-0.01480) (0.01219,-0.01088) (0.00847,-0.00756) (0.00542,-0.00484) (0.00305,-0.00272) (0.00136,-0.00121) (0.00034,-0.00030) (0.00000,-0.00000) (0.00034,0.00030) (0.00136,0.00121) (0.00305,0.00272) (0.00542,0.00484) (0.00847,0.00756) (0.01219,0.01088) (0.01658,0.01480) (0.02163,0.01931) (0.02735,0.02442) (0.03373,0.03011) (0.04077,0.03639) (0.04845,0.04325) (0.05678,0.05068) (0.06575,0.05867) (0.07536,0.06723) (0.08559,0.07634) (0.09644,0.08600) (0.10790,0.09619) (0.11996,0.10691) (0.13263,0.11815) (0.14588,0.12989) (0.15971,0.14213) (0.17411,0.15485) (0.18906,0.16804) (0.20457,0.18167) (0.22061,0.19575) (0.23719,0.21025) (0.25427,0.22515) (0.27187,0.24044) (0.28995,0.25609) (0.30851,0.27209) (0.32754,0.28842) (0.34703,0.30505) (0.36695,0.32196) (0.38731,0.33913) (0.40807,0.35653) (0.42924,0.37413) (0.45078,0.39192) (0.47270,0.40985) (0.49497,0.42791) (0.51759,0.44606) (0.54052,0.46428) (0.56377,0.48252) (0.58730,0.50077) (0.61111,0.51899) (0.63518,0.53715) (0.65950,0.55521) (0.68404,0.57314) (0.70879,0.59090) (0.73373,0.60847) (0.75885,0.62580) (0.78413,0.64287) (0.80955,0.65963) (0.83509,0.67605) (0.86074,0.69209) (0.88647,0.70773) (0.91228,0.72292) (0.93814,0.73762) (0.96404,0.75182) (0.98995,0.76547) (1.01586,0.77853) (1.04176,0.79098) (1.06762,0.80277) (1.09343,0.81389) (1.11916,0.82429) (1.14481,0.83396) (1.17035,0.84285) (1.19577,0.85094) (1.22105,0.85820) (1.24617,0.86461) (1.27111,0.87014) (1.29586,0.87477) (1.32040,0.87847) (1.34472,0.88123) (1.36879,0.88302) (1.39260,0.88383) (1.41613,0.88365) (1.43938,0.88244) (1.46231,0.88022) (1.48492,0.87695) (1.50720,0.87264) (1.52911,0.86727) (1.55066,0.86085) (1.57183,0.85336) (1.59259,0.84481) (1.61294,0.83519) (1.63287,0.82451) (1.65235,0.81277) (1.67139,0.79997) (1.68995,0.78614) (1.70803,0.77127) (1.72563,0.75537) (1.74271,0.73847) (1.75928,0.72058) (1.77533,0.70171) (1.79084,0.68189) (1.80579,0.66113) (1.82019,0.63946) (1.83402,0.61691) (1.84727,0.59351) (1.85993,0.56927) (1.87200,0.54424) (1.88346,0.51844) (1.89431,0.49191) (1.90454,0.46468) (1.91415,0.43679) (1.92312,0.40827) (1.93145,0.37917) (1.93913,0.34953) (1.94617,0.31939) (1.95255,0.28878) (1.95827,0.25776) (1.96332,0.22636) (1.96771,0.19464) (1.97143,0.16263) (1.97448,0.13039) (1.97685,0.09796) (1.97854,0.06539) (1.97956,0.03272) (1.97990,0.00000)};
\draw[flow,dashed] plot coordinates{(2.10695,-0.00000) (2.10623,-0.05537) (2.10406,-0.11059) (2.10045,-0.16554) (2.09540,-0.22007) (2.08892,-0.27405) (2.08101,-0.32733) (2.07167,-0.37979) (2.06091,-0.43129) (2.04873,-0.48171) (2.03515,-0.53092) (2.02018,-0.57880) (2.00383,-0.62523) (1.98610,-0.67010) (1.96700,-0.71330) (1.94657,-0.75472) (1.92479,-0.79427) (1.90170,-0.83186) (1.87730,-0.86740) (1.85162,-0.90080) (1.82467,-0.93200) (1.79647,-0.96093) (1.76703,-0.98753) (1.73639,-1.01174) (1.70456,-1.03353) (1.67155,-1.05286) (1.63741,-1.06970) (1.60214,-1.08403) (1.56577,-1.09583) (1.52833,-1.10512) (1.48984,-1.11189) (1.45033,-1.11616) (1.40982,-1.11796) (1.36835,-1.11734) (1.32594,-1.11432) (1.28263,-1.10898) (1.23843,-1.10138) (1.19339,-1.09161) (1.14753,-1.07975) (1.10088,-1.06592) (1.05347,-1.05024) (1.00535,-1.03283) (0.95653,-1.01385) (0.90706,-0.99347) (0.85697,-0.97187) (0.80629,-0.94924) (0.75506,-0.92582) (0.70331,-0.90184) (0.65108,-0.87758) (0.59841,-0.85331) (0.54532,-0.82935) (0.49186,-0.80604) (0.43806,-0.78373) (0.38396,-0.76280) (0.32960,-0.74361) (0.27501,-0.72655) (0.22024,-0.71200) (0.16531,-0.70027) (0.11027,-0.69167) (0.05515,-0.68642) (0.00000,-0.68465) (-0.05515,-0.68642) (-0.11027,-0.69167) (-0.16531,-0.70027) (-0.22024,-0.71200) (-0.27501,-0.72655) (-0.32960,-0.74361) (-0.38396,-0.76280) (-0.43806,-0.78373) (-0.49186,-0.80604) (-0.54532,-0.82935) (-0.59841,-0.85331) (-0.65108,-0.87758) (-0.70331,-0.90184) (-0.75506,-0.92582) (-0.80629,-0.94924) (-0.85697,-0.97187) (-0.90706,-0.99347) (-0.95653,-1.01385) (-1.00535,-1.03283) (-1.05347,-1.05024) (-1.10088,-1.06592) (-1.14753,-1.07975) (-1.19339,-1.09161) (-1.23843,-1.10138) (-1.28263,-1.10898) (-1.32594,-1.11432) (-1.36835,-1.11734) (-1.40982,-1.11796) (-1.45033,-1.11616) (-1.48984,-1.11189) (-1.52833,-1.10512) (-1.56577,-1.09583) (-1.60214,-1.08403) (-1.63741,-1.06970) (-1.67155,-1.05286) (-1.70456,-1.03353) (-1.73639,-1.01174) (-1.76703,-0.98753) (-1.79647,-0.96093) (-1.82467,-0.93200) (-1.85162,-0.90080) (-1.87730,-0.86740) (-1.90170,-0.83186) (-1.92479,-0.79427) (-1.94657,-0.75472) (-1.96700,-0.71330) (-1.98610,-0.67010) (-2.00383,-0.62523) (-2.02018,-0.57880) (-2.03515,-0.53092) (-2.04873,-0.48171) (-2.06091,-0.43129) (-2.07167,-0.37979) (-2.08101,-0.32733) (-2.08892,-0.27405) (-2.09540,-0.22007) (-2.10045,-0.16554) (-2.10406,-0.11059) (-2.10623,-0.05537) (-2.10695,-0.00000) (-2.10623,0.05537) (-2.10406,0.11059) (-2.10045,0.16554) (-2.09540,0.22007) (-2.08892,0.27405) (-2.08101,0.32733) (-2.07167,0.37979) (-2.06091,0.43129) (-2.04873,0.48171) (-2.03515,0.53092) (-2.02018,0.57880) (-2.00383,0.62523) (-1.98610,0.67010) (-1.96700,0.71330) (-1.94657,0.75472) (-1.92479,0.79427) (-1.90170,0.83186) (-1.87730,0.86740) (-1.85162,0.90080) (-1.82467,0.93200) (-1.79647,0.96093) (-1.76703,0.98753) (-1.73639,1.01174) (-1.70456,1.03353) (-1.67155,1.05286) (-1.63741,1.06970) (-1.60214,1.08403) (-1.56577,1.09583) (-1.52833,1.10512) (-1.48984,1.11189) (-1.45033,1.11616) (-1.40982,1.11796) (-1.36835,1.11734) (-1.32594,1.11432) (-1.28263,1.10898) (-1.23843,1.10138) (-1.19339,1.09161) (-1.14753,1.07975) (-1.10088,1.06592) (-1.05347,1.05024) (-1.00535,1.03283) (-0.95653,1.01385) (-0.90706,0.99347) (-0.85697,0.97187) (-0.80629,0.94924) (-0.75506,0.92582) (-0.70331,0.90184) (-0.65108,0.87758) (-0.59841,0.85331) (-0.54532,0.82935) (-0.49186,0.80604) (-0.43806,0.78373) (-0.38396,0.76280) (-0.32960,0.74361) (-0.27501,0.72655) (-0.22024,0.71200) (-0.16531,0.70027) (-0.11027,0.69167) (-0.05515,0.68642) (-0.00000,0.68465) (0.05515,0.68642) (0.11027,0.69167) (0.16531,0.70027) (0.22024,0.71200) (0.27501,0.72655) (0.32960,0.74361) (0.38396,0.76280) (0.43806,0.78373) (0.49186,0.80604) (0.54532,0.82935) (0.59841,0.85331) (0.65108,0.87758) (0.70331,0.90184) (0.75506,0.92582) (0.80629,0.94924) (0.85697,0.97187) (0.90706,0.99347) (0.95653,1.01385) (1.00535,1.03283) (1.05347,1.05024) (1.10088,1.06592) (1.14753,1.07975) (1.19339,1.09161) (1.23843,1.10138) (1.28263,1.10898) (1.32594,1.11432) (1.36835,1.11734) (1.40982,1.11796) (1.45033,1.11616) (1.48984,1.11189) (1.52833,1.10512) (1.56577,1.09583) (1.60214,1.08403) (1.63741,1.06970) (1.67155,1.05286) (1.70456,1.03353) (1.73639,1.01174) (1.76703,0.98753) (1.79647,0.96093) (1.82467,0.93200) (1.85162,0.90080) (1.87730,0.86740) (1.90170,0.83186) (1.92479,0.79427) (1.94657,0.75472) (1.96700,0.71330) (1.98610,0.67010) (2.00383,0.62523) (2.02018,0.57880) (2.03515,0.53092) (2.04873,0.48171) (2.06091,0.43129) (2.07167,0.37979) (2.08101,0.32733) (2.08892,0.27405) (2.09540,0.22007) (2.10045,0.16554) (2.10406,0.11059) (2.10623,0.05537) (2.10695,0.00000)};
\fill(0,0)circle(1.2pt);\node at(.28,-.48){$S$};
\fill(1.4,0)circle(1.2pt)node[above]{$C_+$};
\fill(-1.4,0)circle(1.2pt)node[above]{$C_-$};
\node[below]at(0,-1.5){粗实线：$h=1/4$；虚线：$h=2/5$};
\end{tikzpicture}
\caption{双阱相图。细实线分别在两个阱中闭合，粗实线是分界轨线，虚线代表高于势垒的跨阱运动。}\label{mp4:fig:l4-e5-s}
\end{figure}

若 $0<h<1/4$，转向条件 $(q^2-1)^2=4h$ 给出右阱的内、外端点
\[
q_{\rm in}=\sqrt{1-2\sqrt h},\qquad
q_{\rm out}=\sqrt{1+2\sqrt h}.
\]
右阱只能在两端点之间往返，左阱由对称性得到另一条闭合轨线。中央区间不满足 $U\le E$，两条轨线不能相通。$h=0$ 时只剩两个静止平衡态。

若 $h>1/4$，势垒处速度满足 $w^2(0)=2h-1/2>0$。只有一对外转向点 $q=\pm\sqrt{1+2\sqrt h}$，相轨线同时包围两中心，上、下支分别向右、向左跨越势垒。

在临界能量 $h=1/4$，
\[
w^2=q^2\left(1-\frac{q^2}{2}\right).
\]
右侧环满足 $0<q\le\sqrt2$。取 $\tau=\omega_bt$，以最右转向时刻为 $\tau=0$。在返回支上，$\dd q/\dd\tau=-q\sqrt{1-q^2/2}$。令$q=\sqrt2\operatorname{sech}z$，$z\ge0$，则
\[
\dd q=-q\tanh z\,\dd z,\qquad
\sqrt{1-q^2/2}=\tanh z,
\]
分离变量后有$\dd\tau=\dd z$。转向点对应$z=0$，故$z=\tau$；再利用运动的时间反演对称性延伸到$\tau<0$，得到
\[
q=\sqrt2\operatorname{sech}\tau,\qquad
w=-\sqrt2\operatorname{sech}\tau\tanh\tau.
\]
对第一式求导得到第二式，将 $\operatorname{sech}^2\tau+\tanh^2\tau=1$ 代入即可验证能量关系。$\tau<0$ 时在上支向右，$\tau>0$ 时在下支向左，两个时间极限都趋向 S。

\begin{figure}[H]\centering
\begin{tikzpicture}[x=1cm,y=1.2cm]
\draw[axis](-3.5,0)--(3.6,0)node[right]{$\tau$};
\draw[axis](0,-.3)--(0,1.9)node[above]{$q$};
\draw[curve,domain=-3.3:3.3,samples=120]plot(\x,{sqrt(2)/cosh(\x)});
\fill(0,{sqrt(2)})circle(1.3pt)node[above right]{$q=\sqrt2$};
\draw[axis](-1.6,{sqrt(2)/cosh(-1.6)})--(-1.3,{sqrt(2)/cosh(-1.3)});
\draw[axis](1.3,{sqrt(2)/cosh(1.3)})--(1.6,{sqrt(2)/cosh(1.6)});
\node at(-2.35,1.35){$w>0$};\node at(2.35,1.35){$w<0$};
\node[below]at(-2.5,0){$\tau\to-\infty$};
\node[below]at(2.5,0){$\tau\to+\infty$};
\end{tikzpicture}
\caption{右侧同宿环的时间参数。质点从鞍点附近出发，转向后再趋近鞍点；曲线两端的 $q=0$ 都只在无限时间达到。}\label{mp4:fig:l4-e5-time}
\end{figure}

图\ref{mp4:fig:l4-e5-time}把同一环展开为时间曲线。接近 $q=0$ 时 $|w|\sim q$。固定一个小的 $q_0>0$，趋近鞍点所需的时间含有
\[
\frac1{\omega_b}\int_\epsilon^{q_0}\frac{\dd q}{q}
=\frac1{\omega_b}[\ln q]_\epsilon^{q_0}
=\frac1{\omega_b}\ln\frac{q_0}\epsilon
\longrightarrow\infty\qquad(\epsilon\to0^+).
\]
右、左两环虽然在相图中共用 S，实际运动不能在有限时间内由一个环穿过 S 进入另一个环。

从 $x=a$ 向左出发，$U(a)=0$，初态能量为 $mv_0^2/2$。进入左阱要求在 $x=0$ 仍有负速度，即
\[
\frac12mv_0^2>\frac{\lambda a^4}{4},
\qquad |v_0|>\frac{a\omega_b}{\sqrt2},\quad v_0<0.
\]
等号的能量足以使允许区间接触势垒顶，但到达该点需要无限时间。因此临界数值是可跨阱初速大小的下确界，实际跨阱条件必须取严格不等号。
\end{solution}

\par\medskip\noindent\begin{minipage}{\textwidth}
\begin{example}[单摆的振动、分界与转动]\label{mp4:example:l4-e6}
无质量刚杆长为 $l$，一端铰接，另一端固定质量为 $m$ 的小球，构成刚杆单摆。刚杆可承受拉力和压力。摆角 $\theta$ 从竖直向下方向起算，正向见图\ref{mp4:fig:l4-e6-q}，$\omega_0=\sqrt{g/l}$。用 $p=\dot\theta/\omega_0$ 作相图，求连续向同一方向转动的能量条件，并推导振动周期积分与接近分界轨线时的变化。
\end{example}

\begin{figure}[H]\centering
\begin{tikzpicture}
\fill(0,1.3)circle(2pt)node[above]{$O$};\draw[guide](0,1.3)--(0,-1.4);\draw[curve](0,1.3)--(1.2,-.7);\fill(1.2,-.7)circle(3pt)node[right]{$m$};\node[above right]at(.6,.3){$l$};\draw[axis](0,.5)arc(-90:-59:.8);\node at(.45,.18){$\theta$};\draw[axis](1.2,-.7)--(1.2,-1.6)node[right]{$mg$};
\end{tikzpicture}
\caption{刚杆单摆。角度从向下竖直线转到刚杆，沿所示箭头为正。}\label{mp4:fig:l4-e6-q}
\end{figure}
\end{minipage}\par\medskip
\begin{solution}
图\ref{mp4:fig:l4-e6-q}中的刚杆使摆球始终保持到悬点的距离为 $l$，不需要另加绳张力非负的约束。以最低点为零势能，
\[
E=\frac12ml^2\dot\theta^2+mgl(1-\cos\theta).
\]
令 $p=\dot\theta/\omega_0$、$\omega_0=\sqrt{g/l}$，能量除以 $mgl$ 后为
\begin{equation}\label{mp4:eq:l4-pendulum}
h=\frac12p^2+1-\cos\theta.
\end{equation}
由 $\ddot\theta=-\omega_0^2\sin\theta$，最低点附近是简谐振动，倒立点附近则有指数离开或趋近的运动。图\ref{mp4:fig:l4-e6-s}中的 C 为中心，$\theta=\pm\pi,p=0$ 为展开相平面上的相邻鞍点。

\begin{figure}[H]\centering
\begin{tikzpicture}
\draw[axis](-2.8,0)--(2.8,0)node[right]{$\theta$};
\draw[axis](0,-2.2)--(0,2.3)node[above]{$p$};
\draw[flow,line width=.4pt] plot coordinates{(0.83468,-0.00000) (0.83440,-0.01943) (0.83354,-0.03885) (0.83211,-0.05825) (0.83011,-0.07763) (0.82754,-0.09697) (0.82440,-0.11627) (0.82070,-0.13553) (0.81644,-0.15472) (0.81162,-0.17384) (0.80624,-0.19288) (0.80031,-0.21184) (0.79383,-0.23069) (0.78681,-0.24945) (0.77924,-0.26808) (0.77114,-0.28659) (0.76252,-0.30495) (0.75337,-0.32318) (0.74371,-0.34124) (0.73353,-0.35912) (0.72286,-0.37683) (0.71168,-0.39434) (0.70002,-0.41165) (0.68788,-0.42873) (0.67527,-0.44559) (0.66220,-0.46220) (0.64867,-0.47855) (0.63470,-0.49463) (0.62029,-0.51043) (0.60546,-0.52593) (0.59021,-0.54112) (0.57456,-0.55598) (0.55851,-0.57051) (0.54208,-0.58468) (0.52528,-0.59848) (0.50812,-0.61191) (0.49061,-0.62494) (0.47277,-0.63756) (0.45460,-0.64976) (0.43612,-0.66153) (0.41734,-0.67285) (0.39828,-0.68371) (0.37894,-0.69410) (0.35934,-0.70400) (0.33950,-0.71340) (0.31942,-0.72230) (0.29912,-0.73067) (0.27862,-0.73852) (0.25793,-0.74583) (0.23706,-0.75259) (0.21603,-0.75879) (0.19485,-0.76442) (0.17354,-0.76949) (0.15211,-0.77397) (0.13057,-0.77787) (0.10895,-0.78117) (0.08725,-0.78389) (0.06549,-0.78600) (0.04368,-0.78751) (0.02185,-0.78842) (0.00000,-0.78872) (-0.02185,-0.78842) (-0.04368,-0.78751) (-0.06549,-0.78600) (-0.08725,-0.78389) (-0.10895,-0.78117) (-0.13057,-0.77787) (-0.15211,-0.77397) (-0.17354,-0.76949) (-0.19485,-0.76442) (-0.21603,-0.75879) (-0.23706,-0.75259) (-0.25793,-0.74583) (-0.27862,-0.73852) (-0.29912,-0.73067) (-0.31942,-0.72230) (-0.33950,-0.71340) (-0.35934,-0.70400) (-0.37894,-0.69410) (-0.39828,-0.68371) (-0.41734,-0.67285) (-0.43612,-0.66153) (-0.45460,-0.64976) (-0.47277,-0.63756) (-0.49061,-0.62494) (-0.50812,-0.61191) (-0.52528,-0.59848) (-0.54208,-0.58468) (-0.55851,-0.57051) (-0.57456,-0.55598) (-0.59021,-0.54112) (-0.60546,-0.52593) (-0.62029,-0.51043) (-0.63470,-0.49463) (-0.64867,-0.47855) (-0.66220,-0.46220) (-0.67527,-0.44559) (-0.68788,-0.42873) (-0.70002,-0.41165) (-0.71168,-0.39434) (-0.72286,-0.37683) (-0.73353,-0.35912) (-0.74371,-0.34124) (-0.75337,-0.32318) (-0.76252,-0.30495) (-0.77114,-0.28659) (-0.77924,-0.26808) (-0.78681,-0.24945) (-0.79383,-0.23069) (-0.80031,-0.21184) (-0.80624,-0.19288) (-0.81162,-0.17384) (-0.81644,-0.15472) (-0.82070,-0.13553) (-0.82440,-0.11627) (-0.82754,-0.09697) (-0.83011,-0.07763) (-0.83211,-0.05825) (-0.83354,-0.03885) (-0.83440,-0.01943) (-0.83468,-0.00000) (-0.83440,0.01943) (-0.83354,0.03885) (-0.83211,0.05825) (-0.83011,0.07763) (-0.82754,0.09697) (-0.82440,0.11627) (-0.82070,0.13553) (-0.81644,0.15472) (-0.81162,0.17384) (-0.80624,0.19288) (-0.80031,0.21184) (-0.79383,0.23069) (-0.78681,0.24945) (-0.77924,0.26808) (-0.77114,0.28659) (-0.76252,0.30495) (-0.75337,0.32318) (-0.74371,0.34124) (-0.73353,0.35912) (-0.72286,0.37683) (-0.71168,0.39434) (-0.70002,0.41165) (-0.68788,0.42873) (-0.67527,0.44559) (-0.66220,0.46220) (-0.64867,0.47855) (-0.63470,0.49463) (-0.62029,0.51043) (-0.60546,0.52593) (-0.59021,0.54112) (-0.57456,0.55598) (-0.55851,0.57051) (-0.54208,0.58468) (-0.52528,0.59848) (-0.50812,0.61191) (-0.49061,0.62494) (-0.47277,0.63756) (-0.45460,0.64976) (-0.43612,0.66153) (-0.41734,0.67285) (-0.39828,0.68371) (-0.37894,0.69410) (-0.35934,0.70400) (-0.33950,0.71340) (-0.31942,0.72230) (-0.29912,0.73067) (-0.27862,0.73852) (-0.25793,0.74583) (-0.23706,0.75259) (-0.21603,0.75879) (-0.19485,0.76442) (-0.17354,0.76949) (-0.15211,0.77397) (-0.13057,0.77787) (-0.10895,0.78117) (-0.08725,0.78389) (-0.06549,0.78600) (-0.04368,0.78751) (-0.02185,0.78842) (-0.00000,0.78872) (0.02185,0.78842) (0.04368,0.78751) (0.06549,0.78600) (0.08725,0.78389) (0.10895,0.78117) (0.13057,0.77787) (0.15211,0.77397) (0.17354,0.76949) (0.19485,0.76442) (0.21603,0.75879) (0.23706,0.75259) (0.25793,0.74583) (0.27862,0.73852) (0.29912,0.73067) (0.31942,0.72230) (0.33950,0.71340) (0.35934,0.70400) (0.37894,0.69410) (0.39828,0.68371) (0.41734,0.67285) (0.43612,0.66153) (0.45460,0.64976) (0.47277,0.63756) (0.49061,0.62494) (0.50812,0.61191) (0.52528,0.59848) (0.54208,0.58468) (0.55851,0.57051) (0.57456,0.55598) (0.59021,0.54112) (0.60546,0.52593) (0.62029,0.51043) (0.63470,0.49463) (0.64867,0.47855) (0.66220,0.46220) (0.67527,0.44559) (0.68788,0.42873) (0.70002,0.41165) (0.71168,0.39434) (0.72286,0.37683) (0.73353,0.35912) (0.74371,0.34124) (0.75337,0.32318) (0.76252,0.30495) (0.77114,0.28659) (0.77924,0.26808) (0.78681,0.24945) (0.79383,0.23069) (0.80031,0.21184) (0.80624,0.19288) (0.81162,0.17384) (0.81644,0.15472) (0.82070,0.13553) (0.82440,0.11627) (0.82754,0.09697) (0.83011,0.07763) (0.83211,0.05825) (0.83354,0.03885) (0.83440,0.01943) (0.83468,0.00000)};
\draw[flow,line width=.4pt] plot coordinates{(1.50796,-0.00000) (1.50745,-0.02539) (1.50590,-0.05079) (1.50332,-0.07621) (1.49970,-0.10166) (1.49506,-0.12716) (1.48940,-0.15272) (1.48271,-0.17833) (1.47501,-0.20402) (1.46630,-0.22978) (1.45658,-0.25562) (1.44587,-0.28155) (1.43416,-0.30757) (1.42147,-0.33368) (1.40781,-0.35988) (1.39318,-0.38616) (1.37759,-0.41252) (1.36107,-0.43895) (1.34361,-0.46545) (1.32522,-0.49199) (1.30594,-0.51858) (1.28575,-0.54519) (1.26469,-0.57180) (1.24275,-0.59840) (1.21997,-0.62496) (1.19635,-0.65145) (1.17191,-0.67785) (1.14667,-0.70413) (1.12064,-0.73026) (1.09384,-0.75621) (1.06629,-0.78193) (1.03801,-0.80740) (1.00903,-0.83257) (0.97934,-0.85740) (0.94899,-0.88186) (0.91799,-0.90589) (0.88636,-0.92947) (0.85412,-0.95253) (0.82130,-0.97504) (0.78791,-0.99696) (0.75398,-1.01823) (0.71954,-1.03882) (0.68460,-1.05867) (0.64920,-1.07775) (0.61334,-1.09600) (0.57707,-1.11339) (0.54041,-1.12987) (0.50337,-1.14541) (0.46599,-1.15997) (0.42829,-1.17351) (0.39029,-1.18599) (0.35203,-1.19738) (0.31352,-1.20766) (0.27480,-1.21680) (0.23590,-1.22476) (0.19683,-1.23154) (0.15763,-1.23711) (0.11831,-1.24146) (0.07892,-1.24458) (0.03947,-1.24645) (0.00000,-1.24708) (-0.03947,-1.24645) (-0.07892,-1.24458) (-0.11831,-1.24146) (-0.15763,-1.23711) (-0.19683,-1.23154) (-0.23590,-1.22476) (-0.27480,-1.21680) (-0.31352,-1.20766) (-0.35203,-1.19738) (-0.39029,-1.18599) (-0.42829,-1.17351) (-0.46599,-1.15997) (-0.50337,-1.14541) (-0.54041,-1.12987) (-0.57707,-1.11339) (-0.61334,-1.09600) (-0.64920,-1.07775) (-0.68460,-1.05867) (-0.71954,-1.03882) (-0.75398,-1.01823) (-0.78791,-0.99696) (-0.82130,-0.97504) (-0.85412,-0.95253) (-0.88636,-0.92947) (-0.91799,-0.90589) (-0.94899,-0.88186) (-0.97934,-0.85740) (-1.00903,-0.83257) (-1.03801,-0.80740) (-1.06629,-0.78193) (-1.09384,-0.75621) (-1.12064,-0.73026) (-1.14667,-0.70413) (-1.17191,-0.67785) (-1.19635,-0.65145) (-1.21997,-0.62496) (-1.24275,-0.59840) (-1.26469,-0.57180) (-1.28575,-0.54519) (-1.30594,-0.51858) (-1.32522,-0.49199) (-1.34361,-0.46545) (-1.36107,-0.43895) (-1.37759,-0.41252) (-1.39318,-0.38616) (-1.40781,-0.35988) (-1.42147,-0.33368) (-1.43416,-0.30757) (-1.44587,-0.28155) (-1.45658,-0.25562) (-1.46630,-0.22978) (-1.47501,-0.20402) (-1.48271,-0.17833) (-1.48940,-0.15272) (-1.49506,-0.12716) (-1.49970,-0.10166) (-1.50332,-0.07621) (-1.50590,-0.05079) (-1.50745,-0.02539) (-1.50796,-0.00000) (-1.50745,0.02539) (-1.50590,0.05079) (-1.50332,0.07621) (-1.49970,0.10166) (-1.49506,0.12716) (-1.48940,0.15272) (-1.48271,0.17833) (-1.47501,0.20402) (-1.46630,0.22978) (-1.45658,0.25562) (-1.44587,0.28155) (-1.43416,0.30757) (-1.42147,0.33368) (-1.40781,0.35988) (-1.39318,0.38616) (-1.37759,0.41252) (-1.36107,0.43895) (-1.34361,0.46545) (-1.32522,0.49199) (-1.30594,0.51858) (-1.28575,0.54519) (-1.26469,0.57180) (-1.24275,0.59840) (-1.21997,0.62496) (-1.19635,0.65145) (-1.17191,0.67785) (-1.14667,0.70413) (-1.12064,0.73026) (-1.09384,0.75621) (-1.06629,0.78193) (-1.03801,0.80740) (-1.00903,0.83257) (-0.97934,0.85740) (-0.94899,0.88186) (-0.91799,0.90589) (-0.88636,0.92947) (-0.85412,0.95253) (-0.82130,0.97504) (-0.78791,0.99696) (-0.75398,1.01823) (-0.71954,1.03882) (-0.68460,1.05867) (-0.64920,1.07775) (-0.61334,1.09600) (-0.57707,1.11339) (-0.54041,1.12987) (-0.50337,1.14541) (-0.46599,1.15997) (-0.42829,1.17351) (-0.39029,1.18599) (-0.35203,1.19738) (-0.31352,1.20766) (-0.27480,1.21680) (-0.23590,1.22476) (-0.19683,1.23154) (-0.15763,1.23711) (-0.11831,1.24146) (-0.07892,1.24458) (-0.03947,1.24645) (-0.00000,1.24708) (0.03947,1.24645) (0.07892,1.24458) (0.11831,1.24146) (0.15763,1.23711) (0.19683,1.23154) (0.23590,1.22476) (0.27480,1.21680) (0.31352,1.20766) (0.35203,1.19738) (0.39029,1.18599) (0.42829,1.17351) (0.46599,1.15997) (0.50337,1.14541) (0.54041,1.12987) (0.57707,1.11339) (0.61334,1.09600) (0.64920,1.07775) (0.68460,1.05867) (0.71954,1.03882) (0.75398,1.01823) (0.78791,0.99696) (0.82130,0.97504) (0.85412,0.95253) (0.88636,0.92947) (0.91799,0.90589) (0.94899,0.88186) (0.97934,0.85740) (1.00903,0.83257) (1.03801,0.80740) (1.06629,0.78193) (1.09384,0.75621) (1.12064,0.73026) (1.14667,0.70413) (1.17191,0.67785) (1.19635,0.65145) (1.21997,0.62496) (1.24275,0.59840) (1.26469,0.57180) (1.28575,0.54519) (1.30594,0.51858) (1.32522,0.49199) (1.34361,0.46545) (1.36107,0.43895) (1.37759,0.41252) (1.39318,0.38616) (1.40781,0.35988) (1.42147,0.33368) (1.43416,0.30757) (1.44587,0.28155) (1.45658,0.25562) (1.46630,0.22978) (1.47501,0.20402) (1.48271,0.17833) (1.48940,0.15272) (1.49506,0.12716) (1.49970,0.10166) (1.50332,0.07621) (1.50590,0.05079) (1.50745,0.02539) (1.50796,0.00000)};
\draw[flow,line width=1pt] plot coordinates{(-2.26195,0.00000) (-2.21671,0.04523) (-2.17147,0.09042) (-2.12623,0.13552) (-2.08099,0.18048) (-2.03575,0.22527) (-1.99051,0.26983) (-1.94527,0.31413) (-1.90004,0.35811) (-1.85480,0.40175) (-1.80956,0.44498) (-1.76432,0.48778) (-1.71908,0.53010) (-1.67384,0.57189) (-1.62860,0.61312) (-1.58336,0.65375) (-1.53812,0.69373) (-1.49288,0.73302) (-1.44765,0.77159) (-1.40241,0.80940) (-1.35717,0.84641) (-1.31193,0.88259) (-1.26669,0.91789) (-1.22145,0.95229) (-1.17621,0.98575) (-1.13097,1.01823) (-1.08573,1.04971) (-1.04050,1.08016) (-0.99526,1.10954) (-0.95002,1.13782) (-0.90478,1.16498) (-0.85954,1.19100) (-0.81430,1.21583) (-0.76906,1.23947) (-0.72382,1.26188) (-0.67858,1.28305) (-0.63335,1.30295) (-0.58811,1.32157) (-0.54287,1.33888) (-0.49763,1.35487) (-0.45239,1.36952) (-0.40715,1.38282) (-0.36191,1.39476) (-0.31667,1.40532) (-0.27143,1.41449) (-0.22619,1.42227) (-0.18096,1.42865) (-0.13572,1.43361) (-0.09048,1.43716) (-0.04524,1.43929) (0.00000,1.44000) (0.04524,1.43929) (0.09048,1.43716) (0.13572,1.43361) (0.18096,1.42865) (0.22619,1.42227) (0.27143,1.41449) (0.31667,1.40532) (0.36191,1.39476) (0.40715,1.38282) (0.45239,1.36952) (0.49763,1.35487) (0.54287,1.33888) (0.58811,1.32157) (0.63335,1.30295) (0.67858,1.28305) (0.72382,1.26188) (0.76906,1.23947) (0.81430,1.21583) (0.85954,1.19100) (0.90478,1.16498) (0.95002,1.13782) (0.99526,1.10954) (1.04050,1.08016) (1.08573,1.04971) (1.13097,1.01823) (1.17621,0.98575) (1.22145,0.95229) (1.26669,0.91789) (1.31193,0.88259) (1.35717,0.84641) (1.40241,0.80940) (1.44765,0.77159) (1.49288,0.73302) (1.53812,0.69373) (1.58336,0.65375) (1.62860,0.61312) (1.67384,0.57189) (1.71908,0.53010) (1.76432,0.48778) (1.80956,0.44498) (1.85480,0.40175) (1.90004,0.35811) (1.94527,0.31413) (1.99051,0.26983) (2.03575,0.22527) (2.08099,0.18048) (2.12623,0.13552) (2.17147,0.09042) (2.21671,0.04523) (2.26195,0.00000)};
\draw[flow,line width=1pt] plot coordinates{(2.26195,-0.00000) (2.21671,-0.04523) (2.17147,-0.09042) (2.12623,-0.13552) (2.08099,-0.18048) (2.03575,-0.22527) (1.99051,-0.26983) (1.94527,-0.31413) (1.90004,-0.35811) (1.85480,-0.40175) (1.80956,-0.44498) (1.76432,-0.48778) (1.71908,-0.53010) (1.67384,-0.57189) (1.62860,-0.61312) (1.58336,-0.65375) (1.53812,-0.69373) (1.49288,-0.73302) (1.44765,-0.77159) (1.40241,-0.80940) (1.35717,-0.84641) (1.31193,-0.88259) (1.26669,-0.91789) (1.22145,-0.95229) (1.17621,-0.98575) (1.13097,-1.01823) (1.08573,-1.04971) (1.04050,-1.08016) (0.99526,-1.10954) (0.95002,-1.13782) (0.90478,-1.16498) (0.85954,-1.19100) (0.81430,-1.21583) (0.76906,-1.23947) (0.72382,-1.26188) (0.67858,-1.28305) (0.63335,-1.30295) (0.58811,-1.32157) (0.54287,-1.33888) (0.49763,-1.35487) (0.45239,-1.36952) (0.40715,-1.38282) (0.36191,-1.39476) (0.31667,-1.40532) (0.27143,-1.41449) (0.22619,-1.42227) (0.18096,-1.42865) (0.13572,-1.43361) (0.09048,-1.43716) (0.04524,-1.43929) (0.00000,-1.44000) (-0.04524,-1.43929) (-0.09048,-1.43716) (-0.13572,-1.43361) (-0.18096,-1.42865) (-0.22619,-1.42227) (-0.27143,-1.41449) (-0.31667,-1.40532) (-0.36191,-1.39476) (-0.40715,-1.38282) (-0.45239,-1.36952) (-0.49763,-1.35487) (-0.54287,-1.33888) (-0.58811,-1.32157) (-0.63335,-1.30295) (-0.67858,-1.28305) (-0.72382,-1.26188) (-0.76906,-1.23947) (-0.81430,-1.21583) (-0.85954,-1.19100) (-0.90478,-1.16498) (-0.95002,-1.13782) (-0.99526,-1.10954) (-1.04050,-1.08016) (-1.08573,-1.04971) (-1.13097,-1.01823) (-1.17621,-0.98575) (-1.22145,-0.95229) (-1.26669,-0.91789) (-1.31193,-0.88259) (-1.35717,-0.84641) (-1.40241,-0.80940) (-1.44765,-0.77159) (-1.49288,-0.73302) (-1.53812,-0.69373) (-1.58336,-0.65375) (-1.62860,-0.61312) (-1.67384,-0.57189) (-1.71908,-0.53010) (-1.76432,-0.48778) (-1.80956,-0.44498) (-1.85480,-0.40175) (-1.90004,-0.35811) (-1.94527,-0.31413) (-1.99051,-0.26983) (-2.03575,-0.22527) (-2.08099,-0.18048) (-2.12623,-0.13552) (-2.17147,-0.09042) (-2.21671,-0.04523) (-2.26195,-0.00000)};
\draw[flow,dashed] plot coordinates{(-2.26195,1.01823) (-2.21671,1.01924) (-2.17147,1.02224) (-2.12623,1.02721) (-2.08099,1.03410) (-2.03575,1.04285) (-1.99051,1.05338) (-1.94527,1.06559) (-1.90004,1.07937) (-1.85480,1.09462) (-1.80956,1.11122) (-1.76432,1.12904) (-1.71908,1.14796) (-1.67384,1.16784) (-1.62860,1.18858) (-1.58336,1.21003) (-1.53812,1.23209) (-1.49288,1.25464) (-1.44765,1.27756) (-1.40241,1.30074) (-1.35717,1.32409) (-1.31193,1.34750) (-1.26669,1.37088) (-1.22145,1.39415) (-1.17621,1.41722) (-1.13097,1.44000) (-1.08573,1.46243) (-1.04050,1.48443) (-0.99526,1.50595) (-0.95002,1.52691) (-0.90478,1.54725) (-0.85954,1.56693) (-0.81430,1.58589) (-0.76906,1.60408) (-0.72382,1.62146) (-0.67858,1.63799) (-0.63335,1.65363) (-0.58811,1.66833) (-0.54287,1.68208) (-0.49763,1.69484) (-0.45239,1.70657) (-0.40715,1.71727) (-0.36191,1.72689) (-0.31667,1.73543) (-0.27143,1.74287) (-0.22619,1.74919) (-0.18096,1.75437) (-0.13572,1.75842) (-0.09048,1.76131) (-0.04524,1.76305) (0.00000,1.76363) (0.04524,1.76305) (0.09048,1.76131) (0.13572,1.75842) (0.18096,1.75437) (0.22619,1.74919) (0.27143,1.74287) (0.31667,1.73543) (0.36191,1.72689) (0.40715,1.71727) (0.45239,1.70657) (0.49763,1.69484) (0.54287,1.68208) (0.58811,1.66833) (0.63335,1.65363) (0.67858,1.63799) (0.72382,1.62146) (0.76906,1.60408) (0.81430,1.58589) (0.85954,1.56693) (0.90478,1.54725) (0.95002,1.52691) (0.99526,1.50595) (1.04050,1.48443) (1.08573,1.46243) (1.13097,1.44000) (1.17621,1.41722) (1.22145,1.39415) (1.26669,1.37088) (1.31193,1.34750) (1.35717,1.32409) (1.40241,1.30074) (1.44765,1.27756) (1.49288,1.25464) (1.53812,1.23209) (1.58336,1.21003) (1.62860,1.18858) (1.67384,1.16784) (1.71908,1.14796) (1.76432,1.12904) (1.80956,1.11122) (1.85480,1.09462) (1.90004,1.07937) (1.94527,1.06559) (1.99051,1.05338) (2.03575,1.04285) (2.08099,1.03410) (2.12623,1.02721) (2.17147,1.02224) (2.21671,1.01924) (2.26195,1.01823)};
\draw[flow,dashed] plot coordinates{(2.26195,-1.01823) (2.21671,-1.01924) (2.17147,-1.02224) (2.12623,-1.02721) (2.08099,-1.03410) (2.03575,-1.04285) (1.99051,-1.05338) (1.94527,-1.06559) (1.90004,-1.07937) (1.85480,-1.09462) (1.80956,-1.11122) (1.76432,-1.12904) (1.71908,-1.14796) (1.67384,-1.16784) (1.62860,-1.18858) (1.58336,-1.21003) (1.53812,-1.23209) (1.49288,-1.25464) (1.44765,-1.27756) (1.40241,-1.30074) (1.35717,-1.32409) (1.31193,-1.34750) (1.26669,-1.37088) (1.22145,-1.39415) (1.17621,-1.41722) (1.13097,-1.44000) (1.08573,-1.46243) (1.04050,-1.48443) (0.99526,-1.50595) (0.95002,-1.52691) (0.90478,-1.54725) (0.85954,-1.56693) (0.81430,-1.58589) (0.76906,-1.60408) (0.72382,-1.62146) (0.67858,-1.63799) (0.63335,-1.65363) (0.58811,-1.66833) (0.54287,-1.68208) (0.49763,-1.69484) (0.45239,-1.70657) (0.40715,-1.71727) (0.36191,-1.72689) (0.31667,-1.73543) (0.27143,-1.74287) (0.22619,-1.74919) (0.18096,-1.75437) (0.13572,-1.75842) (0.09048,-1.76131) (0.04524,-1.76305) (0.00000,-1.76363) (-0.04524,-1.76305) (-0.09048,-1.76131) (-0.13572,-1.75842) (-0.18096,-1.75437) (-0.22619,-1.74919) (-0.27143,-1.74287) (-0.31667,-1.73543) (-0.36191,-1.72689) (-0.40715,-1.71727) (-0.45239,-1.70657) (-0.49763,-1.69484) (-0.54287,-1.68208) (-0.58811,-1.66833) (-0.63335,-1.65363) (-0.67858,-1.63799) (-0.72382,-1.62146) (-0.76906,-1.60408) (-0.81430,-1.58589) (-0.85954,-1.56693) (-0.90478,-1.54725) (-0.95002,-1.52691) (-0.99526,-1.50595) (-1.04050,-1.48443) (-1.08573,-1.46243) (-1.13097,-1.44000) (-1.17621,-1.41722) (-1.22145,-1.39415) (-1.26669,-1.37088) (-1.31193,-1.34750) (-1.35717,-1.32409) (-1.40241,-1.30074) (-1.44765,-1.27756) (-1.49288,-1.25464) (-1.53812,-1.23209) (-1.58336,-1.21003) (-1.62860,-1.18858) (-1.67384,-1.16784) (-1.71908,-1.14796) (-1.76432,-1.12904) (-1.80956,-1.11122) (-1.85480,-1.09462) (-1.90004,-1.07937) (-1.94527,-1.06559) (-1.99051,-1.05338) (-2.03575,-1.04285) (-2.08099,-1.03410) (-2.12623,-1.02721) (-2.17147,-1.02224) (-2.21671,-1.01924) (-2.26195,-1.01823)};
\draw[guide](-2.261946710584651,-2)--(-2.261946710584651,2);\draw[fill=white](-2.261946710584651,0)circle(1.5pt);\node[below left]at(-2.261946710584651,0){$-\pi$};\draw[guide](2.261946710584651,-2)--(2.261946710584651,2);\draw[fill=white](2.261946710584651,0)circle(1.5pt);\node[below right]at(2.261946710584651,0){$\pi$};\fill(0,0)circle(1.2pt)node[below left]{$C$};
\node at(0,-2.55){粗线：$h=2$；虚线：$h=3$};
\end{tikzpicture}
\caption{单摆的相图。上半平面向右，下半平面向左；左右虚线是相同角度的识别边界，粗线在鞍点处只作渐近连接。}\label{mp4:fig:l4-e6-s}
\end{figure}

$0<h<2$ 时，$p=0$ 有两个相邻解 $\theta=\pm\theta_m$，其中 $h=1-\cos\theta_m$。相点在闭合轨线上顺时针绕行，摆球在两端之间往返。$h>2$ 时，对任何角度都有 $p^2=2(h-1+\cos\theta)>0$，角速度不能变号，因此连续向同一方向转动。若从最低点出发，$h=\dot\theta_0^2/(2\omega_0^2)$，连续转动要求
\[
|\dot\theta_0|>2\omega_0.
\]

临界能量 $h=2$ 在 $-\pi\le\theta\le\pi$ 内给出
\[
p=\pm2\cos\frac\theta2.
\]
取上支，分离变量 $\dd\theta/[2\cos(\theta/2)]=\omega_0\dd t$。令$u=\theta/2$，在$-\pi<\theta<\pi$内，
\[
\int\frac{\dd\theta}{2\cos(\theta/2)}
=\int\sec u\,\dd u
=\ln(\sec u+\tan u)
=\ln\tan\left(\frac\pi4+\frac\theta4\right).
\]
以通过最低点的时刻 $t_0$ 确定积分常量，得到 $\tan(\pi/4+\theta/4)=\ee^{\omega_0(t-t_0)}$。取 $u=\omega_0(t-t_0)$，解出 $\theta$ 后用反正切的链式导数，$\dot\theta/\omega_0=4\ee^u/(1+\ee^{2u})=2\operatorname{sech}u$，故
\[
\theta=4\arctan\ee^{\omega_0(t-t_0)}-\pi,
\qquad p=2\operatorname{sech}[\omega_0(t-t_0)].
\]
两端 $\theta=\mp\pi$ 分别对应 $t\to\mp\infty$，任何有限时刻都没有到达鞍点。在展开平面上这是连接相邻鞍点的异宿轨线；若将 $\theta$ 与 $\theta+2\pi$ 识别，相空间成为圆柱，两端表示同一个倒立状态，此支便是圆柱上的同宿轨线。图中左右虚线表示识别边界。

对振动周期，令 $k=\sin(\theta_m/2)\in(0,1)$。上行支的角速度为
\[
\dot\theta=2\omega_0\sqrt{k^2-\sin^2(\theta/2)}.
\]
从最低点到正转向点的时间占周期四分之一。取 $\sin(\theta/2)=k\sin\psi$，则
\[
\dd\theta=\frac{2k\cos\psi}{\sqrt{1-k^2\sin^2\psi}}\,\dd\psi,
\quad
\sqrt{k^2-\sin^2(\theta/2)}=k\cos\psi.
\]
代入得
\begin{equation}\label{mp4:eq:l4-pend-period}
T=\frac4{\omega_0}\int_0^{\pi/2}
\frac{\dd\psi}{\sqrt{1-k^2\sin^2\psi}}.
\end{equation}
小振幅时，被积函数展开为 $1+(k^2/2)\sin^2\psi+O(k^4)$。利用例\ref{mp4:example:l4-e4}中已算出的 $\int_0^{\pi/2}\sin^2\psi\dd\psi=\pi/4$，积分为 $\pi/2+k^2\pi/8+O(k^4)$。再用 $k=\theta_m/2+O(\theta_m^3)$，得到
\[
T=\frac{2\pi}{\omega_0}
\left(1+\frac{k^2}{4}+O(k^4)\right)
=\frac{2\pi}{\omega_0}
\left(1+\frac{\theta_m^2}{16}+O(\theta_m^4)\right).
\]

接近分界时，令 $k'=\sqrt{1-k^2}=\cos(\theta_m/2)$。记式\eqref{mp4:eq:l4-pend-period}中的积分为 $I(k)$，再取 $u=\pi/2-\psi$，则
\[
I(k)=\int_0^{\pi/2}\frac{\dd u}{\sqrt{\sin^2u+k'^2\cos^2u}}.
\]
选取中间尺度 $k'\ll\delta\ll1$，在近端区间将分母近似为 $\sqrt{u^2+k'^2}$，在远端区间近似为 $\sin u$。两段首项分别为
\[
\int_0^\delta\frac{\dd u}{\sqrt{u^2+k'^2}}
=\left[\ln(u+\sqrt{u^2+k'^2})\right]_0^\delta
=\ln\frac{\delta+\sqrt{\delta^2+k'^2}}{k'}
\simeq\ln\frac{2\delta}{k'},
\]
\[
\int_\delta^{\pi/2}\frac{\dd u}{\sin u}
=\left[\ln\tan\frac u2\right]_\delta^{\pi/2}
=-\ln\tan\frac\delta2\simeq\ln\frac2\delta.
\]
先令 $k'\to0$，再使 $\delta\to0$，匹配后得到 $I(k)=\ln(4/k')+o(1)$，从而
\[
T\sim\frac4{\omega_0}\ln\frac4{k'}.
\]
周期的对数发散来自倒立位置附近的长时间停留。图中分界支在鞍点处只有渐近连接，临界运动的往返时间因而不再有限。
\end{solution}

\par\medskip\noindent\begin{minipage}{\textwidth}
\begin{example}[势阱与势垒的合并]\label{mp4:example:l4-e7}
质量为$m$的质点的势能为 $U(x)=\kappa x^3/3-Fx$，$\kappa>0$。$F>0$ 时势能形状见图\ref{mp4:fig:l4-e7-q}。求中心、鞍点、势垒高度及小振动频率，并讨论 $F$ 减到零时发生的变化。求临界能量下右侧轨线的最右端位置。
\end{example}

\begin{figure}[H]\centering
\begin{tikzpicture}
\draw[axis](-2.7,0)--(3,0)node[right]{$x$};
\draw[axis](0,-1.5)--(0,1.8)node[above]{$U$};
\draw[curve,domain=-2:2.1,samples=120]plot(\x,{.65*((\x)^3/3-\x)});\fill(-1,.4333)circle(1.3pt);\fill(1,-.4333)circle(1.3pt);\draw[guide](-1,.4333)--(2,.4333);\node[above]at(-1,.48){势垒};\node[below]at(1,-.48){势阱};
\end{tikzpicture}
\caption{三次势能的局部势阱。左侧越过势垒后可继续向负 x 方向运动。}\label{mp4:fig:l4-e7-q}
\end{figure}
\end{minipage}\par\medskip
\begin{solution}
当 $F>0$ 时，令 $\ell=\sqrt{F/\kappa}$。平衡条件为 $U'=\kappa x^2-F=0$，给出图\ref{mp4:fig:l4-e7-q}中的两个驻点 $x=\pm\ell$。由于 $U''=2\kappa x$，右驻点是势阱底，左驻点是势垒顶。两处势能为
\[
U(\ell)=-\frac23\kappa\ell^3,
\qquad U(-\ell)=\frac23\kappa\ell^3.
\]
相对于右侧阱底的势垒高度为
\[
\Delta U=U(-\ell)-U(\ell)
=\frac43\kappa\ell^3=\frac{4F^{3/2}}{3\sqrt\kappa}.
\]

在 $x=\ell+\xi$ 附近，运动方程的线性项为 $m\ddot\xi=-2\kappa\ell\xi$，故该平衡态为中心，小振动角频率为
\[
\omega_c=\sqrt{\frac{2\kappa\ell}{m}}
=\sqrt{\frac2m}(\kappa F)^{1/4}.
\]
在 $x=-\ell+\xi$ 附近则有 $m\ddot\xi=+2\kappa\ell\xi$，是鞍点。

\begin{figure}[H]\centering
\begin{tikzpicture}
\draw[axis](-2.8,0)--(3,0)node[right]{$q$};
\draw[axis](0,-2.4)--(0,2.4)node[above]{$w$};
\draw[flow,line width=.9pt] plot coordinates{(2.00000,-0.00000) (1.99949,-0.05552) (1.99794,-0.11098) (1.99538,-0.16631) (1.99178,-0.22143) (1.98717,-0.27630) (1.98153,-0.33082) (1.97488,-0.38496) (1.96722,-0.43863) (1.95855,-0.49178) (1.94889,-0.54434) (1.93823,-0.59625) (1.92658,-0.64745) (1.91396,-0.69788) (1.90037,-0.74748) (1.88582,-0.79620) (1.87032,-0.84396) (1.85388,-0.89073) (1.83651,-0.93645) (1.81823,-0.98106) (1.79904,-1.02452) (1.77896,-1.06677) (1.75801,-1.10778) (1.73619,-1.14748) (1.71353,-1.18585) (1.69003,-1.22284) (1.66572,-1.25842) (1.64061,-1.29253) (1.61472,-1.32516) (1.58806,-1.35627) (1.56066,-1.38582) (1.53253,-1.41379) (1.50370,-1.44016) (1.47417,-1.46489) (1.44398,-1.48798) (1.41314,-1.50940) (1.38168,-1.52913) (1.34961,-1.54717) (1.31696,-1.56349) (1.28375,-1.57810) (1.25000,-1.59099) (1.21574,-1.60215) (1.18099,-1.61158) (1.14577,-1.61929) (1.11010,-1.62528) (1.07403,-1.62955) (1.03755,-1.63212) (1.00071,-1.63299) (0.96353,-1.63219) (0.92602,-1.62972) (0.88823,-1.62561) (0.85017,-1.61988) (0.81187,-1.61255) (0.77335,-1.60365) (0.73465,-1.59321) (0.69579,-1.58125) (0.65679,-1.56781) (0.61769,-1.55293) (0.57850,-1.53664) (0.53927,-1.51899) (0.50000,-1.50000) (0.46073,-1.47973) (0.42150,-1.45822) (0.38231,-1.43551) (0.34321,-1.41166) (0.30421,-1.38672) (0.26535,-1.36073) (0.22665,-1.33374) (0.18813,-1.30582) (0.14983,-1.27701) (0.11177,-1.24738) (0.07398,-1.21697) (0.03647,-1.18585) (-0.00071,-1.15409) (-0.03755,-1.12172) (-0.07403,-1.08883) (-0.11010,-1.05547) (-0.14577,-1.02170) (-0.18099,-0.98758) (-0.21574,-0.95318) (-0.25000,-0.91856) (-0.28375,-0.88378) (-0.31696,-0.84891) (-0.34961,-0.81400) (-0.38168,-0.77913) (-0.41314,-0.74435) (-0.44398,-0.70973) (-0.47417,-0.67533) (-0.50370,-0.64120) (-0.53253,-0.60741) (-0.56066,-0.57403) (-0.58806,-0.54110) (-0.61472,-0.50868) (-0.64061,-0.47684) (-0.66572,-0.44563) (-0.69003,-0.41510) (-0.71353,-0.38531) (-0.73619,-0.35630) (-0.75801,-0.32814) (-0.77896,-0.30086) (-0.79904,-0.27452) (-0.81823,-0.24916) (-0.83651,-0.22482) (-0.85388,-0.20155) (-0.87032,-0.17939) (-0.88582,-0.15837) (-0.90037,-0.13854) (-0.91396,-0.11992) (-0.92658,-0.10255) (-0.93823,-0.08645) (-0.94889,-0.07166) (-0.95855,-0.05821) (-0.96722,-0.04610) (-0.97488,-0.03537) (-0.98153,-0.02604) (-0.98717,-0.01811) (-0.99178,-0.01160) (-0.99538,-0.00653) (-0.99794,-0.00291) (-0.99949,-0.00073) (-1.00000,-0.00000) (-0.99949,0.00073) (-0.99794,0.00291) (-0.99538,0.00653) (-0.99178,0.01160) (-0.98717,0.01811) (-0.98153,0.02604) (-0.97488,0.03537) (-0.96722,0.04610) (-0.95855,0.05821) (-0.94889,0.07166) (-0.93823,0.08645) (-0.92658,0.10255) (-0.91396,0.11992) (-0.90037,0.13854) (-0.88582,0.15837) (-0.87032,0.17939) (-0.85388,0.20155) (-0.83651,0.22482) (-0.81823,0.24916) (-0.79904,0.27452) (-0.77896,0.30086) (-0.75801,0.32814) (-0.73619,0.35630) (-0.71353,0.38531) (-0.69003,0.41510) (-0.66572,0.44563) (-0.64061,0.47684) (-0.61472,0.50868) (-0.58806,0.54110) (-0.56066,0.57403) (-0.53253,0.60741) (-0.50370,0.64120) (-0.47417,0.67533) (-0.44398,0.70973) (-0.41314,0.74435) (-0.38168,0.77913) (-0.34961,0.81400) (-0.31696,0.84891) (-0.28375,0.88378) (-0.25000,0.91856) (-0.21574,0.95318) (-0.18099,0.98758) (-0.14577,1.02170) (-0.11010,1.05547) (-0.07403,1.08883) (-0.03755,1.12172) (-0.00071,1.15409) (0.03647,1.18585) (0.07398,1.21697) (0.11177,1.24738) (0.14983,1.27701) (0.18813,1.30582) (0.22665,1.33374) (0.26535,1.36073) (0.30421,1.38672) (0.34321,1.41166) (0.38231,1.43551) (0.42150,1.45822) (0.46073,1.47973) (0.50000,1.50000) (0.53927,1.51899) (0.57850,1.53664) (0.61769,1.55293) (0.65679,1.56781) (0.69579,1.58125) (0.73465,1.59321) (0.77335,1.60365) (0.81187,1.61255) (0.85017,1.61988) (0.88823,1.62561) (0.92602,1.62972) (0.96353,1.63219) (1.00071,1.63299) (1.03755,1.63212) (1.07403,1.62955) (1.11010,1.62528) (1.14577,1.61929) (1.18099,1.61158) (1.21574,1.60215) (1.25000,1.59099) (1.28375,1.57810) (1.31696,1.56349) (1.34961,1.54717) (1.38168,1.52913) (1.41314,1.50940) (1.44398,1.48798) (1.47417,1.46489) (1.50370,1.44016) (1.53253,1.41379) (1.56066,1.38582) (1.58806,1.35627) (1.61472,1.32516) (1.64061,1.29253) (1.66572,1.25842) (1.69003,1.22284) (1.71353,1.18585) (1.73619,1.14748) (1.75801,1.10778) (1.77896,1.06677) (1.79904,1.02452) (1.81823,0.98106) (1.83651,0.93645) (1.85388,0.89073) (1.87032,0.84396) (1.88582,0.79620) (1.90037,0.74748) (1.91396,0.69788) (1.92658,0.64745) (1.93823,0.59625) (1.94889,0.54434) (1.95855,0.49178) (1.96722,0.43863) (1.97488,0.38496) (1.98153,0.33082) (1.98717,0.27630) (1.99178,0.22143) (1.99538,0.16631) (1.99794,0.11098) (1.99949,0.05552) (2.00000,0.00000)};
\draw[flow,dashed] plot coordinates{(-2.30000,2.20106) (-2.28700,2.17575) (-2.27400,2.15051) (-2.26100,2.12532) (-2.24800,2.10020) (-2.23500,2.07514) (-2.22200,2.05015) (-2.20900,2.02521) (-2.19600,2.00034) (-2.18300,1.97553) (-2.17000,1.95078) (-2.15700,1.92610) (-2.14400,1.90147) (-2.13100,1.87691) (-2.11800,1.85242) (-2.10500,1.82799) (-2.09200,1.80362) (-2.07900,1.77931) (-2.06600,1.75507) (-2.05300,1.73090) (-2.04000,1.70678) (-2.02700,1.68273) (-2.01400,1.65875) (-2.00100,1.63483) (-1.98800,1.61098) (-1.97500,1.58719) (-1.96200,1.56346) (-1.94900,1.53980) (-1.93600,1.51620) (-1.92300,1.49267) (-1.91000,1.46921) (-1.89700,1.44581) (-1.88400,1.42248) (-1.87100,1.39921) (-1.85800,1.37601) (-1.84500,1.35288) (-1.83200,1.32981) (-1.81900,1.30681) (-1.80600,1.28388) (-1.79300,1.26101) (-1.78000,1.23821) (-1.76700,1.21548) (-1.75400,1.19281) (-1.74100,1.17022) (-1.72800,1.14769) (-1.71500,1.12523) (-1.70200,1.10283) (-1.68900,1.08051) (-1.67600,1.05825) (-1.66300,1.03606) (-1.65000,1.01394) (-1.63700,0.99189) (-1.62400,0.96991) (-1.61100,0.94800) (-1.59800,0.92616) (-1.58500,0.90439) (-1.57200,0.88269) (-1.55900,0.86105) (-1.54600,0.83949) (-1.53300,0.81800) (-1.52000,0.79658) (-1.50700,0.77523) (-1.49400,0.75395) (-1.48100,0.73274) (-1.46800,0.71161) (-1.45500,0.69054) (-1.44200,0.66955) (-1.42900,0.64863) (-1.41600,0.62778) (-1.40300,0.60700) (-1.39000,0.58630) (-1.37700,0.56567) (-1.36400,0.54511) (-1.35100,0.52462) (-1.33800,0.50421) (-1.32500,0.48388) (-1.31200,0.46361) (-1.29900,0.44342) (-1.28600,0.42331) (-1.27300,0.40326) (-1.26000,0.38330) (-1.24700,0.36341) (-1.23400,0.34359) (-1.22100,0.32385) (-1.20800,0.30418) (-1.19500,0.28459) (-1.18200,0.26508) (-1.16900,0.24564) (-1.15600,0.22628) (-1.14300,0.20700) (-1.13000,0.18779) (-1.11700,0.16866) (-1.10400,0.14961) (-1.09100,0.13063) (-1.07800,0.11173) (-1.06500,0.09291) (-1.05200,0.07417) (-1.03900,0.05551) (-1.02600,0.03693) (-1.01300,0.01842) (-1.00000,0.00000)};
\draw[flow,dashed] plot coordinates{(-1.00000,-0.00000) (-1.01300,-0.01842) (-1.02600,-0.03693) (-1.03900,-0.05551) (-1.05200,-0.07417) (-1.06500,-0.09291) (-1.07800,-0.11173) (-1.09100,-0.13063) (-1.10400,-0.14961) (-1.11700,-0.16866) (-1.13000,-0.18779) (-1.14300,-0.20700) (-1.15600,-0.22628) (-1.16900,-0.24564) (-1.18200,-0.26508) (-1.19500,-0.28459) (-1.20800,-0.30418) (-1.22100,-0.32385) (-1.23400,-0.34359) (-1.24700,-0.36341) (-1.26000,-0.38330) (-1.27300,-0.40326) (-1.28600,-0.42331) (-1.29900,-0.44342) (-1.31200,-0.46361) (-1.32500,-0.48388) (-1.33800,-0.50421) (-1.35100,-0.52462) (-1.36400,-0.54511) (-1.37700,-0.56567) (-1.39000,-0.58630) (-1.40300,-0.60700) (-1.41600,-0.62778) (-1.42900,-0.64863) (-1.44200,-0.66955) (-1.45500,-0.69054) (-1.46800,-0.71161) (-1.48100,-0.73274) (-1.49400,-0.75395) (-1.50700,-0.77523) (-1.52000,-0.79658) (-1.53300,-0.81800) (-1.54600,-0.83949) (-1.55900,-0.86105) (-1.57200,-0.88269) (-1.58500,-0.90439) (-1.59800,-0.92616) (-1.61100,-0.94800) (-1.62400,-0.96991) (-1.63700,-0.99189) (-1.65000,-1.01394) (-1.66300,-1.03606) (-1.67600,-1.05825) (-1.68900,-1.08051) (-1.70200,-1.10283) (-1.71500,-1.12523) (-1.72800,-1.14769) (-1.74100,-1.17022) (-1.75400,-1.19281) (-1.76700,-1.21548) (-1.78000,-1.23821) (-1.79300,-1.26101) (-1.80600,-1.28388) (-1.81900,-1.30681) (-1.83200,-1.32981) (-1.84500,-1.35288) (-1.85800,-1.37601) (-1.87100,-1.39921) (-1.88400,-1.42248) (-1.89700,-1.44581) (-1.91000,-1.46921) (-1.92300,-1.49267) (-1.93600,-1.51620) (-1.94900,-1.53980) (-1.96200,-1.56346) (-1.97500,-1.58719) (-1.98800,-1.61098) (-2.00100,-1.63483) (-2.01400,-1.65875) (-2.02700,-1.68273) (-2.04000,-1.70678) (-2.05300,-1.73090) (-2.06600,-1.75507) (-2.07900,-1.77931) (-2.09200,-1.80362) (-2.10500,-1.82799) (-2.11800,-1.85242) (-2.13100,-1.87691) (-2.14400,-1.90147) (-2.15700,-1.92610) (-2.17000,-1.95078) (-2.18300,-1.97553) (-2.19600,-2.00034) (-2.20900,-2.02521) (-2.22200,-2.05015) (-2.23500,-2.07514) (-2.24800,-2.10020) (-2.26100,-2.12532) (-2.27400,-2.15051) (-2.28700,-2.17575) (-2.30000,-2.20106)};
\fill(-1,0)circle(1.2pt);\node at(-1.6,-.25){$S$};
\fill(1,0)circle(1.2pt)node[below]{$C$};
\fill(2,0)circle(1.2pt)node[below right]{$B$};
\node at(.8,2.1){$h=2/3$};
\end{tikzpicture}
\caption{三次势能的临界相图。S 是鞍点，C 是中心，B 是同宿环的最右转向点；左侧开支与右环不能在有限时间内经 S 互通。}\label{mp4:fig:l4-e7-s}
\end{figure}

图\ref{mp4:fig:l4-e7-s}取 $q=x/\ell$、$w=v/(\ell\omega_s)$，其中 $\omega_s=\sqrt{\kappa\ell/m}$。无量纲能量为
\[
h=\frac12w^2+\frac13q^3-q.
\]
临界能量等于鞍点势能 $h_s=2/3$，所以
\[
w^2=\frac43-\frac23q^3+2q
=\frac23(2-q)(q+1)^2.
\]
$q=-1$ 是二重根，表示鞍点 S；$q=2$ 是单根，表示右端普通转向点 B。因此临界环的最右端为 $x_B=2\ell$。在右环内 $-1<q\le2$，上支从 S 附近向右，到 B 后转入下支，再渐近回到 S。

两类根对应不同的时间性质。在 B 附近 $|w|$ 与 $\sqrt{2-q}$ 成正比，$\int\dd q/|w|$ 有限；在 S 附近 $|w|\sim\sqrt2\,|q+1|$，时间积分对数发散。同一个代数等值线还有 $q<-1$ 的两支开轨线，它们不能在有限时间内穿过 S 与右环互通。图中 $C=(1,0)$ 周围的环内区域，对应束缚在局部势阱内的运动。

当 $F\downarrow0$ 时，两平衡点的间距 $2\sqrt{F/\kappa}$ 趋于零，势垒按 $F^{3/2}$ 消失，小振动频率按 $F^{1/4}$ 趋零。此前使用的 $\ell$ 也趋零，因此不能直接把 $F=0$ 代入归一化坐标来判断稳定性，应回到物理方程。

$F=0$ 时 $m\ddot x=-\kappa x^2$。若从任意小的负位移 $x(0)=-\epsilon$ 静止释放，初始加速度向左；此后 $x\le-\epsilon$，加速度始终不大于 $-\kappa\epsilon^2/m$，质点会离开任意固定的小邻域。因此原点不稳定。$F<0$ 时 $U'=\kappa x^2+|F|>0$，已经没有平衡位置。
\end{solution}

\section{阻尼及综合应用}
耗散使能量随时间变化，相轨线通常不再闭合。阻力和干摩擦按各自适用阶段列式，轨线的连接还须包含可能出现的静止区间。

\subsection{耗散系统的能量变化}
考虑黏性阻尼力 $-b\dot x$，其中 $b>0$，方程为
\[
m\ddot x+b\dot x+U'(x)=0.
\]
相平面上的自治方程仍然成立，只是加速度同时依赖位置和速度：
\[
\dot x=v,\qquad \dot v=-\frac bm v-\frac1mU'(x).
\]
把方程乘以 $v$，总机械能的导数为
\[
\dot E=mv\dot v+U'(x)\dot x
=v[m\dot v+U'(x)]=-bv^2.
\]
因此
\begin{equation}\label{mp4:eq:l4-dissipation}
\dot E=-bv^2\le0,\qquad E=\frac12mv^2+U(x).
\end{equation}
相点因此会穿过能量逐渐降低的等值线。转向瞬间虽有 $v=0$、$\dot E=0$，只要加速度不为零，随后仍会继续耗散。

若 $U$ 有严格极小点，阻尼常使中心变成吸引平衡点；若有极大点，阻尼并不会把它变成稳定点。此时划分两种最终去向的边界，是鞍点的稳定轨线，通常已经不再是某条机械能等值线。

对于有界势阱内的非静止周期解，若周期为 $T$，式\eqref{mp4:eq:l4-dissipation}会要求
\[
E(T)-E(0)=-b\int_0^T v^2\dd t<0,
\]
这与周期解的能量重复矛盾。因此单靠黏性阻尼和位置势能，不会产生非静止闭合相轨线。外力持续供能时则须重新分析能量收支。

\subsection{线性阻尼振子的三种形态}
写成
\begin{equation}\label{mp4:eq:l4-damped}
\ddot x+2\beta\dot x+\omega_0^2x=0,
\qquad\beta=\frac b{2m},\quad\omega_0^2=\frac km.
\end{equation}
取指数解 $x=\ee^{rt}$，代入得
\[
r^2+2\beta r+\omega_0^2=0,
\qquad (r+\beta)^2=\beta^2-\omega_0^2.
\]
两根为 $r_\pm=-\beta\pm\sqrt{\beta^2-\omega_0^2}$。欠阻尼 $\beta<\omega_0$ 时，两根为 $-\beta\pm\mathrm i\omega_d$，其中 $\omega_d=\sqrt{\omega_0^2-\beta^2}$，对应的两实解组合为
\[
x=\ee^{-\beta t}(A\cos\omega_dt+B\sin\omega_dt).
\]
相点不断绕行并趋向原点。临界阻尼 $\beta=\omega_0$ 时特征根重合，令 $x=\ee^{-\omega_0t}u(t)$。保留乘积求导的各项，
\[
\dot x=\ee^{-\omega_0t}(\dot u-\omega_0u),\qquad
\ddot x=\ee^{-\omega_0t}(\ddot u-2\omega_0\dot u+\omega_0^2u).
\]
代回方程，$u$ 与 $\dot u$ 项相消，余下 $\ddot u=0$。积分两次，$\dot u=B$、$u=A+Bt$，所以 $x=(A+Bt)\ee^{-\omega_0t}$。过阻尼 $\beta>\omega_0$ 时两根相异且都小于零，故 $x=C_+\ee^{r_+t}+C_-\ee^{r_-t}$，两个指数项都衰减。常数均由位置和速度初值联立确定，下面分别展开计算。

临界和过阻尼情形都没有正弦振荡因子。若反向初速度足够大，质点仍可能越过平衡位置一次；是否越过须由初态确定的系数判断。

\par\medskip\noindent\begin{minipage}{\textwidth}
\begin{example}[从同一初态比较三种阻尼]\label{mp4:example:l4-e8}
弹簧振子初态 $x(0)=a>0$、$\dot x(0)=0$，随后受黏性阻尼，见图\ref{mp4:fig:l4-e8-q}。分别求欠阻尼、临界阻尼、过阻尼的运动，判断能否越过平衡位置，并说明临界、过阻尼轨线接近原点时的方向。
\end{example}

\begin{figure}[H]\centering
\begin{tikzpicture}
\draw[curve](-2.4,-.35)--(-2.4,.65);\draw(-2.4,.1)--(-2.15,.1)--(-1.97,.35)--(-1.7,-.15)--(-1.43,.35)--(-1.16,-.15)--(-.89,.35)--(-.62,.1)--(-.3,.1);
\draw[fill=white](-.3,-.25)rectangle(.35,.45);\node at(.03,.1){$m$};\draw(-2.6,-.3)--(2.3,-.3);
\draw[axis](-1.1,-.8)--(2.5,-.8)node[right]{$x$};\draw[guide](-1,-.3)--(-1,-1.35);\node[below left]at(-1,-.8){$0$};\node[above]at(-1.4,.55){$k$};
\node at(1.25,.75){$F_b=-b\dot x$};
\draw[guide](.025,-.25)--(.025,-1.35);
\draw[<->,>=Stealth,line width=.35pt](-1,-1.2)--(.025,-1.2)node[midway,below]{$x(0)=a$};\node at(1.55,-1.25){$v(0)=0$};
\end{tikzpicture}
\caption{具有黏性阻尼的弹簧振子。释放时物体位于平衡位置右侧，速度为零。}\label{mp4:fig:l4-e8-q}
\end{figure}
\end{minipage}\par\medskip
\begin{solution}
设黏性阻力为$-2m\beta\dot x$，$\omega_0=\sqrt{k/m}$，则运动方程为$\ddot x+2\beta\dot x+\omega_0^2x=0$。取图\ref{mp4:fig:l4-e8-s}中的变量 $q=x/a$、$p=v/(\omega_0a)$、$\tau=\omega_0t$。三种阻尼的初态均为 $(1,0)$，而 $\ddot x(0)=-\omega_0^2a<0$，所以相点都先竖直向下离开横轴。以后是否穿过纵轴，须由满足同一初值的时间解判断。

\begin{figure}[H]\centering
\begin{tikzpicture}
\begin{scope}[shift={(0.0,0)},scale=0.9]
\draw[axis](-1.2,0)--(2.3,0)node[right]{$q$};
\draw[axis](0,-1.8)--(0,1.3)node[above]{$p$};
\draw[flow] plot coordinates{(1.80000,-0.00000) (1.79720,-0.09974) (1.78887,-0.19718) (1.77516,-0.29206) (1.75621,-0.38414) (1.73219,-0.47319) (1.70328,-0.55900) (1.66965,-0.64136) (1.63151,-0.72008) (1.58907,-0.79500) (1.54255,-0.86595) (1.49216,-0.93280) (1.43815,-0.99542) (1.38076,-1.05371) (1.32022,-1.10755) (1.25679,-1.15689) (1.19073,-1.20165) (1.12230,-1.24179) (1.05174,-1.27727) (0.97933,-1.30808) (0.90532,-1.33422) (0.82998,-1.35569) (0.75357,-1.37252) (0.67635,-1.38476) (0.59856,-1.39246) (0.52047,-1.39568) (0.44233,-1.39451) (0.36437,-1.38903) (0.28684,-1.37935) (0.20996,-1.36558) (0.13397,-1.34785) (0.05907,-1.32628) (-0.01451,-1.30104) (-0.08658,-1.27225) (-0.15694,-1.24010) (-0.22541,-1.20473) (-0.29181,-1.16634) (-0.35599,-1.12508) (-0.41777,-1.08116) (-0.47703,-1.03477) (-0.53362,-0.98608) (-0.58743,-0.93531) (-0.63834,-0.88264) (-0.68626,-0.82829) (-0.73108,-0.77245) (-0.77275,-0.71533) (-0.81118,-0.65713) (-0.84633,-0.59805) (-0.87815,-0.53830) (-0.90661,-0.47806) (-0.93169,-0.41754) (-0.95337,-0.35694) (-0.97167,-0.29644) (-0.98658,-0.23622) (-0.99813,-0.17648) (-1.00636,-0.11737) (-1.01129,-0.05909) (-1.01299,-0.00178) (-1.01151,0.05439) (-1.00693,0.10927) (-0.99930,0.16271) (-0.98873,0.21458) (-0.97530,0.26475) (-0.95911,0.31310) (-0.94027,0.35951) (-0.91888,0.40388) (-0.89507,0.44611) (-0.86896,0.48612) (-0.84067,0.52381) (-0.81034,0.55913) (-0.77810,0.59201) (-0.74408,0.62239) (-0.70844,0.65024) (-0.67130,0.67551) (-0.63283,0.69818) (-0.59316,0.71823) (-0.55243,0.73565) (-0.51081,0.75044) (-0.46843,0.76261) (-0.42545,0.77216) (-0.38200,0.77913) (-0.33823,0.78354) (-0.29429,0.78544) (-0.25031,0.78485) (-0.20643,0.78184) (-0.16279,0.77647) (-0.11951,0.76879) (-0.07673,0.75888) (-0.03456,0.74681) (0.00688,0.73267) (0.04746,0.71653) (0.08709,0.69849) (0.12566,0.67864) (0.16307,0.65709) (0.19923,0.63392) (0.23404,0.60925) (0.26744,0.58318) (0.29933,0.55582) (0.32967,0.52728) (0.35837,0.49767) (0.38539,0.46711) (0.41067,0.43571) (0.43417,0.40359) (0.45586,0.37085) (0.47570,0.33762) (0.49367,0.30400) (0.50974,0.27011) (0.52392,0.23605) (0.53618,0.20195) (0.54653,0.16789) (0.55499,0.13400) (0.56155,0.10037) (0.56623,0.06709) (0.56907,0.03427) (0.57008,0.00201) (0.56931,-0.02963) (0.56678,-0.06053) (0.56254,-0.09064) (0.55664,-0.11986) (0.54914,-0.14813) (0.54007,-0.17537) (0.52951,-0.20152) (0.51752,-0.22653) (0.50417,-0.25033) (0.48951,-0.27289) (0.47363,-0.29414) (0.45659,-0.31406) (0.43848,-0.33261) (0.41936,-0.34975) (0.39933,-0.36547) (0.37846,-0.37973) (0.35683,-0.39254) (0.33452,-0.40387) (0.31162,-0.41372) (0.28821,-0.42209) (0.26438,-0.42898) (0.24020,-0.43441) (0.21575,-0.43837) (0.19112,-0.44090) (0.16640,-0.44201) (0.14164,-0.44173) (0.11695,-0.44008) (0.09238,-0.43709) (0.06802,-0.43281) (0.04393,-0.42727) (0.02019,-0.42052) (-0.00315,-0.41259) (-0.02600,-0.40355) (-0.04832,-0.39343) (-0.07005,-0.38229) (-0.09112,-0.37019) (-0.11149,-0.35718) (-0.13111,-0.34332) (-0.14993,-0.32867) (-0.16791,-0.31329) (-0.18501,-0.29725) (-0.20119,-0.28061) (-0.21642,-0.26343) (-0.23068,-0.24577) (-0.24394,-0.22770) (-0.25618,-0.20929) (-0.26738,-0.19059) (-0.27752,-0.17168) (-0.28660,-0.15261) (-0.29461,-0.13345) (-0.30155,-0.11425) (-0.30741,-0.09509) (-0.31220,-0.07601) (-0.31592,-0.05708) (-0.31860,-0.03834) (-0.32022,-0.01986) (-0.32083,-0.00169) (-0.32042,0.01612) (-0.31903,0.03353) (-0.31667,0.05048) (-0.31338,0.06695) (-0.30919,0.08287) (-0.30411,0.09822) (-0.29820,0.11296) (-0.29147,0.12705) (-0.28398,0.14047) (-0.27575,0.15319) (-0.26684,0.16517) (-0.25727,0.17641) (-0.24709,0.18687) (-0.23635,0.19654) (-0.22510,0.20541) (-0.21336,0.21347) (-0.20120,0.22070) (-0.18866,0.22710) (-0.17578,0.23267) (-0.16262,0.23741) (-0.14921,0.24131) (-0.13561,0.24439) (-0.12185,0.24665) (-0.10800,0.24810) (-0.09408,0.24874) (-0.08015,0.24861) (-0.06625,0.24770) (-0.05242,0.24605) (-0.03871,0.24366) (-0.02515,0.24057) (-0.01178,0.23679) (0.00136,0.23235) (0.01423,0.22727) (0.02681,0.22160) (0.03904,0.21535) (0.05091,0.20855) (0.06239,0.20125) (0.07345,0.19346) (0.08405,0.18523) (0.09418,0.17659) (0.10382,0.16758) (0.11295,0.15822) (0.12154,0.14856) (0.12958,0.13863) (0.13706,0.12847) (0.14396,0.11811) (0.15028,0.10759) (0.15601,0.09695) (0.16114,0.08622) (0.16567,0.07544) (0.16959,0.06464) (0.17291,0.05385) (0.17562,0.04311) (0.17774,0.03245) (0.17926,0.02191) (0.18019,0.01150) (0.18055,0.00127) (0.18034,-0.00876) (0.17957,-0.01857) (0.17827,-0.02812) (0.17643,-0.03739) (0.17408,-0.04636) (0.17124,-0.05501) (0.16793,-0.06332) (0.16416,-0.07126) (0.15995,-0.07882) (0.15534,-0.08599) (0.15033,-0.09275) (0.14496,-0.09908) (0.13924,-0.10499) (0.13321,-0.11044) (0.12688,-0.11545) (0.12029,-0.12000) (0.11345,-0.12408) (0.10640,-0.12770) (0.09916,-0.13085) (0.09175,-0.13353) (0.08421,-0.13574) (0.07656,-0.13749) (0.06882,-0.13877) (0.06102,-0.13960) (0.05319,-0.13998) (0.04535,-0.13992) (0.03753,-0.13942) (0.02975,-0.13850) (0.02202,-0.13717) (0.01439,-0.13545) (0.00686,-0.13333) (-0.00054,-0.13084) (-0.00779,-0.12800) (-0.01487,-0.12482) (-0.02176,-0.12131) (-0.02845,-0.11749) (-0.03491,-0.11339) (-0.04114,-0.10902) (-0.04712,-0.10439) (-0.05283,-0.09954) (-0.05826,-0.09447) (-0.06341,-0.08921) (-0.06825,-0.08378) (-0.07279,-0.07819) (-0.07701,-0.07248) (-0.08090,-0.06665) (-0.08447,-0.06074) (-0.08770,-0.05475) (-0.09060,-0.04871) (-0.09316,-0.04265) (-0.09538,-0.03657) (-0.09725,-0.03050) (-0.09879,-0.02445) (-0.09999,-0.01845) (-0.10086,-0.01251) (-0.10140,-0.00666) (-0.10161,-0.00089) (-0.10150,0.00476) (-0.10108,0.01028) (-0.10035,0.01566) (-0.09933,0.02088) (-0.09802,0.02594) (-0.09643,0.03081) (-0.09457,0.03549) (-0.09245,0.03997) (-0.09010,0.04423) (-0.08751,0.04827) (-0.08469,0.05208) (-0.08168,0.05565) (-0.07847,0.05898) (-0.07508,0.06206) (-0.07152,0.06489) (-0.06781,0.06746) (-0.06397,0.06976) (-0.06000,0.07181) (-0.05593,0.07359) (-0.05177,0.07510) (-0.04753,0.07636) (-0.04322,0.07735) (-0.03887,0.07808) (-0.03448,0.07855) (-0.03007,0.07878) (-0.02566,0.07875) (-0.02126,0.07848) (-0.01688,0.07797) (-0.01253,0.07723) (-0.00823,0.07626) (-0.00399,0.07508) (0.00017,0.07368) (0.00426,0.07209) (0.00824,0.07030) (0.01213,0.06833) (0.01589,0.06619) (0.01954,0.06389) (0.02305,0.06143) (0.02641,0.05883) (0.02963,0.05610) (0.03270,0.05326) (0.03560,0.05030) (0.03833,0.04725) (0.04089,0.04411) (0.04327,0.04089) (0.04546,0.03761) (0.04748,0.03429) (0.04930,0.03092) (0.05094,0.02752) (0.05239,0.02411) (0.05364,0.02069) (0.05470,0.01727) (0.05557,0.01387) (0.05626,0.01049) (0.05675,0.00715) (0.05706,0.00385) (0.05718,0.00060) (0.05713,-0.00258) (0.05689,-0.00569) (0.05649,-0.00872) (0.05592,-0.01166) (0.05519,-0.01451) (0.05430,-0.01725) (0.05326,-0.01989) (0.05207,-0.02241) (0.05075,-0.02482) (0.04929,-0.02710) (0.04772,-0.02924) (0.04602,-0.03126) (0.04422,-0.03314) (0.04231,-0.03487) (0.04031,-0.03647) (0.03823,-0.03792) (0.03607,-0.03922) (0.03384,-0.04038) (0.03155,-0.04138) (0.02921,-0.04224) (0.02682,-0.04295) (0.02440,-0.04351) (0.02195,-0.04393) (0.01948,-0.04420) (0.01700,-0.04433) (0.01452,-0.04432) (0.01204,-0.04417) (0.00958,-0.04389) (0.00713,-0.04348) (0.00471,-0.04294) (0.00232,-0.04227) (-0.00002,-0.04149) (-0.00232,-0.04060) (-0.00457,-0.03960) (-0.00676,-0.03849) (-0.00888,-0.03729) (-0.01093,-0.03600) (-0.01291,-0.03462) (-0.01481,-0.03316) (-0.01662,-0.03162) (-0.01835,-0.03002) (-0.01998,-0.02836) (-0.02152,-0.02664) (-0.02297,-0.02488) (-0.02431,-0.02307) (-0.02555,-0.02123) (-0.02668,-0.01935) (-0.02772,-0.01746) (-0.02864,-0.01555) (-0.02946,-0.01363) (-0.03017,-0.01170) (-0.03077,-0.00978) (-0.03126,-0.00786) (-0.03165,-0.00596) (-0.03193,-0.00408) (-0.03211,-0.00222) (-0.03218,-0.00040) (-0.03215,0.00139) (-0.03202,0.00315) (-0.03180,0.00485) (-0.03148,0.00651) (-0.03107,0.00812) (-0.03057,0.00966) (-0.02999,0.01115) (-0.02933,0.01257) (-0.02858,0.01393) (-0.02777,0.01521) (-0.02688,0.01642) (-0.02593,0.01756) (-0.02492,0.01862) (-0.02385,0.01960) (-0.02272,0.02050) (-0.02155,0.02132) (-0.02034,0.02205) (-0.01908,0.02270) (-0.01780,0.02327) (-0.01648,0.02376) (-0.01514,0.02416) (-0.01377,0.02448) (-0.01240,0.02472) (-0.01101,0.02487) (-0.00961,0.02495) (-0.00822,0.02494) (-0.00682,0.02486) (-0.00543,0.02471) (-0.00405,0.02448) (-0.00269,0.02417) (-0.00135,0.02380) (-0.00003,0.02337) (0.00127,0.02286) (0.00253,0.02230) (0.00376,0.02168) (0.00496,0.02101) (0.00612,0.02028) (0.00723,0.01951) (0.00830,0.01869) (0.00932,0.01782) (0.01030,0.01692) (0.01122,0.01599) (0.01209,0.01502) (0.01290,0.01403) (0.01366,0.01302) (0.01436,0.01198) (0.01500,0.01093) (0.01558,0.00986) (0.01610,0.00878) (0.01656,0.00770) (0.01697,0.00662) (0.01731,0.00554) (0.01759,0.00446) (0.01781,0.00339) (0.01797,0.00233) (0.01807,0.00128) (0.01811,0.00026) (0.01810,-0.00075) (0.01803,-0.00174) (0.01790,-0.00270) (0.01772,-0.00364) (0.01749,-0.00454) (0.01722,-0.00541) (0.01689,-0.00625) (0.01652,-0.00705) (0.01610,-0.00781) (0.01564,-0.00854) (0.01514,-0.00922) (0.01461,-0.00986) (0.01404,-0.01046) (0.01344,-0.01101) (0.01281,-0.01152) (0.01215,-0.01198) (0.01147,-0.01240) (0.01076,-0.01277) (0.01004,-0.01309) (0.00930,-0.01336) (0.00854,-0.01359) (0.00778,-0.01377) (0.00700,-0.01391) (0.00622,-0.01400) (0.00543,-0.01404) (0.00465,-0.01404) (0.00386,-0.01399) (0.00308,-0.01391) (0.00231,-0.01378) (0.00154,-0.01361) (0.00078,-0.01340) (0.00004,-0.01316) (-0.00069,-0.01288) (-0.00140,-0.01256) (-0.00210,-0.01221) (-0.00277,-0.01184) (-0.00342,-0.01143) (-0.00405,-0.01099) (-0.00465,-0.01053) (-0.00523,-0.01005) (-0.00578,-0.00954) (-0.00630,-0.00902) (-0.00679,-0.00847) (-0.00725,-0.00791)};
\fill(1.8,0)circle(1.2pt)node[above]{$P$};
\node at(.5,-2.1){欠阻尼};\end{scope}
\begin{scope}[shift={(3.65,0)},scale=0.9]
\draw[axis](-1.2,0)--(2.3,0)node[right]{$q$};
\draw[axis](0,-1.8)--(0,1.3)node[above]{$p$};
\draw[flow] plot coordinates{(1.80000,-0.00000) (1.79971,-0.03182) (1.79886,-0.06251) (1.79747,-0.09209) (1.79555,-0.12060) (1.79313,-0.14806) (1.79023,-0.17450) (1.78686,-0.19995) (1.78304,-0.22444) (1.77878,-0.24799) (1.77411,-0.27063) (1.76905,-0.29238) (1.76359,-0.31327) (1.75777,-0.33332) (1.75160,-0.35256) (1.74509,-0.37100) (1.73825,-0.38868) (1.73110,-0.40560) (1.72365,-0.42180) (1.71592,-0.43729) (1.70791,-0.45209) (1.69965,-0.46623) (1.69113,-0.47972) (1.68238,-0.49258) (1.67340,-0.50483) (1.66421,-0.51648) (1.65481,-0.52756) (1.64522,-0.53807) (1.63544,-0.54805) (1.62549,-0.55750) (1.61538,-0.56643) (1.60511,-0.57487) (1.59469,-0.58283) (1.58413,-0.59032) (1.57344,-0.59736) (1.56262,-0.60396) (1.55170,-0.61013) (1.54066,-0.61589) (1.52953,-0.62126) (1.51830,-0.62623) (1.50698,-0.63083) (1.49559,-0.63507) (1.48412,-0.63895) (1.47259,-0.64249) (1.46100,-0.64571) (1.44935,-0.64860) (1.43765,-0.65119) (1.42591,-0.65348) (1.41413,-0.65547) (1.40231,-0.65719) (1.39047,-0.65864) (1.37860,-0.65983) (1.36672,-0.66077) (1.35482,-0.66146) (1.34290,-0.66192) (1.33099,-0.66215) (1.31907,-0.66216) (1.30715,-0.66196) (1.29524,-0.66156) (1.28334,-0.66096) (1.27145,-0.66017) (1.25957,-0.65920) (1.24772,-0.65806) (1.23588,-0.65674) (1.22407,-0.65527) (1.21229,-0.65363) (1.20054,-0.65185) (1.18883,-0.64992) (1.17715,-0.64785) (1.16551,-0.64565) (1.15390,-0.64333) (1.14235,-0.64088) (1.13083,-0.63831) (1.11937,-0.63563) (1.10795,-0.63284) (1.09659,-0.62995) (1.08527,-0.62697) (1.07402,-0.62388) (1.06281,-0.62071) (1.05167,-0.61746) (1.04059,-0.61412) (1.02956,-0.61070) (1.01860,-0.60721) (1.00770,-0.60365) (0.99687,-0.60003) (0.98610,-0.59634) (0.97540,-0.59259) (0.96477,-0.58879) (0.95421,-0.58493) (0.94371,-0.58103) (0.93329,-0.57707) (0.92294,-0.57308) (0.91266,-0.56904) (0.90245,-0.56496) (0.89232,-0.56085) (0.88226,-0.55670) (0.87228,-0.55253) (0.86237,-0.54833) (0.85254,-0.54410) (0.84279,-0.53984) (0.83311,-0.53557) (0.82350,-0.53127) (0.81398,-0.52696) (0.80453,-0.52264) (0.79517,-0.51830) (0.78588,-0.51395) (0.77666,-0.50959) (0.76753,-0.50522) (0.75848,-0.50084) (0.74950,-0.49646) (0.74060,-0.49208) (0.73179,-0.48769) (0.72305,-0.48331) (0.71439,-0.47893) (0.70581,-0.47455) (0.69730,-0.47017) (0.68888,-0.46580) (0.68053,-0.46143) (0.67227,-0.45707) (0.66408,-0.45272) (0.65597,-0.44838) (0.64794,-0.44406) (0.63998,-0.43974) (0.63211,-0.43543) (0.62431,-0.43114) (0.61659,-0.42687) (0.60894,-0.42261) (0.60137,-0.41836) (0.59388,-0.41413) (0.58646,-0.40992) (0.57912,-0.40573) (0.57186,-0.40156) (0.56467,-0.39741) (0.55755,-0.39327) (0.55051,-0.38916) (0.54354,-0.38507) (0.53664,-0.38101) (0.52982,-0.37696) (0.52307,-0.37294) (0.51640,-0.36894) (0.50979,-0.36496) (0.50326,-0.36101) (0.49680,-0.35709) (0.49040,-0.35319) (0.48408,-0.34931) (0.47783,-0.34547) (0.47164,-0.34164) (0.46553,-0.33785) (0.45948,-0.33408) (0.45350,-0.33033) (0.44759,-0.32662) (0.44174,-0.32293) (0.43596,-0.31927) (0.43025,-0.31564) (0.42460,-0.31203) (0.41902,-0.30846) (0.41350,-0.30491) (0.40804,-0.30139) (0.40265,-0.29790) (0.39731,-0.29444) (0.39205,-0.29100) (0.38684,-0.28760) (0.38169,-0.28422) (0.37661,-0.28087) (0.37158,-0.27756) (0.36661,-0.27427) (0.36171,-0.27101) (0.35686,-0.26778) (0.35207,-0.26457) (0.34733,-0.26140) (0.34266,-0.25826) (0.33803,-0.25514) (0.33347,-0.25206) (0.32896,-0.24900) (0.32451,-0.24597) (0.32011,-0.24297) (0.31576,-0.24000) (0.31147,-0.23706) (0.30722,-0.23415) (0.30304,-0.23126) (0.29890,-0.22840) (0.29481,-0.22558) (0.29078,-0.22278) (0.28679,-0.22000) (0.28286,-0.21726) (0.27897,-0.21454) (0.27513,-0.21186) (0.27134,-0.20920) (0.26760,-0.20656) (0.26391,-0.20396) (0.26026,-0.20138) (0.25666,-0.19883) (0.25310,-0.19630) (0.24959,-0.19380) (0.24613,-0.19133) (0.24270,-0.18889) (0.23932,-0.18647) (0.23599,-0.18408) (0.23270,-0.18171) (0.22945,-0.17937) (0.22624,-0.17706) (0.22307,-0.17477) (0.21995,-0.17250) (0.21686,-0.17027) (0.21382,-0.16805) (0.21081,-0.16586) (0.20785,-0.16370) (0.20492,-0.16156) (0.20203,-0.15944) (0.19918,-0.15735) (0.19637,-0.15529) (0.19359,-0.15324) (0.19085,-0.15122) (0.18815,-0.14922) (0.18548,-0.14725) (0.18284,-0.14530) (0.18025,-0.14337) (0.17768,-0.14147) (0.17515,-0.13958) (0.17266,-0.13772) (0.17020,-0.13588) (0.16777,-0.13406) (0.16537,-0.13227) (0.16300,-0.13049) (0.16067,-0.12874) (0.15837,-0.12701) (0.15610,-0.12530) (0.15386,-0.12361) (0.15165,-0.12194) (0.14947,-0.12029) (0.14732,-0.11866) (0.14520,-0.11705) (0.14310,-0.11546) (0.14104,-0.11389) (0.13900,-0.11233) (0.13700,-0.11080) (0.13502,-0.10929) (0.13306,-0.10779) (0.13113,-0.10632) (0.12923,-0.10486) (0.12736,-0.10342) (0.12551,-0.10200) (0.12369,-0.10059) (0.12189,-0.09921) (0.12012,-0.09784) (0.11837,-0.09649) (0.11664,-0.09515) (0.11494,-0.09384) (0.11326,-0.09254) (0.11161,-0.09125) (0.10998,-0.08998) (0.10837,-0.08873) (0.10678,-0.08750) (0.10522,-0.08628) (0.10368,-0.08507) (0.10216,-0.08388) (0.10066,-0.08271) (0.09918,-0.08155) (0.09772,-0.08041) (0.09629,-0.07928) (0.09487,-0.07817) (0.09347,-0.07707) (0.09209,-0.07598) (0.09074,-0.07491) (0.08940,-0.07386) (0.08808,-0.07281) (0.08678,-0.07178) (0.08549,-0.07077) (0.08423,-0.06977) (0.08298,-0.06878) (0.08175,-0.06780) (0.08054,-0.06684) (0.07935,-0.06589) (0.07817,-0.06495) (0.07701,-0.06403) (0.07586,-0.06311) (0.07474,-0.06221) (0.07362,-0.06132) (0.07253,-0.06045) (0.07145,-0.05958) (0.07038,-0.05873) (0.06933,-0.05789) (0.06830,-0.05706) (0.06728,-0.05624) (0.06627,-0.05543) (0.06528,-0.05463) (0.06431,-0.05385) (0.06334,-0.05307) (0.06240,-0.05231) (0.06146,-0.05155) (0.06054,-0.05081) (0.05963,-0.05007) (0.05874,-0.04935) (0.05786,-0.04863) (0.05699,-0.04793) (0.05613,-0.04723) (0.05529,-0.04655) (0.05445,-0.04587) (0.05363,-0.04521) (0.05283,-0.04455) (0.05203,-0.04390) (0.05125,-0.04326) (0.05047,-0.04263) (0.04971,-0.04201) (0.04896,-0.04140) (0.04822,-0.04079) (0.04749,-0.04020) (0.04677,-0.03961) (0.04607,-0.03903) (0.04537,-0.03846) (0.04468,-0.03789) (0.04401,-0.03734) (0.04334,-0.03679) (0.04268,-0.03625) (0.04203,-0.03571) (0.04140,-0.03519) (0.04077,-0.03467) (0.04015,-0.03416) (0.03954,-0.03366) (0.03894,-0.03316) (0.03834,-0.03267) (0.03776,-0.03219) (0.03718,-0.03171) (0.03662,-0.03124) (0.03606,-0.03078) (0.03551,-0.03033) (0.03497,-0.02988) (0.03443,-0.02943) (0.03391,-0.02900) (0.03339,-0.02857) (0.03288,-0.02814) (0.03238,-0.02772) (0.03188,-0.02731) (0.03139,-0.02690) (0.03091,-0.02650) (0.03044,-0.02611) (0.02997,-0.02572) (0.02951,-0.02534) (0.02906,-0.02496) (0.02861,-0.02459) (0.02818,-0.02422) (0.02774,-0.02386) (0.02732,-0.02350) (0.02690,-0.02315) (0.02648,-0.02280) (0.02608,-0.02246) (0.02567,-0.02212) (0.02528,-0.02179) (0.02489,-0.02146) (0.02451,-0.02114) (0.02413,-0.02082) (0.02376,-0.02051) (0.02339,-0.02020) (0.02303,-0.01990) (0.02267,-0.01960) (0.02232,-0.01930) (0.02198,-0.01901) (0.02164,-0.01873) (0.02131,-0.01844) (0.02098,-0.01816) (0.02065,-0.01789) (0.02033,-0.01762) (0.02002,-0.01735) (0.01971,-0.01709) (0.01940,-0.01683) (0.01910,-0.01658) (0.01880,-0.01633) (0.01851,-0.01608) (0.01823,-0.01584) (0.01794,-0.01560) (0.01766,-0.01536) (0.01739,-0.01513) (0.01712,-0.01490) (0.01685,-0.01467) (0.01659,-0.01445) (0.01633,-0.01423) (0.01608,-0.01401) (0.01583,-0.01380) (0.01558,-0.01359) (0.01534,-0.01338) (0.01510,-0.01318) (0.01487,-0.01297) (0.01463,-0.01278) (0.01441,-0.01258) (0.01418,-0.01239) (0.01396,-0.01220) (0.01374,-0.01201) (0.01353,-0.01183) (0.01332,-0.01165) (0.01311,-0.01147) (0.01290,-0.01129) (0.01270,-0.01112) (0.01250,-0.01095) (0.01231,-0.01078) (0.01211,-0.01062) (0.01193,-0.01045) (0.01174,-0.01029) (0.01155,-0.01014) (0.01137,-0.00998) (0.01120,-0.00983) (0.01102,-0.00968) (0.01085,-0.00953) (0.01068,-0.00938) (0.01051,-0.00924) (0.01034,-0.00909) (0.01018,-0.00895) (0.01002,-0.00882) (0.00986,-0.00868) (0.00971,-0.00855) (0.00956,-0.00841) (0.00941,-0.00828) (0.00926,-0.00816) (0.00911,-0.00803) (0.00897,-0.00791) (0.00883,-0.00778) (0.00869,-0.00766) (0.00855,-0.00754) (0.00842,-0.00743) (0.00828,-0.00731) (0.00815,-0.00720) (0.00803,-0.00709) (0.00790,-0.00698) (0.00777,-0.00687) (0.00765,-0.00676) (0.00753,-0.00666) (0.00741,-0.00656) (0.00729,-0.00645) (0.00718,-0.00635) (0.00707,-0.00625) (0.00695,-0.00616) (0.00684,-0.00606) (0.00674,-0.00597) (0.00663,-0.00587) (0.00652,-0.00578) (0.00642,-0.00569) (0.00632,-0.00560) (0.00622,-0.00552) (0.00612,-0.00543) (0.00602,-0.00535) (0.00593,-0.00526) (0.00583,-0.00518) (0.00574,-0.00510) (0.00565,-0.00502) (0.00556,-0.00494) (0.00547,-0.00486) (0.00539,-0.00479) (0.00530,-0.00471) (0.00522,-0.00464) (0.00513,-0.00457) (0.00505,-0.00450) (0.00497,-0.00443) (0.00489,-0.00436) (0.00482,-0.00429) (0.00474,-0.00422) (0.00466,-0.00415) (0.00459,-0.00409) (0.00452,-0.00403) (0.00444,-0.00396) (0.00437,-0.00390) (0.00430,-0.00384) (0.00424,-0.00378) (0.00417,-0.00372) (0.00410,-0.00366) (0.00404,-0.00360) (0.00397,-0.00355) (0.00391,-0.00349) (0.00385,-0.00344) (0.00378,-0.00338) (0.00372,-0.00333) (0.00366,-0.00328) (0.00361,-0.00322) (0.00355,-0.00317) (0.00349,-0.00312) (0.00344,-0.00307) (0.00338,-0.00303) (0.00333,-0.00298) (0.00327,-0.00293) (0.00322,-0.00289) (0.00317,-0.00284) (0.00312,-0.00280) (0.00307,-0.00275) (0.00302,-0.00271) (0.00297,-0.00266) (0.00292,-0.00262) (0.00288,-0.00258) (0.00283,-0.00254) (0.00279,-0.00250) (0.00274,-0.00246) (0.00270,-0.00242) (0.00265,-0.00238) (0.00261,-0.00235) (0.00257,-0.00231) (0.00253,-0.00227) (0.00249,-0.00224) (0.00245,-0.00220) (0.00241,-0.00217) (0.00237,-0.00213) (0.00233,-0.00210) (0.00229,-0.00206) (0.00226,-0.00203) (0.00222,-0.00200)};
\fill(1.8,0)circle(1.2pt)node[above]{$P$};
\draw[guide](-.5,0.5)--(1.5,-1.5);\node at(.5,-2.1){临界阻尼};\end{scope}
\begin{scope}[shift={(7.3,0)},scale=0.9]
\draw[axis](-1.2,0)--(2.3,0)node[right]{$q$};
\draw[axis](0,-1.8)--(0,1.3)node[above]{$p$};
\draw[flow] plot coordinates{(1.80000,0.00000) (1.79862,-0.06756) (1.79470,-0.12691) (1.78856,-0.17893) (1.78048,-0.22444) (1.77069,-0.26414) (1.75941,-0.29868) (1.74685,-0.32861) (1.73318,-0.35446) (1.71855,-0.37667) (1.70309,-0.39564) (1.68693,-0.41174) (1.67018,-0.42529) (1.65294,-0.43657) (1.63528,-0.44584) (1.61730,-0.45333) (1.59904,-0.45924) (1.58058,-0.46375) (1.56196,-0.46702) (1.54323,-0.46919) (1.52443,-0.47040) (1.50561,-0.47076) (1.48678,-0.47037) (1.46799,-0.46932) (1.44925,-0.46769) (1.43058,-0.46555) (1.41201,-0.46297) (1.39355,-0.46001) (1.37521,-0.45671) (1.35701,-0.45312) (1.33896,-0.44928) (1.32107,-0.44523) (1.30335,-0.44100) (1.28580,-0.43661) (1.26842,-0.43210) (1.25123,-0.42748) (1.23422,-0.42277) (1.21741,-0.41799) (1.20079,-0.41316) (1.18436,-0.40829) (1.16812,-0.40339) (1.15208,-0.39847) (1.13624,-0.39355) (1.12060,-0.38863) (1.10515,-0.38371) (1.08990,-0.37881) (1.07485,-0.37393) (1.05999,-0.36907) (1.04532,-0.36425) (1.03085,-0.35945) (1.01657,-0.35470) (1.00247,-0.34998) (0.98857,-0.34530) (0.97485,-0.34067) (0.96131,-0.33608) (0.94796,-0.33154) (0.93479,-0.32704) (0.92180,-0.32260) (0.90898,-0.31820) (0.89634,-0.31386) (0.88387,-0.30956) (0.87157,-0.30532) (0.85945,-0.30113) (0.84748,-0.29699) (0.83569,-0.29290) (0.82405,-0.28886) (0.81258,-0.28487) (0.80126,-0.28094) (0.79010,-0.27705) (0.77910,-0.27322) (0.76824,-0.26944) (0.75754,-0.26570) (0.74699,-0.26202) (0.73658,-0.25839) (0.72631,-0.25480) (0.71619,-0.25126) (0.70621,-0.24777) (0.69637,-0.24433) (0.68667,-0.24093) (0.67710,-0.23758) (0.66766,-0.23428) (0.65835,-0.23102) (0.64918,-0.22781) (0.64013,-0.22464) (0.63120,-0.22151) (0.62241,-0.21843) (0.61373,-0.21538) (0.60518,-0.21239) (0.59674,-0.20943) (0.58842,-0.20651) (0.58022,-0.20363) (0.57213,-0.20080) (0.56415,-0.19800) (0.55629,-0.19524) (0.54853,-0.19252) (0.54089,-0.18984) (0.53335,-0.18719) (0.52591,-0.18458) (0.51858,-0.18201) (0.51135,-0.17948) (0.50422,-0.17697) (0.49719,-0.17451) (0.49026,-0.17208) (0.48342,-0.16968) (0.47668,-0.16731) (0.47004,-0.16498) (0.46349,-0.16268) (0.45702,-0.16041) (0.45065,-0.15818) (0.44437,-0.15597) (0.43817,-0.15380) (0.43207,-0.15165) (0.42604,-0.14954) (0.42010,-0.14745) (0.41424,-0.14540) (0.40847,-0.14337) (0.40277,-0.14137) (0.39716,-0.13940) (0.39162,-0.13746) (0.38616,-0.13554) (0.38078,-0.13365) (0.37547,-0.13179) (0.37023,-0.12995) (0.36507,-0.12814) (0.35998,-0.12635) (0.35496,-0.12459) (0.35002,-0.12286) (0.34514,-0.12114) (0.34032,-0.11945) (0.33558,-0.11779) (0.33090,-0.11615) (0.32629,-0.11453) (0.32174,-0.11293) (0.31725,-0.11136) (0.31283,-0.10980) (0.30847,-0.10827) (0.30417,-0.10676) (0.29993,-0.10527) (0.29574,-0.10381) (0.29162,-0.10236) (0.28756,-0.10093) (0.28355,-0.09952) (0.27959,-0.09814) (0.27570,-0.09677) (0.27185,-0.09542) (0.26806,-0.09409) (0.26432,-0.09278) (0.26064,-0.09148) (0.25701,-0.09021) (0.25342,-0.08895) (0.24989,-0.08771) (0.24641,-0.08649) (0.24297,-0.08528) (0.23958,-0.08409) (0.23624,-0.08292) (0.23295,-0.08176) (0.22970,-0.08063) (0.22650,-0.07950) (0.22334,-0.07839) (0.22023,-0.07730) (0.21716,-0.07622) (0.21413,-0.07516) (0.21114,-0.07411) (0.20820,-0.07308) (0.20530,-0.07206) (0.20243,-0.07105) (0.19961,-0.07006) (0.19683,-0.06909) (0.19409,-0.06812) (0.19138,-0.06717) (0.18871,-0.06624) (0.18608,-0.06531) (0.18349,-0.06440) (0.18093,-0.06351) (0.17841,-0.06262) (0.17592,-0.06175) (0.17347,-0.06089) (0.17105,-0.06004) (0.16866,-0.05920) (0.16631,-0.05838) (0.16399,-0.05756) (0.16171,-0.05676) (0.15945,-0.05597) (0.15723,-0.05519) (0.15504,-0.05442) (0.15287,-0.05366) (0.15074,-0.05291) (0.14864,-0.05217) (0.14657,-0.05145) (0.14453,-0.05073) (0.14251,-0.05002) (0.14052,-0.04932) (0.13856,-0.04864) (0.13663,-0.04796) (0.13473,-0.04729) (0.13285,-0.04663) (0.13100,-0.04598) (0.12917,-0.04534) (0.12737,-0.04471) (0.12559,-0.04408) (0.12384,-0.04347) (0.12212,-0.04286) (0.12041,-0.04227) (0.11874,-0.04168) (0.11708,-0.04110) (0.11545,-0.04052) (0.11384,-0.03996) (0.11225,-0.03940) (0.11069,-0.03885) (0.10914,-0.03831) (0.10762,-0.03778) (0.10612,-0.03725) (0.10464,-0.03673) (0.10318,-0.03622) (0.10174,-0.03571) (0.10033,-0.03521) (0.09893,-0.03472) (0.09755,-0.03424) (0.09619,-0.03376) (0.09485,-0.03329) (0.09352,-0.03283) (0.09222,-0.03237) (0.09093,-0.03192) (0.08967,-0.03147) (0.08842,-0.03103) (0.08718,-0.03060) (0.08597,-0.03017) (0.08477,-0.02975) (0.08359,-0.02934) (0.08242,-0.02893) (0.08127,-0.02853) (0.08014,-0.02813) (0.07902,-0.02774) (0.07792,-0.02735) (0.07683,-0.02697) (0.07576,-0.02659) (0.07471,-0.02622) (0.07367,-0.02586) (0.07264,-0.02550) (0.07163,-0.02514) (0.07063,-0.02479) (0.06964,-0.02444) (0.06867,-0.02410) (0.06771,-0.02377) (0.06677,-0.02344) (0.06584,-0.02311) (0.06492,-0.02279) (0.06402,-0.02247) (0.06312,-0.02216) (0.06224,-0.02185) (0.06138,-0.02154) (0.06052,-0.02124) (0.05968,-0.02095) (0.05884,-0.02065) (0.05802,-0.02037) (0.05721,-0.02008) (0.05642,-0.01980) (0.05563,-0.01953) (0.05486,-0.01925) (0.05409,-0.01899) (0.05334,-0.01872) (0.05259,-0.01846) (0.05186,-0.01820) (0.05114,-0.01795) (0.05042,-0.01770) (0.04972,-0.01745) (0.04903,-0.01721) (0.04834,-0.01697) (0.04767,-0.01673) (0.04701,-0.01650) (0.04635,-0.01627) (0.04570,-0.01604) (0.04507,-0.01582) (0.04444,-0.01560) (0.04382,-0.01538) (0.04321,-0.01517) (0.04261,-0.01495) (0.04201,-0.01475) (0.04143,-0.01454) (0.04085,-0.01434) (0.04028,-0.01414) (0.03972,-0.01394) (0.03916,-0.01375) (0.03862,-0.01355) (0.03808,-0.01337) (0.03755,-0.01318) (0.03702,-0.01300) (0.03651,-0.01281) (0.03600,-0.01264) (0.03550,-0.01246) (0.03500,-0.01229) (0.03451,-0.01211) (0.03403,-0.01195) (0.03356,-0.01178) (0.03309,-0.01161) (0.03263,-0.01145) (0.03217,-0.01129) (0.03173,-0.01114) (0.03128,-0.01098) (0.03085,-0.01083) (0.03042,-0.01068) (0.02999,-0.01053) (0.02958,-0.01038) (0.02916,-0.01024) (0.02876,-0.01009) (0.02836,-0.00995) (0.02796,-0.00981) (0.02757,-0.00968) (0.02719,-0.00954) (0.02681,-0.00941) (0.02643,-0.00928) (0.02606,-0.00915) (0.02570,-0.00902) (0.02534,-0.00890) (0.02499,-0.00877) (0.02464,-0.00865) (0.02430,-0.00853) (0.02396,-0.00841) (0.02362,-0.00829) (0.02330,-0.00818) (0.02297,-0.00806) (0.02265,-0.00795) (0.02233,-0.00784) (0.02202,-0.00773) (0.02172,-0.00762) (0.02141,-0.00752) (0.02111,-0.00741) (0.02082,-0.00731) (0.02053,-0.00721) (0.02024,-0.00711) (0.01996,-0.00701) (0.01968,-0.00691) (0.01941,-0.00681) (0.01914,-0.00672) (0.01887,-0.00662) (0.01861,-0.00653) (0.01835,-0.00644) (0.01809,-0.00635) (0.01784,-0.00626) (0.01759,-0.00617) (0.01735,-0.00609) (0.01711,-0.00600) (0.01687,-0.00592) (0.01663,-0.00584) (0.01640,-0.00576) (0.01617,-0.00568) (0.01595,-0.00560) (0.01572,-0.00552) (0.01550,-0.00544) (0.01529,-0.00537) (0.01507,-0.00529) (0.01486,-0.00522) (0.01466,-0.00514) (0.01445,-0.00507) (0.01425,-0.00500) (0.01405,-0.00493) (0.01386,-0.00486) (0.01366,-0.00480) (0.01347,-0.00473) (0.01329,-0.00466) (0.01310,-0.00460) (0.01292,-0.00453) (0.01274,-0.00447) (0.01256,-0.00441) (0.01238,-0.00435) (0.01221,-0.00429) (0.01204,-0.00423) (0.01187,-0.00417) (0.01171,-0.00411) (0.01155,-0.00405) (0.01138,-0.00400) (0.01123,-0.00394) (0.01107,-0.00389) (0.01091,-0.00383) (0.01076,-0.00378) (0.01061,-0.00372) (0.01046,-0.00367) (0.01032,-0.00362) (0.01017,-0.00357) (0.01003,-0.00352) (0.00989,-0.00347) (0.00975,-0.00342) (0.00962,-0.00338) (0.00948,-0.00333) (0.00935,-0.00328) (0.00922,-0.00324) (0.00909,-0.00319) (0.00897,-0.00315) (0.00884,-0.00310) (0.00872,-0.00306) (0.00860,-0.00302) (0.00848,-0.00298) (0.00836,-0.00293) (0.00824,-0.00289) (0.00813,-0.00285) (0.00801,-0.00281) (0.00790,-0.00277) (0.00779,-0.00274) (0.00768,-0.00270) (0.00758,-0.00266) (0.00747,-0.00262) (0.00737,-0.00259) (0.00726,-0.00255) (0.00716,-0.00251) (0.00706,-0.00248) (0.00696,-0.00244) (0.00687,-0.00241) (0.00677,-0.00238) (0.00668,-0.00234) (0.00658,-0.00231) (0.00649,-0.00228) (0.00640,-0.00225) (0.00631,-0.00222) (0.00622,-0.00218) (0.00614,-0.00215) (0.00605,-0.00212) (0.00597,-0.00209) (0.00588,-0.00207) (0.00580,-0.00204) (0.00572,-0.00201) (0.00564,-0.00198) (0.00556,-0.00195) (0.00549,-0.00193) (0.00541,-0.00190) (0.00533,-0.00187) (0.00526,-0.00185) (0.00519,-0.00182) (0.00511,-0.00179) (0.00504,-0.00177) (0.00497,-0.00175) (0.00490,-0.00172) (0.00483,-0.00170) (0.00477,-0.00167) (0.00470,-0.00165) (0.00464,-0.00163) (0.00457,-0.00160) (0.00451,-0.00158) (0.00444,-0.00156) (0.00438,-0.00154) (0.00432,-0.00152) (0.00426,-0.00150) (0.00420,-0.00147) (0.00414,-0.00145) (0.00408,-0.00143) (0.00403,-0.00141) (0.00397,-0.00139) (0.00392,-0.00137) (0.00386,-0.00136) (0.00381,-0.00134) (0.00375,-0.00132) (0.00370,-0.00130) (0.00365,-0.00128) (0.00360,-0.00126) (0.00355,-0.00125) (0.00350,-0.00123) (0.00345,-0.00121) (0.00340,-0.00119) (0.00336,-0.00118) (0.00331,-0.00116) (0.00326,-0.00115) (0.00322,-0.00113) (0.00317,-0.00111) (0.00313,-0.00110) (0.00308,-0.00108) (0.00304,-0.00107) (0.00300,-0.00105) (0.00296,-0.00104) (0.00292,-0.00102) (0.00288,-0.00101) (0.00284,-0.00100) (0.00280,-0.00098) (0.00276,-0.00097) (0.00272,-0.00095) (0.00268,-0.00094) (0.00264,-0.00093) (0.00261,-0.00091) (0.00257,-0.00090) (0.00253,-0.00089) (0.00250,-0.00088) (0.00246,-0.00086) (0.00243,-0.00085) (0.00240,-0.00084) (0.00236,-0.00083) (0.00233,-0.00082) (0.00230,-0.00081) (0.00227,-0.00080) (0.00223,-0.00078) (0.00220,-0.00077) (0.00217,-0.00076) (0.00214,-0.00075) (0.00211,-0.00074) (0.00208,-0.00073) (0.00205,-0.00072) (0.00202,-0.00071) (0.00200,-0.00070) (0.00197,-0.00069) (0.00194,-0.00068) (0.00191,-0.00067) (0.00189,-0.00066) (0.00186,-0.00065) (0.00183,-0.00064)};
\fill(1.8,0)circle(1.2pt)node[above]{$P$};
\draw[guide](-.5,0.1755002001601601)--(1.5,-0.5265006004804803);\node at(.5,-2.1){过阻尼};\end{scope}

\end{tikzpicture}
\caption{同一释放初态的阻尼相图，从左到右取 $\beta/\omega_0=0.18,\,1,\,1.6$。虚线为后两种情形在原点处的极限切线。}\label{mp4:fig:l4-e8-s}
\end{figure}

欠阻尼时，令 $\omega_d=\sqrt{\omega_0^2-\beta^2}$。对通解使用乘积法则，
\[
v=\ee^{-\beta t}\bigl[(-\beta A+\omega_dB)\cos\omega_dt
+(-\beta B-\omega_dA)\sin\omega_dt\bigr].
\]
初态给出 $A=a$、$-\beta A+\omega_dB=0$，即 $B=\beta a/\omega_d$。代回后，正弦项系数为 $-a(\beta^2+\omega_d^2)/\omega_d=-a\omega_0^2/\omega_d$，所以
\[
x=a\ee^{-\beta t}\left(\cos\omega_dt+\frac\beta{\omega_d}\sin\omega_dt\right),
\]
\[
v=-\frac{a\omega_0^2}{\omega_d}\ee^{-\beta t}\sin\omega_dt.
\]
首次经过平衡位置要求位置的括号为零。由于此前 $v<0$，相位应取第二象限，
\[
t_0=\frac{\pi-\arctan(\omega_d/\beta)}{\omega_d}.
\]
此时速度非零，图中的轨线穿过纵轴。以后正弦因子不断变号，运动持续越过平衡位置，振幅逐渐衰减，相轨线向原点螺旋收缩。

临界阻尼时 $x=(A+Bt)\ee^{-\omega_0t}$，其导数为 $v=[B-\omega_0(A+Bt)]\ee^{-\omega_0t}$。由 $x(0)=A=a$、$v(0)=B-\omega_0A=0$，得 $B=\omega_0a$，故
\[
x=a(1+\omega_0t)\ee^{-\omega_0t}>0,\qquad
v=-a\omega_0^2t\ee^{-\omega_0t}<0\quad(t>0).
\]
质点从右侧单调趋近原点，始终不越过。其相轨线斜向趋近原点，且
\[
\frac vx=-\frac{\omega_0^2t}{1+\omega_0t}\longrightarrow-\omega_0.
\]
因此归一化图中的极限切线为 $p=-q$。

过阻尼时 $r_\pm=-\beta\pm\sqrt{\beta^2-\omega_0^2}$，$r_-<r_+<0$。先解初值方程
\[
C_++C_-=a,\qquad r_+C_++r_-C_-=0,
\]
第二式减去第一式的 $r_-$ 倍，得到 $(r_+-r_-)C_+=-r_-a$。再由第一式回代，
\[
C_+=\frac{-r_-a}{r_+-r_-},\qquad
C_-=\frac{r_+a}{r_+-r_-}.
\]
因此
\[
x=a\frac{-r_-\ee^{r_+t}+r_+\ee^{r_-t}}{r_+-r_-},
\quad
v=a\frac{r_+r_-}{r_+-r_-}
(\ee^{r_-t}-\ee^{r_+t}).
\]
对 $t>0$，指数差为负，故 $v<0$。将位置提取正因子后，剩下
\[
-r_-+r_+\ee^{(r_--r_+)t}>-r_-+r_+>0,
\]
所以 $x>0$，也不越过原点。长时间由较慢的 $r_+$ 模态主导，$v/x\to r_+$，图中的极限切线为 $p=(r_+/\omega_0)q$，其斜率绝对值小于 1。

三种运动都满足 $\dot E=-2m\beta v^2\le0$。是否过零还取决于初态：本题临界、过阻尼解的 $x(t)>0$ 已证明它们不越过；若改为足够大的反向初速度，则可能越过一次。
\end{solution}

\subsection{滑动摩擦的分段相图}
物体反向滑动时，摩擦力也随之反向。设摩擦力大小为 $f$，弹簧劲度系数为 $k$，令 $d=f/k$、$\omega_0=\sqrt{k/m}$。在滑动阶段，
\begin{equation}\label{mp4:eq:l4-coulomb}
\begin{cases}
\ddot x+\omega_0^2(x+d)=0,&v>0,\\
\ddot x+\omega_0^2(x-d)=0,&v<0.
\end{cases}
\end{equation}
因此向右滑动的相弧以 $(-d,0)$ 为中心，向左滑动的相弧以 $(d,0)$ 为中心。完整过程由两种不同中心的相弧逐段连接而成。

在 $v=0$ 的位置还须检查静摩擦。下文取最大静摩擦力也等于 $f$，则 $|x|\le d$ 时可以停住；$|x|>d$ 时弹力超过静摩擦上限，物体会继续启动。若最大静摩擦力与滑动摩擦力不同，停止区间也须相应改变。

\par\medskip\noindent\begin{minipage}{\textwidth}
\begin{example}[摩擦振子的停止位置]\label{mp4:example:l4-e9}
水平粗糙面上的弹簧振子，滑动摩擦力大小与最大静摩擦力均为 $f$。令 $d=f/k$、$\omega_0=\sqrt{k/m}$，从 $x_0=5d$ 静止释放，见图\ref{mp4:fig:l4-e9-q}。求逐次转向位置、停止位置、总滑动时间和摩擦生热。
\end{example}

\begin{figure}[H]\centering
\begin{tikzpicture}
\draw[curve](-2.4,-.35)--(-2.4,.65);\draw(-2.4,.1)--(-2.15,.1)--(-1.97,.35)--(-1.7,-.15)--(-1.43,.35)--(-1.16,-.15)--(-.89,.35)--(-.62,.1)--(-.3,.1);
\draw[fill=white](-.3,-.25)rectangle(.35,.45);\node at(.03,.1){$m$};\draw(-2.6,-.3)--(2.3,-.3);
\draw[axis](-1.1,-.8)--(2.5,-.8)node[right]{$x$};\draw[guide](-1,-.3)--(-1,-1.35);\node[below left]at(-1,-.8){$0$};\node[above]at(-1.4,.55){$k$};
\node at(1.4,.8){$f_{\rm s,max}=f_{\rm k}=f$};
\draw[guide](.025,-.25)--(.025,-1.35);
\draw[<->,>=Stealth,line width=.35pt](-1,-1.2)--(.025,-1.2)node[midway,below]{$x_0=5d$};
\end{tikzpicture}
\caption{摩擦振子的释放状态。弹簧原长位置为 $x=0$。}\label{mp4:fig:l4-e9-q}
\end{figure}
\end{minipage}\par\medskip
\begin{solution}
先确定是否启动。初态 $x_0=5d$ 时弹力向左，大小为 $5kd=5f$，超过最大静摩擦力，故物体立即向左滑动。该段 $v<0$，动摩擦向右，方程为
\[
m\ddot x=-kx+f=-k(x-d).
\]
这类固定滑动方向的方程可统一写成 $\ddot x+\omega_0^2(x-c)=0$。取 $q=x-c$，特征方程为 $\lambda^2+\omega_0^2=0$，故 $q=C\cos\omega_0s+D\sin\omega_0s$。若该段从 $x=x_s$、$v=0$ 开始，则 $C=x_s-c$、$\omega_0D=0$，即
\[
x=c+(x_s-c)\cos\omega_0s,\qquad
v=-\omega_0(x_s-c)\sin\omega_0s.
\]
第一段的等效中心在 $c=d$，以释放时刻计时，代入 $x_s=5d$ 得
\[
x=d+4d\cos\omega_0t,\qquad
v=-4\omega_0d\sin\omega_0t,
\quad 0\le t\le\frac\pi{\omega_0}.
\]
图\ref{mp4:fig:l4-e9-s}中的第一段相弧位于下半平面，绕中心 $(d,0)$ 从 A$_0$ 走到 A$_1$。第一次速度再为零时，
\[
t_1=\frac\pi{\omega_0},\qquad x_1=d-4d=-3d.
\]
此处弹力向右且为 $3f$，仍大于静摩擦上限，不能停住。

\begin{figure}[H]\centering
\begin{tikzpicture}
\begin{scope}[shift={(0,0)},scale=0.67]
\draw[axis](-3.8,0)--(5.8,0)node[right]{$x/d$};
\draw[axis](0,-4.8)--(0,2.8)node[above]{$v/(\omega_0d)$};
\fill[gray!25](-1,-.12)rectangle(1,.12);\draw[flow] plot coordinates{(5.00000,-0.00000) (4.99803,-0.12564) (4.99211,-0.25116) (4.98225,-0.37643) (4.96846,-0.50133) (4.95075,-0.62574) (4.92915,-0.74953) (4.90367,-0.87257) (4.87433,-0.99476) (4.84117,-1.11596) (4.80423,-1.23607) (4.76352,-1.35495) (4.71911,-1.47250) (4.67102,-1.58859) (4.61931,-1.70312) (4.56403,-1.81596) (4.50523,-1.92701) (4.44297,-2.03617) (4.37731,-2.14331) (4.30832,-2.24833) (4.23607,-2.35114) (4.16062,-2.45163) (4.08205,-2.54970) (4.00044,-2.64525) (3.91587,-2.73819) (3.82843,-2.82843) (3.73819,-2.91587) (3.64525,-3.00044) (3.54970,-3.08205) (3.45163,-3.16062) (3.35114,-3.23607) (3.24833,-3.30832) (3.14331,-3.37731) (3.03617,-3.44297) (2.92701,-3.50523) (2.81596,-3.56403) (2.70312,-3.61931) (2.58859,-3.67102) (2.47250,-3.71911) (2.35495,-3.76352) (2.23607,-3.80423) (2.11596,-3.84117) (1.99476,-3.87433) (1.87257,-3.90367) (1.74953,-3.92915) (1.62574,-3.95075) (1.50133,-3.96846) (1.37643,-3.98225) (1.25116,-3.99211) (1.12564,-3.99803) (1.00000,-4.00000) (0.87436,-3.99803) (0.74884,-3.99211) (0.62357,-3.98225) (0.49867,-3.96846) (0.37426,-3.95075) (0.25047,-3.92915) (0.12743,-3.90367) (0.00524,-3.87433) (-0.11596,-3.84117) (-0.23607,-3.80423) (-0.35495,-3.76352) (-0.47250,-3.71911) (-0.58859,-3.67102) (-0.70312,-3.61931) (-0.81596,-3.56403) (-0.92701,-3.50523) (-1.03617,-3.44297) (-1.14331,-3.37731) (-1.24833,-3.30832) (-1.35114,-3.23607) (-1.45163,-3.16062) (-1.54970,-3.08205) (-1.64525,-3.00044) (-1.73819,-2.91587) (-1.82843,-2.82843) (-1.91587,-2.73819) (-2.00044,-2.64525) (-2.08205,-2.54970) (-2.16062,-2.45163) (-2.23607,-2.35114) (-2.30832,-2.24833) (-2.37731,-2.14331) (-2.44297,-2.03617) (-2.50523,-1.92701) (-2.56403,-1.81596) (-2.61931,-1.70312) (-2.67102,-1.58859) (-2.71911,-1.47250) (-2.76352,-1.35495) (-2.80423,-1.23607) (-2.84117,-1.11596) (-2.87433,-0.99476) (-2.90367,-0.87257) (-2.92915,-0.74953) (-2.95075,-0.62574) (-2.96846,-0.50133) (-2.98225,-0.37643) (-2.99211,-0.25116) (-2.99803,-0.12564) (-3.00000,-0.00000)};
\draw[flow] plot coordinates{(-3.00000,0.00000) (-2.99901,0.06282) (-2.99605,0.12558) (-2.99112,0.18822) (-2.98423,0.25067) (-2.97538,0.31287) (-2.96457,0.37476) (-2.95183,0.43629) (-2.93717,0.49738) (-2.92059,0.55798) (-2.90211,0.61803) (-2.88176,0.67748) (-2.85955,0.73625) (-2.83551,0.79430) (-2.80965,0.85156) (-2.78201,0.90798) (-2.75261,0.96351) (-2.72148,1.01808) (-2.68866,1.07165) (-2.65416,1.12417) (-2.61803,1.17557) (-2.58031,1.22581) (-2.54103,1.27485) (-2.50022,1.32262) (-2.45794,1.36909) (-2.41421,1.41421) (-2.36909,1.45794) (-2.32262,1.50022) (-2.27485,1.54103) (-2.22581,1.58031) (-2.17557,1.61803) (-2.12417,1.65416) (-2.07165,1.68866) (-2.01808,1.72148) (-1.96351,1.75261) (-1.90798,1.78201) (-1.85156,1.80965) (-1.79430,1.83551) (-1.73625,1.85955) (-1.67748,1.88176) (-1.61803,1.90211) (-1.55798,1.92059) (-1.49738,1.93717) (-1.43629,1.95183) (-1.37476,1.96457) (-1.31287,1.97538) (-1.25067,1.98423) (-1.18822,1.99112) (-1.12558,1.99605) (-1.06282,1.99901) (-1.00000,2.00000) (-0.93718,1.99901) (-0.87442,1.99605) (-0.81178,1.99112) (-0.74933,1.98423) (-0.68713,1.97538) (-0.62524,1.96457) (-0.56371,1.95183) (-0.50262,1.93717) (-0.44202,1.92059) (-0.38197,1.90211) (-0.32252,1.88176) (-0.26375,1.85955) (-0.20570,1.83551) (-0.14844,1.80965) (-0.09202,1.78201) (-0.03649,1.75261) (0.01808,1.72148) (0.07165,1.68866) (0.12417,1.65416) (0.17557,1.61803) (0.22581,1.58031) (0.27485,1.54103) (0.32262,1.50022) (0.36909,1.45794) (0.41421,1.41421) (0.45794,1.36909) (0.50022,1.32262) (0.54103,1.27485) (0.58031,1.22581) (0.61803,1.17557) (0.65416,1.12417) (0.68866,1.07165) (0.72148,1.01808) (0.75261,0.96351) (0.78201,0.90798) (0.80965,0.85156) (0.83551,0.79430) (0.85955,0.73625) (0.88176,0.67748) (0.90211,0.61803) (0.92059,0.55798) (0.93717,0.49738) (0.95183,0.43629) (0.96457,0.37476) (0.97538,0.31287) (0.98423,0.25067) (0.99112,0.18822) (0.99605,0.12558) (0.99901,0.06282) (1.00000,0.00000)};
\fill(5,0)circle(1.2pt)node[above]{$A_0$};
\fill(-3,0)circle(1.2pt)node[above left]{$A_1$};
\fill(1,0)circle(1.2pt)node[above right]{$A_2$};
\node[below]at(-1,0){$-1$};\node at(1.8,-4.65){向左滑动：中心 $+1$};\node at(.3,2.6){向右滑动：中心 $-1$};\end{scope}

\end{tikzpicture}
\caption{摩擦振子的分段相弧。灰色短带为静止允许区间；$A_0$、$A_1$、$A_2$ 依次表示释放、转向与停止状态。}\label{mp4:fig:l4-e9-s}
\end{figure}

第二段 $v>0$，摩擦反向，方程改为 $m\ddot x=-k(x+d)$。令 $s=t-t_1$，以 $x=-3d,v=0$ 为新初值，
\[
x=-d-2d\cos\omega_0s,\qquad
v=2\omega_0d\sin\omega_0s,
\quad 0\le s\le\frac\pi{\omega_0}.
\]
相点绕中心 $(-d,0)$ 沿上半平面的相弧到达 A$_2$，其位置为 $x_2=d$。在此静止所需摩擦力为 $kd=f$，恰好处于允许的静摩擦范围内，因此以后保持静止。总滑动时间为
\[
t_{\rm stop}=\frac{2\pi}{\omega_0},\qquad x_{\rm stop}=d.
\]
图中的灰色短带 $|x|\le d,v=0$ 是静止允许集合。相点一旦在速度为零时进入此带，便结束滑动，不能继续沿下一条假想的动摩擦相弧运动。

生热须按总路程计算。第一段从 $5d$ 到 $-3d$，路程为 $8d$；第二段从 $-3d$ 到 $d$，路程为 $4d$。因此
\[
Q=f(8d+4d)=12fd=12kd^2.
\]
初、末动能均为零，也有
\[
Q=\frac12k(5d)^2-\frac12kd^2=12kd^2.
\]
最终弹簧中仍有 $kd^2/2$ 的势能，由静摩擦平衡剩余弹力。每段移心圆弧都将该段方向固定的摩擦力并入了等效回复力。转向后摩擦改向，圆弧中心随之改变；真实机械能在两段滑动中都持续减少。
\end{solution}

\par\medskip\noindent\begin{minipage}{\textwidth}
\begin{example}[阻尼势垒与稳定分界线]\label{mp4:example:l4-e10}
某不稳定平衡位置附近的运动满足 $\ddot x+2\beta\dot x-\omega_b^2x=0$，$\beta,\omega_b>0$。质点从 $x=-a$ 出发，受到向右冲量后初速为 $V>0$，图\ref{mp4:fig:l4-e10-q}表示这一局部势垒模型。求越过 $x=0$ 所需初速的阈值，并求临界运动及其耗散能量。
\end{example}

\begin{figure}[H]\centering
\begin{tikzpicture}
\draw[axis](-2.5,0)--(2,0)node[right]{$x$};
\draw[axis](0,-1.6)--(0,0.8)node[above]{$U$};
\draw[curve,domain=-1.8:1.8,samples=80]plot(\x,{-.4*(\x)^2});\fill(-1.5,-.9)circle(1.4pt)node[left]{$P$};\draw[axis](-1.5,-.9)--(-.4,-.9)node[midway,above]{$V$};\node[above]at(-1.5,0){$-a$};
\end{tikzpicture}
\caption{局部线性势垒模型。位置力为 $m\omega_b^2x$，阻力与速度反向。}\label{mp4:fig:l4-e10-q}
\end{figure}
\end{minipage}\par\medskip
\begin{solution}
图\ref{mp4:fig:l4-e10-q}中的位置力为 $+m\omega_b^2x$，在势垒顶两侧都指向远离原点的方向，阻尼力则为 $-2m\beta v$。令
\[
s=\sqrt{\beta^2+\omega_b^2},\qquad
r_+=-\beta+s>0,\quad r_-=-\beta-s<0.
\]
这是由特征方程 $\lambda^2+2\beta\lambda-\omega_b^2=0$ 配方为 $(\lambda+\beta)^2=\beta^2+\omega_b^2$ 求出的两根。通解为 $x=C_+\ee^{r_+t}+C_-\ee^{r_-t}$。初值方程
\[
C_++C_-=-a,\qquad r_+C_++r_-C_-=V
\]
给出
\[
C_+=\frac{V+r_-a}{2s},\qquad C_-=-a-C_+.
\]
增长模态系数 $C_+$ 的符号决定最终向哪一侧离开。恰好趋近鞍点要求消去该模态，故临界初速为
\begin{equation}\label{mp4:eq:l4-damped-threshold}
V_c=-r_-a=(\beta+s)a.
\end{equation}
图\ref{mp4:fig:l4-e10-s}中的粗线 $v=r_-x$ 正是衰减模态的相轨线，在 $x=-a$ 处的纵坐标就是 $V_c$。

\begin{figure}[H]\centering
\begin{tikzpicture}
\draw[axis](-2.5,0)--(2,0)node[right]{$x/a$};
\draw[axis](0,-0.7)--(0,2.7)node[above]{$v/(\omega_ba)$};
\draw[flow,line width=1pt] plot coordinates{(-1.70000,2.28485) (-0.02500,0.03360)};
\draw[flow,dashed] plot coordinates{(-1.00000,1.05000) (-0.97729,1.01463) (-0.95535,0.98020) (-0.93416,0.94670) (-0.91369,0.91410) (-0.89393,0.88236) (-0.87486,0.85146) (-0.85646,0.82138) (-0.83871,0.79209) (-0.82160,0.76356) (-0.80511,0.73577) (-0.78922,0.70870) (-0.77392,0.68233) (-0.75919,0.65663) (-0.74503,0.63158) (-0.73140,0.60717) (-0.71831,0.58336) (-0.70573,0.56015) (-0.69365,0.53752) (-0.68207,0.51544) (-0.67097,0.49389) (-0.66034,0.47287) (-0.65016,0.45235) (-0.64043,0.43231) (-0.63114,0.41275) (-0.62227,0.39364) (-0.61381,0.37497) (-0.60576,0.35673) (-0.59811,0.33890) (-0.59085,0.32146) (-0.58397,0.30441) (-0.57745,0.28773) (-0.57130,0.27140) (-0.56551,0.25542) (-0.56006,0.23977) (-0.55496,0.22444) (-0.55019,0.20942) (-0.54574,0.19470) (-0.54162,0.18027) (-0.53781,0.16611) (-0.53431,0.15222) (-0.53111,0.13858) (-0.52821,0.12519) (-0.52560,0.11203) (-0.52328,0.09910) (-0.52124,0.08639) (-0.51947,0.07388) (-0.51798,0.06157) (-0.51676,0.04946) (-0.51581,0.03753) (-0.51511,0.02577) (-0.51467,0.01418) (-0.51448,0.00275) (-0.51455,-0.00853) (-0.51486,-0.01967) (-0.51541,-0.03067) (-0.51621,-0.04154) (-0.51724,-0.05229) (-0.51851,-0.06292) (-0.52001,-0.07344) (-0.52174,-0.08386) (-0.52370,-0.09419) (-0.52588,-0.10442) (-0.52829,-0.11457) (-0.53092,-0.12464) (-0.53377,-0.13464) (-0.53685,-0.14458) (-0.54014,-0.15445) (-0.54364,-0.16427) (-0.54736,-0.17403) (-0.55130,-0.18376) (-0.55545,-0.19344) (-0.55981,-0.20309) (-0.56438,-0.21271) (-0.56917,-0.22231) (-0.57417,-0.23189) (-0.57937,-0.24145) (-0.58479,-0.25101) (-0.59042,-0.26056) (-0.59625,-0.27011) (-0.60230,-0.27967) (-0.60856,-0.28923) (-0.61503,-0.29881) (-0.62171,-0.30840) (-0.62860,-0.31802) (-0.63570,-0.32767) (-0.64302,-0.33734) (-0.65054,-0.34706) (-0.65829,-0.35681) (-0.66624,-0.36660) (-0.67442,-0.37645) (-0.68281,-0.38634) (-0.69142,-0.39629) (-0.70024,-0.40630) (-0.70929,-0.41638) (-0.71857,-0.42652) (-0.72806,-0.43673) (-0.73778,-0.44703) (-0.74773,-0.45740) (-0.75791,-0.46785) (-0.76832,-0.47840)};
\draw[flow,dashed] plot coordinates{(-1.00000,1.70000) (-0.96308,1.65626) (-0.92711,1.61389) (-0.89206,1.57285) (-0.85790,1.53310) (-0.82460,1.49461) (-0.79213,1.45735) (-0.76047,1.42127) (-0.72958,1.38636) (-0.69946,1.35256) (-0.67006,1.31986) (-0.64138,1.28822) (-0.61337,1.25762) (-0.58603,1.22803) (-0.55933,1.19941) (-0.53325,1.17174) (-0.50777,1.14500) (-0.48287,1.11916) (-0.45852,1.09420) (-0.43471,1.07009) (-0.41143,1.04682) (-0.38865,1.02435) (-0.36635,1.00267) (-0.34453,0.98175) (-0.32315,0.96158) (-0.30221,0.94214) (-0.28169,0.92341) (-0.26158,0.90536) (-0.24185,0.88799) (-0.22250,0.87127) (-0.20351,0.85519) (-0.18487,0.83974) (-0.16656,0.82488) (-0.14857,0.81063) (-0.13088,0.79694) (-0.11350,0.78382) (-0.09639,0.77125) (-0.07956,0.75922) (-0.06298,0.74770) (-0.04666,0.73670) (-0.03056,0.72620) (-0.01470,0.71618) (0.00095,0.70664) (0.01640,0.69756) (0.03165,0.68894) (0.04671,0.68077) (0.06160,0.67302) (0.07633,0.66570) (0.09090,0.65880) (0.10532,0.65231) (0.11960,0.64621) (0.13376,0.64051) (0.14779,0.63519) (0.16171,0.63024) (0.17552,0.62566) (0.18924,0.62144) (0.20287,0.61758) (0.21641,0.61406) (0.22989,0.61089) (0.24330,0.60805) (0.25665,0.60554) (0.26994,0.60335) (0.28320,0.60149) (0.29641,0.59993) (0.30959,0.59869) (0.32275,0.59775) (0.33590,0.59711) (0.34903,0.59676) (0.36216,0.59671) (0.37529,0.59694) (0.38842,0.59746) (0.40158,0.59826) (0.41475,0.59934) (0.42795,0.60069) (0.44118,0.60231) (0.45445,0.60420) (0.46777,0.60635) (0.48113,0.60877) (0.49456,0.61145) (0.50804,0.61439) (0.52159,0.61758) (0.53522,0.62103) (0.54892,0.62474) (0.56271,0.62869) (0.57658,0.63290) (0.59056,0.63735) (0.60463,0.64205) (0.61881,0.64700) (0.63310,0.65220) (0.64751,0.65764) (0.66204,0.66333) (0.67669,0.66926) (0.69149,0.67544) (0.70642,0.68185) (0.72149,0.68852) (0.73671,0.69542) (0.75209,0.70257) (0.76763,0.70997) (0.78333,0.71761) (0.79920,0.72549) (0.81525,0.73362)};
\fill(0,0)circle(1.2pt)node[below right]{$S$};
\draw[guide](-1,-.3)--(-1,2.2);\node[below]at(-1,0){$-1$};\node at(1.1,1.35){跨越};\node at(-2,-.55){转回};
\end{tikzpicture}
\caption{阻尼势垒的相图，取 $\beta/\omega_b=0.3$。粗线为稳定分界支；上方轨线跨越，下方轨线在左侧转回。}\label{mp4:fig:l4-e10-s}
\end{figure}

若 $V>V_c$，有 $C_+>0$、$C_-<0$。位置从负值出发，增长项最终占优而变正，必能跨越原点。由 $x(t_*)=0$ 得
\[
t_* =\frac1{2s}\ln\frac{-C_-}{C_+}>0.
\]
此时 $-C_-=a+C_+>C_+$，故所得时间为正。越过时
\[
v(t_*)=(r_+-r_-)C_+\ee^{r_+t_*}>0,
\]
是真正的穿越。

若 $V=V_c$，则
\[
x=-a\ee^{r_-t},\qquad v=-ar_-\ee^{r_-t}.
\]
位置始终为负，只有 $t\to\infty$ 时趋近 S。若 $0<V<V_c$，$C_+<0$；将位置写为
\[
x=\ee^{r_-t}\left[-a+C_+(\ee^{2st}-1)\right]<0,
\]
可见它从不越过原点。初速度为正而长时间速度为负，说明它会在左侧转回。于是实际跨越条件为严格的 $V>V_c$。

临界运动的末态位于势垒顶，动能、势能都趋于零。取 $U=-m\omega_b^2x^2/2$，初态机械能全部耗散，
\begin{align*}
Q&=\frac m2(V_c^2-\omega_b^2a^2)\\
&=\frac{ma^2}{2}\left[(\beta+s)^2-\omega_b^2\right]
=m\beta(\beta+s)a^2.
\end{align*}
直接积阻尼功也得到
\[
Q=2m\beta\int_0^\infty a^2r_-^2\ee^{2r_-t}\dd t
=2m\beta a^2r_-^2\left[\frac{\ee^{2r_-t}}{2r_-}\right]_0^\infty
=-m\beta a^2r_-=m\beta(\beta+s)a^2.
\]
临界运动虽历时无限，速度指数衰减，耗散仍有限。$\beta\to0$ 时 $V_c\to\omega_ba$，恢复无阻尼阈值；有阻尼时，仅有正的初态机械能还不够，须同时补足沿途耗散。若冲量前静止，则阈值冲量为 $mV_c$，实际越过需要更大的冲量。
\end{solution}

\subsection{外力、冲量与相图的使用范围}
短时冲量作用时，若位移可忽略，则相点沿竖直方向跳变：$x^+=x^-$，$v^+=v^-+J/m$。这一跳变不是原自治方程的连续相轨线。冲量结束后，再从新相点沿原方向场运动。改变弹簧系数、突然撤去约束等操作也要先写出新的运动方程，再判断哪些状态量连续。

若方程含显式时间，例如外力 $F_0\cos\Omega t$，仅给 $(x,v)$ 不能确定后续运动，还须给出驱动相位。不同驱动相位下的曲线投影到 $(x,v)$ 平面后可以相交，不违背自治系统的解唯一性。可把驱动相位 $\phi$ 加入状态，补上 $\dot\phi=\Omega$，或在每隔一个驱动周期的相同时刻取点，构造频闪图。本讲前面的“轨线不能相交”限于状态变量已经足以确定运动的自治系统。


\subsection{干摩擦振子的任意释放点}
例\ref{mp4:example:l4-e9}与习题\ref{mp4:exercise:l4-h11}适用于一端固定的水平弹簧振子。若初始静止位置为$A_0$，记$d=\mu mg/k$、$\omega_0=\sqrt{k/m}$，则向左滑动阶段的相圆为
\[
 (x-d)^2+(v/\omega_0)^2=(A_0-d)^2,
\]
向右滑动阶段则把圆心换成$-d$。若一段由静止位置$A$开始，且$|A|>d$，该段中心为$c=d\operatorname{sgn}A$。半周后的端点与$A$关于$c$对称，故
\[
A_{\rm next}=2d\operatorname{sgn}A-A.
\]
当$|A|>2d$，下一端点在原点另一侧，位移绝对值减少$2d$；其中只有$|A_{\rm next}|>d$时才会再启动。若$d<|A|<2d$，最后一段在原点同侧结束，端点已经落入静止区间，不能继续按“振幅减去$2d$”解释。$|A|=2d$时恰停在原点。每段结束都须检查$|A_{\rm next}|\le d$，满足时保持静止；初态本就在$|A_0|\le d$内也不启动。一般释放点的次数与停点公式见习题\ref{mp4:exercise:l4-h11}。






\setcounter{chapter}{4}
\par\medskip\noindent\begin{minipage}{\textwidth}
\begin{example}[两滚筒与光滑平台之间的送杆]\label{mp4:hw:l4-30}
两半径$r$的滚筒中心相距$l$，其左、右各有与筒顶等高的光滑平台，平台边缘到邻近筒心的水平距离也为$l$。长$2l$的均匀杆从左平台以速度$v>0$水平滑入。两筒顶面分别以$+V,-V$运动，$V=\omega_0r$且$V^2=2\mu gl$；杆与筒之间的滑动、极限静摩擦系数均取$\mu$。求杆最终到达右平台的速度$u$，并给出不同初速度对应的过程，见图\ref{mp4:fig:hw:l4-30-q}。
\end{example}

\begin{figure}[H]\centering
\begin{tikzpicture}[>=Stealth,line cap=round,line join=round]
\fill[pattern=north east lines,pattern color=gray](-4,.75)rectangle(-2.25,1);\fill[pattern=north east lines,pattern color=gray](2.25,.75)rectangle(4,1);
\draw[thick](-4,.75)--(-2.25,.75)--(-2.25,1)--(-4,1);\draw[thick](2.25,1)--(4,1);\draw[thick](2.25,1)--(2.25,.75)--(4,.75);
\draw[curve](-.75,.65)circle(.35);\draw[curve](.75,.65)circle(.35);\draw[fill=gray!15](-4,1)rectangle(-1,1.13);
\node at(-3.1,.43){光滑平台};\node at(3.1,.43){光滑平台};
\draw[axis](-3.4,1.45)--(-2.4,1.45)node[midway,above]{$v$};\draw[<->](-4,2.2)--(-1,2.2)node[midway,above]{$2l$};\draw[axis](-1.05,1.35)--(-.45,1.35)node[above]{$V$};\draw[axis](1.05,1.35)--(.45,1.35)node[above]{$-V$};
\foreach \u in {-2.25,-.75,.75,2.25}{\draw[guide](\u,.7)--(\u,-.1);}
\node[below]at(-2.25,-.1){$-3l/2$};\node[below]at(-.75,-.1){$-l/2$};\node[below]at(.75,-.1){$l/2$};\node[below]at(2.25,-.1){$3l/2$};\draw[axis](-3,-.6)--(3,-.6)node[right]{$x$};\node[below]at(0,-.6){$O$};
\end{tikzpicture}
\caption{两侧光滑平台的边缘位于$x=\pm3l/2$，滚筒的支承位置为$x=\pm l/2$。}\label{mp4:fig:hw:l4-30-q}
\end{figure}

\end{minipage}\par\medskip
\begin{solution}
以两筒中点为原点，$x$表示杆质心。令$a=\mu g$、$q=v^2/(al)$、$X=x/l$。杆进入摩擦区时$X=-3/2$，离开时$X=3/2$，此前和此后都在光滑平台上匀速。

先求杆速$0<\dot x<V$时的加速度。左过渡段由左平台边缘和左筒支承，左筒法向力为$mg(X+3/2)$；中段两筒的法向力为$mg(1/2-X)$、$mg(1/2+X)$；右过渡段右筒法向力为$mg(3/2-X)$。因此
\[
 \frac{\ddot x}{a}=\begin{cases}
 X+3/2,&-3/2<X<-1/2,\\
 -2X,&-1/2<X<1/2,\\
 X-3/2,&1/2<X<3/2.
 \end{cases}
\]
为了逐段积分，记$Y=\dot x^2/(al)$，于是$\dd Y/\dd X=2\ddot x/a$。左过渡段以$Y(-3/2)=q$为初值，得到
\[
Y=q+\int_{-3/2}^{X}2(\xi+3/2)\dd\xi
=q+[(\xi+3/2)^2]_{-3/2}^{X}
=q+(X+3/2)^2.
\]
因此进入中段时$Y(-1/2)=q+1$。若仍有$\dot x<V$，中段积分为
\[
Y=q+1+\int_{-1/2}^{X}(-4\xi)\dd\xi
=q+1+[-2\xi^2]_{-1/2}^{X}
=q+\frac32-2X^2.
\]
它在$X=0$取得最大值$q+3/2$。到$X=1/2$时$Y=q+1$；右过渡段再积分得
\[
Y=q+1+\int_{1/2}^{X}2(\xi-3/2)\dd\xi
=q+1+[(\xi-3/2)^2]_{1/2}^{X}
=q+(X-3/2)^2.
\]
若$q\le1/2$，全程$Y\le2=V^2/(al)$，上述分支始终适用。离开摩擦区时$Y(3/2)=q$，所以$u=v$，且过程中杆速保持为正。

若杆速达到$V$，在左过渡段取静摩擦为零即可保持等速。中段的右筒仍对杆施加向左摩擦$\mu N_R$；要维持等速，左筒须提供向右静摩擦$\mu N_R$，其可行条件为
\[
\mu N_R\le\mu N_L
\quad\Longleftrightarrow\quad
\frac12+X\le\frac12-X
\quad\Longleftrightarrow\quad X\le0.
\]
因此图\ref{mp4:fig:hw:l4-30-s}中的等速支最多保持到 $X=0$。过此处以后，杆速降至 $V$ 以下。用中段的同一原函数并取 $Y(0)=2$，得 $Y=2+[-2\xi^2]_0^X=2-2X^2$，到 $X=1/2$ 时为 $3/2$；右过渡段再损失 $1$，故所有到达 $(0,V)$ 的初态最终均有 $u^2=al/2$。

若初速度高于$V$，左过渡段摩擦反向，积分为$Y=q-(X+3/2)^2$。若到中段仍有$Y>2$，两处摩擦都向左，$\ddot x=-a$，所以
\[
Y=q-1+\int_{-1/2}^{X}(-2)\dd\xi
=q-1+[-2\xi]_{-1/2}^{X}=q-2-2X.
\]
令上述各段的$Y=2$，可明确写出首次达到筒速的位置：
\[
X_*=\begin{cases}
-\sqrt{(q-1/2)/2},&1/2<q<1,\\
-3/2+\sqrt{|q-2|},&1\le q\le3,\\
(q-4)/2,&3<q\le5.
\end{cases}
\]
其中$q<2$时从低速侧加速到$V$，$q>2$时从高速侧减速到$V$；$q=2$从进入摩擦区起便可保持等速。当$1/2<q\le4$，均有$X_*\le0$，因此进入上述共同的右半程。

当$4<q\le5$，交会点位于中段右半部，不能维持等速。由$(X_*,Y)=(X_*,2)$起沿低速分支积分，
\[
Y(1/2)=2-2\left(\frac14-X_*^2\right)
=\frac32+2X_*^2.
\]
再扣除右过渡段的$1$，得$u^2/(al)=1/2+2X_*^2$。若$q>5$，杆以大于$V$的速度通过整个中段：左过渡段、中段、右过渡段使$Y$分别减少$1,2,1$，最终为$q-4$。右过渡段只与反向转动的右筒接触，即使杆速降到$V$以下，摩擦方向也不改变。各段合并，
\[
 \boxed{u^2=\begin{cases}
 v^2,&0<v^2\le al/2,\\
 al/2,&al/2<v^2\le4al,\\
 al/2+(v^2-4al)^2/(2al),&4al<v^2\le5al,\\
 v^2-4al,&v^2>5al.
 \end{cases}}
\]
四段在三个阈值处连续。两侧平台维持过渡阶段的支承；缺少平台时，杆越过支承边界后会倾转，水平送杆模型随即失效。


\begin{figure}[H]\centering
\begin{tikzpicture}[>=Stealth,line cap=round,line join=round]

\draw[axis](-3,0)--(3,0)node[right]{$X$};\draw[axis](0,0)--(0,2.6)node[above]{$\dot x/\sqrt{al}$};\draw[guide](-2.8,1.414)--(2.8,1.414)node[right]{$\sqrt2$};\draw[curve](-2.5,.5)..controls(-1.9,.5)and(-1.5,1.414)..(-.8,1.414)--(0,1.414);\draw[curve,domain=0:.5,samples=35]plot(\x,{sqrt(2-2*(\x)^2)});\draw[curve,domain=.5:1.5,samples=50]plot(\x,{sqrt(.5+(1.5-\x)^2)});\draw[curve](1.5,.707)--(2.5,.707);\node at(-1.5,2){等速支};
\begin{scope}[shift={(6.8,0)},xscale=.42,yscale=.8]\draw[axis](0,0)--(7,0)node[right]{$q$};\draw[axis](0,0)--(0,3.5)node[above]{$u^2/(al)$};\draw[curve](0,0)--(.5,.5)--(4,.5);\draw[curve,domain=4:5,samples=40]plot(\x,{.5+(\x-4)^2/2});\draw[curve](5,1)--(7,3);\foreach \x in {.5,4,5}{\draw[guide](\x,0)--(\x,1);\node[below]at(\x,0){$\x$};}\end{scope}

\end{tikzpicture}
\caption{一条经历等速段的相轨线，以及最终速度平方随初速度平方的分段函数}\label{mp4:fig:hw:l4-30-s}
\end{figure}

\end{solution}

\par\medskip\noindent\begin{minipage}{\textwidth}
\begin{example}[窄支座上的长杆与碰撞相图]\label{mp4:hw:l4-31}
均匀杆长$L$、质量$m$，中点置于宽$l\ll L$的固定支座上，支座只在两端与杆接触，接触期间不滑动。初始绕左端抬起小角$\varphi_0$后静止释放，见图\ref{mp4:fig:hw:l4-31-q}。(1)求完全弹性碰撞下主摆动周期的领先项。(2)若每次实际碰撞的法向恢复系数为$e=1-\delta$、$\delta\ll1$，建立转角衰减与碰撞次数的递推关系，用于判定最大转角何时减半，并讨论同一支点连续微小反弹的计数问题。
\end{example}

\begin{figure}[H]\centering
\begin{tikzpicture}[>=Stealth,line cap=round,line join=round]

\draw[fill=gray!15](-.45,-1)rectangle(.45,0);\draw[thick](-3,-.36)--(3,.4870588);\fill(-.45,0)circle(1.6pt)node[below left]{$A$};\fill(.45,0)circle(1.6pt)node[below right]{$B$};\draw[guide](-3,0)--(3,0);\node at(2.2,.18){$\varphi_0$};\draw[<->](-.45,-1.25)--(.45,-1.25)node[midway,below]{$l$};\node at(0,1){$L\gg l$};

\end{tikzpicture}
\caption{长杆在两支座端点之间换接触；碰撞后允许暂时离开两支点}\label{mp4:fig:hw:l4-31-q}
\end{figure}

\end{minipage}\par\medskip
\begin{solution}
在小角度范围内，以杆中点高度$Y$和转角$\varphi$描述构形。转动惯量$I=mL^2/12$，令
\[
 X=\sqrt{I/m}\,\varphi=\frac{L\varphi}{\sqrt{12}},
 \qquad\kappa=\frac{\sqrt3l}{L}.
\]
动能与接触限制成为
\[
 T=\frac m2(\dot X^2+\dot Y^2),\qquad
 Y\ge\kappa|X|,
\]
势能为$mgY$。因而小角度杆运动等价于重力场中一个质点在V形边界上运动。图\ref{mp4:fig:hw:l4-31-s}中的两条边分别对应左、右支点接触；离开两边时表示杆腾空。

最初绕左端落至水平，重力矩的领先项为 $mgl/2$。用 $I=mL^2/12$ 的领先项，角加速度大小为 $\alpha_0=(mgl/2)/I=6gl/L^2$。初角速度为零，两次积分，
\[
\begin{aligned}
\dot\varphi&=-\alpha_0\int_0^t\dd u=-\alpha_0[u]_0^t=-\alpha_0t,\\
\varphi&=\varphi_0-\alpha_0\int_0^t u\dd u
=\varphi_0-\alpha_0\left[\frac{u^2}{2}\right]_0^t
=\varphi_0-\frac12\alpha_0t^2.
\end{aligned}
\]
令 $\varphi(t_1)=0$，得
\[
 |\ddot\varphi|\simeq\frac{6gl}{L^2},\qquad
 t_1\simeq\sqrt{\frac{\varphi_0L^2}{3gl}}.
\]
到水平位置后，弹性碰撞通常使杆腾空，并在另一支点附近多次反弹，不能直接假定角速度反号。下面说明为何这些反弹不改变主摆动周期的领先项。令到达配置平面顶点前的速度为$(-u,-\kappa u)$，其中$u>0$。右边界上的切向加速度给出
\[
u=\frac{g\kappa}{1+\kappa^2}t_1.
\]
对左边界取内法线 $\vct n=(\kappa,1)/\sqrt{1+\kappa^2}$，碰前法向速度为 $\vct u^-\cdot\vct n=-2\kappa u/\sqrt{1+\kappa^2}$。弹性反射使法向速度反号而保持切向速度，故
\[
\vct u^+=\vct u^--2(\vct u^-\cdot\vct n)\vct n
=(-u,-\kappa u)+\frac{4\kappa u}{1+\kappa^2}(\kappa,1).
\]
整理为
\[
\vct u^+=\frac{u}{1+\kappa^2}
(-1+3\kappa^2,\ 3\kappa-\kappa^3).
\]
沿左边界向外的单位切向为$(-1,\kappa)/\sqrt{1+\kappa^2}$，故碰后的切向速率与法向速率分别为
\[
u_\parallel=\frac{u(1-\kappa^2)}{\sqrt{1+\kappa^2}},\qquad
u_\perp=\frac{2\kappa u}{\sqrt{1+\kappa^2}}.
\]
在同一条直边附近，弹性碰撞只反转法向速度；沿边运动始终以$g\kappa/\sqrt{1+\kappa^2}$减速。因此切向速度降为零所需时间为
\[
t_2=\frac{u(1-\kappa^2)}{g\kappa}
=\frac{1-\kappa^2}{1+\kappa^2}t_1
=t_1[1+O(\kappa^2)].
\]
法向反弹速度为$O(\kappa u)$，其对水平配置速度的贡献又多一个$\kappa$因子。故实际转角极值的时刻与$t_2$相差至多$O(\kappa^2t_1)$。在$\kappa\ll1$的领先阶，一个主摆动包含四段这样的慢行程，得到
\[
 \boxed{T\simeq4\sqrt{\frac{\varphi_0L^2}{3gl}}}.
\]
这一结果是窄支座、小转角下的主摆动时间。完整运动还含快速碰撞，任意$l/L$下的全部碰撞状态未必按此时间严格重复。

恢复系数小于1时，可在同一配置平面逐次求解。碰撞某一边$Y=s\kappa X$，$s=\pm1$，内法线为
\[
 \vct n_s=\frac{(-s\kappa,1)}{\sqrt{1+\kappa^2}}.
\]
若碰前、碰后配置速度为$\vct u^-,\vct u^+$，切向动量不变、法向相对速度反向并乘$e$，即
\[
 \boxed{\vct u^+=\vct u^--(1+e)(\vct u^-\cdot\vct n_s)\vct n_s}.
\]
损失的能量为 $m(1-e^2)(\vct u^-\cdot\vct n_s)^2/2$。两次碰撞之间满足 $\ddot X=0$、$\ddot Y=-g$。从碰后初速度积分，先得 $\dot X=u_{xj}$、$\dot Y=u_{yj}-gt$，再以 $(X_j,Y_j)$ 为初位置，得到
\[
 X=X_j+u_{xj}t,\quad
 Y=Y_j+u_{yj}t-gt^2/2.
\]
将它代入边界$Y=s\kappa X$，记$b_s=u_{yj}-s\kappa u_{xj}$、$c_s=Y_j-s\kappa X_j$，交会时间满足
\[
\frac g2t^2-b_st-c_s=0,
\qquad t_s=\frac{b_s+\sqrt{b_s^2+2gc_s}}g.
\]
舍去$t_s=0$以及不满足$s(X_j+u_{xj}t_s)\ge0$的候选，再从余下正根取最早者。碰前速度为$(u_{xj},u_{yj}-gt_s)$，位置由抛体式给出，随后代入速度映射。初值为$X_0=L\varphi_0/\sqrt{12}$、$Y_0=\kappa X_0$、速度为零；初始持续接触段须先沿右边下滑至顶点，不能以自由抛体式穿入边界。

上述递推同时确定碰撞后的能量$E_n$和实际转角极值。由接触边界可得转角的能量上界
\[
 |\varphi|\le\frac{2E_n}{mgl},\qquad
 E_n=E_0-\frac m2(1-e^2)\sum_{j=1}^n (u_{\perp j}^-)^2,
 \quad E_0=\frac{mgl\varphi_0}{2}.
\]
因此$E_n\le E_0/2$足以保证此后的最大转角不超过$\varphi_0/2$；若要求首次实际极值减半，应直接比较递推所得的各次极值。法向碰前速度由轨迹决定，并非每次都与总速度成同一比例，不能把转角直接写成$e^{2n}\varphi_0$。

能量保证转角减半的首次碰撞编号可明确写为
\[
 n_E=\inf\{n\ge0:E_n\le E_0/2\}.
\]
给定几何参数、恢复系数与初态，逐段求根、反射并累计能量损失即可求出这一编号；实际转角极值则由同一轨迹直接取得。若在达到该条件前已经出现无限碰撞聚积，按每次微小碰撞计数的编号取无穷。

在同一条边上，法向反弹速度依次为
\[
 u_\perp,\quad e u_\perp,\quad e^2u_\perp,\quad\ldots,
\]
重力朝向该边的法向分量大小为 $g_\perp=g/\sqrt{1+\kappa^2}$。只要反弹仍留在同一边，起跳法向速度为 $e^ju_\perp$ 的飞行满足 $z_\perp=e^ju_\perp t-g_\perp t^2/2$，正返回根为 $\Delta t_j=2e^ju_\perp/g_\perp$。当 $0<e<1$ 时，
\[
\sum_{j=0}^{N-1}\Delta t_j
=\frac{2u_\perp}{g_\perp}\frac{1-e^N}{1-e}
\longrightarrow\frac{2u_\perp}{g_\perp(1-e)}<\infty.
\]
时间总和有限而碰撞次数无限，这称为碰撞聚积。若在转角减半前出现这种情形，“每次微小接触都计一次”的次数就不再是有限整数。物理计数须约定有限接触时间或最小可分辨法向速度；达到该阈值后转入持续接触，再继续上述递推。由此可按给定的 $l/L,e$ 与计数精度确定次数。只给 $e=1-\delta$，不能推出一个与几何参数和接触计数无关的通用整数。


\begin{figure}[H]\centering
\begin{tikzpicture}[>=Stealth,line cap=round,line join=round]

\draw[axis](-3,0)--(3,0)node[right]{$X$};\draw[axis](0,-.2)--(0,2.7)node[above]{$Y$};\draw[thick](-2.6,1.3)--(0,0)--(2.6,1.3);\node[right]at(2,1.7){$Y=\kappa X$};\node[left]at(-2,1.7){$Y=-\kappa X$};\fill(2,1)circle(1.5pt)node[above left]{$(X_0,\kappa X_0)$};\draw[axis](2,1)--(.7,.35);\draw[curve,domain=0:1,samples=45]plot({-.8*\x},{.4*\x+1.8*\x*(1-\x)});\draw[curve,domain=0:1,samples=45]plot({-.8-.6*\x},{.4+.3*\x+1.6*\x*(1-\x)});\draw[curve,domain=0:1,samples=45]plot({-1.4-.4*\x},{.7+.2*\x+\x*(1-\x)});\node at(0,-.6){持续接触对应沿边运动，腾空对应边界上方的抛体段};

\end{tikzpicture}
\caption{杆运动的配置平面，边界倾角作示意放大：腾空段为抛体轨迹，接触时速度发生弹性或非弹性反射}\label{mp4:fig:hw:l4-31-s}
\end{figure}

\end{solution}

\section{守恒关系与参数平面的允许区域}


动量和能量有时把同一过程限制在两条不同曲线上。例如速度变量$v,\omega$满足
\[
 \frac12mv^2+\frac12I\omega^2=E,\qquad
 av+b\omega=P.
\]
第一式在速度平面中是椭圆，第二式是直线，允许的末态必须同时位于两者上。图\ref{mp4:fig:hw:conservation-intersection}中的交点可能有两个；还须由碰撞前后接近或分离、接触方向及运动阶段选取。令$V=\sqrt m\,v$、$W=\sqrt I\,\omega$，能量曲线化为$V^2+W^2=2E$，但线性守恒关系也必须同时变换。

\begin{figure}[H]\centering
\begin{tikzpicture}[>=Stealth,line cap=round,line join=round]
\draw[axis](-2.2,0)--(2.2,0)node[right]{$v$};\draw[axis](0,-1.8)--(0,2)node[above]{$\omega$};\draw[curve](0,0)ellipse[x radius=1.8,y radius=1.3];\draw[thick](-1.6,1.7)--(1.9,-.75);\fill(-0.8230544,1.1561380)circle(1.5pt)node[above left]{$P_1$};\fill(1.6257392,-0.5580175)circle(1.5pt)node[below right]{$P_2$};\node[right,align=left]at(3,.6){能量：二次约束\\[6pt]动量：线性约束\\[6pt]物理末态还须满足过程条件};
\end{tikzpicture}
\caption{能量与线性动量关系必须同时成立；交点的代数存在性不等于过程可达性}\label{mp4:fig:hw:conservation-intersection}
\end{figure}


\subsection{缩放后的直线距离与两种候选末态}
把速度平面中的能量椭圆变成圆后，线性约束也随之改变。记
\[A=\frac a{\sqrt m},\quad B=\frac b{\sqrt I},
\quad \vct u=(V,W),\quad \vct n=(A,B),
\]
则两条守恒关系写成$|\vct u|^2=2E$、$\vct n\cdot\vct u=P$。假设$A^2+B^2>0$，直线上距原点最近的点为
\[\vct u_c=\frac{P}{A^2+B^2}(A,B),\qquad
|\vct u_c|^2=\frac{P^2}{A^2+B^2}.
\]
直线与圆有交点的条件因此为
\[\boxed{P^2\le2E(A^2+B^2)}.
\]
等号对应相切，严格不等号一般给两个候选点。沿直线的单位方向为 $\vct e=(-B,A)/\sqrt{A^2+B^2}$，令 $\vct u=\vct u_c+t\vct e$。因 $\vct u_c\cdot\vct e=0$，能量方程成为 $|\vct u_c|^2+t^2=2E$，解出 $t$ 的两根，全部代数末态即为
\[\vct u_\pm=\vct u_c
\pm\sqrt{2E-\frac{P^2}{A^2+B^2}}\,\vct e.
\]
将$V,W$分别除以$\sqrt m,\sqrt I$，才恢复实际的$v,\omega$。图\ref{mp4:fig:series-conservation-circle}中圆半径、垂足和弦长分别进入上式，横纵坐标的缩放不能只用于作图而在回代时遗漏。

\begin{figure}[H]\centering
\begin{tikzpicture}[>=Stealth,line cap=round,line join=round]
\draw[axis](-1.4,0)--(2,0)node[right]{$V$};\draw[axis](0,-1.4)--(0,2)node[above]{$W$};
\draw[curve](0,0)circle(1.35);\draw[thick](-.45,1.8)--(1.8,-.45);
\coordinate(K)at(.675,.675);\draw[guide](0,0)--(K);\fill(K)circle(1.4pt)node[above right]{$\vct u_c$};
\fill(0,1.35)circle(1.6pt)node[above left]{$\vct u_+$};\fill(1.35,0)circle(1.6pt)node[below right]{$\vct u_-$};
\draw[->](.675,.675)--(.23,1.12)node[midway,left]{$\vct e$};\node[below left]at(0,0){$O$};
\node[right,align=left]at(2.4,.9){圆：能量固定\\[5pt]直线：线性守恒量固定\\[5pt]垂足决定可行能量下界};
\end{tikzpicture}
\caption{按缩放后的坐标作圆与直线，两个交点仅是代数候选；过程条件继续选择物理末态。}\label{mp4:fig:series-conservation-circle}
\end{figure}

给定线性守恒量$P$时，所允许的最低能量为$E_{\min}=P^2/[2(A^2+B^2)]$。低于此值没有满足两式的状态；等于此值只剩一个状态。高于此值时，两个交点通常分别对应碰撞前、后或接近、分离的分支。题设若允许能量损失，应先把能量等式改成适当的不等式或恢复系数关系，不能继续用同一条固定能量圆。

代数关系给出允许的状态集合。确定实际运动时，还须结合相轨线的时间方向、冲量的接触方向及单侧约束的反力条件。选根时应把这些条件代回原来的物理变量，不能仅以图上某一点位于右半平面就认定它是末态。




把平衡、接触或守恒条件写成等式与不等式，可以在参数平面中划出不同运动状态的区域。边界上的等号表示临界状态；区域的含义仍须由具体力学条件确定。

\par\medskip\noindent\begin{minipage}{\textwidth}
\begin{example}[水平拔河的摩擦容量]\label{mp4:hw:l4-24}
两人在水平地面上以轻绳拔河，质量为$m_1,m_2$，鞋与地面的最大静摩擦系数为$\mu_1,\mu_2$。忽略身体姿态引起的倾覆，假定两人都能施加足够大的肌肉拉力，并以能承受的最大绳张力判断摩擦容量。讨论质量和摩擦系数对胜负的影响，画出临界边界。见图\ref{mp4:fig:hw:l4-24-q}。
\end{example}

\begin{figure}[H]\centering
\begin{tikzpicture}[>=Stealth,line cap=round,line join=round]

\draw[thick](-.5,0)--(6,0);\draw[fill=gray!10](.3,0)rectangle(1.1,1.1);\draw[fill=gray!10](4.2,0)rectangle(5,1.1);\draw[thick](1.1,.7)--(4.2,.7);\draw[axis](1.1,.7)--(2,.7)node[above]{$T$};\draw[axis](4.2,.7)--(3.3,.7)node[above]{$T$};\node[above]at(.7,1.1){$m_1,\mu_1$};\node[above]at(4.6,1.1){$m_2,\mu_2$};\draw[axis](.7,.1)--(-.3,.1)node[above]{$f_1$};\draw[axis](4.6,.1)--(5.6,.1)node[above]{$f_2$};

\end{tikzpicture}
\caption{双方受相同绳张力，各自依靠地面摩擦保持位置}\label{mp4:fig:hw:l4-24-q}
\end{figure}

\end{minipage}\par\medskip
\begin{solution}
两人的法向力分别为$m_1g,m_2g$。保持静止要求
\[
 T\le\mu_1m_1g,\qquad T\le\mu_2m_2g.
\]
在所规定的容量模型中，临界线为
\[
 \boxed{\mu_1m_1=\mu_2m_2}.
\]
若左边较大，第一人能承受第二人无法保持静止的张力；反之第二人占优。图\ref{mp4:fig:hw:l4-24-s}左图以$\mu_1,\mu_2$作坐标，边界为斜率$m_1/m_2$的直线。若改以质量比和摩擦系数比表示，边界成为$ (m_1/m_2)(\mu_1/\mu_2)=1$，即右图的双曲线。

这里比较的是最大静摩擦容量。若双方都不主动收绳，且张力未超过任一上限，两人可以共同静止。开始运动后的加速度还取决于收绳方式、动摩擦和身体运动。


\begin{figure}[H]\centering
\begin{tikzpicture}[>=Stealth,line cap=round,line join=round]

\draw[axis](0,0)--(3.2,0)node[right]{$\mu_1$};\draw[axis](0,0)--(0,2.8)node[above]{$\mu_2$};\draw[curve](0,0)--(2.5,2.5);\node at(1,2.3){第二人占优};\node at(2.1,.6){第一人占优};\begin{scope}[shift={(5.7,0)}]\draw[axis](0,0)--(3,0)node[right]{$m_1/m_2$};\draw[axis](0,0)--(0,2.8)node[above]{$\mu_1/\mu_2$};\draw[curve,domain=.4:2.7,samples=70]plot(\x,{1/\x});\node at(1.9,1.9){第一人占优};\node at(1.5,.25){第二人占优};\end{scope}

\end{tikzpicture}
\caption{同一摩擦容量边界在两种参数坐标中分别为直线与双曲线}\label{mp4:fig:hw:l4-24-s}
\end{figure}

\end{solution}

\par\medskip\noindent\begin{minipage}{\textwidth}
\begin{example}[两种斜面上的拔河与静止区间]\label{mp4:hw:l4-25}
质量$2m$的第一人在$30^\circ$斜面上，质量$m$的第二人在$60^\circ$斜面上，轻绳跨过光滑定滑轮，见图\ref{mp4:fig:hw:l4-25-q}。最大静摩擦系数分别为$\mu_1,\mu_2\ge0$，绳力沿各自斜面。采用上一例的摩擦容量模型，讨论两人的优势区域，并先判断能否共同静止。
\end{example}

\begin{figure}[H]\centering
\begin{tikzpicture}[>=Stealth,line cap=round,line join=round]

\draw[thick](-3,0)--(0,1.732)--(1,0)--cycle;\draw[curve](0,1.9)circle(.13);\draw[thick](-2,.72)--(-.06,2.02)--(.75,.65);\node[left]at(-1.8,1.3){$2m,\mu_1$};\node[right]at(.85,1.2){$m,\mu_2$};\node at(-2.3,.2){$30^\circ$};\node at(.65,.2){$60^\circ$};\node[above]at(0,2.1){$T$};

\end{tikzpicture}
\caption{两斜面上的重力沿坡分量不同，静摩擦可以沿坡向上或向下}\label{mp4:fig:hw:l4-25-q}
\end{figure}

\end{minipage}\par\medskip
\begin{solution}
令$t=T/(mg)$。第一人沿坡向下的重力为$mg$，最大静摩擦为$\sqrt3\mu_1mg$；第二人的相应量为$\sqrt3mg/2$、$\mu_2mg/2$。故各自可保持静止的张力区间为
\[
 I_1=[1-\sqrt3\mu_1,1+\sqrt3\mu_1],\qquad
 I_2=[(\sqrt3-\mu_2)/2,(\sqrt3+\mu_2)/2].
\]
共同静止要求$I_1\cap I_2\cap[0,\infty)\ne\varnothing$。在第一象限，另一方向的区间分离不可能发生，因而条件简化为
\[
 \boxed{2\sqrt3\mu_1+\mu_2\ge2-\sqrt3}.
\]
低于这条下降直线时，第一人的沿坡重力已经大到双方不能共同保持静止，运动趋势是第一人下坡、第二人上坡。该区域内的起动趋势由重力造成。

在可共同静止的区域中，比较各自所能承受的最大张力：
\[
 1+\sqrt3\mu_1\quad\hbox{与}\quad(\sqrt3+\mu_2)/2.
\]
所以另一条临界线为
\[
 \boxed{\mu_2=2\sqrt3\mu_1+2-\sqrt3}.
\]
图\ref{mp4:fig:hw:l4-25-s}中，上升直线上方第二人容量较大；该线下方且在共同静止区内，第一人容量较大。上升边界表示上坡滑动的最大承受张力相同，下降边界则表示两静止区间刚好接触，作图时须按这两种不同的临界过程判断区域。


\begin{figure}[H]\centering
\begin{tikzpicture}[>=Stealth,line cap=round,line join=round]

\draw[axis](0,0)--(3.5,0)node[right]{$\mu_1$};\draw[axis](0,0)--(0,3.5)node[above]{$\mu_2$};\fill[gray!18](0,0)--(.268,0)--(0,.928)--cycle;\draw[curve](0,.928)--(.7,3.353);\draw[curve](0,.928)--(.268,0);\node[right]at(.8,3.15){$\mu_2=2\sqrt3\mu_1+2-\sqrt3$};\node[align=left]at(1.7,1.3){第一人容量较大};\node at(.6,2.5){第二人容量较大};\node[right]at(.05,.3){重力滑动};\node[left]at(0,.928){$2-\sqrt3$};\node[below]at(.268,0){$\dfrac{2-\sqrt3}{2\sqrt3}$};

\end{tikzpicture}
\caption{两条边界围成不同的静止与摩擦容量区域；示意图纵横轴比例分别选取}\label{mp4:fig:hw:l4-25-s}
\end{figure}

\end{solution}

\par\medskip\noindent\begin{minipage}{\textwidth}
\begin{example}[三支脚圆盘的滑动区域]\label{mp4:hw:l4-26}
均匀水平圆盘由相同的三个小支脚$A,B,C$支承，脚到圆心的距离均为$l$，两两夹角$120^\circ$。取$B=(-l,0)$、$A=(l/2,\sqrt3l/2)$、$C=(l/2,-\sqrt3l/2)$。各脚与地面摩擦系数均为$\mu$，支脚高度忽略，水平力引起的竖直载荷转移也忽略。在盘面点$P=(\rho\cos\theta,\rho\sin\theta)$施加水平力，方向为绕圆心逆时针的切向，缓慢增大至开始运动，见图\ref{mp4:fig:hw:l4-26-q}。求恰有一脚保持静止、另两脚开始滑动的作用点区域，以及所有可能的$\rho,\theta$边界。实际作用点还须位于盘面内。
\end{example}

\begin{figure}[H]\centering
\begin{tikzpicture}[>=Stealth,line cap=round,line join=round]

\draw[curve](0,0)circle(2);\foreach \x/\y/\n in {.5/.866/A,-1/0/B,.5/-.866/C}{\fill(\x,\y)circle(2pt)node[above left]{$\n$};\draw[guide](0,0)--(\x,\y);}\fill(0,0)circle(1.3pt)node[below left]{$O$};\coordinate(P)at(1.2,.6);\draw[guide](0,0)--(P)node[midway,above right]{$\rho$};\fill(P)circle(1.5pt)node[right]{$P$};\draw[axis,thick](P)--(.75,1.5)node[above]{$F$};\draw(.5,0)arc(0:26.565:.5);\node at(.75,.16){$\theta$};

\end{tikzpicture}
\caption{三脚等载支承；外力与$OP$垂直并产生逆时针力矩}\label{mp4:fig:hw:l4-26-q}
\end{figure}

\end{minipage}\par\medskip
\begin{solution}
设圆盘与支脚的总质量为$M$。按忽略载荷转移的题设，三脚法向力均为$Mg/3$，记摩擦上限为$f=\mu Mg/3$。在准静态起动临界状态，先假设$B$不动，刚体趋于绕$B$逆时针转动。$A,C$的速度分别垂直于$BA,BC$，所以两脚的极限摩擦为
\[
 \vct f_A=(f/2,-\sqrt3f/2),\qquad
 \vct f_C=(-f/2,-\sqrt3f/2).
\]
关于$B$取矩，得到
\[
 F(\rho+l\cos\theta)=2\sqrt3fl,
 \qquad F=\frac{2\sqrt3fl}{\rho+l\cos\theta}.
\]
由水平合力为零，$B$处所需静摩擦为
\[
 \vct f_B=(F\sin\theta,\sqrt3f-F\cos\theta).
\]
图\ref{mp4:fig:hw:l4-26-s}左图标出这三种方向。由$B$处静摩擦不能超过上限，
\[
|\vct f_B|^2=F^2+3f^2-2\sqrt3fF\cos\theta\le f^2.
\]
令$t=\rho/l+\cos\theta>0$，则$F/f=2\sqrt3/t$。代入上式并乘以$t^2/(2f^2)$，化为
\[
 t^2-6t\cos\theta+6\le0.
\]
二次不等式的两根为 $t_\pm=3\cos\theta\pm\sqrt{9\cos^2\theta-6}$，而 $\rho/l=t-\cos\theta$。判别式要求 $\cos^2\theta\ge2/3$，再由 $F>0$ 即 $t>0$ 排除 $\cos\theta<0$ 的一支。因此，$B$ 不动区域为
\[
 \boxed{|\theta|\le\theta_*:=\arccos\sqrt{2/3}},
\]
\[
 \boxed{2\cos\theta-\sqrt{9\cos^2\theta-6}
 \le\frac\rho l\le
 2\cos\theta+\sqrt{9\cos^2\theta-6}}.
\]
为确定整个区域的径向极值，记 $c=\cos\theta\in[\sqrt{2/3},1]$。根式的链式导数为 $\dd\sqrt{9c^2-6}/\dd c=18c/[2\sqrt{9c^2-6}]$，所以两径向边界的导数分别为 $2-9c/\sqrt{9c^2-6}<0$ 和 $2+9c/\sqrt{9c^2-6}>0$。最小内边界与最大外边界都在 $c=1$ 处取得。因此在未经盘面裁切的候选区域中，
\[
 \rho_{\min}=(2-\sqrt3)l,\quad
 \rho_{\max}=(2+\sqrt3)l,\quad
 \theta_{\max,\min}=\pm35.2644^\circ.
\]
将该区域绕$O$分别转过$120^\circ$、$-120^\circ$，得到另外两脚不动的区域。图右的三个区域互不重叠，最后须与圆盘的实际半径范围相交。边界上不动脚也达到静摩擦极限；在区域外，该脚所需静摩擦超过上限，须转用三脚均滑动的速度场。


\begin{figure}[H]\centering
\begin{tikzpicture}[>=Stealth,line cap=round,line join=round]

\coordinate(B)at(-1,0);\coordinate(A)at(.5,.866);\coordinate(C)at(.5,-.866);\draw[guide](B)--(A)--(C)--cycle;\foreach \p in {A,B,C}{\fill(\p)circle(1.5pt);}\node[above]at(A){$A$};\node[below left]at(B){$B$};\node[right]at(C){$C$};\draw[axis](A)--(1,.0)node[right]{$\vct f_A$};\draw[axis](C)--(0,-1.732)node[below]{$\vct f_C$};\draw[axis](B)--(-1,.6)node[left]{$\vct f_B$};\node at(-.9,-1){$B$静止};
\begin{scope}[shift={(5,0)},scale=.48]
\foreach \ang/\name in {0/B,120/C,240/A}{\begin{scope}[rotate=\ang]\fill[gray!18]plot[domain=-35.2643:35.2643,samples=85]({(2*cos(\x)+sqrt(max(0,9*cos(\x)^2-6)))*cos(\x)},{(2*cos(\x)+sqrt(max(0,9*cos(\x)^2-6)))*sin(\x)})--plot[domain=35.2643:-35.2643,samples=85]({(2*cos(\x)-sqrt(max(0,9*cos(\x)^2-6)))*cos(\x)},{(2*cos(\x)-sqrt(max(0,9*cos(\x)^2-6)))*sin(\x)})--cycle;\draw[curve,domain=-35.2643:35.2643,samples=85]plot({(2*cos(\x)+sqrt(max(0,9*cos(\x)^2-6)))*cos(\x)},{(2*cos(\x)+sqrt(max(0,9*cos(\x)^2-6)))*sin(\x)});\draw[curve,domain=-35.2643:35.2643,samples=85]plot({(2*cos(\x)-sqrt(max(0,9*cos(\x)^2-6)))*cos(\x)},{(2*cos(\x)-sqrt(max(0,9*cos(\x)^2-6)))*sin(\x)});\node at(2,0){$\name$不动};\end{scope}}\fill(0,0)circle(2pt)node[below]{$O$};\end{scope}

\end{tikzpicture}
\caption{指定$B$静止时的摩擦方向，以及轮换三脚得到的三个作用点区域}\label{mp4:fig:hw:l4-26-s}
\end{figure}

\end{solution}

\par\medskip\noindent\begin{minipage}{\textwidth}
\begin{example}[三人拔河的三角形参数图]\label{mp4:hw:l4-27}
三名质量均为$m$的人位于等边三角形的顶点，鞋与地面的摩擦系数为$\mu_1,\mu_2,\mu_3$。采用摩擦容量模型：讨论一人守住位置，另外两人沿通向他的两条边被拉动时，守住者所需的静摩擦；各拉索张力可调，第三边可不受拉。以此划分谁具有单独守住两人的摩擦容量，见图\ref{mp4:fig:hw:l4-27-q}。
\end{example}

\begin{figure}[H]\centering
\begin{tikzpicture}[>=Stealth,line cap=round,line join=round]

\coordinate(A)at(0,2.6);\coordinate(B)at(-1.5,0);\coordinate(C)at(1.5,0);\draw[thick](A)--(B)--(C)--cycle;\fill(A)circle(2pt)node[above]{$1,\mu_1$};\fill(B)circle(2pt)node[below left]{$2,\mu_2$};\fill(C)circle(2pt)node[below right]{$3,\mu_3$};\draw[axis](-1.15,.6)--(-.6,1.55);\draw[axis](1.15,.6)--(.6,1.55);\node at(0,.9){$60^\circ$};

\end{tikzpicture}
\caption{一人守住位置、两人沿相邻边移动的容量比较构形}\label{mp4:fig:hw:l4-27-q}
\end{figure}

\end{minipage}\par\medskip
\begin{solution}
以第一人为守住者。另两人的极限摩擦大小为$\mu_2mg$、$\mu_3mg$，方向夹角为$60^\circ$。整个三人系统的水平合力平衡要求第一人的静摩擦抵消它们的矢量和，故
\[
 f_1=mg\sqrt{\mu_2^2+\mu_3^2+\mu_2\mu_3}.
\]
能守住的条件为
\[
 \boxed{\mu_1^2\ge\mu_2^2+\mu_3^2+\mu_2\mu_3}.
\]
另外两人的判据由轮换下标得到。当$\mu_1+\mu_2+\mu_3>0$时，令$u_i=\mu_i/(\mu_1+\mu_2+\mu_3)$，则$u_i\ge0$且$u_1+u_2+u_3=1$，可把每组参数映射到一个等边三角形。若三系数全为零，三人的摩擦容量均为零，归一化坐标不再定义。图\ref{mp4:fig:hw:l4-27-s}中，三个顶点附近是相应个人的容量优势区；在中间区域$D$，没有一人能单独满足守住条件。

第一人的边界也可写成 $1-2(u_2+u_3)+u_2u_3=0$。令 $u_2=(1-u_1)(1+q)/2$、$u_3=(1-u_1)(1-q)/2$，原边界成为 $u_1^2=(1-u_1)^2(3+q^2)/4$。因 $0\le u_1\le1$，开方须取非负支，得到 $2u_1=(1-u_1)\sqrt{3+q^2}$，移项即得
\[
 u_1=\frac{\sqrt{3+q^2}}{2+\sqrt{3+q^2}},\qquad -1\le q\le1.
\]
这给出可直接作图的曲线参数式。三条边界是双曲线弧，分别在三角形各边中点相接。当$\mu_3=0$时，第一人与第二人的边界退化为$\mu_1=\mu_2$，回到两人拔河。容量判据依赖本题规定的拉动构形；若未指定拉索张力、收绳策略及运动方式，三人的实际胜负并不能仅由三个摩擦系数唯一决定。


\begin{figure}[H]\centering
\begin{tikzpicture}[>=Stealth,line cap=round,line join=round]

\draw[thick](-1.6,0)--(1.6,0)--(0,2.771)--cycle;\node[above]at(0,2.771){$u_1=1$};\node[below left]at(-1.6,0){$u_2=1$};\node[below right]at(1.6,0){$u_3=1$};
\foreach \ang in {0,120,240}{\begin{scope}[shift={(0,.9237)},rotate=\ang]\fill[gray!20](0,1.8473)--plot[domain=-1:1,samples=90]({-1.6*\x*(1-sqrt(3+(\x)^2)/(2+sqrt(3+(\x)^2)))},{2.771*sqrt(3+(\x)^2)/(2+sqrt(3+(\x)^2))-.9237})--cycle;\draw[curve,domain=-1:1,samples=90]plot({-1.6*\x*(1-sqrt(3+(\x)^2)/(2+sqrt(3+(\x)^2)))},{2.771*sqrt(3+(\x)^2)/(2+sqrt(3+(\x)^2))-.9237});\end{scope}}\node at(0,1.95){1};\node at(-.95,.35){2};\node at(.95,.35){3};\node at(0,.9){$D$};\node[right]at(2,1.25){$D$：无单人优势};

\end{tikzpicture}
\caption{归一化摩擦系数的三角形表示；边界由两个阻力的矢量和决定}\label{mp4:fig:hw:l4-27-s}
\end{figure}

\end{solution}

\par\medskip\noindent\begin{minipage}{\textwidth}
\begin{example}[三体衰变的能量允许区域]\label{mp4:hw:l4-28}
一静止粒子衰变为三个质量相等的碎片，总动量为零。用$E_i$表示各碎片能量，并在等边三角形内以重心坐标$u_i=E_i/E$、$\sum u_i=1$标记状态，见图\ref{mp4:fig:hw:l4-28-q}。分别求：(1)非相对论情形，$E$为释放的总动能；(2)碎片静质量可忽略、速度近光速时，$E$为总出射能量。
\end{example}

\begin{figure}[H]\centering
\begin{tikzpicture}[>=Stealth,line cap=round,line join=round]

\coordinate(A)at(0,2.6);\coordinate(B)at(-1.5,0);\coordinate(C)at(1.5,0);\draw[thick](A)--(B)--(C)--cycle;\coordinate(P)at(.15,1);\fill(P)circle(1.5pt)node[right]{$(u_1,u_2,u_3)$};\node[above]at(A){$E_1=E$};\node[below left]at(B){$E_2=E$};\node[below right]at(C){$E_3=E$};\draw[guide](P)--(.15,0);\draw[guide](P)--(-.77,1.334);\draw[guide](P)--(.946,1.04);

\end{tikzpicture}
\caption{三个能量非负且和为定值，因此可用一个三角形表示}\label{mp4:fig:hw:l4-28-q}
\end{figure}

\end{minipage}\par\medskip
\begin{solution}
总动量为零要求三个动量矢量首尾闭合。三条边长$p_1,p_2,p_3$能组成三角形，当且仅当每条不大于另外两条之和。

非相对论情形$p_i=\sqrt{2mE_i}$。记$\vct p_2,\vct p_3$的夹角为$\gamma$，由$\vct p_1=-(\vct p_2+\vct p_3)$，
\[
E_1=E_2+E_3+2\sqrt{E_2E_3}\cos\gamma.
\]
存在实夹角的条件是$(E_1-E_2-E_3)^2\le4E_2E_3$；当$E_2E_3=0$时它退为相应的等式，不必作除法。展开并移项，得到对三下标对称的条件
\[
 2(E_1E_2+E_2E_3+E_3E_1)\ge E_1^2+E_2^2+E_3^2.
\]
利用总动能守恒，化为
\[
 \boxed{u_1^2+u_2^2+u_3^2\le\frac12},\qquad
 \sum_i(u_i-1/3)^2\le\frac16.
\]
在三角形所在平面中，这是一圆盘。以$u_1=0$为例，条件$u_2+u_3=1$给出$u_2^2+u_3^2\ge1/2$，与允许条件合用，只有$u_2=u_3=1/2$成立；故圆盘恰与各边中点相切，见图\ref{mp4:fig:hw:l4-28-s}左图。要找单个碎片的最大动能，利用
\[
u_2^2+u_3^2\ge\frac{(1-u_1)^2}{2},\qquad
u_1^2+\frac{(1-u_1)^2}{2}\le\frac12,
\]
得到$u_1(3u_1-2)\le0$，即$u_1\le2/3$。等号时$u_2=u_3=1/6$，另两个碎片的动量同向，合动量与第一个碎片反向。

质量可忽略时$p_i=E_i/c$，三角形不等式直接成为$E_i\le E-E_i$。所以
\[
 \boxed{0\le u_i\le\frac12,\quad u_1+u_2+u_3=1}.
\]
允许区域是连接原三角形三边中点的倒置小三角形。两种区域的边界都对应三个动量共线。能量与动量大小的关系改变后，同一动量闭合条件便给出不同的能量区域。圆盘内部或小三角形内部的任一点，都能由相应边长组成动量三角形，实现能量、动量同时守恒。


\begin{figure}[H]\centering
\begin{tikzpicture}[>=Stealth,line cap=round,line join=round]

\begin{scope}\draw[thick](-1.5,0)--(1.5,0)--(0,2.598)--cycle;\fill[gray!20](0,.866)circle(.866);\draw[curve](0,.866)circle(.866);\node at(0,1){非相对论};\node at(0,-.45){$\sum u_i^2\le1/2$};\end{scope}
\begin{scope}[shift={(5,0)}]\draw[thick](-1.5,0)--(1.5,0)--(0,2.598)--cycle;\fill[gray!20](0,0)--(.75,1.299)--(-.75,1.299)--cycle;\draw[curve](0,0)--(.75,1.299)--(-.75,1.299)--cycle;\node at(0,.8){近光速};\node at(0,-.45){$u_i\le1/2$};\end{scope}

\end{tikzpicture}
\caption{相同的总能量与总动量约束，在两种能量动量关系下产生不同区域}\label{mp4:fig:hw:l4-28-s}
\end{figure}

\end{solution}

\section{习题与解答}

\par\medskip\noindent\begin{minipage}{\textwidth}
\begin{exercise}[恒力作用下的平移势阱]\label{mp4:exercise:l4-h1}
弹簧振子原来静止在 $x=0$，此处弹簧无形变。$t=0$ 起施加恒定向右的力 $F$，以后保持不变，见图\ref{mp4:fig:l4-h1-q}。求新的平衡位置、两个转向位置，以及首次到达新平衡位置时的速度和时间。
\end{exercise}

\begin{figure}[H]\centering
\begin{tikzpicture}
\draw[curve](-2.4,-.35)--(-2.4,.65);\draw(-2.4,.1)--(-2.15,.1)--(-1.97,.35)--(-1.7,-.15)--(-1.43,.35)--(-1.16,-.15)--(-.89,.35)--(-.62,.1)--(-.3,.1);
\draw[fill=white](-.3,-.25)rectangle(.35,.45);\node at(.03,.1){$m$};\draw(-2.6,-.3)--(2.3,-.3);
\draw[axis](-1.1,-.8)--(2.5,-.8)node[right]{$x$};\draw[guide](0,-.3)--(0,-.9);\node[below]at(0,-.8){$0$};\node[above]at(-1.4,.55){$k$};
\draw[axis](.35,.1)--(1.6,.1)node[above]{$F$};
\end{tikzpicture}
\caption{恒力突然施加于原来静止的弹簧振子。}\label{mp4:fig:l4-h1-q}
\end{figure}
\end{minipage}\par\medskip
\begin{solution}
恒力开始作用后，运动方程为 $m\ddot x=F-kx$。令 $b=F/k$、$\omega=\sqrt{k/m}$，把弹性势能与恒力势能合并，得到
\[
U(x)=\frac12kx^2-Fx
=\frac12k(x-b)^2-\frac12kb^2.
\]
由 $U'(x)=0$ 得新的平衡位置 $x=b$。图\ref{mp4:fig:l4-h1-s}中的相椭圆中心因此从原点移到 $(b,0)$。

有限恒力的突然接通不产生冲量，位置和速度仍为 $x(0)=v(0)=0$。以接通后的时间无关势能计算，初态能量为 $E=0$，所以
\[
\frac12mv^2+\frac12k(x-b)^2=\frac12kb^2.
\]
令 $v=0$ 得 $x-b=\pm b$，两个转向位置为图中的 A：$x=0$ 和 B：$x=2b$。初始合力为 $F>0$，相点从 A 沿上支向右运动。

\begin{figure}[H]\centering
\begin{tikzpicture}
\draw[axis](-0.6,0)--(4.7,0)node[right]{$x/b$};
\draw[axis](0,-1.7)--(0,1.8)node[above]{$v/(\omega b)$};
\draw[flow] plot coordinates{(0.00000,-0.00000) (0.00122,0.04188) (0.00487,0.08371) (0.01096,0.12543) (0.01946,0.16701) (0.03038,0.20838) (0.04370,0.24949) (0.05941,0.29031) (0.07748,0.33076) (0.09789,0.37082) (0.12061,0.41042) (0.14563,0.44953) (0.17291,0.48808) (0.20241,0.52605) (0.23410,0.56337) (0.26795,0.60000) (0.30390,0.63590) (0.34192,0.67103) (0.38197,0.70534) (0.42398,0.73879) (0.46791,0.77135) (0.51371,0.80296) (0.56132,0.83359) (0.61068,0.86321) (0.66174,0.89177) (0.71442,0.91925) (0.76868,0.94561) (0.82443,0.97082) (0.88161,0.99485) (0.94016,1.01766) (1.00000,1.03923) (1.06106,1.05954) (1.12326,1.07855) (1.18653,1.09625) (1.25079,1.11262) (1.31596,1.12763) (1.38197,1.14127) (1.44873,1.15351) (1.51616,1.16435) (1.58418,1.17378) (1.65270,1.18177) (1.72165,1.18832) (1.79094,1.19343) (1.86049,1.19708) (1.93020,1.19927) (2.00000,1.20000) (2.06980,1.19927) (2.13951,1.19708) (2.20906,1.19343) (2.27835,1.18832) (2.34730,1.18177) (2.41582,1.17378) (2.48384,1.16435) (2.55127,1.15351) (2.61803,1.14127) (2.68404,1.12763) (2.74921,1.11262) (2.81347,1.09625) (2.87674,1.07855) (2.93894,1.05954) (3.00000,1.03923) (3.05984,1.01766) (3.11839,0.99485) (3.17557,0.97082) (3.23132,0.94561) (3.28558,0.91925) (3.33826,0.89177) (3.38932,0.86321) (3.43868,0.83359) (3.48629,0.80296) (3.53209,0.77135) (3.57602,0.73879) (3.61803,0.70534) (3.65808,0.67103) (3.69610,0.63590) (3.73205,0.60000) (3.76590,0.56337) (3.79759,0.52605) (3.82709,0.48808) (3.85437,0.44953) (3.87939,0.41042) (3.90211,0.37082) (3.92252,0.33076) (3.94059,0.29031) (3.95630,0.24949) (3.96962,0.20838) (3.98054,0.16701) (3.98904,0.12543) (3.99513,0.08371) (3.99878,0.04188) (4.00000,0.00000) (3.99878,-0.04188) (3.99513,-0.08371) (3.98904,-0.12543) (3.98054,-0.16701) (3.96962,-0.20838) (3.95630,-0.24949) (3.94059,-0.29031) (3.92252,-0.33076) (3.90211,-0.37082) (3.87939,-0.41042) (3.85437,-0.44953) (3.82709,-0.48808) (3.79759,-0.52605) (3.76590,-0.56337) (3.73205,-0.60000) (3.69610,-0.63590) (3.65808,-0.67103) (3.61803,-0.70534) (3.57602,-0.73879) (3.53209,-0.77135) (3.48629,-0.80296) (3.43868,-0.83359) (3.38932,-0.86321) (3.33826,-0.89177) (3.28558,-0.91925) (3.23132,-0.94561) (3.17557,-0.97082) (3.11839,-0.99485) (3.05984,-1.01766) (3.00000,-1.03923) (2.93894,-1.05954) (2.87674,-1.07855) (2.81347,-1.09625) (2.74921,-1.11262) (2.68404,-1.12763) (2.61803,-1.14127) (2.55127,-1.15351) (2.48384,-1.16435) (2.41582,-1.17378) (2.34730,-1.18177) (2.27835,-1.18832) (2.20906,-1.19343) (2.13951,-1.19708) (2.06980,-1.19927) (2.00000,-1.20000) (1.93020,-1.19927) (1.86049,-1.19708) (1.79094,-1.19343) (1.72165,-1.18832) (1.65270,-1.18177) (1.58418,-1.17378) (1.51616,-1.16435) (1.44873,-1.15351) (1.38197,-1.14127) (1.31596,-1.12763) (1.25079,-1.11262) (1.18653,-1.09625) (1.12326,-1.07855) (1.06106,-1.05954) (1.00000,-1.03923) (0.94016,-1.01766) (0.88161,-0.99485) (0.82443,-0.97082) (0.76868,-0.94561) (0.71442,-0.91925) (0.66174,-0.89177) (0.61068,-0.86321) (0.56132,-0.83359) (0.51371,-0.80296) (0.46791,-0.77135) (0.42398,-0.73879) (0.38197,-0.70534) (0.34192,-0.67103) (0.30390,-0.63590) (0.26795,-0.60000) (0.23410,-0.56337) (0.20241,-0.52605) (0.17291,-0.48808) (0.14563,-0.44953) (0.12061,-0.41042) (0.09789,-0.37082) (0.07748,-0.33076) (0.05941,-0.29031) (0.04370,-0.24949) (0.03038,-0.20838) (0.01946,-0.16701) (0.01096,-0.12543) (0.00487,-0.08371) (0.00122,-0.04188) (0.00000,-0.00000)};
\fill(0,0)circle(1.2pt)node[below left]{$A$};
\fill(4,0)circle(1.2pt)node[above right]{$B$};
\fill(2,1.2)circle(1.2pt)node[above]{$C$};
\draw[guide](2,-1.2)--(2,1.2);\node[below right]at(2,0){$1$};\node[below]at(4,0){$2$};
\end{tikzpicture}
\caption{平移后的相椭圆。A 是初态，C 对应首次经过新平衡位置，B 为另一转向点。}\label{mp4:fig:l4-h1-s}
\end{figure}

第一次经过新平衡位置时，对应上支的 C，$x=b$。由能量式求得速度大小为 $\omega b$，再由运动方向取正号，
\[
v_C=\omega b=\frac F{\sqrt{km}}.
\]
令 $\xi=x-b$，方程化为 $\ddot\xi+\omega^2\xi=0$，初值是 $\xi(0)=-b$、$\dot\xi(0)=0$。两特征根为 $\pm\mathrm i\omega$，故 $\xi=C\cos\omega t+D\sin\omega t$；初值分别给出 $C=-b$、$\omega D=0$。因此
\[
x=b(1-\cos\omega t),\qquad v=\omega b\sin\omega t.
\]
$x=b$ 首次成立要求 $\cos\omega t=0$ 且 $v>0$，于是 $t_C=\pi/(2\omega)$。这是从左端点 A 运动到中心竖线的四分之一周期。

若只把弹性势能计入机械能，外力确实做了功：到达 C 时 $W_F=Fb=kb^2$，其中 $kb^2/2$ 成为动能，另 $kb^2/2$ 成为弹性势能。把 $-Fx$ 纳入势能后守恒的总能量，与这一功的计算相容。新平衡位置附近的回复力仍为 $-k\xi$，故平移后的相椭圆仍以周期 $2\pi/\omega$ 绕行。
\end{solution}

\par\medskip\noindent\begin{minipage}{\textwidth}
\begin{exercise}[V 形势阱]\label{mp4:exercise:l4-h2}
质量为$m$的质点在势能 $U=f|x|$ 中运动，$f>0$。最大位移为 $A$，图\ref{mp4:fig:l4-h2-q}给出势能折线。求相轨线和周期，说明质点经过 $x=0$ 时哪些量连续、哪些量发生突变。
\end{exercise}

\begin{figure}[H]\centering
\begin{tikzpicture}
\draw[axis](-2.3,0)--(2.3,0)node[right]{$x$};
\draw[axis](0,-0.2)--(0,2.2)node[above]{$U$};
\draw[curve](-1.9,1.9)--(0,0)--(1.9,1.9);\draw[guide](-2,1.4)--(2,1.4)node[right]{$fA$};\node[below]at(-1.4,0){$-A$};\node[below]at(1.4,0){$A$};
\end{tikzpicture}
\caption{V 形势阱。势能在原点连续，左右斜率不同。}\label{mp4:fig:l4-h2-q}
\end{figure}
\end{minipage}\par\medskip
\begin{solution}
在图\ref{mp4:fig:l4-h2-q}的两侧，位置力分别为
\[
F(x)=\begin{cases}+f,&x<0,\\-f,&x>0.\end{cases}
\]
转向点 $x=\pm A$ 处速度为零，故总能量为 $E=fA$。从能量式解出
\[
v=\pm\sqrt{\frac{2f}{m}(A-|x|)},\qquad |x|\le A.
\]
图\ref{mp4:fig:l4-h2-s}的上支向右、下支向左。左右两段分别是 $v^2$ 与 $x$ 成一次关系的抛物线段，在纵轴上拼接；两个转向点处相点分别沿竖直方向离开横轴。

\begin{figure}[H]\centering
\begin{tikzpicture}
\draw[axis](-2.4,0)--(2.4,0)node[right]{$q$};
\draw[axis](0,-1.7)--(0,1.7)node[above]{$w$};
\draw[flow] plot coordinates{(1.80000,-0.00000) (1.79938,-0.02406) (1.79753,-0.04813) (1.79445,-0.07218) (1.79014,-0.09622) (1.78460,-0.12024) (1.77784,-0.14425) (1.76986,-0.16822) (1.76067,-0.19217) (1.75027,-0.21609) (1.73867,-0.23997) (1.72588,-0.26381) (1.71190,-0.28760) (1.69675,-0.31135) (1.68044,-0.33504) (1.66298,-0.35867) (1.64438,-0.38224) (1.62465,-0.40575) (1.60381,-0.42918) (1.58187,-0.45255) (1.55885,-0.47583) (1.53475,-0.49904) (1.50961,-0.52216) (1.48343,-0.54519) (1.45623,-0.56812) (1.42804,-0.59096) (1.39886,-0.61370) (1.36873,-0.63633) (1.33766,-0.65885) (1.30567,-0.68126) (1.27279,-0.70355) (1.23904,-0.72573) (1.20444,-0.74778) (1.16901,-0.76970) (1.13278,-0.79149) (1.09577,-0.81314) (1.05801,-0.83465) (1.01953,-0.85602) (0.98035,-0.87725) (0.94050,-0.89832) (0.90000,-0.91924) (0.85889,-0.94000) (0.81718,-0.96060) (0.77492,-0.98104) (0.73213,-1.00131) (0.68883,-1.02140) (0.64506,-1.04133) (0.60085,-1.06107) (0.55623,-1.08063) (0.51123,-1.10001) (0.46587,-1.11919) (0.42020,-1.13819) (0.37424,-1.15699) (0.32802,-1.17559) (0.28158,-1.19400) (0.23495,-1.21219) (0.18815,-1.23018) (0.14123,-1.24796) (0.09420,-1.26552) (0.04712,-1.28287) (0.00000,-1.30000) (-0.04712,-1.28287) (-0.09420,-1.26552) (-0.14123,-1.24796) (-0.18815,-1.23018) (-0.23495,-1.21219) (-0.28158,-1.19400) (-0.32802,-1.17559) (-0.37424,-1.15699) (-0.42020,-1.13819) (-0.46587,-1.11919) (-0.51123,-1.10001) (-0.55623,-1.08063) (-0.60085,-1.06107) (-0.64506,-1.04133) (-0.68883,-1.02140) (-0.73213,-1.00131) (-0.77492,-0.98104) (-0.81718,-0.96060) (-0.85889,-0.94000) (-0.90000,-0.91924) (-0.94050,-0.89832) (-0.98035,-0.87725) (-1.01953,-0.85602) (-1.05801,-0.83465) (-1.09577,-0.81314) (-1.13278,-0.79149) (-1.16901,-0.76970) (-1.20444,-0.74778) (-1.23904,-0.72573) (-1.27279,-0.70355) (-1.30567,-0.68126) (-1.33766,-0.65885) (-1.36873,-0.63633) (-1.39886,-0.61370) (-1.42804,-0.59096) (-1.45623,-0.56812) (-1.48343,-0.54519) (-1.50961,-0.52216) (-1.53475,-0.49904) (-1.55885,-0.47583) (-1.58187,-0.45255) (-1.60381,-0.42918) (-1.62465,-0.40575) (-1.64438,-0.38224) (-1.66298,-0.35867) (-1.68044,-0.33504) (-1.69675,-0.31135) (-1.71190,-0.28760) (-1.72588,-0.26381) (-1.73867,-0.23997) (-1.75027,-0.21609) (-1.76067,-0.19217) (-1.76986,-0.16822) (-1.77784,-0.14425) (-1.78460,-0.12024) (-1.79014,-0.09622) (-1.79445,-0.07218) (-1.79753,-0.04813) (-1.79938,-0.02406) (-1.80000,-0.00000) (-1.79938,0.02406) (-1.79753,0.04813) (-1.79445,0.07218) (-1.79014,0.09622) (-1.78460,0.12024) (-1.77784,0.14425) (-1.76986,0.16822) (-1.76067,0.19217) (-1.75027,0.21609) (-1.73867,0.23997) (-1.72588,0.26381) (-1.71190,0.28760) (-1.69675,0.31135) (-1.68044,0.33504) (-1.66298,0.35867) (-1.64438,0.38224) (-1.62465,0.40575) (-1.60381,0.42918) (-1.58187,0.45255) (-1.55885,0.47583) (-1.53475,0.49904) (-1.50961,0.52216) (-1.48343,0.54519) (-1.45623,0.56812) (-1.42804,0.59096) (-1.39886,0.61370) (-1.36873,0.63633) (-1.33766,0.65885) (-1.30567,0.68126) (-1.27279,0.70355) (-1.23904,0.72573) (-1.20444,0.74778) (-1.16901,0.76970) (-1.13278,0.79149) (-1.09577,0.81314) (-1.05801,0.83465) (-1.01953,0.85602) (-0.98035,0.87725) (-0.94050,0.89832) (-0.90000,0.91924) (-0.85889,0.94000) (-0.81718,0.96060) (-0.77492,0.98104) (-0.73213,1.00131) (-0.68883,1.02140) (-0.64506,1.04133) (-0.60085,1.06107) (-0.55623,1.08063) (-0.51123,1.10001) (-0.46587,1.11919) (-0.42020,1.13819) (-0.37424,1.15699) (-0.32802,1.17559) (-0.28158,1.19400) (-0.23495,1.21219) (-0.18815,1.23018) (-0.14123,1.24796) (-0.09420,1.26552) (-0.04712,1.28287) (-0.00000,1.30000) (0.04712,1.28287) (0.09420,1.26552) (0.14123,1.24796) (0.18815,1.23018) (0.23495,1.21219) (0.28158,1.19400) (0.32802,1.17559) (0.37424,1.15699) (0.42020,1.13819) (0.46587,1.11919) (0.51123,1.10001) (0.55623,1.08063) (0.60085,1.06107) (0.64506,1.04133) (0.68883,1.02140) (0.73213,1.00131) (0.77492,0.98104) (0.81718,0.96060) (0.85889,0.94000) (0.90000,0.91924) (0.94050,0.89832) (0.98035,0.87725) (1.01953,0.85602) (1.05801,0.83465) (1.09577,0.81314) (1.13278,0.79149) (1.16901,0.76970) (1.20444,0.74778) (1.23904,0.72573) (1.27279,0.70355) (1.30567,0.68126) (1.33766,0.65885) (1.36873,0.63633) (1.39886,0.61370) (1.42804,0.59096) (1.45623,0.56812) (1.48343,0.54519) (1.50961,0.52216) (1.53475,0.49904) (1.55885,0.47583) (1.58187,0.45255) (1.60381,0.42918) (1.62465,0.40575) (1.64438,0.38224) (1.66298,0.35867) (1.68044,0.33504) (1.69675,0.31135) (1.71190,0.28760) (1.72588,0.26381) (1.73867,0.23997) (1.75027,0.21609) (1.76067,0.19217) (1.76986,0.16822) (1.77784,0.14425) (1.78460,0.12024) (1.79014,0.09622) (1.79445,0.07218) (1.79753,0.04813) (1.79938,0.02406) (1.80000,0.00000)};
\node[below right]at(1.8,0){$1$};\node[above right]at(0,1.3){$1$};
\end{tikzpicture}
\caption{V 形势阱的归一化相轨线，$q=x/A$，$w=v/v_{\max}$。两段轨线在 $q=0$ 处形成折角。}\label{mp4:fig:l4-h2-s}
\end{figure}

从右端 $x=A$ 静止释放，在到达原点前加速度恒为 $-f/m$。以释放时刻为零点，两次积分为
\[
v=\int_0^t\left(-\frac fm\right)\dd u
=-\frac fm[u]_0^t=-\frac fm t,
\qquad
x=A-\frac fm\int_0^t u\dd u
=A-\frac fm\left[\frac{u^2}{2}\right]_0^t.
\]
即
\[
x=A-\frac f{2m}t^2,\qquad v=-\frac fm t.
\]
到达原点的时间与速度为
\[
t_1=\sqrt{\frac{2mA}{f}},\qquad
v(t_1)=-\sqrt{\frac{2fA}{m}}.
\]
进入负半轴后加速度变为 $+f/m$，质点仍向左运动，但速率减小；经过同样的时间 $t_1$ 后在 $x=-A$ 转向。随后两段运动由对称性得到。因此
\[
T=4t_1=4\sqrt{\frac{2mA}{f}},\qquad
v_{\max}=\sqrt{\frac{2fA}{m}}.
\]
这里 $T\propto\sqrt A$：从端点到原点的加速度大小恒定，所需时间由匀加速位移公式确定。

以图中下支从右向左过原点为例，加速度由 $-f/m$ 跳到 $+f/m$。但对任意包含过零时刻的短时间段，冲量的大小不超过 $f\Delta t$，在 $\Delta t\to0$ 时趋于零，故速度没有跃变。位置、速度、势能和动能都连续，力及加速度存在不同的左右极限。

过零时 $v\ne0$，相点方向为 $(\dot x,\dot v)=(v,F/m)$。水平分量连续，竖直分量随力反号，故斜率 $\dd v/\dd x=F/(mv)$ 跳变，在图\ref{mp4:fig:l4-h2-s}中形成折角。运动仍按相同的位置和速度衔接。转向点处则有 $v=0$，相点方向竖直，须沿图中的竖直箭头转入另一支，不能沿横轴折返。
\end{solution}

\par\medskip\noindent\begin{minipage}{\textwidth}
\begin{exercise}[由相椭圆求时间比例]\label{mp4:exercise:l4-h3}
某简谐振子的相轨线横轴截距为 $\pm a$、纵轴截距为 $\pm v_*$，见图\ref{mp4:fig:l4-h3-q}。求角频率；在一个周期中，$|x|\le a/2$ 的时间占比是多少？速率满足 $|v|\ge\sqrt3\,v_*/2$ 的时间占比又是多少？
\end{exercise}

\begin{figure}[H]\centering
\begin{tikzpicture}
\draw[axis](-2.5,0)--(2.5,0)node[right]{$x$};
\draw[axis](0,-1.8)--(0,1.8)node[above]{$v$};
\draw[flow] plot coordinates{(2.00000,-0.00000) (1.99846,-0.04711) (1.99383,-0.09415) (1.98614,-0.14104) (1.97538,-0.18772) (1.96157,-0.23411) (1.94474,-0.28013) (1.92491,-0.32573) (1.90211,-0.37082) (1.87638,-0.41534) (1.84776,-0.45922) (1.81629,-0.50239) (1.78201,-0.54479) (1.74499,-0.58635) (1.70528,-0.62700) (1.66294,-0.66668) (1.61803,-0.70534) (1.57063,-0.74291) (1.52081,-0.77934) (1.46865,-0.81456) (1.41421,-0.84853) (1.35760,-0.88119) (1.29890,-0.91249) (1.23819,-0.94238) (1.17557,-0.97082) (1.11114,-0.99776) (1.04500,-1.02317) (0.97724,-1.04700) (0.90798,-1.06921) (0.83732,-1.08977) (0.76537,-1.10866) (0.69223,-1.12583) (0.61803,-1.14127) (0.54288,-1.15495) (0.46689,-1.16684) (0.39018,-1.17694) (0.31287,-1.18523) (0.23507,-1.19168) (0.15692,-1.19630) (0.07852,-1.19907) (0.00000,-1.20000) (-0.07852,-1.19907) (-0.15692,-1.19630) (-0.23507,-1.19168) (-0.31287,-1.18523) (-0.39018,-1.17694) (-0.46689,-1.16684) (-0.54288,-1.15495) (-0.61803,-1.14127) (-0.69223,-1.12583) (-0.76537,-1.10866) (-0.83732,-1.08977) (-0.90798,-1.06921) (-0.97724,-1.04700) (-1.04500,-1.02317) (-1.11114,-0.99776) (-1.17557,-0.97082) (-1.23819,-0.94238) (-1.29890,-0.91249) (-1.35760,-0.88119) (-1.41421,-0.84853) (-1.46865,-0.81456) (-1.52081,-0.77934) (-1.57063,-0.74291) (-1.61803,-0.70534) (-1.66294,-0.66668) (-1.70528,-0.62700) (-1.74499,-0.58635) (-1.78201,-0.54479) (-1.81629,-0.50239) (-1.84776,-0.45922) (-1.87638,-0.41534) (-1.90211,-0.37082) (-1.92491,-0.32573) (-1.94474,-0.28013) (-1.96157,-0.23411) (-1.97538,-0.18772) (-1.98614,-0.14104) (-1.99383,-0.09415) (-1.99846,-0.04711) (-2.00000,-0.00000) (-1.99846,0.04711) (-1.99383,0.09415) (-1.98614,0.14104) (-1.97538,0.18772) (-1.96157,0.23411) (-1.94474,0.28013) (-1.92491,0.32573) (-1.90211,0.37082) (-1.87638,0.41534) (-1.84776,0.45922) (-1.81629,0.50239) (-1.78201,0.54479) (-1.74499,0.58635) (-1.70528,0.62700) (-1.66294,0.66668) (-1.61803,0.70534) (-1.57063,0.74291) (-1.52081,0.77934) (-1.46865,0.81456) (-1.41421,0.84853) (-1.35760,0.88119) (-1.29890,0.91249) (-1.23819,0.94238) (-1.17557,0.97082) (-1.11114,0.99776) (-1.04500,1.02317) (-0.97724,1.04700) (-0.90798,1.06921) (-0.83732,1.08977) (-0.76537,1.10866) (-0.69223,1.12583) (-0.61803,1.14127) (-0.54288,1.15495) (-0.46689,1.16684) (-0.39018,1.17694) (-0.31287,1.18523) (-0.23507,1.19168) (-0.15692,1.19630) (-0.07852,1.19907) (-0.00000,1.20000) (0.07852,1.19907) (0.15692,1.19630) (0.23507,1.19168) (0.31287,1.18523) (0.39018,1.17694) (0.46689,1.16684) (0.54288,1.15495) (0.61803,1.14127) (0.69223,1.12583) (0.76537,1.10866) (0.83732,1.08977) (0.90798,1.06921) (0.97724,1.04700) (1.04500,1.02317) (1.11114,0.99776) (1.17557,0.97082) (1.23819,0.94238) (1.29890,0.91249) (1.35760,0.88119) (1.41421,0.84853) (1.46865,0.81456) (1.52081,0.77934) (1.57063,0.74291) (1.61803,0.70534) (1.66294,0.66668) (1.70528,0.62700) (1.74499,0.58635) (1.78201,0.54479) (1.81629,0.50239) (1.84776,0.45922) (1.87638,0.41534) (1.90211,0.37082) (1.92491,0.32573) (1.94474,0.28013) (1.96157,0.23411) (1.97538,0.18772) (1.98614,0.14104) (1.99383,0.09415) (1.99846,0.04711) (2.00000,0.00000)};
\node[below right]at(2,0){$a$};\node[below left]at(-2,0){$-a$};\node[above right]at(0,1.2){$v_*$};
\end{tikzpicture}
\caption{已知截距的相椭圆。图的横、纵轴具有不同量纲。}\label{mp4:fig:l4-h3-q}
\end{figure}
\end{minipage}\par\medskip
\begin{solution}
图\ref{mp4:fig:l4-h3-q}的横轴截距给出振幅 $a$，纵轴截距给出最大速率 $v_*$。相椭圆为
\[
\frac{x^2}{a^2}+\frac{v^2}{v_*^2}=1.
\]
简谐振子的能量关系是 $v^2=\omega^2(a^2-x^2)$，比较两式可得
\[
\omega=\frac{v_*}{a},\qquad T=\frac{2\pi a}{v_*}.
\]
若取 $x=a\sin\varphi$、$v=v_*\cos\varphi$，在 $v\ne0$ 处由 $v=\dot x$ 得 $\dot\varphi=v_*/a=\omega$；这一关系连续延伸到转向时刻。相位匀速增加，因此一段运动占据的时间比例等于其相位总长度除以 $2\pi$。

\begin{figure}[H]\centering
\begin{tikzpicture}
\draw[axis](-2.5,0)--(2.5,0)node[right]{$x/a$};
\draw[axis](0,-2.2)--(0,2.2)node[above]{$v/v_*$};
\draw[curve] plot coordinates{(1.80000,-0.00000) (1.79861,-0.07067) (1.79445,-0.14123) (1.78752,-0.21157) (1.77784,-0.28158) (1.76541,-0.35116) (1.75027,-0.42020) (1.73242,-0.48859) (1.71190,-0.55623) (1.68874,-0.62301) (1.66298,-0.68883) (1.63466,-0.75359) (1.60381,-0.81718) (1.57049,-0.87952) (1.53475,-0.94050) (1.49665,-1.00003) (1.45623,-1.05801) (1.41357,-1.11437) (1.36873,-1.16901) (1.32178,-1.22184) (1.27279,-1.27279) (1.22184,-1.32178) (1.16901,-1.36873) (1.11437,-1.41357) (1.05801,-1.45623) (1.00003,-1.49665) (0.94050,-1.53475) (0.87952,-1.57049) (0.81718,-1.60381) (0.75359,-1.63466) (0.68883,-1.66298) (0.62301,-1.68874) (0.55623,-1.71190) (0.48859,-1.73242) (0.42020,-1.75027) (0.35116,-1.76541) (0.28158,-1.77784) (0.21157,-1.78752) (0.14123,-1.79445) (0.07067,-1.79861) (0.00000,-1.80000) (-0.07067,-1.79861) (-0.14123,-1.79445) (-0.21157,-1.78752) (-0.28158,-1.77784) (-0.35116,-1.76541) (-0.42020,-1.75027) (-0.48859,-1.73242) (-0.55623,-1.71190) (-0.62301,-1.68874) (-0.68883,-1.66298) (-0.75359,-1.63466) (-0.81718,-1.60381) (-0.87952,-1.57049) (-0.94050,-1.53475) (-1.00003,-1.49665) (-1.05801,-1.45623) (-1.11437,-1.41357) (-1.16901,-1.36873) (-1.22184,-1.32178) (-1.27279,-1.27279) (-1.32178,-1.22184) (-1.36873,-1.16901) (-1.41357,-1.11437) (-1.45623,-1.05801) (-1.49665,-1.00003) (-1.53475,-0.94050) (-1.57049,-0.87952) (-1.60381,-0.81718) (-1.63466,-0.75359) (-1.66298,-0.68883) (-1.68874,-0.62301) (-1.71190,-0.55623) (-1.73242,-0.48859) (-1.75027,-0.42020) (-1.76541,-0.35116) (-1.77784,-0.28158) (-1.78752,-0.21157) (-1.79445,-0.14123) (-1.79861,-0.07067) (-1.80000,-0.00000) (-1.79861,0.07067) (-1.79445,0.14123) (-1.78752,0.21157) (-1.77784,0.28158) (-1.76541,0.35116) (-1.75027,0.42020) (-1.73242,0.48859) (-1.71190,0.55623) (-1.68874,0.62301) (-1.66298,0.68883) (-1.63466,0.75359) (-1.60381,0.81718) (-1.57049,0.87952) (-1.53475,0.94050) (-1.49665,1.00003) (-1.45623,1.05801) (-1.41357,1.11437) (-1.36873,1.16901) (-1.32178,1.22184) (-1.27279,1.27279) (-1.22184,1.32178) (-1.16901,1.36873) (-1.11437,1.41357) (-1.05801,1.45623) (-1.00003,1.49665) (-0.94050,1.53475) (-0.87952,1.57049) (-0.81718,1.60381) (-0.75359,1.63466) (-0.68883,1.66298) (-0.62301,1.68874) (-0.55623,1.71190) (-0.48859,1.73242) (-0.42020,1.75027) (-0.35116,1.76541) (-0.28158,1.77784) (-0.21157,1.78752) (-0.14123,1.79445) (-0.07067,1.79861) (-0.00000,1.80000) (0.07067,1.79861) (0.14123,1.79445) (0.21157,1.78752) (0.28158,1.77784) (0.35116,1.76541) (0.42020,1.75027) (0.48859,1.73242) (0.55623,1.71190) (0.62301,1.68874) (0.68883,1.66298) (0.75359,1.63466) (0.81718,1.60381) (0.87952,1.57049) (0.94050,1.53475) (1.00003,1.49665) (1.05801,1.45623) (1.11437,1.41357) (1.16901,1.36873) (1.22184,1.32178) (1.27279,1.27279) (1.32178,1.22184) (1.36873,1.16901) (1.41357,1.11437) (1.45623,1.05801) (1.49665,1.00003) (1.53475,0.94050) (1.57049,0.87952) (1.60381,0.81718) (1.63466,0.75359) (1.66298,0.68883) (1.68874,0.62301) (1.71190,0.55623) (1.73242,0.48859) (1.75027,0.42020) (1.76541,0.35116) (1.77784,0.28158) (1.78752,0.21157) (1.79445,0.14123) (1.79861,0.07067) (1.80000,0.00000)};
\draw[flow,line width=1.5pt] plot coordinates{(-0.90000,1.55885) (-0.88363,1.56818) (-0.86716,1.57735) (-0.85059,1.58635) (-0.83393,1.59517) (-0.81718,1.60381) (-0.80034,1.61228) (-0.78342,1.62057) (-0.76640,1.62869) (-0.74931,1.63662) (-0.73213,1.64438) (-0.71487,1.65196) (-0.69753,1.65935) (-0.68011,1.66657) (-0.66262,1.67360) (-0.64506,1.68044) (-0.62743,1.68711) (-0.60973,1.69359) (-0.59196,1.69988) (-0.57413,1.70598) (-0.55623,1.71190) (-0.53827,1.71763) (-0.52026,1.72318) (-0.50218,1.72853) (-0.48406,1.73369) (-0.46587,1.73867) (-0.44764,1.74345) (-0.42936,1.74804) (-0.41103,1.75244) (-0.39266,1.75665) (-0.37424,1.76067) (-0.35578,1.76449) (-0.33729,1.76812) (-0.31875,1.77155) (-0.30018,1.77479) (-0.28158,1.77784) (-0.26295,1.78069) (-0.24429,1.78335) (-0.22560,1.78581) (-0.20689,1.78807) (-0.18815,1.79014) (-0.16939,1.79201) (-0.15062,1.79369) (-0.13183,1.79517) (-0.11302,1.79645) (-0.09420,1.79753) (-0.07538,1.79842) (-0.05654,1.79911) (-0.03770,1.79961) (-0.01885,1.79990) (0.00000,1.80000) (0.01885,1.79990) (0.03770,1.79961) (0.05654,1.79911) (0.07538,1.79842) (0.09420,1.79753) (0.11302,1.79645) (0.13183,1.79517) (0.15062,1.79369) (0.16939,1.79201) (0.18815,1.79014) (0.20689,1.78807) (0.22560,1.78581) (0.24429,1.78335) (0.26295,1.78069) (0.28158,1.77784) (0.30018,1.77479) (0.31875,1.77155) (0.33729,1.76812) (0.35578,1.76449) (0.37424,1.76067) (0.39266,1.75665) (0.41103,1.75244) (0.42936,1.74804) (0.44764,1.74345) (0.46587,1.73867) (0.48406,1.73369) (0.50218,1.72853) (0.52026,1.72318) (0.53827,1.71763) (0.55623,1.71190) (0.57413,1.70598) (0.59196,1.69988) (0.60973,1.69359) (0.62743,1.68711) (0.64506,1.68044) (0.66262,1.67360) (0.68011,1.66657) (0.69753,1.65935) (0.71487,1.65196) (0.73213,1.64438) (0.74931,1.63662) (0.76640,1.62869) (0.78342,1.62057) (0.80034,1.61228) (0.81718,1.60381) (0.83393,1.59517) (0.85059,1.58635) (0.86716,1.57735) (0.88363,1.56818) (0.90000,1.55885)};
\draw[flow,line width=1.5pt] plot coordinates{(0.90000,-1.55885) (0.88363,-1.56818) (0.86716,-1.57735) (0.85059,-1.58635) (0.83393,-1.59517) (0.81718,-1.60381) (0.80034,-1.61228) (0.78342,-1.62057) (0.76640,-1.62869) (0.74931,-1.63662) (0.73213,-1.64438) (0.71487,-1.65196) (0.69753,-1.65935) (0.68011,-1.66657) (0.66262,-1.67360) (0.64506,-1.68044) (0.62743,-1.68711) (0.60973,-1.69359) (0.59196,-1.69988) (0.57413,-1.70598) (0.55623,-1.71190) (0.53827,-1.71763) (0.52026,-1.72318) (0.50218,-1.72853) (0.48406,-1.73369) (0.46587,-1.73867) (0.44764,-1.74345) (0.42936,-1.74804) (0.41103,-1.75244) (0.39266,-1.75665) (0.37424,-1.76067) (0.35578,-1.76449) (0.33729,-1.76812) (0.31875,-1.77155) (0.30018,-1.77479) (0.28158,-1.77784) (0.26295,-1.78069) (0.24429,-1.78335) (0.22560,-1.78581) (0.20689,-1.78807) (0.18815,-1.79014) (0.16939,-1.79201) (0.15062,-1.79369) (0.13183,-1.79517) (0.11302,-1.79645) (0.09420,-1.79753) (0.07538,-1.79842) (0.05654,-1.79911) (0.03770,-1.79961) (0.01885,-1.79990) (0.00000,-1.80000) (-0.01885,-1.79990) (-0.03770,-1.79961) (-0.05654,-1.79911) (-0.07538,-1.79842) (-0.09420,-1.79753) (-0.11302,-1.79645) (-0.13183,-1.79517) (-0.15062,-1.79369) (-0.16939,-1.79201) (-0.18815,-1.79014) (-0.20689,-1.78807) (-0.22560,-1.78581) (-0.24429,-1.78335) (-0.26295,-1.78069) (-0.28158,-1.77784) (-0.30018,-1.77479) (-0.31875,-1.77155) (-0.33729,-1.76812) (-0.35578,-1.76449) (-0.37424,-1.76067) (-0.39266,-1.75665) (-0.41103,-1.75244) (-0.42936,-1.74804) (-0.44764,-1.74345) (-0.46587,-1.73867) (-0.48406,-1.73369) (-0.50218,-1.72853) (-0.52026,-1.72318) (-0.53827,-1.71763) (-0.55623,-1.71190) (-0.57413,-1.70598) (-0.59196,-1.69988) (-0.60973,-1.69359) (-0.62743,-1.68711) (-0.64506,-1.68044) (-0.66262,-1.67360) (-0.68011,-1.66657) (-0.69753,-1.65935) (-0.71487,-1.65196) (-0.73213,-1.64438) (-0.74931,-1.63662) (-0.76640,-1.62869) (-0.78342,-1.62057) (-0.80034,-1.61228) (-0.81718,-1.60381) (-0.83393,-1.59517) (-0.85059,-1.58635) (-0.86716,-1.57735) (-0.88363,-1.56818) (-0.90000,-1.55885)};
\draw[guide](-.9,-1.9)--(-.9,1.9);\draw[guide](.9,-1.9)--(.9,1.9);\node[below right]at(.9,0){$1/2$};\node[below left]at(-.9,0){$-1/2$};
\end{tikzpicture}
\caption{按截距归一化后得到的圆。加粗弧对应两个等价的不等式，按相位计共占一周的三分之一。}\label{mp4:fig:l4-h3-s}
\end{figure}

图\ref{mp4:fig:l4-h3-s}的两条竖直虚线为 $x=\pm a/2$，所选区域包含上、下各一段相轨线。条件 $|x|\le a/2$ 等价于 $|\sin\varphi|\le1/2$。在一个周期内可取两段相位区间
\[
-\frac\pi6\le\varphi\le\frac\pi6,
\qquad
\frac{5\pi}6\le\varphi\le\frac{7\pi}6.
\]
它们分别对应向右和向左通过中央区域，每段相位长 $\pi/3$。总时间及占比为
\[
\Delta t=\frac{2\pi}{3\omega},\qquad
\frac{\Delta t}{T}=\frac13.
\]

第二个条件可直接由同一相椭圆换算：
\[
|v|\ge\frac{\sqrt3}{2}v_*
\ \Longleftrightarrow\ 
\frac{v^2}{v_*^2}\ge\frac34
\ \Longleftrightarrow\ 
\frac{x^2}{a^2}\le\frac14.
\]
所以它选中图中的同两段轨线，时间占比仍为 $1/3$。端点是否计入不影响时间比例。

这里的 $\varphi$ 是归一化圆 $(x/a,v/v_*)$ 上的参数角。原图两轴的比例尺改变时，椭圆弧长和几何圆心角随之改变，时间比例仍由上述相位区间确定。也可对中央区域的四个对称半段直接积分：
\[
\Delta t=4\int_0^{a/2}\frac{\dd x}{\omega\sqrt{a^2-x^2}}
=\frac4\omega\left[\arcsin\frac xa\right]_0^{a/2}
=\frac{2\pi}{3\omega}.
\]
\end{solution}

\par\medskip\noindent\begin{minipage}{\textwidth}
\begin{exercise}[有解离阈值的势阱（选做）]\label{mp4:exercise:l4-h4}
质量为$m$的质点在 $U(x)=D(1-\ee^{-x/a})^2$ 中运动，$D,a>0$，势能见图\ref{mp4:fig:l4-h4-q}。对 $0<E<D$，求转向位置和周期；再讨论 $E=D$ 时向右运动能否在有限时间内到达无穷远。
\end{exercise}

\begin{figure}[H]\centering
\begin{tikzpicture}
\draw[axis](-1.5,0)--(3.7,0)node[right]{$x/a$};
\draw[axis](0,-0.2)--(0,2.5)node[above]{$U/D$};
\draw[curve,domain=-.9:3.3,samples=130]plot(\x,{(1-exp(-\x))^2});\draw[guide](-1.1,1)--(3.4,1)node[right]{$1$};\draw[guide](-.9,.5)--(2.7,.5)node[right]{$E/D$};
\end{tikzpicture}
\caption{一侧趋于有限高度的势阱。右侧水平虚线为势能的渐近值 D。}\label{mp4:fig:l4-h4-q}
\end{figure}
\end{minipage}\par\medskip
\begin{solution}
记 $h=E/D$，题设限定 $0<h<1$。在转向点，$\ee^{-x/a}>0$，且
\[
1-\ee^{-x/a}=\pm\sqrt h.
\]
取负号得到左转向点，取正号得到右转向点：
\[
x_-=-a\ln(1+\sqrt h)<0,\qquad
x_+=-a\ln(1-\sqrt h)>0.
\]
两者之间 $U(x)<E$，形成一个有限的可达区间。图\ref{mp4:fig:l4-h4-s}中每条闭合轨线的上支向右、下支向左；左右转向位置不对称。

周期为整个允许区间往返一次的时间，因而
\[
T=2\sqrt{\frac m{2D}}\int_{x_-}^{x_+}
\frac{\dd x}{\sqrt{h-(1-\ee^{-x/a})^2}}.
\]
作 $z=\ee^{-x/a}$，则 $\dd x=-a\,\dd z/z$。从左到右运动时 $z$ 从 $1+\sqrt h$ 减到 $1-\sqrt h$，交换积分上下限后，
\[
T=2a\sqrt{\frac m{2D}}\int_{1-\sqrt h}^{1+\sqrt h}
\frac{\dd z}{z\sqrt{h-(1-z)^2}}.
\]

\begin{figure}[H]\centering
\begin{tikzpicture}
\draw[axis](-1.3,0)--(4.2,0)node[right]{$q$};
\draw[axis](0,-1.6)--(0,1.7)node[above]{$w$};
\draw[flow] plot coordinates{(-0.53480,0.00000) (-0.53460,0.02221) (-0.53398,0.04440) (-0.53296,0.06654) (-0.53153,0.08862) (-0.52969,0.11062) (-0.52744,0.13250) (-0.52477,0.15425) (-0.52170,0.17585) (-0.51822,0.19728) (-0.51432,0.21851) (-0.51001,0.23952) (-0.50528,0.26030) (-0.50014,0.28083) (-0.49458,0.30107) (-0.48860,0.32102) (-0.48221,0.34065) (-0.47539,0.35995) (-0.46815,0.37889) (-0.46048,0.39745) (-0.45239,0.41563) (-0.44387,0.43339) (-0.43492,0.45073) (-0.42553,0.46762) (-0.41572,0.48405) (-0.40547,0.50000) (-0.39477,0.51546) (-0.38364,0.53041) (-0.37206,0.54484) (-0.36004,0.55872) (-0.34757,0.57206) (-0.33465,0.58483) (-0.32128,0.59703) (-0.30745,0.60864) (-0.29316,0.61964) (-0.27840,0.63004) (-0.26319,0.63981) (-0.24751,0.64895) (-0.23135,0.65745) (-0.21473,0.66530) (-0.19763,0.67250) (-0.18005,0.67903) (-0.16199,0.68489) (-0.14345,0.69008) (-0.12443,0.69458) (-0.10491,0.69840) (-0.08491,0.70153) (-0.06442,0.70397) (-0.04344,0.70571) (-0.02197,0.70676) (-0.00000,0.70711) (0.02246,0.70676) (0.04542,0.70571) (0.06886,0.70397) (0.09280,0.70153) (0.11723,0.69840) (0.14214,0.69458) (0.16753,0.69008) (0.19340,0.68489) (0.21974,0.67903) (0.24655,0.67250) (0.27381,0.66530) (0.30152,0.65745) (0.32965,0.64895) (0.35821,0.63981) (0.38716,0.63004) (0.41650,0.61964) (0.44620,0.60864) (0.47624,0.59703) (0.50659,0.58483) (0.53722,0.57206) (0.56809,0.55872) (0.59916,0.54484) (0.63039,0.53041) (0.66174,0.51546) (0.69315,0.50000) (0.72455,0.48405) (0.75589,0.46762) (0.78710,0.45073) (0.81808,0.43339) (0.84878,0.41563) (0.87908,0.39745) (0.90889,0.37889) (0.93812,0.35995) (0.96664,0.34065) (0.99435,0.32102) (1.02112,0.30107) (1.04683,0.28083) (1.07134,0.26030) (1.09453,0.23952) (1.11626,0.21851) (1.13641,0.19728) (1.15484,0.17585) (1.17143,0.15425) (1.18607,0.13250) (1.19866,0.11062) (1.20909,0.08862) (1.21729,0.06654) (1.22319,0.04440) (1.22676,0.02221) (1.22795,0.00000) (1.22676,-0.02221) (1.22319,-0.04440) (1.21729,-0.06654) (1.20909,-0.08862) (1.19866,-0.11062) (1.18607,-0.13250) (1.17143,-0.15425) (1.15484,-0.17585) (1.13641,-0.19728) (1.11626,-0.21851) (1.09453,-0.23952) (1.07134,-0.26030) (1.04683,-0.28083) (1.02112,-0.30107) (0.99435,-0.32102) (0.96664,-0.34065) (0.93812,-0.35995) (0.90889,-0.37889) (0.87908,-0.39745) (0.84878,-0.41563) (0.81808,-0.43339) (0.78710,-0.45073) (0.75589,-0.46762) (0.72455,-0.48405) (0.69315,-0.50000) (0.66174,-0.51546) (0.63039,-0.53041) (0.59916,-0.54484) (0.56809,-0.55872) (0.53722,-0.57206) (0.50659,-0.58483) (0.47624,-0.59703) (0.44620,-0.60864) (0.41650,-0.61964) (0.38716,-0.63004) (0.35821,-0.63981) (0.32965,-0.64895) (0.30152,-0.65745) (0.27381,-0.66530) (0.24655,-0.67250) (0.21974,-0.67903) (0.19340,-0.68489) (0.16753,-0.69008) (0.14214,-0.69458) (0.11723,-0.69840) (0.09280,-0.70153) (0.06886,-0.70397) (0.04542,-0.70571) (0.02246,-0.70676) (0.00000,-0.70711) (-0.02197,-0.70676) (-0.04344,-0.70571) (-0.06442,-0.70397) (-0.08491,-0.70153) (-0.10491,-0.69840) (-0.12443,-0.69458) (-0.14345,-0.69008) (-0.16199,-0.68489) (-0.18005,-0.67903) (-0.19763,-0.67250) (-0.21473,-0.66530) (-0.23135,-0.65745) (-0.24751,-0.64895) (-0.26319,-0.63981) (-0.27840,-0.63004) (-0.29316,-0.61964) (-0.30745,-0.60864) (-0.32128,-0.59703) (-0.33465,-0.58483) (-0.34757,-0.57206) (-0.36004,-0.55872) (-0.37206,-0.54484) (-0.38364,-0.53041) (-0.39477,-0.51546) (-0.40547,-0.50000) (-0.41572,-0.48405) (-0.42553,-0.46762) (-0.43492,-0.45073) (-0.44387,-0.43339) (-0.45239,-0.41563) (-0.46048,-0.39745) (-0.46815,-0.37889) (-0.47539,-0.35995) (-0.48221,-0.34065) (-0.48860,-0.32102) (-0.49458,-0.30107) (-0.50014,-0.28083) (-0.50528,-0.26030) (-0.51001,-0.23952) (-0.51432,-0.21851) (-0.51822,-0.19728) (-0.52170,-0.17585) (-0.52477,-0.15425) (-0.52744,-0.13250) (-0.52969,-0.11062) (-0.53153,-0.08862) (-0.53296,-0.06654) (-0.53398,-0.04440) (-0.53460,-0.02221) (-0.53480,-0.00000)};
\draw[flow,dashed] plot coordinates{(-0.69315,0.00000) (-0.65122,0.39688) (-0.60928,0.54395) (-0.56735,0.64570) (-0.52542,0.72269) (-0.48349,0.78324) (-0.44156,0.83176) (-0.39963,0.87101) (-0.35770,0.90281) (-0.31576,0.92851) (-0.27383,0.94909) (-0.23190,0.96534) (-0.18997,0.97787) (-0.14804,0.98719) (-0.10611,0.99371) (-0.06418,0.99780) (-0.02224,0.99975) (0.01969,0.99981) (0.06162,0.99821) (0.10355,0.99515) (0.14548,0.99079) (0.18741,0.98529) (0.22935,0.97877) (0.27128,0.97136) (0.31321,0.96317) (0.35514,0.95428) (0.39707,0.94478) (0.43900,0.93474) (0.48093,0.92425) (0.52287,0.91335) (0.56480,0.90210) (0.60673,0.89056) (0.64866,0.87877) (0.69059,0.86676) (0.73252,0.85459) (0.77445,0.84228) (0.81639,0.82986) (0.85832,0.81736) (0.90025,0.80481) (0.94218,0.79223) (0.98411,0.77964) (1.02604,0.76706) (1.06797,0.75450) (1.10991,0.74199) (1.15184,0.72953) (1.19377,0.71715) (1.23570,0.70484) (1.27763,0.69262) (1.31956,0.68050) (1.36149,0.66849) (1.40343,0.65659) (1.44536,0.64482) (1.48729,0.63317) (1.52922,0.62165) (1.57115,0.61027) (1.61308,0.59903) (1.65502,0.58793) (1.69695,0.57698) (1.73888,0.56618) (1.78081,0.55553) (1.82274,0.54503) (1.86467,0.53468) (1.90660,0.52449) (1.94854,0.51445) (1.99047,0.50457) (2.03240,0.49485) (2.07433,0.48528) (2.11626,0.47586) (2.15819,0.46660) (2.20012,0.45750) (2.24206,0.44855) (2.28399,0.43974) (2.32592,0.43110) (2.36785,0.42260) (2.40978,0.41425) (2.45171,0.40604) (2.49364,0.39799) (2.53558,0.39007) (2.57751,0.38230) (2.61944,0.37467) (2.66137,0.36719) (2.70330,0.35983) (2.74523,0.35262) (2.78716,0.34554) (2.82910,0.33859) (2.87103,0.33177) (2.91296,0.32508) (2.95489,0.31852) (2.99682,0.31208) (3.03875,0.30577) (3.08069,0.29958) (3.12262,0.29350) (3.16455,0.28755) (3.20648,0.28170) (3.24841,0.27598) (3.29034,0.27036) (3.33227,0.26486) (3.37421,0.25946) (3.41614,0.25416) (3.45807,0.24898) (3.50000,0.24389)};
\draw[flow,dashed] plot coordinates{(3.50000,-0.24389) (3.45807,-0.24898) (3.41614,-0.25416) (3.37421,-0.25946) (3.33227,-0.26486) (3.29034,-0.27036) (3.24841,-0.27598) (3.20648,-0.28170) (3.16455,-0.28755) (3.12262,-0.29350) (3.08069,-0.29958) (3.03875,-0.30577) (2.99682,-0.31208) (2.95489,-0.31852) (2.91296,-0.32508) (2.87103,-0.33177) (2.82910,-0.33859) (2.78716,-0.34554) (2.74523,-0.35262) (2.70330,-0.35983) (2.66137,-0.36719) (2.61944,-0.37467) (2.57751,-0.38230) (2.53558,-0.39007) (2.49364,-0.39799) (2.45171,-0.40604) (2.40978,-0.41425) (2.36785,-0.42260) (2.32592,-0.43110) (2.28399,-0.43974) (2.24206,-0.44855) (2.20012,-0.45750) (2.15819,-0.46660) (2.11626,-0.47586) (2.07433,-0.48528) (2.03240,-0.49485) (1.99047,-0.50457) (1.94854,-0.51445) (1.90660,-0.52449) (1.86467,-0.53468) (1.82274,-0.54503) (1.78081,-0.55553) (1.73888,-0.56618) (1.69695,-0.57698) (1.65502,-0.58793) (1.61308,-0.59903) (1.57115,-0.61027) (1.52922,-0.62165) (1.48729,-0.63317) (1.44536,-0.64482) (1.40343,-0.65659) (1.36149,-0.66849) (1.31956,-0.68050) (1.27763,-0.69262) (1.23570,-0.70484) (1.19377,-0.71715) (1.15184,-0.72953) (1.10991,-0.74199) (1.06797,-0.75450) (1.02604,-0.76706) (0.98411,-0.77964) (0.94218,-0.79223) (0.90025,-0.80481) (0.85832,-0.81736) (0.81639,-0.82986) (0.77445,-0.84228) (0.73252,-0.85459) (0.69059,-0.86676) (0.64866,-0.87877) (0.60673,-0.89056) (0.56480,-0.90210) (0.52287,-0.91335) (0.48093,-0.92425) (0.43900,-0.93474) (0.39707,-0.94478) (0.35514,-0.95428) (0.31321,-0.96317) (0.27128,-0.97136) (0.22935,-0.97877) (0.18741,-0.98529) (0.14548,-0.99079) (0.10355,-0.99515) (0.06162,-0.99821) (0.01969,-0.99981) (-0.02224,-0.99975) (-0.06418,-0.99780) (-0.10611,-0.99371) (-0.14804,-0.98719) (-0.18997,-0.97787) (-0.23190,-0.96534) (-0.27383,-0.94909) (-0.31576,-0.92851) (-0.35770,-0.90281) (-0.39963,-0.87101) (-0.44156,-0.83176) (-0.48349,-0.78324) (-0.52542,-0.72269) (-0.56735,-0.64570) (-0.60928,-0.54395) (-0.65122,-0.39688) (-0.69315,-0.00000)};
\node at(2.9,1){$h=1$};\node at(2,-1.25){$h=1/2$};\node[below]at(-.535,0){$x_-/a$};\node[above right]at(1.228,0){$x_+/a$};
\end{tikzpicture}
\caption{闭合轨线与临界开轨线。横轴 $q=x/a$，纵轴 $w=v/\sqrt{2D/m}$；虚线对应 $E=D$。}\label{mp4:fig:l4-h4-s}
\end{figure}

再令 $z=1+\sqrt h\cos\psi$。根号等于 $\sqrt h\sin\psi$，$\dd z=-\sqrt h\sin\psi\,\dd\psi$，端点分别对应 $\psi=\pi,0$，故
\[
T=2a\sqrt{\frac m{2D}}\int_0^\pi
\frac{\dd\psi}{1+\sqrt h\cos\psi}.
\]
用半角代换 $u=\tan(\psi/2)$，$\dd\psi=2\dd u/(1+u^2)$、$\cos\psi=(1-u^2)/(1+u^2)$，上下限 $0,\pi$ 变为 $0,\infty$，得到
\begin{align*}
\int_0^\pi\frac{\dd\psi}{1+\sqrt h\cos\psi}
&=2\int_0^\infty\frac{\dd u}{(1+\sqrt h)+(1-\sqrt h)u^2}\\
&=\frac{2}{\sqrt{(1+\sqrt h)(1-\sqrt h)}}
\left[\arctan\left(u\sqrt{\frac{1-\sqrt h}{1+\sqrt h}}\right)\right]_0^\infty\\
&=\frac\pi{\sqrt{(1+\sqrt h)(1-\sqrt h)}}.
\end{align*}
因此
\[
T=\pi a\sqrt{\frac{2m}{D}}\frac1{\sqrt{1-E/D}}.
\]
在阱底附近 $U(x)=Dx^2/a^2+O(x^3)$，线性角频率为 $\omega_c=\sqrt{2D/(ma^2)}$；上式在 $E\to0$ 时确实趋于 $2\pi/\omega_c$。当 $E\uparrow D$，图中右转向点向无穷远移动，周期按 $(1-E/D)^{-1/2}$ 发散。

临界能量 $E=D$ 时，左转向点仍有限，等于 $-a\ln2$，向右则不再有有限转向点。远区有
\[
D-U=D\left(2\ee^{-x/a}-\ee^{-2x/a}\right)
\sim2D\ee^{-x/a},
\]
所以右行速度满足 $v\sim2\sqrt{D/m}\ee^{-x/(2a)}$。在远区取一个固定的 $x_0$，由 $\dd t/\dd x=1/v$ 积分，其随 $x$ 增长的首项来自
\[
\frac12\sqrt{\frac mD}\int_{x_0}^x\ee^{\xi/(2a)}\dd\xi
=a\sqrt{\frac mD}[\ee^{\xi/(2a)}]_{x_0}^x.
\]
因此速度和时间的渐近式为
\[
v\sim2\sqrt{\frac Dm}\ee^{-x/(2a)},\qquad
t\sim a\sqrt{\frac mD}\ee^{x/(2a)}.
\]
尽管 $x$ 可以无限增大，到达无穷远仍需无限时间，速度同时趋于零。反解给出 $x\sim2a\ln[t\sqrt{D/m}/a]$，表示临界逃逸在远区只按对数规律前进。
\end{solution}

\par\medskip\noindent\begin{minipage}{\textwidth}
\begin{exercise}[倒抛物线势垒]\label{mp4:exercise:l4-h5}
在 $U=-m\lambda^2x^2/2$ 的区域内，质点从 $x=a>0$ 以向左速度 $v_0<0$ 出发，见图\ref{mp4:fig:l4-h5-q}。判断它何时能越过 $x=0$，何时在右侧转回，并求临界情形的运动。
\end{exercise}

\begin{figure}[H]\centering
\begin{tikzpicture}
\draw[axis](-2,0)--(2.7,0)node[right]{$x$};
\draw[axis](0,-1.6)--(0,0.8)node[above]{$U$};
\draw[curve,domain=-1.7:1.7,samples=90]plot(\x,{-.45*(\x)^2});\fill(1.4,-.882)circle(1.5pt)node[right]{$P$};\draw[axis](1.4,-.882)--(.5,-.882)node[midway,below]{$v_0$};\node[above]at(1.4,0){$a$};
\end{tikzpicture}
\caption{倒抛物线势垒。质点从右侧 P 向左出发。}\label{mp4:fig:l4-h5-q}
\end{figure}
\end{minipage}\par\medskip
\begin{solution}
以图\ref{mp4:fig:l4-h5-q}的势垒顶为零势能，初态能量为
\[
E=\frac m2(v_0^2-\lambda^2a^2).
\]
整个运动满足 $v^2-\lambda^2x^2=2E/m$。初态 $x=a>0,v_0<0$ 位于第四象限；图\ref{mp4:fig:l4-h5-s}中的负斜率直线 $v=-\lambda x$ 是向鞍点趋近的临界支。

\begin{figure}[H]\centering
\begin{tikzpicture}
\draw[axis](-2.4,0)--(2.5,0)node[right]{$x/a$};
\draw[axis](0,-2.2)--(0,2.2)node[above]{$v/(\lambda a)$};
\draw[flow,line width=1pt] plot coordinates{(1.80000,-1.80000) (0.03000,-0.03000)};
\draw[flow] plot coordinates{(1.86448,-1.76530) (1.80213,-1.69931) (1.74211,-1.63552) (1.68434,-1.57385) (1.62877,-1.51422) (1.57530,-1.45656) (1.52387,-1.40078) (1.47442,-1.34682) (1.42688,-1.29460) (1.38119,-1.24406) (1.33729,-1.19513) (1.29512,-1.14775) (1.25463,-1.10186) (1.21577,-1.05740) (1.17848,-1.01431) (1.14272,-0.97253) (1.10844,-0.93201) (1.07560,-0.89270) (1.04416,-0.85455) (1.01406,-0.81751) (0.98528,-0.78153) (0.95778,-0.74655) (0.93152,-0.71255) (0.90647,-0.67947) (0.88259,-0.64727) (0.85985,-0.61591) (0.83823,-0.58535) (0.81770,-0.55554) (0.79822,-0.52646) (0.77978,-0.49806) (0.76236,-0.47030) (0.74592,-0.44316) (0.73044,-0.41659) (0.71591,-0.39056) (0.70232,-0.36503) (0.68963,-0.33998) (0.67783,-0.31537) (0.66692,-0.29116) (0.65686,-0.26734) (0.64766,-0.24386) (0.63930,-0.22070) (0.63177,-0.19782) (0.62506,-0.17520) (0.61915,-0.15281) (0.61405,-0.13061) (0.60975,-0.10858) (0.60623,-0.08670) (0.60350,-0.06493) (0.60156,-0.04324) (0.60039,-0.02160) (0.60000,0.00000) (0.60039,0.02160) (0.60156,0.04324) (0.60350,0.06493) (0.60623,0.08670) (0.60975,0.10858) (0.61405,0.13061) (0.61915,0.15281) (0.62506,0.17520) (0.63177,0.19782) (0.63930,0.22070) (0.64766,0.24386) (0.65686,0.26734) (0.66692,0.29116) (0.67783,0.31537) (0.68963,0.33998) (0.70232,0.36503) (0.71591,0.39056) (0.73044,0.41659) (0.74592,0.44316) (0.76236,0.47030) (0.77978,0.49806) (0.79822,0.52646) (0.81770,0.55554) (0.83823,0.58535) (0.85985,0.61591) (0.88259,0.64727) (0.90647,0.67947) (0.93152,0.71255) (0.95778,0.74655) (0.98528,0.78153) (1.01406,0.81751) (1.04416,0.85455) (1.07560,0.89270) (1.10844,0.93201) (1.14272,0.97253) (1.17848,1.01431) (1.21577,1.05740) (1.25463,1.10186) (1.29512,1.14775) (1.33729,1.19513) (1.38119,1.24406) (1.42688,1.29460) (1.47442,1.34682) (1.52387,1.40078) (1.57530,1.45656) (1.62877,1.51422) (1.68434,1.57385) (1.74211,1.63552) (1.80213,1.69931) (1.86448,1.76530)};
\draw[flow,dashed] plot coordinates{(1.50946,-1.81066) (1.47177,-1.77935) (1.43472,-1.74884) (1.39831,-1.71909) (1.36252,-1.69010) (1.32732,-1.66186) (1.29271,-1.63435) (1.25867,-1.60756) (1.22519,-1.58148) (1.19225,-1.55610) (1.15983,-1.53141) (1.12792,-1.50739) (1.09651,-1.48403) (1.06559,-1.46133) (1.03513,-1.43927) (1.00514,-1.41785) (0.97558,-1.39705) (0.94645,-1.37687) (0.91775,-1.35730) (0.88944,-1.33832) (0.86153,-1.31994) (0.83400,-1.30214) (0.80684,-1.28491) (0.78003,-1.26825) (0.75357,-1.25214) (0.72744,-1.23659) (0.70163,-1.22159) (0.67613,-1.20712) (0.65093,-1.19319) (0.62601,-1.17978) (0.60137,-1.16690) (0.57700,-1.15452) (0.55288,-1.14266) (0.52900,-1.13130) (0.50536,-1.12044) (0.48194,-1.11007) (0.45873,-1.10020) (0.43573,-1.09081) (0.41292,-1.08190) (0.39028,-1.07346) (0.36783,-1.06550) (0.34553,-1.05801) (0.32339,-1.05099) (0.30139,-1.04443) (0.27952,-1.03833) (0.25777,-1.03269) (0.23614,-1.02750) (0.21461,-1.02277) (0.19318,-1.01849) (0.17183,-1.01466) (0.15056,-1.01127) (0.12936,-1.00833) (0.10821,-1.00584) (0.08711,-1.00379) (0.06605,-1.00218) (0.04502,-1.00101) (0.02400,-1.00029) (0.00300,-1.00000) (-0.01800,-1.00016) (-0.03901,-1.00076) (-0.06004,-1.00180) (-0.08109,-1.00328) (-0.10218,-1.00521) (-0.12331,-1.00757) (-0.14450,-1.01039) (-0.16575,-1.01364) (-0.18707,-1.01735) (-0.20848,-1.02150) (-0.22998,-1.02610) (-0.25158,-1.03116) (-0.27329,-1.03667) (-0.29512,-1.04264) (-0.31709,-1.04907) (-0.33919,-1.05596) (-0.36144,-1.06332) (-0.38385,-1.07114) (-0.40643,-1.07944) (-0.42919,-1.08821) (-0.45214,-1.09747) (-0.47529,-1.10720) (-0.49865,-1.11743) (-0.52222,-1.12815) (-0.54603,-1.13936) (-0.57008,-1.15108) (-0.59438,-1.16331) (-0.61894,-1.17605) (-0.64378,-1.18931) (-0.66890,-1.20309) (-0.69431,-1.21740) (-0.72003,-1.23225) (-0.74607,-1.24765) (-0.77244,-1.26359) (-0.79914,-1.28009) (-0.82620,-1.29716) (-0.85363,-1.31479) (-0.88143,-1.33301) (-0.90962,-1.35182) (-0.93821,-1.37122) (-0.96721,-1.39122) (-0.99665,-1.41184) (-1.02652,-1.43309)};
\fill(0,0)circle(1.2pt)node[above left]{$S$};
\draw[guide](1,-1.9)--(1,1.8);\node[above right]at(1,0){$1$};\node[above]at(1.7,1.7){右返};\node[left]at(-1,-1.5){跨越};
\end{tikzpicture}
\caption{倒抛物线势垒的相图。粗线是临界稳定支；上方轨线在右侧转向，下方轨线越过纵轴。}\label{mp4:fig:l4-h5-s}
\end{figure}

若 $v_0<-\lambda a$，则 $E>0$，在 $x=0$ 时仍有 $v=-\sqrt{2E/m}<0$，可以越过势垒顶。若 $-\lambda a<v_0<0$，则 $E<0$，相轨线不能接触纵轴；令 $v=0$ 得右侧转向点
\[
x_T=\sqrt{a^2-\frac{v_0^2}{\lambda^2}}>0.
\]
到达此处时加速度 $\lambda^2x_T>0$，相点从下支转入上支，随后向右离开。

还可由时间解把三种情形逐一确定。方程 $\ddot x-\lambda^2x=0$ 的特征根为 $\pm\lambda$，通解为 $x=C_+\ee^{\lambda t}+C_-\ee^{-\lambda t}$。初值给出 $C_++C_-=a$、$\lambda(C_+-C_-)=v_0$，所以 $C_\pm=(a\pm v_0/\lambda)/2$。将指数和差写成双曲函数，初值解为
\[
x=a\cosh\lambda t+\frac{v_0}{\lambda}\sinh\lambda t,
\quad
v=\lambda a\sinh\lambda t+v_0\cosh\lambda t.
\]
越过情形由 $x(t_*)=0$ 得
\[
t_*=\frac1\lambda\operatorname{artanh}
\frac{\lambda a}{|v_0|},\qquad |v_0|>\lambda a.
\]
右返情形由 $v(t_T)=0$ 得
\[
t_T=\frac1\lambda\operatorname{artanh}
\frac{|v_0|}{\lambda a},\qquad |v_0|<\lambda a.
\]
各式中反双曲正切的自变量都在 $(0,1)$，对应有限时间。

当 $v_0=-\lambda a$，增长指数项恰好消失，解退化为
\[
x=a\ee^{-\lambda t},\qquad v=-\lambda a\ee^{-\lambda t}.
\]
它对所有有限 $t$ 均有 $x>0$，只在 $t\to\infty$ 时到达原点的极限。上面两条时间公式在初速趋近临界值时也趋于无穷大，和图中相轨线渐近进入鞍点的性质一致。
\end{solution}

\par\medskip\noindent\begin{minipage}{\textwidth}
\begin{exercise}[线性化失效的稳定点]\label{mp4:exercise:l4-h6}
质量为$m$的质点的势能为 $U(x)=\alpha x^6/6$，$\alpha>0$，见图\ref{mp4:fig:l4-h6-q}。说明原点的稳定性；振幅为 $A$ 时，求周期的积分表达式及它对 $A$ 的依赖。
\end{exercise}

\begin{figure}[H]\centering
\begin{tikzpicture}
\draw[axis](-2,0)--(2,0)node[right]{$x$};
\draw[axis](0,-0.2)--(0,2)node[above]{$U$};
\draw[curve,domain=-1.35:1.35,samples=110]plot(\x,{.25*(\x)^6});\node[below left]at(0,0){$O$};
\end{tikzpicture}
\caption{六次势阱。原点附近的势能增长比二次势阱更缓。}\label{mp4:fig:l4-h6-q}
\end{figure}
\end{minipage}\par\medskip
\begin{solution}
平衡方程为 $U'=\alpha x^5=0$，唯一平衡位置是原点，且 $U(x)>U(0)=0$ 对任意 $x\ne0$ 都成立。尽管 $U''(0)=0$，原点仍是严格势能极小点。

给定位移界限 $\delta_x$ 和速度界限 $\delta_v$，若初态能量满足
\[
E<\min\left(\frac\alpha6\delta_x^6,\frac m2\delta_v^2\right),
\]
则能量守恒使任意时刻都有 $|x|<\delta_x$、$|v|<\delta_v$，所以相空间平衡态稳定。图\ref{mp4:fig:l4-h6-s}中的闭合轨线围住原点，且可随能量减小而任意收缩。系统无耗散，非零能量轨线不会渐近进入原点。

\begin{figure}[H]\centering
\begin{tikzpicture}
\draw[axis](-2.4,0)--(2.4,0)node[right]{$q$};
\draw[axis](0,-1.7)--(0,1.7)node[above]{$w$};
\draw[flow] plot coordinates{(0.97500,-0.00000) (0.97452,-0.00971) (0.97308,-0.01938) (0.97067,-0.02897) (0.96731,-0.03845) (0.96300,-0.04778) (0.95773,-0.05692) (0.95152,-0.06585) (0.94437,-0.07453) (0.93629,-0.08292) (0.92728,-0.09102) (0.91736,-0.09878) (0.90653,-0.10620) (0.89481,-0.11325) (0.88221,-0.11991) (0.86873,-0.12618) (0.85440,-0.13204) (0.83922,-0.13750) (0.82322,-0.14255) (0.80640,-0.14719) (0.78879,-0.15143) (0.77040,-0.15527) (0.75125,-0.15873) (0.73136,-0.16183) (0.71074,-0.16457) (0.68943,-0.16698) (0.66743,-0.16907) (0.64478,-0.17088) (0.62149,-0.17242) (0.59758,-0.17371) (0.57309,-0.17479) (0.54803,-0.17567) (0.52243,-0.17638) (0.49632,-0.17695) (0.46971,-0.17739) (0.44264,-0.17772) (0.41513,-0.17797) (0.38722,-0.17816) (0.35892,-0.17828) (0.33027,-0.17837) (0.30129,-0.17843) (0.27202,-0.17846) (0.24247,-0.17849) (0.21269,-0.17850) (0.18270,-0.17850) (0.15252,-0.17850) (0.12220,-0.17851) (0.09176,-0.17851) (0.06122,-0.17851) (0.03063,-0.17851) (0.00000,-0.17851) (-0.03063,-0.17851) (-0.06122,-0.17851) (-0.09176,-0.17851) (-0.12220,-0.17851) (-0.15252,-0.17850) (-0.18270,-0.17850) (-0.21269,-0.17850) (-0.24247,-0.17849) (-0.27202,-0.17846) (-0.30129,-0.17843) (-0.33027,-0.17837) (-0.35892,-0.17828) (-0.38722,-0.17816) (-0.41513,-0.17797) (-0.44264,-0.17772) (-0.46971,-0.17739) (-0.49632,-0.17695) (-0.52243,-0.17638) (-0.54803,-0.17567) (-0.57309,-0.17479) (-0.59758,-0.17371) (-0.62149,-0.17242) (-0.64478,-0.17088) (-0.66743,-0.16907) (-0.68943,-0.16698) (-0.71074,-0.16457) (-0.73136,-0.16183) (-0.75125,-0.15873) (-0.77040,-0.15527) (-0.78879,-0.15143) (-0.80640,-0.14719) (-0.82322,-0.14255) (-0.83922,-0.13750) (-0.85440,-0.13204) (-0.86873,-0.12618) (-0.88221,-0.11991) (-0.89481,-0.11325) (-0.90653,-0.10620) (-0.91736,-0.09878) (-0.92728,-0.09102) (-0.93629,-0.08292) (-0.94437,-0.07453) (-0.95152,-0.06585) (-0.95773,-0.05692) (-0.96300,-0.04778) (-0.96731,-0.03845) (-0.97067,-0.02897) (-0.97308,-0.01938) (-0.97452,-0.00971) (-0.97500,-0.00000) (-0.97452,0.00971) (-0.97308,0.01938) (-0.97067,0.02897) (-0.96731,0.03845) (-0.96300,0.04778) (-0.95773,0.05692) (-0.95152,0.06585) (-0.94437,0.07453) (-0.93629,0.08292) (-0.92728,0.09102) (-0.91736,0.09878) (-0.90653,0.10620) (-0.89481,0.11325) (-0.88221,0.11991) (-0.86873,0.12618) (-0.85440,0.13204) (-0.83922,0.13750) (-0.82322,0.14255) (-0.80640,0.14719) (-0.78879,0.15143) (-0.77040,0.15527) (-0.75125,0.15873) (-0.73136,0.16183) (-0.71074,0.16457) (-0.68943,0.16698) (-0.66743,0.16907) (-0.64478,0.17088) (-0.62149,0.17242) (-0.59758,0.17371) (-0.57309,0.17479) (-0.54803,0.17567) (-0.52243,0.17638) (-0.49632,0.17695) (-0.46971,0.17739) (-0.44264,0.17772) (-0.41513,0.17797) (-0.38722,0.17816) (-0.35892,0.17828) (-0.33027,0.17837) (-0.30129,0.17843) (-0.27202,0.17846) (-0.24247,0.17849) (-0.21269,0.17850) (-0.18270,0.17850) (-0.15252,0.17850) (-0.12220,0.17851) (-0.09176,0.17851) (-0.06122,0.17851) (-0.03063,0.17851) (-0.00000,0.17851) (0.03063,0.17851) (0.06122,0.17851) (0.09176,0.17851) (0.12220,0.17851) (0.15252,0.17850) (0.18270,0.17850) (0.21269,0.17850) (0.24247,0.17849) (0.27202,0.17846) (0.30129,0.17843) (0.33027,0.17837) (0.35892,0.17828) (0.38722,0.17816) (0.41513,0.17797) (0.44264,0.17772) (0.46971,0.17739) (0.49632,0.17695) (0.52243,0.17638) (0.54803,0.17567) (0.57309,0.17479) (0.59758,0.17371) (0.62149,0.17242) (0.64478,0.17088) (0.66743,0.16907) (0.68943,0.16698) (0.71074,0.16457) (0.73136,0.16183) (0.75125,0.15873) (0.77040,0.15527) (0.78879,0.15143) (0.80640,0.14719) (0.82322,0.14255) (0.83922,0.13750) (0.85440,0.13204) (0.86873,0.12618) (0.88221,0.11991) (0.89481,0.11325) (0.90653,0.10620) (0.91736,0.09878) (0.92728,0.09102) (0.93629,0.08292) (0.94437,0.07453) (0.95152,0.06585) (0.95773,0.05692) (0.96300,0.04778) (0.96731,0.03845) (0.97067,0.02897) (0.97308,0.01938) (0.97452,0.00971) (0.97500,0.00000)};
\draw[flow] plot coordinates{(1.35000,-0.00000) (1.34933,-0.02577) (1.34734,-0.05143) (1.34401,-0.07690) (1.33935,-0.10206) (1.33338,-0.12682) (1.32609,-0.15110) (1.31749,-0.17479) (1.30759,-0.19783) (1.29640,-0.22013) (1.28393,-0.24161) (1.27019,-0.26222) (1.25520,-0.28191) (1.23897,-0.30061) (1.22152,-0.31830) (1.20286,-0.33494) (1.18301,-0.35051) (1.16200,-0.36500) (1.13984,-0.37840) (1.11656,-0.39072) (1.09217,-0.40197) (1.06671,-0.41217) (1.04019,-0.42137) (1.01265,-0.42958) (0.98411,-0.43685) (0.95459,-0.44325) (0.92414,-0.44881) (0.89277,-0.45360) (0.86052,-0.45768) (0.82742,-0.46112) (0.79351,-0.46398) (0.75881,-0.46632) (0.72337,-0.46821) (0.68721,-0.46971) (0.65037,-0.47088) (0.61289,-0.47177) (0.57480,-0.47244) (0.53615,-0.47292) (0.49697,-0.47326) (0.45730,-0.47349) (0.41717,-0.47364) (0.37664,-0.47374) (0.33573,-0.47379) (0.29449,-0.47382) (0.25296,-0.47384) (0.21119,-0.47385) (0.16920,-0.47385) (0.12705,-0.47385) (0.08477,-0.47385) (0.04240,-0.47385) (0.00000,-0.47385) (-0.04240,-0.47385) (-0.08477,-0.47385) (-0.12705,-0.47385) (-0.16920,-0.47385) (-0.21119,-0.47385) (-0.25296,-0.47384) (-0.29449,-0.47382) (-0.33573,-0.47379) (-0.37664,-0.47374) (-0.41717,-0.47364) (-0.45730,-0.47349) (-0.49697,-0.47326) (-0.53615,-0.47292) (-0.57480,-0.47244) (-0.61289,-0.47177) (-0.65037,-0.47088) (-0.68721,-0.46971) (-0.72337,-0.46821) (-0.75881,-0.46632) (-0.79351,-0.46398) (-0.82742,-0.46112) (-0.86052,-0.45768) (-0.89277,-0.45360) (-0.92414,-0.44881) (-0.95459,-0.44325) (-0.98411,-0.43685) (-1.01265,-0.42958) (-1.04019,-0.42137) (-1.06671,-0.41217) (-1.09217,-0.40197) (-1.11656,-0.39072) (-1.13984,-0.37840) (-1.16200,-0.36500) (-1.18301,-0.35051) (-1.20286,-0.33494) (-1.22152,-0.31830) (-1.23897,-0.30061) (-1.25520,-0.28191) (-1.27019,-0.26222) (-1.28393,-0.24161) (-1.29640,-0.22013) (-1.30759,-0.19783) (-1.31749,-0.17479) (-1.32609,-0.15110) (-1.33338,-0.12682) (-1.33935,-0.10206) (-1.34401,-0.07690) (-1.34734,-0.05143) (-1.34933,-0.02577) (-1.35000,-0.00000) (-1.34933,0.02577) (-1.34734,0.05143) (-1.34401,0.07690) (-1.33935,0.10206) (-1.33338,0.12682) (-1.32609,0.15110) (-1.31749,0.17479) (-1.30759,0.19783) (-1.29640,0.22013) (-1.28393,0.24161) (-1.27019,0.26222) (-1.25520,0.28191) (-1.23897,0.30061) (-1.22152,0.31830) (-1.20286,0.33494) (-1.18301,0.35051) (-1.16200,0.36500) (-1.13984,0.37840) (-1.11656,0.39072) (-1.09217,0.40197) (-1.06671,0.41217) (-1.04019,0.42137) (-1.01265,0.42958) (-0.98411,0.43685) (-0.95459,0.44325) (-0.92414,0.44881) (-0.89277,0.45360) (-0.86052,0.45768) (-0.82742,0.46112) (-0.79351,0.46398) (-0.75881,0.46632) (-0.72337,0.46821) (-0.68721,0.46971) (-0.65037,0.47088) (-0.61289,0.47177) (-0.57480,0.47244) (-0.53615,0.47292) (-0.49697,0.47326) (-0.45730,0.47349) (-0.41717,0.47364) (-0.37664,0.47374) (-0.33573,0.47379) (-0.29449,0.47382) (-0.25296,0.47384) (-0.21119,0.47385) (-0.16920,0.47385) (-0.12705,0.47385) (-0.08477,0.47385) (-0.04240,0.47385) (-0.00000,0.47385) (0.04240,0.47385) (0.08477,0.47385) (0.12705,0.47385) (0.16920,0.47385) (0.21119,0.47385) (0.25296,0.47384) (0.29449,0.47382) (0.33573,0.47379) (0.37664,0.47374) (0.41717,0.47364) (0.45730,0.47349) (0.49697,0.47326) (0.53615,0.47292) (0.57480,0.47244) (0.61289,0.47177) (0.65037,0.47088) (0.68721,0.46971) (0.72337,0.46821) (0.75881,0.46632) (0.79351,0.46398) (0.82742,0.46112) (0.86052,0.45768) (0.89277,0.45360) (0.92414,0.44881) (0.95459,0.44325) (0.98411,0.43685) (1.01265,0.42958) (1.04019,0.42137) (1.06671,0.41217) (1.09217,0.40197) (1.11656,0.39072) (1.13984,0.37840) (1.16200,0.36500) (1.18301,0.35051) (1.20286,0.33494) (1.22152,0.31830) (1.23897,0.30061) (1.25520,0.28191) (1.27019,0.26222) (1.28393,0.24161) (1.29640,0.22013) (1.30759,0.19783) (1.31749,0.17479) (1.32609,0.15110) (1.33338,0.12682) (1.33935,0.10206) (1.34401,0.07690) (1.34734,0.05143) (1.34933,0.02577) (1.35000,0.00000)};
\draw[flow] plot coordinates{(1.80000,-0.00000) (1.79911,-0.06108) (1.79645,-0.12191) (1.79201,-0.18227) (1.78581,-0.24192) (1.77784,-0.30062) (1.76812,-0.35816) (1.75665,-0.41433) (1.74345,-0.46893) (1.72853,-0.52178) (1.71190,-0.57271) (1.69359,-0.62157) (1.67360,-0.66823) (1.65196,-0.71256) (1.62869,-0.75449) (1.60381,-0.79393) (1.57735,-0.83084) (1.54934,-0.86518) (1.51979,-0.89694) (1.48875,-0.92615) (1.45623,-0.95282) (1.42228,-0.97701) (1.38692,-0.99879) (1.35020,-1.01826) (1.31214,-1.03551) (1.27279,-1.05066) (1.23218,-1.06384) (1.19036,-1.07520) (1.14736,-1.08488) (1.10323,-1.09302) (1.05801,-1.09980) (1.01175,-1.10535) (0.96449,-1.10983) (0.91627,-1.11339) (0.86716,-1.11616) (0.81718,-1.11827) (0.76640,-1.11985) (0.71487,-1.12099) (0.66262,-1.12180) (0.60973,-1.12235) (0.55623,-1.12271) (0.50218,-1.12294) (0.44764,-1.12307) (0.39266,-1.12314) (0.33729,-1.12318) (0.28158,-1.12319) (0.22560,-1.12320) (0.16939,-1.12320) (0.11302,-1.12320) (0.05654,-1.12320) (0.00000,-1.12320) (-0.05654,-1.12320) (-0.11302,-1.12320) (-0.16939,-1.12320) (-0.22560,-1.12320) (-0.28158,-1.12319) (-0.33729,-1.12318) (-0.39266,-1.12314) (-0.44764,-1.12307) (-0.50218,-1.12294) (-0.55623,-1.12271) (-0.60973,-1.12235) (-0.66262,-1.12180) (-0.71487,-1.12099) (-0.76640,-1.11985) (-0.81718,-1.11827) (-0.86716,-1.11616) (-0.91627,-1.11339) (-0.96449,-1.10983) (-1.01175,-1.10535) (-1.05801,-1.09980) (-1.10323,-1.09302) (-1.14736,-1.08488) (-1.19036,-1.07520) (-1.23218,-1.06384) (-1.27279,-1.05066) (-1.31214,-1.03551) (-1.35020,-1.01826) (-1.38692,-0.99879) (-1.42228,-0.97701) (-1.45623,-0.95282) (-1.48875,-0.92615) (-1.51979,-0.89694) (-1.54934,-0.86518) (-1.57735,-0.83084) (-1.60381,-0.79393) (-1.62869,-0.75449) (-1.65196,-0.71256) (-1.67360,-0.66823) (-1.69359,-0.62157) (-1.71190,-0.57271) (-1.72853,-0.52178) (-1.74345,-0.46893) (-1.75665,-0.41433) (-1.76812,-0.35816) (-1.77784,-0.30062) (-1.78581,-0.24192) (-1.79201,-0.18227) (-1.79645,-0.12191) (-1.79911,-0.06108) (-1.80000,-0.00000) (-1.79911,0.06108) (-1.79645,0.12191) (-1.79201,0.18227) (-1.78581,0.24192) (-1.77784,0.30062) (-1.76812,0.35816) (-1.75665,0.41433) (-1.74345,0.46893) (-1.72853,0.52178) (-1.71190,0.57271) (-1.69359,0.62157) (-1.67360,0.66823) (-1.65196,0.71256) (-1.62869,0.75449) (-1.60381,0.79393) (-1.57735,0.83084) (-1.54934,0.86518) (-1.51979,0.89694) (-1.48875,0.92615) (-1.45623,0.95282) (-1.42228,0.97701) (-1.38692,0.99879) (-1.35020,1.01826) (-1.31214,1.03551) (-1.27279,1.05066) (-1.23218,1.06384) (-1.19036,1.07520) (-1.14736,1.08488) (-1.10323,1.09302) (-1.05801,1.09980) (-1.01175,1.10535) (-0.96449,1.10983) (-0.91627,1.11339) (-0.86716,1.11616) (-0.81718,1.11827) (-0.76640,1.11985) (-0.71487,1.12099) (-0.66262,1.12180) (-0.60973,1.12235) (-0.55623,1.12271) (-0.50218,1.12294) (-0.44764,1.12307) (-0.39266,1.12314) (-0.33729,1.12318) (-0.28158,1.12319) (-0.22560,1.12320) (-0.16939,1.12320) (-0.11302,1.12320) (-0.05654,1.12320) (-0.00000,1.12320) (0.05654,1.12320) (0.11302,1.12320) (0.16939,1.12320) (0.22560,1.12320) (0.28158,1.12319) (0.33729,1.12318) (0.39266,1.12314) (0.44764,1.12307) (0.50218,1.12294) (0.55623,1.12271) (0.60973,1.12235) (0.66262,1.12180) (0.71487,1.12099) (0.76640,1.11985) (0.81718,1.11827) (0.86716,1.11616) (0.91627,1.11339) (0.96449,1.10983) (1.01175,1.10535) (1.05801,1.09980) (1.10323,1.09302) (1.14736,1.08488) (1.19036,1.07520) (1.23218,1.06384) (1.27279,1.05066) (1.31214,1.03551) (1.35020,1.01826) (1.38692,0.99879) (1.42228,0.97701) (1.45623,0.95282) (1.48875,0.92615) (1.51979,0.89694) (1.54934,0.86518) (1.57735,0.83084) (1.60381,0.79393) (1.62869,0.75449) (1.65196,0.71256) (1.67360,0.66823) (1.69359,0.62157) (1.71190,0.57271) (1.72853,0.52178) (1.74345,0.46893) (1.75665,0.41433) (1.76812,0.35816) (1.77784,0.30062) (1.78581,0.24192) (1.79201,0.18227) (1.79645,0.12191) (1.79911,0.06108) (1.80000,0.00000)};
\fill(0,0)circle(1.2pt);
\end{tikzpicture}
\caption{六次势阱的一族相轨线。横轴 $q=x/L$，纵轴 $w=v/(L^3\sqrt{\alpha/(3m)})$，三条轨线使用相同的固定尺度 $L$。}\label{mp4:fig:l4-h6-s}
\end{figure}

振幅为 $A$ 时，转向点给出 $E=\alpha A^6/6$。代入能量式，
\[
v^2=\frac\alpha{3m}(A^6-x^6),\qquad
v_{\max}=A^3\sqrt{\frac\alpha{3m}}.
\]
图中各条轨线使用同一固定尺度 $L$，因而其截距随振幅改变；若每条轨线分别用自身的 $A$ 和 $v_{\max}$ 归一化，则都成为 $w^2+q^6=1$。

周期由一个四分之一行程积分得到：
\begin{align*}
T&=4\sqrt{\frac{3m}{\alpha}}\int_0^A
\frac{\dd x}{\sqrt{A^6-x^6}}\\
&=4\sqrt{\frac{3m}{\alpha}}A^{-2}
\int_0^1\frac{\dd u}{\sqrt{1-u^6}}.
\end{align*}
作变量代换 $x=Au$ 后，$\dd x=A\dd u$，$\sqrt{A^6-x^6}=A^3\sqrt{1-u^6}$，上下限 $0,A$ 变为 $0,1$，因此振幅因子为 $A^{-2}$。在 $u=1$ 附近，$1-u^6\sim6(1-u)$，积分有限；因此
\[
T\propto A^{-2},\qquad \omega_{\rm osc}=\frac{2\pi}T\propto A^2.
\]
振幅减半，周期变为四倍。回复力的首个非零项是 $-\alpha x^5$。只保留一次项会得到 $\ddot x=0$，而本题的稳定性和振动周期都须保留五次回复力后求得。
还可从运动方程直接确认振幅标度。令$x=Aq$、$t=\sqrt{m/\alpha}\,A^{-2}s$，方程$m\ddot x+\alpha x^5=0$化为
\[
 \frac{\dd^2q}{\dd s^2}+q^5=0.
\]
归一化后的振幅与周期都是定值，恢复时间单位就得到$T\propto A^{-2}$。图\ref{mp4:fig:l4-h6-s}中，小振幅相轨线沿速度方向收缩，运动随之减慢。周期趋于无穷与平衡点稳定可以同时成立：稳定性要求邻近初态保持邻近，并不要求周期有界。
\end{solution}

\par\medskip\noindent\begin{minipage}{\textwidth}
\begin{exercise}[从侧向位置启动单摆]\label{mp4:exercise:l4-h7}
刚杆单摆从 $\theta_0=\pi/2$ 的位置获得正向角速度 $\dot\theta_0$，定义与图\ref{mp4:fig:l4-h7-q}一致。求能越过上方倒立位置的角速度条件，并说明临界运动在倒立位置附近怎样趋近。
\end{exercise}

\begin{figure}[H]\centering
\begin{tikzpicture}
\fill(-1.2,0)circle(2pt)node[above]{$O$};\draw[curve](-1.2,0)--(1.2,0)node[midway,above]{$l$};\fill(1.2,0)circle(3pt)node[right]{$m$};\draw[guide](-1.2,0)--(-1.2,-1.6);\draw[axis](1.2,0)--(1.2,1.1)node[right]{$\vct v_0$};\draw[axis](-1.2,-.7)arc(-90:0:.7);\node at(-.45,-.65){$\pi/2$};
\end{tikzpicture}
\caption{刚杆处于水平位置，正向初速度使摆球上升。}\label{mp4:fig:l4-h7-q}
\end{figure}
\end{minipage}\par\medskip
\begin{solution}
图\ref{mp4:fig:l4-h7-q}中，摆球从水平位置向上运动。采用例\ref{mp4:example:l4-e6}的无量纲变量，初始位置 $\theta_0=\pi/2$ 对应势能 $mgl$，所以
\[
h=1+\frac12p_0^2,\qquad
p_0=\frac{\dot\theta_0}{\omega_0}>0.
\]
倒立位置 $\theta=\pi$ 的势能为 $2mgl$。真正越过该位置时角速度仍须为正，故要求 $h>2$，即
\[
\dot\theta_0>\sqrt2\,\omega_0.
\]
若 $0<p_0<\sqrt2$，转向角满足 $\cos\theta_m=1-h=-p_0^2/2$，因而 $\pi/2<\theta_m<\pi$；摆球会在到达倒立位置之前反向。图\ref{mp4:fig:l4-h7-s}中 $P=(\pi/2,\sqrt2)$ 是临界初态，位于上方分界支上。

\begin{figure}[H]\centering
\begin{tikzpicture}
\draw[axis](-2.8,0)--(2.9,0)node[right]{$\theta$};
\draw[axis](0,-0.4)--(0,2.7)node[above]{$p$};
\draw[flow,line width=.9pt] plot coordinates{(-2.26195,0.00000) (-2.21671,0.06282) (-2.17147,0.12558) (-2.12623,0.18822) (-2.08099,0.25067) (-2.03575,0.31287) (-1.99051,0.37476) (-1.94527,0.43629) (-1.90004,0.49738) (-1.85480,0.55798) (-1.80956,0.61803) (-1.76432,0.67748) (-1.71908,0.73625) (-1.67384,0.79430) (-1.62860,0.85156) (-1.58336,0.90798) (-1.53812,0.96351) (-1.49288,1.01808) (-1.44765,1.07165) (-1.40241,1.12417) (-1.35717,1.17557) (-1.31193,1.22581) (-1.26669,1.27485) (-1.22145,1.32262) (-1.17621,1.36909) (-1.13097,1.41421) (-1.08573,1.45794) (-1.04050,1.50022) (-0.99526,1.54103) (-0.95002,1.58031) (-0.90478,1.61803) (-0.85954,1.65416) (-0.81430,1.68866) (-0.76906,1.72148) (-0.72382,1.75261) (-0.67858,1.78201) (-0.63335,1.80965) (-0.58811,1.83551) (-0.54287,1.85955) (-0.49763,1.88176) (-0.45239,1.90211) (-0.40715,1.92059) (-0.36191,1.93717) (-0.31667,1.95183) (-0.27143,1.96457) (-0.22619,1.97538) (-0.18096,1.98423) (-0.13572,1.99112) (-0.09048,1.99605) (-0.04524,1.99901) (0.00000,2.00000) (0.04524,1.99901) (0.09048,1.99605) (0.13572,1.99112) (0.18096,1.98423) (0.22619,1.97538) (0.27143,1.96457) (0.31667,1.95183) (0.36191,1.93717) (0.40715,1.92059) (0.45239,1.90211) (0.49763,1.88176) (0.54287,1.85955) (0.58811,1.83551) (0.63335,1.80965) (0.67858,1.78201) (0.72382,1.75261) (0.76906,1.72148) (0.81430,1.68866) (0.85954,1.65416) (0.90478,1.61803) (0.95002,1.58031) (0.99526,1.54103) (1.04050,1.50022) (1.08573,1.45794) (1.13097,1.41421) (1.17621,1.36909) (1.22145,1.32262) (1.26669,1.27485) (1.31193,1.22581) (1.35717,1.17557) (1.40241,1.12417) (1.44765,1.07165) (1.49288,1.01808) (1.53812,0.96351) (1.58336,0.90798) (1.62860,0.85156) (1.67384,0.79430) (1.71908,0.73625) (1.76432,0.67748) (1.80956,0.61803) (1.85480,0.55798) (1.90004,0.49738) (1.94527,0.43629) (1.99051,0.37476) (2.03575,0.31287) (2.08099,0.25067) (2.12623,0.18822) (2.17147,0.12558) (2.21671,0.06282) (2.26195,0.00000)};
\draw[flow,dashed] plot coordinates{(-2.26195,0.77460) (-2.21671,0.77714) (-2.17147,0.78471) (-2.12623,0.79714) (-2.08099,0.81415) (-2.03575,0.83540) (-1.99051,0.86049) (-1.94527,0.88901) (-1.90004,0.92054) (-1.85480,0.95464) (-1.80956,0.99094) (-1.76432,1.02906) (-1.71908,1.06867) (-1.67384,1.10946) (-1.62860,1.15115) (-1.58336,1.19349) (-1.53812,1.23626) (-1.49288,1.27925) (-1.44765,1.32229) (-1.40241,1.36519) (-1.35717,1.40782) (-1.31193,1.45004) (-1.26669,1.49172) (-1.22145,1.53275) (-1.17621,1.57303) (-1.13097,1.61245) (-1.08573,1.65093) (-1.04050,1.68839) (-0.99526,1.72475) (-0.95002,1.75994) (-0.90478,1.79389) (-0.85954,1.82654) (-0.81430,1.85784) (-0.76906,1.88773) (-0.72382,1.91616) (-0.67858,1.94308) (-0.63335,1.96846) (-0.58811,1.99226) (-0.54287,2.01443) (-0.49763,2.03495) (-0.45239,2.05379) (-0.40715,2.07091) (-0.36191,2.08629) (-0.31667,2.09992) (-0.27143,2.11177) (-0.22619,2.12182) (-0.18096,2.13006) (-0.13572,2.13649) (-0.09048,2.14108) (-0.04524,2.14384) (0.00000,2.14476) (0.04524,2.14384) (0.09048,2.14108) (0.13572,2.13649) (0.18096,2.13006) (0.22619,2.12182) (0.27143,2.11177) (0.31667,2.09992) (0.36191,2.08629) (0.40715,2.07091) (0.45239,2.05379) (0.49763,2.03495) (0.54287,2.01443) (0.58811,1.99226) (0.63335,1.96846) (0.67858,1.94308) (0.72382,1.91616) (0.76906,1.88773) (0.81430,1.85784) (0.85954,1.82654) (0.90478,1.79389) (0.95002,1.75994) (0.99526,1.72475) (1.04050,1.68839) (1.08573,1.65093) (1.13097,1.61245) (1.17621,1.57303) (1.22145,1.53275) (1.26669,1.49172) (1.31193,1.45004) (1.35717,1.40782) (1.40241,1.36519) (1.44765,1.32229) (1.49288,1.27925) (1.53812,1.23626) (1.58336,1.19349) (1.62860,1.15115) (1.67384,1.10946) (1.71908,1.06867) (1.76432,1.02906) (1.80956,0.99094) (1.85480,0.95464) (1.90004,0.92054) (1.94527,0.88901) (1.99051,0.86049) (2.03575,0.83540) (2.08099,0.81415) (2.12623,0.79714) (2.17147,0.78471) (2.21671,0.77714) (2.26195,0.77460)};
\fill(1.1309733552923256,1.4142135623730951)circle(1.2pt)node[below left]{$P$};
\draw[guide](1.131,0)--(1.131,1.414);\node[below]at(1.131,0){$\pi/2$};\node[below]at(2.262,0){$\pi$};\draw[fill=white](2.262,0)circle(1.5pt);
\end{tikzpicture}
\caption{从侧向位置出发的临界相点 P。上方虚线代表足够大的初始角速度，能够越过右端倒立位置。}\label{mp4:fig:l4-h7-s}
\end{figure}

临界时 $h=2$，上行角速度为 $\dot\theta=2\omega_0\cos(\theta/2)$。使用例\ref{mp4:example:l4-e6}已求出的原函数，并保留本题初角，
\[
\int_{\pi/2}^{\theta}\frac{\dd u}{2\cos(u/2)}
=\left[\ln\tan\left(\frac\pi4+\frac u4\right)\right]_{\pi/2}^{\theta}
=\omega_0\int_0^t\dd s=\omega_0t.
\]
移项并取指数，得到
\[
\tan\left(\frac\pi4+\frac\theta4\right)
=\tan\frac{3\pi}8\,\ee^{\omega_0t}
=(1+\sqrt2)\ee^{\omega_0t},
\]
即
\[
\theta(t)=4\arctan\left[(1+\sqrt2)\ee^{\omega_0t}\right]-\pi.
\]
在 $t=0$ 时该式给出 $\theta=\pi/2$，且初始角速度为 $\sqrt2\omega_0$，符合 P 点。

令 $\delta=\pi-\theta$。当 $t\to\infty$ 时，利用 $\arctan z=\pi/2-z^{-1}+O(z^{-3})$，有
\[
\delta=\frac4{1+\sqrt2}\ee^{-\omega_0t}
+O(\ee^{-3\omega_0t}).
\]
因此角度差和角速度均按指数规律衰减。从局部方程也可看出 $\dot\delta=-2\omega_0\sin(\delta/2)\sim-\omega_0\delta$。到达任何非零小角度差所需时间有限，恰好到达 $\delta=0$ 却需无限时间。临界初速只能使相点渐近鞍点，越过倒立位置必须取严格大于阈值的初速。
\end{solution}

\par\medskip\noindent\begin{minipage}{\textwidth}
\begin{exercise}[旋转圆环中的分岔（选做）]\label{mp4:exercise:l4-h8}
半径为 $R$ 的光滑竖直圆环绕竖直直径以恒定角速度 $\Omega$ 转动，小珠套在圆环上。$\theta$ 从最低点起算，见图\ref{mp4:fig:l4-h8-q}。在随环坐标中，有效势能为
\[
U_{\rm eff}=mgR(1-\cos\theta)-\frac12m\Omega^2R^2\sin^2\theta.
\]
以 $\nu=\Omega^2R/g$ 为参数，求平衡点及其稳定性、小振动频率，并讨论 $\nu=1$。
\end{exercise}

\begin{figure}[H]\centering
\begin{tikzpicture}
\draw[curve](0,0)circle(1.5);\draw[guide](0,-1.8)--(0,1.8);\draw[axis](0,1.8)--(0,2.3)node[right]{$\vct\Omega$};\draw[curve](0,0)--(.964,-1.149);\fill(.964,-1.149)circle(2.5pt)node[right]{$m$};\draw[axis](0,-.6)arc(-90:-50:.6);\node at(.38,-.7){$\theta$};\fill(0,0)circle(1.3pt)node[below left]{$O$};\draw[guide](0,0)--(-1.06066,1.06066)node[midway,above left]{$R$};
\end{tikzpicture}
\caption{绕竖直直径转动的圆环。摆角从最低半径转到小珠所在半径。}\label{mp4:fig:l4-h8-q}
\end{figure}
\end{minipage}\par\medskip
\begin{solution}
先从给定的有效势能求随环运动的方程。令 $\omega_0=\sqrt{g/R}$、$u=U_{\rm eff}/(mgR)$，则
\[
u(\theta)=1-\cos\theta-\frac\nu2\sin^2\theta,
\qquad u'=\sin\theta-\nu\sin\theta\cos\theta
=\sin\theta(1-\nu\cos\theta),
\]
\[
\ddot\theta+\omega_0^2\sin\theta(1-\nu\cos\theta)=0.
\]
图\ref{mp4:fig:l4-h8-q}中的小珠沿环运动，等效转动惯量为 $mR^2$。在平衡角 $\theta_e$ 附近，令 $\theta=\theta_e+\eta$，线性方程为 $\ddot\eta+\omega_0^2u''(\theta_e)\eta=0$，因此 $u''>0$ 时有小振动，$u''<0$ 时有指数增长模态。

每个 $2\pi$ 区间内总有 $\theta=0,\pi$ 两个平衡角。当 $\nu>1$，还存在 $\theta=\pm\theta_e$，满足 $\cos\theta_e=1/\nu$。图\ref{mp4:fig:l4-h8-s}中的两个势阱底便对应这两个非零平衡。

\begin{figure}[H]\centering
\begin{tikzpicture}
\draw[axis](-2.4,0)--(2.4,0)node[right]{$\theta$};
\draw[axis](0,-0.6)--(0,1.1)node[above]{$u$};
\draw[curve,domain=-1.8:1.8,samples=120]plot(\x,{2*(1-cos(\x r)-.75*sin(\x r)^2)});\draw[curve,dashed,domain=-1.8:1.8,samples=120]plot(\x,{2*(1-cos(\x r)-.5*sin(\x r)^2)});\fill(.84107,-.16667)circle(1.2pt)node[below]{$\theta_e$};\fill(-.84107,-.16667)circle(1.2pt)node[below]{$-\theta_e$};
\end{tikzpicture}
\caption{旋转圆环在临界参数附近的有效势能。实线取 $\nu=3/2$，虚线取 $\nu=1$；新平衡点位于两个极小值处。}\label{mp4:fig:l4-h8-s}
\end{figure}

对二阶导数逐点计算，
\[
u''=\cos\theta(1-\nu\cos\theta)+\nu\sin^2\theta.
\]
在最低点 $\theta=0$，$u''=1-\nu$：$0\le\nu<1$ 时为中心，小振动角频率为 $\omega_c=\omega_0\sqrt{1-\nu}$；$\nu>1$ 时成为鞍点。在最高点 $\theta=\pi$，$u''=-1-\nu<0$，始终不稳定。

对新增平衡 $\pm\theta_e$，第一项因平衡条件而为零，所以
\[
u''(\pm\theta_e)=\nu\sin^2\theta_e=\nu-\frac1\nu>0,
\qquad \omega_e=\omega_0\sqrt{\nu-\frac1\nu}.
\]
图\ref{mp4:fig:l4-h8-branches}将最低点附近的平衡角随参数展开；两条非零稳定分支在 $\nu=1$ 之后出现。

\begin{figure}[H]\centering
\begin{tikzpicture}[x=1.8cm,y=1.05cm]
\draw[axis](-.1,-1.65)--(2.9,-1.65)node[right]{$\nu$};
\draw[axis](0,-1.8)--(0,1.65)node[above]{$\theta_e$};
\draw[curve](0,0)--(1,0);\draw[guide](1,0)--(2.7,0);
\draw[curve,domain=1:2.7,samples=100]plot(\x,{acos(1/\x)*pi/180});
\draw[curve,domain=1:2.7,samples=100]plot(\x,{-acos(1/\x)*pi/180});
\draw[guide](1,-1.65)--(1,0);\node[below]at(1,-1.65){$1$};
\node[left]at(0,0){$0$};\fill(1,0)circle(1.5pt);
\node[right]at(2.85,1.18){稳定};\node[right]at(2.85,-1.18){稳定};
\node[above]at(2.1,0){不稳定};
\end{tikzpicture}
\caption{最低点附近的平衡分支，角度取 $|\theta|<\pi$。实线表示稳定平衡，虚线表示不稳定平衡；临界点本身须由高阶势能判断。}\label{mp4:fig:l4-h8-branches}
\end{figure}

在 $\nu=1$ 时，二次项消失，但有效势能可精确写为
\[
u(\theta)=\frac12(1-\cos\theta)^2
=\frac18\theta^4+O(\theta^6).
\]
原点仍是严格极小值，平衡态保持稳定；线性频率为零时，振动须由四次势能项计算。若小振动的角振幅为 $A\ll1$，保留四次势能，有
\[
\frac12mR^2\dot\theta^2\simeq\frac{mgR}{8}(A^4-\theta^4),
\qquad
\dot\theta^2\simeq\frac{\omega_0^2}{4}(A^4-\theta^4).
\]
取正角速度的四分之一行程，令 $\theta=Az$、$\dd\theta=A\dd z$，根式提出 $A^2$，上下限变为 $0,1$，得到
\[
T\sim4\int_0^A\frac{\dd\theta}{(\omega_0/2)\sqrt{A^4-\theta^4}}
=\frac8{\omega_0A}\int_0^1\frac{\dd z}{\sqrt{1-z^4}}.
\]
因此临界参数处的振动频率与角振幅成正比，随振幅趋零而趋零。另令 $\nu=1+\epsilon$，$0<\epsilon\ll1$，在 $\cos\theta_e=1/\nu$ 中分别展开两边：$1-\theta_e^2/2+O(\theta_e^4)=1-\epsilon+O(\epsilon^2)$。比较首项得 $\theta_e\sim\sqrt{2\epsilon}$，说明两分支如何从最低点分出。

这里守恒的是随环方程的积分
\[
J=\frac12mR^2\dot\theta^2+U_{\rm eff}.
\]
惯性系机械能还包含方位运动动能，满足 $E_{\rm inertial}=J+\Omega L_z$，其中 $L_z=m\Omega R^2\sin^2\theta$。当小珠沿环移动时 $L_z$ 改变，电机以功率 $\Omega\dot L_z$ 供能，所以不能把 J 当作惯性系机械能。
\end{solution}

\par\medskip\noindent\begin{minipage}{\textwidth}
\begin{exercise}[欠阻尼系统的第一次最大位移]\label{mp4:exercise:l4-h9}
欠阻尼振子满足式\eqref{mp4:eq:l4-damped}，初态 $x(0)=0$、$v(0)=V>0$。初始冲量在图\ref{mp4:fig:l4-h9-q}中表示。求第一次最大位移及到达时间，并求以后同侧峰值的比值。
\end{exercise}

\begin{figure}[H]\centering
\begin{tikzpicture}
\draw[axis](-2,0)--(2,0)node[right]{$x$};
\draw[axis](0,-0.4)--(0,2)node[above]{$v$};
\draw[axis,line width=.9pt](0,0)--(0,1.4);\fill(0,1.4)circle(1.5pt)node[right]{$P:(0,V)$};\node[below left]at(0,0){$O$};
\end{tikzpicture}
\caption{冲量使静止在平衡位置的物体获得正速度。竖直线段表示相点的瞬时跳变。}\label{mp4:fig:l4-h9-q}
\end{figure}
\end{minipage}\par\medskip
\begin{solution}
欠阻尼条件为 $0<\beta<\omega_0$，令 $\omega_d=\sqrt{\omega_0^2-\beta^2}$。使用例\ref{mp4:example:l4-e8}中已展开的通解及其速度式，$x(0)=0$ 给出余弦系数 $A=0$，$v(0)=V$ 给出 $-\beta A+\omega_dB=V$，即 $B=V/\omega_d$。因此
\[
x=\frac V{\omega_d}\ee^{-\beta t}\sin\omega_dt,
\quad
v=V\ee^{-\beta t}\left(\cos\omega_dt-\frac\beta{\omega_d}\sin\omega_dt\right).
\]
图\ref{mp4:fig:l4-h9-s}中的初态 P 在纵轴正半轴，质点首先向右运动，相点进入第一象限。第一次最大位移在 B 处，对应 $v=0$ 且加速度向左。

\begin{figure}[H]\centering
\begin{tikzpicture}
\draw[axis](-1.7,0)--(2,0)node[right]{$x\omega_0/V$};
\draw[axis](0,-1.5)--(0,2.2)node[above]{$v/V$};
\draw[flow] plot coordinates{(0.00000,1.80000) (0.06067,1.76863) (0.12025,1.73576) (0.17869,1.70146) (0.23593,1.66578) (0.29194,1.62880) (0.34668,1.59057) (0.40009,1.55117) (0.45214,1.51065) (0.50280,1.46908) (0.55203,1.42653) (0.59980,1.38306) (0.64607,1.33874) (0.69082,1.29363) (0.73403,1.24780) (0.77567,1.20132) (0.81571,1.15424) (0.85415,1.10663) (0.89096,1.05855) (0.92612,1.01007) (0.95964,0.96125) (0.99149,0.91215) (1.02166,0.86284) (1.05016,0.81336) (1.07697,0.76379) (1.10209,0.71418) (1.12553,0.66459) (1.14729,0.61507) (1.16736,0.56568) (1.18576,0.51647) (1.20248,0.46750) (1.21755,0.41882) (1.23097,0.37049) (1.24275,0.32254) (1.25290,0.27503) (1.26145,0.22801) (1.26841,0.18151) (1.27380,0.13560) (1.27764,0.09030) (1.27995,0.04566) (1.28075,0.00172) (1.28008,-0.04148) (1.27794,-0.08390) (1.27438,-0.12551) (1.26942,-0.16628) (1.26308,-0.20617) (1.25541,-0.24515) (1.24642,-0.28319) (1.23616,-0.32027) (1.22466,-0.35635) (1.21194,-0.39142) (1.19805,-0.42545) (1.18302,-0.45842) (1.16689,-0.49030) (1.14970,-0.52109) (1.13147,-0.55076) (1.11226,-0.57930) (1.09209,-0.60669) (1.07101,-0.63293) (1.04907,-0.65800) (1.02628,-0.68189) (1.00271,-0.70460) (0.97838,-0.72612) (0.95335,-0.74644) (0.92764,-0.76557) (0.90130,-0.78349) (0.87438,-0.80022) (0.84690,-0.81574) (0.81892,-0.83007) (0.79047,-0.84321) (0.76159,-0.85516) (0.73233,-0.86592) (0.70273,-0.87552) (0.67281,-0.88394) (0.64263,-0.89122) (0.61222,-0.89734) (0.58162,-0.90234) (0.55088,-0.90622) (0.52001,-0.90899) (0.48908,-0.91068) (0.45810,-0.91129) (0.42712,-0.91085) (0.39617,-0.90938) (0.36529,-0.90688) (0.33452,-0.90339) (0.30387,-0.89892) (0.27340,-0.89350) (0.24313,-0.88714) (0.21309,-0.87987) (0.18331,-0.87172) (0.15382,-0.86271) (0.12465,-0.85286) (0.09583,-0.84220) (0.06739,-0.83075) (0.03935,-0.81854) (0.01174,-0.80560) (-0.01542,-0.79196) (-0.04211,-0.77764) (-0.06829,-0.76267) (-0.09396,-0.74707) (-0.11909,-0.73088) (-0.14366,-0.71413) (-0.16764,-0.69685) (-0.19104,-0.67905) (-0.21381,-0.66078) (-0.23596,-0.64205) (-0.25747,-0.62291) (-0.27832,-0.60338) (-0.29849,-0.58348) (-0.31799,-0.56325) (-0.33679,-0.54272) (-0.35489,-0.52191) (-0.37228,-0.50085) (-0.38895,-0.47957) (-0.40489,-0.45810) (-0.42010,-0.43647) (-0.43457,-0.41471) (-0.44829,-0.39283) (-0.46128,-0.37088) (-0.47351,-0.34886) (-0.48500,-0.32682) (-0.49574,-0.30478) (-0.50572,-0.28276) (-0.51497,-0.26079) (-0.52346,-0.23888) (-0.53121,-0.21708) (-0.53822,-0.19539) (-0.54450,-0.17384) (-0.55004,-0.15246) (-0.55487,-0.13126) (-0.55897,-0.11027) (-0.56237,-0.08951) (-0.56506,-0.06899) (-0.56706,-0.04875) (-0.56838,-0.02878) (-0.56902,-0.00912) (-0.56900,0.01022) (-0.56833,0.02922) (-0.56702,0.04786) (-0.56508,0.06614) (-0.56253,0.08404) (-0.55937,0.10153) (-0.55563,0.11862) (-0.55131,0.13528) (-0.54643,0.15150) (-0.54101,0.16728) (-0.53507,0.18260) (-0.52860,0.19745) (-0.52164,0.21183) (-0.51420,0.22572) (-0.50630,0.23912) (-0.49795,0.25202) (-0.48917,0.26441) (-0.47998,0.27629) (-0.47039,0.28765) (-0.46042,0.29849) (-0.45010,0.30881) (-0.43943,0.31860) (-0.42844,0.32786) (-0.41714,0.33659) (-0.40556,0.34478) (-0.39370,0.35244) (-0.38160,0.35957) (-0.36926,0.36616) (-0.35670,0.37223) (-0.34395,0.37776) (-0.33102,0.38277) (-0.31793,0.38726) (-0.30469,0.39123) (-0.29133,0.39468) (-0.27786,0.39762) (-0.26430,0.40006) (-0.25066,0.40199) (-0.23697,0.40344) (-0.22323,0.40439) (-0.20947,0.40487) (-0.19571,0.40487) (-0.18195,0.40441) (-0.16821,0.40350) (-0.15452,0.40214) (-0.14087,0.40034) (-0.12730,0.39811) (-0.11381,0.39546) (-0.10041,0.39240) (-0.08713,0.38895) (-0.07397,0.38511) (-0.06094,0.38089) (-0.04807,0.37631) (-0.03536,0.37137) (-0.02282,0.36609) (-0.01047,0.36048) (0.00169,0.35455) (0.01364,0.34831) (0.02537,0.34178) (0.03687,0.33497) (0.04814,0.32789) (0.05917,0.32055) (0.06994,0.31297) (0.08045,0.30516) (0.09069,0.29713) (0.10065,0.28889) (0.11033,0.28047) (0.11972,0.27186) (0.12881,0.26309) (0.13761,0.25416) (0.14610,0.24509) (0.15427,0.23590) (0.16214,0.22659) (0.16968,0.21717) (0.17690,0.20767) (0.18380,0.19809) (0.19037,0.18844) (0.19661,0.17874) (0.20253,0.16900) (0.20811,0.15922) (0.21335,0.14944) (0.21827,0.13964) (0.22285,0.12985) (0.22710,0.12007) (0.23101,0.11033) (0.23460,0.10062) (0.23786,0.09096) (0.24079,0.08135) (0.24339,0.07182) (0.24567,0.06237) (0.24763,0.05300) (0.24927,0.04373) (0.25061,0.03456) (0.25163,0.02551) (0.25234,0.01659) (0.25276,0.00779) (0.25287,-0.00086) (0.25270,-0.00937) (0.25224,-0.01772) (0.25150,-0.02592) (0.25048,-0.03394) (0.24919,-0.04179) (0.24764,-0.04946) (0.24583,-0.05695) (0.24377,-0.06424) (0.24146,-0.07134) (0.23892,-0.07823) (0.23614,-0.08492) (0.23315,-0.09140) (0.22993,-0.09767) (0.22651,-0.10372) (0.22288,-0.10954) (0.21906,-0.11515) (0.21505,-0.12052) (0.21087,-0.12567) (0.20651,-0.13059) (0.20199,-0.13527) (0.19731,-0.13973) (0.19249,-0.14394) (0.18753,-0.14792) (0.18244,-0.15166) (0.17722,-0.15517) (0.17189,-0.15844) (0.16645,-0.16147) (0.16091,-0.16427) (0.15528,-0.16683) (0.14957,-0.16915) (0.14378,-0.17125) (0.13792,-0.17311) (0.13201,-0.17474) (0.12605,-0.17615) (0.12004,-0.17732) (0.11399,-0.17828) (0.10791,-0.17902) (0.10182,-0.17953) (0.09571,-0.17984) (0.08959,-0.17993) (0.08348,-0.17981) (0.07737,-0.17949) (0.07127,-0.17897) (0.06520,-0.17825) (0.05915,-0.17735) (0.05314,-0.17625) (0.04717,-0.17497) (0.04125,-0.17351) (0.03537,-0.17188) (0.02956,-0.17007) (0.02381,-0.16810) (0.01813,-0.16598) (0.01253,-0.16370) (0.00700,-0.16127) (0.00156,-0.15869) (-0.00379,-0.15598) (-0.00904,-0.15313) (-0.01420,-0.15016) (-0.01925,-0.14706) (-0.02420,-0.14385) (-0.02903,-0.14053) (-0.03375,-0.13710) (-0.03835,-0.13357) (-0.04283,-0.12995) (-0.04719,-0.12624) (-0.05142,-0.12245) (-0.05552,-0.11859) (-0.05948,-0.11465) (-0.06331,-0.11065) (-0.06700,-0.10658) (-0.07056,-0.10247) (-0.07397,-0.09830) (-0.07724,-0.09410) (-0.08037,-0.08985) (-0.08335,-0.08558) (-0.08619,-0.08128) (-0.08888,-0.07696) (-0.09142,-0.07262) (-0.09382,-0.06827) (-0.09606,-0.06392) (-0.09816,-0.05957) (-0.10012,-0.05522) (-0.10192,-0.05088) (-0.10358,-0.04656) (-0.10509,-0.04226) (-0.10645,-0.03798) (-0.10767,-0.03373) (-0.10874,-0.02951) (-0.10968,-0.02533) (-0.11047,-0.02120) (-0.11112,-0.01710) (-0.11163,-0.01306) (-0.11201,-0.00907) (-0.11225,-0.00514) (-0.11236,-0.00126) (-0.11233,0.00255) (-0.11218,0.00629) (-0.11191,0.00996) (-0.11151,0.01356) (-0.11099,0.01708) (-0.11035,0.02052) (-0.10959,0.02388) (-0.10872,0.02716) (-0.10775,0.03035) (-0.10666,0.03346) (-0.10547,0.03647) (-0.10418,0.03939) (-0.10279,0.04221) (-0.10131,0.04494) (-0.09974,0.04757) (-0.09808,0.05011) (-0.09633,0.05254) (-0.09451,0.05487) (-0.09260,0.05710) (-0.09063,0.05923) (-0.08858,0.06125) (-0.08646,0.06317) (-0.08428,0.06498) (-0.08204,0.06669) (-0.07975,0.06829) (-0.07740,0.06979) (-0.07500,0.07118) (-0.07256,0.07247) (-0.07008,0.07365) (-0.06756,0.07473) (-0.06500,0.07571) (-0.06241,0.07658) (-0.05979,0.07735) (-0.05715,0.07801) (-0.05449,0.07858) (-0.05181,0.07905) (-0.04911,0.07942) (-0.04641,0.07969) (-0.04370,0.07986) (-0.04098,0.07994) (-0.03826,0.07993) (-0.03554,0.07983) (-0.03283,0.07964) (-0.03013,0.07935) (-0.02744,0.07899) (-0.02476,0.07853) (-0.02210,0.07800) (-0.01946,0.07739) (-0.01684,0.07669) (-0.01424,0.07592) (-0.01168,0.07508) (-0.00914,0.07417) (-0.00663,0.07318) (-0.00416,0.07213) (-0.00173,0.07101) (0.00067,0.06983) (0.00302,0.06859) (0.00533,0.06730) (0.00759,0.06594) (0.00981,0.06454) (0.01198,0.06308) (0.01410,0.06158) (0.01617,0.06003) (0.01818,0.05844) (0.02014,0.05681) (0.02205,0.05514) (0.02389,0.05344) (0.02568,0.05170) (0.02741,0.04993) (0.02907,0.04814) (0.03068,0.04632) (0.03222,0.04448) (0.03371,0.04262) (0.03512,0.04074) (0.03648,0.03884) (0.03776,0.03694) (0.03899,0.03502) (0.04014,0.03310) (0.04124,0.03117) (0.04226,0.02923) (0.04323,0.02730) (0.04412,0.02537) (0.04495,0.02344) (0.04571,0.02151) (0.04641,0.01960) (0.04705,0.01769) (0.04762,0.01580) (0.04812,0.01392) (0.04856,0.01205) (0.04894,0.01021) (0.04926,0.00838) (0.04951,0.00657) (0.04970,0.00479) (0.04984,0.00303) (0.04991,0.00130) (0.04993,-0.00041) (0.04988,-0.00208) (0.04979,-0.00373) (0.04963,-0.00534) (0.04942,-0.00692) (0.04916,-0.00847) (0.04885,-0.00998) (0.04848,-0.01145) (0.04807,-0.01288) (0.04761,-0.01428) (0.04710,-0.01563) (0.04654,-0.01695) (0.04595,-0.01822) (0.04531,-0.01945) (0.04462,-0.02064) (0.04390,-0.02179) (0.04314,-0.02289) (0.04235,-0.02394) (0.04152,-0.02495) (0.04065,-0.02592) (0.03975,-0.02683) (0.03883,-0.02771) (0.03787,-0.02853) (0.03689,-0.02931) (0.03588,-0.03004) (0.03484,-0.03073) (0.03379,-0.03137) (0.03271,-0.03196) (0.03162,-0.03251) (0.03050,-0.03301) (0.02937,-0.03346) (0.02823,-0.03387) (0.02707,-0.03423) (0.02590,-0.03454) (0.02472,-0.03481) (0.02353,-0.03504) (0.02234,-0.03522) (0.02114,-0.03536) (0.01993,-0.03546) (0.01873,-0.03551) (0.01752,-0.03552) (0.01631,-0.03550) (0.01511,-0.03543) (0.01390,-0.03532) (0.01271,-0.03517) (0.01151,-0.03499) (0.01033,-0.03477) (0.00915,-0.03451) (0.00798,-0.03421) (0.00682,-0.03389) (0.00568,-0.03353) (0.00454,-0.03313) (0.00342,-0.03271) (0.00232,-0.03226) (0.00123,-0.03177) (0.00016,-0.03126) (-0.00089,-0.03072) (-0.00193,-0.03015) (-0.00294,-0.02956) (-0.00394,-0.02895) (-0.00491,-0.02831) (-0.00586,-0.02765) (-0.00679,-0.02697) (-0.00770,-0.02627) (-0.00858,-0.02556) (-0.00944,-0.02482) (-0.01027,-0.02407) (-0.01107,-0.02331) (-0.01185,-0.02253) (-0.01260,-0.02173) (-0.01333,-0.02093) (-0.01403,-0.02012) (-0.01470,-0.01929) (-0.01534,-0.01846) (-0.01595,-0.01762) (-0.01654,-0.01678) (-0.01709,-0.01593) (-0.01762,-0.01507) (-0.01812,-0.01422)};
\fill(0,1.8)circle(1.2pt)node[right]{$P$};
\fill(1.2807560385024006,0)circle(1.2pt)node[above right]{$B$};

\end{tikzpicture}
\caption{初始相点 P 到第一次最大位移 B，再继续向内绕行。图取 $\beta/\omega_0=1/4$。}\label{mp4:fig:l4-h9-s}
\end{figure}

速度为零的条件是 $\tan\omega_dt=\omega_d/\beta$。从 $t=0$ 起，第一个根位于 $(0,\pi/2)$，所以
\[
\varphi=\arctan\frac{\omega_d}{\beta},\qquad
t_B=\frac\varphi{\omega_d}.
\]
因为 $\beta^2+\omega_d^2=\omega_0^2$，该角满足
\[
\sin\varphi=\frac{\omega_d}{\omega_0},\qquad
\cos\varphi=\frac\beta{\omega_0}.
\]
代回位置式得到
\[
x_B=\frac V{\omega_0}\exp\left[-\frac\beta{\omega_d}
\arctan\frac{\omega_d}{\beta}\right].
\]
B 处 $\ddot x=-\omega_0^2x_B<0$，证明这是最大值而非最小值。

以后各次极值的相位为 $\varphi+n\pi$，其中 $n=0,1,2,\ldots$。相应时间与位移为
\[
t_n=\frac{\varphi+n\pi}{\omega_d},\qquad
x_n=(-1)^n\frac V{\omega_0}
\exp\left[-\frac\beta{\omega_d}(\varphi+n\pi)\right].
\]
同侧峰值对应 $n$ 相差 2，间隔为 $2\pi/\omega_d$，后一个与前一个的比为
\[
\frac{x_{n+2}}{x_n}=\ee^{-2\pi\beta/\omega_d}.
\]
相邻异侧极值则相隔 $\pi/\omega_d$，位移异号，大小比为 $\ee^{-\pi\beta/\omega_d}$。图中相轨线逐圈缩小，与同侧峰值按固定比例衰减相符。

在无阻尼极限 $\beta\to0$，有 $t_B\to\pi/(2\omega_0)$、$x_B\to V/\omega_0$，峰值比趋于 1。趋近临界阻尼时，对含 $\omega_d$ 的完整表达式取极限，得到 $t_B\to1/\omega_0$、$x_B\to V/(\mathrm e\omega_0)$，与临界解 $x=Vt\ee^{-\omega_0t}$ 一致。
\end{solution}

\par\medskip\noindent\begin{minipage}{\textwidth}
\begin{exercise}[临界阻尼的越过条件]\label{mp4:exercise:l4-h10}
临界阻尼振子初态为 $x(0)=a>0$、$v(0)=-u$，$u>0$，见图\ref{mp4:fig:l4-h10-q}。求会越过 $x=0$ 的初速条件及越过时间；区分阈值两侧的相轨线。
\end{exercise}

\begin{figure}[H]\centering
\begin{tikzpicture}
\draw[axis](-0.7,0)--(2.3,0)node[right]{$x$};
\draw[axis](0,-2)--(0,1)node[above]{$v$};
\fill(1.5,-1)circle(1.2pt)node[right]{$P$};
\draw[guide](1.5,-1)--(1.5,0);\node[above]at(1.5,0){$a$};\node[left]at(0,-1){$-u$};
\end{tikzpicture}
\caption{临界阻尼振子的初态位于第四象限，$u$ 为初速度大小。}\label{mp4:fig:l4-h10-q}
\end{figure}
\end{minipage}\par\medskip
\begin{solution}
临界阻尼方程的通解为 $x=(A+Bt)\ee^{-\omega_0t}$。初位移给出 $A=a$，初速度给出 $B-\omega_0a=-u$，因此
\[
x=[a+(\omega_0a-u)t]\ee^{-\omega_0t},
\]
\[
v=[-u-\omega_0(\omega_0a-u)t]\ee^{-\omega_0t}.
\]
指数因子恒正，位置是否过零完全由一次因子 $a+(\omega_0a-u)t$ 决定。

若 $u>\omega_0a$，该因子在正时间首次为零，
\[
t_*=\frac a{u-\omega_0a},\qquad
v(t_*)=-(u-\omega_0a)\ee^{-\omega_0t_*}<0.
\]
质点以非零负速度越过原点。若 $u=\omega_0a$，位置为 $a\ee^{-\omega_0t}>0$，只渐近原点；若 $0<u<\omega_0a$，一次因子始终为正，同样不会越过。因此过零条件为严格的 $u>\omega_0a$。

\begin{figure}[H]\centering
\begin{tikzpicture}
\draw[axis](-1,0)--(2.0,0)node[right]{$q$};
\draw[axis](0,-4.1)--(0,1)node[above]{$p$};
\draw[flow] plot coordinates{(1.20000,-0.60000) (1.18125,-0.59971) (1.16252,-0.59888) (1.14383,-0.59752) (1.12518,-0.59569) (1.10660,-0.59340) (1.08810,-0.59068) (1.06969,-0.58758) (1.05138,-0.58410) (1.03319,-0.58028) (1.01512,-0.57615) (0.99718,-0.57172) (0.97939,-0.56701) (0.96174,-0.56206) (0.94426,-0.55687) (0.92694,-0.55147) (0.90980,-0.54588) (0.89283,-0.54011) (0.87604,-0.53417) (0.85944,-0.52809) (0.84304,-0.52188) (0.82683,-0.51555) (0.81082,-0.50912) (0.79501,-0.50259) (0.77940,-0.49598) (0.76401,-0.48931) (0.74882,-0.48258) (0.73385,-0.47579) (0.71909,-0.46897) (0.70454,-0.46212) (0.69020,-0.45524) (0.67609,-0.44835) (0.66218,-0.44146) (0.64850,-0.43456) (0.63502,-0.42767) (0.62177,-0.42079) (0.60872,-0.41393) (0.59589,-0.40710) (0.58328,-0.40029) (0.57088,-0.39352) (0.55868,-0.38678) (0.54670,-0.38009) (0.53493,-0.37344) (0.52336,-0.36684) (0.51200,-0.36030) (0.50084,-0.35381) (0.48989,-0.34737) (0.47913,-0.34100) (0.46857,-0.33470) (0.45821,-0.32845) (0.44804,-0.32228) (0.43807,-0.31617) (0.42828,-0.31014) (0.41868,-0.30417) (0.40927,-0.29828) (0.40004,-0.29247) (0.39099,-0.28673) (0.38212,-0.28106) (0.37342,-0.27548) (0.36490,-0.26997) (0.35655,-0.26454) (0.34837,-0.25919) (0.34035,-0.25391) (0.33250,-0.24872) (0.32480,-0.24360) (0.31727,-0.23857) (0.30989,-0.23361) (0.30267,-0.22873) (0.29560,-0.22394) (0.28867,-0.21922) (0.28189,-0.21458) (0.27526,-0.21001) (0.26877,-0.20553) (0.26241,-0.20112) (0.25620,-0.19679) (0.25011,-0.19253) (0.24416,-0.18835) (0.23834,-0.18425) (0.23265,-0.18022) (0.22708,-0.17626) (0.22163,-0.17238) (0.21630,-0.16857) (0.21109,-0.16483) (0.20600,-0.16116) (0.20102,-0.15756) (0.19615,-0.15403) (0.19139,-0.15056) (0.18674,-0.14717) (0.18219,-0.14384) (0.17775,-0.14057) (0.17341,-0.13738) (0.16916,-0.13424) (0.16502,-0.13117) (0.16097,-0.12816) (0.15701,-0.12521) (0.15314,-0.12232) (0.14936,-0.11949) (0.14567,-0.11672) (0.14207,-0.11400) (0.13854,-0.11135) (0.13511,-0.10874) (0.13175,-0.10620) (0.12847,-0.10370) (0.12527,-0.10126) (0.12214,-0.09887) (0.11909,-0.09654) (0.11610,-0.09425) (0.11319,-0.09201) (0.11035,-0.08982) (0.10758,-0.08768) (0.10487,-0.08559) (0.10223,-0.08354) (0.09965,-0.08153) (0.09713,-0.07957) (0.09468,-0.07766) (0.09228,-0.07578) (0.08994,-0.07395) (0.08766,-0.07216) (0.08543,-0.07041) (0.08326,-0.06870) (0.08114,-0.06703) (0.07907,-0.06539) (0.07705,-0.06379) (0.07508,-0.06223) (0.07316,-0.06071) (0.07129,-0.05922) (0.06946,-0.05776) (0.06768,-0.05634) (0.06594,-0.05495) (0.06424,-0.05359) (0.06259,-0.05226) (0.06097,-0.05097) (0.05940,-0.04970) (0.05787,-0.04847) (0.05637,-0.04726) (0.05491,-0.04608) (0.05349,-0.04493) (0.05210,-0.04381) (0.05075,-0.04271) (0.04943,-0.04164) (0.04815,-0.04060) (0.04690,-0.03958) (0.04568,-0.03858) (0.04449,-0.03761) (0.04333,-0.03666) (0.04219,-0.03573) (0.04109,-0.03483) (0.04002,-0.03395) (0.03897,-0.03309) (0.03795,-0.03225) (0.03695,-0.03143) (0.03598,-0.03063) (0.03504,-0.02985) (0.03412,-0.02909) (0.03322,-0.02834) (0.03235,-0.02762) (0.03149,-0.02691) (0.03066,-0.02622) (0.02986,-0.02555) (0.02907,-0.02490) (0.02830,-0.02426) (0.02755,-0.02363) (0.02682,-0.02302) (0.02611,-0.02243) (0.02542,-0.02185) (0.02475,-0.02129) (0.02409,-0.02074) (0.02345,-0.02020) (0.02283,-0.01968) (0.02222,-0.01917) (0.02163,-0.01867) (0.02105,-0.01819) (0.02049,-0.01771) (0.01995,-0.01725) (0.01941,-0.01680) (0.01890,-0.01637) (0.01839,-0.01594) (0.01790,-0.01552) (0.01742,-0.01512) (0.01695,-0.01472) (0.01650,-0.01434) (0.01606,-0.01396) (0.01563,-0.01359) (0.01521,-0.01324) (0.01480,-0.01289) (0.01440,-0.01255) (0.01402,-0.01222) (0.01364,-0.01190) (0.01327,-0.01159) (0.01291,-0.01128) (0.01257,-0.01098) (0.01223,-0.01069) (0.01190,-0.01041) (0.01158,-0.01014) (0.01126,-0.00987) (0.01096,-0.00961) (0.01066,-0.00935) (0.01038,-0.00910) (0.01009,-0.00886) (0.00982,-0.00863) (0.00956,-0.00840) (0.00930,-0.00817) (0.00904,-0.00796) (0.00880,-0.00774) (0.00856,-0.00754) (0.00833,-0.00734) (0.00810,-0.00714) (0.00788,-0.00695) (0.00767,-0.00677) (0.00746,-0.00658) (0.00726,-0.00641) (0.00706,-0.00624) (0.00687,-0.00607) (0.00668,-0.00591) (0.00650,-0.00575) (0.00632,-0.00559) (0.00615,-0.00544) (0.00598,-0.00530) (0.00582,-0.00516) (0.00566,-0.00502) (0.00550,-0.00488) (0.00535,-0.00475) (0.00521,-0.00462) (0.00506,-0.00450) (0.00492,-0.00438) (0.00479,-0.00426) (0.00466,-0.00414) (0.00453,-0.00403) (0.00441,-0.00392) (0.00428,-0.00382) (0.00417,-0.00371) (0.00405,-0.00361) (0.00394,-0.00352) (0.00383,-0.00342) (0.00373,-0.00333) (0.00363,-0.00324) (0.00353,-0.00315) (0.00343,-0.00306) (0.00333,-0.00298) (0.00324,-0.00290) (0.00315,-0.00282) (0.00307,-0.00274) (0.00298,-0.00267) (0.00290,-0.00260) (0.00282,-0.00253) (0.00274,-0.00246) (0.00267,-0.00239) (0.00259,-0.00232) (0.00252,-0.00226) (0.00245,-0.00220) (0.00238,-0.00214) (0.00232,-0.00208) (0.00225,-0.00202) (0.00219,-0.00197) (0.00213,-0.00191) (0.00207,-0.00186) (0.00201,-0.00181) (0.00196,-0.00176) (0.00190,-0.00171) (0.00185,-0.00167) (0.00180,-0.00162) (0.00175,-0.00158) (0.00170,-0.00153) (0.00165,-0.00149) (0.00161,-0.00145) (0.00156,-0.00141) (0.00152,-0.00137) (0.00148,-0.00133) (0.00144,-0.00130) (0.00140,-0.00126) (0.00136,-0.00123) (0.00132,-0.00119) (0.00128,-0.00116) (0.00125,-0.00113) (0.00121,-0.00110) (0.00118,-0.00107) (0.00114,-0.00104) (0.00111,-0.00101) (0.00108,-0.00098) (0.00105,-0.00095) (0.00102,-0.00093) (0.00099,-0.00090) (0.00097,-0.00088) (0.00094,-0.00085) (0.00091,-0.00083) (0.00089,-0.00081) (0.00086,-0.00078) (0.00084,-0.00076) (0.00081,-0.00074) (0.00079,-0.00072) (0.00077,-0.00070) (0.00075,-0.00068) (0.00073,-0.00066) (0.00071,-0.00064) (0.00069,-0.00063) (0.00067,-0.00061) (0.00065,-0.00059) (0.00063,-0.00057) (0.00061,-0.00056) (0.00060,-0.00054) (0.00058,-0.00053) (0.00056,-0.00051) (0.00055,-0.00050) (0.00053,-0.00049) (0.00052,-0.00047) (0.00050,-0.00046) (0.00049,-0.00045) (0.00047,-0.00043) (0.00046,-0.00042) (0.00045,-0.00041) (0.00044,-0.00040) (0.00042,-0.00039) (0.00041,-0.00038) (0.00040,-0.00037) (0.00039,-0.00036) (0.00038,-0.00035) (0.00037,-0.00034) (0.00036,-0.00033) (0.00035,-0.00032) (0.00034,-0.00031) (0.00033,-0.00030)};
\draw[flow,line width=1pt] plot coordinates{(1.20000,-1.20000) (1.16308,-1.16308) (1.12730,-1.12730) (1.09261,-1.09261) (1.05900,-1.05900) (1.02641,-1.02641) (0.99483,-0.99483) (0.96423,-0.96423) (0.93456,-0.93456) (0.90581,-0.90581) (0.87794,-0.87794) (0.85093,-0.85093) (0.82475,-0.82475) (0.79937,-0.79937) (0.77478,-0.77478) (0.75094,-0.75094) (0.72784,-0.72784) (0.70544,-0.70544) (0.68374,-0.68374) (0.66270,-0.66270) (0.64231,-0.64231) (0.62255,-0.62255) (0.60340,-0.60340) (0.58483,-0.58483) (0.56684,-0.56684) (0.54940,-0.54940) (0.53250,-0.53250) (0.51611,-0.51611) (0.50023,-0.50023) (0.48484,-0.48484) (0.46993,-0.46993) (0.45547,-0.45547) (0.44146,-0.44146) (0.42787,-0.42787) (0.41471,-0.41471) (0.40195,-0.40195) (0.38958,-0.38958) (0.37760,-0.37760) (0.36598,-0.36598) (0.35472,-0.35472) (0.34381,-0.34381) (0.33323,-0.33323) (0.32298,-0.32298) (0.31304,-0.31304) (0.30341,-0.30341) (0.29407,-0.29407) (0.28502,-0.28502) (0.27626,-0.27626) (0.26776,-0.26776) (0.25952,-0.25952) (0.25153,-0.25153) (0.24379,-0.24379) (0.23629,-0.23629) (0.22902,-0.22902) (0.22198,-0.22198) (0.21515,-0.21515) (0.20853,-0.20853) (0.20211,-0.20211) (0.19589,-0.19589) (0.18987,-0.18987) (0.18403,-0.18403) (0.17836,-0.17836) (0.17288,-0.17288) (0.16756,-0.16756) (0.16240,-0.16240) (0.15741,-0.15741) (0.15256,-0.15256) (0.14787,-0.14787) (0.14332,-0.14332) (0.13891,-0.13891) (0.13464,-0.13464) (0.13049,-0.13049) (0.12648,-0.12648) (0.12259,-0.12259) (0.11882,-0.11882) (0.11516,-0.11516) (0.11162,-0.11162) (0.10818,-0.10818) (0.10485,-0.10485) (0.10163,-0.10163) (0.09850,-0.09850) (0.09547,-0.09547) (0.09253,-0.09253) (0.08969,-0.08969) (0.08693,-0.08693) (0.08425,-0.08425) (0.08166,-0.08166) (0.07915,-0.07915) (0.07671,-0.07671) (0.07435,-0.07435) (0.07207,-0.07207) (0.06985,-0.06985) (0.06770,-0.06770) (0.06562,-0.06562) (0.06360,-0.06360) (0.06164,-0.06164) (0.05974,-0.05974) (0.05791,-0.05791) (0.05612,-0.05612) (0.05440,-0.05440) (0.05272,-0.05272) (0.05110,-0.05110) (0.04953,-0.04953) (0.04801,-0.04801) (0.04653,-0.04653) (0.04510,-0.04510) (0.04371,-0.04371) (0.04237,-0.04237) (0.04106,-0.04106) (0.03980,-0.03980) (0.03857,-0.03857) (0.03739,-0.03739) (0.03624,-0.03624) (0.03512,-0.03512) (0.03404,-0.03404) (0.03299,-0.03299) (0.03198,-0.03198) (0.03100,-0.03100) (0.03004,-0.03004) (0.02912,-0.02912) (0.02822,-0.02822) (0.02735,-0.02735) (0.02651,-0.02651) (0.02570,-0.02570) (0.02491,-0.02491) (0.02414,-0.02414) (0.02340,-0.02340) (0.02268,-0.02268) (0.02198,-0.02198) (0.02130,-0.02130) (0.02065,-0.02065) (0.02001,-0.02001) (0.01940,-0.01940) (0.01880,-0.01880) (0.01822,-0.01822) (0.01766,-0.01766) (0.01712,-0.01712) (0.01659,-0.01659) (0.01608,-0.01608) (0.01559,-0.01559) (0.01511,-0.01511) (0.01464,-0.01464) (0.01419,-0.01419) (0.01375,-0.01375) (0.01333,-0.01333) (0.01292,-0.01292) (0.01252,-0.01252) (0.01214,-0.01214) (0.01176,-0.01176) (0.01140,-0.01140) (0.01105,-0.01105) (0.01071,-0.01071) (0.01038,-0.01038) (0.01006,-0.01006) (0.00975,-0.00975) (0.00945,-0.00945) (0.00916,-0.00916) (0.00888,-0.00888) (0.00861,-0.00861) (0.00834,-0.00834) (0.00809,-0.00809) (0.00784,-0.00784) (0.00760,-0.00760) (0.00736,-0.00736) (0.00714,-0.00714) (0.00692,-0.00692) (0.00670,-0.00670) (0.00650,-0.00650) (0.00630,-0.00630) (0.00610,-0.00610) (0.00592,-0.00592) (0.00573,-0.00573) (0.00556,-0.00556) (0.00539,-0.00539) (0.00522,-0.00522) (0.00506,-0.00506) (0.00490,-0.00490) (0.00475,-0.00475) (0.00461,-0.00461) (0.00447,-0.00447) (0.00433,-0.00433) (0.00419,-0.00419) (0.00407,-0.00407) (0.00394,-0.00394) (0.00382,-0.00382) (0.00370,-0.00370) (0.00359,-0.00359) (0.00348,-0.00348) (0.00337,-0.00337) (0.00327,-0.00327) (0.00317,-0.00317) (0.00307,-0.00307) (0.00297,-0.00297) (0.00288,-0.00288) (0.00279,-0.00279) (0.00271,-0.00271) (0.00262,-0.00262) (0.00254,-0.00254) (0.00247,-0.00247) (0.00239,-0.00239) (0.00232,-0.00232) (0.00225,-0.00225) (0.00218,-0.00218) (0.00211,-0.00211) (0.00204,-0.00204) (0.00198,-0.00198) (0.00192,-0.00192) (0.00186,-0.00186) (0.00180,-0.00180) (0.00175,-0.00175) (0.00169,-0.00169) (0.00164,-0.00164) (0.00159,-0.00159) (0.00154,-0.00154) (0.00150,-0.00150) (0.00145,-0.00145) (0.00141,-0.00141) (0.00136,-0.00136) (0.00132,-0.00132) (0.00128,-0.00128) (0.00124,-0.00124) (0.00120,-0.00120) (0.00116,-0.00116) (0.00113,-0.00113) (0.00109,-0.00109) (0.00106,-0.00106) (0.00103,-0.00103) (0.00100,-0.00100) (0.00097,-0.00097) (0.00094,-0.00094) (0.00091,-0.00091) (0.00088,-0.00088) (0.00085,-0.00085) (0.00083,-0.00083) (0.00080,-0.00080) (0.00078,-0.00078) (0.00075,-0.00075) (0.00073,-0.00073) (0.00071,-0.00071) (0.00068,-0.00068) (0.00066,-0.00066) (0.00064,-0.00064) (0.00062,-0.00062) (0.00060,-0.00060) (0.00059,-0.00059) (0.00057,-0.00057) (0.00055,-0.00055) (0.00053,-0.00053) (0.00052,-0.00052) (0.00050,-0.00050) (0.00049,-0.00049) (0.00047,-0.00047) (0.00046,-0.00046) (0.00044,-0.00044) (0.00043,-0.00043) (0.00042,-0.00042) (0.00040,-0.00040) (0.00039,-0.00039) (0.00038,-0.00038) (0.00037,-0.00037) (0.00036,-0.00036) (0.00034,-0.00034) (0.00033,-0.00033) (0.00032,-0.00032) (0.00031,-0.00031) (0.00030,-0.00030) (0.00029,-0.00029) (0.00029,-0.00029) (0.00028,-0.00028) (0.00027,-0.00027) (0.00026,-0.00026) (0.00025,-0.00025) (0.00024,-0.00024) (0.00024,-0.00024) (0.00023,-0.00023) (0.00022,-0.00022) (0.00022,-0.00022) (0.00021,-0.00021) (0.00020,-0.00020) (0.00020,-0.00020) (0.00019,-0.00019) (0.00018,-0.00018) (0.00018,-0.00018) (0.00017,-0.00017) (0.00017,-0.00017) (0.00016,-0.00016) (0.00016,-0.00016) (0.00015,-0.00015) (0.00015,-0.00015) (0.00014,-0.00014) (0.00014,-0.00014) (0.00013,-0.00013) (0.00013,-0.00013) (0.00013,-0.00013) (0.00012,-0.00012) (0.00012,-0.00012) (0.00012,-0.00012) (0.00011,-0.00011) (0.00011,-0.00011) (0.00011,-0.00011) (0.00010,-0.00010) (0.00010,-0.00010) (0.00010,-0.00010) (0.00009,-0.00009) (0.00009,-0.00009) (0.00009,-0.00009) (0.00008,-0.00008) (0.00008,-0.00008) (0.00008,-0.00008) (0.00008,-0.00008) (0.00007,-0.00007) (0.00007,-0.00007) (0.00007,-0.00007) (0.00007,-0.00007) (0.00007,-0.00007) (0.00006,-0.00006) (0.00006,-0.00006) (0.00006,-0.00006) (0.00006,-0.00006) (0.00006,-0.00006) (0.00005,-0.00005)};
\draw[flow] plot coordinates{(1.20000,-3.60000) (1.09039,-3.41655) (0.98638,-3.24098) (0.88775,-3.07297) (0.79425,-2.91224) (0.70566,-2.75849) (0.62177,-2.61144) (0.54238,-2.47083) (0.46728,-2.33640) (0.39629,-2.20791) (0.32923,-2.08510) (0.26591,-1.96777) (0.20619,-1.85568) (0.14988,-1.74863) (0.09685,-1.64640) (0.04693,-1.54882) (0.00000,-1.45567) (-0.04409,-1.36680) (-0.08547,-1.28201) (-0.12426,-1.20115) (-0.16058,-1.12405) (-0.19455,-1.05056) (-0.22627,-0.98052) (-0.25586,-0.91380) (-0.28342,-0.85026) (-0.30904,-0.78976) (-0.33281,-0.73218) (-0.35483,-0.67740) (-0.37518,-0.62529) (-0.39394,-0.57575) (-0.41119,-0.52867) (-0.42700,-0.48394) (-0.44146,-0.44146) (-0.45462,-0.40113) (-0.46655,-0.36287) (-0.47732,-0.32658) (-0.48698,-0.29219) (-0.49560,-0.25960) (-0.50322,-0.22874) (-0.50991,-0.19953) (-0.51571,-0.17190) (-0.52067,-0.14579) (-0.52484,-0.12112) (-0.52825,-0.09782) (-0.53096,-0.07585) (-0.53301,-0.05514) (-0.53442,-0.03563) (-0.53525,-0.01727) (-0.53551,0.00000) (-0.53526,0.01622) (-0.53451,0.03144) (-0.53330,0.04571) (-0.53166,0.05907) (-0.52962,0.07157) (-0.52720,0.08324) (-0.52442,0.09413) (-0.52132,0.10426) (-0.51791,0.11369) (-0.51422,0.12243) (-0.51027,0.13053) (-0.50607,0.13802) (-0.50165,0.14492) (-0.49702,0.15127) (-0.49220,0.15709) (-0.48721,0.16240) (-0.48206,0.16724) (-0.47676,0.17163) (-0.47133,0.17559) (-0.46579,0.17915) (-0.46014,0.18232) (-0.45440,0.18512) (-0.44857,0.18759) (-0.44268,0.18972) (-0.43672,0.19154) (-0.43071,0.19308) (-0.42465,0.19433) (-0.41857,0.19533) (-0.41245,0.19608) (-0.40631,0.19660) (-0.40016,0.19691) (-0.39401,0.19700) (-0.38785,0.19691) (-0.38170,0.19663) (-0.37556,0.19619) (-0.36944,0.19559) (-0.36334,0.19484) (-0.35727,0.19394) (-0.35122,0.19292) (-0.34521,0.19178) (-0.33924,0.19053) (-0.33330,0.18917) (-0.32741,0.18772) (-0.32157,0.18617) (-0.31578,0.18455) (-0.31004,0.18284) (-0.30435,0.18107) (-0.29872,0.17923) (-0.29315,0.17734) (-0.28764,0.17539) (-0.28219,0.17339) (-0.27680,0.17135) (-0.27148,0.16928) (-0.26622,0.16716) (-0.26103,0.16502) (-0.25591,0.16285) (-0.25085,0.16066) (-0.24587,0.15845) (-0.24095,0.15622) (-0.23611,0.15398) (-0.23133,0.15173) (-0.22662,0.14947) (-0.22199,0.14721) (-0.21742,0.14495) (-0.21293,0.14268) (-0.20850,0.14042) (-0.20415,0.13816) (-0.19987,0.13591) (-0.19566,0.13367) (-0.19151,0.13143) (-0.18744,0.12921) (-0.18344,0.12700) (-0.17950,0.12480) (-0.17564,0.12262) (-0.17184,0.12045) (-0.16811,0.11830) (-0.16445,0.11617) (-0.16085,0.11406) (-0.15732,0.11196) (-0.15385,0.10989) (-0.15045,0.10784) (-0.14711,0.10582) (-0.14384,0.10381) (-0.14062,0.10183) (-0.13747,0.09987) (-0.13438,0.09794) (-0.13135,0.09603) (-0.12838,0.09414) (-0.12547,0.09228) (-0.12261,0.09045) (-0.11981,0.08864) (-0.11707,0.08686) (-0.11438,0.08510) (-0.11175,0.08337) (-0.10917,0.08166) (-0.10665,0.07998) (-0.10417,0.07833) (-0.10175,0.07670) (-0.09938,0.07510) (-0.09706,0.07353) (-0.09478,0.07198) (-0.09256,0.07045) (-0.09038,0.06896) (-0.08825,0.06748) (-0.08616,0.06604) (-0.08412,0.06461) (-0.08212,0.06322) (-0.08017,0.06184) (-0.07826,0.06050) (-0.07639,0.05917) (-0.07456,0.05787) (-0.07277,0.05660) (-0.07102,0.05535) (-0.06931,0.05412) (-0.06764,0.05291) (-0.06600,0.05173) (-0.06440,0.05057) (-0.06284,0.04944) (-0.06131,0.04832) (-0.05982,0.04723) (-0.05836,0.04616) (-0.05694,0.04511) (-0.05554,0.04408) (-0.05418,0.04307) (-0.05285,0.04208) (-0.05155,0.04111) (-0.05028,0.04016) (-0.04904,0.03923) (-0.04783,0.03832) (-0.04665,0.03743) (-0.04549,0.03656) (-0.04436,0.03570) (-0.04326,0.03487) (-0.04218,0.03405) (-0.04113,0.03325) (-0.04010,0.03246) (-0.03910,0.03170) (-0.03812,0.03095) (-0.03717,0.03021) (-0.03623,0.02949) (-0.03532,0.02879) (-0.03443,0.02810) (-0.03357,0.02743) (-0.03272,0.02677) (-0.03189,0.02613) (-0.03109,0.02550) (-0.03030,0.02488) (-0.02953,0.02428) (-0.02878,0.02369) (-0.02805,0.02312) (-0.02734,0.02256) (-0.02664,0.02201) (-0.02596,0.02147) (-0.02530,0.02095) (-0.02465,0.02043) (-0.02402,0.01993) (-0.02341,0.01944) (-0.02281,0.01896) (-0.02222,0.01850) (-0.02165,0.01804) (-0.02109,0.01760) (-0.02055,0.01716) (-0.02002,0.01673) (-0.01950,0.01632) (-0.01900,0.01591) (-0.01851,0.01552) (-0.01803,0.01513) (-0.01756,0.01475) (-0.01711,0.01438) (-0.01666,0.01402) (-0.01623,0.01367) (-0.01581,0.01333) (-0.01540,0.01299) (-0.01500,0.01267) (-0.01461,0.01235) (-0.01423,0.01204) (-0.01385,0.01173) (-0.01349,0.01144) (-0.01314,0.01115) (-0.01280,0.01086) (-0.01246,0.01059) (-0.01213,0.01032) (-0.01182,0.01006) (-0.01150,0.00980) (-0.01120,0.00955) (-0.01091,0.00931) (-0.01062,0.00907) (-0.01034,0.00884) (-0.01007,0.00861) (-0.00980,0.00839) (-0.00954,0.00817) (-0.00929,0.00796) (-0.00905,0.00776) (-0.00881,0.00756) (-0.00857,0.00736) (-0.00835,0.00718) (-0.00813,0.00699) (-0.00791,0.00681) (-0.00770,0.00663) (-0.00749,0.00646) (-0.00730,0.00629) (-0.00710,0.00613) (-0.00691,0.00597) (-0.00673,0.00582) (-0.00655,0.00566) (-0.00637,0.00552) (-0.00620,0.00537) (-0.00604,0.00523) (-0.00588,0.00510) (-0.00572,0.00496) (-0.00557,0.00483) (-0.00542,0.00471) (-0.00527,0.00458) (-0.00513,0.00446) (-0.00499,0.00435) (-0.00486,0.00423) (-0.00473,0.00412) (-0.00460,0.00401) (-0.00448,0.00391) (-0.00436,0.00380) (-0.00424,0.00370) (-0.00413,0.00361) (-0.00401,0.00351) (-0.00391,0.00342) (-0.00380,0.00333) (-0.00370,0.00324) (-0.00360,0.00315) (-0.00350,0.00307) (-0.00341,0.00299) (-0.00331,0.00291) (-0.00322,0.00283) (-0.00314,0.00276) (-0.00305,0.00268) (-0.00297,0.00261) (-0.00289,0.00254) (-0.00281,0.00248) (-0.00273,0.00241) (-0.00266,0.00234) (-0.00259,0.00228) (-0.00252,0.00222) (-0.00245,0.00216) (-0.00238,0.00210) (-0.00232,0.00205) (-0.00225,0.00199) (-0.00219,0.00194) (-0.00213,0.00189) (-0.00207,0.00184) (-0.00202,0.00179) (-0.00196,0.00174) (-0.00191,0.00169) (-0.00186,0.00165) (-0.00181,0.00160) (-0.00176,0.00156) (-0.00171,0.00152) (-0.00166,0.00148) (-0.00162,0.00144) (-0.00157,0.00140) (-0.00153,0.00136) (-0.00149,0.00132) (-0.00145,0.00129) (-0.00141,0.00125) (-0.00137,0.00122) (-0.00133,0.00119) (-0.00129,0.00115) (-0.00126,0.00112) (-0.00122,0.00109) (-0.00119,0.00106) (-0.00116,0.00103) (-0.00113,0.00101) (-0.00109,0.00098) (-0.00106,0.00095) (-0.00104,0.00093)};
\draw(1.2,.04)--(1.2,-.04)node[above,yshift=3pt]{$1$};
\draw(-.04,-1.2)--(.04,-1.2);\node[right]at(.06,-1.2){$-1$};
\node[right]at(1.2,-.6){$\eta=1/2$};
\node[right]at(1.2,-1.2){$\eta=1$};
\node[right]at(1.2,-3.6){$\eta=3$};
\end{tikzpicture}
\caption{临界阻尼相图。曲线旁标出初速比 $\eta=u/(\omega_0a)$，粗线 $\eta=1$ 对应临界初态。}\label{mp4:fig:l4-h10-s}
\end{figure}

在图\ref{mp4:fig:l4-h10-s}中，$q=x/a$、$p=v/(\omega_0a)$，所有初态位于 $q=1$ 的竖线上。阈值初态为 $(1,-1)$，它沿直线 $p=-q$ 渐近原点。初态在该点下方时，反向速度更大，轨线穿过纵轴进入第三象限；在该点上方时始终留在 $q>0$ 一侧。

对越过情形，令 $c=u-\omega_0a>0$。速度可写为 $v=(-u+\omega_0ct)\ee^{-\omega_0t}$，首次负向最大位移发生在
\[
t_{\min}=\frac u{\omega_0c}=t_*+\frac1{\omega_0},
\qquad
x_{\min}=-\frac c{\omega_0}\exp\left(-\frac u c\right)<0.
\]
此后 $v>0$、$x<0$，从负侧趋近原点。位置中的一次因子只有一个零点，因此不会再次穿回正侧。

当 $u\downarrow\omega_0a$ 时，$t_*\to\infty$，$x_{\min}\to0^-$，可越过轨线逐渐贴近阈值直线。图中的阈值直线将不同初态分开；临界阻尼系统是否过零，仍须结合初速度判断。
令 $w=v+\omega_0x$，用原运动方程逐项计算，
\[
\dot w=\ddot x+\omega_0\dot x
=-2\omega_0\dot x-\omega_0^2x+\omega_0\dot x
=-\omega_0(v+\omega_0x)=-\omega_0w.
\]
乘积分因子 $\ee^{\omega_0t}$，得到 $\dd(\ee^{\omega_0t}w)/\dd t=0$。由 $w(0)=\omega_0a-u$ 确定常量，故
\[
 w(t)=(\omega_0a-u)\ee^{-\omega_0t}.
\]
由于 $w$ 的符号保持不变，$w=0$ 即 $v=-\omega_0x$ 是一条不变直线。图\ref{mp4:fig:l4-h10-s}中的阈值初态沿此线趋近原点，两侧的轨线均不能穿过它。
\end{solution}

\par\medskip\noindent\begin{minipage}{\textwidth}
\begin{exercise}[摩擦振子的残余位移]\label{mp4:exercise:l4-h11}
沿用例\ref{mp4:example:l4-e9}的摩擦模型，从 $x_0=9d/2$ 静止释放，图\ref{mp4:fig:l4-h11-q}给出初态。求各转向位置、总滑动时间和生热。再对一般 $x_0=A>d$ 给出滑动半周次数的公式。
\end{exercise}

\begin{figure}[H]\centering
\begin{tikzpicture}
\draw[axis](-1.8,0)--(3.5,0)node[right]{$x$};
\draw[axis](0,-0.7)--(0,1.1)node[above]{$v$};
\fill[gray!25](-.6,-.1)rectangle(.6,.1);\node[below]at(-.6,0){$-d$};\node[below]at(.6,0){$d$};\fill(2.7,0)circle(1.5pt)node[above]{$A_0:9d/2$};
\end{tikzpicture}
\caption{摩擦振子的初态。灰带为静止允许区间 $-d\le x\le d$。}\label{mp4:fig:l4-h11-q}
\end{figure}
\end{minipage}\par\medskip
\begin{solution}
本题 $x_0=9d/2>d$，初始弹力大于静摩擦上限，开始向左滑动。沿用例\ref{mp4:example:l4-e9}的移心方法，第一段中心为 $d$，相对初位移为 $7d/2$，所以
\[
x=d+\frac{7d}{2}\cos\omega_0t,
\quad v=-\frac{7\omega_0d}{2}\sin\omega_0t.
\]
在 $t_1=\pi/\omega_0$ 处第一次转向，$x_1=-5d/2$。所需静摩擦力为 $5f/2>f$，物体继续向右启动。

第二段中心改为 $-d$。令 $s=t-t_1$，新初值相对于该中心为 $-3d/2$，于是
\[
x=-d-\frac{3d}{2}\cos\omega_0s,
\quad v=\frac{3\omega_0d}{2}\sin\omega_0s.
\]
再经过 $\pi/\omega_0$ 到达 $x_2=d/2$、$v=0$。此处弹力为 $f/2$，可由静摩擦平衡。图\ref{mp4:fig:l4-h11-s}中第二段相弧在静止允许带内终止。

\begin{figure}[H]\centering
\begin{tikzpicture}
\begin{scope}[shift={(0,0)},scale=0.72]
\draw[axis](-3.2,0)--(5.2,0)node[right]{$x/d$};
\draw[axis](0,-4.1)--(0,2.3)node[above]{$v/(\omega_0d)$};
\fill[gray!25](-1,-.12)rectangle(1,.12);\draw[flow] plot coordinates{(4.50000,-0.00000) (4.49827,-0.10994) (4.49309,-0.21977) (4.48447,-0.32938) (4.47240,-0.43867) (4.45691,-0.54752) (4.43801,-0.65583) (4.41571,-0.76350) (4.39004,-0.87041) (4.36103,-0.97647) (4.32870,-1.08156) (4.29308,-1.18558) (4.25422,-1.28844) (4.21214,-1.39002) (4.16689,-1.49023) (4.11852,-1.58897) (4.06707,-1.68614) (4.01260,-1.78164) (3.95515,-1.87539) (3.89478,-1.96729) (3.83156,-2.05725) (3.76554,-2.14517) (3.69680,-2.23098) (3.62539,-2.31459) (3.55139,-2.39591) (3.47487,-2.47487) (3.39591,-2.55139) (3.31459,-2.62539) (3.23098,-2.69680) (3.14517,-2.76554) (3.05725,-2.83156) (2.96729,-2.89478) (2.87539,-2.95515) (2.78164,-3.01260) (2.68614,-3.06707) (2.58897,-3.11852) (2.49023,-3.16689) (2.39002,-3.21214) (2.28844,-3.25422) (2.18558,-3.29308) (2.08156,-3.32870) (1.97647,-3.36103) (1.87041,-3.39004) (1.76350,-3.41571) (1.65583,-3.43801) (1.54752,-3.45691) (1.43867,-3.47240) (1.32938,-3.48447) (1.21977,-3.49309) (1.10994,-3.49827) (1.00000,-3.50000) (0.89006,-3.49827) (0.78023,-3.49309) (0.67062,-3.48447) (0.56133,-3.47240) (0.45248,-3.45691) (0.34417,-3.43801) (0.23650,-3.41571) (0.12959,-3.39004) (0.02353,-3.36103) (-0.08156,-3.32870) (-0.18558,-3.29308) (-0.28844,-3.25422) (-0.39002,-3.21214) (-0.49023,-3.16689) (-0.58897,-3.11852) (-0.68614,-3.06707) (-0.78164,-3.01260) (-0.87539,-2.95515) (-0.96729,-2.89478) (-1.05725,-2.83156) (-1.14517,-2.76554) (-1.23098,-2.69680) (-1.31459,-2.62539) (-1.39591,-2.55139) (-1.47487,-2.47487) (-1.55139,-2.39591) (-1.62539,-2.31459) (-1.69680,-2.23098) (-1.76554,-2.14517) (-1.83156,-2.05725) (-1.89478,-1.96729) (-1.95515,-1.87539) (-2.01260,-1.78164) (-2.06707,-1.68614) (-2.11852,-1.58897) (-2.16689,-1.49023) (-2.21214,-1.39002) (-2.25422,-1.28844) (-2.29308,-1.18558) (-2.32870,-1.08156) (-2.36103,-0.97647) (-2.39004,-0.87041) (-2.41571,-0.76350) (-2.43801,-0.65583) (-2.45691,-0.54752) (-2.47240,-0.43867) (-2.48447,-0.32938) (-2.49309,-0.21977) (-2.49827,-0.10994) (-2.50000,-0.00000)};
\draw[flow] plot coordinates{(-2.50000,0.00000) (-2.49926,0.04712) (-2.49704,0.09419) (-2.49334,0.14116) (-2.48817,0.18800) (-2.48153,0.23465) (-2.47343,0.28107) (-2.46388,0.32721) (-2.45287,0.37303) (-2.44044,0.41849) (-2.42658,0.46353) (-2.41132,0.50811) (-2.39466,0.55219) (-2.37663,0.59572) (-2.35724,0.63867) (-2.33651,0.68099) (-2.31446,0.72263) (-2.29111,0.76356) (-2.26649,0.80374) (-2.24062,0.84313) (-2.21353,0.88168) (-2.18523,0.91936) (-2.15577,0.95614) (-2.12517,0.99197) (-2.09345,1.02682) (-2.06066,1.06066) (-2.02682,1.09345) (-1.99197,1.12517) (-1.95614,1.15577) (-1.91936,1.18523) (-1.88168,1.21353) (-1.84313,1.24062) (-1.80374,1.26649) (-1.76356,1.29111) (-1.72263,1.31446) (-1.68099,1.33651) (-1.63867,1.35724) (-1.59572,1.37663) (-1.55219,1.39466) (-1.50811,1.41132) (-1.46353,1.42658) (-1.41849,1.44044) (-1.37303,1.45287) (-1.32721,1.46388) (-1.28107,1.47343) (-1.23465,1.48153) (-1.18800,1.48817) (-1.14116,1.49334) (-1.09419,1.49704) (-1.04712,1.49926) (-1.00000,1.50000) (-0.95288,1.49926) (-0.90581,1.49704) (-0.85884,1.49334) (-0.81200,1.48817) (-0.76535,1.48153) (-0.71893,1.47343) (-0.67279,1.46388) (-0.62697,1.45287) (-0.58151,1.44044) (-0.53647,1.42658) (-0.49189,1.41132) (-0.44781,1.39466) (-0.40428,1.37663) (-0.36133,1.35724) (-0.31901,1.33651) (-0.27737,1.31446) (-0.23644,1.29111) (-0.19626,1.26649) (-0.15687,1.24062) (-0.11832,1.21353) (-0.08064,1.18523) (-0.04386,1.15577) (-0.00803,1.12517) (0.02682,1.09345) (0.06066,1.06066) (0.09345,1.02682) (0.12517,0.99197) (0.15577,0.95614) (0.18523,0.91936) (0.21353,0.88168) (0.24062,0.84313) (0.26649,0.80374) (0.29111,0.76356) (0.31446,0.72263) (0.33651,0.68099) (0.35724,0.63867) (0.37663,0.59572) (0.39466,0.55219) (0.41132,0.50811) (0.42658,0.46353) (0.44044,0.41849) (0.45287,0.37303) (0.46388,0.32721) (0.47343,0.28107) (0.48153,0.23465) (0.48817,0.18800) (0.49334,0.14116) (0.49704,0.09419) (0.49926,0.04712) (0.50000,0.00000)};
\fill(4.5,0)circle(1.2pt)node[above]{$A_0$};
\fill(-2.5,0)circle(1.2pt)node[above left]{$A_1$};
\fill(0.5,0)circle(1.2pt)node[above right]{$A_2$};
\node[below]at(-1,0){$-1$};\node[below]at(1,0){$1$};\end{scope}

\end{tikzpicture}
\caption{两次滑动后停在 $A_2$。灰带内的停点可以不在区间端点，也可以不在原点。}\label{mp4:fig:l4-h11-s}
\end{figure}

故总滑动时间与生热分别为
\[
t_{\rm stop}=\frac{2\pi}{\omega_0},\qquad
Q=f\left(7d+3d\right)=10kd^2.
\]
也可用初、末弹性势能相减验证：$Q=k[(9d/2)^2-(d/2)^2]/2=10kd^2$。

对一般初位移 $x_0=A>d$，设第 $n$ 段开始前仍未停止。奇数段向左，中心为 $d$；偶数段向右，中心为 $-d$。一个滑动半周的两个端点关于该段中心对称，因此
\[
x_n=2(-1)^{n-1}d-x_{n-1}.
\]
令 $y_n=(-1)^nx_n$，递推式成为 $y_n=y_{n-1}-2d$。从 $y_0=A$ 逐次相加，$y_n=A-2nd$，故
\[
x_n=(-1)^n(A-2nd).
\]
这一公式只用于实际发生的滑动段；每到一个速度为零的端点，必须先判断能否由静摩擦维持静止。

停止条件为 $|x_N|\le d$，即 $-d\le A-2Nd\le d$。满足它的最小非负整数为
\[
N=\left\lceil\frac{A-d}{2d}\right\rceil\quad(A>d),
\]
其中向上取整自动包含落在静止区间端点的情形。由该取整关系，$A-2Nd$ 必在 $(-d,d]$ 内；此前的半周末仍在静止区间之外。因此
\[
x_{\rm stop}=(-1)^N(A-2Nd),\qquad
t_{\rm stop}=\frac{N\pi}{\omega_0},\qquad
Q=\frac k2(A^2-x_{\rm stop}^2).
\]
最后一段即使没有跨过 $x=0$，仍是相对于移心位置的完整半周，耗时依然为 $\pi/\omega_0$。若初态已在 $|x_0|\le d$ 且静止，则不启动，单独取 $N=0$。
在边界初态$A=(2j+1)d$、$j\ge1$时，取整公式给出$N=j$，停止位置为$x_{\rm stop}=(-1)^jd$，静摩擦恰达到上限，仍可维持静止。因此半开区间包含右端$d$，不包含左端$-d$。若$A$略高于这一边界，会多出一个幅度很小、耗时仍为$\pi/\omega_0$的滑动半周；停止位置连续变化，而理想模型中的停止时间会跳变。
\end{solution}

\par\medskip\noindent\begin{minipage}{\textwidth}
\begin{exercise}[弱阻尼单摆的启动阈值（选做）]\label{mp4:exercise:l4-h12}
刚杆单摆在最低点获得正向角速度，随后满足 $\ddot\theta+2\beta\dot\theta+\omega_0^2\sin\theta=0$，其中 $\omega_0=\sqrt{g/l}$，$0<\beta\ll\omega_0$。图\ref{mp4:fig:l4-h12-q}表示最低点启动。求刚好渐近上方倒立位置所需初始角速度的一阶近似，并解释它为何高于无阻尼阈值。
\end{exercise}

\begin{figure}[H]\centering
\begin{tikzpicture}
\fill(0,1.2)circle(2pt)node[above]{$O$};\draw[curve](0,1.2)--(0,-1.2)node[midway,left]{$l$};\fill(0,-1.2)circle(3pt)node[left]{$m$};\draw[axis](0,-1.2)--(1.5,-1.2)node[above]{$\vct v_0$};\draw[guide,domain=-90:30,samples=60]plot({2.4*cos(\x)},{1.2+2.4*sin(\x)});\draw[axis]({2.4*cos(-70)},{1.2+2.4*sin(-70)})arc(-70:-50:2.4);\node at(2.7,-.8){$\theta$ 正向};
\end{tikzpicture}
\caption{单摆由最低点向正角度方向启动。实际越过倒立位置需要比临界启动更大的角速度。}\label{mp4:fig:l4-h12-q}
\end{figure}
\end{minipage}\par\medskip
\begin{solution}
令 $I=ml^2$。将运动方程乘以 $I\dot\theta$，得到惯性系机械能的变化率
\[
\frac{\dd E}{\dd t}=-2I\beta\dot\theta^2,
\qquad E=\frac12I\dot\theta^2+I\omega_0^2(1-\cos\theta).
\]
临界运动从 $\theta=0$ 单调上升，最后以零角速度渐近 $\theta=\pi$。最终能量为 $2I\omega_0^2$，因此初始角速度 $\dot\theta_c$ 满足
\[
\frac12I\dot\theta_c^2=2I\omega_0^2+Q,
\qquad Q=2I\beta\int_0^\infty\dot\theta^2\dd t.
\]
图\ref{mp4:fig:l4-h12-s}的临界有阻尼轨线在无阻尼分界线上方，表示从同一角度继续上升还需保留补偿以后耗散的能量。

\begin{figure}[H]\centering
\begin{tikzpicture}
\draw[axis](-0.4,0)--(3,0)node[right]{$\theta$};
\draw[axis](0,-0.4)--(0,3.1)node[above]{$p$};
\draw[flow,line width=.9pt] plot coordinates{(0.00000,2.37629) (0.00136,2.37590) (0.00272,2.37550) (0.00408,2.37510) (0.00545,2.37469) (0.00681,2.37429) (0.00817,2.37389) (0.00953,2.37348) (0.01089,2.37307) (0.01225,2.37267) (0.01361,2.37226) (0.01497,2.37185) (0.01634,2.37144) (0.01770,2.37102) (0.01906,2.37061) (0.02042,2.37019) (0.02178,2.36978) (0.02314,2.36936) (0.02450,2.36894) (0.02587,2.36852) (0.02723,2.36810) (0.02859,2.36768) (0.02995,2.36726) (0.03131,2.36683) (0.03267,2.36641) (0.03403,2.36598) (0.03540,2.36556) (0.03676,2.36513) (0.03812,2.36470) (0.03948,2.36427) (0.04084,2.36383) (0.04220,2.36340) (0.04356,2.36297) (0.04492,2.36253) (0.04629,2.36209) (0.04765,2.36166) (0.04901,2.36122) (0.05037,2.36078) (0.05173,2.36033) (0.05309,2.35989) (0.05445,2.35945) (0.05582,2.35900) (0.05718,2.35856) (0.05854,2.35811) (0.05990,2.35766) (0.06126,2.35721) (0.06262,2.35676) (0.06398,2.35631) (0.06534,2.35585) (0.06671,2.35540) (0.06807,2.35494) (0.06943,2.35449) (0.07079,2.35403) (0.07215,2.35357) (0.07351,2.35311) (0.07487,2.35265) (0.07624,2.35219) (0.07760,2.35172) (0.07896,2.35126) (0.08032,2.35079) (0.08168,2.35032) (0.08304,2.34985) (0.08440,2.34938) (0.08577,2.34891) (0.08713,2.34844) (0.08849,2.34797) (0.08985,2.34749) (0.09121,2.34702) (0.09257,2.34654) (0.09393,2.34606) (0.09529,2.34558) (0.09666,2.34510) (0.09802,2.34462) (0.09938,2.34414) (0.10074,2.34366) (0.10210,2.34317) (0.10346,2.34268) (0.10482,2.34220) (0.10619,2.34171) (0.10755,2.34122) (0.10891,2.34073) (0.11027,2.34023) (0.11163,2.33974) (0.11299,2.33925) (0.11435,2.33875) (0.11571,2.33825) (0.11708,2.33776) (0.11844,2.33726) (0.11980,2.33676) (0.12116,2.33625) (0.12252,2.33575) (0.12388,2.33525) (0.12524,2.33474) (0.12661,2.33424) (0.12797,2.33373) (0.12933,2.33322) (0.13069,2.33271) (0.13205,2.33220) (0.13341,2.33169) (0.13477,2.33117) (0.13614,2.33066) (0.13750,2.33014) (0.13886,2.32963) (0.14022,2.32911) (0.14158,2.32859) (0.14294,2.32807) (0.14430,2.32755) (0.14566,2.32702) (0.14703,2.32650) (0.14839,2.32597) (0.14975,2.32545) (0.15111,2.32492) (0.15247,2.32439) (0.15383,2.32386) (0.15519,2.32333) (0.15656,2.32280) (0.15792,2.32226) (0.15928,2.32173) (0.16064,2.32119) (0.16200,2.32065) (0.16336,2.32011) (0.16472,2.31958) (0.16609,2.31903) (0.16745,2.31849) (0.16881,2.31795) (0.17017,2.31740) (0.17153,2.31686) (0.17289,2.31631) (0.17425,2.31576) (0.17561,2.31521) (0.17698,2.31466) (0.17834,2.31411) (0.17970,2.31356) (0.18106,2.31301) (0.18242,2.31245) (0.18378,2.31189) (0.18514,2.31134) (0.18651,2.31078) (0.18787,2.31022) (0.18923,2.30966) (0.19059,2.30909) (0.19195,2.30853) (0.19331,2.30797) (0.19467,2.30740) (0.19603,2.30683) (0.19740,2.30626) (0.19876,2.30570) (0.20012,2.30512) (0.20148,2.30455) (0.20284,2.30398) (0.20420,2.30341) (0.20556,2.30283) (0.20693,2.30225) (0.20829,2.30167) (0.20965,2.30110) (0.21101,2.30052) (0.21237,2.29993) (0.21373,2.29935) (0.21509,2.29877) (0.21646,2.29818) (0.21782,2.29760) (0.21918,2.29701) (0.22054,2.29642) (0.22190,2.29583) (0.22326,2.29524) (0.22462,2.29465) (0.22598,2.29405) (0.22735,2.29346) (0.22871,2.29286) (0.23007,2.29227) (0.23143,2.29167) (0.23279,2.29107) (0.23415,2.29047) (0.23551,2.28986) (0.23688,2.28926) (0.23824,2.28866) (0.23960,2.28805) (0.24096,2.28745) (0.24232,2.28684) (0.24368,2.28623) (0.24504,2.28562) (0.24640,2.28501) (0.24777,2.28439) (0.24913,2.28378) (0.25049,2.28316) (0.25185,2.28255) (0.25321,2.28193) (0.25457,2.28131) (0.25593,2.28069) (0.25730,2.28007) (0.25866,2.27945) (0.26002,2.27883) (0.26138,2.27820) (0.26274,2.27757) (0.26410,2.27695) (0.26546,2.27632) (0.26683,2.27569) (0.26819,2.27506) (0.26955,2.27443) (0.27091,2.27379) (0.27227,2.27316) (0.27363,2.27252) (0.27499,2.27189) (0.27635,2.27125) (0.27772,2.27061) (0.27908,2.26997) (0.28044,2.26933) (0.28180,2.26869) (0.28316,2.26804) (0.28452,2.26740) (0.28588,2.26675) (0.28725,2.26610) (0.28861,2.26545) (0.28997,2.26480) (0.29133,2.26415) (0.29269,2.26350) (0.29405,2.26285) (0.29541,2.26219) (0.29677,2.26154) (0.29814,2.26088) (0.29950,2.26022) (0.30086,2.25956) (0.30222,2.25890) (0.30358,2.25824) (0.30494,2.25757) (0.30630,2.25691) (0.30767,2.25624) (0.30903,2.25558) (0.31039,2.25491) (0.31175,2.25424) (0.31311,2.25357) (0.31447,2.25290) (0.31583,2.25223) (0.31720,2.25155) (0.31856,2.25088) (0.31992,2.25020) (0.32128,2.24952) (0.32264,2.24884) (0.32400,2.24816) (0.32536,2.24748) (0.32672,2.24680) (0.32809,2.24612) (0.32945,2.24543) (0.33081,2.24475) (0.33217,2.24406) (0.33353,2.24337) (0.33489,2.24268) (0.33625,2.24199) (0.33762,2.24130) (0.33898,2.24060) (0.34034,2.23991) (0.34170,2.23921) (0.34306,2.23852) (0.34442,2.23782) (0.34578,2.23712) (0.34714,2.23642) (0.34851,2.23572) (0.34987,2.23501) (0.35123,2.23431) (0.35259,2.23360) (0.35395,2.23290) (0.35531,2.23219) (0.35667,2.23148) (0.35804,2.23077) (0.35940,2.23006) (0.36076,2.22935) (0.36212,2.22863) (0.36348,2.22792) (0.36484,2.22720) (0.36620,2.22648) (0.36757,2.22577) (0.36893,2.22505) (0.37029,2.22432) (0.37165,2.22360) (0.37301,2.22288) (0.37437,2.22215) (0.37573,2.22143) (0.37709,2.22070) (0.37846,2.21997) (0.37982,2.21924) (0.38118,2.21851) (0.38254,2.21778) (0.38390,2.21705) (0.38526,2.21631) (0.38662,2.21558) (0.38799,2.21484) (0.38935,2.21410) (0.39071,2.21336) (0.39207,2.21262) (0.39343,2.21188) (0.39479,2.21114) (0.39615,2.21040) (0.39751,2.20965) (0.39888,2.20890) (0.40024,2.20816) (0.40160,2.20741) (0.40296,2.20666) (0.40432,2.20591) (0.40568,2.20516) (0.40704,2.20440) (0.40841,2.20365) (0.40977,2.20289) (0.41113,2.20213) (0.41249,2.20138) (0.41385,2.20062) (0.41521,2.19986) (0.41657,2.19909) (0.41794,2.19833) (0.41930,2.19757) (0.42066,2.19680) (0.42202,2.19603) (0.42338,2.19527) (0.42474,2.19450) (0.42610,2.19373) (0.42746,2.19295) (0.42883,2.19218) (0.43019,2.19141) (0.43155,2.19063) (0.43291,2.18986) (0.43427,2.18908) (0.43563,2.18830) (0.43699,2.18752) (0.43836,2.18674) (0.43972,2.18596) (0.44108,2.18517) (0.44244,2.18439) (0.44380,2.18360) (0.44516,2.18281) (0.44652,2.18203) (0.44788,2.18124) (0.44925,2.18045) (0.45061,2.17965) (0.45197,2.17886) (0.45333,2.17807) (0.45469,2.17727) (0.45605,2.17647) (0.45741,2.17568) (0.45878,2.17488) (0.46014,2.17408) (0.46150,2.17328) (0.46286,2.17247) (0.46422,2.17167) (0.46558,2.17086) (0.46694,2.17006) (0.46831,2.16925) (0.46967,2.16844) (0.47103,2.16763) (0.47239,2.16682) (0.47375,2.16601) (0.47511,2.16519) (0.47647,2.16438) (0.47783,2.16356) (0.47920,2.16275) (0.48056,2.16193) (0.48192,2.16111) (0.48328,2.16029) (0.48464,2.15946) (0.48600,2.15864) (0.48736,2.15782) (0.48873,2.15699) (0.49009,2.15617) (0.49145,2.15534) (0.49281,2.15451) (0.49417,2.15368) (0.49553,2.15285) (0.49689,2.15201) (0.49826,2.15118) (0.49962,2.15034) (0.50098,2.14951) (0.50234,2.14867) (0.50370,2.14783) (0.50506,2.14699) (0.50642,2.14615) (0.50778,2.14531) (0.50915,2.14447) (0.51051,2.14362) (0.51187,2.14277) (0.51323,2.14193) (0.51459,2.14108) (0.51595,2.14023) (0.51731,2.13938) (0.51868,2.13853) (0.52004,2.13767) (0.52140,2.13682) (0.52276,2.13597) (0.52412,2.13511) (0.52548,2.13425) (0.52684,2.13339) (0.52820,2.13253) (0.52957,2.13167) (0.53093,2.13081) (0.53229,2.12994) (0.53365,2.12908) (0.53501,2.12821) (0.53637,2.12735) (0.53773,2.12648) (0.53910,2.12561) (0.54046,2.12474) (0.54182,2.12386) (0.54318,2.12299) (0.54454,2.12212) (0.54590,2.12124) (0.54726,2.12036) (0.54863,2.11949) (0.54999,2.11861) (0.55135,2.11773) (0.55271,2.11685) (0.55407,2.11596) (0.55543,2.11508) (0.55679,2.11419) (0.55815,2.11331) (0.55952,2.11242) (0.56088,2.11153) (0.56224,2.11064) (0.56360,2.10975) (0.56496,2.10886) (0.56632,2.10796) (0.56768,2.10707) (0.56905,2.10617) (0.57041,2.10528) (0.57177,2.10438) (0.57313,2.10348) (0.57449,2.10258) (0.57585,2.10168) (0.57721,2.10077) (0.57857,2.09987) (0.57994,2.09896) (0.58130,2.09806) (0.58266,2.09715) (0.58402,2.09624) (0.58538,2.09533) (0.58674,2.09442) (0.58810,2.09351) (0.58947,2.09259) (0.59083,2.09168) (0.59219,2.09076) (0.59355,2.08985) (0.59491,2.08893) (0.59627,2.08801) (0.59763,2.08709) (0.59900,2.08617) (0.60036,2.08524) (0.60172,2.08432) (0.60308,2.08339) (0.60444,2.08247) (0.60580,2.08154) (0.60716,2.08061) (0.60852,2.07968) (0.60989,2.07875) (0.61125,2.07782) (0.61261,2.07688) (0.61397,2.07595) (0.61533,2.07501) (0.61669,2.07407) (0.61805,2.07314) (0.61942,2.07220) (0.62078,2.07126) (0.62214,2.07031) (0.62350,2.06937) (0.62486,2.06843) (0.62622,2.06748) (0.62758,2.06654) (0.62894,2.06559) (0.63031,2.06464) (0.63167,2.06369) (0.63303,2.06274) (0.63439,2.06178) (0.63575,2.06083) (0.63711,2.05988) (0.63847,2.05892) (0.63984,2.05796) (0.64120,2.05700) (0.64256,2.05605) (0.64392,2.05508) (0.64528,2.05412) (0.64664,2.05316) (0.64800,2.05220) (0.64937,2.05123) (0.65073,2.05026) (0.65209,2.04930) (0.65345,2.04833) (0.65481,2.04736) (0.65617,2.04639) (0.65753,2.04541) (0.65889,2.04444) (0.66026,2.04347) (0.66162,2.04249) (0.66298,2.04151) (0.66434,2.04053) (0.66570,2.03956) (0.66706,2.03858) (0.66842,2.03759) (0.66979,2.03661) (0.67115,2.03563) (0.67251,2.03464) (0.67387,2.03366) (0.67523,2.03267) (0.67659,2.03168) (0.67795,2.03069) (0.67931,2.02970) (0.68068,2.02871) (0.68204,2.02771) (0.68340,2.02672) (0.68476,2.02572) (0.68612,2.02473) (0.68748,2.02373) (0.68884,2.02273) (0.69021,2.02173) (0.69157,2.02073) (0.69293,2.01972) (0.69429,2.01872) (0.69565,2.01772) (0.69701,2.01671) (0.69837,2.01570) (0.69974,2.01469) (0.70110,2.01368) (0.70246,2.01267) (0.70382,2.01166) (0.70518,2.01065) (0.70654,2.00964) (0.70790,2.00862) (0.70926,2.00760) (0.71063,2.00659) (0.71199,2.00557) (0.71335,2.00455) (0.71471,2.00353) (0.71607,2.00250) (0.71743,2.00148) (0.71879,2.00046) (0.72016,1.99943) (0.72152,1.99840) (0.72288,1.99738) (0.72424,1.99635) (0.72560,1.99532) (0.72696,1.99428) (0.72832,1.99325) (0.72968,1.99222) (0.73105,1.99118) (0.73241,1.99015) (0.73377,1.98911) (0.73513,1.98807) (0.73649,1.98703) (0.73785,1.98599) (0.73921,1.98495) (0.74058,1.98391) (0.74194,1.98286) (0.74330,1.98182) (0.74466,1.98077) (0.74602,1.97972) (0.74738,1.97867) (0.74874,1.97763) (0.75011,1.97657) (0.75147,1.97552) (0.75283,1.97447) (0.75419,1.97341) (0.75555,1.97236) (0.75691,1.97130) (0.75827,1.97024) (0.75963,1.96919) (0.76100,1.96813) (0.76236,1.96706) (0.76372,1.96600) (0.76508,1.96494) (0.76644,1.96387) (0.76780,1.96281) (0.76916,1.96174) (0.77053,1.96067) (0.77189,1.95960) (0.77325,1.95853) (0.77461,1.95746) (0.77597,1.95639) (0.77733,1.95532) (0.77869,1.95424) (0.78005,1.95316) (0.78142,1.95209) (0.78278,1.95101) (0.78414,1.94993) (0.78550,1.94885) (0.78686,1.94777) (0.78822,1.94668) (0.78958,1.94560) (0.79095,1.94452) (0.79231,1.94343) (0.79367,1.94234) (0.79503,1.94125) (0.79639,1.94016) (0.79775,1.93907) (0.79911,1.93798) (0.80048,1.93689) (0.80184,1.93579) (0.80320,1.93470) (0.80456,1.93360) (0.80592,1.93251) (0.80728,1.93141) (0.80864,1.93031) (0.81000,1.92921) (0.81137,1.92810) (0.81273,1.92700) (0.81409,1.92590) (0.81545,1.92479) (0.81681,1.92368) (0.81817,1.92258) (0.81953,1.92147) (0.82090,1.92036) (0.82226,1.91925) (0.82362,1.91814) (0.82498,1.91702) (0.82634,1.91591) (0.82770,1.91479) (0.82906,1.91368) (0.83043,1.91256) (0.83179,1.91144) (0.83315,1.91032) (0.83451,1.90920) (0.83587,1.90808) (0.83723,1.90695) (0.83859,1.90583) (0.83995,1.90470) (0.84132,1.90358) (0.84268,1.90245) (0.84404,1.90132) (0.84540,1.90019) (0.84676,1.89906) (0.84812,1.89793) (0.84948,1.89679) (0.85085,1.89566) (0.85221,1.89452) (0.85357,1.89339) (0.85493,1.89225) (0.85629,1.89111) (0.85765,1.88997) (0.85901,1.88883) (0.86037,1.88769) (0.86174,1.88654) (0.86310,1.88540) (0.86446,1.88425) (0.86582,1.88311) (0.86718,1.88196) (0.86854,1.88081) (0.86990,1.87966) (0.87127,1.87851) (0.87263,1.87736) (0.87399,1.87620) (0.87535,1.87505) (0.87671,1.87389) (0.87807,1.87274) (0.87943,1.87158) (0.88080,1.87042) (0.88216,1.86926) (0.88352,1.86810) (0.88488,1.86694) (0.88624,1.86578) (0.88760,1.86461) (0.88896,1.86345) (0.89032,1.86228) (0.89169,1.86111) (0.89305,1.85994) (0.89441,1.85877) (0.89577,1.85760) (0.89713,1.85643) (0.89849,1.85526) (0.89985,1.85408) (0.90122,1.85291) (0.90258,1.85173) (0.90394,1.85055) (0.90530,1.84938) (0.90666,1.84820) (0.90802,1.84702) (0.90938,1.84583) (0.91074,1.84465) (0.91211,1.84347) (0.91347,1.84228) (0.91483,1.84110) (0.91619,1.83991) (0.91755,1.83872) (0.91891,1.83753) (0.92027,1.83634) (0.92164,1.83515) (0.92300,1.83396) (0.92436,1.83276) (0.92572,1.83157) (0.92708,1.83037) (0.92844,1.82918) (0.92980,1.82798) (0.93117,1.82678) (0.93253,1.82558) (0.93389,1.82438) (0.93525,1.82318) (0.93661,1.82197) (0.93797,1.82077) (0.93933,1.81956) (0.94069,1.81836) (0.94206,1.81715) (0.94342,1.81594) (0.94478,1.81473) (0.94614,1.81352) (0.94750,1.81231) (0.94886,1.81109) (0.95022,1.80988) (0.95159,1.80866) (0.95295,1.80745) (0.95431,1.80623) (0.95567,1.80501) (0.95703,1.80379) (0.95839,1.80257) (0.95975,1.80135) (0.96111,1.80013) (0.96248,1.79890) (0.96384,1.79768) (0.96520,1.79645) (0.96656,1.79522) (0.96792,1.79400) (0.96928,1.79277) (0.97064,1.79154) (0.97201,1.79031) (0.97337,1.78907) (0.97473,1.78784) (0.97609,1.78661) (0.97745,1.78537) (0.97881,1.78413) (0.98017,1.78290) (0.98154,1.78166) (0.98290,1.78042) (0.98426,1.77918) (0.98562,1.77794) (0.98698,1.77669) (0.98834,1.77545) (0.98970,1.77420) (0.99106,1.77296) (0.99243,1.77171) (0.99379,1.77046) (0.99515,1.76921) (0.99651,1.76796) (0.99787,1.76671) (0.99923,1.76546) (1.00059,1.76420) (1.00196,1.76295) (1.00332,1.76169) (1.00468,1.76044) (1.00604,1.75918) (1.00740,1.75792) (1.00876,1.75666) (1.01012,1.75540) (1.01148,1.75414) (1.01285,1.75288) (1.01421,1.75161) (1.01557,1.75035) (1.01693,1.74908) (1.01829,1.74781) (1.01965,1.74655) (1.02101,1.74528) (1.02238,1.74401) (1.02374,1.74274) (1.02510,1.74146) (1.02646,1.74019) (1.02782,1.73892) (1.02918,1.73764) (1.03054,1.73636) (1.03191,1.73509) (1.03327,1.73381) (1.03463,1.73253) (1.03599,1.73125) (1.03735,1.72997) (1.03871,1.72868) (1.04007,1.72740) (1.04143,1.72611) (1.04280,1.72483) (1.04416,1.72354) (1.04552,1.72225) (1.04688,1.72096) (1.04824,1.71968) (1.04960,1.71838) (1.05096,1.71709) (1.05233,1.71580) (1.05369,1.71451) (1.05505,1.71321) (1.05641,1.71191) (1.05777,1.71062) (1.05913,1.70932) (1.06049,1.70802) (1.06185,1.70672) (1.06322,1.70542) (1.06458,1.70412) (1.06594,1.70281) (1.06730,1.70151) (1.06866,1.70020) (1.07002,1.69890) (1.07138,1.69759) (1.07275,1.69628) (1.07411,1.69497) (1.07547,1.69366) (1.07683,1.69235) (1.07819,1.69104) (1.07955,1.68972) (1.08091,1.68841) (1.08228,1.68709) (1.08364,1.68578) (1.08500,1.68446) (1.08636,1.68314) (1.08772,1.68182) (1.08908,1.68050) (1.09044,1.67918) (1.09180,1.67786) (1.09317,1.67653) (1.09453,1.67521) (1.09589,1.67388) (1.09725,1.67256) (1.09861,1.67123) (1.09997,1.66990) (1.10133,1.66857) (1.10270,1.66724) (1.10406,1.66591) (1.10542,1.66458) (1.10678,1.66324) (1.10814,1.66191) (1.10950,1.66057) (1.11086,1.65924) (1.11222,1.65790) (1.11359,1.65656) (1.11495,1.65522) (1.11631,1.65388) (1.11767,1.65254) (1.11903,1.65120) (1.12039,1.64985) (1.12175,1.64851) (1.12312,1.64716) (1.12448,1.64582) (1.12584,1.64447) (1.12720,1.64312) (1.12856,1.64177) (1.12992,1.64042) (1.13128,1.63907) (1.13265,1.63772) (1.13401,1.63636) (1.13537,1.63501) (1.13673,1.63365) (1.13809,1.63230) (1.13945,1.63094) (1.14081,1.62958) (1.14217,1.62822) (1.14354,1.62686) (1.14490,1.62550) (1.14626,1.62414) (1.14762,1.62277) (1.14898,1.62141) (1.15034,1.62005) (1.15170,1.61868) (1.15307,1.61731) (1.15443,1.61594) (1.15579,1.61457) (1.15715,1.61320) (1.15851,1.61183) (1.15987,1.61046) (1.16123,1.60909) (1.16260,1.60771) (1.16396,1.60634) (1.16532,1.60496) (1.16668,1.60359) (1.16804,1.60221) (1.16940,1.60083) (1.17076,1.59945) (1.17212,1.59807) (1.17349,1.59669) (1.17485,1.59530) (1.17621,1.59392) (1.17757,1.59254) (1.17893,1.59115) (1.18029,1.58976) (1.18165,1.58838) (1.18302,1.58699) (1.18438,1.58560) (1.18574,1.58421) (1.18710,1.58282) (1.18846,1.58142) (1.18982,1.58003) (1.19118,1.57864) (1.19254,1.57724) (1.19391,1.57584) (1.19527,1.57445) (1.19663,1.57305) (1.19799,1.57165) (1.19935,1.57025) (1.20071,1.56885) (1.20207,1.56745) (1.20344,1.56604) (1.20480,1.56464) (1.20616,1.56324) (1.20752,1.56183) (1.20888,1.56042) (1.21024,1.55902) (1.21160,1.55761) (1.21297,1.55620) (1.21433,1.55479) (1.21569,1.55338) (1.21705,1.55196) (1.21841,1.55055) (1.21977,1.54914) (1.22113,1.54772) (1.22249,1.54630) (1.22386,1.54489) (1.22522,1.54347) (1.22658,1.54205) (1.22794,1.54063) (1.22930,1.53921) (1.23066,1.53779) (1.23202,1.53636) (1.23339,1.53494) (1.23475,1.53352) (1.23611,1.53209) (1.23747,1.53066) (1.23883,1.52924) (1.24019,1.52781) (1.24155,1.52638) (1.24291,1.52495) (1.24428,1.52352) (1.24564,1.52209) (1.24700,1.52065) (1.24836,1.51922) (1.24972,1.51779) (1.25108,1.51635) (1.25244,1.51491) (1.25381,1.51348) (1.25517,1.51204) (1.25653,1.51060) (1.25789,1.50916) (1.25925,1.50772) (1.26061,1.50627) (1.26197,1.50483) (1.26334,1.50339) (1.26470,1.50194) (1.26606,1.50050) (1.26742,1.49905) (1.26878,1.49760) (1.27014,1.49615) (1.27150,1.49470) (1.27286,1.49325) (1.27423,1.49180) (1.27559,1.49035) (1.27695,1.48890) (1.27831,1.48744) (1.27967,1.48599) (1.28103,1.48453) (1.28239,1.48308) (1.28376,1.48162) (1.28512,1.48016) (1.28648,1.47870) (1.28784,1.47724) (1.28920,1.47578) (1.29056,1.47432) (1.29192,1.47285) (1.29328,1.47139) (1.29465,1.46992) (1.29601,1.46846) (1.29737,1.46699) (1.29873,1.46552) (1.30009,1.46406) (1.30145,1.46259) (1.30281,1.46112) (1.30418,1.45965) (1.30554,1.45817) (1.30690,1.45670) (1.30826,1.45523) (1.30962,1.45375) (1.31098,1.45228) (1.31234,1.45080) (1.31371,1.44932) (1.31507,1.44785) (1.31643,1.44637) (1.31779,1.44489) (1.31915,1.44341) (1.32051,1.44192) (1.32187,1.44044) (1.32323,1.43896) (1.32460,1.43747) (1.32596,1.43599) (1.32732,1.43450) (1.32868,1.43301) (1.33004,1.43153) (1.33140,1.43004) (1.33276,1.42855) (1.33413,1.42706) (1.33549,1.42557) (1.33685,1.42407) (1.33821,1.42258) (1.33957,1.42109) (1.34093,1.41959) (1.34229,1.41810) (1.34365,1.41660) (1.34502,1.41510) (1.34638,1.41360) (1.34774,1.41210) (1.34910,1.41060) (1.35046,1.40910) (1.35182,1.40760) (1.35318,1.40610) (1.35455,1.40460) (1.35591,1.40309) (1.35727,1.40159) (1.35863,1.40008) (1.35999,1.39857) (1.36135,1.39706) (1.36271,1.39556) (1.36408,1.39405) (1.36544,1.39254) (1.36680,1.39102) (1.36816,1.38951) (1.36952,1.38800) (1.37088,1.38649) (1.37224,1.38497) (1.37360,1.38346) (1.37497,1.38194) (1.37633,1.38042) (1.37769,1.37890) (1.37905,1.37738) (1.38041,1.37586) (1.38177,1.37434) (1.38313,1.37282) (1.38450,1.37130) (1.38586,1.36978) (1.38722,1.36825) (1.38858,1.36673) (1.38994,1.36520) (1.39130,1.36368) (1.39266,1.36215) (1.39402,1.36062) (1.39539,1.35909) (1.39675,1.35756) (1.39811,1.35603) (1.39947,1.35450) (1.40083,1.35297) (1.40219,1.35143) (1.40355,1.34990) (1.40492,1.34837) (1.40628,1.34683) (1.40764,1.34529) (1.40900,1.34376) (1.41036,1.34222) (1.41172,1.34068) (1.41308,1.33914) (1.41445,1.33760) (1.41581,1.33606) (1.41717,1.33451) (1.41853,1.33297) (1.41989,1.33143) (1.42125,1.32988) (1.42261,1.32834) (1.42397,1.32679) (1.42534,1.32524) (1.42670,1.32370) (1.42806,1.32215) (1.42942,1.32060) (1.43078,1.31905) (1.43214,1.31750) (1.43350,1.31594) (1.43487,1.31439) (1.43623,1.31284) (1.43759,1.31128) (1.43895,1.30973) (1.44031,1.30817) (1.44167,1.30661) (1.44303,1.30506) (1.44439,1.30350) (1.44576,1.30194) (1.44712,1.30038) (1.44848,1.29882) (1.44984,1.29726) (1.45120,1.29569) (1.45256,1.29413) (1.45392,1.29257) (1.45529,1.29100) (1.45665,1.28943) (1.45801,1.28787) (1.45937,1.28630) (1.46073,1.28473) (1.46209,1.28316) (1.46345,1.28159) (1.46482,1.28002) (1.46618,1.27845) (1.46754,1.27688) (1.46890,1.27531) (1.47026,1.27373) (1.47162,1.27216) (1.47298,1.27058) (1.47434,1.26901) (1.47571,1.26743) (1.47707,1.26585) (1.47843,1.26427) (1.47979,1.26270) (1.48115,1.26112) (1.48251,1.25954) (1.48387,1.25795) (1.48524,1.25637) (1.48660,1.25479) (1.48796,1.25320) (1.48932,1.25162) (1.49068,1.25003) (1.49204,1.24845) (1.49340,1.24686) (1.49477,1.24527) (1.49613,1.24369) (1.49749,1.24210) (1.49885,1.24051) (1.50021,1.23892) (1.50157,1.23732) (1.50293,1.23573) (1.50429,1.23414) (1.50566,1.23254) (1.50702,1.23095) (1.50838,1.22935) (1.50974,1.22776) (1.51110,1.22616) (1.51246,1.22456) (1.51382,1.22297) (1.51519,1.22137) (1.51655,1.21977) (1.51791,1.21817) (1.51927,1.21656) (1.52063,1.21496) (1.52199,1.21336) (1.52335,1.21176) (1.52471,1.21015) (1.52608,1.20855) (1.52744,1.20694) (1.52880,1.20533) (1.53016,1.20373) (1.53152,1.20212) (1.53288,1.20051) (1.53424,1.19890) (1.53561,1.19729) (1.53697,1.19568) (1.53833,1.19407) (1.53969,1.19245) (1.54105,1.19084) (1.54241,1.18923) (1.54377,1.18761) (1.54514,1.18600) (1.54650,1.18438) (1.54786,1.18276) (1.54922,1.18115) (1.55058,1.17953) (1.55194,1.17791) (1.55330,1.17629) (1.55466,1.17467) (1.55603,1.17305) (1.55739,1.17142) (1.55875,1.16980) (1.56011,1.16818) (1.56147,1.16655) (1.56283,1.16493) (1.56419,1.16330) (1.56556,1.16168) (1.56692,1.16005) (1.56828,1.15842) (1.56964,1.15679) (1.57100,1.15516) (1.57236,1.15353) (1.57372,1.15190) (1.57508,1.15027) (1.57645,1.14864) (1.57781,1.14701) (1.57917,1.14537) (1.58053,1.14374) (1.58189,1.14210) (1.58325,1.14047) (1.58461,1.13883) (1.58598,1.13719) (1.58734,1.13555) (1.58870,1.13392) (1.59006,1.13228) (1.59142,1.13064) (1.59278,1.12900) (1.59414,1.12735) (1.59551,1.12571) (1.59687,1.12407) (1.59823,1.12243) (1.59959,1.12078) (1.60095,1.11914) (1.60231,1.11749) (1.60367,1.11585) (1.60503,1.11420) (1.60640,1.11255) (1.60776,1.11090) (1.60912,1.10925) (1.61048,1.10760) (1.61184,1.10595) (1.61320,1.10430) (1.61456,1.10265) (1.61593,1.10100) (1.61729,1.09934) (1.61865,1.09769) (1.62001,1.09604) (1.62137,1.09438) (1.62273,1.09272) (1.62409,1.09107) (1.62545,1.08941) (1.62682,1.08775) (1.62818,1.08609) (1.62954,1.08444) (1.63090,1.08278) (1.63226,1.08111) (1.63362,1.07945) (1.63498,1.07779) (1.63635,1.07613) (1.63771,1.07447) (1.63907,1.07280) (1.64043,1.07114) (1.64179,1.06947) (1.64315,1.06781) (1.64451,1.06614) (1.64588,1.06447) (1.64724,1.06280) (1.64860,1.06113) (1.64996,1.05947) (1.65132,1.05780) (1.65268,1.05612) (1.65404,1.05445) (1.65540,1.05278) (1.65677,1.05111) (1.65813,1.04944) (1.65949,1.04776) (1.66085,1.04609) (1.66221,1.04441) (1.66357,1.04274) (1.66493,1.04106) (1.66630,1.03938) (1.66766,1.03770) (1.66902,1.03603) (1.67038,1.03435) (1.67174,1.03267) (1.67310,1.03099) (1.67446,1.02931) (1.67582,1.02762) (1.67719,1.02594) (1.67855,1.02426) (1.67991,1.02258) (1.68127,1.02089) (1.68263,1.01921) (1.68399,1.01752) (1.68535,1.01583) (1.68672,1.01415) (1.68808,1.01246) (1.68944,1.01077) (1.69080,1.00908) (1.69216,1.00739) (1.69352,1.00570) (1.69488,1.00401) (1.69625,1.00232) (1.69761,1.00063) (1.69897,0.99894) (1.70033,0.99725) (1.70169,0.99555) (1.70305,0.99386) (1.70441,0.99216) (1.70577,0.99047) (1.70714,0.98877) (1.70850,0.98707) (1.70986,0.98538) (1.71122,0.98368) (1.71258,0.98198) (1.71394,0.98028) (1.71530,0.97858) (1.71667,0.97688) (1.71803,0.97518) (1.71939,0.97348) (1.72075,0.97178) (1.72211,0.97007) (1.72347,0.96837) (1.72483,0.96666) (1.72619,0.96496) (1.72756,0.96325) (1.72892,0.96155) (1.73028,0.95984) (1.73164,0.95813) (1.73300,0.95643) (1.73436,0.95472) (1.73572,0.95301) (1.73709,0.95130) (1.73845,0.94959) (1.73981,0.94788) (1.74117,0.94617) (1.74253,0.94446) (1.74389,0.94274) (1.74525,0.94103) (1.74662,0.93932) (1.74798,0.93760) (1.74934,0.93589) (1.75070,0.93417) (1.75206,0.93246) (1.75342,0.93074) (1.75478,0.92902) (1.75614,0.92731) (1.75751,0.92559) (1.75887,0.92387) (1.76023,0.92215) (1.76159,0.92043) (1.76295,0.91871) (1.76431,0.91699) (1.76567,0.91527) (1.76704,0.91354) (1.76840,0.91182) (1.76976,0.91010) (1.77112,0.90837) (1.77248,0.90665) (1.77384,0.90492) (1.77520,0.90320) (1.77656,0.90147) (1.77793,0.89974) (1.77929,0.89802) (1.78065,0.89629) (1.78201,0.89456) (1.78337,0.89283) (1.78473,0.89110) (1.78609,0.88937) (1.78746,0.88764) (1.78882,0.88591) (1.79018,0.88418) (1.79154,0.88244) (1.79290,0.88071) (1.79426,0.87898) (1.79562,0.87724) (1.79699,0.87551) (1.79835,0.87377) (1.79971,0.87204) (1.80107,0.87030) (1.80243,0.86856) (1.80379,0.86683) (1.80515,0.86509) (1.80651,0.86335) (1.80788,0.86161) (1.80924,0.85987) (1.81060,0.85813) (1.81196,0.85639) (1.81332,0.85465) (1.81468,0.85291) (1.81604,0.85117) (1.81741,0.84942) (1.81877,0.84768) (1.82013,0.84594) (1.82149,0.84419) (1.82285,0.84245) (1.82421,0.84070) (1.82557,0.83895) (1.82694,0.83721) (1.82830,0.83546) (1.82966,0.83371) (1.83102,0.83197) (1.83238,0.83022) (1.83374,0.82847) (1.83510,0.82672) (1.83646,0.82497) (1.83783,0.82322) (1.83919,0.82147) (1.84055,0.81971) (1.84191,0.81796) (1.84327,0.81621) (1.84463,0.81446) (1.84599,0.81270) (1.84736,0.81095) (1.84872,0.80919) (1.85008,0.80744) (1.85144,0.80568) (1.85280,0.80393) (1.85416,0.80217) (1.85552,0.80041) (1.85688,0.79865) (1.85825,0.79690) (1.85961,0.79514) (1.86097,0.79338) (1.86233,0.79162) (1.86369,0.78986) (1.86505,0.78810) (1.86641,0.78634) (1.86778,0.78457) (1.86914,0.78281) (1.87050,0.78105) (1.87186,0.77929) (1.87322,0.77752) (1.87458,0.77576) (1.87594,0.77399) (1.87731,0.77223) (1.87867,0.77046) (1.88003,0.76870) (1.88139,0.76693) (1.88275,0.76516) (1.88411,0.76339) (1.88547,0.76163) (1.88683,0.75986) (1.88820,0.75809) (1.88956,0.75632) (1.89092,0.75455) (1.89228,0.75278) (1.89364,0.75101) (1.89500,0.74924) (1.89636,0.74746) (1.89773,0.74569) (1.89909,0.74392) (1.90045,0.74215) (1.90181,0.74037) (1.90317,0.73860) (1.90453,0.73682) (1.90589,0.73505) (1.90725,0.73327) (1.90862,0.73150) (1.90998,0.72972) (1.91134,0.72794) (1.91270,0.72617) (1.91406,0.72439) (1.91542,0.72261) (1.91678,0.72083) (1.91815,0.71905) (1.91951,0.71727) (1.92087,0.71549) (1.92223,0.71371) (1.92359,0.71193) (1.92495,0.71015) (1.92631,0.70837) (1.92768,0.70658) (1.92904,0.70480) (1.93040,0.70302) (1.93176,0.70123) (1.93312,0.69945) (1.93448,0.69766) (1.93584,0.69588) (1.93720,0.69409) (1.93857,0.69231) (1.93993,0.69052) (1.94129,0.68873) (1.94265,0.68695) (1.94401,0.68516) (1.94537,0.68337) (1.94673,0.68158) (1.94810,0.67979) (1.94946,0.67800) (1.95082,0.67622) (1.95218,0.67442) (1.95354,0.67263) (1.95490,0.67084) (1.95626,0.66905) (1.95762,0.66726) (1.95899,0.66547) (1.96035,0.66367) (1.96171,0.66188) (1.96307,0.66009) (1.96443,0.65829) (1.96579,0.65650) (1.96715,0.65470) (1.96852,0.65291) (1.96988,0.65111) (1.97124,0.64932) (1.97260,0.64752) (1.97396,0.64572) (1.97532,0.64393) (1.97668,0.64213) (1.97805,0.64033) (1.97941,0.63853) (1.98077,0.63673) (1.98213,0.63493) (1.98349,0.63313) (1.98485,0.63133) (1.98621,0.62953) (1.98757,0.62773) (1.98894,0.62593) (1.99030,0.62413) (1.99166,0.62233) (1.99302,0.62052) (1.99438,0.61872) (1.99574,0.61692) (1.99710,0.61512) (1.99847,0.61331) (1.99983,0.61151) (2.00119,0.60970) (2.00255,0.60790) (2.00391,0.60609) (2.00527,0.60429) (2.00663,0.60248) (2.00799,0.60067) (2.00936,0.59886) (2.01072,0.59706) (2.01208,0.59525) (2.01344,0.59344) (2.01480,0.59163) (2.01616,0.58982) (2.01752,0.58801) (2.01889,0.58620) (2.02025,0.58439) (2.02161,0.58258) (2.02297,0.58077) (2.02433,0.57896) (2.02569,0.57715) (2.02705,0.57534) (2.02842,0.57353) (2.02978,0.57171) (2.03114,0.56990) (2.03250,0.56809) (2.03386,0.56627) (2.03522,0.56446) (2.03658,0.56264) (2.03794,0.56083) (2.03931,0.55901) (2.04067,0.55720) (2.04203,0.55538) (2.04339,0.55357) (2.04475,0.55175) (2.04611,0.54993) (2.04747,0.54812) (2.04884,0.54630) (2.05020,0.54448) (2.05156,0.54266) (2.05292,0.54084) (2.05428,0.53903) (2.05564,0.53721) (2.05700,0.53539) (2.05836,0.53357) (2.05973,0.53175) (2.06109,0.52993) (2.06245,0.52810) (2.06381,0.52628) (2.06517,0.52446) (2.06653,0.52264) (2.06789,0.52082) (2.06926,0.51899) (2.07062,0.51717) (2.07198,0.51535) (2.07334,0.51352) (2.07470,0.51170) (2.07606,0.50988) (2.07742,0.50805) (2.07879,0.50623) (2.08015,0.50440) (2.08151,0.50258) (2.08287,0.50075) (2.08423,0.49892) (2.08559,0.49710) (2.08695,0.49527) (2.08831,0.49344) (2.08968,0.49162) (2.09104,0.48979) (2.09240,0.48796) (2.09376,0.48613) (2.09512,0.48430) (2.09648,0.48247) (2.09784,0.48064) (2.09921,0.47882) (2.10057,0.47699) (2.10193,0.47516) (2.10329,0.47332) (2.10465,0.47149) (2.10601,0.46966) (2.10737,0.46783) (2.10873,0.46600) (2.11010,0.46417) (2.11146,0.46234) (2.11282,0.46050) (2.11418,0.45867) (2.11554,0.45684) (2.11690,0.45500) (2.11826,0.45317) (2.11963,0.45134) (2.12099,0.44950) (2.12235,0.44767) (2.12371,0.44583) (2.12507,0.44400) (2.12643,0.44216) (2.12779,0.44033) (2.12916,0.43849) (2.13052,0.43665) (2.13188,0.43482) (2.13324,0.43298) (2.13460,0.43114) (2.13596,0.42931) (2.13732,0.42747) (2.13868,0.42563) (2.14005,0.42379) (2.14141,0.42195) (2.14277,0.42012) (2.14413,0.41828) (2.14549,0.41644) (2.14685,0.41460) (2.14821,0.41276) (2.14958,0.41092) (2.15094,0.40908) (2.15230,0.40724) (2.15366,0.40540) (2.15502,0.40356) (2.15638,0.40171) (2.15774,0.39987) (2.15911,0.39803) (2.16047,0.39619) (2.16183,0.39435) (2.16319,0.39250) (2.16455,0.39066) (2.16591,0.38882) (2.16727,0.38697) (2.16863,0.38513) (2.17000,0.38329) (2.17136,0.38144) (2.17272,0.37960) (2.17408,0.37775) (2.17544,0.37591) (2.17680,0.37406) (2.17816,0.37222) (2.17953,0.37037) (2.18089,0.36853) (2.18225,0.36668) (2.18361,0.36484) (2.18497,0.36299) (2.18633,0.36114) (2.18769,0.35930) (2.18905,0.35745) (2.19042,0.35560) (2.19178,0.35375) (2.19314,0.35191) (2.19450,0.35006) (2.19586,0.34821) (2.19722,0.34636) (2.19858,0.34451) (2.19995,0.34267) (2.20131,0.34082) (2.20267,0.33897) (2.20403,0.33712) (2.20539,0.33527) (2.20675,0.33342) (2.20811,0.33157) (2.20948,0.32972) (2.21084,0.32787) (2.21220,0.32602) (2.21356,0.32417) (2.21492,0.32231) (2.21628,0.32046) (2.21764,0.31861) (2.21900,0.31676) (2.22037,0.31491) (2.22173,0.31306) (2.22309,0.31120) (2.22445,0.30935) (2.22581,0.30750) (2.22717,0.30564) (2.22853,0.30379) (2.22990,0.30194) (2.23126,0.30008) (2.23262,0.29823) (2.23398,0.29638) (2.23534,0.29452) (2.23670,0.29267) (2.23806,0.29081) (2.23942,0.28896) (2.24079,0.28710) (2.24215,0.28525) (2.24351,0.28339) (2.24487,0.28154) (2.24623,0.27968) (2.24759,0.27783) (2.24895,0.27597) (2.25032,0.27412) (2.25168,0.27226) (2.25304,0.27040) (2.25440,0.26855) (2.25576,0.26669) (2.25712,0.26483) (2.25848,0.26298) (2.25985,0.26112) (2.26121,0.25926) (2.26257,0.25740) (2.26393,0.25555) (2.26529,0.25369) (2.26665,0.25183) (2.26801,0.24997) (2.26937,0.24811) (2.27074,0.24625) (2.27210,0.24440) (2.27346,0.24254) (2.27482,0.24068) (2.27618,0.23882) (2.27754,0.23696) (2.27890,0.23510) (2.28027,0.23324) (2.28163,0.23138) (2.28299,0.22952) (2.28435,0.22766) (2.28571,0.22580) (2.28707,0.22394) (2.28843,0.22208) (2.28979,0.22022) (2.29116,0.21836) (2.29252,0.21650) (2.29388,0.21464) (2.29524,0.21278) (2.29660,0.21091) (2.29796,0.20905) (2.29932,0.20719) (2.30069,0.20533) (2.30205,0.20347) (2.30341,0.20161) (2.30477,0.19974) (2.30613,0.19788) (2.30749,0.19602) (2.30885,0.19416) (2.31022,0.19229) (2.31158,0.19043) (2.31294,0.18857) (2.31430,0.18671) (2.31566,0.18484) (2.31702,0.18298) (2.31838,0.18112) (2.31974,0.17925) (2.32111,0.17739) (2.32247,0.17553) (2.32383,0.17366) (2.32519,0.17180) (2.32655,0.16993) (2.32791,0.16807) (2.32927,0.16621) (2.33064,0.16434) (2.33200,0.16248) (2.33336,0.16061) (2.33472,0.15875) (2.33608,0.15688) (2.33744,0.15502) (2.33880,0.15315) (2.34016,0.15129) (2.34153,0.14942) (2.34289,0.14756) (2.34425,0.14569) (2.34561,0.14383) (2.34697,0.14196) (2.34833,0.14010) (2.34969,0.13823) (2.35106,0.13637) (2.35242,0.13450) (2.35378,0.13263) (2.35514,0.13077) (2.35650,0.12890) (2.35786,0.12704) (2.35922,0.12517) (2.36059,0.12330) (2.36195,0.12144) (2.36331,0.11957) (2.36467,0.11770) (2.36603,0.11584) (2.36739,0.11397) (2.36875,0.11210) (2.37011,0.11024) (2.37148,0.10837) (2.37284,0.10650) (2.37420,0.10464) (2.37556,0.10277) (2.37692,0.10090) (2.37828,0.09903) (2.37964,0.09717) (2.38101,0.09530) (2.38237,0.09343) (2.38373,0.09156) (2.38509,0.08970) (2.38645,0.08783) (2.38781,0.08596) (2.38917,0.08409) (2.39053,0.08223) (2.39190,0.08036) (2.39326,0.07849) (2.39462,0.07662) (2.39598,0.07475) (2.39734,0.07289) (2.39870,0.07102) (2.40006,0.06915) (2.40143,0.06728) (2.40279,0.06541) (2.40415,0.06355) (2.40551,0.06168) (2.40687,0.05981) (2.40823,0.05794) (2.40959,0.05607) (2.41096,0.05420) (2.41232,0.05234) (2.41368,0.05047) (2.41504,0.04860) (2.41640,0.04673) (2.41776,0.04486) (2.41912,0.04299) (2.42048,0.04112) (2.42185,0.03926) (2.42321,0.03739) (2.42457,0.03552) (2.42593,0.03365) (2.42729,0.03178) (2.42865,0.02991) (2.43001,0.02804) (2.43138,0.02618) (2.43274,0.02431) (2.43410,0.02244) (2.43546,0.02057) (2.43682,0.01870) (2.43818,0.01683) (2.43954,0.01496) (2.44091,0.01309) (2.44227,0.01122) (2.44363,0.00936) (2.44499,0.00749) (2.44635,0.00562) (2.44771,0.00375) (2.44907,0.00188) (2.45043,0.00001)};
\draw[flow,dashed] plot coordinates{(0.00000,1.90000) (0.02450,1.89977) (0.04901,1.89906) (0.07351,1.89789) (0.09802,1.89625) (0.12252,1.89414) (0.14703,1.89157) (0.17153,1.88853) (0.19604,1.88502) (0.22054,1.88104) (0.24504,1.87661) (0.26955,1.87171) (0.29405,1.86635) (0.31856,1.86052) (0.34306,1.85424) (0.36757,1.84750) (0.39207,1.84031) (0.41658,1.83266) (0.44108,1.82456) (0.46558,1.81601) (0.49009,1.80701) (0.51459,1.79756) (0.53910,1.78767) (0.56360,1.77734) (0.58811,1.76658) (0.61261,1.75537) (0.63711,1.74373) (0.66162,1.73167) (0.68612,1.71917) (0.71063,1.70625) (0.73513,1.69291) (0.75964,1.67915) (0.78414,1.66498) (0.80865,1.65040) (0.83315,1.63541) (0.85765,1.62002) (0.88216,1.60422) (0.90666,1.58803) (0.93117,1.57145) (0.95567,1.55448) (0.98018,1.53713) (1.00468,1.51940) (1.02919,1.50129) (1.05369,1.48282) (1.07819,1.46398) (1.10270,1.44477) (1.12720,1.42521) (1.15171,1.40530) (1.17621,1.38504) (1.20072,1.36444) (1.22522,1.34350) (1.24973,1.32223) (1.27423,1.30064) (1.29873,1.27872) (1.32324,1.25649) (1.34774,1.23395) (1.37225,1.21111) (1.39675,1.18796) (1.42126,1.16452) (1.44576,1.14080) (1.47027,1.11679) (1.49477,1.09251) (1.51927,1.06796) (1.54378,1.04314) (1.56828,1.01807) (1.59279,0.99275) (1.61729,0.96718) (1.64180,0.94137) (1.66630,0.91533) (1.69081,0.88907) (1.71531,0.86258) (1.73981,0.83588) (1.76432,0.80898) (1.78882,0.78188) (1.81333,0.75458) (1.83783,0.72710) (1.86234,0.69944) (1.88684,0.67160) (1.91134,0.64360) (1.93585,0.61544) (1.96035,0.58713) (1.98486,0.55868) (2.00936,0.53008) (2.03387,0.50136) (2.05837,0.47251) (2.08288,0.44355) (2.10738,0.41447) (2.13188,0.38530) (2.15639,0.35602) (2.18089,0.32667) (2.20540,0.29723) (2.22990,0.26771) (2.25441,0.23813) (2.27891,0.20850) (2.30342,0.17881) (2.32792,0.14907) (2.35242,0.11930) (2.37693,0.08950) (2.40143,0.05968) (2.42594,0.02984) (2.45044,0.00000)};
\fill(2.4504422698000385,0)circle(1.2pt)node[above right]{$S$};
\fill(0,2.3762941943399194)circle(1.2pt)node[above right]{$P_c$};
\node[left]at(0,1.9){$2$};\node[below]at(2.45,0){$\pi$};\node at(1.65,2.1){有阻尼};
\end{tikzpicture}
\caption{弱阻尼与无阻尼的上行分界支。虚线为无阻尼，实线取 $\beta/\omega_0=0.12$，由相方程积分得到；两支都只在无限时间趋近 S。}\label{mp4:fig:l4-h12-s}
\end{figure}

因临界上行过程中 $\dot\theta>0$，可将时间积分改为角度积分：$\dot\theta^2\dd t=\dot\theta\,\dd\theta$。耗散式前已有一阶小参数 $\beta$，为求首阶项，可将积分中的速度取为无阻尼分界支
\[
\dot\theta^{(0)}=2\omega_0\cos\frac\theta2,
\qquad 0\le\theta<\pi.
\]
于是
\begin{align*}
Q&=2I\beta\int_0^\pi\dot\theta\,\dd\theta\\
&=4I\beta\omega_0\int_0^\pi\cos\frac\theta2\,\dd\theta
+o(I\beta\omega_0)\\
&=4I\beta\omega_0\left[2\sin\frac\theta2\right]_0^\pi
+o(I\beta\omega_0)\\
&=8I\beta\omega_0+o(I\beta\omega_0).
\end{align*}
写 $\dot\theta_c=2\omega_0+\delta$，其中 $\delta=O(\beta)$。将初始动能展开，
\[
\frac I2(2\omega_0+\delta)^2
=2I\omega_0^2+2I\omega_0\delta+O(I\delta^2),
\]
比较一阶项得 $2I\omega_0\delta=8I\beta\omega_0+o(I\beta\omega_0)$，所以
\[
\boxed{\dot\theta_c=2\omega_0+4\beta+o(\beta).}
\]

还可求临界运动在中途各角度应具有的能量。记 $\epsilon=\beta/\omega_0$、$h=E/(I\omega_0^2)$，从当前角度到倒立位置的未来耗散为
\[
2I\beta\int_\theta^\pi\dot\theta^{(0)}\dd\theta
=4I\beta\omega_0\left[2\sin\frac u2\right]_\theta^\pi
=8I\beta\omega_0\left(1-\sin\frac\theta2\right).
\]
因此到一阶精度，
\[
h(\theta)=2+8\epsilon\left(1-\sin\frac\theta2\right)+o(\epsilon),
\]
如图\ref{mp4:fig:l4-h12-energy}。图中 $h-2$ 是之后仍会耗散的能量，而不是已经耗散的能量；它从释放时的 $8\epsilon$ 单调减到零。

\begin{figure}[H]\centering
\begin{tikzpicture}[x=1.65cm,y=4cm]
\draw[axis](-.1,0)--(3.5,0)node[right]{$\theta$};
\draw[axis](0,-.06)--(0,.55)node[above]{$h-2$};
\draw[curve,domain=0:pi,samples=100]plot(\x,{.4*(1-sin(\x*90/pi))});
\fill(0,.4)circle(.8pt)node[above right]{$P_c$};
\fill(pi,0)circle(.8pt)node[above right]{$S$};
\node[left]at(0,.4){$8\epsilon$};\node[below]at(pi,0){$\pi$};
\node[below left]at(0,0){$0$};
\draw[guide](0,.4)--(1.1,.4);
\node at(1.9,.4){$h(\theta)-2$};
\end{tikzpicture}
\caption{弱阻尼临界运动的剩余耗散能量，按一阶式作图，取 $\epsilon=0.05$。$P_c$ 对应最低点的临界初态，S 对应倒立平衡；摆球上升时，剩余耗散能量逐渐减到零。}\label{mp4:fig:l4-h12-energy}
\end{figure}

耗散积分的收敛性还可由倒立位置附近的时间解看出。令 $\delta_\theta=\pi-\theta$，倒立位置附近满足
\[
\ddot\delta_\theta+2\beta\dot\delta_\theta-\omega_0^2\delta_\theta\simeq0.
\]
临界轨线只保留衰减根 $-\beta-\sqrt{\beta^2+\omega_0^2}$，角度差和速度都指数衰减，$\int\dot\theta^2\dd t$ 有限。所求等号仍只给出渐近倒立位置的初态，实际越过需严格超过该临界值；一阶公式适用于 $\beta/\omega_0\ll1$。
\end{solution}

\chapter{自由度和能量解法}

约束使各物体的位置相互关联，整个系统常可用少数几个坐标描述。由这些坐标写出动能、势能和外力做功，就能消去许多约束力。所得运动仍须满足约束成立的条件，例如绳子绷紧、支持力非负，以及静摩擦力不超过上限。


\section{自由度与坐标选择}
独立坐标由几何约束的秩确定，特殊构形须另作讨论。各质点位置写成独立坐标的函数后，动能中的速度关联和可消去变量才有明确依据。

\subsection{约束方程与独立坐标}
自由质点在三维空间有三个位置自由度；自由刚体还须确定姿态，因而有三个平动、三个转动自由度，共六个。平面刚体只需两个平动坐标和一个转角，共三个自由度。用刚体结点的多个坐标表示同一物体时，点间距离约束已经包含其中，不能再次把每个结点当成独立质点重复计数。
设系统原来用 $n$ 个位置变量 $x_1,\ldots,x_n$ 描述，有 $r$ 个独立的几何约束
\[
f_\alpha(x_1,\ldots,x_n,t)=0,\qquad \alpha=1,\ldots,r.
\]
在约束方程的雅可比矩阵秩等于 $r$ 的普通构形附近，自由度为 $n-r$。约束是否独立，须由方程间的关系判断。同一根绳的长度方程重复列写，不会进一步减少自由度；机构到达特殊构形时，约束矩阵的秩也可能改变。

若选取 $s=n-r$ 个广义坐标 $q_1,\ldots,q_s$，各质点位置写成 $\vct r_i=\vct r_i(q_1,\ldots,q_s,t)$，则
\begin{equation}\label{mp5:eq:l5-velocity}
\vct v_i=\sum_j\frac{\partial\vct r_i}{\partial q_j}\dot q_j+
\frac{\partial\vct r_i}{\partial t}.
\end{equation}
第一项来自独立坐标的变化，第二项来自约束本身的运动。在固定时刻改变构形，虚位移为
\begin{equation}\label{mp5:eq:l5-virtual}
\delta\vct r_i=\sum_j\frac{\partial\vct r_i}{\partial q_j}\delta q_j,
\qquad\delta t=0.
\end{equation}
所以动约束下的虚位移一般不等于 $\vct v_i\delta t$。例如珠子套在转动直杆上，固定某一时刻后允许的虚位移沿杆，而实际速度还包括杆带来的横向分量。


\begin{figure}[H]\centering
\begin{tikzpicture}
\draw[curve,domain=-1.6:1.6,samples=80]plot(\x,{.35*(\x)^2});\fill(.8,.224)circle(1.4pt)node[above left]{$P$};\draw[axis](.8,.224)--(1.8,.784)node[right]{$\delta\vct r$};\node at(-.9,1.5){$f(x,y)=0$};\begin{scope}[shift={(6,0)},scale=1]
\draw[curve](-1.3,0)--(1.7,0);\draw[guide](-1.3,-.5)--(1.7,.654);\fill(0,0)circle(1.3pt)node[below]{$O$};\fill(1.2,0)circle(1.4pt)node[below]{$P$};\draw[axis](1.2,0)--(2.2,0)node[right]{$\delta\vct r$};\draw[axis](1.2,0)--(1.2,.9)node[right]{$\Omega r$};\node at(.1,1.5){固定 $t$：沿瞬时杆向};\end{scope}

\end{tikzpicture}
\caption{左图中允许虚位移沿约束曲线的切线；右图在固定时刻取定瞬时杆向，实际速度还包含约束转动产生的横向分量。}\label{mp5:fig:l5-tangent}
\end{figure}

\subsection{从动能得到惯性系数}
固定约束满足 $\partial\vct r_i/\partial t=0$，动能可以写成二次型
\begin{equation}\label{mp5:eq:l5-matrix}
T=\frac12\sum_{j,k}M_{jk}(\vct q)\dot q_j\dot q_k,
\qquad
M_{jk}=\sum_i m_i\frac{\partial\vct r_i}{\partial q_j}\cdot
\frac{\partial\vct r_i}{\partial q_k}.
\end{equation}
刚体还要计入绕质心的转动动能。交叉项来自不同广义速度对同一质点速度的贡献。坐标独立，并不要求这些速度贡献彼此正交。坐标没有冗余且处在普通构形时，动能二次型正定。

对于理想约束，约束力对允许虚位移的总功为零。将牛顿方程与虚位移作内积、求和，约束力因而从广义运动方程中消去。惯性项可用动能表示：先由式\eqref{mp5:eq:l5-velocity}得到
\[
\frac{\partial\vct v_i}{\partial\dot q_j}
=\frac{\partial\vct r_i}{\partial q_j},\qquad
\frac{\dd}{\dd t}\frac{\partial\vct r_i}{\partial q_j}
=\sum_k\frac{\partial^2\vct r_i}{\partial q_j\partial q_k}\dot q_k
+\frac{\partial^2\vct r_i}{\partial q_j\partial t}
=\frac{\partial\vct v_i}{\partial q_j}.
\]
因此对 $\partial T/\partial\dot q_j=\sum_i m_i\vct v_i\cdot\partial\vct r_i/\partial q_j$ 求时间导数，
\[
\begin{aligned}
\frac{\dd}{\dd t}\frac{\partial T}{\partial\dot q_j}
&=\sum_i m_i\vct a_i\cdot\frac{\partial\vct r_i}{\partial q_j}
+\sum_i m_i\vct v_i\cdot\frac{\dd}{\dd t}\frac{\partial\vct r_i}{\partial q_j}\\
&=\sum_i m_i\vct a_i\cdot\frac{\partial\vct r_i}{\partial q_j}
+\sum_i m_i\vct v_i\cdot\frac{\partial\vct v_i}{\partial q_j}.
\end{aligned}
\]
最后一项正是 $\partial T/\partial q_j$。将它移项，得到
\[
\sum_i m_i\vct a_i\cdot\frac{\partial\vct r_i}{\partial q_j}
=\frac{\dd}{\dd t}\frac{\partial T}{\partial\dot q_j}
-\frac{\partial T}{\partial q_j},
\]
代入广义运动方程，即得拉格朗日方程
\begin{equation}\label{mp5:eq:l5-lagrange}
\frac{\dd}{\dd t}\frac{\partial L}{\partial\dot q_j}
-\frac{\partial L}{\partial q_j}=Q_j^{\rm nc},
\qquad L=T-U.
\end{equation}
其中 $Q_j^{\rm nc}=\sum_i\vct F_i^{\rm nc}\cdot\partial\vct r_i/\partial q_j$ 是未收入势能的广义力。这一形式也适用于显含时间的理想约束，但此时不能自动认定惯性系机械能守恒。

若自由度超过一，单独一个能量积分通常不能确定全部运动。还须使用其余运动方程，或找到动量、角动量等额外积分。若 $L$ 不含某个坐标 $q_j$ 本身、且 $Q_j^{\rm nc}=0$，则相应动量 $p_j=\partial L/\partial\dot q_j$ 守恒；这种坐标称为循环坐标。

\par\medskip\noindent\begin{minipage}{\textwidth}
\begin{example}[复合阿特伍德装置的两个自由度]\label{mp5:example:l5-e1}
图\ref{mp5:fig:l5-e1-q}中两根轻绳分别跨过定滑轮 O 和动滑轮 C。定滑轮绳的另一端悬挂 $3m$，动滑轮两侧悬挂 $m$、$2m$。滑轮质量及摩擦忽略，绳始终绷紧。竖直向下为正，略去不影响速度的固定长度后，取 $y_C=q$、$y_1=q+s$、$y_2=q-s$。求三个重物的加速度，并用张力关系验证。
\end{example}

\begin{figure}[H]\centering
\begin{tikzpicture}
\draw[curve](-2.4,3.5)--(3,3.5);
\draw[fill=white,curve](0,2.8)circle(0.45);\fill(0,2.8)circle(1pt);
\node[above right]at(0.04,2.84){$O$};
\draw(0,2.8)--(0,3.5);\draw[fill=white,curve](0.45,0.65)circle(0.35);\fill(0.45,0.65)circle(1pt);
\node[above right]at(0.49,0.6900000000000001){$C$};
\draw[curve](-.45,.875)--(-.45,2.8)arc(180:0:.45)--(.45,.65);\draw[curve](.1,-1.05)--(.1,.65)arc(180:0:.35)--(.8,-.55);\draw[fill=white](-0.75,0.32499999999999996)rectangle(-0.15000000000000002,0.875);\node at(-0.45,0.6){$3m$};
\draw[fill=white](-0.19999999999999998,-1.6)rectangle(0.4,-1.0499999999999998);\node at(0.1,-1.325){$m$};
\draw[fill=white](0.5,-1.1)rectangle(1.1,-0.5499999999999999);\node at(0.8,-0.825){$2m$};
\draw[axis](2.6,2.8)--(2.6,-2)node[below]{$y$};\draw[guide](.8,.65)--(1.65,.65);\draw[<->,>=Stealth](1.55,2.8)--(1.55,.65)node[midway,right]{$q$};\node[left]at(2.6,2.8){$0$};\node at(-1.7,-.25){$y_0$};\node[left]at(-.05,-1.85){$y_1=q+s$};\node at(1,-2.7){$y_2=q-s$};
\end{tikzpicture}
\caption{复合滑轮与向下的位置坐标。坐标关系中略去了固定长度；两根绳分处不同轴向平面，图中作投影。}\label{mp5:fig:l5-e1-q}
\end{figure}
\end{minipage}\par\medskip
\begin{solution}
定滑轮绳长给出 $y_0=-q+\text{常量}$，下方绳长使 $y_1+y_2=2q+\text{常量}$。两个独立坐标为 $q,s$。下绳两端相对 C 的运动相反，但相对地面的速度还包含 C 的平动。

将三个重物的动能和重力势能分别相加，
\begin{align*}
T&=\frac12(3m)\dot q^2+\frac12m(\dot q+\dot s)^2
+\frac12(2m)(\dot q-\dot s)^2\\
&=\frac m2(6\dot q^2-2\dot q\dot s+3\dot s^2),\\
U&=-g(3my_0+my_1+2my_2)=mg s+\text{常量}.
\end{align*}
重力在 $q$ 方向的广义力为零：$q$ 增大时，$3m$ 上升，$m$ 与 $2m$ 一起下降，三项重力虚功正好相消。$s$ 增大时，$m$ 下降而 $2m$ 上升，净重力虚功为 $-mg\,\delta s$。势能只含 $s$，重力便通过这一相对坐标驱动整个系统。

将求导写出，动量与势能梯度分别为
\[
\frac{\partial T}{\partial\dot q}=m(6\dot q-\dot s),\qquad
\frac{\partial T}{\partial\dot s}=m(-\dot q+3\dot s),
\qquad U_q=0,\quad U_s=mg.
\]
交叉项的系数 $-2m$ 来自下方两重物不同的质量及相反的相对速度；它使 $q,s$ 两个坐标的加速度耦合。

对 $q,s$ 分别使用式\eqref{mp5:eq:l5-lagrange}，得到下列方程；相应加速度方向见图\ref{mp5:fig:l5-e1-s}。
\[
6\ddot q-\ddot s=0,\qquad-\ddot q+3\ddot s=-g.
\]

\begin{figure}[H]\centering
\begin{tikzpicture}
\draw[curve](-2.4,3.5)--(3,3.5);
\draw[fill=white,curve](0,2.8)circle(0.45);\fill(0,2.8)circle(1pt);
\node[above right]at(0.04,2.84){$O$};
\draw(0,2.8)--(0,3.5);\draw[fill=white,curve](0.45,0.65)circle(0.35);\fill(0.45,0.65)circle(1pt);
\node[above right]at(0.49,0.6900000000000001){$C$};
\draw[curve](-.45,.875)--(-.45,2.8)arc(180:0:.45)--(.45,.65);\draw[curve](.1,-1.05)--(.1,.65)arc(180:0:.35)--(.8,-.55);\draw[fill=white](-0.75,0.32499999999999996)rectangle(-0.15000000000000002,0.875);\node at(-0.45,0.6){$3m$};
\draw[fill=white](-0.19999999999999998,-1.6)rectangle(0.4,-1.0499999999999998);\node at(0.1,-1.325){$m$};
\draw[fill=white](0.5,-1.1)rectangle(1.1,-0.5499999999999999);\node at(0.8,-0.825){$2m$};
\draw[axis](2.6,2.8)--(2.6,-2)node[below]{$y$};\draw[axis](-1.45,.6)--(-1.45,-.05)node[left]{$g/17$};\draw[axis](-.2,-1.3)--(-.2,-.25);\node[left]at(-.2,-1.7){$7g/17$};\draw[axis](1.35,-.8)--(1.35,-1.55)node[below]{$5g/17$};
\end{tikzpicture}
\caption{三个重物相对地面的加速度：$3m$ 与 $2m$ 向下，$m$ 向上。箭头只表示方向。}\label{mp5:fig:l5-e1-s}
\end{figure}

由第一式 $\ddot s=6\ddot q$，代入第二式得
\[
\ddot q=-\frac g{17},\qquad\ddot s=-\frac{6g}{17}.
\]
故向下为正的三个加速度为
\[
\boxed{\ddot y_0=\frac g{17},\qquad
\ddot y_1=-\frac{7g}{17},\qquad
\ddot y_2=\frac{5g}{17}.}
\]
设下绳张力为 $T_1$，由 $m$、$2m$ 两物体分别得到
\[
T_1=m\left(g+\frac{7g}{17}\right)
=2m\left(g-\frac{5g}{17}\right)=\frac{24}{17}mg.
\]
无质量动滑轮要求上绳张力 $T_0=2T_1=48mg/17$。于是 $3mg-T_0=3m(g/17)$，与左侧重物的加速度相符。两个张力均为正，满足题设的绷紧条件。

还可将方程写成质量矩阵形式：
\[
\begin{pmatrix}6&-1\\-1&3\end{pmatrix}
\begin{pmatrix}\ddot q\\\ddot s\end{pmatrix}
=\begin{pmatrix}0\\-g\end{pmatrix},\qquad
\begin{pmatrix}6&-1\\-1&3\end{pmatrix}^{-1}
=\frac1{17}\begin{pmatrix}3&1\\1&6\end{pmatrix}.
\]
矩阵行列式为 17，两个坐标确实独立。图\ref{mp5:fig:l5-e1-s}的加速度还满足 $\ddot y_0=-\ddot q$、$\ddot y_1+\ddot y_2=2\ddot q$。这些关系来自绳长，与三个重物各自的受力方程同时成立。上绳的 $3m$ 与下绳两端的总质量虽然相等，下方相对运动仍使 C 向上加速，因此不能先将动滑轮当成静止的定滑轮。

\end{solution}

\par\medskip\noindent\begin{minipage}{\textwidth}
\begin{example}[转动直杆上的弹簧滑珠]\label{mp5:example:l5-e2}
光滑直杆在水平面内绕固定点 $O$ 以恒定角速度 $\Omega\ge0$ 转动，滑珠质量为 $m$。以杆上的有向坐标 $q$ 表示珠子位置；弹簧对珠子的沿杆作用力为 $-kq$，自然长度按零处理，见图\ref{mp5:fig:l5-e2-q}。求 $q$ 的运动方程，判断何时有稳定振动，并区分机械能与守恒积分。
\end{example}

\begin{figure}[H]\centering
\begin{tikzpicture}
\draw[curve](-2,0)--(2.4,0);\fill(0,0)circle(1.3pt)node[below]{$O$};\draw[axis](0,-.55)--(2.5,-.55)node[right]{$q$};\draw[guide](0,-.1)--(0,-.65);\draw[fill=white](1.3,0)circle(.13)node[above right]{$m$};\draw[axis](1.3,0)--(.45,0)node[midway,above,yshift=3pt]{$-kq$};\draw[axis](-.75,0)arc(180:240:.75);\node at(-.7,-.9){$\Omega$};\node at(-1.1,.7){俯视图};
\end{tikzpicture}
\caption{水平面内转动直杆的俯视图。$q$ 是可正可负的有向坐标，滑珠可通过 $O$。}\label{mp5:fig:l5-e2-q}
\end{figure}
\end{minipage}\par\medskip
\begin{solution}
令沿杆单位向量为 $\vct e(t)$，横向单位向量为 $\vct e_\perp=\vct e_z\times\vct e$，则 $\vct r=q\vct e$，
\[
\vct v=\dot q\vct e+\Omega q\vct e_\perp,
\qquad
T=\frac m2(\dot q^2+\Omega^2q^2),\qquad U=\frac12kq^2.
\]
由 $\dot{\vct e}=\Omega\vct e_\perp$、$\dot{\vct e}_\perp=-\Omega\vct e$，再求一次时间导数，
\[
\vct a=(\ddot q-\Omega^2q)\vct e+2\Omega\dot q\vct e_\perp.
\]
沿杆方向只有弹簧力，故 $m(\ddot q-\Omega^2q)=-kq$；横向则由导杆提供 $2m\Omega\dot q$。由此可见，$\ddot q$ 是杆上有向坐标的二阶导数，实际沿杆加速度为 $\ddot q-\Omega^2q$；二者的差别产生了方程中的 $m\Omega^2q$ 项。

对 $q$ 使用式\eqref{mp5:eq:l5-lagrange}，
\begin{equation}\label{mp5:eq:l5-rotating-rod}
m\ddot q=(m\Omega^2-k)q.
\end{equation}
当 $\Omega<\omega_s=\sqrt{k/m}$ 时，角频率为 $\sqrt{\omega_s^2-\Omega^2}$；等号时 $\ddot q=0$，一般初态作沿杆匀速运动；$\Omega>\omega_s$ 时有指数增长模态。

\begin{figure}[H]\centering
\begin{tikzpicture}
\begin{scope}[shift={(0,0)},scale=0.9]
\draw[curve](-1.5,0)--(2.3,0);\fill(.4,0)circle(1.5pt)node[below]{$P$};\draw[axis](.4,0)--(1.8,0)node[below]{$\dot q\vct e$};\draw[axis](.4,0)--(.4,1.2)node[left]{$\Omega q\vct e_\perp$};\draw[axis](.4,0)--(1.8,1.2)node[above]{$\vct v$};\draw[guide](.4,1.2)--(1.8,1.2)--(1.8,0);\draw[axis,line width=.8pt](.4,0)--(.4,1.8)node[right]{$N_\perp$};\end{scope}
\begin{scope}[shift={(6.2,0)},scale=1]
\node[align=center]at(0,1){$L_z=m\Omega q^2$\\[5pt]$\tau_z=2m\Omega q\dot q$};\node[align=center]at(0,-.2){$P=\Omega\tau_z$\\[5pt]$\dot E_{\rm in}=P$};\end{scope}

\end{tikzpicture}
\caption{取 $q>0,\dot q>0$。左图在珠子处画出速度分量与横向作用力，右图列出导杆力矩与功率的关系。}\label{mp5:fig:l5-e2-s}
\end{figure}

图\ref{mp5:fig:l5-e2-s}中实际速度有横向分量，导杆给珠子的横向力为 $N_\perp=2m\Omega\dot q$。它对瞬时沿杆虚位移不做功，却对实际速度做功，功率为
\[
P=N_\perp\Omega q=2m\Omega^2q\dot q.
\]
惯性系机械能
\[
E_{\rm in}=\frac m2\dot q^2+\frac12(k+m\Omega^2)q^2
\]
对它求导，并用式\eqref{mp5:eq:l5-rotating-rod}消去 $\ddot q$，
\[
\begin{aligned}
\dot E_{\rm in}
&=m\dot q\ddot q+(k+m\Omega^2)q\dot q\\
&=(m\Omega^2-k)q\dot q+(k+m\Omega^2)q\dot q
=2m\Omega^2q\dot q=P.
\end{aligned}
\]
所以机械能一般不守恒。将动力学方程写为 $m\ddot q+(k-m\Omega^2)q=0$，再乘以 $\dot q$，得到 $\dd[m\dot q^2/2+(k-m\Omega^2)q^2/2]/\dd t=0$，故守恒的是
\[
\boxed{J=\frac m2\dot q^2+\frac12(k-m\Omega^2)q^2.}
\]
维持恒定 $\Omega$ 的装置补给或吸收机械能。也可从角动量 $L_z=m\Omega q^2$ 得出：装置力矩为 $\dot L_z$，功率 $\Omega\dot L_z=2m\Omega^2q\dot q$。
若给定 $q(0)=q_0$、$\dot q(0)=u_0$，稳定区间内的特征根为 $\pm\mathrm i\omega$。通解 $q=C\cos\omega t+D\sin\omega t$ 的初值为 $C=q_0$、$\omega D=u_0$，所以
\[
q=q_0\cos\omega t+\frac{u_0}{\omega}\sin\omega t,
\qquad \omega=\sqrt{k/m-\Omega^2}.
\]
临界时 $\ddot q=0$，积分得 $\dot q=u_0$、$q=q_0+\int_0^t u_0\dd s=q_0+u_0t$。超临界时，令 $\lambda=\sqrt{\Omega^2-k/m}$，特征根为 $\pm\lambda$，通解 $q=C_1\ee^{\lambda t}+C_2\ee^{-\lambda t}$ 的初值满足
\[
C_1+C_2=q_0,\qquad \lambda(C_1-C_2)=u_0.
\]
相加、相减后，$C_1=(q_0+u_0/\lambda)/2$、$C_2=(q_0-u_0/\lambda)/2$。一般扰动含增长项，只有 $u_0=-\lambda q_0$ 的特殊初态完全消去这一项。

机械能与守恒积分之间满足
\[
J=E_{\rm in}-\Omega L_z,
\qquad \frac{\dd}{\dd t}(E_{\rm in}-\Omega L_z)=P-\Omega\dot L_z=0.
\]
固定角速度的驱动装置与珠子交换能量和角动量，所保持的是上式的组合。图\ref{mp5:fig:l5-e2-s}中的横向力垂直于允许虚位移，却不垂直于实际速度，因而这两个功的判断必须使用各自的位移。

\end{solution}

\par\medskip\noindent\begin{minipage}{\textwidth}
\begin{example}[同一刚杆的两种坐标]\label{mp5:example:l5-e3}
均匀刚杆 AB 长 $L$、质量 $m$，两端装在相互垂直的光滑双侧直槽内，图\ref{mp5:fig:l5-e3-q}取 $x_A>0,y_B>0$ 的构形。令杆与竖直方向夹角为 $\theta\in(0,\pi/2)$。求以 $\theta$ 为坐标的动能和运动方程，再改用 $x=x_A$，解释接近水平位置时惯性系数的变化。
\end{example}

\begin{figure}[H]\centering
\begin{tikzpicture}
\draw[curve](-.35,0)--(3.5,0);\draw[curve](0,-.35)--(0,3.4);\draw[line width=1pt](3,0)--(0,3);\fill(3,0)circle(1.3pt)node[below right]{$A$};
\fill(0,3)circle(1.3pt)node[above left]{$B$};
\fill(0,0)circle(1.3pt)node[below left]{$O$};
\fill(1.5,1.5)circle(1.3pt)node[above right]{$G$};
\draw[guide](0,0)--(1.1,0);\draw[axis](0,2.4)arc(-90:-45:.6);\node at(.5,2.2){$\theta$};\node[above]at(.8,2.25){$L$};\draw[<->,>=Stealth](0,-.65)--(3,-.65)node[midway,below]{$x=L\sin\theta$};\draw[guide](3,0)--(3,-.75);
\end{tikzpicture}
\caption{双侧直槽限制 $A,B$ 的位置，两端反力均可随运动改变方向。}\label{mp5:fig:l5-e3-q}
\end{figure}
\end{minipage}\par\medskip
\begin{solution}
图\ref{mp5:fig:l5-e3-s}中
\[
A=(L\sin\theta,0),\quad B=(0,L\cos\theta),\quad
G=\left(\frac L2\sin\theta,\frac L2\cos\theta\right).
\]
对质心坐标求导，$\vct v_G=(L\dot\theta/2)(\cos\theta,-\sin\theta)$，因此 $v_G^2=(L^2\dot\theta^2/4)(\cos^2\theta+\sin^2\theta)=L^2\dot\theta^2/4$，而杆的角速度大小为 $|\dot\theta|$。质心平动与绕质心转动两部分都要计入：
\[
T=\frac12m v_G^2+\frac12\frac{mL^2}{12}\dot\theta^2
=\frac{mL^2}{6}\dot\theta^2,
\qquad U=\frac12mgL\cos\theta.
\]

\begin{figure}[H]\centering
\begin{tikzpicture}
\draw[curve](-.35,0)--(3.5,0);\draw[curve](0,-.35)--(0,3.4);\draw[line width=1pt](3,0)--(0,3);\fill(3,0)circle(1.3pt)node[below right]{$A$};
\fill(0,3)circle(1.3pt)node[above left]{$B$};
\fill(0,0)circle(1.3pt)node[below left]{$O$};
\fill(1.5,1.5)circle(1.3pt)node[above right]{$G$};
\draw[guide](0,0)--(1.1,0);\draw[axis](0,2.4)arc(-90:-45:.6);\node at(.5,2.2){$\theta$};\node[above]at(.8,2.25){$L$};\draw[axis](3,0)--(3.8,0)node[above]{$v_A$};\draw[axis](0,3)--(0,2.2)node[left]{$v_B$};\draw[axis](1.5,1.5)--(1.9,1.1)node[right]{$\vct v_G$};
\end{tikzpicture}
\caption{端点与质心的速度，取 $\dot\theta>0$。端点速度沿各自导槽，质心速度由两端速度的矢量平均得到。}\label{mp5:fig:l5-e3-s}
\end{figure}

拉格朗日方程为
\[
\frac{mL^2}{3}\ddot\theta-\frac{mgL}{2}\sin\theta=0,
\qquad\boxed{\ddot\theta=\frac{3g}{2L}\sin\theta.}
\]
在图\ref{mp5:fig:l5-e3-s}中，$\vct v_G=(\vct v_A+\vct v_B)/2$，但杆上各点速度并不相同，因此只计质心平动会漏掉 $mL^2\dot\theta^2/24$。以 $T=\tfrac12M_\theta\dot\theta^2$ 定义广义惯性系数，可得 $M_\theta=mL^2/3$。势能导数为 $U_\theta=-mgL\sin\theta/2$，所以拉格朗日方程中的惯性项与重力项分别来自这两个量。角加速度为正，指向 $\theta$ 增大的方向。若杆初速度也沿这一方向，它就加速趋向水平；若初速度相反，则先减速。

改用 $x=L\sin\theta$，有 $\dot\theta=\dot x/[L\sqrt{1-x^2/L^2}]$，故
\[
T=\frac{m\dot x^2}{6(1-x^2/L^2)},\qquad
U=\frac12mg\sqrt{L^2-x^2}.
\]
对应的惯性系数为 $M_x=m/[3(1-x^2/L^2)]$。当 $x\uparrow L$，$M_x$ 发散，同时 $\dot x=L\cos\theta\dot\theta$ 趋零；动能仍可保持有限。几何上，水平位置附近的 $x$ 对 $\theta$ 的一阶变化不敏感。这是坐标在该构形退化，不能据此认定刚杆的真实质量或转动惯量发散。若要跨过水平位置继续追踪运动，须换成在该处仍有效的坐标，并重新说明导槽允许的分支。
在 $0<x<L$ 内也可直接对新坐标列方程。所需导数为
\[
\begin{aligned}
M_x&=\frac{mL^2}{3(L^2-x^2)},\\
\frac{\dd M_x}{\dd x}
&=\frac{mL^2}{3}(-1)(L^2-x^2)^{-2}(-2x)
=\frac{2mL^2x}{3(L^2-x^2)^2},\\
U_x&=\frac{mg}{2}\frac{-2x}{2\sqrt{L^2-x^2}}
=-\frac{mgx}{2\sqrt{L^2-x^2}}.
\end{aligned}
\]
代入 $M_x\ddot x+\tfrac12M_x'\dot x^2+U_x=0$，整理得
\[
\ddot x+\frac{x\dot x^2}{L^2-x^2}
-\frac{3gx\sqrt{L^2-x^2}}{2L^2}=0.
\]
对 $x=L\sin\theta$ 直接求两次导数，
\[
\ddot x=L\cos\theta\ddot\theta-L\sin\theta\dot\theta^2
=\frac{3g}{2}\sin\theta\cos\theta-L\sin\theta\dot\theta^2.
\]
用 $\sin\theta=x/L$、$\cos\theta=\sqrt{L^2-x^2}/L$、$\dot\theta=\dot x/\sqrt{L^2-x^2}$ 代回后，也得到同一方程。

图\ref{mp5:fig:l5-e3-coordinate}画出位置映射和它的斜率。令 $\varepsilon=\pi/2-\theta$，则 $L-x\simeq L\varepsilon^2/2$、$\dot x\simeq L\varepsilon\dot\theta$，同时 $M_x\simeq m/(3\varepsilon^2)$。两者在动能中的乘积仍给出 $T\simeq mL^2\dot\theta^2/6$，坐标系数的发散并未产生额外能量。
\begin{figure}[H]\centering
\begin{tikzpicture}[scale=.95]
\draw[axis](-.1,0)--(3.5,0)node[right]{$\theta$};\draw[axis](0,-.1)--(0,1.8)node[above]{$x/L$};
\draw[curve,domain=0:90,samples=90]plot({3*\x/90},{1.4*sin(\x)});
\draw[guide](3,0)--(3,1.4)--(0,1.4);\node[left]at(0,1.4){$1$};\node[below]at(3,0){$\pi/2$};
\begin{scope}[shift={(5.5,0)}]
\draw[axis](-.1,0)--(3.5,0)node[right]{$\theta$};\draw[axis](0,-.1)--(0,1.8)node[above]{$L^{-1}\dd x/\dd\theta$};
\draw[curve,domain=0:90,samples=90]plot({3*\x/90},{1.4*cos(\x)});
\node[left]at(0,1.4){$1$};\node[below]at(3,0){$\pi/2$};\fill(3,0)circle(1.2pt);
\end{scope}
\end{tikzpicture}
\caption{位置映射 $x/L=\sin\theta$ 及其斜率。水平构形处斜率为零，$x$ 因而不能继续区分两侧的角度分支。}\label{mp5:fig:l5-e3-coordinate}
\end{figure}

\end{solution}

\section{虚功与静力平衡}
虚功采用同一时刻满足约束的允许位移。各力作用点的位移由约束联系，约束反力是否退出虚功方程则取决于理想约束条件。

\subsection{固定时刻的允许虚位移}
虚位移是相邻允许构形之间的微小变化，取 $\delta t=0$，不必与实际运动方向相同。对固定的光滑曲线，支持力垂直于允许虚位移，因此不做虚功；理想铰链、无伸长轻杆等约束的内力，在对整个系统求和时也可以消去。

若系统有 $s$ 个独立双侧坐标，静力平衡要求
\begin{equation}\label{mp5:eq:l5-vwork}
\delta W=\sum_i\vct F_i\cdot\delta\vct r_i
=\sum_{j=1}^s Q_j\delta q_j=0.
\end{equation}
由于各 $\delta q_j$ 可独立取正、负，故每个 $Q_j=0$。力来自势能时，$Q_j=-\partial U/\partial q_j$，平衡条件就是势能对独立坐标驻定。


\begin{figure}[H]\centering
\begin{tikzpicture}
\draw[curve,domain=-1.6:1.8,samples=90]plot(\x,{.35*(\x)^2});\fill(.8,.224)circle(1.4pt)node[below left]{$P$};\draw[axis](.8,.224)--(1.8,.784)node[right]{$\delta\vct r$};\draw[axis](.8,.224)--(.24,1.224)node[left]{$N$};\draw[axis](.8,.224)--(.8,-.7)node[right]{$\vct F$};\node at(-1.1,1.6){$N\,\delta r\cos90^\circ=0$};
\end{tikzpicture}
\caption{光滑约束的支持力 $N$ 垂直于允许虚位移；外力的切向分量进入虚功方程。}\label{mp5:fig:l5-virtual-normal}
\end{figure}

\subsection{稳定性与单侧接触}
对固定约束下的保守系统，若所用坐标有效且动能二次型正定，势能的严格局部极小点给出稳定平衡。若二阶项非退化，单自由度用 $U''(q_e)>0$ 判断；多自由度须检查二阶型 $\delta^2U$ 在所有允许方向上是否正定，不能只分别检查若干对角元素。

若有硬挡块 $q\ge0$，在 $q=0$ 处只允许 $\delta q\ge0$。边界平衡不必满足 $U'(0)=0$，而应先检查单侧变化 $\delta U\ge0$，并求出挡块反力是否非负。对保守、无摩擦的单侧约束，除接触反力外，其余力的虚功相应满足 $\delta W\le0$。这些不等式只取允许脱离接触的方向。

摩擦力通常做虚功。物体虽在实际平衡时没有位移，比较一个可能的滑动构形仍须计入摩擦力的虚功。临界状态可按将要发生的运动确定摩擦方向，再用 $|f|\le\mu N$ 判断；没有达到临界时，不应提前令 $f=\mu N$。

\par\medskip\noindent\begin{minipage}{\textwidth}
\begin{example}[光滑墙面与粗糙地面上的梯子]\label{mp5:example:l5-e4}
均匀梯子长 $L$、质量 $m$，与水平面的夹角为 $\theta$。墙面光滑，地面与梯子间的静摩擦系数为 $\mu$，见图\ref{mp5:fig:l5-e4-q}。用虚功法求平衡所需摩擦力和允许角度，并求墙面反力。
\end{example}

\begin{figure}[H]\centering
\begin{tikzpicture}
\draw[curve](0,-.35)--(0,3.9);\draw[curve](-.3,0)--(3.8,0);\draw[line width=1pt](3,0)--(0,3);\fill(3,0)circle(1.3pt)node[below right]{$A$};
\fill(0,3)circle(1.3pt)node[above left]{$B$};
\fill(1.5,1.5)circle(1.3pt)node[above right]{$G$};
\draw(2.4,0)arc(180:135:.6);\node at(2.25,.28){$\theta$};\node[above]at(1,2.2){$L$};\node at(3.1,-.6){$\mu$};\node[left]at(-.08,1.6){光滑};
\end{tikzpicture}
\caption{梯脚 $A$ 在粗糙地面上，顶端 $B$ 贴着光滑墙面，$\theta$ 为梯子与地面的夹角。}\label{mp5:fig:l5-e4-q}
\end{figure}
\end{minipage}\par\medskip
\begin{solution}
取 $\theta$ 为独立坐标，几何关系为
\[
x_A=L\cos\theta,\qquad y_G=\frac L2\sin\theta.
\]
在同一个固定时刻取任意允许的 $\delta\theta$，则
\[
\delta x_A=-L\sin\theta\,\delta\theta,\qquad
\delta y_B=L\cos\theta\,\delta\theta,\qquad
\delta y_G=\frac L2\cos\theta\,\delta\theta.
\]
若 $\delta\theta<0$，梯脚向外、梯顶向下，重心下降。这一几何趋势确定了平衡所需的摩擦方向。虚功方程中的 $\delta\theta$ 仍可取任意符号，它描述满足约束的试探位移。

图\ref{mp5:fig:l5-e4-s}中，顶端的允许位移竖直，墙面反力 $N_B$ 水平；梯脚的允许位移水平，地面支持力 $N_A$ 竖直。这两项虚功均为零。设地面对梯脚的摩擦力大小为 $f_A$，方向向左，则
\[
\delta W=-mg\,\delta y_G-f_A\,\delta x_A
=\left(-\frac{mgL}{2}\cos\theta+f_A L\sin\theta\right)\delta\theta.
\]

\begin{figure}[H]\centering
\begin{tikzpicture}
\draw[curve](0,-.35)--(0,3.9);\draw[curve](-.3,0)--(3.8,0);\draw[line width=1pt](3,0)--(0,3);\fill(3,0)circle(1.3pt)node[below right]{$A$};
\fill(0,3)circle(1.3pt)node[above left]{$B$};
\fill(1.5,1.5)circle(1.3pt)node[above right]{$G$};
\draw(2.4,0)arc(180:135:.6);\node at(2.25,.28){$\theta$};\node[above]at(1,2.2){$L$};\draw[axis](0,3)--(1.1,3)node[above]{$N_B$};\draw[axis](3,0)--(3,1.05)node[right]{$N_A$};\draw[axis](3,0)--(2,0)node[midway,below]{$f_A$};\draw[axis](1.5,1.5)--(1.5,.6)node[left]{$mg$};
\end{tikzpicture}
\caption{梯子受力：重力作用在 $G$，墙面反力作用在 $B$，地面支持力与摩擦力作用在 $A$。}\label{mp5:fig:l5-e4-s}
\end{figure}

由 $\delta W=0$ 得
\[
f_A=\frac{mg}{2}\cot\theta.
\]
竖直平衡给出 $N_A=mg$，水平平衡给出 $N_B=f_A$。因此维持静止须有
\[
\frac{mg}{2}\cot\theta\le\mu mg,
\qquad\boxed{\tan\theta\ge\frac1{2\mu}.}
\]
等号处梯脚将向右滑、梯顶将向下滑，摩擦力向左，与图中一致。若角度更大，所需摩擦力小于上限；这时把它直接写成 $\mu mg$ 会排除本来存在的平衡构形。
还可用绕梯脚 A 的力矩核对虚功结果：
\[
N_B L\sin\theta-\frac{mgL}{2}\cos\theta=0.
\]
由此得到 $N_B=mg\cot\theta/2$，再由水平平衡恢复 $f_A=N_B$。墙面支持力和地面支持力没有进入虚功方程，是因为它们与允许位移正交；求出摩擦力后，仍需由整体平衡方程恢复这两个反力。

当 $\mu>0$，最小允许角为 $\theta_{\min}=\arctan[1/(2\mu)]$。角度增大时 $\cot\theta$ 下降，所需静摩擦随之减小。若 $\mu=0$，任何 $0<\theta<\pi/2$ 的斜放构形都不能静止；竖直构形只能作为端点极限另行讨论。

\end{solution}

\par\medskip\noindent\begin{minipage}{\textwidth}
\begin{example}[剪式千斤顶的力与位移]\label{mp5:example:l5-e5}
图\ref{mp5:fig:l5-e5-q}的四根轻杆长度均为 $L$，构成对称菱形机构。下铰点 D 固定，左右铰点 A、B 由水平丝杠拉近，各受向内的力 $P$；顶点 C 承受向下荷载 $W$。设半宽为 $a$、半高为 $h$，$a^2+h^2=L^2$。忽略摩擦，求 P 与杆内力。
\end{example}

\begin{figure}[H]\centering
\begin{tikzpicture}
\draw[line width=.8pt](-1.5,1.5)--(0,3)--(1.5,1.5)--(0,0)--cycle;\fill(-1.5,1.5)circle(1.3pt)node[left]{$A$};
\fill(1.5,1.5)circle(1.3pt)node[right]{$B$};
\fill(0,3)circle(1.3pt)node[above left]{$C$};
\fill(0,0)circle(1.3pt)node[left]{$D$};
\draw(-.28,-.4)--(0,-.05)--(.28,-.4)--cycle;\draw(-.45,-.4)--(.45,-.4);\node[above left]at(-.8,2.4){$L$};\draw[axis](-1.5,1.5)--(-.45,1.5)node[midway,above]{$P$};\draw[axis](1.5,1.5)--(.45,1.5)node[midway,above]{$P$};\draw[axis](0,3)--(0,2.05)node[right]{$W$};\draw[guide](-1.5,1.5)--(-1.5,-.9);\draw[guide](1.5,1.5)--(1.5,-.9);\draw[<->,>=Stealth](-1.5,-.8)--(1.5,-.8)node[midway,below]{$2a$};\draw[guide](0,0)--(2.35,0);\draw[guide](0,3)--(2.35,3);\draw[<->,>=Stealth](2.2,0)--(2.2,3)node[midway,right]{$2h$};
\end{tikzpicture}
\caption{对称剪式机构。左右两端各受大小为 $P$ 的向内作用力，顶点高度为 $2h$。}\label{mp5:fig:l5-e5-q}
\end{figure}
\end{minipage}\par\medskip
\begin{solution}
以下取下铰点 D 为位置固定的原点，四个铰点的坐标为
\[
A=(-a,h),\quad B=(a,h),\quad C=(0,2h),\quad D=(0,0).
\]
四根杆的长度约束都归结为 $a^2+h^2=L^2$，故普通对称构形只有一个自由度。A、B 的竖直坐标也随 $a$ 改变，丝杠端点并不是沿固定水平直线运动。

由几何关系 $h=\sqrt{L^2-a^2}$，$\dd h/\dd a=(-2a)/(2\sqrt{L^2-a^2})=-a/h$，故
\[
\delta x_A=-\delta a,\quad\delta x_B=\delta a,
\quad\delta H=2\delta h=-\frac{2a}{h}\delta a.
\]
图\ref{mp5:fig:l5-e5-s}右图取 $\delta a<0$，$A,B$ 向内、向上移动，$C$ 上升，且 $\delta y_A=\delta y_B=\delta h>0$。两端力 $P$ 水平，其虚功只取决于水平位移分量；固定铰点 $D$ 的反力不做虚功。因此
\[
\delta W=P(-\delta a)-P\delta a-W\delta H
=\left(-2P+\frac{2Wa}{h}\right)\delta a=0.
\]

\begin{figure}[H]\centering
\begin{tikzpicture}
\begin{scope}[shift={(0,1)},scale=1]
\fill(0,0)circle(1.5pt)node[left]{$C$};\draw[axis](0,0)--(-.7,.7)node[above left]{$S$};\draw[axis](0,0)--(.7,.7)node[above right]{$S$};\draw[axis](0,0)--(0,-1.4)node[right]{$W$};\node at(0,-2){$2Sh/L=W$};\end{scope}
\begin{scope}[shift={(5,0)},scale=0.9]
\draw[line width=.8pt](-1.5,1.5)--(0,3)--(1.5,1.5)--(0,0)--cycle;\fill(-1.5,1.5)circle(1.3pt)node[left]{$A$};
\fill(1.5,1.5)circle(1.3pt)node[right]{$B$};
\fill(0,3)circle(1.3pt)node[above left]{$C$};
\fill(0,0)circle(1.3pt)node[left]{$D$};
\draw(-.28,-.4)--(0,-.05)--(.28,-.4)--cycle;\draw(-.45,-.4)--(.45,-.4);\node[above left]at(-.8,2.4){$L$};\draw[guide](-1.5,1.5)--(-1,1.5)--(-1,2);\draw[guide](1.5,1.5)--(1,1.5)--(1,2);
\draw[axis,line width=.8pt](-1.5,1.5)--(-1,2);\node[above left]at(-1.3,1.85){$\delta\vct r_A$};
\draw[axis,line width=.8pt](1.5,1.5)--(1,2);\node[above right]at(1.3,1.85){$\delta\vct r_B$};
\node[align=center]at(0,-1.05){$\delta x_A=-\delta a,\quad\delta x_B=\delta a$\\[3pt]$\delta y_A=\delta y_B=\delta h=\tfrac12\delta H>0$};\draw[axis](0,3)--(0,4)node[right]{$\delta H$};\node at(0,.8){$\delta a<0$};\end{scope}

\end{tikzpicture}
\caption{左图为顶点 $C$ 的受力；右图中 $A,B$ 的完整虚位移均斜向上，虚线表示水平、竖直分量，$\delta H=2\delta h$。}\label{mp5:fig:l5-e5-s}
\end{figure}

所以
\[
\boxed{P=W\frac ah.}
\]
设各杆的压缩内力大小为 $S$。图中顶端两杆的竖直分量平衡荷载：$2S(h/L)=W$，故
\[
\boxed{S=\frac{WL}{2h}.}
\]
在左侧铰点，两个杆力的水平分量合为 $2Sa/L$，也给出 $P=Wa/h$。$P$ 是丝杠对每一端施加的力；两端的功合计为 $P[-\delta(2a)]$，没有额外的二倍力。

当 $h\to0$，给定荷载所需的 $P$ 与杆内力都发散。$a\to0$ 时虽有 $P\to0$，机构却趋于共线构形；上述坐标和虚位移关系只用于 $a,h>0$ 的普通构形。
四根杆内力相等还可逐个结点证明。设 AC、AD 的压缩力分别为 $S_u,S_l$。在 A 点，丝杠没有竖直分力，竖直平衡要求 $(S_l-S_u)h/L=0$，因 $h>0$，故 $S_l=S_u$；左右对称再给出另两杆相同的内力。于是 C 点的 $2Sh/L=W$ 可以用于全部四杆，而 D 点的总竖直反力为 W。

图\ref{mp5:fig:l5-e5-s}右图中的输入、输出功可改写为
\[
2P(-\delta a)=W\delta H,
\qquad \frac{\delta H}{-2\delta a}=\frac ah.
\]
丝杠缩短的距离是 $-2\delta a$，顶部上升距离是 $\delta H$。若机构接近平摊，$a/h$ 很大，丝杠小幅缩短会使顶部上升较多，相应地需要很大的端力；这与杆压缩力 $S\propto1/h$ 同时增大的结果一致。

\end{solution}

\par\medskip\noindent\begin{minipage}{\textwidth}
\begin{example}[先滑动还是先倾覆]\label{mp5:example:l5-e6}
均匀长方体质量为 $M$，剖面内的底宽为 $2a$，高为 $h$。在左侧面离地高度为 $H$ 的点 $P$ 施加水平向右的力 $F$，$0<H\le h$，见图\ref{mp5:fig:l5-e6-q}。地面静摩擦系数为 $\mu$，从零逐渐增大 $F$。求静止条件，以及先滑动或先倾覆的判据。
\end{example}

\begin{figure}[H]\centering
\begin{tikzpicture}
\draw[curve](-1.7,0)--(1.8,0);\draw[curve](-1,0)rectangle(1,3);\fill(0,1.5)circle(1.3pt)node[left]{$G$};
\fill(-1,2.4)circle(1.3pt)node[above left]{$P$};
\draw[axis](-1,2.4)--(.35,2.4)node[above]{$F$};\draw[guide](-1,0)--(-1,-.75);\draw[guide](1,0)--(1,-.75);\draw[<->,>=Stealth](-1,-.6)--(1,-.6)node[midway,below]{$2a$};\draw[guide](-1,2.4)--(-1.55,2.4);\draw[<->,>=Stealth](-1.4,0)--(-1.4,2.4)node[midway,left]{$H$};\draw[guide](1,3)--(1.75,3);\draw[<->,>=Stealth](1.6,0)--(1.6,3)node[midway,right]{$h$};
\end{tikzpicture}
\caption{长方体受水平推力，重心 $G$ 位于几何中心，点 $P$ 的高度为 $H$。}\label{mp5:fig:l5-e6-q}
\end{figure}
\end{minipage}\par\medskip
\begin{solution}
设地面支持力合力为 $N$，作用线相对底面中心的水平坐标为 $x_N$；摩擦力大小为 $f$，方向向左。图\ref{mp5:fig:l5-e6-s}左图的平衡方程给出
\[
N=Mg,\qquad f=F,\qquad Nx_N=FH.
\]
地面只能推、不能拉，因此支持力合力的作用线必须位于底面内：$|x_N|\le a$。与摩擦条件一起，
\[
F\le\mu Mg,\qquad F\le\frac{Mga}{H}.
\]

\begin{figure}[H]\centering
\begin{tikzpicture}
\begin{scope}[shift={(0,0)},scale=0.9]
\draw[curve](-1.7,0)--(1.8,0);\draw[curve](-1,0)rectangle(1,3);\fill(0,1.5)circle(1.3pt)node[left]{$G$};
\draw[axis](-1,2.4)--(.35,2.4)node[above]{$F$};\draw[axis](0,1.5)--(0,.7)node[left]{$Mg$};\fill(.5,0)circle(1.2pt);\draw[axis](.5,0)--(.5,.8)node[right]{$N$};\draw[axis](.5,0)--(-.35,0)node[midway,below]{$f$};\draw[guide](0,0)--(0,-.6);\draw[guide](.5,0)--(.5,-.6);\draw[<->,>=Stealth](0,-.5)--(.5,-.5)node[midway,below]{$x_N$};\end{scope}
\begin{scope}[shift={(5.3,0)},scale=0.9]
\draw[curve](-1.7,0)--(1.8,0);\draw[curve](-1,0)rectangle(1,3);\fill(0,1.5)circle(1.3pt)node[left]{$G$};
\fill(1,0)circle(1.3pt)node[below right]{$C$};
\fill(-1,2.4)circle(1.3pt)node[left]{$P$};
\draw[axis](0,1.5)--(.6,1.9)node[above left]{$\delta\vct r_G$};\draw[guide](0,1.5)--(.6,1.5)--(.6,1.9);\draw[guide](.6,1.5)--(1.5,1.5);\draw[guide](.6,1.9)--(1.5,1.9);\draw[<->,>=Stealth](1.4,1.5)--(1.4,1.9)node[midway,right]{$\delta y_G$};\draw[axis](-1,2.4)--(-.04,2.4)node[midway,above]{$\delta x_P$};\draw[axis](.45,0)arc(180:150:.55);\node at(.35,.38){$\delta\varphi$};\end{scope}

\end{tikzpicture}
\caption{左图将底面接触作用合成为 $N,f$；右图画出绕右边缘 $C$ 的虚位移及重心位移的竖直分量。}\label{mp5:fig:l5-e6-s}
\end{figure}

为明确支持力的非负条件，可把底面作用等效分配到左、右边缘，记向上的等效力为 $N_L,N_R$。整体的竖直力和力矩要求
\[
N_L+N_R=Mg,\qquad a(N_R-N_L)=FH,
\]
从而
\[
N_L=\frac{Mg}{2}-\frac{FH}{2a},\qquad
N_R=\frac{Mg}{2}+\frac{FH}{2a}.
\]
$N_L,N_R\ge0$ 等价于支持合力落在底面内。图\ref{mp5:fig:l5-e6-edge}取 $FH/(Mga)=3/5$ 作示意：右侧承载增大，左侧仍未卸载。临界 $FH=Mga$ 时左侧卸载，右侧总支持力仍为 $Mg$。
\begin{figure}[H]\centering
\begin{tikzpicture}[scale=.95]
\draw[thick](-.5,0)--(4.5,0);\draw[curve](.5,0)rectangle(3.5,2.7);
\fill(2,1.35)circle(1.3pt)node[right]{$G$};\draw[axis](2,1.35)--(2,.6)node[right]{$Mg$};
\draw[axis](.5,0)--(.5,.4)node[left]{$N_L$};\draw[axis](3.5,0)--(3.5,1.6)node[right]{$N_R$};
\fill(.5,2)circle(1.2pt)node[above left]{$P$};\draw[axis](.5,2)--(1.6,2)node[above]{$F$};\draw[axis](2,0)--(.9,0)node[midway,below]{$f$};
\draw[|-|](.5,-.5)--(2,-.5)node[midway,below]{$a$};\draw[|-|](2,-.5)--(3.5,-.5)node[midway,below]{$a$};
\node[anchor=west]at(5,2){$N_L+N_R=Mg$};\node[anchor=west]at(5,1.2){$a(N_R-N_L)=FH$};
\end{tikzpicture}
\caption{将支持合力等效分配到两侧边缘。图取 $N_L:N_R=1:4$，两支持力箭头按此比例绘制；其余箭头只表示方向。}\label{mp5:fig:l5-e6-edge}
\end{figure}

当支持力作用线到达右边缘 C，可考虑绕 C 顺时针的允许虚转角 $\delta\varphi>0$。右图给出一阶几何关系
\[
\delta x_P=H\delta\varphi,\qquad
\delta y_G=a\delta\varphi.
\]
底面其余点向上离地，属于单侧接触允许的方向。若 C 不滑动，该处支持力与摩擦力均不做虚功，于是
\[
\delta W=(FH-Mga)\delta\varphi.
\]
$FH>Mga$ 时外力沿倾覆方向的虚功为正，静止不能维持。两种阈值为
\[
F_{\rm slide}=\mu Mg,\qquad F_{\rm tip}=\frac{Mga}{H}.
\]
所以 $\mu<a/H$ 时先滑动，$\mu>a/H$ 时先倾覆；等号时二者同时达到临界。这个判据还要求接触反力能够由底面上的非负压力分布实现，不需要预先假定压力均匀。
静摩擦条件应在倾覆候选点处再核对一次。若先达到 $F_{\rm tip}=Mga/H$，该点所需的 $f/N=a/H$；只有 $\mu\ge a/H$ 才能使 C 在倾覆瞬间保持不滑。反过来，在滑动候选点 $F_{\rm slide}=\mu Mg$，支持合力偏移为 $x_N=\mu H$；只有 $\mu H\le a$ 时底面尚能提供这一合力。两种检验给出同一先后判据。

\end{solution}


\par\medskip\noindent\begin{minipage}{\textwidth}
\begin{example}[双滑块圆盘机构的平衡力]\label{mp5:hw:l5-37}
半径$r$的圆盘绕固定中心$O$转动，盘上$C,D$两铰点满足$OC\perp OD$。$D$通过长$r$的杆连接到水平滑块$E$，$C$通过长$r$的杆连接到竖直滑块$B$，后者携带重力为$G$的物体$A$。水平导轨为过$O$的$x$轴，竖直导轨为$x=r$。取$\alpha=\angle DOE\in(0,\pi/2)$，各点位于图\ref{mp5:fig:hw:l5-37-q}所示分支，求作用在$E$上、以向左为正的平衡力$F$。忽略杆和盘的重力及所有摩擦。
\end{example}

\begin{figure}[H]\centering
\begin{tikzpicture}[>=Stealth,line cap=round,line join=round]

\draw[curve](0,0)circle(1.4);\coordinate(D)at(-.99,.99);\coordinate(C)at(.99,.99);\coordinate(E)at(-1.98,0);\coordinate(B)at(1.4,2.329);\draw[thick](E)--(D);\draw[thick](C)--(B)--(1.4,3);\draw[guide](-2.8,0)--(2.2,0);\draw[guide](1.4,-.5)--(1.4,3.1);\draw[fill=gray!15](1.1,2.9)rectangle(1.7,3.4);\node at(1.4,3.15){$A$};\draw[axis](1.9,3.1)--(1.9,2.3)node[right]{$G$};\draw[axis](E)--(-2.9,0)node[above]{$F$};\draw[guide](D)--(0,0)--(C);\foreach \p in {B,C,D,E}{\fill(\p)circle(1.4pt)node[above left]{$\p$};}\fill(0,0)circle(1.4pt)node[below]{$O$};\draw(-.6,0)arc(180:135:.6);\node at(-.78,.23){$\alpha$};\node at(-2.3,1.1){$DE=r$};\node[right]at(1.7,1.6){$BC=r$};

\end{tikzpicture}
\caption{圆盘的两根连杆分别驱动水平与竖直滑块}\label{mp5:fig:hw:l5-37-q}
\end{figure}

\end{minipage}\par\medskip
\begin{solution}
取唯一的形状坐标$\alpha$，各点位置均可由它表示。当前分支的盘上点坐标为
\[
 D=(-r\cos\alpha,r\sin\alpha),\qquad
 C=(r\sin\alpha,r\cos\alpha).
\]
由$DE=r$且$E$在$x$轴左侧，
\[
 x_E=-2r\cos\alpha,\qquad\dd x_E=2r\sin\alpha\,\dd\alpha.
\]
设$BC$与竖直线夹角为$\beta$，则水平投影给出
\[
 \sin\beta=1-\sin\alpha,\qquad
 \cos\beta\,\dd\beta=-\cos\alpha\,\dd\alpha.
\]
竖直滑块高度为$y_B=r\cos\alpha+r\cos\beta$，因此
\[
 \dd y_B=r\left[-\sin\alpha
            +(1-\sin\alpha)\frac{\cos\alpha}{\cos\beta}\right]\dd\alpha.
\]
图\ref{mp5:fig:hw:l5-37-s}中的水平力做功$-F\dd x_E$，重力做功$-G\dd y_B$，约束反力均不做虚功。故
\[
 \boxed{\frac FG=\frac{\sin\alpha-
 (1-\sin\alpha)\cos\alpha/\sqrt{1-(1-\sin\alpha)^2}}
 {2\sin\alpha}}.
\]
当$\alpha=\pi/6$时，$\beta=\pi/6$且$\dd y_B/\dd\alpha=0$，所需水平力为零。$\alpha>\pi/6$时$F>0$，方向向左；$\alpha<\pi/6$时结果为负，实际应向右作用。端点$\alpha\to0$接近连杆的奇异构形，不能把分母为零的式子直接用于端点。


\begin{figure}[H]\centering
\begin{tikzpicture}[>=Stealth,line cap=round,line join=round]

\draw[thick](-1.3,0)--(1.3,0);\fill(0,0)circle(1.5pt)node[below]{$E$};\draw[axis,thick](0,.2)--(.9,.2)node[right]{$\delta x_E$};\draw[axis](0,-.3)--(-.9,-.3)node[left]{$F$};\begin{scope}[shift={(4.7,.8)}]\draw[thick](0,-1.5)--(0,1);\fill(0,0)circle(1.5pt)node[right]{$B$};\draw[axis,thick](.4,0)--(.4,-.9)node[right]{$-\delta y_B$};\draw[axis](-.4,0)--(-.4,-.9)node[left]{$G$};\end{scope}\node at(1.8,-1.3){图示$\alpha>\pi/6$、$\delta\alpha>0$时的虚位移};

\end{tikzpicture}
\caption{两个滑块的虚位移由同一几何参数关联；外力功与重力功相抵}\label{mp5:fig:hw:l5-37-s}
\end{figure}

\end{solution}

\par\medskip\noindent\begin{minipage}{\textwidth}
\begin{example}[两根铰接杆的地面反力]\label{mp5:hw:l5-40}
两根长$l$、质量$m$的均匀杆首尾铰接，上杆上端固定于$O$，下杆下端$P$与粗糙水平地面接触。两杆分别与竖直向下方向成角$\theta_1,\theta_2\in(0,\pi/2)$，均向右倾斜，见图\ref{mp5:fig:hw:l5-40-q}。求地面法向力$N$与以向右为正的摩擦力$f$，并给出接触、静摩擦条件。
\end{example}

\begin{figure}[H]\centering
\begin{tikzpicture}[>=Stealth,line cap=round,line join=round]

\coordinate(O)at(0,3.346);\coordinate(A)at(.518,1.414);\coordinate(P)at(1.932,0);\draw[thick](-.7,0)--(3,0);\draw[thick](O)--(A)--(P);\draw[guide](O)--(0,1.2);\draw[guide](A)--(.518,-.3);\fill(O)circle(1.5pt)node[above]{$O$};\fill(A)circle(1.5pt)node[right]{$A$};\fill(P)circle(1.5pt)node[below]{$P$};\draw[axis](P)--(1.932,1)node[right]{$N$};\draw[axis](P)--(2.8,0)node[above]{$f$};\node[right]at(.1,2.6){$\theta_1$};\node[right]at(.72,.8){$\theta_2$};\node[left]at(.2,2){$l,m$};\node[above right]at(1.3,1.15){$l,m$};

\end{tikzpicture}
\caption{给定两杆姿态，先把地面反力作为未知外力求出，再检验其可实现性}\label{mp5:fig:hw:l5-40-q}
\end{figure}

\end{minipage}\par\medskip
\begin{solution}
取$O$为原点，竖直向上为$y$正向。若暂时释放端点的地面约束，端点坐标为
\[
 x_P=l(\sin\theta_1+\sin\theta_2),\qquad
 y_P=-l(\cos\theta_1+\cos\theta_2).
\]
两杆重力势能为
\[
 U_g=-\frac{mgl}{2}(3\cos\theta_1+\cos\theta_2).
\]
将$N,f$视为维持此构形的待定外力，图\ref{mp5:fig:hw:l5-40-s}中两种独立虚位移给出$\delta U_g=f\delta x_P+N\delta y_P$，即
\[
 (3mg/2-N)\sin\theta_1-f\cos\theta_1=0,
\]
\[
 (mg/2-N)\sin\theta_2-f\cos\theta_2=0.
\]
将两式分别除以$\sin\theta_1,\sin\theta_2$后相减，消去$N$，得到
\[
mg=f(\cot\theta_1-\cot\theta_2)
=f\frac{\sin(\theta_2-\theta_1)}{\sin\theta_1\sin\theta_2}.
\]
若$\sin(\theta_2-\theta_1)\ne0$，先求$f$再代回任一平衡方程，得
\[
 \boxed{f=\frac{mg\sin\theta_1\sin\theta_2}
                   {\sin(\theta_2-\theta_1)}},
\]
\[
 \boxed{N=\frac{mg[\cos\theta_1\sin\theta_2
                    -3\sin\theta_1\cos\theta_2]}
                   {2\sin(\theta_2-\theta_1)}}.
\]
$f$的符号直接确定维持该姿态所需的摩擦方向。设地面的静摩擦系数为$\mu$，实际接触须同时满足$N\ge0$和$|f|\le\mu N$；特别地，$N=0$时也必须有$f=0$。在本题角度范围内，非共线解的$f\ne0$，所以可实现的平衡必有$N>0$。不满足这些条件时，题示姿态无法靠该地面静止维持。若$\theta_1=\theta_2$，原两方程相减会给出$mg\sin\theta_1=0$，与$0<\theta_1<\pi/2$矛盾；同向共线姿态不存在有限反力的平衡解。

暂时释放约束后，约束力便以维持构形所需的广义外力进入方程。若从一开始只取$\delta y_P=0$的实际允许虚位移，地面法向力的功为零，它会从方程中消去，便无法单独求出$N$。


\begin{figure}[H]\centering
\begin{tikzpicture}[>=Stealth,line cap=round,line join=round]

\draw[thick](0,2.8)--(.45,1.2)--(1.65,0);\fill(1.65,0)circle(1.5pt)node[below]{$P$};\draw[axis,thick](1.65,0)--(2.65,0)node[right]{$\delta x_P$};\draw[axis,thick](1.65,0)--(1.65,1.1)node[above]{$\delta y_P$};\draw[axis](.225,2)--(.225,1.3)node[left]{$mg$};\draw[axis](1.05,.6)--(1.05,-.1)node[left]{$mg$};\node[right,align=left]at(3.5,1.6){暂时释放端点约束\\[5pt]$f\delta x_P+N\delta y_P=\delta U_g$\\[5pt]求出反力后恢复$y_P=\text{常量}$};

\end{tikzpicture}
\caption{为求反力而取端点的两种独立虚位移，两杆重力作用于各自中点}\label{mp5:fig:hw:l5-40-s}
\end{figure}

\end{solution}

\section{约束系统的能量解法}
势能沿允许构形的变化决定平衡和回复趋势，动能给出运动的惯性尺度。两者使用同一组独立坐标，分别处理稳定性与有限振幅运动。

\subsection{位置相关的等效质量}
对于不随时间变化的理想约束，约束力的总功率为零；若其余力有势能，机械能 $T+U$ 守恒。若约束本身运动，理想约束只保证虚功为零，实际功率可以不为零，例\ref{mp5:example:l5-e2}中的转动直杆就是这种情况。

单自由度系统的动能常为 $T=M(q)\dot q^2/2$。其中 $M(q)>0$ 是这个坐标下的等效质量，未必等于某个物体的质量。由拉格朗日方程直接得到
\begin{equation}\label{mp5:eq:l5-variable-mass}
M(q)\ddot q+\frac12M'(q)\dot q^2+U'(q)=Q^{\rm nc}.
\end{equation}
推导时先算 $\dd(M\dot q)/\dd t=M\ddot q+M'\dot q^2$，再减去 $\partial T/\partial q=M'\dot q^2/2$。$M(q)$ 的导数项须保留；这一推导没有除以 $\dot q$，故式\eqref{mp5:eq:l5-variable-mass}在转向点同样有效。

无非保守力时，
\[
E=\frac12M(q)\dot q^2+U(q),\qquad
\dot q=\pm\sqrt{\frac{2[E-U(q)]}{M(q)}}.
\]
若普通转向点 $q_-,q_+$ 之间是同一个可达区间，去回两支的时间均为
\[
\int_{q_-}^{q_+}\sqrt{\frac{M(q)}{2[E-U(q)]}}\dd q,
\]
相加后一个周期为
\begin{equation}\label{mp5:eq:l5-period}
T_{\rm osc}=\sqrt2\int_{q_-}^{q_+}\sqrt{\frac{M(q)}{E-U(q)}}\,\dd q.
\end{equation}
在 $U'(q_e)=0$、$U''(q_e)>0$ 的普通平衡点附近，令 $q=q_e+\xi$，保留一阶小量，得
\begin{equation}\label{mp5:eq:l5-small}
M(q_e)\ddot\xi+U''(q_e)\xi=0,
\qquad \omega^2=\frac{U''(q_e)}{M(q_e)}.
\end{equation}
$M'\dot q^2$ 是二阶小量；$M(q)-M(q_e)$ 为一阶，与 $\ddot\xi$ 相乘也是二阶，故二者均不进入线性方程。若 $U''=0$，则须保留更高阶项，不能用零频率替代完整的稳定性分析。

也可在 $M(q)>0$ 的区间内重新标记位置。取常量 $M_0>0$，定义
\[
s(q)=\int_{q_e}^q\sqrt{\frac{M(u)}{M_0}}\dd u,
\qquad
\dot s=\frac{\dd s}{\dd q}\dot q
=\sqrt{\frac{M(q)}{M_0}}\dot q.
\]
于是 $T=M(q)\dot q^2/2=M_0\dot s^2/2$。势能同时改写为 $U(q(s))$，运动时间仍由同一条实际轨迹决定；这种换元只把位置相关的惯性系数移入了坐标映射。


\begin{figure}[H]\centering
\begin{tikzpicture}
\draw[axis](-2.1,0)--(2.1,0)node[right]{$q$};
\draw[axis](0,-0.2)--(0,2.1)node[above]{$U$};
\draw[curve] plot coordinates{(-1.90000,1.80500) (-1.86200,1.73352) (-1.82400,1.66349) (-1.78600,1.59490) (-1.74800,1.52775) (-1.71000,1.46205) (-1.67200,1.39779) (-1.63400,1.33498) (-1.59600,1.27361) (-1.55800,1.21368) (-1.52000,1.15520) (-1.48200,1.09816) (-1.44400,1.04257) (-1.40600,0.98842) (-1.36800,0.93571) (-1.33000,0.88445) (-1.29200,0.83463) (-1.25400,0.78626) (-1.21600,0.73933) (-1.17800,0.69384) (-1.14000,0.64980) (-1.10200,0.60720) (-1.06400,0.56605) (-1.02600,0.52634) (-0.98800,0.48807) (-0.95000,0.45125) (-0.91200,0.41587) (-0.87400,0.38194) (-0.83600,0.34945) (-0.79800,0.31840) (-0.76000,0.28880) (-0.72200,0.26064) (-0.68400,0.23393) (-0.64600,0.20866) (-0.60800,0.18483) (-0.57000,0.16245) (-0.53200,0.14151) (-0.49400,0.12202) (-0.45600,0.10397) (-0.41800,0.08736) (-0.38000,0.07220) (-0.34200,0.05848) (-0.30400,0.04621) (-0.26600,0.03538) (-0.22800,0.02599) (-0.19000,0.01805) (-0.15200,0.01155) (-0.11400,0.00650) (-0.07600,0.00289) (-0.03800,0.00072) (0.00000,0.00000) (0.03800,0.00072) (0.07600,0.00289) (0.11400,0.00650) (0.15200,0.01155) (0.19000,0.01805) (0.22800,0.02599) (0.26600,0.03538) (0.30400,0.04621) (0.34200,0.05848) (0.38000,0.07220) (0.41800,0.08736) (0.45600,0.10397) (0.49400,0.12202) (0.53200,0.14151) (0.57000,0.16245) (0.60800,0.18483) (0.64600,0.20866) (0.68400,0.23393) (0.72200,0.26064) (0.76000,0.28880) (0.79800,0.31840) (0.83600,0.34945) (0.87400,0.38194) (0.91200,0.41587) (0.95000,0.45125) (0.98800,0.48807) (1.02600,0.52634) (1.06400,0.56605) (1.10200,0.60720) (1.14000,0.64980) (1.17800,0.69384) (1.21600,0.73933) (1.25400,0.78626) (1.29200,0.83463) (1.33000,0.88445) (1.36800,0.93571) (1.40600,0.98842) (1.44400,1.04257) (1.48200,1.09816) (1.52000,1.15520) (1.55800,1.21368) (1.59600,1.27361) (1.63400,1.33498) (1.67200,1.39779) (1.71000,1.46205) (1.74800,1.52775) (1.78600,1.59490) (1.82400,1.66349) (1.86200,1.73352) (1.90000,1.80500)};
\draw[guide](-1.5,0)--(-1.5,1.125)--(1.5,1.125)--(1.5,0);\node[right]at(2,1.125){$E$};\fill(-1.5,1.125)circle(1.2pt)node[above left]{$q_-$};\fill(1.5,1.125)circle(1.2pt)node[below right]{$q_+$};\node[below left]at(0,0){$q_e$};\begin{scope}[shift={(6,0)},scale=1]
\node[align=center]at(0,1.5){同一 $q$、同一 $E-U$};\draw[axis](-1,.8)--(.3,.8)node[right]{$\dot q_1$};\draw[axis](-1,.1)--(-.35,.1)node[right]{$\dot q_2$};\node[align=center]at(.1,-.7){$M_2=4M_1$\quad $\dot q_2=\dot q_1/2$};\end{scope}

\end{tikzpicture}
\caption{同一势阱中，不同的惯性系数会改变运动时间。引入 $s=\int\sqrt{M(q)/M_0}\,\dd q$ 后，动能化为 $M_0\dot s^2/2$；势能须同时写成 $U(q(s))$。}\label{mp5:fig:l5-metric}
\end{figure}

\subsection{滚动约束与循环坐标}
刚体在固定平面上纯滚动，接触点瞬时速度为零，静摩擦力因而不做实际功；它仍可能不为零，并决定刚体的角加速度。写出能量解后，须用平动、转动方程求摩擦力，检验 $|f|\le\mu N$。若承载面也运动，接触点在地面系中未必静止，但两接触体上的摩擦功率仍可在整个系统内相互抵消。

若系统有循环坐标 $\phi$，相应动量 $J$ 守恒。以中心力运动为例，
\[
L=\frac m2\dot r^2+\frac m2r^2\dot\phi^2-U(r),\qquad
J=mr^2\dot\phi.
\]
将守恒关系代入能量，得到径向问题
\begin{equation}\label{mp5:eq:l5-central-reduction}
E=\frac m2\dot r^2+U_{\rm eff}(r),\qquad
U_{\rm eff}=U(r)+\frac{J^2}{2mr^2}.
\end{equation}
这里固定的是 $J$。若要从拉格朗日函数约化，须先作勒让德变换 $\mathcal R=L-J\dot\phi$，再代入 $\dot\phi=J/(mr^2)$，得
\[
\mathcal R=\frac m2\dot r^2+\frac{J^2}{2mr^2}-U(r)-\frac{J^2}{mr^2}
=\frac m2\dot r^2-U(r)-\frac{J^2}{2mr^2}.
\]
直接把 $\dot\phi$ 代回 $L$ 再对 $r$ 变分，会把离心项的符号算错。固定角速度与固定角动量也是不同条件：前者通常需要外界提供力矩和功，后者允许角速度随构形变化。

\par\medskip\noindent\begin{minipage}{\textwidth}
\begin{example}[纯滚动的能量与摩擦]\label{mp5:example:l5-e7}
质量 $m$、半径 $R$ 的圆柱绕自身轴的转动惯量为 $I=\eta mR^2$。它从静止开始沿倾角 $\alpha\in(0,\pi/2)$ 的固定粗糙斜面纯滚动，见图\ref{mp5:fig:l5-e7-q}。求沿斜面下降距离 $s$ 后的速度、加速度及维持纯滚动所需的最小静摩擦因数。
\end{example}

\begin{figure}[H]\centering
\begin{tikzpicture}
\draw[curve](-2,1.1547)--(3,-1.7321);\draw[guide](0,0)--(2,0);\draw(1,0)arc(0:-30:1);\node at(1.3,-.3){$\alpha$};\draw[curve](0.49999999999999994,0.8660254037844387)circle(1.0);\draw(0.49999999999999994,0.8660254037844387)--(0,0);\fill(0.49999999999999994,0.8660254037844387)circle(1.3pt)node[above right]{$C$};
\fill(0,0)circle(1.3pt)node[below left]{$Q$};
\node at(.4,2.35){$m,\ I=\eta mR^2$};\node[right]at(.27,.45){$R$};\draw[axis](1.5,-1.45)--(2.7,-2.1428)node[right]{$s$};
\end{tikzpicture}
\caption{$s$ 沿斜面向下为正，$C$ 是圆柱质心，$Q$ 是瞬时接触点。}\label{mp5:fig:l5-e7-q}
\end{figure}
\end{minipage}\par\medskip
\begin{solution}
图\ref{mp5:fig:l5-e7-q}中质心下降高度为 $s\sin\alpha$，纯滚动给出角速度大小 $|\omega|=\dot s/R$。于是
\[
T=\frac12m\dot s^2+\frac12I\frac{\dot s^2}{R^2}
=\frac12m(1+\eta)\dot s^2.
\]
取初态势能为零，$U=-mgs\sin\alpha$。接触点 Q 相对固定斜面的速度为零，静摩擦力功率为零，故
\[
\frac12m(1+\eta)\dot s^2=mgs\sin\alpha,
\qquad \boxed{v^2=\frac{2gs\sin\alpha}{1+\eta}}.
\]

\begin{figure}[H]\centering
\begin{tikzpicture}
\draw[curve](-2,1.1547)--(3,-1.7321);\draw[guide](0,0)--(2,0);\draw(1,0)arc(0:-30:1);\node at(1.3,-.3){$\alpha$};\draw[curve](0.49999999999999994,0.8660254037844387)circle(1.0);\draw(0.49999999999999994,0.8660254037844387)--(0,0);\fill(0.49999999999999994,0.8660254037844387)circle(1.3pt)node[above left]{$C$};
\fill(0,0)circle(1.3pt)node[below left]{$Q$};
\draw[axis](0.49999999999999994,0.8660254037844387)--(0.49999999999999994,-0.5839745962155612)node[right]{$mg$};\draw[axis](0,0)--(1.1,1.90526)node[above]{$N$};\draw[axis](0,0)--(-1.0392,.6)node[left]{$f$};\draw[axis](.5,.86603)--(1.2794,.41603)node[right]{$a$};
\end{tikzpicture}
\caption{圆柱受力与质心加速度。$N$ 沿斜面法线，$f$ 沿坡向上；箭头只表示方向。}\label{mp5:fig:l5-e7-s}
\end{figure}

图\ref{mp5:fig:l5-e7-s}把摩擦力画在 Q 处，方向沿坡向上。沿坡方向与绕质心分别有
\[
mg\sin\alpha-f=ma,\qquad fR=I\frac aR.
\]
由转动方程得$f=\eta ma$，代回平动方程，$mg\sin\alpha=m(1+\eta)a$，所以
\[
\boxed{a=\frac{g\sin\alpha}{1+\eta}},\qquad
f=\frac{\eta mg\sin\alpha}{1+\eta},\qquad N=mg\cos\alpha.
\]
这一联立过程没有约去速度，因此也适用于$v=0$的释放瞬间。不发生相对滑动须满足
\[
\boxed{\mu\ge\frac{\eta}{1+\eta}\tan\alpha.}
\]
均匀实心圆柱 $\eta=1/2$，有 $a=2g\sin\alpha/3$、$\mu_{\min}=\tan\alpha/3$。摩擦力减小质心加速度，同时提供自转所需力矩；计算运动后的接触条件时仍须保留它。

静摩擦的功率还可拆成两部分核对。质心平动动能变化率为 $(mg\sin\alpha-f)v$，绕质心的转动动能变化率为 $fR|\omega|=fv$，相加后摩擦项抵消。因此重力功率全部进入总动能；若只研究质心平动，摩擦所做的负功并不能当成整个刚体的热损失。图\ref{mp5:fig:l5-e7-s}中的摩擦方向由产生下行自转所需的力矩确定，与静摩擦不耗散能量并不矛盾。

条件 $f/N\le\mu$ 是对所得纯滚动解的可实现性检验。本题 f 与 N 都是常量，只须核对一个比值；若路径倾角或外力随位置变化，则应沿全过程求所需比值的最大值。

\end{solution}

\par\medskip\noindent\begin{minipage}{\textwidth}
\begin{example}[弹簧驱动的滚动振子]\label{mp5:example:l5-e8}
圆盘质量为 $m$、半径为 $R$，$I=\eta mR^2$。劲度系数为 $k$ 的水平弹簧连接圆盘轴心，圆盘在固定水平面上纯滚动，$x$ 从弹簧自然长度对应的质心位置起算，见图\ref{mp5:fig:l5-e8-q}。从 $x=A>0$ 静止释放，求振动周期、摩擦方向及纯滚动条件。
\end{example}

\begin{figure}[H]\centering
\begin{tikzpicture}
\draw[curve](-3.4,0)--(3.1,0);\draw[curve](-3,0)--(-3,2.4);\draw[curve](.7,1.2)circle(1.2);\fill(0.7,1.2)circle(1.3pt)node[above right]{$C$};
\fill(0.7,0)circle(1.3pt)node[below right]{$Q$};
\draw(-3,1.2)--(-2.7,1.2)--(-2.55,1.4)--(-2.25,1)--(-1.95,1.4)--(-1.65,1)--(-1.35,1.4)--(-1.05,1)--(-.9,1.2)--(.7,1.2);\node[above]at(-1.8,1.5){$k$};\draw[guide](-.4,-.1)--(-.4,-.65);\draw[guide](.7,-.1)--(.7,-.65);\draw[axis](-.4,-.55)--(2,-.55)node[right]{$x$};\node[below]at(-.4,-.55){$0$};\node at(1,2.8){$m,\ I$};
\end{tikzpicture}
\caption{弹簧连接轴心，$x$ 从弹簧自然长度对应的位置起算，$Q$ 为地面接触点。}\label{mp5:fig:l5-e8-q}
\end{figure}
\end{minipage}\par\medskip
\begin{solution}
令逆时针角速度为正，则纯滚动约束是 $\omega_z=-\dot x/R$。平动和转动动能相加，得到
\[
T=\frac12m(1+\eta)\dot x^2,\qquad U=\frac12kx^2.
\]
对图\ref{mp5:fig:l5-e8-s}中的力取水平向右为正，平动与转动方程分别为
\[
m\ddot x=-kx+f,\qquad I\dot\omega_z=Rf,
\qquad \dot\omega_z=-\frac{\ddot x}{R}.
\]
由后两式得 $f=-\eta m\ddot x$，代回第一式，才得到含等效质量 $m(1+\eta)$ 的方程。弹簧接在轴心，对轴心不产生力矩；维持相应自转加速度的力矩全部由地面摩擦提供。

故 $m(1+\eta)\ddot x+kx=0$。令 $\omega^2=k/[m(1+\eta)]$，特征根为 $\pm\mathrm i\omega$，通解为 $x=C\cos\omega t+D\sin\omega t$。初位移给出 $C=A$，初速度给出 $\omega D=0$，所以
\[
x=A\cos\omega t,\quad
\boxed{\omega^2=\frac{k}{m(1+\eta)},\qquad
T_{\rm osc}=2\pi\sqrt{\frac{m(1+\eta)}k}.}
\]

\begin{figure}[H]\centering
\begin{tikzpicture}
\draw[curve](-3.4,0)--(3.1,0);\draw[curve](-3,0)--(-3,2.4);\draw[curve](.7,1.2)circle(1.2);\fill(0.7,1.2)circle(1.3pt)node[above right]{$C$};
\fill(0.7,0)circle(1.3pt)node[below right]{$Q$};
\draw(-3,1.2)--(-2.7,1.2)--(-2.55,1.4)--(-2.25,1)--(-1.95,1.4)--(-1.65,1)--(-1.35,1.4)--(-1.05,1)--(-.9,1.2)--(.7,1.2);\node[above]at(-1.8,1.5){$k$};\draw[axis](.7,1.2)--(-.5,1.2)node[above]{$-kx$};\draw[axis](.7,0)--(1.75,0)node[below]{$f>0$};\draw[axis](.7,1.2)--(.7,.75)node[right]{$mg$};\draw[axis](.7,0)--(.7,.45)node[left]{$N$};\draw[axis](1.55,2.05)arc(45:110:1.2);\node at(.1,2.7){$\dot\omega_z>0$};
\end{tikzpicture}
\caption{$x>0$ 时，弹簧向左拉、摩擦向右，圆盘的角加速度逆时针。$N$ 与 $mg$ 的箭头等长。}\label{mp5:fig:l5-e8-s}
\end{figure}

将简谐运动的加速度$\ddot x=-kx/[m(1+\eta)]$代入已经求出的$f=-\eta m\ddot x$，得到
\[
f=\frac{\eta}{1+\eta}kx.
\]
图\ref{mp5:fig:l5-e8-s}中$x>0$，摩擦向右，质心加速度向左；$x<0$时二者反向。摩擦抵抗接触点的相对滑动趋势，不能根据质心速度直接判定方向。

竖直方向 $N=mg$，整个周期中 $|f|$ 的最大值为 $\eta kA/(1+\eta)$，所以
\[
\boxed{\mu\ge\frac{\eta kA}{(1+\eta)mg}.}
\]
均匀圆盘对应 $\omega^2=2k/(3m)$、$\mu\ge kA/(3mg)$。若不满足这一条件，上述简谐解只能用于实际仍保持纯滚动的阶段，随后须改用滑动摩擦模型。
这一方向判断在来回运动中也成立。同一个 $x>0$ 位置，圆盘去、回两次的速度相反，所需摩擦却都向右，因为两个时刻的加速度相同。通过 $x=0$ 时速度最大，而弹簧力和静摩擦均为零；在 $x=\pm A$ 的转向点，速度为零，所需摩擦的绝对值反而最大。全周期所需的静摩擦上限因而在转向点取得。

静摩擦足够时，能量可写成
\[
\frac12m(1+\eta)\dot x^2+\frac12kx^2=\frac12kA^2,
\qquad v_{\max}=A\sqrt{\frac{k}{m(1+\eta)}}.
\]
最大静摩擦约束也可表示为允许振幅
\[
A\le A_{\max}=\frac{(1+\eta)\mu mg}{\eta k}.
\]
这一上限限定了简谐滚动解的振幅范围。超过上限后会出现相对滑动，须重新求各阶段运动及其周期。

\end{solution}

\par\medskip\noindent\begin{minipage}{\textwidth}
\begin{example}[可动斜面上的滚动圆柱]\label{mp5:example:l5-e9}
质量为 $M$ 的斜面置于光滑水平地面，倾角为 $\alpha$。质量为 $m$、半径为 $R$、转动惯量 $I=\eta mR^2$ 的圆柱在斜面上纯滚动，全部物体从静止开始运动。用 $X$ 表示斜面的水平位移，用 $s$ 表示圆柱相对斜面沿坡向下的位移，见图\ref{mp5:fig:l5-e9-q}。求两坐标的加速度及纯滚动条件。
\end{example}

\begin{figure}[H]\centering
\begin{tikzpicture}
\draw[curve](-.5,0)--(4.2,0);\draw[curve](0,0)--(0,2.07846)--(3.6,0)--cycle;\node at(1,.45){$M$};\draw[curve](1.525,1.9485571585149868)circle(0.65);\fill(1.525,1.9485571585149868)circle(1.3pt)node[above right]{$C$};
\fill(1.2,1.3856406460551016)circle(1.3pt)node[below left]{$Q$};
\node at(1.6,3.5){$m,\ I=\eta mR^2$};\draw[axis](2.4,3.1)--(3.5,2.46491)node[right]{$s$};\draw[axis](0,-.5)--(2.5,-.5)node[right]{$X$};\draw(2.8,0)arc(180:150:.8);\node at(2.6,.3){$\alpha$};
\end{tikzpicture}
\caption{斜面水平平移，$X$ 为其地面系位移，$s$ 为圆柱相对斜面沿坡向下的位移。}\label{mp5:fig:l5-e9-q}
\end{figure}
\end{minipage}\par\medskip
\begin{solution}
图\ref{mp5:fig:l5-e9-q}中的位置关系给出
\[
\vct v_C=(\dot X+\dot s\cos\alpha,\,-\dot s\sin\alpha),
\qquad \omega_z=-\dot s/R.
\]
计入斜面平动，
\begin{align*}
T&=\frac12(M+m)\dot X^2+m\cos\alpha\,\dot X\dot s
+\frac12m(1+\eta)\dot s^2,\\
U&=-mgs\sin\alpha.
\end{align*}
分别对两个独立坐标求拉格朗日方程，得到
\[
(M+m)\ddot X+m\cos\alpha\,\ddot s=0,
\qquad
m\cos\alpha\,\ddot X+m(1+\eta)\ddot s=mg\sin\alpha.
\]
第一式来自水平方向没有外力，第二式则把斜面平移对圆柱相对运动的影响保留在交叉项中。即使斜面初始静止，也不能在运动开始以后直接令 $\ddot X=0$。

$X$ 是循环坐标。初始水平动量为零，故
\[
(M+m)\dot X+m\cos\alpha\,\dot s=0.
\]
令 $\rho=m/(M+m)\in(0,1)$、$D=1+\eta-\rho\cos^2\alpha$。将动能配方，
\[
T=\frac{M+m}{2}(\dot X+\rho\cos\alpha\,\dot s)^2
+\frac m2(1+\eta-\rho\cos^2\alpha)\dot s^2.
\]
零水平动量使第一个平方项为零。由于 $D\ge1+\eta-\rho>0$，约化动能始终正定，所得加速度的分母也始终为正。
因此
\[
\dot X=-\rho\cos\alpha\,\dot s,\qquad T_{\rm red}=\frac12mD\dot s^2.
\]

\begin{figure}[H]\centering
\begin{tikzpicture}
\draw[axis](0,0)--(2,0)node[midway,above]{$\dot s\cos\alpha$};\draw[axis](2,0)--(2,-1.1547)node[right]{$-\dot s\sin\alpha$};\draw[axis](2,-1.1547)--(1,-1.1547)node[midway,below]{$\dot X<0$};\draw[axis](0,0)--(1,-1.1547)node[midway,left]{$\vct v_C$};\draw[guide](0,0)--(2,-1.1547);\node at(1,1){速度首尾相接};\begin{scope}[shift={(6,0.1)},scale=1]
\draw[curve](-1.7,.9815)--(1.7,-.9815);\draw[curve](.4,.69282)circle(.8);\fill(0,0)circle(1.3pt)node[below]{$Q$};\fill(.4,.69282)circle(1.3pt)node[above left]{$C$};\draw[axis](0,0)--(.6,1.0392)node[above]{$N$};\draw[axis](0,0)--(-.866,.5)node[left]{$f$};\draw[axis](.4,.69282)--(.4,-.5)node[right]{$mg$};\end{scope}

\end{tikzpicture}
\caption{左图将速度分量首尾相接，比例取 $M=m,\eta=1/2,\alpha=30^\circ$；右图的受力箭头只示方向，摩擦沿坡向上。}\label{mp5:fig:l5-e9-s}
\end{figure}

图\ref{mp5:fig:l5-e9-s}左图把质心速度的三部分首尾相接。斜面向左运动，减小圆柱在地面系中的水平速度。由第一条动力学方程，$\ddot X=-\rho\cos\alpha\ddot s$，代入第二条得到 $mD\ddot s=mg\sin\alpha$；这一消元在释放瞬间也成立，不需要约去速度。再用全系统机械能守恒求速度，得到
\[
\dot s^2=\frac{2gs\sin\alpha}{D},\quad
\boxed{\ddot s=\frac{g\sin\alpha}{D},\qquad
\ddot X=-\frac{\rho g\sin\alpha\cos\alpha}{D}.}
\]
单独圆柱的机械能一般不守恒，因为斜面上的接触点在地面系中有速度；将斜面包括进系统后，接触力的总功率才抵消。

右图中法向单位矢量为 $\vct n=(\sin\alpha,\cos\alpha)$。质心加速度在这个方向的分量是 $\ddot X\sin\alpha$，故
\[
N-mg\cos\alpha=m\ddot X\sin\alpha,
\quad
N=\frac{mg\cos\alpha(1+\eta-\rho)}D>0.
\]
绕质心的角加速度大小为 $\ddot s/R$，所以沿坡向上的摩擦为
\[
f=\frac I{R^2}\ddot s=\frac{\eta mg\sin\alpha}{D}.
\]
由 $f\le\mu N$ 得
\[
\boxed{\mu\ge\frac{\eta\tan\alpha}{1+\eta-\rho}.}
\]
$M\to\infty$ 时 $\rho\to0$，恢复固定斜面的结果；有限质量斜面使相对下滑加速度增大，也提高维持纯滚动所需的静摩擦因数。
圆柱的地面加速度分量进一步为
\[
a_{Cx}=(1-\rho)\cos\alpha\,\frac{g\sin\alpha}{D},\qquad
a_{Cy}=-\sin\alpha\,\frac{g\sin\alpha}{D}.
\]
它与相对下滑加速度 $\ddot s$ 不同。

若再检查斜面与水平地面的接触，设地面总支持力为 $N_0$，整体竖直动量方程给出
\[
N_0-(M+m)g=-m\ddot s\sin\alpha,
\qquad
N_0=\frac{[M+(M+m)\eta]g}{D}>0.
\]
因此题设运动中地面接触不会先失效；圆柱与斜面的接触已由前述 $N>0$ 保证，剩下的限制就是所需静摩擦是否在允许范围内。

\end{solution}

\par\medskip\noindent\begin{minipage}{\textwidth}
\begin{example}[自由转动圆环与角动量约化（选做）]\label{mp5:example:l5-e10}
半径为 $R$ 的竖直圆环可绕竖直直径无摩擦转动。除滑珠外，转动部分对该轴的惯量为常量 $I_0>0$。质量为 $m$ 的滑珠套在圆环上，角坐标 $\theta$ 自最低点起算，圆环方位角为 $\phi$，见图\ref{mp5:fig:l5-e10-q}。无外界竖直轴力矩，总角动量为 $J$。求约化势能、最低点的稳定条件，以及偏离最低点的平衡与小振动频率。
\end{example}

\begin{figure}[H]\centering
\begin{tikzpicture}
\draw[curve](0,0)circle(1.8);\draw[axis](0,-2.2)--(0,2.7)node[above]{$z$};\draw[guide](0,0)--(0,-1.8);\draw(0,0)--(1.4744736797201852,-1.032437585431883);\fill(0,0)circle(1.3pt)node[above left]{$O$};
\fill(1.4744736797201852,-1.032437585431883)circle(1.3pt)node[right]{$P$};
\draw(0,-.65)arc(-90:-35:.65);\node at(.35,-.75){$\theta$};\node[above]at(.85,-.45){$R$};\draw[curve](0,2.2)ellipse(.8 and .14);\node[right]at(.9,2.2){$I_0$};\draw[axis](.8,2.2)arc(0:160:.8 and .14);\node[left]at(-1,2.35){$\dot\phi$};
\end{tikzpicture}
\caption{圆环可自由改变转速，$I_0$ 为除滑珠外的轴向惯量；$\theta$ 描述珠子在环内的位置，$\phi$ 为环面方位角。}\label{mp5:fig:l5-e10-q}
\end{figure}
\end{minipage}\par\medskip
\begin{solution}
滑珠到竖直轴的距离为 $R\sin\theta$。记
\[
D(\theta)=I_0+mR^2\sin^2\theta.
\]
图\ref{mp5:fig:l5-e10-q}中的两个角速度产生正交的速度分量，故
\[
T=\frac12mR^2\dot\theta^2+\frac12D(\theta)\dot\phi^2,
\qquad U=mgR(1-\cos\theta).
\]
$\phi$ 是循环坐标，$J=D\dot\phi$。先保留两个角坐标，珠子的方程为
\[
mR^2\ddot\theta-\frac12D'(\theta)\dot\phi^2+mgR\sin\theta=0.
\]
由 $J=D\dot\phi$ 代入此方程，便得到 $mR^2\ddot\theta+V_J'(\theta)=0$。若从拉格朗日函数直接消元，应使用固定角动量下的劳斯函数
\[
\mathcal R=\left.(T-U-J\dot\phi)\right|_{\dot\phi=J/D}
=\frac12mR^2\dot\theta^2-U-\frac{J^2}{2D}.
\]
这里 $J^2/(2D)$ 在有效势能中取正号。固定角动量的约化须从循环动量方程出发，或先构造劳斯函数；直接在原拉格朗日函数中替换速度，会漏掉 $-J\dot\phi$ 对符号的影响。

在固定 $J$ 的条件下，能量约化为
\[
E=\frac12mR^2\dot\theta^2+V_J(\theta),\qquad
\boxed{V_J=mgR(1-\cos\theta)+\frac{J^2}{2D(\theta)}.}
\]
固定 $J$ 求导时，角动量项中 $D(\theta)$ 随位置改变，故
\[
\frac{\dd}{\dd\theta}\frac{J^2}{2D}
=-\frac{J^2D'}{2D^2},\qquad
D'=2mR^2\sin\theta\cos\theta.
\]
与重力势能的导数相加并提出 $\sin\theta$，
\[
V_J'=\sin\theta\left[mgR-
\frac{J^2mR^2\cos\theta}{D^2}\right].
\]
最低点 $\theta=0$ 的曲率为
\[
V_J''(0)=mgR-\frac{J^2mR^2}{I_0^2}.
\]
令 $J_c=I_0\sqrt{g/R}$，则 $|J|<J_c$ 时最低点稳定，小振动频率满足
\[
\omega_0^2=\frac gR-\frac{J^2}{I_0^2}.
\]

\begin{figure}[H]\centering
\begin{tikzpicture}
\begin{scope}[xscale=1.7,yscale=2]\draw[axis](-1.7,0)--(1.7,0)node[right]{$\theta$};
\draw[axis](0,-0.35)--(0,1.2)node[above]{$u$};
\draw[curve] plot coordinates{(-1.60000,0.90675) (-1.58400,0.89071) (-1.56800,0.87470) (-1.55200,0.85873) (-1.53600,0.84278) (-1.52000,0.82688) (-1.50400,0.81103) (-1.48800,0.79522) (-1.47200,0.77946) (-1.45600,0.76376) (-1.44000,0.74813) (-1.42400,0.73255) (-1.40800,0.71705) (-1.39200,0.70162) (-1.37600,0.68627) (-1.36000,0.67100) (-1.34400,0.65582) (-1.32800,0.64073) (-1.31200,0.62573) (-1.29600,0.61083) (-1.28000,0.59604) (-1.26400,0.58135) (-1.24800,0.56677) (-1.23200,0.55231) (-1.21600,0.53797) (-1.20000,0.52375) (-1.18400,0.50966) (-1.16800,0.49570) (-1.15200,0.48188) (-1.13600,0.46820) (-1.12000,0.45466) (-1.10400,0.44127) (-1.08800,0.42804) (-1.07200,0.41496) (-1.05600,0.40204) (-1.04000,0.38928) (-1.02400,0.37669) (-1.00800,0.36428) (-0.99200,0.35204) (-0.97600,0.33997) (-0.96000,0.32809) (-0.94400,0.31640) (-0.92800,0.30489) (-0.91200,0.29358) (-0.89600,0.28246) (-0.88000,0.27155) (-0.86400,0.26083) (-0.84800,0.25032) (-0.83200,0.24002) (-0.81600,0.22993) (-0.80000,0.22005) (-0.78400,0.21039) (-0.76800,0.20095) (-0.75200,0.19173) (-0.73600,0.18273) (-0.72000,0.17395) (-0.70400,0.16540) (-0.68800,0.15708) (-0.67200,0.14899) (-0.65600,0.14113) (-0.64000,0.13350) (-0.62400,0.12610) (-0.60800,0.11893) (-0.59200,0.11200) (-0.57600,0.10530) (-0.56000,0.09883) (-0.54400,0.09259) (-0.52800,0.08659) (-0.51200,0.08081) (-0.49600,0.07526) (-0.48000,0.06994) (-0.46400,0.06485) (-0.44800,0.05998) (-0.43200,0.05533) (-0.41600,0.05089) (-0.40000,0.04668) (-0.38400,0.04267) (-0.36800,0.03888) (-0.35200,0.03529) (-0.33600,0.03190) (-0.32000,0.02870) (-0.30400,0.02571) (-0.28800,0.02290) (-0.27200,0.02027) (-0.25600,0.01783) (-0.24000,0.01556) (-0.22400,0.01346) (-0.20800,0.01153) (-0.19200,0.00977) (-0.17600,0.00816) (-0.16000,0.00671) (-0.14400,0.00541) (-0.12800,0.00425) (-0.11200,0.00324) (-0.09600,0.00237) (-0.08000,0.00164) (-0.06400,0.00105) (-0.04800,0.00059) (-0.03200,0.00026) (-0.01600,0.00007) (0.00000,0.00000) (0.01600,0.00007) (0.03200,0.00026) (0.04800,0.00059) (0.06400,0.00105) (0.08000,0.00164) (0.09600,0.00237) (0.11200,0.00324) (0.12800,0.00425) (0.14400,0.00541) (0.16000,0.00671) (0.17600,0.00816) (0.19200,0.00977) (0.20800,0.01153) (0.22400,0.01346) (0.24000,0.01556) (0.25600,0.01783) (0.27200,0.02027) (0.28800,0.02290) (0.30400,0.02571) (0.32000,0.02870) (0.33600,0.03190) (0.35200,0.03529) (0.36800,0.03888) (0.38400,0.04267) (0.40000,0.04668) (0.41600,0.05089) (0.43200,0.05533) (0.44800,0.05998) (0.46400,0.06485) (0.48000,0.06994) (0.49600,0.07526) (0.51200,0.08081) (0.52800,0.08659) (0.54400,0.09259) (0.56000,0.09883) (0.57600,0.10530) (0.59200,0.11200) (0.60800,0.11893) (0.62400,0.12610) (0.64000,0.13350) (0.65600,0.14113) (0.67200,0.14899) (0.68800,0.15708) (0.70400,0.16540) (0.72000,0.17395) (0.73600,0.18273) (0.75200,0.19173) (0.76800,0.20095) (0.78400,0.21039) (0.80000,0.22005) (0.81600,0.22993) (0.83200,0.24002) (0.84800,0.25032) (0.86400,0.26083) (0.88000,0.27155) (0.89600,0.28246) (0.91200,0.29358) (0.92800,0.30489) (0.94400,0.31640) (0.96000,0.32809) (0.97600,0.33997) (0.99200,0.35204) (1.00800,0.36428) (1.02400,0.37669) (1.04000,0.38928) (1.05600,0.40204) (1.07200,0.41496) (1.08800,0.42804) (1.10400,0.44127) (1.12000,0.45466) (1.13600,0.46820) (1.15200,0.48188) (1.16800,0.49570) (1.18400,0.50966) (1.20000,0.52375) (1.21600,0.53797) (1.23200,0.55231) (1.24800,0.56677) (1.26400,0.58135) (1.28000,0.59604) (1.29600,0.61083) (1.31200,0.62573) (1.32800,0.64073) (1.34400,0.65582) (1.36000,0.67100) (1.37600,0.68627) (1.39200,0.70162) (1.40800,0.71705) (1.42400,0.73255) (1.44000,0.74813) (1.45600,0.76376) (1.47200,0.77946) (1.48800,0.79522) (1.50400,0.81103) (1.52000,0.82688) (1.53600,0.84278) (1.55200,0.85873) (1.56800,0.87470) (1.58400,0.89071) (1.60000,0.90675)};
\draw[curve,dashed] plot coordinates{(-1.60000,0.77931) (-1.58400,0.76323) (-1.56800,0.74720) (-1.55200,0.73125) (-1.53600,0.71536) (-1.52000,0.69955) (-1.50400,0.68381) (-1.48800,0.66816) (-1.47200,0.65259) (-1.45600,0.63711) (-1.44000,0.62172) (-1.42400,0.60643) (-1.40800,0.59125) (-1.39200,0.57617) (-1.37600,0.56121) (-1.36000,0.54636) (-1.34400,0.53163) (-1.32800,0.51702) (-1.31200,0.50255) (-1.29600,0.48820) (-1.28000,0.47400) (-1.26400,0.45994) (-1.24800,0.44603) (-1.23200,0.43226) (-1.21600,0.41866) (-1.20000,0.40521) (-1.18400,0.39193) (-1.16800,0.37881) (-1.15200,0.36587) (-1.13600,0.35311) (-1.12000,0.34053) (-1.10400,0.32814) (-1.08800,0.31594) (-1.07200,0.30393) (-1.05600,0.29212) (-1.04000,0.28052) (-1.02400,0.26912) (-1.00800,0.25794) (-0.99200,0.24697) (-0.97600,0.23622) (-0.96000,0.22569) (-0.94400,0.21539) (-0.92800,0.20532) (-0.91200,0.19548) (-0.89600,0.18588) (-0.88000,0.17652) (-0.86400,0.16740) (-0.84800,0.15853) (-0.83200,0.14990) (-0.81600,0.14153) (-0.80000,0.13341) (-0.78400,0.12555) (-0.76800,0.11794) (-0.75200,0.11059) (-0.73600,0.10351) (-0.72000,0.09668) (-0.70400,0.09011) (-0.68800,0.08381) (-0.67200,0.07776) (-0.65600,0.07198) (-0.64000,0.06646) (-0.62400,0.06120) (-0.60800,0.05620) (-0.59200,0.05145) (-0.57600,0.04696) (-0.56000,0.04271) (-0.54400,0.03872) (-0.52800,0.03496) (-0.51200,0.03145) (-0.49600,0.02817) (-0.48000,0.02512) (-0.46400,0.02230) (-0.44800,0.01969) (-0.43200,0.01729) (-0.41600,0.01510) (-0.40000,0.01310) (-0.38400,0.01129) (-0.36800,0.00966) (-0.35200,0.00819) (-0.33600,0.00689) (-0.32000,0.00574) (-0.30400,0.00474) (-0.28800,0.00386) (-0.27200,0.00311) (-0.25600,0.00246) (-0.24000,0.00192) (-0.22400,0.00147) (-0.20800,0.00110) (-0.19200,0.00081) (-0.17600,0.00058) (-0.16000,0.00040) (-0.14400,0.00026) (-0.12800,0.00016) (-0.11200,0.00010) (-0.09600,0.00005) (-0.08000,0.00003) (-0.06400,0.00001) (-0.04800,0.00000) (-0.03200,0.00000) (-0.01600,0.00000) (0.00000,0.00000) (0.01600,0.00000) (0.03200,0.00000) (0.04800,0.00000) (0.06400,0.00001) (0.08000,0.00003) (0.09600,0.00005) (0.11200,0.00010) (0.12800,0.00016) (0.14400,0.00026) (0.16000,0.00040) (0.17600,0.00058) (0.19200,0.00081) (0.20800,0.00110) (0.22400,0.00147) (0.24000,0.00192) (0.25600,0.00246) (0.27200,0.00311) (0.28800,0.00386) (0.30400,0.00474) (0.32000,0.00574) (0.33600,0.00689) (0.35200,0.00819) (0.36800,0.00966) (0.38400,0.01129) (0.40000,0.01310) (0.41600,0.01510) (0.43200,0.01729) (0.44800,0.01969) (0.46400,0.02230) (0.48000,0.02512) (0.49600,0.02817) (0.51200,0.03145) (0.52800,0.03496) (0.54400,0.03872) (0.56000,0.04271) (0.57600,0.04696) (0.59200,0.05145) (0.60800,0.05620) (0.62400,0.06120) (0.64000,0.06646) (0.65600,0.07198) (0.67200,0.07776) (0.68800,0.08381) (0.70400,0.09011) (0.72000,0.09668) (0.73600,0.10351) (0.75200,0.11059) (0.76800,0.11794) (0.78400,0.12555) (0.80000,0.13341) (0.81600,0.14153) (0.83200,0.14990) (0.84800,0.15853) (0.86400,0.16740) (0.88000,0.17652) (0.89600,0.18588) (0.91200,0.19548) (0.92800,0.20532) (0.94400,0.21539) (0.96000,0.22569) (0.97600,0.23622) (0.99200,0.24697) (1.00800,0.25794) (1.02400,0.26912) (1.04000,0.28052) (1.05600,0.29212) (1.07200,0.30393) (1.08800,0.31594) (1.10400,0.32814) (1.12000,0.34053) (1.13600,0.35311) (1.15200,0.36587) (1.16800,0.37881) (1.18400,0.39193) (1.20000,0.40521) (1.21600,0.41866) (1.23200,0.43226) (1.24800,0.44603) (1.26400,0.45994) (1.28000,0.47400) (1.29600,0.48820) (1.31200,0.50255) (1.32800,0.51702) (1.34400,0.53163) (1.36000,0.54636) (1.37600,0.56121) (1.39200,0.57617) (1.40800,0.59125) (1.42400,0.60643) (1.44000,0.62172) (1.45600,0.63711) (1.47200,0.65259) (1.48800,0.66816) (1.50400,0.68381) (1.52000,0.69955) (1.53600,0.71536) (1.55200,0.73125) (1.56800,0.74720) (1.58400,0.76323) (1.60000,0.77931)};
\draw[curve,densely dotted] plot coordinates{(-1.60000,0.53941) (-1.58400,0.52325) (-1.56800,0.50721) (-1.55200,0.49129) (-1.53600,0.47551) (-1.52000,0.45986) (-1.50400,0.44435) (-1.48800,0.42898) (-1.47200,0.41376) (-1.45600,0.39869) (-1.44000,0.38378) (-1.42400,0.36903) (-1.40800,0.35444) (-1.39200,0.34003) (-1.37600,0.32579) (-1.36000,0.31173) (-1.34400,0.29785) (-1.32800,0.28416) (-1.31200,0.27067) (-1.29600,0.25738) (-1.28000,0.24429) (-1.26400,0.23141) (-1.24800,0.21874) (-1.23200,0.20629) (-1.21600,0.19407) (-1.20000,0.18207) (-1.18400,0.17031) (-1.16800,0.15879) (-1.15200,0.14751) (-1.13600,0.13648) (-1.12000,0.12570) (-1.10400,0.11518) (-1.08800,0.10493) (-1.07200,0.09494) (-1.05600,0.08523) (-1.04000,0.07579) (-1.02400,0.06664) (-1.00800,0.05777) (-0.99200,0.04919) (-0.97600,0.04091) (-0.96000,0.03293) (-0.94400,0.02525) (-0.92800,0.01788) (-0.91200,0.01082) (-0.89600,0.00407) (-0.88000,-0.00236) (-0.86400,-0.00847) (-0.84800,-0.01426) (-0.83200,-0.01972) (-0.81600,-0.02486) (-0.80000,-0.02967) (-0.78400,-0.03415) (-0.76800,-0.03830) (-0.75200,-0.04212) (-0.73600,-0.04562) (-0.72000,-0.04878) (-0.70400,-0.05161) (-0.68800,-0.05412) (-0.67200,-0.05631) (-0.65600,-0.05817) (-0.64000,-0.05972) (-0.62400,-0.06096) (-0.60800,-0.06189) (-0.59200,-0.06253) (-0.57600,-0.06287) (-0.56000,-0.06292) (-0.54400,-0.06270) (-0.52800,-0.06221) (-0.51200,-0.06146) (-0.49600,-0.06047) (-0.48000,-0.05924) (-0.46400,-0.05779) (-0.44800,-0.05614) (-0.43200,-0.05430) (-0.41600,-0.05228) (-0.40000,-0.05011) (-0.38400,-0.04779) (-0.36800,-0.04535) (-0.35200,-0.04280) (-0.33600,-0.04017) (-0.32000,-0.03748) (-0.30400,-0.03474) (-0.28800,-0.03197) (-0.27200,-0.02921) (-0.25600,-0.02646) (-0.24000,-0.02375) (-0.22400,-0.02110) (-0.20800,-0.01853) (-0.19200,-0.01606) (-0.17600,-0.01370) (-0.16000,-0.01149) (-0.14400,-0.00942) (-0.12800,-0.00753) (-0.11200,-0.00583) (-0.09600,-0.00432) (-0.08000,-0.00302) (-0.06400,-0.00194) (-0.04800,-0.00110) (-0.03200,-0.00049) (-0.01600,-0.00012) (0.00000,0.00000) (0.01600,-0.00012) (0.03200,-0.00049) (0.04800,-0.00110) (0.06400,-0.00194) (0.08000,-0.00302) (0.09600,-0.00432) (0.11200,-0.00583) (0.12800,-0.00753) (0.14400,-0.00942) (0.16000,-0.01149) (0.17600,-0.01370) (0.19200,-0.01606) (0.20800,-0.01853) (0.22400,-0.02110) (0.24000,-0.02375) (0.25600,-0.02646) (0.27200,-0.02921) (0.28800,-0.03197) (0.30400,-0.03474) (0.32000,-0.03748) (0.33600,-0.04017) (0.35200,-0.04280) (0.36800,-0.04535) (0.38400,-0.04779) (0.40000,-0.05011) (0.41600,-0.05228) (0.43200,-0.05430) (0.44800,-0.05614) (0.46400,-0.05779) (0.48000,-0.05924) (0.49600,-0.06047) (0.51200,-0.06146) (0.52800,-0.06221) (0.54400,-0.06270) (0.56000,-0.06292) (0.57600,-0.06287) (0.59200,-0.06253) (0.60800,-0.06189) (0.62400,-0.06096) (0.64000,-0.05972) (0.65600,-0.05817) (0.67200,-0.05631) (0.68800,-0.05412) (0.70400,-0.05161) (0.72000,-0.04878) (0.73600,-0.04562) (0.75200,-0.04212) (0.76800,-0.03830) (0.78400,-0.03415) (0.80000,-0.02967) (0.81600,-0.02486) (0.83200,-0.01972) (0.84800,-0.01426) (0.86400,-0.00847) (0.88000,-0.00236) (0.89600,0.00407) (0.91200,0.01082) (0.92800,0.01788) (0.94400,0.02525) (0.96000,0.03293) (0.97600,0.04091) (0.99200,0.04919) (1.00800,0.05777) (1.02400,0.06664) (1.04000,0.07579) (1.05600,0.08523) (1.07200,0.09494) (1.08800,0.10493) (1.10400,0.11518) (1.12000,0.12570) (1.13600,0.13648) (1.15200,0.14751) (1.16800,0.15879) (1.18400,0.17031) (1.20000,0.18207) (1.21600,0.19407) (1.23200,0.20629) (1.24800,0.21874) (1.26400,0.23141) (1.28000,0.24429) (1.29600,0.25738) (1.31200,0.27067) (1.32800,0.28416) (1.34400,0.29785) (1.36000,0.31173) (1.37600,0.32579) (1.39200,0.34003) (1.40800,0.35444) (1.42400,0.36903) (1.44000,0.38378) (1.45600,0.39869) (1.47200,0.41376) (1.48800,0.42898) (1.50400,0.44435) (1.52000,0.45986) (1.53600,0.47551) (1.55200,0.49129) (1.56800,0.50721) (1.58400,0.52325) (1.60000,0.53941)};
\node[above left]at(0,0){$0$};\draw[curve](2,1.05)--(2.25,1.05)node[right]{$j=0.7$};\draw[curve,dashed](2,.8)--(2.25,.8)node[right]{$j=1$};\draw[curve,densely dotted](2,.55)--(2.25,.55)node[right]{$j=1.4$};\end{scope}
\end{tikzpicture}
\caption{取 $\mu=mR^2/I_0=1$、$j=J/J_c$，纵轴 $u=[V_J(\theta)-V_J(0)]/(mgR)$。虚线在原点有四阶极小值，点线有两个偏离原点的极小值。}\label{mp5:fig:l5-e10-s}
\end{figure}

图\ref{mp5:fig:l5-e10-s}取 $\mu=mR^2/I_0=1$，画出不同角动量下的约化势能。判断临界情形须保留四次项。对一般$\mu=mR^2/I_0$，先由$\sin\theta=\theta-\theta^3/6+O(\theta^5)$得到
\[
\sin^2\theta=\theta^2-\frac13\theta^4+O(\theta^6).
\]
令$z=\mu\sin^2\theta$，在$1/(1+z)=1-z+z^2+O(z^3)$中保留至$\theta^4$：$-z$贡献$-\mu\theta^2+\mu\theta^4/3$，$z^2$贡献$\mu^2\theta^4$。因此两项势能所需的展开分别为
\[
\begin{aligned}
mgR(1-\cos\theta)&=mgR\left(\frac{\theta^2}{2}-\frac{\theta^4}{24}\right)+O(\theta^6),\\
\frac1{D(\theta)}&=\frac1{I_0}\left[1-\mu\theta^2+
\left(\mu^2+\frac\mu3\right)\theta^4\right]+O(\theta^6).
\end{aligned}
\]
代入$J^2=I_0^2g/R$，角动量项的二次系数为$-mgR/2$，四次系数为$mgR(\mu/2+1/6)$。它与重力势能相加，二次项相消，得到
\[
\left.V_J(\theta)-V_J(0)\right|_{|J|=J_c}
=mgR\left(\frac18+\frac\mu2\right)\theta^4+O(\theta^6),
\]
四次项为正，所以临界最低点仍稳定，此时回复作用从三阶项开始，线性简谐近似不再适用。$|J|>J_c$ 时最低点不稳定。

这里的平衡指 $\theta$ 不变、环面匀速转动的相对平衡。非零平衡满足 $mgR=J^2mR^2\cos\theta_e/D_e^2$，因此必有 $\cos\theta_e>0$。在 $0<\theta<\pi/2$ 上，
\[
\frac{\dd}{\dd\theta}\frac{D^2}{\cos\theta}
=\frac{2DD'}{\cos\theta}+\frac{D^2\sin\theta}{\cos^2\theta}>0.
\]
该函数从 $I_0^2$ 严格增大到无穷，故且仅当 $|J|>J_c$，存在一对对称平衡 $\pm\theta_e$。令 $\Omega_e=J/D_e$，也可写成
\[
\cos\theta_e=\frac g{\Omega_e^2R}.
\]
在$V_J'=\sin\theta\,B(\theta)$中，记$B=mgR-J^2mR^2\cos\theta/D^2$。由于$B(\theta_e)=0$，有$V_J''(\theta_e)=\sin\theta_e B'(\theta_e)$，而
\[
B'=J^2mR^2\left(\frac{\sin\theta}{D^2}
+\frac{2\cos\theta D'}{D^3}\right),\qquad
D'=2mR^2\sin\theta\cos\theta.
\]
代入并用$J^2/D_e^2=\Omega_e^2$，得
\[
V_J''(\theta_e)=mR^2\Omega_e^2\sin^2\theta_e
\left(1+\frac{4mR^2\cos^2\theta_e}{D_e}\right)>0.
\]
所以这两个平衡稳定，沿环的小振动频率为
\[
\boxed{\omega_\theta^2=\Omega_e^2\sin^2\theta_e
\left(1+\frac{4mR^2\cos^2\theta_e}{D_e}\right).}
\]
扰动时珠子改变轴向惯量，环的转速随之变化，括号内的附加项正来自角动量守恒。最高点 $\theta=\pi$ 的曲率为 $-mgR-J^2mR^2/I_0^2<0$，始终不稳定。

在非零相对平衡附近，$\delta\dot\phi=-JD'_e\delta\theta/D_e^2$。珠子偏移会改变环速，所以沿环的振动刚度含有上述附加项。$\omega_\theta$ 描述珠子沿环偏角的振动，整体转动方位仍可任意改变，并无独立的方位回复作用。

\end{solution}

\section{柔性系统与杆架的虚功}
软绳和杆架的位形变化常比各处内力更容易求出。长度约束联系各载荷点的虚位移，总虚功给出平衡；局部拉力再由截取部分求解。

求平衡力时，先由约束确定允许的微小位移，再计算各力的虚功。求稳定性还须列全独立的构形变化。对有外部电源或电流源的系统，能量核算还要包括外源做功。

\par\medskip\noindent\begin{minipage}{\textwidth}
\begin{example}[固定圆柱上的四分之一圆周软绳]\label{mp5:hw:l5-32a}
质量$m$、长度$\pi R/2$的均匀软绳贴在光滑固定圆柱的外表面，恰占从最高点$A$到右侧点$B$的四分之一圆弧。$B$端自由，在$A$端施加水平向左的力$F$保持平衡，见图\ref{mp5:fig:hw:l5-32a-q}。求$F$，并用绳上张力核实。
\end{example}

\begin{figure}[H]\centering
\begin{tikzpicture}[>=Stealth,line cap=round,line join=round]

\draw[curve](0,0)circle(1.7);\draw[very thick](0,1.7)arc(90:0:1.7);\fill(0,1.7)circle(1.5pt)node[above right]{$A$};\fill(1.7,0)circle(1.5pt)node[right]{$B$};\draw[axis](0,1.7)--(-1.5,1.7)node[midway,above]{$F$};\draw[guide](0,0)--(1.202,1.202)node[midway,above left]{$R$};\node[below left]at(0,0){$O$};\draw(.55,0)arc(0:90:.55);\node at(.85,.13){$\pi/2$};

\end{tikzpicture}
\caption{软绳在圆柱右上方，拉力沿最高点处的切线}\label{mp5:fig:hw:l5-32a-q}
\end{figure}

\end{minipage}\par\medskip
\begin{solution}
设线密度$\lambda=2m/(\pi R)$。令绳整体沿圆周向右下方移动微小弧长$\delta s=R\delta\phi$，图\ref{mp5:fig:hw:l5-32a-s}中的每个材料微元都沿局部切线移动。以从最高点向右量起的角为$\theta$，重力虚功为
\[
 \delta W_g=\int_0^{\pi/2}\lambda R\dd\theta\,g\sin\theta\,\delta s
            =\lambda Rg\delta s[-\cos\theta]_0^{\pi/2}
            =\lambda Rg\,\delta s.
\]
光滑支持不做虚功，外拉力做功$-F\delta s$。故
\[
 \boxed{F=\lambda Rg=\frac{2mg}{\pi}}.
\]
再取弧长微元列切向平衡：$\dd T/\dd s=-\lambda g\sin\theta$。用 $\dd s=R\dd\theta$，并从自由端 $T(\pi/2)=0$ 积分，
\[
 T(\theta)=-\lambda gR\int_{\pi/2}^\theta\sin u\dd u
 =\lambda gR[\cos u]_{\pi/2}^\theta=\lambda gR\cos\theta.
\]
最高点张力正是 $F$。若圆柱单位弧长的向外支持力为 $q_n$，法向平衡为 $q_n-\lambda g\cos\theta-T/R=0$，所以 $q_n=2\lambda g\cos\theta\ge0$，所设外侧接触可成立。各微元的位移方向随位置改变，因此与 $F$ 比较的应是重力虚功的积分。


\begin{figure}[H]\centering
\begin{tikzpicture}[>=Stealth,line cap=round,line join=round]
\coordinate(O)at(0,-1.8);\coordinate(P)at({1.8*sin(40)},{-1.8+1.8*cos(40)});
\draw[curve](0,0)arc(90:0:1.8);\node[left]at(0,0){$A$};\fill(O)circle(1.2pt)node[below left]{$O$};
\draw[guide](O)--(0,0);\draw[guide](O)--(P);\draw(0,-1.25)arc(90:50:.55);\node at(.23,-1.12){$\theta$};
\fill(P)circle(1.5pt)node[above right]{$\dd m$};\draw[axis](P)--++({sin(50)},{-cos(50)})node[right]{$\delta\vct r$};\draw[axis](P)--++(0,-1)node[below]{$g\dd m$};
\node[right,align=left]at(2.7,-.55){$\delta W_g=g\dd m\,\delta s\sin\theta$\\[6pt]$\theta$从最高点半径量起};
\end{tikzpicture}
\caption{微元的重力与切向虚位移。$\theta$从半径$OA$量至$OP$，虚位移的向下分量为$\delta s\sin\theta$。}\label{mp5:fig:hw:l5-32a-s}
\end{figure}

\end{solution}

\par\medskip\noindent\begin{minipage}{\textwidth}
\begin{example}[圆形轨道上车列的内部拉力]\label{mp5:hw:l5-32b}
质量$m$、弧长$l$的均匀车列在半径$R$的光滑竖直圆形轨道上以恒定速率$v$运动。用从最低点向上量起的圆心角$\theta$标记位置，令$\theta_0$处的车钩拉力为$T_0$，见图\ref{mp5:fig:hw:l5-32b-q}。求内部拉力分布；再取尾部在最低点且无外拉力的情形。车列由轨道约束在圆周上，维持给定速率的驱动力只施于所讨论车列的端点，内部没有额外的切向驱动力。
\end{example}

\begin{figure}[H]\centering
\begin{tikzpicture}[>=Stealth,line cap=round,line join=round]

\draw[curve](0,0)circle(1.8);\draw[very thick](0,-1.8)arc(-90:135:1.8);\foreach \ang in {-75,-45,-15,15,45,75,105,135}{\fill({1.8*cos(\ang)},{1.8*sin(\ang)})circle(2pt);}\draw[guide](0,0)--(1.559,-.9)node[midway,above]{$R$};\draw[guide](0,0)--(0,-1.8);\draw(0,-.65)arc(-90:-30:.65);\node at(.5,-.8){$\theta$};\draw[axis](1.559,-.9)--(1.959,-.207)node[right]{$v$};\node[left]at(0,-1.8){$\theta_0=0$};

\end{tikzpicture}
\caption{车列的切向速率恒定，内部车钩力仍随重力切向分量改变}\label{mp5:fig:hw:l5-32b-q}
\end{figure}

\end{minipage}\par\medskip
\begin{solution}
线密度为$\lambda=m/l$。恒速使切向加速度为零，取相邻两车钩之间的微元，沿运动方向平衡：
\[
 T(\theta+\dd\theta)-T(\theta)-\lambda R\dd\theta\,g\sin\theta=0.
\]
除以 $\dd\theta$，得到 $T'=\lambda gR\sin\theta$。以 $T(\theta_0)=T_0$ 为边界条件，
\[
T(\theta)-T_0=\lambda gR\int_{\theta_0}^\theta\sin u\dd u
=\lambda gR[-\cos u]_{\theta_0}^\theta.
\]
所以
\[
 \boxed{T(\theta)=T_0+\lambda gR(\cos\theta_0-\cos\theta)}.
\]
若自由尾部在最低点，$T_0=0,\theta_0=0$，则
\[
 \boxed{T(\theta)=\frac{mgR}{l}(1-\cos\theta)}.
\]
也可让从 $\theta_0$ 至 $\theta$ 的所截车列整体虚移弧长 $\delta s$。该段重力势能为
\[
U_g=-\lambda gR^2\int_{\theta_0}^\theta\cos u\dd u
=-\lambda gR^2[\sin u]_{\theta_0}^\theta.
\]
两端角度同时增加 $\delta s/R$，所以
\[
 \delta U_g=\lambda gR(\cos\theta_0-\cos\theta)\delta s
           =[T(\theta)-T_0]\delta s.
\]
这从势能变化直接给出了同一张力差。图\ref{mp5:fig:hw:l5-32b-s}中的力差平衡微元重力的切向分量，因而这一张力差与$v$无关。

速率进入法向方程。若$Q$为轨道单位弧长沿向心方向的作用力，则张力的曲率合力为$T/R$，故
\[
 Q+\frac TR-\lambda g\cos\theta=\lambda\frac{v^2}{R}.
\]
单侧轨道须满足$Q\ge0$；带锁轮的轨道则可双向施力。绝对拉力还需由一个端点张力或驱动边界条件确定；质量、长度和速率本身不足以给出这一常数。


\begin{figure}[H]\centering
\begin{tikzpicture}[>=Stealth,line cap=round,line join=round]

\coordinate(P)at(0,0);\draw[curve](-1,-.15)..controls(0,-.4)and(.8,0)..(1.1,.5);\draw[axis](-.7,-.2)--(-1.6,-.35)node[left]{$T$};\draw[axis](.7,.2)--(1.55,.7)node[right]{$T+\dd T$};\draw[axis](P)--(0,-1.3)node[below]{$\lambda g\dd s$};\node at(0,1.2){$\dd T=\lambda g\sin\theta\,\dd s$};\node[right,align=left]at(3,.1){切向：$a_t=0$\\[6pt]法向：$a_n=v^2/R$};

\end{tikzpicture}
\caption{车列微元的两端拉力与重力；恒速条件只消去切向惯性项}\label{mp5:fig:hw:l5-32b-s}
\end{figure}

\end{solution}

\par\medskip\noindent\begin{minipage}{\textwidth}
\begin{example}[两级菱形绞架与圆柱]\label{mp5:hw:l5-33}
图\ref{mp5:fig:hw:l5-33-q}的对称绞架悬于$O$。上端两短杆$OA,OB$和下端两短杆$DP,EP$均长$l$、质量$m$；跨过中间交点的两根长杆$AE,BD$均长$2l$、质量$2m$，交点及端点均可转动。弹簧连接$A,B$，原长$2l$。下菱形内托住半径$r$、质量$M$的光滑圆柱，只与下方两短杆接触。所有杆与竖直线的锐角同为$\theta$。已知$\theta=45^\circ$平衡，求劲度系数$k$；取$0<r<l/2$保证该构形附近的接触关系不变。
\end{example}

\begin{figure}[H]\centering
\begin{tikzpicture}[>=Stealth,line cap=round,line join=round]

\coordinate(O)at(0,0);\coordinate(A)at(-.707,-.707);\coordinate(B)at(.707,-.707);\coordinate(D)at(-.707,-2.121);\coordinate(E)at(.707,-2.121);\coordinate(P)at(0,-2.828);\draw[thick](-1,0)--(1,0);\draw[thick](O)--(A)--(E)--(P)--(D)--(B)--(O);\draw[guide](0,0)--(0,-2.828);\draw[curve](0,-2.404)circle(.3);\draw[curve](A)--(-.55,-.707)--(-.45,-.62)--(-.3,-.79)--(-.15,-.62)--(0,-.79)--(.15,-.62)--(.3,-.79)--(.45,-.62)--(.55,-.707)--(B);\foreach \p in {O,A,B,D,E,P}{\fill(\p)circle(1.2pt);}\node[above]at(O){$O$};\node[left]at(A){$A$};\node[right]at(B){$B$};\node[left]at(D){$D$};\node[right]at(E){$E$};\node[below]at(P){$P$};\node[right]at(1.8,-.8){短杆：$l,m$};\node[right]at(1.8,-1.4){长杆：$2l,2m$};\node[right]at(1.8,-2.3){圆柱：$r,M$};

\draw[guide](A)--++(0,.85);\draw($(A)+(0,.4)$)arc(90:45:.4);\node at(-.46,-.2){$\theta$};
\end{tikzpicture}
\caption{两根长杆连续跨过中间铰点；四根短杆分别位于顶部与底部}\label{mp5:fig:hw:l5-33-q}
\end{figure}

\end{minipage}\par\medskip
\begin{solution}
只需一个形状坐标$\theta$。以向下距离$Y$计高，$A,B$在$Y=l\cos\theta$，中间交点在$2l\cos\theta$，$D,E$在$3l\cos\theta$，底点$P$在$4l\cos\theta$。

两根长杆质心均在中间交点，总质量$4m$；顶部两短杆质心高度为$l\cos\theta/2$，底部两短杆为$7l\cos\theta/2$。各重力势能相加，
\[
 U_{\rm rods}=-16mgl\cos\theta.
\]
图\ref{mp5:fig:hw:l5-33-s}中，圆柱中心高于底点$P$的距离为$r/\sin\theta$，故
\[
 Y_M=4l\cos\theta-r\csc\theta.
\]
弹簧长度为$2l\sin\theta$，于是
\[
 U=-16mgl\cos\theta-Mg(4l\cos\theta-r\csc\theta)
       +\frac k2(2l\sin\theta-2l)^2.
\]
弹簧势能的链式导数为
\[
\frac{\dd}{\dd\theta}\frac k2(2l\sin\theta-2l)^2
=k(2l\sin\theta-2l)(2l\cos\theta).
\]
再用 $\dd(\csc\theta)/\dd\theta=-\cos\theta/\sin^2\theta$，各项相加，
\[
 U'=(16m+4M)gl\sin\theta
      -Mgr\frac{\cos\theta}{\sin^2\theta}
      +4kl^2(\sin\theta-1)\cos\theta.
\]
代入 $\sin\theta=\cos\theta=1/\sqrt2$，平衡条件化为
\[
0=\frac{(16m+4M)gl}{\sqrt2}-\sqrt2Mgr
+2kl^2(1-\sqrt2).
\]
将弹簧项移到另一边，解得
\[
 \boxed{k=\frac{[(8m+2M)l-Mr]g}{(2-\sqrt2)l^2}}.
\]
每根长杆跨过中间铰点，两半段的总质量仍为$2m$，应按整根杆的质心计入重力势能。圆柱在底菱形中只压住下方两杆，还要求$r<l\sin\theta\cos\theta$；在平衡构形即为题给$r<l/2$。达到等号时会同时接触上方两杆，允许虚位移与受力模型均需改变。


\begin{figure}[H]\centering
\begin{tikzpicture}[>=Stealth,line cap=round,line join=round]

\draw[thick](-1.4,1.4)--(0,0)--(1.4,1.4);\draw[curve](0,.849)circle(.6);\draw[guide](0,0)--(-1.8,0);\draw[guide](0,.849)--(-1.8,.849);\draw[<->](-1.7,0)--(-1.7,.849)node[midway,left]{$r/\sin\theta$};\fill(0,0)circle(1.4pt)node[below]{$P$};\fill(0,.849)circle(1.4pt)node[right]{$C_M$};\draw[guide](0,0)--(0,1.7);\draw(0,.45)arc(90:135:.45);\node at(-.22,.59){$\theta$};\node[right,align=left]at(2.4,.8){$Y_P=4l\cos\theta$\\[6pt]$Y_M=Y_P-r/\sin\theta$\\[6pt]$AB=2l\sin\theta$};

\end{tikzpicture}
\caption{圆柱与下杆相切，$\theta$为杆与竖直线的夹角。圆心到$P$的高度差为$r/\sin\theta$。}\label{mp5:fig:hw:l5-33-s}
\end{figure}

\end{solution}

\par\medskip\noindent\begin{minipage}{\textwidth}
\begin{example}[载重绞架的形变与整体倾斜]\label{mp5:hw:l5-34}
图\ref{mp5:fig:hw:l5-34-q}中，两根长$3l$、质量$3m$的均匀杆从顶铰点$T$伸向左右载重端，各端悬质量$M$。在长杆距$T$为$l$处，分别以长$l$、质量$m$的短杆连接到下铰点$O$，$O$在支承台上固定。弹簧连接$O,T$，原长$l/\sqrt2$。对称构形中，长杆与竖直方向夹角$\theta=45^\circ$时平衡。(1)求$k$；(2)同时考虑形状变化与整体倾斜，求严格势能稳定条件。
\end{example}

\begin{figure}[H]\centering
\begin{tikzpicture}[>=Stealth,line cap=round,line join=round]

\coordinate(O)at(0,0);\coordinate(T)at(0,1.414);\coordinate(A)at(-.707,.707);\coordinate(B)at(.707,.707);\coordinate(L)at(-2.121,-.707);\coordinate(R)at(2.121,-.707);\draw[fill=gray!20](-.5,-.35)rectangle(.5,-.05);\draw[thick](L)--(T)--(R);\draw[thick](A)--(O)--(B);\draw[curve](0,0)--(0,.2)--(-.08,.3)--(.08,.45)--(-.08,.6)--(.08,.75)--(-.08,.9)--(.08,1.05)--(0,1.2)--(T);\fill(O)circle(1.4pt)node[left]{$O$};\fill(T)circle(1.4pt)node[above]{$T$};\fill(L)circle(2pt)node[below]{$M$};\fill(R)circle(2pt)node[below]{$M$};\node[left]at(-1.3,.3){$3l,3m$};\node[right]at(.45,.2){$l,m$};\draw[guide](R)--++(0,1.1);\draw($(R)+(0,.5)$)arc(90:135:.5);\node at(1.86,-.07){$\theta$};

\end{tikzpicture}
\caption{$O$为固定铰点。长杆与竖直线的夹角为$\theta$；绞架还可绕$O$整体倾斜。}\label{mp5:fig:hw:l5-34-q}
\end{figure}

\end{minipage}\par\medskip
\begin{solution}
先沿对称形变计算。$T$在支点上方$2l\cos\theta$，每根长杆质心位于支点上方$l\cos\theta/2$；短杆质心高度同样为$l\cos\theta/2$。两载重端在支点下方$l\cos\theta$。于是
\[
 U_g=(4m-2M)gl\cos\theta,
\]
\[
 U=2(2m-M)gl\cos\theta
        +\frac k2(2l\cos\theta-l/\sqrt2)^2.
\]
平衡条件为
\[
 U'=2(M-2m)gl\sin\theta
        -kl^2\sin\theta(4\cos\theta-\sqrt2)=0.
\]
在$\theta=\pi/4$处，
\[
 \boxed{k=\frac{\sqrt2(M-2m)g}{l}}.
\]
绞架还允许独立的整体倾斜。令整体绕$O$偏离竖直对称位置的角为$\phi$，其重力势能成为
\[
 U_g(\theta,\phi)=2(2m-M)gl\cos\theta\cos\phi,
\]
而弹簧长度只取决于 $\theta$。图\ref{mp5:fig:hw:l5-34-s}中的两个方向须同时考察。对两坐标分别求二阶偏导，
\[
\begin{aligned}
U_{\theta\theta}={}&2(M-2m)gl\cos\theta\cos\phi\\
&-kl^2\cos\theta(4\cos\theta-\sqrt2)+4kl^2\sin^2\theta,\\
U_{\phi\phi}&=2(M-2m)gl\cos\theta\cos\phi,\\
U_{\theta\phi}&=2(2m-M)gl\sin\theta\sin\phi.
\end{aligned}
\]
代入上式的 $k$，在 $(\theta,\phi)=(\pi/4,0)$，
\[
 U_{\theta\theta}=2\sqrt2(M-2m)gl,\quad
 U_{\phi\phi}=\sqrt2(M-2m)gl,\quad U_{\theta\phi}=0.
\]
故严格势能极小条件为
\[
 \boxed{M>2m}.
\]
此时$k>0$，且体系质心位于支点下方。$M=2m$时需$k=0$，势能不再有严格极小；$M<2m$则所需$k<0$，普通被动弹簧不能实现题设平衡。沿$\theta$的二阶导数只判断形变稳定性，整体倾斜须另行检验。


\begin{figure}[H]\centering
\begin{tikzpicture}[>=Stealth,line cap=round,line join=round]

\draw[thick](-1,1)--(0,2)--(1,1);\draw[thick](-.5,1.5)--(0,1)--(.5,1.5);\fill(0,1)circle(1.4pt)node[below]{$O$};\draw[axis](-.9,.5)arc(220:320:1.2);\node at(0,-.1){形状变化$\delta\theta$};
\begin{scope}[shift={(4.6,0)},rotate=-15]\draw[thick](-1,1)--(0,2)--(1,1);\draw[thick](-.5,1.5)--(0,1)--(.5,1.5);\fill(0,1)circle(1.4pt)node[below]{$O$};\end{scope}\draw[guide](4.86,.96)--(4.86,2.5);\node at(5,.1){整体倾斜$\delta\phi$};

\end{tikzpicture}
\caption{形状变化与整体倾斜是两个独立方向，稳定性须同时检验}\label{mp5:fig:hw:l5-34-s}
\end{figure}

\end{solution}

\par\medskip\noindent\begin{minipage}{\textwidth}
\begin{example}[零原长弹簧与抛物线约束]\label{mp5:hw:l5-35}
取$x$水平、$y$竖直向下。质量$m$的小球位于光滑轨道$y=y_0+ax^2$上，$a>0,y_0>0$；零原长、劲度系数$k$的弹簧连接小球和原点$O$，见图\ref{mp5:fig:hw:l5-35-q}。求平衡位置与沿轨道的稳定性；若轨道只能向上支持小球，还须检验接触条件。
\end{example}

\begin{figure}[H]\centering
\begin{tikzpicture}[>=Stealth,line cap=round,line join=round]

\draw[axis](-2.2,0)--(2.2,0)node[right]{$x$};\draw[axis](0,.3)--(0,-3.3)node[below]{$y$};\draw[curve,domain=-1.7:1.7,samples=80]plot(\x,{-1-.5*(\x)^2});\fill(1,-1.5)circle(2pt)node[right]{$m$};\draw[curve](0,0)--(.15,-.2)--(.1,-.4)--(.4,-.45)--(.35,-.7)--(.65,-.75)--(.6,-1)--(.9,-1.05)--(1,-1.5);\node[above left]at(0,0){$O$};\node[left]at(0,-1){$y_0$};\node[right]at(1.7,-2.7){$y=y_0+ax^2$};

\end{tikzpicture}
\caption{向下为正的高度坐标，以及从原点连接到小球的零原长弹簧}\label{mp5:fig:hw:l5-35-q}
\end{figure}

\end{minipage}\par\medskip
\begin{solution}
用$x$作独立坐标，重力势能为$-mgy$，弹性势能为$k(x^2+y^2)/2$。消去$y$后，令
\[
 A=k(1+2ay_0)-2amg,
\]
则势能可写成
\[
 U=\text{常量}+\frac A2x^2+\frac{ka^2}{2}x^4,
 \qquad U'=x(A+2ka^2x^2).
\]
故$x=0$始终为约束平衡点；$A<0$时另外出现
\[
 \boxed{x_\pm=\pm\sqrt{\frac{2amg-k(1+2ay_0)}{2ka^2}}}.
\]
再次求导，$U''=A+6ka^2x^2$。在原点 $U''=A$；在两侧平衡点，用 $2ka^2x_\pm^2=-A$，得到 $U''=A-3A=-2A>0$。因此 $A>0$ 时原点稳定，$A<0$ 时原点不稳定而两侧稳定，$A=0$ 时势能首项为正四次项，原点仍稳定但没有非零线性振动频率。图\ref{mp5:fig:hw:l5-35-s}画出这些势能形态。

沿轨道的势能极小还须满足实际接触条件。向上的单位法线在本坐标中为
\[
 \vct n=\frac{(2ax,-1)}{\sqrt{1+4a^2x^2}}.
\]
重力和弹簧力之和为 $\vct F_0=(-kx,mg-ky)$。平衡时法向分量须由支持力抵消，故
\[
N=-\vct F_0\cdot\vct n
=\frac{2kax^2+mg-ky}{\sqrt{1+4a^2x^2}}
=\frac{mg-ky_0+ka x^2}{\sqrt{1+4a^2x^2}}.
\]
原点须有$mg\ge ky_0$；若不满足，小球会被弹簧拉离轨道。两侧平衡点代入平衡关系后，$N=k\sqrt{1+4a^2x_\pm^2}/(2a)>0$，接触可成立。

沿轨道运动的动能为$m(1+4a^2x^2)\dot x^2/2$。在严格稳定点$x_e$附近，小振动频率因此为
\[
 \omega^2=\frac{U''(x_e)}{m(1+4a^2x_e^2)}.
\]
分母中的几何因子来自弧长与横坐标位移的关系。


\begin{figure}[H]\centering
\begin{tikzpicture}[>=Stealth,line cap=round,line join=round]

\draw[axis](-2.1,0)--(2.1,0)node[right]{$x$};\draw[axis](0,-.7)--(0,2.4)node[above]{$U$};\draw[curve,domain=-1.4:1.4,samples=90]plot(\x,{.4*(\x)^2+.3*(\x)^4});\draw[curve,dashed,domain=-1.4:1.4,samples=90]plot(\x,{.3*(\x)^4});\draw[curve,dotted,domain=-1.7:1.7,samples=90]plot(\x,{-.65*(\x)^2+.3*(\x)^4});\node[right]at(2.4,1.6){实线：$A>0$};\node[right]at(2.4,.9){虚线：$A=0$};\node[right]at(2.4,.2){点线：$A<0$};

\end{tikzpicture}
\caption{二次项变号产生两个新的稳定点；临界四次势阱仍有稳定性}\label{mp5:fig:hw:l5-35-s}
\end{figure}

\end{solution}

\section{电磁系统中的虚功与外源功}
固定电荷、固定电势和固定电流对应不同的外源功。电磁力或力矩由构形变化的完整能量收支求出，储能变化与外源供能须一并计入。

\par\medskip\noindent\begin{minipage}{\textwidth}
\begin{example}[平行板的电压力与电流片的磁压力]\label{mp5:hw:l5-36}
(1)面积$S$、间距$d$的真空平行板电容器两板带电$\pm Q$，求相互吸引力，忽略边缘场。(2)两无限大导体片相距$d$，流过等大反向的面电流，面电流密度大小为$K$，求单位面积上维持平衡所需的外力。图\ref{mp5:fig:hw:l5-36-q}标出电荷、电流与场的方向。可用能量密度$\varepsilon_0E^2/2$、$B^2/(2\mu_0)$。
\end{example}

\begin{figure}[H]\centering
\begin{tikzpicture}[>=Stealth,line cap=round,line join=round]

\draw[thick](0,-1)--(0,1.4);\draw[thick](1.8,-1)--(1.8,1.4);\node[above]at(0,1.4){$+Q$};\node[above]at(1.8,1.4){$-Q$};\draw[axis](.4,.4)--(1.4,.4)node[midway,above]{$E$};\draw[axis](0,-.4)--(.65,-.4);\draw[axis](1.8,-.4)--(1.15,-.4);\node at(.9,-1.5){电场：两板相吸};
\begin{scope}[shift={(5,0)}]\draw[thick](0,-1)--(0,1.4);\draw[thick](1.8,-1)--(1.8,1.4);\node at(0,.6){$\odot$};\node at(1.8,.6){$\otimes$};\node[above]at(0,1.4){$K$};\node[above]at(1.8,1.4){$-K$};\draw[axis](.9,-.5)--(.9,.9)node[right]{$B$};\draw[axis](0,-.4)--(-.65,-.4);\draw[axis](1.8,-.4)--(2.45,-.4);\node at(.9,-1.5){磁场：两片相斥};\end{scope}

\end{tikzpicture}
\caption{电场指向负板；反向面电流之间的磁场相加，磁力使两片分离}\label{mp5:fig:hw:l5-36-q}
\end{figure}

\end{minipage}\par\medskip
\begin{solution}
电容器电场为$E=Q/(\varepsilon_0S)$，固定电荷时场能
\[
 U_E=\frac{\varepsilon_0E^2}{2}Sd
     =\frac{Q^2d}{2\varepsilon_0S}.
\]
让一板沿增大 $d$ 方向虚移 $\delta d$，固定电荷时 $\delta U_E=[Q^2/(2\varepsilon_0S)]\delta d$。外力做功 $F_{\rm ext}\delta d=\delta U_E$，故外力须向外拉，平衡的电吸引力大小为
\[
 \boxed{F_E=\frac{Q^2}{2\varepsilon_0S}},\qquad
 p_E=\frac{F_E}{S}=\frac{\varepsilon_0E^2}{2}.
\]
直接计算本板受力时，应使用另一板的场$E_{\rm other}=E/2$。本板自身的场不产生整板所受的外力。

磁场中，两电流片外侧的场抵消，片间$B=\mu_0K$。一片在另一片磁场中的单位面积受力大小为
\[
 \boxed{p_B=K\frac{\mu_0K}{2}=\frac{\mu_0K^2}{2}
                       =\frac{B^2}{2\mu_0}}.
\]
图\ref{mp5:fig:hw:l5-36-s}中，维持平衡的外力应分别指向两片之间，抵消排斥。若在恒定电流条件下用场能计算，片间场能$U_B=\mu_0K^2Sd/2$随$d$增加而增加，磁力却也做正功。可把电流片与远处的回线合成供电回路，以$\mathcal L$表示随$d$改变的电感，$I$为保持不变的回路电流。磁链$\Psi=\mathcal L I$改变时，电源克服感应电动势所做的功为
\[
\delta W_{\rm source}=I\,\dd\Psi=I^2\dd\mathcal L,
\qquad \delta U_B=\frac12I^2\dd\mathcal L.
\]
因此$\delta W_{\rm source}=2\delta U_B$，其中一半成为新增场能，另一半成为机械功。固定电流时不可遗漏外源，而直接使用$-\partial U_B/\partial d$判断磁力方向。


\begin{figure}[H]\centering
\begin{tikzpicture}[>=Stealth,line cap=round,line join=round]

\draw[thick](0,-.7)--(0,1.6);\draw[thick](2,-.7)--(2,1.6);\draw[axis](2,.3)--(3,.3)node[right]{$\delta d$};\draw[guide](2.5,-.7)--(2.5,1.6);\draw[axis](0,1)--(.7,1)node[above left]{$F_{\rm ext}$};\draw[axis](2,1)--(1.3,1)node[above right]{$F_{\rm ext}$};\node[right,align=left]at(4.1,.5){恒电流时：\\[5pt]$\delta U_B=p_BS\,\delta d$\\$\delta W_B=p_BS\,\delta d$\\$\delta W_{\rm source}=2p_BS\,\delta d$};

\end{tikzpicture}
\caption{磁压力对应两片分离的机械功，维持恒电流的电源还要补充场能}\label{mp5:fig:hw:l5-36-s}
\end{figure}

\end{solution}

\par\medskip\noindent\begin{minipage}{\textwidth}
\begin{example}[可转动半圆金属片的电转矩]\label{mp5:hw:l5-38}
上下固定极板及中间可转金属片均为半径$R$的半圆薄片，上下间距$h\ll R$，中片恰位于中面，可绕过圆心的竖直轴转动，质量为$m$。完全重叠时定义$\theta=0$，转出后$|\theta|<\pi$，见图\ref{mp5:fig:hw:l5-38-q}。忽略边缘场，分别讨论：(1)两外片短接到同一电势$U$，中片接地；(2)两外片各带固定电荷$Q$，中片总电荷$-2Q$。求电转矩、平衡与稳定性，以及小幅摆动的周期或维持当前位置所需的力矩。
\end{example}

\begin{figure}[H]\centering
\begin{tikzpicture}[>=Stealth,line cap=round,line join=round]

\fill[gray!18](0,0)--({1.6*cos(30)},{1.6*sin(30)})arc(30:180:1.6)--cycle;\draw[thick](-1.6,0)arc(180:0:1.6)--cycle;\draw[curve,dashed]({1.6*cos(30)},{1.6*sin(30)})arc(30:210:1.6)--cycle;\draw[axis](0,0)--(2,0);\draw(.65,0)arc(0:30:.65);\node at(.9,.22){$\theta$};\fill(0,0)circle(1.3pt);\node at(0,2.1){重叠面积$S=R^2(\pi-|\theta|)/2$};\begin{scope}[shift={(5,0)}]\draw[thick](-1,1.2)--(1,1.2);\draw[thick](-1,0)--(1,0);\draw[thick](-1,-1.2)--(1,-1.2);\node[right]at(1,1.2){外片};\node[right]at(1,0){中片，接地};\node[right]at(1,-1.2){外片};\draw[<->](-1.3,0)--(-1.3,1.2)node[midway,left]{$h/2$};\draw[<->](-1.3,-1.2)--(-1.3,0)node[midway,left]{$h/2$};\end{scope}

\end{tikzpicture}
\caption{半圆片的重叠扇形与三片电容器的两个等厚间隙}\label{mp5:fig:hw:l5-38-q}
\end{figure}

\end{minipage}\par\medskip
\begin{solution}
先取$0<\theta<\pi$。每个间隙的电容为
\[
 c(\theta)=\frac{\varepsilon_0S}{h/2}
          =\frac{\varepsilon_0R^2(\pi-\theta)}h.
\]
两外片同电势时，两个间隙相当于并联，$C=2c$。恒电势的机械有效势应包含电源能量，取$\Phi_U=-CU^2/2$，从而
\[
 \Phi_U=\text{常量}+K_U|\theta|,
 \qquad K_U=\frac{\varepsilon_0R^2U^2}{h}.
\]
在 $\theta\ne0$ 处，$\dd|\theta|/\dd\theta=\operatorname{sgn}\theta$，故电转矩 $\tau_U=-\dd\Phi_U/\dd\theta$ 为
\[
 \boxed{\tau_U=-K_U\operatorname{sgn}\theta}.
\]
在 $\theta\ne0$ 处不是自由平衡，要保持该位置，外力矩须为 $+K_U\operatorname{sgn}\theta$。$\theta=0$ 是尖点势阱的稳定平衡。图\ref{mp5:fig:hw:l5-38-s}左图说明恢复转矩大小恒定，不能套用简谐公式。中片对圆心轴转动惯量 $I=mR^2/2$。从静止振幅 $A$ 释放，在 $0<\theta<A$ 内有 $\ddot\theta=-K_U/I$。以释放为零时刻，积分得 $\dot\theta=-(K_U/I)t$、$\theta=A-K_Ut^2/(2I)$，也可直接采用习题\ref{mp4:exercise:l4-h2}已展开的恒加速度积分，将 $m,f,x$ 换为 $I,K_U,\theta$。令 $\theta=0$，四分之一周期为 $\sqrt{2IA/K_U}$，所以
\[
 \boxed{T_U=4\sqrt{\frac{2IA}{K_U}}
            =4\sqrt{\frac{mhA}{\varepsilon_0U^2}}}.
\]
在题设忽略边缘场的模型中，$T_U\propto\sqrt A$，振幅趋零时周期也趋零。

固定电荷时，每个间隙外片各有$Q$，总场能为
\[
 \Phi_Q=2\frac{Q^2}{2c}
       =\frac{Q^2h}{\varepsilon_0R^2(\pi-|\theta|)}.
\]
对 $\theta$ 求导，$\dd(\pi-|\theta|)^{-1}/\dd\theta=-(\pi-|\theta|)^{-2}(-\operatorname{sgn}\theta)$。电转矩取势能导数的负号，故
\[
 \boxed{\tau_Q=-\frac{Q^2h\,\operatorname{sgn}\theta}
 {\varepsilon_0R^2(\pi-|\theta|)^2}}.
\]
它同样使完全重叠位置稳定，转出越多，所需保持力矩越大。对于振幅$A$，精确周期可用
\[
 T_Q=4\sqrt{I/2}\int_0^A
 \frac{\dd\theta}{\sqrt{\Phi_Q(A)-\Phi_Q(\theta)}}.
\]
这个积分可以继续求出。记 $\Lambda=Q^2h/(\varepsilon_0R^2)$、$b=\pi-A>0$，则
\[
\Phi_Q(A)-\Phi_Q(\theta)
=\frac{\Lambda(A-\theta)}{(\pi-A)(\pi-\theta)}.
\]
令 $u=\sqrt{A-\theta}$，$\dd\theta=-2u\dd u$，两端 $\theta=0,A$ 对应 $u=\sqrt A,0$。交换积分上下限，得到
\[
T_Q=4\sqrt{\frac{Ib}{2\Lambda}}\,2\int_0^{\sqrt A}\sqrt{b+u^2}\dd u.
\]
用分部积分所得的根式原函数（参见习题\ref{mp1:pr:p10}），
\[
2\int_0^{\sqrt A}\sqrt{b+u^2}\dd u
=\left[u\sqrt{b+u^2}+b\ln\frac{u+\sqrt{b+u^2}}{\sqrt b}\right]_0^{\sqrt A}.
\]
代入两端，精确周期为
\[
\boxed{T_Q=4\sqrt{\frac{I(\pi-A)}{2\Lambda}}\left[
\sqrt{\pi A}+(\pi-A)\ln\frac{\sqrt A+\sqrt\pi}{\sqrt{\pi-A}}\right].}
\]
当 $A\ll1$，势能在原点附近为 $\Phi_Q(0)+K_Q|\theta|+O(\theta^2)$，其中 $K_Q=\Lambda/\pi^2$，故恢复为 $T_Q\simeq4\sqrt{2IA/K_Q}$。

以上$Q$指每一外片的电荷。若改用两外片总电荷$Q_\Sigma=2Q$，转矩分子应写为$Q_\Sigma^2h/4$，两种电荷定义对应的系数相差四倍。电压条件也须明确：若$U$改指上下外片之间的电势差，且中片处于它们的中点电势，结论会改变。被中片覆盖的区域内，两个半间隙各承担$U/2$；未覆盖的区域则仍有上下外片之间的场。两部分场能相加为$\varepsilon_0S_0U^2/(2h)$，$S_0=\pi R^2/2$，与转角无关。在这一对称接法与零厚度近似下，电转矩为零。只计算覆盖区而漏掉未覆盖区，会错误地得到一个缩小后的非零转矩。


\begin{figure}[H]\centering
\begin{tikzpicture}[>=Stealth,line cap=round,line join=round]

\draw[axis](-1.8,0)--(1.8,0)node[right]{$\theta$};\draw[axis](0,-.2)--(0,2.3)node[above]{$\Phi_U$};\draw[curve](-1.6,1.6)--(0,0)--(1.6,1.6);\draw[axis](-1.1,1.1)--(-.65,.65);\draw[axis](1.1,1.1)--(.65,.65);\node at(0,-.5){恒电势：尖点势阱};\begin{scope}[shift={(5,0)}]\draw[axis](-1.8,0)--(1.8,0)node[right]{$\theta$};\draw[axis](0,-.2)--(0,2.3)node[above]{$\Phi_Q-\Phi_Q(0)$};\draw[curve,domain=-1.6:0,samples=65]plot(\x,{5*(1/(pi+\x)-1/pi)});\draw[curve,domain=0:1.6,samples=65]plot(\x,{5*(1/(pi-\x)-1/pi)});\fill(0,0)circle(1.1pt);\node at(0,-.5){固定电荷：尖点势阱};\end{scope}

\end{tikzpicture}
\caption{两种接法的势能在$\theta=0$处均有尖点。固定电荷时，曲线两侧的斜率随$|\theta|$增大。}\label{mp5:fig:hw:l5-38-s}
\end{figure}

\end{solution}

\par\medskip\noindent\begin{minipage}{\textwidth}
\begin{example}[液体介质进入竖直与水平电容器]\label{mp5:hw:l5-39}
真空平行板电容器面积$ab$、板距$d$，带电量$Q$。液体密度$\rho$、相对介电常数$\varepsilon_r>1$，外部液面保持不变，忽略表面张力与边缘场。(1)极板竖直，竖直高度$a$、宽$b$，下边缘在外液面，求液体在板间上升高度$z\in[0,a]$；(2)极板水平，面积$ab$，下板在外液面，液体在板间充入厚度$z\in[0,d]$，求$z$。再分别给出接理想恒压源$U$时的结果，见图\ref{mp5:fig:hw:l5-39-q}。
\end{example}

\begin{figure}[H]\centering
\begin{tikzpicture}[>=Stealth,line cap=round,line join=round]

\draw[thick](0,0)--(0,2.8);\draw[thick](1,0)--(1,2.8);\fill[gray!25](.03,0)rectangle(.97,1.1);\draw[guide](-.6,0)--(1.6,0);\draw[<->](1.3,0)--(1.3,1.1)node[midway,right]{$z$};\draw[<->](-.4,0)--(-.4,2.8)node[midway,left]{$a$};\node at(.5,3.2){竖直：横向并联};\node at(.5,-.5){板宽$b$，间距$d$};\begin{scope}[shift={(5,0)}]\draw[thick](-1,0)--(2.2,0);\draw[thick](-1,2)--(2.2,2);\fill[gray!25](-.95,.02)rectangle(2.15,.8);\draw[guide](-1.6,0)--(2.8,0);\draw[<->](2.5,0)--(2.5,.8)node[midway,right]{$z$};\draw[<->](-1.3,0)--(-1.3,2)node[midway,left]{$d$};\node at(.6,3.2){水平：纵向串联};\node at(.6,-.5){极板面积$ab$};\end{scope}

\end{tikzpicture}
\caption{两种放置方式对应并联与串联的介质填充；高度从外部液面量起}\label{mp5:fig:hw:l5-39-q}
\end{figure}

\end{minipage}\par\medskip
\begin{solution}
竖直极板时，上部真空和下部液体的电容并联。记$k_r=\varepsilon_r-1$，
\[
 C_v(z)=\frac{\varepsilon_0b}{d}(a+k_rz),\qquad
 U_g=\frac12\rho gbdz^2.
\]
固定电荷下，总势能为 $\rho gbdz^2/2+Q^2d/[2\varepsilon_0b(a+k_rz)]$。对 $z$ 求导，平衡条件为
\[
\rho gbdz-\frac{Q^2dk_r}{2\varepsilon_0b(a+k_rz)^2}=0.
\]
将分母移到另一边并约去公共的 $d$，得到
\[
 \boxed{z(a+k_rz)^2=\frac{Q^2k_r}{2\varepsilon_0b^2\rho g}}.
\]
左端的导数为 $(a+k_rz)^2+2k_rz(a+k_rz)=(a+k_rz)(a+3k_rz)>0$，因此有唯一非负根；只有根小于 $a$ 时才是部分填充解。总势能二阶导数为
\[
\rho gbd+\frac{Q^2dk_r^2}{\varepsilon_0b(a+k_rz)^3}>0,
\]
所以内部平衡稳定。恒电压时使用有效势能 $U_g-C_vU^2/2$，其导数为 $\rho gbdz-\varepsilon_0bk_rU^2/(2d)$，令其为零，得到
\[
 \boxed{z=\frac{\varepsilon_0k_rU^2}{2\rho gd^2}}.
\]
超过$a$时，理想模型的填充高度到顶，须按完全填充边界处理。

水平极板时，两种介质沿电场方向串联。令$\gamma=1-1/\varepsilon_r$、$S=ab$，
\[
 C_h(z)=\frac{\varepsilon_0S}{d-\gamma z},\qquad
 U_g=\frac12\rho gSz^2.
\]
固定电荷的场能 $Q^2(d-\gamma z)/(2\varepsilon_0S)$ 对 $z$ 呈线性，总势能的导数为 $\rho gSz-Q^2\gamma/(2\varepsilon_0S)$。令其为零，得到
\[
 \boxed{z=\frac{Q^2\gamma}{2\varepsilon_0\rho gS^2}}.
\]
二阶导数为 $\rho gS>0$，该部分填充解稳定，且在固定 $Q,S$ 下与 $d$ 无关，但必须满足 $z<d$；若算出的高度达到或超过 $d$，则在本题高度范围内取完全填充边界 $z=d$。

恒电压时，有效势能及其导数为
\[
\Phi=\frac12\rho gSz^2-\frac{\varepsilon_0SU^2}{2(d-\gamma z)},
\qquad
\Phi'=\rho gSz-\frac{\varepsilon_0S\gamma U^2}{2(d-\gamma z)^2}.
\]
令 $\Phi'=0$，得到
\[
 \boxed{z(d-\gamma z)^2=\frac{\varepsilon_0\gamma U^2}{2\rho g}}.
\]
该平衡方程可能有多个正根，还须逐根判断稳定性。先在固定电压下求二阶导数，再用平衡关系替换电压因子：
\[
\Phi''=\rho gS-\frac{\varepsilon_0S\gamma^2U^2}{(d-\gamma z)^3}
=\rho gS\left(1-\frac{2\gamma z}{d-\gamma z}\right)
=\rho gS\frac{d-3\gamma z}{d-\gamma z}.
\]
因此内部根仅在$z<d/(3\gamma)$时稳定；大于此值的根不稳定。令
\[
B=\frac{\varepsilon_0\gamma U^2}{2\rho g},\qquad
f(z)=z(d-\gamma z)^2,
\]
则 $f'=(d-\gamma z)^2-2\gamma z(d-\gamma z)=(d-\gamma z)(d-3\gamma z)$，而整个物理区间内 $d-\gamma z>0$。若 $\gamma\le1/3$，$f$ 在 $[0,d]$ 上递增；$B<f(d)$ 时有一个稳定的内部根，$B\ge f(d)$ 时达到完全填充边界。

若$\gamma>1/3$，最大值出现在$z_c=d/(3\gamma)<d$，且
\[
f(z_c)=\frac{4d^3}{27\gamma}.
\]
当$0<B<f(d)$时只有一个内部根，位于稳定支；当$f(d)<B<f(z_c)$时有两个内部根，较小者稳定，较大者不稳定。完全填充边界此时也可稳定，因为
$\Phi'(d)=\rho gS[f(d)-B]/[d(1-\gamma)]^2<0$，从边界向内退缩会增加势能。因此仅凭给定电压，不能在两个稳定填充状态之间作唯一选择，还需说明初始状态或加压过程。

在$B=f(z_c)$的内部临界点，虽然$\Phi'=\Phi''=0$，但
\[
\Phi'''=-\varepsilon_0S\gamma^2U^2(-3)(d-\gamma z)^{-4}(-\gamma)
=-\frac{3\varepsilon_0S\gamma^3U^2}{(d-\gamma z)^4}<0.
\]
该点是具有非零三次项的驻点，不是势能极小值；稳定支与不稳定支在此合并。继续增大电压便没有部分填充平衡，液体升至$z=d$。边界$B=f(d)$须按单侧位移判断：$\gamma>1/3$时边界不是稳定极小，而$\gamma\le1/3$时它是稳定支到达顶端的终点。

图\ref{mp5:fig:hw:l5-39-s}也可从界面电压力理解四个结果。竖直极板处电场平行液面，向上电压力为$\varepsilon_0(\varepsilon_r-1)E^2/2$；水平极板处位移矢量法向连续，电压力为$D^2(1-1/\varepsilon_r)/(2\varepsilon_0)$。它们分别与静水压$\rho gz$平衡。固定电荷与固定电压的场量随$z$变化方式不同，这正是两类方程的差别。


\begin{figure}[H]\centering
\begin{tikzpicture}[>=Stealth,line cap=round,line join=round]

\draw[curve](-1.2,0)--(1.2,0);\fill[gray!15](-1.2,-.8)rectangle(1.2,0);\draw[axis](-.8,.35)--(.8,.35)node[right]{$E$};\draw[axis](0,0)--(0,.9)node[above]{$p_{Ev}$};\draw[axis](.5,0)--(.5,-.7)node[below]{$\rho gz$};\node at(0,-1.3){$p_{Ev}=\varepsilon_0(\varepsilon_r-1)E^2/2$};\begin{scope}[shift={(5.6,0)}]\draw[curve](-1.2,0)--(1.2,0);\fill[gray!15](-1.2,-.8)rectangle(1.2,0);\draw[axis](-.7,-.6)--(-.7,.7)node[above]{$D$};\draw[axis](0,0)--(0,.9)node[above]{$p_{Eh}$};\draw[axis](.5,0)--(.5,-.7)node[below]{$\rho gz$};\node at(0,-1.3){$p_{Eh}=D^2(1-1/\varepsilon_r)/(2\varepsilon_0)$};\end{scope}

\end{tikzpicture}
\caption{界面电压力与静水压平衡；电场切向连续与电位移法向连续对应不同表达式}\label{mp5:fig:hw:l5-39-s}
\end{figure}

\end{solution}

\section{习题与解答}

\par\medskip\noindent\begin{minipage}{\textwidth}
\begin{exercise}[一个动滑轮的等效惯性]\label{mp5:exercise:l5-h1}
一根轻绳按图\ref{mp5:fig:l5-h1-q}绕过动滑轮 C 和定滑轮 O，绳一端固定，另一端挂质量 $m$。动滑轮轴上挂质量 $M$，滑轮无质量。以 C 的向下位移 $q$ 为坐标，求加速度和绳张力，说明为什么等效惯性中会出现 $4m$。
\end{exercise}

\begin{figure}[H]\centering
\begin{tikzpicture}
\draw[curve](-2.1,3.5)--(1.4,3.5);
\draw[fill=white,curve](-1.15,1)circle(0.35);\fill(-1.15,1)circle(1pt);
\node[above right]at(-1.1099999999999999,1.04){$C$};
\draw[fill=white,curve](-0.45,2.8)circle(0.35);\fill(-0.45,2.8)circle(1pt);
\node[above right]at(-0.41000000000000003,2.84){$O$};
\draw(-.45,2.8)--(-.45,3.5);\draw[curve](-1.5,3.5)--(-1.5,1)arc(180:360:.35)--(-.8,2.8)arc(180:0:.35)--(-.1,.275);\draw[curve](-1.15,1)--(-1.15,-.275);\draw[fill=white](-1.45,-0.8250000000000001)rectangle(-0.8499999999999999,-0.275);\node at(-1.15,-0.55){$M$};
\draw[fill=white](-0.4,-0.275)rectangle(0.19999999999999998,0.275);\node at(-0.1,0){$m$};
\draw[axis](.7,2)--(.7,-.6)node[below]{$q$ 正向};
\end{tikzpicture}
\caption{一根绳的两个变长分支共同连接动滑轮，绳的自由端再经定滑轮转向。}\label{mp5:fig:l5-h1-q}
\end{figure}
\end{minipage}\par\medskip
\begin{solution}
C 向下移动 $q$，与它相连的两段绳总长增加 $2q$，故自由端重物 $m$ 向上移动 $2q$。图\ref{mp5:fig:l5-h1-s}中的速度关系是
\[
v_C=\dot q,\qquad v_m=-2\dot q.
\]
绳长约束为 $2y_C+y_m=\text{常量}$。再对上述速度关系求导，得到 $a_m=-2\ddot q$。位移比、速度比和加速度比都由这条线性约束确定。

因此
\[
T=\frac12(M+4m)\dot q^2,\qquad U=-(M-2m)gq+\text{常量}.
\]

\begin{figure}[H]\centering
\begin{tikzpicture}
\draw[curve](-2.1,3.5)--(1.4,3.5);
\draw[fill=white,curve](-1.15,1)circle(0.35);\fill(-1.15,1)circle(1pt);
\node[above right]at(-1.1099999999999999,1.04){$C$};
\draw[fill=white,curve](-0.45,2.8)circle(0.35);\fill(-0.45,2.8)circle(1pt);
\node[above right]at(-0.41000000000000003,2.84){$O$};
\draw(-.45,2.8)--(-.45,3.5);\draw[curve](-1.5,3.5)--(-1.5,1)arc(180:360:.35)--(-.8,2.8)arc(180:0:.35)--(-.1,.275);\draw[curve](-1.15,1)--(-1.15,-.275);\draw[fill=white](-1.45,-0.8250000000000001)rectangle(-0.8499999999999999,-0.275);\node at(-1.15,-0.55){$M$};
\draw[fill=white](-0.4,-0.275)rectangle(0.19999999999999998,0.275);\node at(-0.1,0){$m$};
\draw[axis](-1.8,-.5)--(-1.8,-1.1)node[left]{$\dot q$};\draw[axis](.45,0)--(.45,1.2)node[right]{$2\dot q$};
\end{tikzpicture}
\caption{取 $\dot q>0$。动滑轮向下，自由端重物向上，速度箭头按 $1:2$ 绘制。}\label{mp5:fig:l5-h1-s}
\end{figure}

分别求广义惯性项与广义重力，
\[
\frac{\dd}{\dd t}\frac{\partial T}{\partial\dot q}=(M+4m)\ddot q,
\qquad Q_q=-U_q=(M-2m)g.
\]
两者相等，得到
\[
\boxed{\ddot q=\frac{M-2m}{M+4m}g.}
\]
对自由端 $m$，向下加速度为 $-2\ddot q$，所以
\[
mg-\mathcal T=-2m\ddot q,
\qquad\boxed{\mathcal T=\frac{3mM}{M+4m}g.}
\]
将动滑轮与重物 $M$ 作为整体，$Mg-2\mathcal T=M\ddot q$，所得结果相同。$M=2m$ 时各物体加速度为零；$M>2m$ 时 $C$ 的加速度向下，$M<2m$ 时向上。自由端的速率是 $|\dot q|$ 的两倍，因此它对等效惯性的贡献为 $4m$。
分子中的$2m$来自重力虚功的位移比，与分母中速度比平方产生的$4m$不同。$M,m>0$时$\mathcal T>0$，满足绳保持绷紧的条件。$M=2m$时的零加速度并不要求速度为零；若释放前已有满足绳长关系的速度，理想模型中可继续匀速运动。

\end{solution}

\par\medskip\noindent\begin{minipage}{\textwidth}
\begin{exercise}[双摆的两个小振动模态（选做）]\label{mp5:exercise:l5-h2}
两个质量均为 $m$ 的小球由两根长为 $l$ 的无质量刚杆串接成双摆，悬点固定，见图\ref{mp5:fig:l5-h2-q}。两根杆与竖直向下方向的有向夹角分别为 $\theta_1,\theta_2$。求完整动能，再在小角度下求两个固有频率和相应振幅比。
\end{exercise}

\begin{figure}[H]\centering
\begin{tikzpicture}
\fill(0,0)circle(1.3pt)node[above]{$O$};
\draw[curve](0,0)--(0.49442719099991583,-1.5216904260722457)--(0.16176848569150087,-3.0867265872463348);\fill(0.49442719099991583,-1.5216904260722457)circle(2.2pt);\fill(0.16176848569150087,-3.0867265872463348)circle(2.2pt);\draw[guide](0,0)--(0,-1.8);\draw[guide](0.49442719099991583,-1.5216904260722457)--(0.49442719099991583,-3.2216904260722456);\node[right]at(0.5944271909999158,-1.5216904260722457){$m$};\node[right]at(0.2617684856915009,-3.0867265872463348){$m$};\node[left]at(.05,-.7){$l$};\node[right]at(0.7080978383457084,-2.30420850665929){$l$};\draw[axis](0,-.8)arc(-90:-72:.8);\node[right]at(.32,-.58){$\theta_1$};\draw[axis](0.49442719099991583,-2.3216904260722457)arc(-90:-102:.8);\node[left]at(0.22442719099991582,-2.2216904260722456){$\theta_2$};
\end{tikzpicture}
\caption{双摆的两个独立角坐标。下杆角度相对固定竖直方向定义，第二根竖直虚线穿过上球。}\label{mp5:fig:l5-h2-q}
\end{figure}
\end{minipage}\par\medskip
\begin{solution}
令两球的位置分别为
\[
\vct r_1=(l\sin\theta_1,-l\cos\theta_1),\qquad
\vct r_2=\vct r_1+(l\sin\theta_2,-l\cos\theta_2).
\]
逐项求导，
\[
\vct v_1=l\dot\theta_1(\cos\theta_1,\sin\theta_1),\qquad
\vct v_2=\vct v_1+l\dot\theta_2(\cos\theta_2,\sin\theta_2).
\]
下球速度平方为 $l^2[\dot\theta_1^2+\dot\theta_2^2+2\dot\theta_1\dot\theta_2\cos(\theta_1-\theta_2)]$。再加上上球动能，得到
\[
T=\frac{ml^2}{2}\left[2\dot\theta_1^2+\dot\theta_2^2
+2\dot\theta_1\dot\theta_2\cos(\theta_1-\theta_2)\right],
\]
\[
U=mgl(3-2\cos\theta_1-\cos\theta_2).
\]
下球速度是上球速度与下杆转动速度之和。两矢量的夹角为 $\theta_1-\theta_2$，所以交叉项必须保留。作小振动近似时，$\dot\theta_1\dot\theta_2$ 已为二阶小量，而 $\cos(\theta_1-\theta_2)-1$ 为二阶小量；二者相乘是四阶项，在二次动能中舍去。势能则由 $1-\cos\theta_j\simeq\theta_j^2/2$ 展开。

图\ref{mp5:fig:l5-h2-s}给出两个小振动模态。将动能与势能保留到二阶，
\[
T_2=\frac{ml^2}{2}(2\dot\theta_1^2+2\dot\theta_1\dot\theta_2+\dot\theta_2^2),
\quad U_2=\frac{mgl}{2}(2\theta_1^2+\theta_2^2).
\]

\begin{figure}[H]\centering
\begin{tikzpicture}
\begin{scope}[shift={(0,0)},scale=1]
\fill(0,0)circle(1.3pt)node[above]{$O$};
\draw[curve](0,0)--(0.33265870530841496,-1.5650361611740893)--(0.7996672473949139,-3.0953639109121127);\fill(0.33265870530841496,-1.5650361611740893)circle(2.2pt);\fill(0.7996672473949139,-3.0953639109121127)circle(2.2pt);\draw[guide](0,0)--(0,-1.8);\draw[guide](0.33265870530841496,-1.5650361611740893)--(0.33265870530841496,-3.2650361611740895);\node[right]at(0.432658705308415,-1.5650361611740893){$m$};\node[right]at(0.8996672473949139,-3.0953639109121127){$m$};\node[left]at(.05,-.7){$l$};\node[right]at(0.9461629763516645,-2.330200036043101){$l$};\node at(.3,-3.7){$\theta_2=+\sqrt2\,\theta_1$};\end{scope}
\begin{scope}[shift={(4.8,0)},scale=1]
\fill(0,0)circle(1.3pt)node[above]{$O$};
\draw[curve](0,0)--(0.33265870530841496,-1.5650361611740893)--(-0.134349836778084,-3.0953639109121127);\fill(0.33265870530841496,-1.5650361611740893)circle(2.2pt);\fill(-0.134349836778084,-3.0953639109121127)circle(2.2pt);\draw[guide](0,0)--(0,-1.8);\draw[guide](0.33265870530841496,-1.5650361611740893)--(0.33265870530841496,-3.2650361611740895);\node[right]at(0.432658705308415,-1.5650361611740893){$m$};\node[right]at(-0.03434983677808398,-3.0953639109121127){$m$};\node[left]at(.05,-.7){$l$};\node[right]at(0.47915443426516546,-2.330200036043101){$l$};\node at(.3,-3.7){$\theta_2=-\sqrt2\,\theta_1$};\end{scope}

\end{tikzpicture}
\caption{双摆的小振动构形。左图为低频同向模态，右图为高频反向模态；角度按相应振幅比绘制。}\label{mp5:fig:l5-h2-s}
\end{figure}

二次动能对应的质量矩阵不是对角矩阵。分别求导，
\[
\begin{aligned}
\frac{\partial T_2}{\partial\dot\theta_1}&=ml^2(2\dot\theta_1+\dot\theta_2),&
\frac{\partial U_2}{\partial\theta_1}&=2mgl\theta_1,\\
\frac{\partial T_2}{\partial\dot\theta_2}&=ml^2(\dot\theta_1+\dot\theta_2),&
\frac{\partial U_2}{\partial\theta_2}&=mgl\theta_2.
\end{aligned}
\]
代入拉格朗日方程并除以 $ml^2$，小振动方程为
\[
2\ddot\theta_1+\ddot\theta_2+\frac{2g}{l}\theta_1=0,
\qquad
\ddot\theta_1+\ddot\theta_2+\frac gl\theta_2=0.
\]
因此两条方程都含有两个角加速度，求固有频率时必须联立这两个方程。

令 $\theta_j=A_j\cos\omega t$、$\lambda=\omega^2l/g$，方程组为
\[
\begin{pmatrix}2-2\lambda&-\lambda\\-\lambda&1-\lambda\end{pmatrix}
\begin{pmatrix}A_1\\A_2\end{pmatrix}=0.
\]
非零振幅要求行列式为零，即 $(2-2\lambda)(1-\lambda)-\lambda^2=\lambda^2-4\lambda+2=0$。两根为 $\lambda=[4\pm\sqrt{16-8}]/2=2\pm\sqrt2$，所以
\[
\boxed{\omega_\pm^2=(2\pm\sqrt2)\frac gl.}
\]
由矩阵第二行，$A_2/A_1=\lambda/(1-\lambda)$；代入两个根，分别得到 $\sqrt2$ 与 $-\sqrt2$，与图\ref{mp5:fig:l5-h2-s}中同向、反向的角度约定一致。一般小运动是两个模态的叠加：
\[
\begin{aligned}
\theta_1&=C_-\cos(\omega_-t+\delta_-)+C_+\cos(\omega_+t+\delta_+),\\
\theta_2&=\sqrt2C_-\cos(\omega_-t+\delta_-)-\sqrt2C_+\cos(\omega_+t+\delta_+).
\end{aligned}
\]
只有初始角度和角速度都满足同一模态的振幅比，运动中才会保持该构形比例。

\end{solution}

\par\medskip\noindent\begin{minipage}{\textwidth}
\begin{exercise}[圆环上的弹簧珠]\label{mp5:exercise:l5-h3}
小珠质量为 $m$，在水平面内沿半径为 $R$ 的光滑圆环运动。圆环上的固定点 A 与小珠 P 之间接劲度系数 $k$、自然长度 $b$ 的轻弹簧，$0<b<2R$。从 $OA$ 逆时针转向 $OP$ 的有向角为 $\theta\in(0,2\pi)$，见图\ref{mp5:fig:l5-h3-q}。求区间内的平衡位置、稳定性及稳定位置附近的小振动频率。
\end{exercise}

\begin{figure}[H]\centering
\begin{tikzpicture}
\draw[curve](0,0)circle(1.7);\draw[guide](0,0)--(1.7,0);\draw[curve](0,0)--(0.43999237667428526,1.6420739046914161);\draw[curve](1.7,0)--(0.43999237667428526,1.6420739046914161);\fill(0,0)circle(1.3pt)node[below left]{$O$};
\fill(1.7,0)circle(1.3pt)node[below right]{$A$};
\fill(0.43999237667428526,1.6420739046914161)circle(1.3pt)node[above]{$P$};
\draw[axis](.6,0)arc(0:75:.6);\node at(.65,.45){$\theta$};\node[above right]at(.2,.9){$R$};\node[right]at(1.0699961883371425,0.8210369523457081){$\ell$};
\end{tikzpicture}
\caption{水平圆环的俯视图。弹簧长度 $\ell$ 是弦 $AP$ 的长度，$\theta$ 是从 $OA$ 逆时针转向 $OP$ 的有向角。}\label{mp5:fig:l5-h3-q}
\end{figure}
\end{minipage}\par\medskip
\begin{solution}
在给定的角度分支上，$\ell=2R\sin(\theta/2)>0$。动能和势能为
\[
T=\frac12mR^2\dot\theta^2,
\qquad U=\frac k2[2R\sin(\theta/2)-b]^2.
\]
图\ref{mp5:fig:l5-h3-q}的三角形 OAP 给出弦长关系
\[
\ell^2=2R^2(1-\cos\theta),\quad
\ell'=R\cos(\theta/2),\quad
\ell''=-\frac R2\sin(\theta/2).
\]
故弹簧的切向广义力为 $Q_\theta=-k(\ell-b)\ell'$。平衡既可能来自弹簧力为零，也可能来自弹簧力沿圆环法线、切向分量为零。后一种情形对应 $\theta=\pi$，故求解时须保留 $\cos(\theta/2)=0$ 的分支。

由
\[
U'=k(\ell-b)R\cos(\theta/2)
\]
可得两类平衡位置。图\ref{mp5:fig:l5-h3-s}中的 P$_\pm$ 满足 $\ell=b$：
\[
\theta_e=2\arcsin\frac b{2R},\qquad 2\pi-\theta_e.
\]
另有直径另一端 B，$\theta=\pi$。

\begin{figure}[H]\centering
\begin{tikzpicture}
\draw[curve](0,0)circle(1.7);\fill(1.7,0)circle(1.4pt)node[right]{$A$};\fill(-1.7,0)circle(1.4pt)node[left]{$B$};\fill(.85,1.47224)circle(1.4pt)node[above]{$P_+$};\fill(.85,-1.47224)circle(1.4pt)node[below]{$P_-$};\draw[guide](1.7,0)--(.85,1.47224);\draw[guide](1.7,0)--(.85,-1.47224);\fill(0,0)circle(1pt)node[left]{$O$};\draw[guide](0,0)--(.85,1.47224);\draw[axis](.55,0)arc(0:60:.55);\node at(.7,.45){$\theta_e$};
\end{tikzpicture}
\caption{取 $b=R$。$P_+,P_-$ 是弹簧处于自然长度时的稳定位置，$B$ 不稳定；两条弦分别表示两种构形。}\label{mp5:fig:l5-h3-s}
\end{figure}

完整的二阶导数为
\[
U''=k\big[(\ell')^2+(\ell-b)\ell''\big].
\]
在两个自然长度位置，第二项为零。令$\xi=\theta-\theta_e$，则
\[
U''=kR^2\cos^2(\theta_e/2)>0,
\qquad mR^2\ddot\xi+kR^2\left(1-\frac{b^2}{4R^2}\right)\xi=0.
\]
故小振动频率为
\[
\boxed{\omega=\sqrt{\frac km}\sqrt{1-\frac{b^2}{4R^2}}.}
\]
它们都是稳定平衡点。在 B，
\[
U''(\pi)=-\frac{kR}{2}(2R-b)<0,
\]
故不稳定。在 $B$ 处，弹簧力沿直径，由圆环反力平衡，因而切向广义力为零。端点 $\theta=0,2\pi$ 对应 P 与固定点 A 重合，弹簧方向在该处退化，本题已经排除该构形。
在B，$U''$的第一项为零，第二项为负。此时弹簧仍被拉伸；小珠向任一侧偏离B都会缩短弹簧、降低势能。$b\to2R^-$时稳定位置靠近B，频率趋零；给定的$0<b<2R$保证两个稳定位置仍彼此分开。

\end{solution}

\par\medskip\noindent\begin{minipage}{\textwidth}
\begin{exercise}[竖直运动悬点的单摆]\label{mp5:exercise:l5-h4}
无质量刚杆长 $l$，末端小球质量 $m$。悬点 O 的竖直坐标为给定函数 $Y(t)$，竖直向上为正；摆角从竖直向下起算，见图\ref{mp5:fig:l5-h4-q}。求 $\theta$ 的运动方程。若 $\ddot Y=A$ 为常量，写出相对悬点运动的能量积分并判断最低位置的稳定性。
\end{exercise}

\begin{figure}[H]\centering
\begin{tikzpicture}
\draw[axis](-1.2,-1.5)--(-1.2,2)node[above]{$Y$};\fill(0,1)circle(2pt)node[above left]{$O$};\draw[guide](0,1)--(0,-1.5);\draw[curve](0,1)--(1.2,-1.078)node[midway,right]{$l$};\fill(1.2,-1.078)circle(2.5pt)node[right]{$m$};\draw[axis](.6,1)--(.6,1.9)node[right]{$\ddot Y$};\draw[axis](0,.4)arc(-90:-60:.6);\node at(.37,.22){$\theta$};
\end{tikzpicture}
\caption{悬点竖直运动，$Y,\dot Y,\ddot Y$ 均取向上为正，$\theta$ 从竖直向下方向量起。}\label{mp5:fig:l5-h4-q}
\end{figure}
\end{minipage}\par\medskip
\begin{solution}
图\ref{mp5:fig:l5-h4-q}中，摆球位置为 $(l\sin\theta,Y-l\cos\theta)$，求导得 $\vct v=(l\cos\theta\dot\theta,\dot Y+l\sin\theta\dot\theta)$。将两分量平方相加，得到
\[
 T=\frac m2[l^2\dot\theta^2+\dot Y^2
 +2l\dot Y\sin\theta\dot\theta],\qquad
 U=mg(Y-l\cos\theta).
\]
$Y(t)$是外界给定的运动，只有$\theta$是待求坐标。逐项计算：
\[
 \begin{aligned}
 T_{\dot\theta}&=ml^2\dot\theta+ml\dot Y\sin\theta,\\
 \frac{\dd}{\dd t}T_{\dot\theta}
 &=ml^2\ddot\theta+ml\ddot Y\sin\theta
 +ml\dot Y\cos\theta\dot\theta,\\
 T_\theta&=ml\dot Y\cos\theta\dot\theta,\qquad
 U_\theta=mgl\sin\theta.
 \end{aligned}
\]
代入$\dd T_{\dot\theta}/\dd t-T_\theta+U_\theta=0$，
含$\dot Y\dot\theta$的两项相消，得到
\[
 \boxed{\ddot\theta+\frac{g+\ddot Y}{l}\sin\theta=0}.
\]
\begin{figure}[H]\centering
\begin{tikzpicture}
\begin{scope}[shift={(0,0)},scale=1]
\fill(0,1)circle(1.5pt)node[above]{$O$};\draw[curve](0,1)--(1,-.732);\fill(1,-.732)circle(2pt)node[right]{$P$};\draw[axis](1,-.732)--(1,-1.65)node[right]{$mg$};\node at(.2,-2.1){重力};\end{scope}
\begin{scope}[shift={(4.8,0)},scale=1]
\fill(0,1)circle(1.5pt)node[above]{$O$};\draw[curve](0,1)--(1,-.732);\fill(1,-.732)circle(2pt)node[right]{$P$};\draw[axis](1,-.732)--(1,-1.65)node[right]{$mA$};\node at(.2,-2.1){惯性力};\end{scope}

\end{tikzpicture}
\caption{取 $A>0$，随悬点平动的参考系中，重力与惯性力均竖直向下。两种力分别画在相同构形中，箭头只示方向。}\label{mp5:fig:l5-h4-s}
\end{figure}
图\ref{mp5:fig:l5-h4-s}中悬点向上加速时，随悬点平动的参考系中惯性力向下。若$\ddot Y=A$为常量，将摆角方程乘以$ml^2\dot\theta$，有
\[
 \frac{\dd}{\dd t}\left[
 \frac12ml^2\dot\theta^2+m(g+A)l(1-\cos\theta)\right]=0.
\]
因而相对运动的积分为
\[
 \boxed{E_{\rm rel}=
 \frac12ml^2\dot\theta^2+m(g+A)l(1-\cos\theta)}
 =\frac12ml^2\dot\theta_0^2+m(g+A)l(1-\cos\theta_0).
\]
若$g+A>0$，最低位置$\theta=0$附近的线性化方程为
$\ddot\theta+(g+A)\theta/l=0$，故该位置稳定，频率为
$\sqrt{(g+A)/l}$。
若 $A=-g$，方程为 $\ddot\theta=0$，故 $\dot\theta=\dot\theta_0$，再次积分得到 $\theta=\theta_0+\int_0^t\dot\theta_0\dd s=\theta_0+\dot\theta_0t$。任意静止方向均为平衡，没有回复作用；但每个静止点对一般相空间扰动并不稳定，因为任意小的非零初始角速度都会使偏角随时间积累。仅把初始角度稍作改变、仍保持角速度为零，才会停在另一个平衡位置。
若$A<-g$，最低位置不稳定。令$\psi=\theta-\pi$，上方位置附近满足
$\ddot\psi+[-(g+A)/l]\psi=0$，稳定频率为$\sqrt{-(g+A)/l}$。

该积分对应相对悬点的运动。悬点装置能够做功，惯性系中的$T+U$一般不守恒。若$\ddot Y$随时间改变，相对势能显含时间；上述形式的时间导数还含$ml\dddot Y(1-\cos\theta)$，一般不再为零。
\end{solution}

\par\medskip\noindent\begin{minipage}{\textwidth}
\begin{exercise}[墙面也有摩擦的梯子（选做）]\label{mp5:exercise:l5-h5}
沿用例\ref{mp5:example:l5-e4}的梯子，地面摩擦系数为 $\mu>0$，墙面摩擦系数为 $\nu\ge0$，见图\ref{mp5:fig:l5-h5-q}。考虑梯脚将向外、梯顶将向下的滑动趋势，求存在静力平衡的最小倾角条件。两处静摩擦力均须保留各自的取值范围。
\end{exercise}

\begin{figure}[H]\centering
\begin{tikzpicture}
\draw[curve](0,-.35)--(0,3.9);\draw[curve](-.3,0)--(3.8,0);\draw[line width=1pt](3,0)--(0,3);\fill(3,0)circle(1.3pt)node[below right]{$A$};
\fill(0,3)circle(1.3pt)node[above left]{$B$};
\fill(1.5,1.5)circle(1.3pt)node[above right]{$G$};
\draw(2.4,0)arc(180:135:.6);\node at(2.25,.28){$\theta$};\node[above]at(1,2.2){$L$};\node at(3.1,-.6){$\mu$};\node[left]at(-.08,1.6){$\nu$};
\end{tikzpicture}
\caption{两个粗糙接触面。地面摩擦阻止梯脚外滑，墙面摩擦阻止梯顶下滑。}\label{mp5:fig:l5-h5-q}
\end{figure}
\end{minipage}\par\medskip
\begin{solution}
按图\ref{mp5:fig:l5-h5-s}左图，墙面法向力 $N_B$ 向右、摩擦力 $f_B$ 向上；地面摩擦力 $f_A$ 向左。平衡关系为
\[
f_A=N_B,\quad N_A=mg-f_B,\quad
N_B\tan\theta+f_B=\frac{mg}{2}.
\]
其中力矩方程可由绕 A 点取矩得到：
\[
N_B L\sin\theta+f_B L\cos\theta
-\frac{mgL}{2}\cos\theta=0.
\]
在 $0<\theta<\pi/2$ 内除以 $L\cos\theta$ 后即得上式。墙面摩擦承担一部分重量，既直接贡献支撑力矩，又减小地面正压力，所以两处摩擦上限不能独立地同时代入再随意消去反力。

令 $n=N_B/(mg)$、$r=f_B/(mg)$，允许的接触力满足
\[
n\ge0,\quad r\ge0,\quad r\le\nu n,\quad n\le\mu(1-r).
\]

\begin{figure}[H]\centering
\begin{tikzpicture}
\begin{scope}[shift={(0,0)},scale=0.8]
\draw[curve](0,-.35)--(0,3.9);\draw[curve](-.3,0)--(3.8,0);\draw[line width=1pt](3,0)--(0,3);\fill(3,0)circle(1.3pt)node[below right]{$A$};
\fill(0,3)circle(1.3pt)node[above left]{$B$};
\fill(1.5,1.5)circle(1.3pt)node[above right]{$G$};
\draw(2.4,0)arc(180:135:.6);\node at(2.25,.28){$\theta$};\node[above]at(1,2.2){$L$};\draw[axis](0,3)--(1.1,3)node[above]{$N_B$};\draw[axis](3,0)--(3,1.05)node[right]{$N_A$};\draw[axis](3,0)--(2,0)node[midway,below]{$f_A$};\draw[axis](1.5,1.5)--(1.5,.6)node[left]{$mg$};\draw[axis](0,3)--(0,3.8)node[left]{$f_B$};\end{scope}
\begin{scope}[shift={(6,0)},scale=4]
\draw[axis](-0.08,0)--(0.75,0)node[right]{$n$};
\draw[axis](0,-0.1)--(0,0.7)node[above]{$r$};
\fill[gray!18](0,0)--(.5,0)--(.4,.2)--cycle;\draw[curve](0,0)--(.5,0)--(.4,.2)--cycle;\draw[guide](0,0)--(.65,.325)node[right]{$r=\nu n$};\draw[curve,line width=.4pt](0,.5)--(.65,.0125);\fill(.4,.2)circle(.35pt);\draw[guide](.4,.2)--(.62,-.12);\node[below]at(.62,-.12){$(n_*,r_*)$};\node[right]at(.15,.64){$n\tan\theta+r=1/2$};\end{scope}

\end{tikzpicture}
\caption{左图为接触力，右图取 $\mu=\nu=1/2$。可行区域的斜顶点对应两处摩擦均达到上限，粗线对应临界倾角。}\label{mp5:fig:l5-h5-s}
\end{figure}

图右取 $\mu=\nu=1/2$，画出接触力的可行三角形。平衡还要求直线 $n\tan\theta+r=1/2$ 与可行区域相交。对 $r<1$，由 $n\le\mu(1-r)$ 可得
\[
\tan\theta=\frac{1/2-r}{n}\ge\frac{1/2-r}{\mu(1-r)}
\quad(0\le r\le1/2).
\]
又由 $r\le\nu n\le\mu\nu(1-r)$，得 $r\le\mu\nu/(1+\mu\nu)$。前述角度下界的商导数为
\[
\frac{\dd}{\dd r}\frac{1/2-r}{\mu(1-r)}
=\frac{-(1-r)+(1/2-r)}{\mu(1-r)^2}
=-\frac1{2\mu(1-r)^2}<0.
\]
因此下界随 $r$ 增大而减小。当 $\mu\nu<1$ 时，最大可行的 $r$ 在两条摩擦边界交点处取得，对应最小倾角：
\[
n_*=\frac{\mu}{1+\mu\nu},\qquad
r_*=\frac{\mu\nu}{1+\mu\nu}<\frac12.
\]
代回得
\[
\boxed{\tan\theta\ge\frac{1-\mu\nu}{2\mu}\qquad(\mu\nu<1).}
\]
若 $\mu\nu\ge1$，可行区已经能容纳 $r=1/2$ 的极限，任意 $\theta>0$ 都有满足上述方向约定的静力平衡。也可合写为 $\tan\theta\ge\max\{0,(1-\mu\nu)/(2\mu)\}$，并保持 $\theta>0$。

取 $\nu=0$ 即还原光滑墙面结果。在 $\mu\nu<1$ 的最小角度临界构形中，两处摩擦同时达到上限；角度更大时，接触力可有余量。
所得角度条件还可用一组明确的接触力证明其充分性。令 $t=\tan\theta>0$，取墙面摩擦达到边界 $f_B=\nu N_B$，则
\[
N_B=f_A=\frac{mg}{2(t+\nu)},\qquad
f_B=\frac{\nu mg}{2(t+\nu)},\qquad
N_A=\frac{(2t+\nu)mg}{2(t+\nu)}.
\]
这些力均非负，并满足三个静力方程。地面摩擦条件化为
\[
\frac{f_A}{N_A}=\frac1{2t+\nu}\le\mu
\quad\Longleftrightarrow\quad
t\ge\frac{1-\mu\nu}{2\mu}.
\]
因此满足角度下界时确实能找到允许的反力；当 $\mu\nu\ge1$，右端非正，所有普通正倾角都可行。在严格大于最小角度时，这组力的地面摩擦尚有余量；静力解也不必唯一。

\end{solution}

\par\medskip\noindent\begin{minipage}{\textwidth}
\begin{exercise}[双杆机构与压缩弹簧]\label{mp5:exercise:l5-h6}
两根长为 $L$ 的无质量刚杆下端共同连接质量为 $M$ 的物体，上端各接一个轻环，轻环可沿水平光滑直槽移动。机构左右对称，见图\ref{mp5:fig:l5-h6-q}。两环之间有劲度系数 $k$、自然长度 $\ell_0=3L/2$ 的弹簧。已知 $Mg=\sqrt3 kL$。以两环间距的一半 $x$ 为坐标，求平衡构形、杆内力及稳定性，讨论 $0<x<L$ 的对称分支。
\end{exercise}

\begin{figure}[H]\centering
\begin{tikzpicture}
\draw[curve](-2.5,0)--(2.5,0);\draw[curve](-1.25,0)--(0,-2.1650635094610964)--(1.25,0);\draw[fill=white](-1.25,0)circle(.13);\draw[fill=white](1.25,0)circle(.13);\fill(0,-2.1650635094610964)circle(2.5pt)node[right]{$M$};\draw(-1.25,.08)--(-1.05,.4)--(-.75,.1)--(-.45,.4)--(-.15,.1)--(.15,.4)--(.45,.1)--(.75,.4)--(1.05,.1)--(1.25,.08);\node[above]at(0,.55){$k,\ell_0$};\node[left]at(-.72,-1.1){$L$};\node[right]at(.72,-1.1){$L$};\draw[guide](0,0)--(0,-2.1650635094610964);\draw[guide](-1.25,0)--(-1.25,1.4);\draw[guide](1.25,0)--(1.25,1.4);\draw[<->,>=Stealth](-1.25,1.3)--(1.25,1.3)node[midway,above]{$2x$};\draw[guide](0,-2.1650635094610964)--(1.9,-2.1650635094610964);\draw[<->,>=Stealth](1.8,0)--(1.8,-2.1650635094610964)node[midway,right]{$h$};
\end{tikzpicture}
\caption{对称双杆机构。弹簧可伸长或压缩，悬挂物的深度 $h$ 随环距变化。}\label{mp5:fig:l5-h6-q}
\end{figure}
\end{minipage}\par\medskip
\begin{solution}
图\ref{mp5:fig:l5-h6-q}中 $h=\sqrt{L^2-x^2}$，弹簧长度为 $2x$。忽略常量，势能为
\[
U(x)=-Mg\sqrt{L^2-x^2}+\frac k2(2x-\ell_0)^2.
\]
重力势能的导数为 $Mgx/\sqrt{L^2-x^2}$，弹簧项的链式导数为 $k(2x-\ell_0)\cdot2$。二次求导时，
\[
\frac{\dd}{\dd x}\frac{x}{\sqrt{L^2-x^2}}
=\frac1{\sqrt{L^2-x^2}}+\frac{x^2}{(L^2-x^2)^{3/2}}
=\frac{L^2}{(L^2-x^2)^{3/2}}.
\]
因此
\[
U'=\frac{Mgx}{\sqrt{L^2-x^2}}+2k(2x-\ell_0),
\quad U''=\frac{MgL^2}{(L^2-x^2)^{3/2}}+4k>0.
\]
图\ref{mp5:fig:l5-h6-s}的势能图也表明该分支只有一个极小点。代入给定参数，$x=L/2$ 时 $U'=Mg/\sqrt3-kL=0$，于是
\[
\boxed{x_e=L/2,\qquad h_e=\sqrt3 L/2.}
\]

\begin{figure}[H]\centering
\begin{tikzpicture}
\begin{scope}[xscale=2.8,yscale=1.3]\draw[axis](-0.1,0)--(1.2,0)node[right]{$q$};
\draw[axis](0,-2)--(0,1)node[above]{$U/(kL^2)$};
\draw[curve] plot coordinates{(0.00000000,-0.60705081) (0.00654494,-0.62656285) (0.01308960,-0.64582853) (0.01963369,-0.66484705) (0.02617695,-0.68361766) (0.03271908,-0.70213962) (0.03925982,-0.72041224) (0.04579887,-0.73843486) (0.05233596,-0.75620686) (0.05887080,-0.77372762) (0.06540313,-0.79099660) (0.07193265,-0.80801326) (0.07845910,-0.82477710) (0.08498218,-0.84128766) (0.09150162,-0.85754449) (0.09801714,-0.87354721) (0.10452846,-0.88929544) (0.11103531,-0.90478885) (0.11753740,-0.92002714) (0.12403445,-0.93501002) (0.13052619,-0.94973728) (0.13701234,-0.96420869) (0.14349262,-0.97842408) (0.14996676,-0.99238333) (0.15643447,-1.00608630) (0.16289547,-1.01953293) (0.16934950,-1.03272317) (0.17579628,-1.04565701) (0.18223553,-1.05833446) (0.18866696,-1.07075558) (0.19509032,-1.08292044) (0.20150532,-1.09482915) (0.20791169,-1.10648187) (0.21430915,-1.11787877) (0.22069744,-1.12902004) (0.22707626,-1.13990594) (0.23344536,-1.15053672) (0.23980446,-1.16091269) (0.24615329,-1.17103417) (0.25249158,-1.18090153) (0.25881905,-1.19051515) (0.26513543,-1.19987545) (0.27144045,-1.20898288) (0.27773385,-1.21783793) (0.28401534,-1.22644110) (0.29028468,-1.23479293) (0.29654157,-1.24289398) (0.30278577,-1.25074486) (0.30901699,-1.25834618) (0.31523498,-1.26569862) (0.32143947,-1.27280283) (0.32763018,-1.27965955) (0.33380686,-1.28626950) (0.33996924,-1.29263345) (0.34611706,-1.29875220) (0.35225005,-1.30462657) (0.35836795,-1.31025741) (0.36447050,-1.31564559) (0.37055744,-1.32079202) (0.37662850,-1.32569763) (0.38268343,-1.33036337) (0.38872197,-1.33479022) (0.39474386,-1.33897920) (0.40074883,-1.34293134) (0.40673664,-1.34664768) (0.41270703,-1.35012933) (0.41865974,-1.35337738) (0.42459451,-1.35639297) (0.43051110,-1.35917725) (0.43640924,-1.36173141) (0.44228869,-1.36405666) (0.44814919,-1.36615421) (0.45399050,-1.36802532) (0.45981236,-1.36967127) (0.46561452,-1.37109335) (0.47139674,-1.37229288) (0.47715876,-1.37327121) (0.48290034,-1.37402968) (0.48862124,-1.37456970) (0.49432121,-1.37489267) (0.50000000,-1.37500000) (0.50565737,-1.37489315) (0.51129309,-1.37457359) (0.51690690,-1.37404279) (0.52249856,-1.37330228) (0.52806785,-1.37235357) (0.53361452,-1.37119820) (0.53913832,-1.36983775) (0.54463904,-1.36827378) (0.55011642,-1.36650791) (0.55557023,-1.36454174) (0.56100025,-1.36237692) (0.56640624,-1.36001509) (0.57178796,-1.35745792) (0.57714519,-1.35470709) (0.58247770,-1.35176431) (0.58778525,-1.34863129) (0.59306763,-1.34530976) (0.59832460,-1.34180146) (0.60355594,-1.33810816) (0.60876143,-1.33423163) (0.61394084,-1.33017365) (0.61909395,-1.32593604) (0.62422054,-1.32152059) (0.62932039,-1.31692915) (0.63439328,-1.31216355) (0.63943900,-1.30722565) (0.64445733,-1.30211730) (0.64944805,-1.29684038) (0.65441095,-1.29139677) (0.65934582,-1.28578838) (0.66425244,-1.28001711) (0.66913061,-1.27408488) (0.67398011,-1.26799360) (0.67880075,-1.26174523) (0.68359230,-1.25534169) (0.68835458,-1.24878495) (0.69308736,-1.24207696) (0.69779046,-1.23521969) (0.70246367,-1.22821511) (0.70710678,-1.22106521) (0.71171961,-1.21377198) (0.71630194,-1.20633741) (0.72085360,-1.19876350) (0.72537437,-1.19105226) (0.72986407,-1.18320569) (0.73432251,-1.17522581) (0.73874949,-1.16711465) (0.74314483,-1.15887422) (0.74750833,-1.15050655) (0.75183981,-1.14201368) (0.75613908,-1.13339763) (0.76040597,-1.12466045) (0.76464028,-1.11580416) (0.76884183,-1.10683081) (0.77301045,-1.09774244) (0.77714596,-1.08854109) (0.78124818,-1.07922879) (0.78531693,-1.06980760) (0.78935204,-1.06027956) (0.79335334,-1.05064670) (0.79732065,-1.04091107) (0.80125381,-1.03107470) (0.80515265,-1.02113964) (0.80901699,-1.01110791) (0.81284668,-1.00098155) (0.81664156,-0.99076260) (0.82040144,-0.98045307) (0.82412619,-0.97005500) (0.82781563,-0.95957039) (0.83146961,-0.94900128) (0.83508798,-0.93834966) (0.83867057,-0.92761754) (0.84221723,-0.91680693) (0.84572782,-0.90591982) (0.84920218,-0.89495820) (0.85264016,-0.88392405) (0.85604162,-0.87281936) (0.85940641,-0.86164608) (0.86273439,-0.85040618) (0.86602540,-0.83910162) (0.86927932,-0.82773433) (0.87249601,-0.81630627) (0.87567532,-0.80481936) (0.87881711,-0.79327552) (0.88192126,-0.78167666) (0.88498764,-0.77002468) (0.88801610,-0.75832148) (0.89100652,-0.74656893) (0.89395878,-0.73476891) (0.89687274,-0.72292328) (0.89974828,-0.71103388) (0.90258528,-0.69910255) (0.90538362,-0.68713113) (0.90814317,-0.67512141) (0.91086382,-0.66307520) (0.91354546,-0.65099430) (0.91618796,-0.63888047) (0.91879121,-0.62673547) (0.92135511,-0.61456106) (0.92387953,-0.60235896) (0.92636438,-0.59013091) (0.92880955,-0.57787860) (0.93121493,-0.56560372) (0.93358043,-0.55330795) (0.93590593,-0.54099295) (0.93819134,-0.52866037) (0.94043656,-0.51631183) (0.94264149,-0.50394895) (0.94480605,-0.49157333) (0.94693013,-0.47918653) (0.94901365,-0.46679014) (0.95105652,-0.45438569) (0.95305864,-0.44197471) (0.95501994,-0.42955872) (0.95694034,-0.41713920) (0.95881973,-0.40471764) (0.96065806,-0.39229549) (0.96245524,-0.37987420) (0.96421118,-0.36745517) (0.96592583,-0.35503981) (0.96759909,-0.34262951) (0.96923091,-0.33022563) (0.97082121,-0.31782950) (0.97236992,-0.30544247) (0.97387698,-0.29306582) (0.97534232,-0.28070085) (0.97676588,-0.26834881) (0.97814760,-0.25601096) (0.97948742,-0.24368850) (0.98078528,-0.23138266) (0.98204113,-0.21909460) (0.98325491,-0.20682549) (0.98442657,-0.19457646) (0.98555606,-0.18234863) (0.98664333,-0.17014310) (0.98768834,-0.15796095) (0.98869104,-0.14580322) (0.98965139,-0.13367094) (0.99056934,-0.12156512) (0.99144486,-0.10948675) (0.99227791,-0.09743680) (0.99306846,-0.08541619) (0.99381646,-0.07342586) (0.99452190,-0.06146669) (0.99518473,-0.04953957) (0.99580493,-0.03764533) (0.99638247,-0.02578480) (0.99691733,-0.01395880) (0.99740949,-0.00216810) (0.99785892,0.00958655) (0.99826561,0.02130440) (0.99862953,0.03298476) (0.99895068,0.04462692) (0.99922904,0.05623023) (0.99946459,0.06779405) (0.99965732,0.07931776) (0.99980724,0.09080076) (0.99991433,0.10224250) (0.99997858,0.11364242) (1.00000000,0.12500000)};
\fill(0.5,-1.375)circle(1.3pt)node[below]{$E$};
\draw[guide](.5,0)--(.5,-1.375);\node[above]at(.5,0){$1/2$};\draw[guide](1,-1.9)--(1,.5);\node[below right]at(1,0){$1$};\node[draw,circle,inner sep=1.3pt,fill=white]at(0,-.6070508075688772){};\node[draw,circle,inner sep=1.3pt,fill=white]at(1,.125){};\end{scope}
\end{tikzpicture}
\caption{取 $q=x/L$，纵轴为 $U/(kL^2)$。在 $0<q<1$ 内，唯一极小点为 $q=1/2$；两端空心点不属于所取分支。}\label{mp5:fig:l5-h6-s}
\end{figure}

在 $x=0$ 处 $U'=-3kL<0$，当 $x\uparrow L$ 时 $U'\to+\infty$，再结合 $U''>0$，可知上面的平衡点存在且唯一，并且稳定。

悬挂物的重力由两杆张力的竖直分量平衡。设杆内张力为 $S$，则 $2S h_e/L=Mg$，所以 $S=Mg/\sqrt3=kL$。弹簧压缩量为 $\ell_0-2x_e=L/2$，对每个环向外推 $kL/2$，恰好平衡杆力向内的水平分量 $Sx_e/L=kL/2$。
用结点受力还能直接得到平衡方程。设每根杆的张力为 S、压缩弹簧对每个环的外推力为 $F_s$。图\ref{mp5:fig:l5-h6-force}中，下端物体竖直平衡与右环水平平衡分别要求
\[
2S\frac hL=Mg,\qquad F_s=S\frac xL.
\]
故 $F_s=Mgx/(2h)$。压缩弹簧的力为 $F_s=k(\ell_0-2x)$，联立得到
\[
\frac{Mgx}{h}+2k(2x-\ell_0)=0,
\]
正好与 $U'=0$ 相同。平衡时 $F_s>0$，因而 $2x<\ell_0$。弹簧处于压缩状态，向外推环，抵消杆力向内的水平分量。
\begin{figure}[H]\centering
\begin{tikzpicture}[scale=.95]
\coordinate(C)at(0,0);\fill(C)circle(1.5pt)node[below left]{$C$};
\draw[axis](C)--(-.6,1.0392)node[above left]{$S$};\draw[axis](C)--(.6,1.0392)node[above right]{$S$};\draw[axis](C)--(0,-.9)node[below]{$Mg$};
\node at(0,1.8){悬挂物};
\begin{scope}[shift={(4.7,.8)}]
\draw[thick](-1.3,0)--(2,0);\draw[fill=white](0,0)circle(.12);\node[above left]at(0,.12){右环};
\draw[axis](0,0)--(-.6,-1.0392)node[below]{$S$};\draw[axis](0,0)--(1.3,0)node[above]{$F_s$};\draw[axis](0,0)--(0,1.1)node[right]{$N_{\text{槽}}$};
\end{scope}
\end{tikzpicture}
\caption{平衡构形下的两个隔离图。杆力沿杆轴，弹簧压缩后向外推右环；光滑双侧槽提供竖直反力。}\label{mp5:fig:l5-h6-force}
\end{figure}

由势能的二阶展开还可求出小位移对应的势能增量。在平衡点$x_e=L/2$处，
\[
 U''(x_e)=\frac{20}{3}k,
 \qquad U(x_e+\delta x)-U(x_e)=\frac{10}{3}k(\delta x)^2+O((\delta x)^3).
\]
系数为正，所以沿题设对称分支的小偏移会增加势能。图\ref{mp5:fig:l5-h6-force}中，两杆受拉，张力由竖直平衡确定。受压弹簧向外推环，其方向与上述势能极小点相符。
\end{solution}

\par\medskip\noindent\begin{minipage}{\textwidth}
\begin{exercise}[硬挡块处的边界平衡]\label{mp5:exercise:l5-h7}
滑块被限制在 $x\ge0$，$x=0$ 是硬挡块位置。弹簧沿 $x$ 方向的力为 $-k(x-a)$，其中 $a>0$；另有大小为 $F$ 的恒力向左，$F\ge0$，无摩擦。图\ref{mp5:fig:l5-h7-q}画出允许位置。求平衡位置、挡块反力与 $F$ 的关系，并用势能解释边界平衡为何不必满足 $U\prime(0)=0$。
\end{exercise}

\begin{figure}[H]\centering
\begin{tikzpicture}
\fill[gray!15](-1.1,-.45)rectangle(0,1.2);\draw[curve](0,-.45)--(0,1.2);\draw[axis](-.3,0)--(3,0)node[right]{$x$};\node[below left]at(0,0){$0$};\draw[fill=white](1,.25)rectangle(1.7,.9);\node at(1.35,.57){$m$};\draw[axis](1,.6)--(.2,.6)node[above]{$F$};\draw[guide](2.2,0)--(2.2,1.1);\node[below]at(2.2,0){$a$};
\end{tikzpicture}
\caption{阴影区域 $x<0$ 不可进入；没有恒力时，弹簧的平衡位置为 $a$。}\label{mp5:fig:l5-h7-q}
\end{figure}
\end{minipage}\par\medskip
\begin{solution}
设挡块对滑块的反力 $N$ 向右。单侧接触与静力方程应同时满足
\[
x\ge0,\qquad N\ge0,\qquad Nx=0,
\qquad -k(x-a)-F+N=0.
\]
若 $x>0$，反力必须为零；若 $x=0$，反力可为正。这两种分支分别求解，再与允许范围比较，便能避免把挡块外的平衡点当成实际位置。

总势能为 $U=\tfrac12k(x-a)^2+Fx$。若暂不考虑挡块，势能极小点为 $x_*=a-F/k$；再把取值范围限制为 $x\ge0$，得到
\[
\boxed{x_e=\max\{0,a-F/k\}.}
\]
图\ref{mp5:fig:l5-h7-s}分别给出内部极小点与边界极小点。若 $F<ka$，平衡位置在内部，反力为零；若 $F\ge ka$，滑块贴住挡块，水平方向平衡给出
\[
\boxed{N=F-ka\ge0.}
\]

\begin{figure}[H]\centering
\begin{tikzpicture}
\begin{scope}[xscale=2,yscale=1.5]\draw[axis](-0.2,0)--(2,0)node[right]{$q$};
\draw[axis](0,-0.3)--(0,1.8)node[above]{$u$};
\draw[curve] plot coordinates{(0.00000,0.00000) (0.01700,-0.00836) (0.03400,-0.01642) (0.05100,-0.02420) (0.06800,-0.03169) (0.08500,-0.03889) (0.10200,-0.04580) (0.11900,-0.05242) (0.13600,-0.05875) (0.15300,-0.06480) (0.17000,-0.07055) (0.18700,-0.07602) (0.20400,-0.08119) (0.22100,-0.08608) (0.23800,-0.09068) (0.25500,-0.09499) (0.27200,-0.09901) (0.28900,-0.10274) (0.30600,-0.10618) (0.32300,-0.10934) (0.34000,-0.11220) (0.35700,-0.11478) (0.37400,-0.11706) (0.39100,-0.11906) (0.40800,-0.12077) (0.42500,-0.12219) (0.44200,-0.12332) (0.45900,-0.12416) (0.47600,-0.12471) (0.49300,-0.12498) (0.51000,-0.12495) (0.52700,-0.12464) (0.54400,-0.12403) (0.56100,-0.12314) (0.57800,-0.12196) (0.59500,-0.12049) (0.61200,-0.11873) (0.62900,-0.11668) (0.64600,-0.11434) (0.66300,-0.11172) (0.68000,-0.10880) (0.69700,-0.10560) (0.71400,-0.10210) (0.73100,-0.09832) (0.74800,-0.09425) (0.76500,-0.08989) (0.78200,-0.08524) (0.79900,-0.08030) (0.81600,-0.07507) (0.83300,-0.06956) (0.85000,-0.06375) (0.86700,-0.05766) (0.88400,-0.05127) (0.90100,-0.04460) (0.91800,-0.03764) (0.93500,-0.03039) (0.95200,-0.02285) (0.96900,-0.01502) (0.98600,-0.00690) (1.00300,0.00150) (1.02000,0.01020) (1.03700,0.01918) (1.05400,0.02846) (1.07100,0.03802) (1.08800,0.04787) (1.10500,0.05801) (1.12200,0.06844) (1.13900,0.07916) (1.15600,0.09017) (1.17300,0.10146) (1.19000,0.11305) (1.20700,0.12492) (1.22400,0.13709) (1.24100,0.14954) (1.25800,0.16228) (1.27500,0.17531) (1.29200,0.18863) (1.30900,0.20224) (1.32600,0.21614) (1.34300,0.23032) (1.36000,0.24480) (1.37700,0.25956) (1.39400,0.27462) (1.41100,0.28996) (1.42800,0.30559) (1.44500,0.32151) (1.46200,0.33772) (1.47900,0.35422) (1.49600,0.37101) (1.51300,0.38808) (1.53000,0.40545) (1.54700,0.42310) (1.56400,0.44105) (1.58100,0.45928) (1.59800,0.47780) (1.61500,0.49661) (1.63200,0.51571) (1.64900,0.53510) (1.66600,0.55478) (1.68300,0.57474) (1.70000,0.59500)};
\draw[curve,dashed] plot coordinates{(0.00000,0.00000) (0.01350,0.00684) (0.02700,0.01386) (0.04050,0.02107) (0.05400,0.02846) (0.06750,0.03603) (0.08100,0.04378) (0.09450,0.05172) (0.10800,0.05983) (0.12150,0.06813) (0.13500,0.07661) (0.14850,0.08528) (0.16200,0.09412) (0.17550,0.10315) (0.18900,0.11236) (0.20250,0.12175) (0.21600,0.13133) (0.22950,0.14109) (0.24300,0.15102) (0.25650,0.16115) (0.27000,0.17145) (0.28350,0.18194) (0.29700,0.19260) (0.31050,0.20346) (0.32400,0.21449) (0.33750,0.22570) (0.35100,0.23710) (0.36450,0.24868) (0.37800,0.26044) (0.39150,0.27239) (0.40500,0.28451) (0.41850,0.29682) (0.43200,0.30931) (0.44550,0.32199) (0.45900,0.33484) (0.47250,0.34788) (0.48600,0.36110) (0.49950,0.37450) (0.51300,0.38808) (0.52650,0.40185) (0.54000,0.41580) (0.55350,0.42993) (0.56700,0.44424) (0.58050,0.45874) (0.59400,0.47342) (0.60750,0.48828) (0.62100,0.50332) (0.63450,0.51855) (0.64800,0.53395) (0.66150,0.54954) (0.67500,0.56531) (0.68850,0.58127) (0.70200,0.59740) (0.71550,0.61372) (0.72900,0.63022) (0.74250,0.64690) (0.75600,0.66377) (0.76950,0.68082) (0.78300,0.69804) (0.79650,0.71546) (0.81000,0.73305) (0.82350,0.75083) (0.83700,0.76878) (0.85050,0.78693) (0.86400,0.80525) (0.87750,0.82375) (0.89100,0.84244) (0.90450,0.86131) (0.91800,0.88036) (0.93150,0.89960) (0.94500,0.91901) (0.95850,0.93861) (0.97200,0.95839) (0.98550,0.97836) (0.99900,0.99850) (1.01250,1.01883) (1.02600,1.03934) (1.03950,1.06003) (1.05300,1.08090) (1.06650,1.10196) (1.08000,1.12320) (1.09350,1.14462) (1.10700,1.16622) (1.12050,1.18801) (1.13400,1.20998) (1.14750,1.23213) (1.16100,1.25446) (1.17450,1.27698) (1.18800,1.29967) (1.20150,1.32255) (1.21500,1.34561) (1.22850,1.36886) (1.24200,1.39228) (1.25550,1.41589) (1.26900,1.43968) (1.28250,1.46365) (1.29600,1.48781) (1.30950,1.51215) (1.32300,1.53666) (1.33650,1.56137) (1.35000,1.58625)};
\fill(0.5,-0.125)circle(1.3pt)node[below]{$E_1$};
\fill(0,0)circle(1.3pt)node[above left]{$E_2$};
\draw[curve](2.1,1.55)--(2.35,1.55)node[right]{$F/(ka)=1/2$};\draw[curve,dashed](2.1,1.25)--(2.35,1.25)node[right]{$F/(ka)=3/2$};\end{scope}
\end{tikzpicture}
\caption{取 $q=x/a$、$u=[U(x)-U(0)]/(ka^2)$。实线在内部点 $E_1$ 取最小值，虚线在边界点 $E_2$ 取最小值。}\label{mp5:fig:l5-h7-s}
\end{figure}

在$F>ka$时，$U'(0)=F-ka>0$，并非零。内部平衡点允许左右两向虚位移，故须有$U'(x_e)=0$；挡块处只允许$\delta x\ge0$，相应的一阶条件是$\delta U=U'(0)\delta x\ge0$。所需的向右反力正是$N=U'(0)=F-ka$。

在挡块处展开势能，有精确关系
\[
U(x)-U(0)=(F-ka)x+\frac12kx^2.
\]
当 $F>ka$，两个增量项对 $x>0$ 均为正；当 $F=ka$，二次项仍使势能增加。内部平衡 $x_e>0$ 则满足 $U(x)-U(x_e)=k(x-x_e)^2/2$。所以三个参数区间都对应允许位置范围内的势能最小值，只是支承反力是否参与平衡有所不同。

图\ref{mp5:fig:l5-h7-branches}将位置与反力随 F 的变化并列画出。在 $F=ka$ 处，平衡位置连续到达挡块，反力也从零连续开始增长；两条平衡分支在该点连续衔接。
\begin{figure}[H]\centering
\begin{tikzpicture}[scale=.95]
\draw[axis](-.15,0)--(3.8,0)node[right]{$F/(ka)$};\draw[axis](0,-.1)--(0,1.9)node[above]{$x_e/a$};
\draw[curve](0,1.5)--(1.7,0)--(3.4,0);\node[left]at(0,1.5){$1$};\node[below]at(1.7,0){$1$};\node[below]at(3.4,0){$2$};
\begin{scope}[shift={(5.5,0)}]
\draw[axis](-.15,0)--(3.8,0)node[right]{$F/(ka)$};\draw[axis](0,-.1)--(0,1.9)node[above]{$N/(ka)$};
\draw[curve](0,0)--(1.7,0)--(3.4,1.5);\draw[guide](0,1.5)--(3.4,1.5);\node[left]at(0,1.5){$1$};\node[below]at(1.7,0){$1$};\node[below]at(3.4,0){$2$};
\end{scope}
\end{tikzpicture}
\caption{平衡位置和挡块反力的两个分支。超过 $F=ka$ 后，增加的向左外力由挡块反力承担。}\label{mp5:fig:l5-h7-branches}
\end{figure}

这里讨论静态接触；若发生有限速度撞击挡块，后续运动还须给定碰撞规律。
\end{solution}

\par\medskip\noindent\begin{minipage}{\textwidth}
\begin{exercise}[斜向拉力下的滑动与倾覆]\label{mp5:exercise:l5-h8}
长方体的质量、底宽和接触条件同例\ref{mp5:example:l5-e6}。力 $F$ 改为作用在竖直中心线上、离地高度为 $H$ 的点 $P$，与水平面夹角为 $\gamma\in[0,\pi/2)$，方向向右上方，见图\ref{mp5:fig:l5-h8-q}。求滑动和向右倾覆阈值，比较两者先后，并检查地面支持力。
\end{exercise}

\begin{figure}[H]\centering
\begin{tikzpicture}
\draw[curve](-1.7,0)--(2.5,0);\draw[curve](-1,0)rectangle(1,2.8);\fill(0,1.4)circle(1.3pt)node[left]{$G$};
\fill(0,2.1)circle(1.3pt)node[above left]{$P$};
\draw[axis](0,2.1)--(1.8,3.13923)node[above right]{$F$};\draw[guide](0,2.1)--(1.1,2.1);\draw(.45,2.1)arc(0:30:.45);\node at(.67,2.3){$\gamma$};\draw[guide](-1,0)--(-1,-.7);\draw[guide](1,0)--(1,-.7);\draw[<->,>=Stealth](-1,-.6)--(1,-.6)node[midway,below]{$2a$};\draw[guide](0,2.1)--(-1.5,2.1);\draw[<->,>=Stealth](-1.4,0)--(-1.4,2.1)node[midway,left]{$H$};
\end{tikzpicture}
\caption{斜向拉力作用于中心线上的 $P$，其高度为 $H$；此处力的作用线与水平推力题不同。}\label{mp5:fig:l5-h8-q}
\end{figure}
\end{minipage}\par\medskip
\begin{solution}
图\ref{mp5:fig:l5-h8-s}中竖直分力减小地面支持力，水平分力需要由摩擦平衡：
\[
N=Mg-F\sin\gamma,\qquad f=F\cos\gamma.
\]
取底面中心为矩心，P 的水平坐标为零，因此
\[
Nx_N=FH\cos\gamma.
\]
达到滑动边界 $f=\mu N$ 时，
\[
F_{\rm slide}=\frac{\mu Mg}{\cos\gamma+\mu\sin\gamma}.
\]
达到倾覆边界 $x_N=a$ 时，
\[
F_{\rm tip}=\frac{aMg}{H\cos\gamma+a\sin\gamma}.
\]

\begin{figure}[H]\centering
\begin{tikzpicture}
\draw[curve](-1.7,0)--(2.5,0);\draw[curve](-1,0)rectangle(1,2.8);\fill(0,1.4)circle(1.3pt)node[left]{$G$};
\fill(0,2.1)circle(1.3pt)node[above left]{$P$};
\draw[axis](0,2.1)--(1.8,3.13923)node[above right]{$F$};\draw[guide](0,2.1)--(1.1,2.1);\draw(.45,2.1)arc(0:30:.45);\node at(.67,2.3){$\gamma$};\draw[axis](0,2.1)--(1.8,2.1)node[below]{$F\cos\gamma$};\draw[axis](0,2.1)--(0,3.13923)node[left]{$F\sin\gamma$};\draw[guide](0,3.13923)--(1.8,3.13923)--(1.8,2.1);\draw[axis](0,1.4)--(0,.65)node[left]{$Mg$};\draw[axis](.65,0)--(.65,.65)node[right]{$N$};\draw[axis](.65,0)--(-.25,0)node[midway,below]{$f$};\draw[guide](0,0)--(0,-.7);\draw[guide](.65,0)--(.65,-.7);\draw[<->,>=Stealth](0,-.6)--(.65,-.6)node[midway,below]{$x_N$};
\end{tikzpicture}
\caption{斜向拉力在 $P$ 分解。$N,f$ 为底面接触合力的分量，$x_N$ 从底面中心量至支持力作用线。}\label{mp5:fig:l5-h8-s}
\end{figure}

在给定角度范围内 $\cos\gamma>0$。比较两个分式并约去公共项，有
\[
F_{\rm slide}<F_{\rm tip}\quad\Longleftrightarrow\quad\mu H<a.
\]
所以先后判据仍为 $\mu<a/H$ 时先滑动，$\mu>a/H$ 时先倾覆，等号时同时临界。拉力角度改变了两个阈值，二者先后仍只由 $\mu H$ 与 a 的大小决定。

当 $\sin\gamma>0$，由 $N=0$ 得到的形式阈值为 $Mg/\sin\gamma$。滑动、倾覆阈值均小于它，因此从零增大 $F$ 时会先出现滑动或倾覆，此时 $N>0$。$\gamma=0$ 时没有这项竖直失接触阈值。
也可直接考察每个临界点是否满足另一项接触条件。在滑动候选点，
\[
N_{\rm slide}=\frac{Mg\cos\gamma}{\cos\gamma+\mu\sin\gamma}>0,
\qquad x_{N,\rm slide}=\mu H.
\]
若 $\mu H<a$，合力作用线仍位于底面内，故确实先滑动。在倾覆候选点，
\[
N_{\rm tip}=\frac{MgH\cos\gamma}{H\cos\gamma+a\sin\gamma}>0,
\qquad \frac{f_{\rm tip}}{N_{\rm tip}}=\frac aH.
\]
若 $\mu H>a$，摩擦还有余量，可以先到达底面右边缘。图\ref{mp5:fig:l5-h8-s}中的竖直分力虽然同时降低了两种临界力，但在比较它们的先后时，相同的竖直卸载项抵消了。初次失稳之后，原来的静力分支已不能继续使用。

取右侧底边为矩心，还可独立求倾覆候选值。重力的逆时针力矩为$Mga$；$P$点相对该边的位置为$(-a,H)$，拉力的顺时针力矩为$F(a\sin\gamma+H\cos\gamma)$。临界平衡给出
\[
 Mga=F(a\sin\gamma+H\cos\gamma),
\]
与上式的$F_{\rm tip}$相同。图\ref{mp5:fig:l5-h8-s}中，支持力合力到达右边缘时，还须满足静摩擦条件。力矩平衡给出倾覆候选值；若此前已经滑动，就应改用滑动后的受力模型。
\end{solution}

\par\medskip\noindent\begin{minipage}{\textwidth}
\begin{exercise}[球面外侧的脱离点]\label{mp5:exercise:l5-h9}
质点从固定光滑球面的外侧、与竖直向上方向夹角 $\theta_0\in(0,\pi/2)$ 的位置由静止释放，半径为 $R$，见图\ref{mp5:fig:l5-h9-q}。求脱离球面的位置和速度，并写出脱离后的轨迹。最后取 $\theta_0\to0^+$ 的极限，说明它与恰好静止在球顶有何区别。
\end{exercise}

\begin{figure}[H]\centering
\begin{tikzpicture}
\draw[curve](0,0)circle(1.8);\draw[axis](-.2,0)--(3.4,0)node[right]{$x$};\draw[axis](0,-.3)--(0,2.3)node[above]{$y$};\draw[guide](0,0)--(0.9538546756197689,1.5264865730815667);\fill(0.9538546756197689,1.5264865730815667)circle(1.3pt)node[above right]{$P$};
\fill(0,0)circle(1.3pt)node[below left]{$O$};
\draw[guide](0,0)--(.4354594120794019,1.7465323072967938);\draw(0,.42)arc(90:76:.42);\node[left]at(-.2,.42){$\theta_0$};\draw[guide](-.18,.42)--(.05,.42);\draw(0,.6)arc(90:58.0:.6);\node at(.25,.8){$\theta$};\draw[axis](0.9538546756197689,1.5264865730815667)--(0.9538546756197689,0.6764865730815667)node[right]{$mg$};\draw[axis](0.9538546756197689,1.5264865730815667)--(1.3512941237946725,2.1625226451988864)node[right]{$N$};\fill(0.4354594120794019,1.7465323072967938)circle(1.3pt)node[above left]{$P_0$};

\end{tikzpicture}
\caption{球面的竖直截面。释放角 $\theta_0$ 与当前位置角 $\theta$ 均自竖直向上方向量起，支持力沿外法线。}\label{mp5:fig:l5-h9-q}
\end{figure}
\end{minipage}\par\medskip
\begin{solution}
图\ref{mp5:fig:l5-h9-q}中质点高度为 $R\cos\theta$。接触期间支持力不做功，故
\[
v^2=2gR(\cos\theta_0-\cos\theta).
\]
沿球心方向写牛顿方程，重力有向内分量 $mg\cos\theta$，支持力向外：
\[
mg\cos\theta-N=\frac{mv^2}R,
\quad N=mg(3\cos\theta-2\cos\theta_0).
\]
单侧球面不能拉住质点，必须 $N\ge0$。首次达到零时脱离，于是
\[
\boxed{\cos\theta_d=\frac23\cos\theta_0,\qquad
v_d^2=\frac23gR\cos\theta_0.}
\]

\begin{figure}[H]\centering
\begin{tikzpicture}
\draw[curve](0,0)circle(1.8);\draw[axis](-.2,0)--(3.4,0)node[right]{$x$};\draw[axis](0,-.3)--(0,2.3)node[above]{$y$};\draw[guide](0,0)--(1.341640786499874,1.2);\fill(1.341640786499874,1.2)circle(1.3pt)node[above right]{$D$};
\fill(0,0)circle(1.3pt)node[below left]{$O$};
\draw(0,.6)arc(90:41.81031489577859:.6);\node at(.25,.8){$\theta_d$};\draw[curve] plot coordinates{(1.34164,1.20000) (1.35077,1.18972) (1.35990,1.17928) (1.36903,1.16868) (1.37816,1.15793) (1.38728,1.14702) (1.39641,1.13595) (1.40554,1.12473) (1.41467,1.11335) (1.42380,1.10182) (1.43293,1.09013) (1.44206,1.07828) (1.45119,1.06628) (1.46031,1.05412) (1.46944,1.04180) (1.47857,1.02933) (1.48770,1.01670) (1.49683,1.00392) (1.50596,0.99098) (1.51509,0.97788) (1.52421,0.96463) (1.53334,0.95122) (1.54247,0.93765) (1.55160,0.92393) (1.56073,0.91005) (1.56986,0.89602) (1.57899,0.88183) (1.58812,0.86748) (1.59724,0.85298) (1.60637,0.83832) (1.61550,0.82350) (1.62463,0.80853) (1.63376,0.79340) (1.64289,0.77812) (1.65202,0.76268) (1.66115,0.74708) (1.67027,0.73133) (1.67940,0.71542) (1.68853,0.69935) (1.69766,0.68313) (1.70679,0.66675) (1.71592,0.65022) (1.72505,0.63353) (1.73418,0.61668) (1.74330,0.59968) (1.75243,0.58252) (1.76156,0.56520) (1.77069,0.54773) (1.77982,0.53010) (1.78895,0.51232) (1.79808,0.49438) (1.80720,0.47628) (1.81633,0.45803) (1.82546,0.43962) (1.83459,0.42105) (1.84372,0.40233) (1.85285,0.38345) (1.86198,0.36442) (1.87111,0.34523) (1.88023,0.32588) (1.88936,0.30638) (1.89849,0.28672) (1.90762,0.26690) (1.91675,0.24693) (1.92588,0.22680) (1.93501,0.20652) (1.94414,0.18608) (1.95326,0.16548) (1.96239,0.14473) (1.97152,0.12382) (1.98065,0.10275) (1.98978,0.08153) (1.99891,0.06015) (2.00804,0.03862) (2.01717,0.01693) (2.02629,-0.00492) (2.03542,-0.02692) (2.04455,-0.04908) (2.05368,-0.07140) (2.06281,-0.09387) (2.07194,-0.11650) (2.08107,-0.13928) (2.09019,-0.16222) (2.09932,-0.18532) (2.10845,-0.20857) (2.11758,-0.23198) (2.12671,-0.25555) (2.13584,-0.27927) (2.14497,-0.30315) (2.15410,-0.32718) (2.16322,-0.35137) (2.17235,-0.37572) (2.18148,-0.40022) (2.19061,-0.42488) (2.19974,-0.44970) (2.20887,-0.47467) (2.21800,-0.49980) (2.22713,-0.52508) (2.23625,-0.55052) (2.24538,-0.57612) (2.25451,-0.60187) (2.26364,-0.62778) (2.27277,-0.65385) (2.28190,-0.68007) (2.29103,-0.70645) (2.30016,-0.73298) (2.30928,-0.75967) (2.31841,-0.78652) (2.32754,-0.81352) (2.33667,-0.84068) (2.34580,-0.86800) (2.35493,-0.89547) (2.36406,-0.92310) (2.37318,-0.95088) (2.38231,-0.97882) (2.39144,-1.00692) (2.40057,-1.03517) (2.40970,-1.06358) (2.41883,-1.09214) (2.42796,-1.12087) (2.43709,-1.14974)};
\draw[axis](1.341640786499874,1.2)--(1.8528485068380292,0.6284523933505917)node[right]{$\vct v_d$};\node at(2.5,1.9){$N_D=0$};
\end{tikzpicture}
\caption{脱离点 $D$ 与后续抛物线，取 $\theta_0\to0^+$ 的极限。$D$ 处速度沿圆的切线，抛物线从该点连续接出。}\label{mp5:fig:l5-h9-s}
\end{figure}

图\ref{mp5:fig:l5-h9-s}以球心为原点，D 为脱离点。切线方向向右下方，故
\[
x_d=R\sin\theta_d,\quad y_d=R\cos\theta_d,\quad
v_{xd}=v_d\cos\theta_d,\quad v_{yd}=-v_d\sin\theta_d.
\]
从脱离瞬间起令时间为 $\tau$，质点作抛体运动：
\[
\boxed{x=x_d+v_d\cos\theta_d\,\tau,\quad
y=y_d-v_d\sin\theta_d\,\tau-\tfrac12g\tau^2.}
\]
这段轨迹在球面外，不能把接触公式延伸到 $N<0$ 的区域。取 $\theta_0\to0^+$，得到 $\cos\theta_d\to2/3$、$v_d^2\to2gR/3$；图中按这一极限绘制。若初态恰是球顶且速度严格为零，理想方程的解会一直停在不稳定平衡点，不能把极限过程说成它必然自行滑落。
分离是否真实发生，可在图\ref{mp5:fig:l5-h9-s}的抛体轨迹上直接计算间隙。以球心为原点，将上式的 $x(\tau),y(\tau)$ 平方相加，按时间幂次整理，
\[
\begin{aligned}
x^2+y^2-R^2={}&(x_d^2+y_d^2-R^2)
+2(x_dv_{xd}+y_dv_{yd})\tau\\
&+(v_d^2-gy_d)\tau^2-gv_{yd}\tau^3+\frac14g^2\tau^4.
\end{aligned}
\]
前三项分别由球面初位置、切向初速度和 $v_d^2=gR\cos\theta_d=gy_d$ 消去，得到
\[
x^2+y^2-R^2
=g v_d\sin\theta_d\,\tau^3+\frac14g^2\tau^4>0
\qquad(\tau>0).
\]
零次项满足位于球面，一次项因速度沿切线为零，二次项又因临界法向方程相消；此后的正间隙由三次项开始出现。因此抛体段确实位于球面外侧，而不是仅仅把 $N=0$ 当作分离的充分条件。

题设 $\theta_0>0$ 使初始切向重力 $mg\sin\theta_0$ 为正，保证质点从释放时便开始下滑。

\end{solution}

\par\medskip\noindent\begin{minipage}{\textwidth}
\begin{exercise}[中心力轨道与径向势阱（选做）]\label{mp5:exercise:l5-h10}
质量为 $m$ 的质点受中心吸引力 $\vct F=-\kappa\vct e_r/r^2$，$\kappa>0$，对中心的角动量 $J>0$，见图\ref{mp5:fig:l5-h10-q}。求稳定圆轨道及其径向小振动频率，并由径向方程推导轨道形状，说明能量怎样决定椭圆、抛物线与双曲线。
\end{exercise}

\begin{figure}[H]\centering
\begin{tikzpicture}
\draw[axis](-1.8,0)--(2.3,0)node[right]{$x$};
\draw[axis](0,-1.1)--(0,1.9)node[above]{$y$};
\fill(0,0)circle(1.3pt)node[below left]{$O$};
\fill(1.25,0.8)circle(1.3pt)node[above right]{$P$};
\draw[guide](0,0)--(1.25,.8);\node[left]at(.35,.4){$r$};\draw[axis](1.25,.8)--(.5,.32);\node[above]at(.85,.9){$\vct F$};\draw[axis](1.25,.8)--(1.96,1.254)node[right]{$\dot r\,\vct e_r$};\draw[axis](1.25,.8)--(.77,1.55)node[left]{$r\dot\phi\,\vct e_\phi$};\draw(.55,0)arc(0:32.62:.55);\node at(.74,.2){$\phi$};
\end{tikzpicture}
\caption{力心 $O$、极径 $r$ 及质点的速度分量。径向、切向速度均从同一位置 $P$ 画出。}\label{mp5:fig:l5-h10-q}
\end{figure}
\end{minipage}\par\medskip
\begin{solution}
势能取 $U=-\kappa/r$。固定角动量后，
\[
E=\frac m2\dot r^2+\frac{J^2}{2mr^2}-\frac\kappa r,
\qquad V_{\rm eff}=\frac{J^2}{2mr^2}-\frac\kappa r.
\]
圆轨道要求径向保持不变，因此
\[
V_{\rm eff}'=-\frac{J^2}{mr^3}+\frac\kappa{r^2}=0
\quad\Longleftrightarrow\quad J^2=m\kappa r_0.
\]
先解出$r_0$，再代入$V_{\rm eff}$，得
\[
\boxed{r_0=\frac{J^2}{m\kappa},\qquad
E_{\min}=-\frac{m\kappa^2}{2J^2}.}
\]
条件$J^2=m\kappa r_0$也等价于$mv_\phi^2/r_0=\kappa/r_0^2$，因为$v_\phi=J/(mr_0)$。径向小扰动应在同一个固定$J$的条件下取$r=r_0+\xi$；角速度会随$r$改变，不能把它额外固定。

再次求导，$V_{\rm eff}''=3J^2/(mr^4)-2\kappa/r^3$。在 $r_0$ 处用 $J^2=m\kappa r_0$，得 $V_{\rm eff}''=3\kappa/r_0^3-2\kappa/r_0^3=\kappa/r_0^3>0$。因此
\[
\omega_r^2=\frac\kappa{mr_0^3}.
\]
圆轨道的转动角速度为 $\Omega=J/(mr_0^2)$，代入 $J^2=m\kappa r_0$，有 $\Omega^2=\omega_r^2$。径向小振动与绕中心转动同频。

\begin{figure}[H]\centering
\begin{tikzpicture}
\begin{scope}[shift={(0,0)},scale=0.85]
\draw[axis](-3.4,0)--(1.5,0)node[right]{$x$};
\draw[axis](0,-2)--(0,2)node[above]{$y$};
\draw[flow] plot coordinates{(1.00000,0.00000) (0.99985,0.02094) (0.99942,0.04189) (0.99868,0.06283) (0.99766,0.08378) (0.99634,0.10472) (0.99473,0.12566) (0.99282,0.14661) (0.99062,0.16755) (0.98812,0.18849) (0.98532,0.20944) (0.98223,0.23038) (0.97883,0.25132) (0.97514,0.27226) (0.97114,0.29320) (0.96683,0.31414) (0.96222,0.33508) (0.95730,0.35601) (0.95206,0.37695) (0.94652,0.39788) (0.94065,0.41881) (0.93447,0.43973) (0.92797,0.46065) (0.92114,0.48156) (0.91399,0.50247) (0.90651,0.52337) (0.89869,0.54427) (0.89054,0.56515) (0.88204,0.58603) (0.87321,0.60689) (0.86402,0.62775) (0.85448,0.64859) (0.84459,0.66941) (0.83434,0.69022) (0.82372,0.71101) (0.81273,0.73178) (0.80137,0.75253) (0.78963,0.77326) (0.77750,0.79396) (0.76499,0.81463) (0.75208,0.83527) (0.73877,0.85587) (0.72505,0.87644) (0.71093,0.89697) (0.69638,0.91745) (0.68142,0.93789) (0.66602,0.95827) (0.65018,0.97861) (0.63391,0.99888) (0.61718,1.01909) (0.60000,1.03923) (0.58235,1.05930) (0.56424,1.07929) (0.54565,1.09920) (0.52657,1.11901) (0.50700,1.13874) (0.48693,1.15835) (0.46635,1.17786) (0.44526,1.19726) (0.42364,1.21653) (0.40149,1.23566) (0.37880,1.25466) (0.35557,1.27351) (0.33178,1.29220) (0.30743,1.31072) (0.28250,1.32906) (0.25699,1.34721) (0.23090,1.36516) (0.20421,1.38290) (0.17691,1.40041) (0.14901,1.41769) (0.12048,1.43471) (0.09132,1.45147) (0.06153,1.46795) (0.03109,1.48413) (0.00000,1.50000) (-0.03175,1.51554) (-0.06416,1.53073) (-0.09724,1.54556) (-0.13100,1.56001) (-0.16544,1.57405) (-0.20057,1.58767) (-0.23639,1.60084) (-0.27291,1.61354) (-0.31013,1.62575) (-0.34805,1.63744) (-0.38668,1.64860) (-0.42601,1.65919) (-0.46604,1.66918) (-0.50679,1.67856) (-0.54823,1.68728) (-0.59038,1.69533) (-0.63322,1.70268) (-0.67675,1.70928) (-0.72097,1.71511) (-0.76586,1.72014) (-0.81141,1.72433) (-0.85762,1.72766) (-0.90446,1.73007) (-0.95193,1.73155) (-1.00000,1.73205) (-1.04866,1.73154) (-1.09788,1.72998) (-1.14763,1.72733) (-1.19790,1.72355) (-1.24865,1.71861) (-1.29984,1.71248) (-1.35145,1.70510) (-1.40342,1.69645) (-1.45573,1.68648) (-1.50833,1.67517) (-1.56117,1.66247) (-1.61419,1.64836) (-1.66735,1.63279) (-1.72058,1.61573) (-1.77382,1.59715) (-1.82701,1.57703) (-1.88009,1.55534) (-1.93297,1.53205) (-1.98558,1.50714) (-2.03786,1.48059) (-2.08970,1.45238) (-2.14104,1.42250) (-2.19178,1.39095) (-2.24184,1.35771) (-2.29112,1.32278) (-2.33954,1.28617) (-2.38698,1.24788) (-2.43336,1.20793) (-2.47859,1.16633) (-2.52255,1.12311) (-2.56516,1.07829) (-2.60631,1.03191) (-2.64590,0.98400) (-2.68384,0.93461) (-2.72004,0.88379) (-2.75440,0.83160) (-2.78683,0.77810) (-2.81724,0.72334) (-2.84555,0.66742) (-2.87169,0.61040) (-2.89557,0.55236) (-2.91714,0.49340) (-2.93632,0.43360) (-2.95306,0.37306) (-2.96731,0.31188) (-2.97903,0.25016) (-2.98818,0.18800) (-2.99474,0.12552) (-2.99868,0.06281) (-3.00000,0.00000) (-2.99868,-0.06281) (-2.99474,-0.12552) (-2.98818,-0.18800) (-2.97903,-0.25016) (-2.96731,-0.31188) (-2.95306,-0.37306) (-2.93632,-0.43360) (-2.91714,-0.49340) (-2.89557,-0.55236) (-2.87169,-0.61040) (-2.84555,-0.66742) (-2.81724,-0.72334) (-2.78683,-0.77810) (-2.75440,-0.83160) (-2.72004,-0.88379) (-2.68384,-0.93461) (-2.64590,-0.98400) (-2.60631,-1.03191) (-2.56516,-1.07829) (-2.52255,-1.12311) (-2.47859,-1.16633) (-2.43336,-1.20793) (-2.38698,-1.24788) (-2.33954,-1.28617) (-2.29112,-1.32278) (-2.24184,-1.35771) (-2.19178,-1.39095) (-2.14104,-1.42250) (-2.08970,-1.45238) (-2.03786,-1.48059) (-1.98558,-1.50714) (-1.93297,-1.53205) (-1.88009,-1.55534) (-1.82701,-1.57703) (-1.77382,-1.59715) (-1.72058,-1.61573) (-1.66735,-1.63279) (-1.61419,-1.64836) (-1.56117,-1.66247) (-1.50833,-1.67517) (-1.45573,-1.68648) (-1.40342,-1.69645) (-1.35145,-1.70510) (-1.29984,-1.71248) (-1.24865,-1.71861) (-1.19790,-1.72355) (-1.14763,-1.72733) (-1.09788,-1.72998) (-1.04866,-1.73154) (-1.00000,-1.73205) (-0.95193,-1.73155) (-0.90446,-1.73007) (-0.85762,-1.72766) (-0.81141,-1.72433) (-0.76586,-1.72014) (-0.72097,-1.71511) (-0.67675,-1.70928) (-0.63322,-1.70268) (-0.59038,-1.69533) (-0.54823,-1.68728) (-0.50679,-1.67856) (-0.46604,-1.66918) (-0.42601,-1.65919) (-0.38668,-1.64860) (-0.34805,-1.63744) (-0.31013,-1.62575) (-0.27291,-1.61354) (-0.23639,-1.60084) (-0.20057,-1.58767) (-0.16544,-1.57405) (-0.13100,-1.56001) (-0.09724,-1.54556) (-0.06416,-1.53073) (-0.03175,-1.51554) (-0.00000,-1.50000) (0.03109,-1.48413) (0.06153,-1.46795) (0.09132,-1.45147) (0.12048,-1.43471) (0.14901,-1.41769) (0.17691,-1.40041) (0.20421,-1.38290) (0.23090,-1.36516) (0.25699,-1.34721) (0.28250,-1.32906) (0.30743,-1.31072) (0.33178,-1.29220) (0.35557,-1.27351) (0.37880,-1.25466) (0.40149,-1.23566) (0.42364,-1.21653) (0.44526,-1.19726) (0.46635,-1.17786) (0.48693,-1.15835) (0.50700,-1.13874) (0.52657,-1.11901) (0.54565,-1.09920) (0.56424,-1.07929) (0.58235,-1.05930) (0.60000,-1.03923) (0.61718,-1.01909) (0.63391,-0.99888) (0.65018,-0.97861) (0.66602,-0.95827) (0.68142,-0.93789) (0.69638,-0.91745) (0.71093,-0.89697) (0.72505,-0.87644) (0.73877,-0.85587) (0.75208,-0.83527) (0.76499,-0.81463) (0.77750,-0.79396) (0.78963,-0.77326) (0.80137,-0.75253) (0.81273,-0.73178) (0.82372,-0.71101) (0.83434,-0.69022) (0.84459,-0.66941) (0.85448,-0.64859) (0.86402,-0.62775) (0.87321,-0.60689) (0.88204,-0.58603) (0.89054,-0.56515) (0.89869,-0.54427) (0.90651,-0.52337) (0.91399,-0.50247) (0.92114,-0.48156) (0.92797,-0.46065) (0.93447,-0.43973) (0.94065,-0.41881) (0.94652,-0.39788) (0.95206,-0.37695) (0.95730,-0.35601) (0.96222,-0.33508) (0.96683,-0.31414) (0.97114,-0.29320) (0.97514,-0.27226) (0.97883,-0.25132) (0.98223,-0.23038) (0.98532,-0.20944) (0.98812,-0.18849) (0.99062,-0.16755) (0.99282,-0.14661) (0.99473,-0.12566) (0.99634,-0.10472) (0.99766,-0.08378) (0.99868,-0.06283) (0.99942,-0.04189) (0.99985,-0.02094) (1.00000,-0.00000)};
\fill(0,0)circle(1.3pt)node[below]{$O$};
\fill(-1,0)circle(1.3pt)node[above]{$C$};
\fill(1,0)circle(1.3pt)node[above right]{$P$};
\fill(-3,0)circle(1.3pt)node[above left]{$A$};
\node at(-.9,-2.35){$e=1/2,\ p=3/2$};\end{scope}
\begin{scope}[shift={(5,0)},scale=1.05]
\draw[axis](0,0)--(3.5,0)node[right]{$r/r_0$};
\draw[axis](0,-0.65)--(0,1.4)node[above]{$V_{\rm eff}/(\kappa/r_0)$};
\draw[curve] plot coordinates{(0.42000,0.45351) (0.43440,0.34764) (0.44880,0.25419) (0.46320,0.17152) (0.47760,0.09820) (0.49200,0.03305) (0.50640,-0.02496) (0.52080,-0.07669) (0.53520,-0.12289) (0.54960,-0.16421) (0.56400,-0.20120) (0.57840,-0.23435) (0.59280,-0.26408) (0.60720,-0.29076) (0.62160,-0.31471) (0.63600,-0.33622) (0.65040,-0.35554) (0.66480,-0.37289) (0.67920,-0.38846) (0.69360,-0.40243) (0.70800,-0.41495) (0.72240,-0.42617) (0.73680,-0.43620) (0.75120,-0.44515) (0.76560,-0.45313) (0.78000,-0.46022) (0.79440,-0.46651) (0.80880,-0.47206) (0.82320,-0.47694) (0.83760,-0.48120) (0.85200,-0.48491) (0.86640,-0.48811) (0.88080,-0.49084) (0.89520,-0.49315) (0.90960,-0.49506) (0.92400,-0.49662) (0.93840,-0.49785) (0.95280,-0.49877) (0.96720,-0.49942) (0.98160,-0.49982) (0.99600,-0.49999) (1.01040,-0.49995) (1.02480,-0.49971) (1.03920,-0.49929) (1.05360,-0.49871) (1.06800,-0.49797) (1.08240,-0.49710) (1.09680,-0.49611) (1.11120,-0.49499) (1.12560,-0.49377) (1.14000,-0.49246) (1.15440,-0.49106) (1.16880,-0.48957) (1.18320,-0.48801) (1.19760,-0.48639) (1.21200,-0.48470) (1.22640,-0.48296) (1.24080,-0.48117) (1.25520,-0.47933) (1.26960,-0.47745) (1.28400,-0.47554) (1.29840,-0.47359) (1.31280,-0.47161) (1.32720,-0.46961) (1.34160,-0.46758) (1.35600,-0.46554) (1.37040,-0.46347) (1.38480,-0.46139) (1.39920,-0.45930) (1.41360,-0.45720) (1.42800,-0.45508) (1.44240,-0.45296) (1.45680,-0.45084) (1.47120,-0.44871) (1.48560,-0.44658) (1.50000,-0.44444) (1.51440,-0.44231) (1.52880,-0.44018) (1.54320,-0.43805) (1.55760,-0.43592) (1.57200,-0.43380) (1.58640,-0.43168) (1.60080,-0.42957) (1.61520,-0.42746) (1.62960,-0.42537) (1.64400,-0.42327) (1.65840,-0.42119) (1.67280,-0.41912) (1.68720,-0.41705) (1.70160,-0.41500) (1.71600,-0.41295) (1.73040,-0.41092) (1.74480,-0.40889) (1.75920,-0.40688) (1.77360,-0.40488) (1.78800,-0.40288) (1.80240,-0.40091) (1.81680,-0.39894) (1.83120,-0.39698) (1.84560,-0.39504) (1.86000,-0.39311) (1.87440,-0.39119) (1.88880,-0.38929) (1.90320,-0.38739) (1.91760,-0.38551) (1.93200,-0.38364) (1.94640,-0.38179) (1.96080,-0.37995) (1.97520,-0.37812) (1.98960,-0.37630) (2.00400,-0.37450) (2.01840,-0.37271) (2.03280,-0.37093) (2.04720,-0.36917) (2.06160,-0.36742) (2.07600,-0.36568) (2.09040,-0.36395) (2.10480,-0.36224) (2.11920,-0.36054) (2.13360,-0.35886) (2.14800,-0.35718) (2.16240,-0.35552) (2.17680,-0.35387) (2.19120,-0.35223) (2.20560,-0.35061) (2.22000,-0.34900) (2.23440,-0.34740) (2.24880,-0.34581) (2.26320,-0.34424) (2.27760,-0.34267) (2.29200,-0.34112) (2.30640,-0.33958) (2.32080,-0.33805) (2.33520,-0.33654) (2.34960,-0.33503) (2.36400,-0.33354) (2.37840,-0.33206) (2.39280,-0.33059) (2.40720,-0.32913) (2.42160,-0.32769) (2.43600,-0.32625) (2.45040,-0.32483) (2.46480,-0.32341) (2.47920,-0.32201) (2.49360,-0.32062) (2.50800,-0.31923) (2.52240,-0.31786) (2.53680,-0.31650) (2.55120,-0.31515) (2.56560,-0.31381) (2.58000,-0.31248) (2.59440,-0.31116) (2.60880,-0.30985) (2.62320,-0.30855) (2.63760,-0.30726) (2.65200,-0.30598) (2.66640,-0.30471) (2.68080,-0.30345) (2.69520,-0.30220) (2.70960,-0.30096) (2.72400,-0.29972) (2.73840,-0.29850) (2.75280,-0.29729) (2.76720,-0.29608) (2.78160,-0.29488) (2.79600,-0.29370) (2.81040,-0.29252) (2.82480,-0.29135) (2.83920,-0.29019) (2.85360,-0.28903) (2.86800,-0.28789) (2.88240,-0.28675) (2.89680,-0.28562) (2.91120,-0.28450) (2.92560,-0.28339) (2.94000,-0.28229) (2.95440,-0.28119) (2.96880,-0.28011) (2.98320,-0.27903) (2.99760,-0.27796) (3.01200,-0.27689) (3.02640,-0.27584) (3.04080,-0.27479) (3.05520,-0.27374) (3.06960,-0.27271) (3.08400,-0.27168) (3.09840,-0.27066) (3.11280,-0.26965) (3.12720,-0.26865) (3.14160,-0.26765) (3.15600,-0.26666) (3.17040,-0.26567) (3.18480,-0.26470) (3.19920,-0.26373) (3.21360,-0.26276) (3.22800,-0.26180) (3.24240,-0.26085) (3.25680,-0.25991) (3.27120,-0.25897) (3.28560,-0.25804) (3.30000,-0.25712)};
\fill(1,-0.5)circle(1.3pt)node[below]{$E_0$};
\draw[guide](.65,-.4)--(3.25,-.4);\node[right]at(3.25,-.4){$E<0$};\node[below]at(1,0){$1$};\end{scope}

\end{tikzpicture}
\caption{左图为椭圆轨道，箭头按 $J>0$ 沿逆时针方向绘制。$O$ 位于焦点，$C$ 是几何中心，$P,A$ 分别为近、远心点；右图为径向约化势能。}\label{mp5:fig:l5-h10-s}
\end{figure}

图\ref{mp5:fig:l5-h10-s}左图中，力心 $O$ 位于椭圆的一个焦点，$C$ 为几何中心。求精确轨道时，令 $u=1/r$，对 $\phi$ 的导数记为撇号。由 $\dot\phi=Ju^2/m$ 得
\[
\dot r=-u^{-2}u'\dot\phi=-\frac Jm u',\qquad
\ddot r=-\frac Jm u''\dot\phi=-\frac{J^2}{m^2}u^2u''.
\]
将它们代入径向方程 $m(\ddot r-r\dot\phi^2)=-\kappa/r^2$，约去非零因子，得到
\[
u''+u=\frac{m\kappa}{J^2}.
\]
取 $p=J^2/(m\kappa)$，特解为 $u=1/p$，齐次方程 $w''+w=0$ 的两根为 $\pm\mathrm i$，所以 $u=1/p+C\cos\phi+D\sin\phi$。选择最近点方向为 $\phi=0$，$u'(0)=0$ 给出 $D=0$；最近点的 $u$ 最大，故 $C\ge0$。令 $e=pC\ge0$，轨道为
\[
\boxed{r=\frac p{1+e\cos\phi},\qquad
p=\frac{J^2}{m\kappa},\qquad e\ge0.}
\]
再把 $u=(1+e\cos\phi)/p$、$u'=-e\sin\phi/p$ 代入能量，得
\begin{align*}
E&=\frac{J^2}{2m}\bigl[(u')^2+u^2\bigr]-\kappa u\\
&=\frac{J^2}{2mp^2}\Bigl[e^2\sin^2\phi+(1+e\cos\phi)^2\\
&\hspace{5em}{}-2(1+e\cos\phi)\Bigr]\\
&=\frac{m\kappa^2}{2J^2}(e^2-1).
\end{align*}
所以 $e^2=1+2EJ^2/(m\kappa^2)$。右图的径向势阱表明 $E\ge E_{\min}$；$E=E_{\min}$ 是圆，$E_{\min}<E<0$ 是椭圆，$E=0$ 是抛物线，$E>0$ 是双曲线。后两种轨道只取 $1+e\cos\phi>0$ 的物理分支。椭圆近心点、远心点分别为 $p/(1+e)$、$p/(1-e)$，与图中 P、A 对应。
径向转折点还可由 $E=V_{\rm eff}$ 求出：
\[
Er^2+\kappa r-\frac{J^2}{2m}=0.
\]
在 $E_{\min}<E<0$ 时，它有两个正根，分别为 $r_p=p/(1+e)$、$r_a=p/(1-e)$；两根相等时就是圆轨道。以这两根的平均值定义椭圆半长轴 $a_{\rm orb}$，有
\[
a_{\rm orb}=\frac{r_p+r_a}{2}=\frac p{1-e^2}
=-\frac\kappa{2E}.
\]
这与图\ref{mp5:fig:l5-h10-s}中力心为焦点、近远心距不相等的几何关系一致。$E=0$ 时唯一有限正转折点为 $p/2$，另一端延伸到无穷；$E>0$ 时同样只有一个正的最近距离，远端速度趋于 $\sqrt{2E/m}$。当 $J>0$，离心项在 $r\to0$ 时发散，轨道不会到达力心。

\end{solution}

\par\medskip\noindent\begin{minipage}{\textwidth}
\begin{exercise}[双杆机构的小振动]\label{mp5:exercise:l5-h11}
将习题\ref{mp5:exercise:l5-h6}中的两个轻环改为质量均为 $m$ 的小环，其余条件不变：两杆长 $L$，无质量，悬挂质量为 $M$，弹簧劲度系数为 $k$、自然长度 $3L/2$，且 $Mg=\sqrt3kL$。运动保持左右对称，见图\ref{mp5:fig:l5-h11-q}。求平衡位置附近的小振动频率，说明怎样计入悬挂质量的动能。
\end{exercise}

\begin{figure}[H]\centering
\begin{tikzpicture}
\draw[curve](-2.5,0)--(2.5,0);\draw[curve](-1.25,0)--(0,-2.1650635094610964)--(1.25,0);\draw[fill=white](-1.25,0)circle(.13);\draw[fill=white](1.25,0)circle(.13);\fill(0,-2.1650635094610964)circle(2.5pt)node[right]{$M$};\draw(-1.25,.08)--(-1.05,.4)--(-.75,.1)--(-.45,.4)--(-.15,.1)--(.15,.4)--(.45,.1)--(.75,.4)--(1.05,.1)--(1.25,.08);\node[above]at(0,.55){$k,\ell_0$};\node[left]at(-.72,-1.1){$L$};\node[right]at(.72,-1.1){$L$};\draw[guide](0,0)--(0,-2.1650635094610964);\node[above left]at(-1.25,.05){$m$};\node[above right]at(1.25,.05){$m$};\draw[guide](-1.25,0)--(-1.25,1.4);\draw[guide](1.25,0)--(1.25,1.4);\draw[<->,>=Stealth](-1.25,1.3)--(1.25,1.3)node[midway,above]{$2x$};\draw[guide](0,-2.1650635094610964)--(1.9,-2.1650635094610964);\draw[<->,>=Stealth](1.8,0)--(1.8,-2.1650635094610964)node[midway,right]{$h$};
\end{tikzpicture}
\caption{两个有质量的小环沿水平直槽运动，悬挂物在对称轴上运动，杆长保持不变。}\label{mp5:fig:l5-h11-q}
\end{figure}
\end{minipage}\par\medskip
\begin{solution}
两环的高度不变，其重力势能是常量。因此势能与习题\ref{mp5:exercise:l5-h6}相同，平衡位置仍为 $x_e=L/2$，$h_e=\sqrt3L/2$。环的质量只改变动能，下面由几何约束计算各部分速度。

图\ref{mp5:fig:l5-h11-s}中两环的水平速度大小都是 $|\dot x|$。取竖直向上为正，悬挂质量坐标 $y_M=-h=-\sqrt{L^2-x^2}$，所以
\[
\dot y_M=\frac xh\dot x.
\]
于是
\[
T=\frac12m\dot x^2+\frac12m\dot x^2
+\frac12M\frac{x^2}{h^2}\dot x^2
=\frac12\left(2m+M\frac{x^2}{h^2}\right)\dot x^2.
\]

\begin{figure}[H]\centering
\begin{tikzpicture}
\draw[curve](-2.5,0)--(2.5,0);\draw[curve](-1.25,0)--(0,-2.1650635094610964)--(1.25,0);\draw[fill=white](-1.25,0)circle(.13);\draw[fill=white](1.25,0)circle(.13);\fill(0,-2.1650635094610964)circle(2.5pt)node[right]{$M$};\draw(-1.25,.08)--(-1.05,.4)--(-.75,.1)--(-.45,.4)--(-.15,.1)--(.15,.4)--(.45,.1)--(.75,.4)--(1.05,.1)--(1.25,.08);\node[above]at(0,.55){$k,\ell_0$};\node[left]at(-.72,-1.1){$L$};\node[right]at(.72,-1.1){$L$};\draw[guide](0,0)--(0,-2.1650635094610964);\node[above left]at(-1.25,.05){$m$};\node[above right]at(1.25,.05){$m$};\draw[axis](-1.25,0)--(-2.05,0)node[above]{$\dot x$};\draw[axis](1.25,0)--(2.05,0)node[above]{$\dot x$};\draw[axis](0,-2.1650635094610964)--(0,-1.7031832941093956);\node[above,fill=white,inner sep=1pt]at(0,-1.57){$\dot y_M$};
\end{tikzpicture}
\caption{取 $\dot x>0$，两环向外、悬挂物向上。图取平衡构形，竖直与水平速度箭头的长度比为 $1:\sqrt3$。}\label{mp5:fig:l5-h11-s}
\end{figure}

在平衡位置，图中速度比为 $\dot y_M/\dot x=1/\sqrt3$，所以
\[
M_{\rm eff}(x_e)=2m+\frac M3.
\]
势能二阶导数由前题的表达式得到
\[
U''(x_e)=\frac{MgL^2}{h_e^3}+4k
=\frac{8k}3+4k=\frac{20k}3.
\]
因此
\[
\boxed{\omega^2=\frac{U''(x_e)}{M_{\rm eff}(x_e)}
=\frac{20k}{6m+M}.}
\]
悬挂物对等效质量的贡献为 $M/3$，来自平衡位置的速度比 $1/\sqrt3$。等效质量虽随 $x$ 变化，其导数项在本题的小振动线性方程中是二阶小量。
若暂不作小振动近似，令 $\mathcal M(x)=2m+Mx^2/(L^2-x^2)$，其导数为
\[
\mathcal M'(x)=M\frac{2x(L^2-x^2)-x^2(-2x)}{(L^2-x^2)^2}
=\frac{2MxL^2}{(L^2-x^2)^2}.
\]
完整的一自由度方程为
\[
\mathcal M(x)\ddot x+\frac{MxL^2}{(L^2-x^2)^2}\dot x^2
+\frac{Mgx}{\sqrt{L^2-x^2}}+2k(2x-\ell_0)=0.
\]
令 $x=x_e+\xi$，其中 $\xi,\dot\xi,\ddot\xi$ 都按振幅的一阶量计算。则 $\mathcal M'\dot\xi^2/2$ 和 $[\mathcal M(x)-\mathcal M(x_e)]\ddot\xi$ 都是二阶量；势能导数的一阶项为 $U''(x_e)\xi$。因此线性方程确为
\[
\left(2m+\frac M3\right)\ddot\xi+\frac{20k}{3}\xi=0,
\qquad T_{\rm osc}=2\pi\sqrt{\frac{6m+M}{20k}}.
\]
图\ref{mp5:fig:l5-h11-s}的速度比在位移的一阶关系中也表现为 $y_M-y_{M,e}=\xi/\sqrt3$。悬挂物与右环的位移同相，左环的水平位移则与右环相反。

\end{solution}

\par\medskip\noindent\begin{minipage}{\textwidth}
\begin{exercise}[线轴拉动方向的反转（选做）]\label{mp5:exercise:l5-h12}
线轴质量为 $m$，外半径为 $R$，内轴半径为 $r<R$，转动惯量 $I=\eta mR^2$，置于粗糙水平地面。细绳从内轴下侧切向引出，拉力大小为 $F$，与水平向右方向成角 $\beta\in[0,\pi/2)$，见图\ref{mp5:fig:l5-h12-q}。在线轴保持纯滚动时，求加速度、方向反转的临界角及接触条件；用虚功说明为什么较陡的拉力可以使线轴向左运动。
\end{exercise}

\begin{figure}[H]\centering
\begin{tikzpicture}
\draw[curve](-2,0)--(3.3,0);\draw[curve](0,1.2)circle(1.2);\draw[curve](0,1.2)circle(.72);\fill(0,1.2)circle(1.3pt)node[above right]{$C$};
\fill(0,0)circle(1.3pt)node[below left]{$Q$};
\fill(0.6235382907247957,0.8399999999999999)circle(1.3pt)node[below right]{$P$};
\draw(0,1.2)--(0.6235382907247957,0.8399999999999999)node[midway,below left]{$r$};\draw[axis](0.6235382907247957,0.8399999999999999)--(1.8235382907247957,2.9184609690826524)node[above right]{$F$};\draw[guide](0.6235382907247957,0.8399999999999999)--(1.6235382907247957,0.8399999999999999);\draw(1.1235382907247957,0.8399999999999999)arc(0:60:.5);\node at(1.2735382907247956,1.14){$\beta$};\draw[axis](-.5,-.6)--(2.2,-.6)node[right]{$x$};\draw(0,1.2)--(-1.0392,1.8)node[midway,above]{$R$};
\end{tikzpicture}
\caption{绳从内轴下侧切向引出，$P$ 为此刻的切点。绳与外轮缘在轴向错开，图中作平面投影。}\label{mp5:fig:l5-h12-q}
\end{figure}
\end{minipage}\par\medskip
\begin{solution}
取水平向右为正，逆时针角速度为正。纯滚动给出 $\omega_z=-\dot x/R$。图\ref{mp5:fig:l5-h12-q}中绳方向单位矢量为 $\vct t=(\cos\beta,\sin\beta)$，此刻物质点 P 相对 C 的位置为
\[
\vct r_{P/C}=(r\sin\beta,-r\cos\beta).
\]
因此该物质点的瞬时速度为
\[
\vct v_P=\dot x\,\vct e_x+\omega_z\vct e_z\times\vct r_{P/C}
=\dot x\left(\vct e_x-\frac rR\vct t\right).
\]
拉力功率由力作用的物质点速度计算，
\[
\mathcal P_F=F\vct t\cdot\vct v_P
=F\dot x\left(\cos\beta-\frac rR\right).
\]

\begin{figure}[H]\centering
\begin{tikzpicture}
\begin{scope}[shift={(0,0)},scale=1]
\draw[axis](0,0)--(2,0)node[midway,above]{$\delta x\,\vct e_x$};\draw[axis](2,0)--(1.4,-1.03923)node[right]{$-\frac rR\delta x\,\vct t$};\draw[axis](0,0)--(1.4,-1.03923)node[midway,left]{$\delta\vct r_P$};\node at(1,.9){瞬时物质点 P};\end{scope}
\begin{scope}[shift={(5.2,0)},scale=1.2]
\draw[axis](-0.1,0)--(1.8,0)node[right]{$\beta$};
\draw[axis](0,-0.8)--(0,0.8)node[above]{$a/a_*$};
\draw[curve] plot coordinates{(0.00000,0.40000) (0.01571,0.39988) (0.03142,0.39951) (0.04712,0.39889) (0.06283,0.39803) (0.07854,0.39692) (0.09425,0.39556) (0.10996,0.39396) (0.12566,0.39211) (0.14137,0.39002) (0.15708,0.38769) (0.17279,0.38511) (0.18850,0.38229) (0.20420,0.37922) (0.21991,0.37592) (0.23562,0.37237) (0.25133,0.36858) (0.26704,0.36456) (0.28274,0.36029) (0.29845,0.35579) (0.31416,0.35106) (0.32987,0.34609) (0.34558,0.34088) (0.36128,0.33544) (0.37699,0.32978) (0.39270,0.32388) (0.40841,0.31775) (0.42412,0.31140) (0.43982,0.30483) (0.45553,0.29803) (0.47124,0.29101) (0.48695,0.28377) (0.50265,0.27631) (0.51836,0.26863) (0.53407,0.26074) (0.54978,0.25264) (0.56549,0.24433) (0.58119,0.23581) (0.59690,0.22708) (0.61261,0.21815) (0.62832,0.20902) (0.64403,0.19968) (0.65973,0.19016) (0.67544,0.18043) (0.69115,0.17051) (0.70686,0.16041) (0.72257,0.15011) (0.73827,0.13963) (0.75398,0.12897) (0.76969,0.11813) (0.78540,0.10711) (0.80111,0.09591) (0.81681,0.08455) (0.83252,0.07301) (0.84823,0.06131) (0.86394,0.04945) (0.87965,0.03742) (0.89535,0.02524) (0.91106,0.01291) (0.92677,0.00042) (0.94248,-0.01221) (0.95819,-0.02499) (0.97389,-0.03792) (0.98960,-0.05098) (1.00531,-0.06417) (1.02102,-0.07750) (1.03673,-0.09096) (1.05243,-0.10454) (1.06814,-0.11825) (1.08385,-0.13207) (1.09956,-0.14601) (1.11527,-0.16006) (1.13097,-0.17422) (1.14668,-0.18849) (1.16239,-0.20285) (1.17810,-0.21732) (1.19381,-0.23188) (1.20951,-0.24653) (1.22522,-0.26126) (1.24093,-0.27608) (1.25664,-0.29098) (1.27235,-0.30596) (1.28805,-0.32101) (1.30376,-0.33613) (1.31947,-0.35131) (1.33518,-0.36655) (1.35088,-0.38186) (1.36659,-0.39721) (1.38230,-0.41262) (1.39801,-0.42807) (1.41372,-0.44357) (1.42942,-0.45910) (1.44513,-0.47467) (1.46084,-0.49027) (1.47655,-0.50589) (1.49226,-0.52154) (1.50796,-0.53721) (1.52367,-0.55289) (1.53938,-0.56859) (1.55509,-0.58429) (1.57080,-0.60000)};
\fill(0.9272952180016123,0)circle(1.3pt);\draw[guide](0.9272952180016123,.05)--(0.9272952180016123,.25);\node[above]at(0.9272952180016123,.25){$\beta_c$};\node at(.3,.6){向右};\node at(1.3,-.65){向左};\node[below]at(1.57,0){$\pi/2$};\end{scope}

\end{tikzpicture}
\caption{左图为 $\delta x>0$ 的虚位移分解，取 $\beta=60^\circ,r/R=0.6$；右图取 $a_*=F/[m(1+\eta)]$，加速度在 $\beta_c$ 处变号。}\label{mp5:fig:l5-h12-s}
\end{figure}

图\ref{mp5:fig:l5-h12-s}左图画出允许的瞬时虚位移，沿绳的投影为 $\delta\ell=(\cos\beta-r/R)\delta x$。所以广义力 $Q_x=F(\cos\beta-r/R)$，而 $T=m(1+\eta)\dot x^2/2$，得到
\[
\boxed{a=\frac{F(\cos\beta-r/R)}{m(1+\eta)},\qquad
\beta_c=\arccos(r/R).}
\]
从静止出发，$\beta<\beta_c$ 时向右滚，$\beta>\beta_c$ 时向左滚；临界角时可保持静止。若初速不为零，上述结论只说明加速度的方向，不能直接当作速度方向。

令地面摩擦的向右分量为 $f$。平动方程与绕 C 的转动方程是
\[
ma=F\cos\beta+f,\qquad
-\frac{Ia}R=rF+Rf.
\]
联立可得
\[
f=-\frac{\eta\cos\beta+r/R}{1+\eta}F,\qquad
N=mg-F\sin\beta.
\]
所以摩擦始终向左；维持地面接触和纯滚动需要
\[
\boxed{N>0,\qquad
\frac{\eta\cos\beta+r/R}{1+\eta}F\le\mu(mg-F\sin\beta).}
\]
在临界角，$f=-F\cos\beta=-Fr/R$，平动和转动均可平衡。大于临界角并从静止释放时，$\dot x<0$ 且 $\cos\beta-r/R<0$，拉力功率仍为正。

几何切点随构形移动，与它重合的物质点也不断更换。计算瞬时功率时使用上面的物质点速度，不能把几何切点的轨迹速度代入。
用地面接触点 Q 为矩心还可核对方向反转条件。拉力的力矩为
\[
\tau_{Q,z}=F(r-R\cos\beta).
\]
对该瞬间位于 Q 的固定地面点写角动量定理，有 $\tau_{Q,z}=I\dot\omega_z-mRa$。代入 $a=-R\dot\omega_z$，得到 $(I+mR^2)\dot\omega_z=F(r-R\cos\beta)$，恢复前述 a。$\beta=\beta_c$ 时拉力作用线经过 Q，正是这个力矩变号的位置。

若地面静摩擦因数为 $\mu$，还可把接触不等式解为拉力上限：
\[
F\le\frac{(1+\eta)\mu mg}
{\eta\cos\beta+r/R+(1+\eta)\mu\sin\beta}.
\]
对于 $\mu>0$，上式已使 $N>0$；当 $\sin\beta>0$ 时，其上限严格小于 $mg/\sin\beta$。因此力过大时首先需要重新判断摩擦能否维持纯滚动，而不能仅检查有没有被提离地面。图\ref{mp5:fig:l5-h12-s}中的加速度方向图，只有在这些接触条件满足时才对应实际运动。

\end{solution}
\end{document}
